% 08_feature_importance.tex -- the rehauled paper's section, as bullets to expand into prose (2026-10-15).
% Sources and pointers behind every number: writeup/PAPER_REHAUL_OUTLINE_2026-10-15.md, section 8.
% House rules: every number in the final prose comes from a generated macro (writeup/generated/*.tex), never
% typed; park (comment out with a dated note) rather than delete.
% Macros to use here: rehaul_fi_numbers.tex (\rhfi..., writeup/make_rehaul_fi.py from results/rehaul_fi/, computed by
% experiments/rehaul_fi_compute.py); floats in 08b_fi_measures_floats.tex. The old study's measures (permutation,
% MDI, drop-column; appendix_campaign_16h_featimp.tex) stay in 08_feature_importance_floats.tex; their bullets are
% parked at the end of this file.
% 2026-10-15 measures pass: the user's plan (lasso counts, ridge sum |beta|, per-refit SHAP ranks aggregated,
% SHAP interactions global and per refit with the H-statistic cross-check) on the paper's models of record.
% 2026-10-15 (notation pass, writeup/NOTATION_GUIDE_2026-10-15.md): symbols and wording only; no number, macro or
% label changed. Old -> new: unit -> predictor group; column -> predictor; backbone -> HAR benchmark block;
% pf_j -> omega_j; alpha -> lambda; session -> trading day; fit space -> the scale of y_t; W_k -> \mathcal W_{t_k};
% T_k -> out-of-sample block \mathcal O_k; N (scored sessions) -> |\mathcal O|, \mathcal O the union of the blocks
% (not the guide's T_oos: this 2018-2024 period is not the paper's 2020-2024 test period); x_t -> \bm x_t;
% \beta^{(k)}_j -> \hat\beta_{k,j} (a parenthesized superscript is the guide's HAR horizon);
% b_{k,A} -> \lVert\hat{\bm\beta}_{k,A}\rVert_1 (b is the boosting round); s_{k,A} -> S_{k,A} (s is the periodicity
% factor); B_A (Borda) -> \mathrm{BC}_A; tie sizes \tau_{kg} -> |g|, g in \mathcal G_k (tau is the network step);
% H^2 points x_i and background b -> days t, t' in \mathcal H_k (i is the intraday interval).
% The generated floats (make_rehaul_fi.py) still print the old symbols ($T_k$, $\beta^{(k)}_j$, "unit", code names).
% 2026-10-15 (consolidation): the floats now print this section's symbols. Selection count N_A, N_j -> K_A, K_j (N is
% the one-minute count N_{t,i}, the symbol of the numobs series in the tables); group interaction Phi_{t,AB} ->
% phi_{t,AB} (Phi is the standard normal distribution function); tie sets \mathcal G_k -> \mathcal E_k (\mathcal G is
% the interval grid, Data); the ridge's standardized share S^{sd}_A (as in tab:rh_fi_summary, tab:rh_fi_linear);
% "buckets" -> categories. The refit index stays k, subscripted (guide 6b); HAR horizons are the superscript (k).
% 2026-10-15 (consolidation): parked -- the number files are loaded once in the preamble of paper_rehaul.tex.
% \ifdefined\rhfiNRefits\else% AUTO-GENERATED by writeup/make_rehaul_fi.py -- do not edit; re-run the script.
% \rhfi... macros for writeup/rehaul/08_feature_importance.tex (results/rehaul_fi/).
\providecommand{\rhfiNRefits}{}\renewcommand{\rhfiNRefits}{147}  % refits, one every 10 sessions
\providecommand{\rhfiNOos}{}\renewcommand{\rhfiNOos}{1{,}469}  % scored held-out sessions
\providecommand{\rhfiP}{}\renewcommand{\rhfiP}{628}  % columns, all features
\providecommand{\rhfiNUnits}{}\renewcommand{\rhfiNUnits}{49}  % units: 47 exogenous series + HAR ladder + calendar
\providecommand{\rhfiNExog}{}\renewcommand{\rhfiNExog}{47}  % exogenous series
\providecommand{\rhfiFirst}{}\renewcommand{\rhfiFirst}{25~June~2018}  % first scored session
\providecommand{\rhfiLast}{}\renewcommand{\rhfiLast}{30~April~2024}  % last scored session
\providecommand{\rhfiWindow}{}\renewcommand{\rhfiWindow}{2{,}000}  % sessions in a training window
\providecommand{\rhfiRefitEvery}{}\renewcommand{\rhfiRefitEvery}{10}  % sessions between refits
\providecommand{\rhfiBackbonePf}{}\renewcommand{\rhfiBackbonePf}{0.01}  % backbone penalty factor
\providecommand{\rhfiNBuckets}{}\renewcommand{\rhfiNBuckets}{8}  % exogenous buckets
\providecommand{\rhfiOldBlock}{}\renewcommand{\rhfiOldBlock}{250}  % sessions between the old arm's penalty re-choices
\providecommand{\rhfiNPairs}{}\renewcommand{\rhfiNPairs}{1{,}176}  % series pairs
\providecommand{\rhfiNProbeCols}{}\renewcommand{\rhfiNProbeCols}{631}  % columns of a probe fit
\providecommand{\rhfiTopWPairs}{}\renewcommand{\rhfiTopWPairs}{50}  % pairs in the interaction W
\providecommand{\rhfiProbeLevels}{}\renewcommand{\rhfiProbeLevels}{20}  % values of the few-valued noise column
\providecommand{\rhfiRidgeAlpha}{}\renewcommand{\rhfiRidgeAlpha}{\ensuremath{10^{3}}}
\providecommand{\rhfiLassoAlpha}{}\renewcommand{\rhfiLassoAlpha}{\ensuremath{10^{-2}}}
\providecommand{\rhfiLassoNzMed}{}\renewcommand{\rhfiLassoNzMed}{30}  % median nonzero coefficients per refit
\providecommand{\rhfiLassoNzMin}{}\renewcommand{\rhfiLassoNzMin}{27}
\providecommand{\rhfiLassoNzMax}{}\renewcommand{\rhfiLassoNzMax}{35}
\providecommand{\rhfiLassoExogNzMed}{}\renewcommand{\rhfiLassoExogNzMed}{14}  % ... exogenous columns
\providecommand{\rhfiLassoColsEver}{}\renewcommand{\rhfiLassoColsEver}{60}  % columns nonzero in >= 1 refit
\providecommand{\rhfiLassoColsMajority}{}\renewcommand{\rhfiLassoColsMajority}{29}  % columns nonzero in >= half the refits
\providecommand{\rhfiLassoExogEver}{}\renewcommand{\rhfiLassoExogEver}{27}  % exogenous series with a nonzero column in >= 1 refit
\providecommand{\rhfiLassoExogMajority}{}\renewcommand{\rhfiLassoExogMajority}{11}  % ... in >= half the refits
\providecommand{\rhfiLassoExogAlways}{}\renewcommand{\rhfiLassoExogAlways}{3}  % ... in every refit
\providecommand{\rhfiLassoAlwaysList}{}\renewcommand{\rhfiLassoAlwaysList}{\texttt{sumret}, \texttt{sumabsret\_vwstock}, \texttt{sumautocov}}  % exogenous series selected in every refit
\providecommand{\rhfiLassoProbeSel}{}\renewcommand{\rhfiLassoProbeSel}{1}  % continuous probe selected, of the probe fits
\providecommand{\rhfiLassoProbeSelLevels}{}\renewcommand{\rhfiLassoProbeSelLevels}{0}
\providecommand{\rhfiLassoProbeSelBinary}{}\renewcommand{\rhfiLassoProbeSelBinary}{0}
\providecommand{\rhfiNProbe}{}\renewcommand{\rhfiNProbe}{49}  % probe refits (every third)
\providecommand{\rhfiLassoW}{}\renewcommand{\rhfiLassoW}{0.74}  % Kendall W of the per-refit series ranks by sum |beta|
\providecommand{\rhfiLassoBackboneShare}{}\renewcommand{\rhfiLassoBackboneShare}{88}  % median share of sum |beta| on the backbone, %
\providecommand{\rhfiOldLassoSpearman}{}\renewcommand{\rhfiOldLassoSpearman}{0.70}  % Spearman, selection share: pre-tuned vs re-tuned one-penalty lasso
\providecommand{\rhfiOldLassoNzMed}{}\renewcommand{\rhfiOldLassoNzMed}{23}  % re-tuned lasso: median nonzero per refit
\providecommand{\rhfiOldLassoExogMajority}{}\renewcommand{\rhfiOldLassoExogMajority}{18}
\providecommand{\rhfiRidgeTopOne}{}\renewcommand{\rhfiRidgeTopOne}{Calendar}
\providecommand{\rhfiRidgeTopOneShare}{}\renewcommand{\rhfiRidgeTopOneShare}{31.0}  % % of sum |beta|
\providecommand{\rhfiRidgeTopTwo}{}\renewcommand{\rhfiRidgeTopTwo}{HAR ladder}
\providecommand{\rhfiRidgeTopTwoShare}{}\renewcommand{\rhfiRidgeTopTwoShare}{22.8}  % % of sum |beta|
\providecommand{\rhfiRidgeTopThree}{}\renewcommand{\rhfiRidgeTopThree}{\texttt{sumret}}
\providecommand{\rhfiRidgeTopThreeShare}{}\renewcommand{\rhfiRidgeTopThreeShare}{3.6}  % % of sum |beta|
\providecommand{\rhfiRidgeTopFour}{}\renewcommand{\rhfiRidgeTopFour}{\texttt{sumabsret\_vwstock}}
\providecommand{\rhfiRidgeTopFourShare}{}\renewcommand{\rhfiRidgeTopFourShare}{2.2}  % % of sum |beta|
\providecommand{\rhfiRidgeSdTopOne}{}\renewcommand{\rhfiRidgeSdTopOne}{HAR ladder}
\providecommand{\rhfiRidgeSdTopOneShare}{}\renewcommand{\rhfiRidgeSdTopOneShare}{30.6}  % % of sum |beta| sd
\providecommand{\rhfiRidgeSdTopTwo}{}\renewcommand{\rhfiRidgeSdTopTwo}{Calendar}
\providecommand{\rhfiRidgeSdTopTwoShare}{}\renewcommand{\rhfiRidgeSdTopTwoShare}{9.3}  % % of sum |beta| sd
\providecommand{\rhfiRidgeSdTopThree}{}\renewcommand{\rhfiRidgeSdTopThree}{\texttt{sumret}}
\providecommand{\rhfiRidgeSdTopThreeShare}{}\renewcommand{\rhfiRidgeSdTopThreeShare}{4.5}  % % of sum |beta| sd
\providecommand{\rhfiRidgeTopExog}{}\renewcommand{\rhfiRidgeTopExog}{\texttt{sumret}}
\providecommand{\rhfiRidgeTopExogShare}{}\renewcommand{\rhfiRidgeTopExogShare}{3.6}
\providecommand{\rhfiRidgeBackboneShare}{}\renewcommand{\rhfiRidgeBackboneShare}{54}  % median share of sum |beta| on the backbone, %
\providecommand{\rhfiRidgeW}{}\renewcommand{\rhfiRidgeW}{0.85}
\providecommand{\rhfiRidgeAllSelected}{}\renewcommand{\rhfiRidgeAllSelected}{49}  % units with N_A = K under the ridge
\providecommand{\rhfiLassoProbeRate}{}\renewcommand{\rhfiLassoProbeRate}{2.0}  % continuous probe's selection rate, %
\providecommand{\rhfiLassoColsAboveProbe}{}\renewcommand{\rhfiLassoColsAboveProbe}{54}  % columns selected more often than the continuous probe
\providecommand{\rhfiLassoSeriesAboveProbe}{}\renewcommand{\rhfiLassoSeriesAboveProbe}{23}  % exogenous series with such a column
\providecommand{\rhfiLassoTopCol}{}\renewcommand{\rhfiLassoTopCol}{\texttt{is\_month\_end}}
\providecommand{\rhfiRidgeSpearmanSd}{}\renewcommand{\rhfiRidgeSpearmanSd}{0.85}  % Spearman, series sum|beta| vs sum|beta| sd
\providecommand{\rhfiRidgeProbeBeat}{}\renewcommand{\rhfiRidgeProbeBeat}{107}  % columns with |beta| above the continuous probe's in most probe fits
\providecommand{\rhfiRidgeProbeBeatSeries}{}\renewcommand{\rhfiRidgeProbeBeatSeries}{38}
\providecommand{\rhfiScaleValueN}{}\renewcommand{\rhfiScaleValueN}{282}
\providecommand{\rhfiScaleValueSd}{}\renewcommand{\rhfiScaleValueSd}{0.84}  % median window sd, value ladders
\providecommand{\rhfiScaleValueSdLo}{}\renewcommand{\rhfiScaleValueSdLo}{0.03}
\providecommand{\rhfiScaleValueSdHi}{}\renewcommand{\rhfiScaleValueSdHi}{1.76}
\providecommand{\rhfiScaleValueIqr}{}\renewcommand{\rhfiScaleValueIqr}{0.94}  % median window IQR, value ladders
\providecommand{\rhfiScaleFlagN}{}\renewcommand{\rhfiScaleFlagN}{324}
\providecommand{\rhfiScaleFlagSd}{}\renewcommand{\rhfiScaleFlagSd}{0.022}  % median window sd, availability / activity ladders
\providecommand{\rhfiScaleFlagSdHi}{}\renewcommand{\rhfiScaleFlagSdHi}{0.17}
\providecommand{\rhfiScaleCalSd}{}\renewcommand{\rhfiScaleCalSd}{0.39}  % median window sd, calendar
\providecommand{\rhfiScaleHarSd}{}\renewcommand{\rhfiScaleHarSd}{1.07}  % median window sd, HAR ladder
\providecommand{\rhfiLgbmW}{}\renewcommand{\rhfiLgbmW}{0.57}  % lgbm: Kendall W, per-refit series ranks
\providecommand{\rhfiLgbmWCols}{}\renewcommand{\rhfiLgbmWCols}{0.91}  % lgbm: Kendall W, column ranks
\providecommand{\rhfiLgbmTopOne}{}\renewcommand{\rhfiLgbmTopOne}{HAR ladder}
\providecommand{\rhfiLgbmTopOneRank}{}\renewcommand{\rhfiLgbmTopOneRank}{1.0}  % mean rank
\providecommand{\rhfiLgbmTopTwo}{}\renewcommand{\rhfiLgbmTopTwo}{\texttt{sumret}}
\providecommand{\rhfiLgbmTopTwoRank}{}\renewcommand{\rhfiLgbmTopTwoRank}{2.9}  % mean rank
\providecommand{\rhfiLgbmTopThree}{}\renewcommand{\rhfiLgbmTopThree}{\texttt{sumabsret}}
\providecommand{\rhfiLgbmTopThreeRank}{}\renewcommand{\rhfiLgbmTopThreeRank}{4.8}  % mean rank
\providecommand{\rhfiLgbmSpBorda}{}\renewcommand{\rhfiLgbmSpBorda}{0.90}  % Spearman, mean rank vs top-10 Borda
\providecommand{\rhfiLgbmSpGlobal}{}\renewcommand{\rhfiLgbmSpGlobal}{0.96}  % Spearman, mean rank vs global mean |phi|
\providecommand{\rhfiLgbmSpLocal}{}\renewcommand{\rhfiLgbmSpLocal}{1.000}  % Spearman, mean ranks local refit vs stored
\providecommand{\rhfiLgbmProbeRank}{}\renewcommand{\rhfiLgbmProbeRank}{83}  % median rank of the continuous probe, 631 columns
\providecommand{\rhfiLgbmProbeCols}{}\renewcommand{\rhfiLgbmProbeCols}{60}  % columns above the continuous probe in most probe fits
\providecommand{\rhfiLgbmProbeSeries}{}\renewcommand{\rhfiLgbmProbeSeries}{27}  % exogenous series with such a column
\providecommand{\rhfiLgbmTopCol}{}\renewcommand{\rhfiLgbmTopCol}{\texttt{har\_ma\_1}}  % column with the lowest mean rank
\providecommand{\rhfiLgbmTopColRank}{}\renewcommand{\rhfiLgbmTopColRank}{1.3}
\providecommand{\rhfiLgbmTopColTwo}{}\renewcommand{\rhfiLgbmTopColTwo}{\texttt{har\_ma\_5}}
\providecommand{\rhfiLgbmTopColTwoRank}{}\renewcommand{\rhfiLgbmTopColTwoRank}{1.8}
\providecommand{\rhfiLgbmPredGap}{}\renewcommand{\rhfiLgbmPredGap}{0}  % max |local refit - stored refit| forecast
\providecommand{\rhfiLgbmInterW}{}\renewcommand{\rhfiLgbmInterW}{0.45}  % Kendall W of per-refit pair ranks, top-50 pairs
\providecommand{\rhfiLgbmInterRatio}{}\renewcommand{\rhfiLgbmInterRatio}{90}  % sum of series-pair interactions / sum main, %
\providecommand{\rhfiLgbmInterSpGlobal}{}\renewcommand{\rhfiLgbmInterSpGlobal}{0.99}
\providecommand{\rhfiLgbmInterFloor}{}\renewcommand{\rhfiLgbmInterFloor}{0.49}  % probe floor x 1e3
\providecommand{\rhfiLgbmInterAboveFloor}{}\renewcommand{\rhfiLgbmInterAboveFloor}{116}  % series pairs above the probe floor
\providecommand{\rhfiLgbmInterTopA}{}\renewcommand{\rhfiLgbmInterTopA}{HAR ladder}
\providecommand{\rhfiLgbmInterTopB}{}\renewcommand{\rhfiLgbmInterTopB}{\texttt{sumabsret}}
\providecommand{\rhfiLgbmInterTopVal}{}\renewcommand{\rhfiLgbmInterTopVal}{12.6}  % x 1e3
\providecommand{\rhfiLgbmHTop}{}\renewcommand{\rhfiLgbmHTop}{0.33}  % H^2 of the top series pair (pooled), %
\providecommand{\rhfiLgbmHMin}{}\renewcommand{\rhfiLgbmHMin}{0.02}  % min H^2 over the top-10 series pairs, %
\providecommand{\rhfiLgbmHMax}{}\renewcommand{\rhfiLgbmHMax}{0.33}  % max H^2 over the top-10 series pairs, %
\providecommand{\rhfiLgbmHMinTop}{}\renewcommand{\rhfiLgbmHMinTop}{0.19}  % top pair: interaction PD var / smaller main PD var
\providecommand{\rhfiLgbmHMinLo}{}\renewcommand{\rhfiLgbmHMinLo}{0.05}  % ... min over the top-10 series pairs
\providecommand{\rhfiLgbmHMinHi}{}\renewcommand{\rhfiLgbmHMinHi}{0.54}  % ... max
\providecommand{\rhfiLgbmHColTop}{}\renewcommand{\rhfiLgbmHColTop}{0.36}  % H^2 of the top column pair, %
\providecommand{\rhfiLgbmHMinColTop}{}\renewcommand{\rhfiLgbmHMinColTop}{0.02}  % ... min-normalised
\providecommand{\rhfiLgbmMainTopVal}{}\renewcommand{\rhfiLgbmMainTopVal}{195}  % largest global mean |phi_A| x 1e3
\providecommand{\rhfiLgbmWithinTop}{}\renewcommand{\rhfiLgbmWithinTop}{HAR ladder}
\providecommand{\rhfiLgbmWithinTopVal}{}\renewcommand{\rhfiLgbmWithinTopVal}{18.7}
\providecommand{\rhfiXgbW}{}\renewcommand{\rhfiXgbW}{0.60}  % xgb: Kendall W, per-refit series ranks
\providecommand{\rhfiXgbWCols}{}\renewcommand{\rhfiXgbWCols}{0.90}  % xgb: Kendall W, column ranks
\providecommand{\rhfiXgbTopOne}{}\renewcommand{\rhfiXgbTopOne}{HAR ladder}
\providecommand{\rhfiXgbTopOneRank}{}\renewcommand{\rhfiXgbTopOneRank}{1.0}  % mean rank
\providecommand{\rhfiXgbTopTwo}{}\renewcommand{\rhfiXgbTopTwo}{\texttt{sumret}}
\providecommand{\rhfiXgbTopTwoRank}{}\renewcommand{\rhfiXgbTopTwoRank}{3.0}  % mean rank
\providecommand{\rhfiXgbTopThree}{}\renewcommand{\rhfiXgbTopThree}{\texttt{sumabsret}}
\providecommand{\rhfiXgbTopThreeRank}{}\renewcommand{\rhfiXgbTopThreeRank}{3.9}  % mean rank
\providecommand{\rhfiXgbSpBorda}{}\renewcommand{\rhfiXgbSpBorda}{0.87}  % Spearman, mean rank vs top-10 Borda
\providecommand{\rhfiXgbSpGlobal}{}\renewcommand{\rhfiXgbSpGlobal}{0.97}  % Spearman, mean rank vs global mean |phi|
\providecommand{\rhfiXgbSpLocal}{}\renewcommand{\rhfiXgbSpLocal}{0.993}  % Spearman, mean ranks local refit vs stored
\providecommand{\rhfiXgbProbeRank}{}\renewcommand{\rhfiXgbProbeRank}{110}  % median rank of the continuous probe, 631 columns
\providecommand{\rhfiXgbProbeCols}{}\renewcommand{\rhfiXgbProbeCols}{78}  % columns above the continuous probe in most probe fits
\providecommand{\rhfiXgbProbeSeries}{}\renewcommand{\rhfiXgbProbeSeries}{35}  % exogenous series with such a column
\providecommand{\rhfiXgbTopCol}{}\renewcommand{\rhfiXgbTopCol}{\texttt{har\_ma\_1}}  % column with the lowest mean rank
\providecommand{\rhfiXgbTopColRank}{}\renewcommand{\rhfiXgbTopColRank}{1.3}
\providecommand{\rhfiXgbTopColTwo}{}\renewcommand{\rhfiXgbTopColTwo}{\texttt{har\_ma\_5}}
\providecommand{\rhfiXgbTopColTwoRank}{}\renewcommand{\rhfiXgbTopColTwoRank}{1.8}
\providecommand{\rhfiXgbPredGap}{}\renewcommand{\rhfiXgbPredGap}{0.17}  % max |local refit - stored refit| forecast
\providecommand{\rhfiXgbInterW}{}\renewcommand{\rhfiXgbInterW}{0.47}  % Kendall W of per-refit pair ranks, top-50 pairs
\providecommand{\rhfiXgbInterRatio}{}\renewcommand{\rhfiXgbInterRatio}{74}  % sum of series-pair interactions / sum main, %
\providecommand{\rhfiXgbInterSpGlobal}{}\renewcommand{\rhfiXgbInterSpGlobal}{0.99}
\providecommand{\rhfiXgbInterFloor}{}\renewcommand{\rhfiXgbInterFloor}{0.41}  % probe floor x 1e3
\providecommand{\rhfiXgbInterAboveFloor}{}\renewcommand{\rhfiXgbInterAboveFloor}{115}  % series pairs above the probe floor
\providecommand{\rhfiXgbInterTopA}{}\renewcommand{\rhfiXgbInterTopA}{HAR ladder}
\providecommand{\rhfiXgbInterTopB}{}\renewcommand{\rhfiXgbInterTopB}{\texttt{sumabsret}}
\providecommand{\rhfiXgbInterTopVal}{}\renewcommand{\rhfiXgbInterTopVal}{15.2}  % x 1e3
\providecommand{\rhfiXgbHTop}{}\renewcommand{\rhfiXgbHTop}{0.37}  % H^2 of the top series pair (pooled), %
\providecommand{\rhfiXgbHMin}{}\renewcommand{\rhfiXgbHMin}{0.02}  % min H^2 over the top-10 series pairs, %
\providecommand{\rhfiXgbHMax}{}\renewcommand{\rhfiXgbHMax}{0.68}  % max H^2 over the top-10 series pairs, %
\providecommand{\rhfiXgbHMinTop}{}\renewcommand{\rhfiXgbHMinTop}{0.39}  % top pair: interaction PD var / smaller main PD var
\providecommand{\rhfiXgbHMinLo}{}\renewcommand{\rhfiXgbHMinLo}{0.08}  % ... min over the top-10 series pairs
\providecommand{\rhfiXgbHMinHi}{}\renewcommand{\rhfiXgbHMinHi}{0.80}  % ... max
\providecommand{\rhfiXgbHColTop}{}\renewcommand{\rhfiXgbHColTop}{0.39}  % H^2 of the top column pair, %
\providecommand{\rhfiXgbHMinColTop}{}\renewcommand{\rhfiXgbHMinColTop}{0.02}  % ... min-normalised
\providecommand{\rhfiXgbMainTopVal}{}\renewcommand{\rhfiXgbMainTopVal}{192}  % largest global mean |phi_A| x 1e3
\providecommand{\rhfiXgbWithinTop}{}\renewcommand{\rhfiXgbWithinTop}{HAR ladder}
\providecommand{\rhfiXgbWithinTopVal}{}\renewcommand{\rhfiXgbWithinTopVal}{19.1}
\providecommand{\rhfiRfW}{}\renewcommand{\rhfiRfW}{0.62}  % rf: Kendall W, per-refit series ranks
\providecommand{\rhfiRfWCols}{}\renewcommand{\rhfiRfWCols}{0.93}  % rf: Kendall W, column ranks
\providecommand{\rhfiRfTopOne}{}\renewcommand{\rhfiRfTopOne}{HAR ladder}
\providecommand{\rhfiRfTopOneRank}{}\renewcommand{\rhfiRfTopOneRank}{1.0}  % mean rank
\providecommand{\rhfiRfTopTwo}{}\renewcommand{\rhfiRfTopTwo}{\texttt{sumabsret}}
\providecommand{\rhfiRfTopTwoRank}{}\renewcommand{\rhfiRfTopTwoRank}{2.6}  % mean rank
\providecommand{\rhfiRfTopThree}{}\renewcommand{\rhfiRfTopThree}{\texttt{sumret}}
\providecommand{\rhfiRfTopThreeRank}{}\renewcommand{\rhfiRfTopThreeRank}{4.0}  % mean rank
\providecommand{\rhfiRfSpBorda}{}\renewcommand{\rhfiRfSpBorda}{0.92}  % Spearman, mean rank vs top-10 Borda
\providecommand{\rhfiRfSpGlobal}{}\renewcommand{\rhfiRfSpGlobal}{0.98}  % Spearman, mean rank vs global mean |phi|
\providecommand{\rhfiRfSpLocal}{}\renewcommand{\rhfiRfSpLocal}{0.996}  % Spearman, mean ranks local refit vs stored
\providecommand{\rhfiRfProbeRank}{}\renewcommand{\rhfiRfProbeRank}{97}  % median rank of the continuous probe, 631 columns
\providecommand{\rhfiRfProbeCols}{}\renewcommand{\rhfiRfProbeCols}{71}  % columns above the continuous probe in most probe fits
\providecommand{\rhfiRfProbeSeries}{}\renewcommand{\rhfiRfProbeSeries}{29}  % exogenous series with such a column
\providecommand{\rhfiRfTopCol}{}\renewcommand{\rhfiRfTopCol}{\texttt{har\_ma\_1}}  % column with the lowest mean rank
\providecommand{\rhfiRfTopColRank}{}\renewcommand{\rhfiRfTopColRank}{1.3}
\providecommand{\rhfiRfTopColTwo}{}\renewcommand{\rhfiRfTopColTwo}{\texttt{har\_ma\_5}}
\providecommand{\rhfiRfTopColTwoRank}{}\renewcommand{\rhfiRfTopColTwoRank}{1.7}
\providecommand{\rhfiRfPredGap}{}\renewcommand{\rhfiRfPredGap}{0.11}  % max |local refit - stored refit| forecast
\providecommand{\rhfiRfInterW}{}\renewcommand{\rhfiRfInterW}{0.57}  % Kendall W of per-refit pair ranks, top-50 pairs
\providecommand{\rhfiRfInterRatio}{}\renewcommand{\rhfiRfInterRatio}{72}  % sum of series-pair interactions / sum main, %
\providecommand{\rhfiRfInterSpGlobal}{}\renewcommand{\rhfiRfInterSpGlobal}{1.00}
\providecommand{\rhfiRfInterFloor}{}\renewcommand{\rhfiRfInterFloor}{0.50}  % probe floor x 1e3
\providecommand{\rhfiRfInterAboveFloor}{}\renewcommand{\rhfiRfInterAboveFloor}{52}  % series pairs above the probe floor
\providecommand{\rhfiRfInterTopA}{}\renewcommand{\rhfiRfInterTopA}{HAR ladder}
\providecommand{\rhfiRfInterTopB}{}\renewcommand{\rhfiRfInterTopB}{\texttt{sumabsret}}
\providecommand{\rhfiRfInterTopVal}{}\renewcommand{\rhfiRfInterTopVal}{23.8}  % x 1e3
\providecommand{\rhfiRfHTop}{}\renewcommand{\rhfiRfHTop}{0.85}  % H^2 of the top series pair (pooled), %
\providecommand{\rhfiRfHMin}{}\renewcommand{\rhfiRfHMin}{0.02}  % min H^2 over the top-10 series pairs, %
\providecommand{\rhfiRfHMax}{}\renewcommand{\rhfiRfHMax}{0.85}  % max H^2 over the top-10 series pairs, %
\providecommand{\rhfiRfHMinTop}{}\renewcommand{\rhfiRfHMinTop}{1.50}  % top pair: interaction PD var / smaller main PD var
\providecommand{\rhfiRfHMinLo}{}\renewcommand{\rhfiRfHMinLo}{0.17}  % ... min over the top-10 series pairs
\providecommand{\rhfiRfHMinHi}{}\renewcommand{\rhfiRfHMinHi}{5.46}  % ... max
\providecommand{\rhfiRfHColTop}{}\renewcommand{\rhfiRfHColTop}{1.16}  % H^2 of the top column pair, %
\providecommand{\rhfiRfHMinColTop}{}\renewcommand{\rhfiRfHMinColTop}{0.06}  % ... min-normalised
\providecommand{\rhfiRfMainTopVal}{}\renewcommand{\rhfiRfMainTopVal}{224}  % largest global mean |phi_A| x 1e3
\providecommand{\rhfiRfWithinTop}{}\renewcommand{\rhfiRfWithinTop}{HAR ladder}
\providecommand{\rhfiRfWithinTopVal}{}\renewcommand{\rhfiRfWithinTopVal}{34.3}
\providecommand{\rhfiSpTreesMin}{}\renewcommand{\rhfiSpTreesMin}{0.83}  % Spearman of mean ranks, trees
\providecommand{\rhfiSpTreesMax}{}\renewcommand{\rhfiSpTreesMax}{0.96}
\providecommand{\rhfiSpTreeLinMin}{}\renewcommand{\rhfiSpTreeLinMin}{0.17}  % Spearman of mean ranks, tree vs linear
\providecommand{\rhfiSpTreeLinMax}{}\renewcommand{\rhfiSpTreeLinMax}{0.55}
\providecommand{\rhfiSpLassoRidge}{}\renewcommand{\rhfiSpLassoRidge}{0.47}
\providecommand{\rhfiHRefits}{}\renewcommand{\rhfiHRefits}{15}  % refits for H
\providecommand{\rhfiHRows}{}\renewcommand{\rhfiHRows}{100}  % evaluation points per refit
\providecommand{\rhfiHTopSeries}{}\renewcommand{\rhfiHTopSeries}{10}
\providecommand{\rhfiFitSecRf}{}\renewcommand{\rhfiFitSecRf}{212}  % mean seconds per forest refit (11 processes sharing the CPU)
\providecommand{\rhfiShapSecRf}{}\renewcommand{\rhfiShapSecRf}{182}  % mean seconds per forest refit, SHAP + interactions
\providecommand{\rhfiShapSecLgbm}{}\renewcommand{\rhfiShapSecLgbm}{0.9}  % ... LightGBM
\providecommand{\rhfiShapSecXgb}{}\renewcommand{\rhfiShapSecXgb}{0.4}  % ... XGBoost
\providecommand{\rhfiWallTreesMin}{}\renewcommand{\rhfiWallTreesMin}{133}  % wall minutes, tree stage on 11 processes
\providecommand{\rhfiWallHMin}{}\renewcommand{\rhfiWallHMin}{7}  % wall minutes, H stage
\providecommand{\rhfiWallLinMin}{}\renewcommand{\rhfiWallLinMin}{7}  % wall minutes, linear stage
\providecommand{\rhfiWorkers}{}\renewcommand{\rhfiWorkers}{11}
\fi
% 2026-10-15 (finish): "refit" -> "re-estimation" throughout the prose (guide terms); the index k and every \rhfi
% macro name (\rhfiNRefits, \rhfiRefitEvery, \rhfiHRefits) unchanged; the reproduction bullet of "Known issues" says
% "reproduce" instead of "re-estimate" to stay unambiguous; "periodicity-adjusted" -> "diurnally adjusted" (user).
% The generated tables' notes (make_rehaul_fi.py) still say "refits".
% 2026-10-23 (user, RR / E19 / new-notes 10b and 22): the live text below is on the FINAL pooled runs (W = 3,000 trading
% days, first evaluation day 2 January 2014). Macros \rhfip...: generated/rehaul_fip_numbers.tex, written by
% writeup/make_rehaul_fi_pooled.py from results/tree_shap_w3000_2014start_2026-10-23 (TreeSHAP, RULE_shap.txt) and
% results/fi_by_window_2014start_2026-10-22 (linear coefficients, RULE_fiw.txt); it also re-checks the origin refits
% against the merged every-bar walks of the three trees. Floats: tab:rh_fip_shares, tab:rh_fip_nonstd, fig:rh_fip_rank,
% fig:rh_fip_time, fig:rh_fip_cond (08_feature_importance_floats.tex); tab:rh_fip_window, tab:rh_fip_inter,
% fig:rh_fip_shapes (08b_fi_measures_floats.tex). The number file is loaded below (guarded) because the preamble of
% paper_rehaul.tex is not edited in this pass; \providecommand + \renewcommand make a second load harmless. The
% closing-interval body stays parked in \iffalse ... \fi (those models are not rerun; user).
% 2026-10-26 (integration and pruning pass, JFEc page budget; writeup/INTEGRATION_SPEC_2026-10-26.md, VI agent):
% Variable importance is promoted from the last subsection of the results to its own section (\section, label
% sec:rh_fi kept), about three pages: what the models use, the feature-count ladder with the ridge--lasso comparison,
% the late-starting series while live, use against accuracy. The rest moved to the online appendix, rehaul/A_fi.tex
% (labels app:rh_fi_compute, app:rh_fi_rank, app:rh_fi_cond, app:rh_fi_shapes, app:rh_fi_window,
% app:rh_fi_scaling). Audit fixes here: results-7 (elastic-net evidence, in A_fi), results-8 (holiday sessions in the
% explained bars), results-9 (FOMC statement days; generator change flagged to NUMBERS, note in A_fi), results-20
% (tuning samples). Replaced text is parked after the section's itemize below.

% 2026-10-15 (notation pass): parked -- old title; "variable importance" is the language of Gu, Kelly and Xiu (2020).
% \section{Feature importance}
% 2026-10-15 (user): parked -- "the feature importance needs to move to the end of the experiments section ... as a
% subsection"; included after 06_results in paper_rehaul.tex, so it is the last subsection of Experiments and results.
% \section{Variable importance}
% 2026-10-17 (user): parked -- closing-interval material: this subsection reports variable importance for the
% closing-interval models only, so its body is parked in \iffalse ... \fi below. Note E19 (the pooled intraday
% proposal, not about the parked text only) is kept visible here, moved verbatim out of the Notation bullet.
% 2026-10-26 (integration, JFEc): parked -- the subsection heading; Variable importance is now its own section,
% after the results and before the economic application (FRONT wires the order in paper_rehaul.tex). Was:
% "\subsection{Variable importance}"
% 2026-10-28 (author: Sections 6 and 7 are one section): the former Section 7 is Section 6.3, a subsection of Results (its
% subsections are \subsubsection's, labels unchanged); paper_rehaul.tex inputs Robustness (rehaul/06b_robustness.tex,
% Section 6.4) after this file. Was: \section{Which predictors carry the gain?}
% 2026-10-28 (author: conventional titles). Was: \subsection{Which predictors carry the gain?}
\subsection{Predictor-group ablations}
\label{sec:rh_fi}
\label{sec:rh_results_ablation}
% 2026-10-28 (author: v2 of Sections 6-7 approved preliminarily): the section is v2 (drafts/08_predictors_DRAFT.tex
% and drafts/MAP_06_07.md, 26 Oct), opened by the consolidated statement of 27 Oct (formerly Section 6.2, whose text
% is reused in the opening, Sections 7.2-7.4). Changes from the draft: Section 7.3 follows the follow-up tests of
% 27 Oct (a few series carry the trees' remainder, not a factor; the draft's "spread across groups, used nonlinearly"
% is retired); the redundancy bullet reports the counts of the RULE readings. The previous body is parked verbatim
% after the section; the every-bar attribution, ladder and ridge-against-lasso bullets moved to Appendix A_fi.tex.
% 2026-10-28 (author: v2 of Sections 6-7 approved preliminarily): section title (v2: attribution is one of three kinds of evidence). Was:
% \section{Variable importance}
% \label{sec:rh_fi}
\ifdefined\rhfipNOrigins\else% AUTO-GENERATED by writeup/make_rehaul_fi_pooled.py -- do not edit; re-run the script.
\providecommand{\ser}[2]{#2}\providecommand{\serx}[3]{#3}
% \rhfip... macros for writeup/rehaul/08_feature_importance.tex (results/tree_shap_w3000_2014start_fomcfix_2026-10-26, the
% FOMC-patched design, its linear reference inputs/fiw_gamma_fomcfix.npz; the window study's ...W... macros:
% results/fi_by_window_2014start_2026-10-22, pre-patch; the final pooled runs, W = 3,000 trading days).
\providecommand{\rhfipNOrigins}{}\renewcommand{\rhfipNOrigins}{126}  % monthly origins (re-estimations explained)
\providecommand{\rhfipNRows}{}\renewcommand{\rhfipNRows}{33{,}970}  % bars explained, every bar of the evaluation span
\providecommand{\rhfipOriginEvery}{}\renewcommand{\rhfipOriginEvery}{21}  % trading days between origins
\providecommand{\rhfipLastBlock}{}\renewcommand{\rhfipLastBlock}{18}  % trading days of the last origin's block
\providecommand{\rhfipFirstOrigin}{}\renewcommand{\rhfipFirstOrigin}{2~January~2014}  % first origin (10:00 bar)
\providecommand{\rhfipLastOrigin}{}\renewcommand{\rhfipLastOrigin}{5~April~2024}  % last origin
\providecommand{\rhfipLastDay}{}\renewcommand{\rhfipLastDay}{30~April~2024}  % day of the last bar explained
\providecommand{\rhfipWindowDays}{}\renewcommand{\rhfipWindowDays}{3{,}000}  % rolling window, trading days
\providecommand{\rhfipWindowBars}{}\renewcommand{\rhfipWindowBars}{39{,}000}  % rolling window, bars
\providecommand{\rhfipNCols}{}\renewcommand{\rhfipNCols}{564}  % predictors of the full set
\providecommand{\rhfipNGroups}{}\renewcommand{\rhfipNGroups}{48}  % predictor groups: 46 exogenous series + HAR components + calendar dummies
\providecommand{\rhfipNExog}{}\renewcommand{\rhfipNExog}{46}  % exogenous series
\providecommand{\rhfipNPerSeries}{}\renewcommand{\rhfipNPerSeries}{12}  % predictors per exogenous series: value and availability averages at six horizons
\providecommand{\rhfipNStd}{}\renewcommand{\rhfipNStd}{12}  % standard groups
\providecommand{\rhfipNNonStd}{}\renewcommand{\rhfipNNonStd}{36}  % non-standard series
\providecommand{\rhfipNStdExog}{}\renewcommand{\rhfipNStdExog}{10}  % standard exogenous series
\providecommand{\rhfipTuneLast}{}\renewcommand{\rhfipTuneLast}{31~December~2013}  % last day of the last tuning fold (fold 3)
\providecommand{\rhfipRidgePen}{}\renewcommand{\rhfipRidgePen}{\ensuremath{10^{3}}}
\providecommand{\rhfipLassoPen}{}\renewcommand{\rhfipLassoPen}{\ensuremath{10^{-4}}}
\providecommand{\rhfipEnetPen}{}\renewcommand{\rhfipEnetPen}{\ensuremath{10^{-3}}}
\providecommand{\rhfipHarPf}{}\renewcommand{\rhfipHarPf}{0.01}  % penalty factor, HAR block
\providecommand{\rhfipBootBlock}{}\renewcommand{\rhfipBootBlock}{12}  % bootstrap block, origins
\providecommand{\rhfipBootB}{}\renewcommand{\rhfipBootB}{2{,}000}
\providecommand{\rhfipCiLevel}{}\renewcommand{\rhfipCiLevel}{95}  % percentile intervals of both rules (RULE_shap.txt, RULE_fiw.txt)
\providecommand{\rhfipShapeMin}{}\renewcommand{\rhfipShapeMin}{0.5}  % M3 selection: share of at least this many percent
\providecommand{\rhfipTopN}{}\renewcommand{\rhfipTopN}{10}  % length of the top lists compared (m2_rank.csv top10_overlap)
\providecommand{\rhfipSlopeUnit}{}\renewcommand{\rhfipSlopeUnit}{1{,}000}  % trading days per unit of the window slope
\providecommand{\rhfipNMom}{}\renewcommand{\rhfipNMom}{seven}  % the futures' realized moments (standard)
\providecommand{\rhfipLgbmGateEq}{}\renewcommand{\rhfipLgbmGateEq}{126}  % origins bit-equal to the merged every-bar walk
\providecommand{\rhfipXgbGateEq}{}\renewcommand{\rhfipXgbGateEq}{126}  % origins bit-equal to the merged every-bar walk
\providecommand{\rhfipRfGateEq}{}\renewcommand{\rhfipRfGateEq}{126}  % origins bit-equal to the merged every-bar walk
\providecommand{\rhfipRfGateAtShap}{}\renewcommand{\rhfipRfGateAtShap}{126}  % forest origins the walk had written when the SHAP gate ran
\providecommand{\rhfipXgbAddMax}{}\renewcommand{\rhfipXgbAddMax}{\ensuremath{3.7\times 10^{-6}}}  % XGBoost additivity, float32
\providecommand{\rhfipXgbAddTol}{}\renewcommand{\rhfipXgbAddTol}{\ensuremath{10^{-5}}}
\providecommand{\rhfipAddMaxDouble}{}\renewcommand{\rhfipAddMaxDouble}{\ensuremath{5.1\times 10^{-14}}}
\providecommand{\rhfipLinGateMax}{}\renewcommand{\rhfipLinGateMax}{\ensuremath{4.4\times 10^{-16}}}  % linear replay vs walk forecasts, W = 3,000
\providecommand{\rhfipInterRows}{}\renewcommand{\rhfipInterRows}{1{,}008}  % M5 rows
\providecommand{\rhfipInterPerOrigin}{}\renewcommand{\rhfipInterPerOrigin}{8}  % M5 rows per origin
\providecommand{\rhfipFomcRows}{}\renewcommand{\rhfipFomcRows}{1{,}014}
\providecommand{\rhfipFomcDaysShap}{}\renewcommand{\rhfipFomcDaysShap}{78}  % FOMC statement days in the SHAP subset (M4; releases.parquet only)
\providecommand{\rhfipFomcDays}{}\renewcommand{\rhfipFomcDays}{81}  % scheduled FOMC statement days in the span (loader calendar)
\providecommand{\rhfipFomcDaysMissing}{}\renewcommand{\rhfipFomcDaysMissing}{3}  % statement days the SHAP subset leaves out (counted with the other days; the patched rerun's reducer keeps releases.parquet's days)
\providecommand{\rhfipNHolSessions}{}\renewcommand{\rhfipNHolSessions}{64}  % NYSE-holiday futures sessions among the explained bars (not scored in Table 1)
\providecommand{\rhfipNHolBars}{}\renewcommand{\rhfipNHolBars}{443}  % their regular-hours bars
\providecommand{\rhfipVixLo}{}\renewcommand{\rhfipVixLo}{13.96}
\providecommand{\rhfipVixHi}{}\renewcommand{\rhfipVixHi}{19.02}
\providecommand{\rhfipLgbmHar}{}\renewcommand{\rhfipLgbmHar}{51.7}
\providecommand{\rhfipLgbmHarLo}{}\renewcommand{\rhfipLgbmHarLo}{49.5}
\providecommand{\rhfipLgbmHarHi}{}\renewcommand{\rhfipLgbmHarHi}{53.7}
\providecommand{\rhfipLgbmMom}{}\renewcommand{\rhfipLgbmMom}{22.3}
\providecommand{\rhfipLgbmNs}{}\renewcommand{\rhfipLgbmNs}{21.8}
\providecommand{\rhfipLgbmNsLo}{}\renewcommand{\rhfipLgbmNsLo}{20.4}
\providecommand{\rhfipLgbmNsHi}{}\renewcommand{\rhfipLgbmNsHi}{23.4}
\providecommand{\rhfipLgbmNsEx}{}\renewcommand{\rhfipLgbmNsEx}{45.9}
\providecommand{\rhfipLgbmNsExLo}{}\renewcommand{\rhfipLgbmNsExLo}{44.2}
\providecommand{\rhfipLgbmNsExHi}{}\renewcommand{\rhfipLgbmNsExHi}{47.7}
\providecommand{\rhfipXgbHar}{}\renewcommand{\rhfipXgbHar}{52.0}
\providecommand{\rhfipXgbHarLo}{}\renewcommand{\rhfipXgbHarLo}{50.0}
\providecommand{\rhfipXgbHarHi}{}\renewcommand{\rhfipXgbHarHi}{53.8}
\providecommand{\rhfipXgbMom}{}\renewcommand{\rhfipXgbMom}{23.6}
\providecommand{\rhfipXgbNs}{}\renewcommand{\rhfipXgbNs}{20.6}
\providecommand{\rhfipXgbNsLo}{}\renewcommand{\rhfipXgbNsLo}{19.2}
\providecommand{\rhfipXgbNsHi}{}\renewcommand{\rhfipXgbNsHi}{22.2}
\providecommand{\rhfipXgbNsEx}{}\renewcommand{\rhfipXgbNsEx}{43.6}
\providecommand{\rhfipXgbNsExLo}{}\renewcommand{\rhfipXgbNsExLo}{41.8}
\providecommand{\rhfipXgbNsExHi}{}\renewcommand{\rhfipXgbNsExHi}{45.7}
\providecommand{\rhfipRfHar}{}\renewcommand{\rhfipRfHar}{57.6}
\providecommand{\rhfipRfHarLo}{}\renewcommand{\rhfipRfHarLo}{55.4}
\providecommand{\rhfipRfHarHi}{}\renewcommand{\rhfipRfHarHi}{59.6}
\providecommand{\rhfipRfMom}{}\renewcommand{\rhfipRfMom}{20.7}
\providecommand{\rhfipRfNs}{}\renewcommand{\rhfipRfNs}{18.6}
\providecommand{\rhfipRfNsLo}{}\renewcommand{\rhfipRfNsLo}{17.2}
\providecommand{\rhfipRfNsHi}{}\renewcommand{\rhfipRfNsHi}{20.3}
\providecommand{\rhfipRfNsEx}{}\renewcommand{\rhfipRfNsEx}{44.2}
\providecommand{\rhfipRfNsExLo}{}\renewcommand{\rhfipRfNsExLo}{42.1}
\providecommand{\rhfipRfNsExHi}{}\renewcommand{\rhfipRfNsExHi}{46.5}
\providecommand{\rhfipRidgeNsEx}{}\renewcommand{\rhfipRidgeNsEx}{68.8}
\providecommand{\rhfipRidgeNs}{}\renewcommand{\rhfipRidgeNs}{55.6}
\providecommand{\rhfipRidgeHar}{}\renewcommand{\rhfipRidgeHar}{17.5}
\providecommand{\rhfipLassoNsEx}{}\renewcommand{\rhfipLassoNsEx}{82.5}
\providecommand{\rhfipLassoNs}{}\renewcommand{\rhfipLassoNs}{76.5}
\providecommand{\rhfipLassoHar}{}\renewcommand{\rhfipLassoHar}{6.7}
\providecommand{\rhfipEnetNsEx}{}\renewcommand{\rhfipEnetNsEx}{60.0}
\providecommand{\rhfipEnetNs}{}\renewcommand{\rhfipEnetNs}{42.9}
\providecommand{\rhfipEnetHar}{}\renewcommand{\rhfipEnetHar}{26.4}
\providecommand{\rhfipHarMin}{}\renewcommand{\rhfipHarMin}{51.7}
\providecommand{\rhfipHarMax}{}\renewcommand{\rhfipHarMax}{57.6}
\providecommand{\rhfipMomMin}{}\renewcommand{\rhfipMomMin}{20.7}
\providecommand{\rhfipMomMax}{}\renewcommand{\rhfipMomMax}{23.6}
\providecommand{\rhfipNsExMin}{}\renewcommand{\rhfipNsExMin}{43.6}
\providecommand{\rhfipNsExMax}{}\renewcommand{\rhfipNsExMax}{45.9}
\providecommand{\rhfipSentMin}{}\renewcommand{\rhfipSentMin}{0.6}
\providecommand{\rhfipSentMax}{}\renewcommand{\rhfipSentMax}{0.8}
\providecommand{\rhfipVdMin}{}\renewcommand{\rhfipVdMin}{0.2}
\providecommand{\rhfipVdMax}{}\renewcommand{\rhfipVdMax}{0.3}
\providecommand{\rhfipFomcBucketMin}{}\renewcommand{\rhfipFomcBucketMin}{1.4}
\providecommand{\rhfipFomcBucketMax}{}\renewcommand{\rhfipFomcBucketMax}{2.3}
\providecommand{\rhfipVvixMin}{}\renewcommand{\rhfipVvixMin}{0.2}
\providecommand{\rhfipVvixMax}{}\renewcommand{\rhfipVvixMax}{0.3}
\providecommand{\rhfipVixThreeMMin}{}\renewcommand{\rhfipVixThreeMMin}{0.3}
\providecommand{\rhfipVixThreeMMax}{}\renewcommand{\rhfipVixThreeMMax}{0.7}
\providecommand{\rhfipFomcDayMin}{}\renewcommand{\rhfipFomcDayMin}{0.1}
\providecommand{\rhfipFomcDayMax}{}\renewcommand{\rhfipFomcDayMax}{0.1}
\providecommand{\rhfipNsOne}{}\renewcommand{\rhfipNsOne}{\ensuremath{\ser{sumabsret_vwstock}{\mathrm{RAV}^{\mathrm{VW}}}}}
\providecommand{\rhfipNsOneMin}{}\renewcommand{\rhfipNsOneMin}{2.4}
\providecommand{\rhfipNsOneMax}{}\renewcommand{\rhfipNsOneMax}{3.9}
\providecommand{\rhfipNsOneLgbm}{}\renewcommand{\rhfipNsOneLgbm}{3.5}
\providecommand{\rhfipNsOneXgb}{}\renewcommand{\rhfipNsOneXgb}{3.9}
\providecommand{\rhfipNsOneRf}{}\renewcommand{\rhfipNsOneRf}{2.4}
\providecommand{\rhfipNsTwo}{}\renewcommand{\rhfipNsTwo}{\ensuremath{\ser{sumabsret_ewstock}{\mathrm{RAV}^{\mathrm{EW}}}}}
\providecommand{\rhfipNsTwoMin}{}\renewcommand{\rhfipNsTwoMin}{1.9}
\providecommand{\rhfipNsTwoMax}{}\renewcommand{\rhfipNsTwoMax}{3.9}
\providecommand{\rhfipNsTwoLgbm}{}\renewcommand{\rhfipNsTwoLgbm}{3.9}
\providecommand{\rhfipNsTwoXgb}{}\renewcommand{\rhfipNsTwoXgb}{3.8}
\providecommand{\rhfipNsTwoRf}{}\renewcommand{\rhfipNsTwoRf}{1.9}
\providecommand{\rhfipNsThree}{}\renewcommand{\rhfipNsThree}{\ensuremath{\ser{fomc_until_inv}{\mathrm{FOMC}^{\mathrm{next}}}}}
\providecommand{\rhfipNsThreeMin}{}\renewcommand{\rhfipNsThreeMin}{1.1}
\providecommand{\rhfipNsThreeMax}{}\renewcommand{\rhfipNsThreeMax}{2.0}
\providecommand{\rhfipNsThreeLgbm}{}\renewcommand{\rhfipNsThreeLgbm}{2.0}
\providecommand{\rhfipNsThreeXgb}{}\renewcommand{\rhfipNsThreeXgb}{2.0}
\providecommand{\rhfipNsThreeRf}{}\renewcommand{\rhfipNsThreeRf}{1.1}
\providecommand{\rhfipNsFour}{}\renewcommand{\rhfipNsFour}{\ensuremath{\ser{ofi_vwstock}{\mathrm{OFI}^{\mathrm{VW}}}}}
\providecommand{\rhfipNsFourMin}{}\renewcommand{\rhfipNsFourMin}{1.4}
\providecommand{\rhfipNsFourMax}{}\renewcommand{\rhfipNsFourMax}{1.8}
\providecommand{\rhfipNsFourLgbm}{}\renewcommand{\rhfipNsFourLgbm}{1.8}
\providecommand{\rhfipNsFourXgb}{}\renewcommand{\rhfipNsFourXgb}{1.7}
\providecommand{\rhfipNsFourRf}{}\renewcommand{\rhfipNsFourRf}{1.4}
\providecommand{\rhfipNsFive}{}\renewcommand{\rhfipNsFive}{\ensuremath{\ser{effspread_vwstock}{\mathrm{EffSpread}^{\mathrm{VW}}}}}
\providecommand{\rhfipNsFiveMin}{}\renewcommand{\rhfipNsFiveMin}{1.2}
\providecommand{\rhfipNsFiveMax}{}\renewcommand{\rhfipNsFiveMax}{1.3}
\providecommand{\rhfipNsFiveLgbm}{}\renewcommand{\rhfipNsFiveLgbm}{1.2}
\providecommand{\rhfipNsFiveXgb}{}\renewcommand{\rhfipNsFiveXgb}{1.2}
\providecommand{\rhfipNsFiveRf}{}\renewcommand{\rhfipNsFiveRf}{1.3}
\providecommand{\rhfipTreeTopOne}{}\renewcommand{\rhfipTreeTopOne}{\ensuremath{\ser{sumabsret}{\mathrm{RAV}}}}
\providecommand{\rhfipTreeTopTwo}{}\renewcommand{\rhfipTreeTopTwo}{\ensuremath{\ser{sumret}{r}}}
\providecommand{\rhfipLgbmTopOneEx}{}\renewcommand{\rhfipLgbmTopOneEx}{20.3}
\providecommand{\rhfipXgbTopOneEx}{}\renewcommand{\rhfipXgbTopOneEx}{24.7}
\providecommand{\rhfipRfTopOneEx}{}\renewcommand{\rhfipRfTopOneEx}{29.2}
\providecommand{\rhfipRidgeTopOne}{}\renewcommand{\rhfipRidgeTopOne}{\ensuremath{\ser{fomc_until_inv}{\mathrm{FOMC}^{\mathrm{next}}}}}
\providecommand{\rhfipRidgeTopOneEx}{}\renewcommand{\rhfipRidgeTopOneEx}{9.0}
\providecommand{\rhfipLgbmNUsedAll}{}\renewcommand{\rhfipLgbmNUsedAll}{42}  % exogenous series split on at every origin
\providecommand{\rhfipLgbmUsedMin}{}\renewcommand{\rhfipLgbmUsedMin}{58}
\providecommand{\rhfipXgbNUsedAll}{}\renewcommand{\rhfipXgbNUsedAll}{43}  % exogenous series split on at every origin
\providecommand{\rhfipXgbUsedMin}{}\renewcommand{\rhfipXgbUsedMin}{72}
\providecommand{\rhfipRfNUsedAll}{}\renewcommand{\rhfipRfNUsedAll}{46}  % exogenous series split on at every origin
\providecommand{\rhfipRfUsedMin}{}\renewcommand{\rhfipRfUsedMin}{126}
\providecommand{\rhfipLgbmDRidge}{}\renewcommand{\rhfipLgbmDRidge}{\ensuremath{-}22.9}
\providecommand{\rhfipLgbmDRidgeLo}{}\renewcommand{\rhfipLgbmDRidgeLo}{\ensuremath{-}24.8}
\providecommand{\rhfipLgbmDRidgeHi}{}\renewcommand{\rhfipLgbmDRidgeHi}{\ensuremath{-}21.1}
\providecommand{\rhfipLgbmDLasso}{}\renewcommand{\rhfipLgbmDLasso}{\ensuremath{-}36.6}
\providecommand{\rhfipLgbmDLassoLo}{}\renewcommand{\rhfipLgbmDLassoLo}{\ensuremath{-}39.1}
\providecommand{\rhfipLgbmDLassoHi}{}\renewcommand{\rhfipLgbmDLassoHi}{\ensuremath{-}34.4}
\providecommand{\rhfipLgbmDEnet}{}\renewcommand{\rhfipLgbmDEnet}{\ensuremath{-}14.1}
\providecommand{\rhfipLgbmDEnetLo}{}\renewcommand{\rhfipLgbmDEnetLo}{\ensuremath{-}15.7}
\providecommand{\rhfipLgbmDEnetHi}{}\renewcommand{\rhfipLgbmDEnetHi}{\ensuremath{-}12.7}
\providecommand{\rhfipXgbDRidge}{}\renewcommand{\rhfipXgbDRidge}{\ensuremath{-}25.1}
\providecommand{\rhfipXgbDRidgeLo}{}\renewcommand{\rhfipXgbDRidgeLo}{\ensuremath{-}27.0}
\providecommand{\rhfipXgbDRidgeHi}{}\renewcommand{\rhfipXgbDRidgeHi}{\ensuremath{-}23.3}
\providecommand{\rhfipXgbDLasso}{}\renewcommand{\rhfipXgbDLasso}{\ensuremath{-}38.9}
\providecommand{\rhfipXgbDLassoLo}{}\renewcommand{\rhfipXgbDLassoLo}{\ensuremath{-}41.4}
\providecommand{\rhfipXgbDLassoHi}{}\renewcommand{\rhfipXgbDLassoHi}{\ensuremath{-}36.6}
\providecommand{\rhfipXgbDEnet}{}\renewcommand{\rhfipXgbDEnet}{\ensuremath{-}16.3}
\providecommand{\rhfipXgbDEnetLo}{}\renewcommand{\rhfipXgbDEnetLo}{\ensuremath{-}17.8}
\providecommand{\rhfipXgbDEnetHi}{}\renewcommand{\rhfipXgbDEnetHi}{\ensuremath{-}14.9}
\providecommand{\rhfipRfDRidge}{}\renewcommand{\rhfipRfDRidge}{\ensuremath{-}24.5}
\providecommand{\rhfipRfDRidgeLo}{}\renewcommand{\rhfipRfDRidgeLo}{\ensuremath{-}26.3}
\providecommand{\rhfipRfDRidgeHi}{}\renewcommand{\rhfipRfDRidgeHi}{\ensuremath{-}22.9}
\providecommand{\rhfipRfDLasso}{}\renewcommand{\rhfipRfDLasso}{\ensuremath{-}38.3}
\providecommand{\rhfipRfDLassoLo}{}\renewcommand{\rhfipRfDLassoLo}{\ensuremath{-}40.5}
\providecommand{\rhfipRfDLassoHi}{}\renewcommand{\rhfipRfDLassoHi}{\ensuremath{-}36.3}
\providecommand{\rhfipRfDEnet}{}\renewcommand{\rhfipRfDEnet}{\ensuremath{-}15.7}
\providecommand{\rhfipRfDEnetLo}{}\renewcommand{\rhfipRfDEnetLo}{\ensuremath{-}17.0}
\providecommand{\rhfipRfDEnetHi}{}\renewcommand{\rhfipRfDEnetHi}{\ensuremath{-}14.4}
\providecommand{\rhfipNDLess}{}\renewcommand{\rhfipNDLess}{27}  % tree-minus-linear D with the 95% interval below 0
\providecommand{\rhfipNDAll}{}\renewcommand{\rhfipNDAll}{27}  % D comparisons: 3 trees x 3 linear models x 3 denominators
\providecommand{\rhfipLgbmSpRidge}{}\renewcommand{\rhfipLgbmSpRidge}{0.57}
\providecommand{\rhfipLgbmTopRidge}{}\renewcommand{\rhfipLgbmTopRidge}{6}
\providecommand{\rhfipLgbmSpLasso}{}\renewcommand{\rhfipLgbmSpLasso}{0.49}
\providecommand{\rhfipLgbmTopLasso}{}\renewcommand{\rhfipLgbmTopLasso}{5}
\providecommand{\rhfipLgbmSpEnet}{}\renewcommand{\rhfipLgbmSpEnet}{0.66}
\providecommand{\rhfipLgbmTopEnet}{}\renewcommand{\rhfipLgbmTopEnet}{6}
\providecommand{\rhfipXgbSpRidge}{}\renewcommand{\rhfipXgbSpRidge}{0.55}
\providecommand{\rhfipXgbTopRidge}{}\renewcommand{\rhfipXgbTopRidge}{6}
\providecommand{\rhfipXgbSpLasso}{}\renewcommand{\rhfipXgbSpLasso}{0.47}
\providecommand{\rhfipXgbTopLasso}{}\renewcommand{\rhfipXgbTopLasso}{5}
\providecommand{\rhfipXgbSpEnet}{}\renewcommand{\rhfipXgbSpEnet}{0.63}
\providecommand{\rhfipXgbTopEnet}{}\renewcommand{\rhfipXgbTopEnet}{6}
\providecommand{\rhfipRfSpRidge}{}\renewcommand{\rhfipRfSpRidge}{0.53}
\providecommand{\rhfipRfTopRidge}{}\renewcommand{\rhfipRfTopRidge}{5}
\providecommand{\rhfipRfSpLasso}{}\renewcommand{\rhfipRfSpLasso}{0.43}
\providecommand{\rhfipRfTopLasso}{}\renewcommand{\rhfipRfTopLasso}{4}
\providecommand{\rhfipRfSpEnet}{}\renewcommand{\rhfipRfSpEnet}{0.60}
\providecommand{\rhfipRfTopEnet}{}\renewcommand{\rhfipRfTopEnet}{5}
\providecommand{\rhfipSpLgbmXgb}{}\renewcommand{\rhfipSpLgbmXgb}{0.99}
\providecommand{\rhfipSpLgbmRf}{}\renewcommand{\rhfipSpLgbmRf}{0.90}
\providecommand{\rhfipSpXgbRf}{}\renewcommand{\rhfipSpXgbRf}{0.89}
\providecommand{\rhfipRidgeRankFomcUntil}{}\renewcommand{\rhfipRidgeRankFomcUntil}{1}
\providecommand{\rhfipTreeRankFomcUntilMin}{}\renewcommand{\rhfipTreeRankFomcUntilMin}{6}
\providecommand{\rhfipTreeRankFomcUntilMax}{}\renewcommand{\rhfipTreeRankFomcUntilMax}{10}
\providecommand{\rhfipRidgeRankFomcDay}{}\renewcommand{\rhfipRidgeRankFomcDay}{9}
\providecommand{\rhfipTreeRankFomcDayMin}{}\renewcommand{\rhfipTreeRankFomcDayMin}{39}
\providecommand{\rhfipTreeRankFomcDayMax}{}\renewcommand{\rhfipTreeRankFomcDayMax}{43}
\providecommand{\rhfipFomcUntil}{}\renewcommand{\rhfipFomcUntil}{\ensuremath{\ser{fomc_until_inv}{\mathrm{FOMC}^{\mathrm{next}}}}}
\providecommand{\rhfipFomcDay}{}\renewcommand{\rhfipFomcDay}{\ensuremath{\ser{fomc_day}{\mathrm{FOMC}^{\mathrm{day}}}}}
\providecommand{\rhfipAbsVw}{}\renewcommand{\rhfipAbsVw}{\ensuremath{\ser{sumabsret_vwstock}{\mathrm{RAV}^{\mathrm{VW}}}}}
\providecommand{\rhfipOfi}{}\renewcommand{\rhfipOfi}{\ensuremath{\ser{ofi_vwstock}{\mathrm{OFI}^{\mathrm{VW}}}}}
\providecommand{\rhfipLgbmNeff}{}\renewcommand{\rhfipLgbmNeff}{11.0}
\providecommand{\rhfipLgbmNHalf}{}\renewcommand{\rhfipLgbmNHalf}{4}
\providecommand{\rhfipLgbmNEighty}{}\renewcommand{\rhfipLgbmNEighty}{13}
\providecommand{\rhfipXgbNeff}{}\renewcommand{\rhfipXgbNeff}{9.1}
\providecommand{\rhfipXgbNHalf}{}\renewcommand{\rhfipXgbNHalf}{4}
\providecommand{\rhfipXgbNEighty}{}\renewcommand{\rhfipXgbNEighty}{11}
\providecommand{\rhfipRfNeff}{}\renewcommand{\rhfipRfNeff}{8.9}
\providecommand{\rhfipRfNHalf}{}\renewcommand{\rhfipRfNHalf}{4}
\providecommand{\rhfipRfNEighty}{}\renewcommand{\rhfipRfNEighty}{16}
\providecommand{\rhfipRidgeNeff}{}\renewcommand{\rhfipRidgeNeff}{28.5}
\providecommand{\rhfipRidgeNHalf}{}\renewcommand{\rhfipRidgeNHalf}{11}
\providecommand{\rhfipRidgeNEighty}{}\renewcommand{\rhfipRidgeNEighty}{26}
\providecommand{\rhfipLassoNeff}{}\renewcommand{\rhfipLassoNeff}{5.1}
\providecommand{\rhfipLassoNHalf}{}\renewcommand{\rhfipLassoNHalf}{2}
\providecommand{\rhfipLassoNEighty}{}\renewcommand{\rhfipLassoNEighty}{11}
\providecommand{\rhfipEnetNeff}{}\renewcommand{\rhfipEnetNeff}{16.2}
\providecommand{\rhfipEnetNHalf}{}\renewcommand{\rhfipEnetNHalf}{6}
\providecommand{\rhfipEnetNEighty}{}\renewcommand{\rhfipEnetNEighty}{15}
\providecommand{\rhfipYearFirst}{}\renewcommand{\rhfipYearFirst}{2014}
\providecommand{\rhfipYearLast}{}\renewcommand{\rhfipYearLast}{2024}
\providecommand{\rhfipLgbmNsYearMin}{}\renewcommand{\rhfipLgbmNsYearMin}{18.1}
\providecommand{\rhfipLgbmNsYearMinYr}{}\renewcommand{\rhfipLgbmNsYearMinYr}{2021}
\providecommand{\rhfipLgbmNsYearMax}{}\renewcommand{\rhfipLgbmNsYearMax}{25.6}
\providecommand{\rhfipLgbmNsYearMaxYr}{}\renewcommand{\rhfipLgbmNsYearMaxYr}{2015}
\providecommand{\rhfipLgbmHarYearMax}{}\renewcommand{\rhfipLgbmHarYearMax}{57.4}
\providecommand{\rhfipLgbmHarYearMaxYr}{}\renewcommand{\rhfipLgbmHarYearMaxYr}{2021}
\providecommand{\rhfipLgbmHarYearMin}{}\renewcommand{\rhfipLgbmHarYearMin}{44.6}
\providecommand{\rhfipLgbmHarYearMinYr}{}\renewcommand{\rhfipLgbmHarYearMinYr}{2017}
\providecommand{\rhfipXgbNsYearMin}{}\renewcommand{\rhfipXgbNsYearMin}{17.2}
\providecommand{\rhfipXgbNsYearMinYr}{}\renewcommand{\rhfipXgbNsYearMinYr}{2021}
\providecommand{\rhfipXgbNsYearMax}{}\renewcommand{\rhfipXgbNsYearMax}{24.3}
\providecommand{\rhfipXgbNsYearMaxYr}{}\renewcommand{\rhfipXgbNsYearMaxYr}{2015}
\providecommand{\rhfipXgbHarYearMax}{}\renewcommand{\rhfipXgbHarYearMax}{57.0}
\providecommand{\rhfipXgbHarYearMaxYr}{}\renewcommand{\rhfipXgbHarYearMaxYr}{2021}
\providecommand{\rhfipXgbHarYearMin}{}\renewcommand{\rhfipXgbHarYearMin}{45.8}
\providecommand{\rhfipXgbHarYearMinYr}{}\renewcommand{\rhfipXgbHarYearMinYr}{2017}
\providecommand{\rhfipRfNsYearMin}{}\renewcommand{\rhfipRfNsYearMin}{15.2}
\providecommand{\rhfipRfNsYearMinYr}{}\renewcommand{\rhfipRfNsYearMinYr}{2021}
\providecommand{\rhfipRfNsYearMax}{}\renewcommand{\rhfipRfNsYearMax}{22.9}
\providecommand{\rhfipRfNsYearMaxYr}{}\renewcommand{\rhfipRfNsYearMaxYr}{2015}
\providecommand{\rhfipRfHarYearMax}{}\renewcommand{\rhfipRfHarYearMax}{62.3}
\providecommand{\rhfipRfHarYearMaxYr}{}\renewcommand{\rhfipRfHarYearMaxYr}{2019}
\providecommand{\rhfipRfHarYearMin}{}\renewcommand{\rhfipRfHarYearMin}{50.4}
\providecommand{\rhfipRfHarYearMinYr}{}\renewcommand{\rhfipRfHarYearMinYr}{2017}
\providecommand{\rhfipNsYearMinYr}{}\renewcommand{\rhfipNsYearMinYr}{2021}  % the year of every tree's lowest non-standard share
\providecommand{\rhfipDipFirst}{}\renewcommand{\rhfipDipFirst}{2018}  % first of the years in which every tree's non-standard share is below its pooled share
\providecommand{\rhfipDipLast}{}\renewcommand{\rhfipDipLast}{2022}
\providecommand{\rhfipRidgeNsExYearMin}{}\renewcommand{\rhfipRidgeNsExYearMin}{65.1}
\providecommand{\rhfipRidgeNsExYearMinYr}{}\renewcommand{\rhfipRidgeNsExYearMinYr}{2021}
\providecommand{\rhfipRidgeNsExYearMax}{}\renewcommand{\rhfipRidgeNsExYearMax}{72.4}
\providecommand{\rhfipRidgeNsExYearMaxYr}{}\renewcommand{\rhfipRidgeNsExYearMaxYr}{2017}
\providecommand{\rhfipLassoNsExYearMin}{}\renewcommand{\rhfipLassoNsExYearMin}{77.1}
\providecommand{\rhfipLassoNsExYearMinYr}{}\renewcommand{\rhfipLassoNsExYearMinYr}{2021}
\providecommand{\rhfipLassoNsExYearMax}{}\renewcommand{\rhfipLassoNsExYearMax}{87.1}
\providecommand{\rhfipLassoNsExYearMaxYr}{}\renewcommand{\rhfipLassoNsExYearMaxYr}{2017}
\providecommand{\rhfipEnetNsExYearMin}{}\renewcommand{\rhfipEnetNsExYearMin}{56.1}
\providecommand{\rhfipEnetNsExYearMinYr}{}\renewcommand{\rhfipEnetNsExYearMinYr}{2019}
\providecommand{\rhfipEnetNsExYearMax}{}\renewcommand{\rhfipEnetNsExYearMax}{65.2}
\providecommand{\rhfipEnetNsExYearMaxYr}{}\renewcommand{\rhfipEnetNsExYearMaxYr}{2014}
\providecommand{\rhfipLgbmNsExYearMin}{}\renewcommand{\rhfipLgbmNsExYearMin}{42.3}
\providecommand{\rhfipLgbmNsExYearMax}{}\renewcommand{\rhfipLgbmNsExYearMax}{50.1}
\providecommand{\rhfipXgbNsExYearMin}{}\renewcommand{\rhfipXgbNsExYearMin}{39.4}
\providecommand{\rhfipXgbNsExYearMax}{}\renewcommand{\rhfipXgbNsExYearMax}{48.9}
\providecommand{\rhfipRfNsExYearMin}{}\renewcommand{\rhfipRfNsExYearMin}{39.5}
\providecommand{\rhfipRfNsExYearMax}{}\renewcommand{\rhfipRfNsExYearMax}{50.9}
\providecommand{\rhfipLgbmNShape}{}\renewcommand{\rhfipLgbmNShape}{11}  % non-standard series at >= 0.5% (M3)
\providecommand{\rhfipLgbmNMono}{}\renewcommand{\rhfipLgbmNMono}{7}
\providecommand{\rhfipLgbmNIrr}{}\renewcommand{\rhfipLgbmNIrr}{3}
\providecommand{\rhfipLgbmNStep}{}\renewcommand{\rhfipLgbmNStep}{1}
\providecommand{\rhfipLgbmNFlat}{}\renewcommand{\rhfipLgbmNFlat}{0}
\providecommand{\rhfipLgbmNThresh}{}\renewcommand{\rhfipLgbmNThresh}{0}
\providecommand{\rhfipLgbmNU}{}\renewcommand{\rhfipLgbmNU}{0}
\providecommand{\rhfipLgbmNAsym}{}\renewcommand{\rhfipLgbmNAsym}{9}
\providecommand{\rhfipXgbNShape}{}\renewcommand{\rhfipXgbNShape}{10}  % non-standard series at >= 0.5% (M3)
\providecommand{\rhfipXgbNMono}{}\renewcommand{\rhfipXgbNMono}{5}
\providecommand{\rhfipXgbNIrr}{}\renewcommand{\rhfipXgbNIrr}{3}
\providecommand{\rhfipXgbNStep}{}\renewcommand{\rhfipXgbNStep}{0}
\providecommand{\rhfipXgbNFlat}{}\renewcommand{\rhfipXgbNFlat}{1}
\providecommand{\rhfipXgbNThresh}{}\renewcommand{\rhfipXgbNThresh}{1}
\providecommand{\rhfipXgbNU}{}\renewcommand{\rhfipXgbNU}{0}
\providecommand{\rhfipXgbNAsym}{}\renewcommand{\rhfipXgbNAsym}{8}
\providecommand{\rhfipRfNShape}{}\renewcommand{\rhfipRfNShape}{13}  % non-standard series at >= 0.5% (M3)
\providecommand{\rhfipRfNMono}{}\renewcommand{\rhfipRfNMono}{10}
\providecommand{\rhfipRfNIrr}{}\renewcommand{\rhfipRfNIrr}{3}
\providecommand{\rhfipRfNStep}{}\renewcommand{\rhfipRfNStep}{0}
\providecommand{\rhfipRfNFlat}{}\renewcommand{\rhfipRfNFlat}{0}
\providecommand{\rhfipRfNThresh}{}\renewcommand{\rhfipRfNThresh}{0}
\providecommand{\rhfipRfNU}{}\renewcommand{\rhfipRfNU}{0}
\providecommand{\rhfipRfNAsym}{}\renewcommand{\rhfipRfNAsym}{13}
\providecommand{\rhfipNUThreshAll}{}\renewcommand{\rhfipNUThreshAll}{1}  % U-shaped, inverted-U or threshold profiles in any tree
\providecommand{\rhfipAbsVwLgbmLo}{}\renewcommand{\rhfipAbsVwLgbmLo}{\ensuremath{-}1.4}
\providecommand{\rhfipAbsVwLgbmHi}{}\renewcommand{\rhfipAbsVwLgbmHi}{3.9}
\providecommand{\rhfipAbsVwXgbLo}{}\renewcommand{\rhfipAbsVwXgbLo}{\ensuremath{-}1.5}
\providecommand{\rhfipAbsVwXgbHi}{}\renewcommand{\rhfipAbsVwXgbHi}{4.3}
\providecommand{\rhfipAbsVwRfLo}{}\renewcommand{\rhfipAbsVwRfLo}{\ensuremath{-}0.9}
\providecommand{\rhfipAbsVwRfHi}{}\renewcommand{\rhfipAbsVwRfHi}{2.3}
\providecommand{\rhfipOfiLgbmLo}{}\renewcommand{\rhfipOfiLgbmLo}{0.8}
\providecommand{\rhfipOfiLgbmHi}{}\renewcommand{\rhfipOfiLgbmHi}{\ensuremath{-}1.1}
\providecommand{\rhfipOfiXgbLo}{}\renewcommand{\rhfipOfiXgbLo}{0.8}
\providecommand{\rhfipOfiXgbHi}{}\renewcommand{\rhfipOfiXgbHi}{\ensuremath{-}0.9}
\providecommand{\rhfipOfiRfLo}{}\renewcommand{\rhfipOfiRfLo}{0.8}
\providecommand{\rhfipOfiRfHi}{}\renewcommand{\rhfipOfiRfHi}{\ensuremath{-}0.7}
\providecommand{\rhfipUntilLgbmLo}{}\renewcommand{\rhfipUntilLgbmLo}{\ensuremath{-}0.3}
\providecommand{\rhfipUntilLgbmHi}{}\renewcommand{\rhfipUntilLgbmHi}{2.9}
\providecommand{\rhfipUntilXgbLo}{}\renewcommand{\rhfipUntilXgbLo}{0.2}
\providecommand{\rhfipUntilXgbHi}{}\renewcommand{\rhfipUntilXgbHi}{2.4}
\providecommand{\rhfipUntilRfLo}{}\renewcommand{\rhfipUntilRfLo}{\ensuremath{-}0.0}
\providecommand{\rhfipUntilRfHi}{}\renewcommand{\rhfipUntilRfHi}{1.5}
\providecommand{\rhfipUntilHorizon}{}\renewcommand{\rhfipUntilHorizon}{25}  % horizon of fomc_until_inv's dominant column
\providecommand{\rhfipUntilHorizonXgb}{}\renewcommand{\rhfipUntilHorizonXgb}{1}  % the same in XGBoost (differs from the other trees')
\providecommand{\rhfipBarFirst}{}\renewcommand{\rhfipBarFirst}{10:00}
\providecommand{\rhfipBarLast}{}\renewcommand{\rhfipBarLast}{16:00}
\providecommand{\rhfipBarPeakA}{}\renewcommand{\rhfipBarPeakA}{14:30}
\providecommand{\rhfipBarPeakB}{}\renewcommand{\rhfipBarPeakB}{15:00}
\providecommand{\rhfipNBars}{}\renewcommand{\rhfipNBars}{13}
\providecommand{\rhfipLgbmNsFomc}{}\renewcommand{\rhfipLgbmNsFomc}{42.6}
\providecommand{\rhfipLgbmNsOther}{}\renewcommand{\rhfipLgbmNsOther}{20.9}
\providecommand{\rhfipLgbmUntilFomc}{}\renewcommand{\rhfipLgbmUntilFomc}{25.4}
\providecommand{\rhfipLgbmUntilOther}{}\renewcommand{\rhfipLgbmUntilOther}{1.06}
\providecommand{\rhfipLgbmFomcBucketFomc}{}\renewcommand{\rhfipLgbmFomcBucketFomc}{27.3}
\providecommand{\rhfipLgbmFomcBucketOther}{}\renewcommand{\rhfipLgbmFomcBucketOther}{1.3}
\providecommand{\rhfipLgbmHarFomc}{}\renewcommand{\rhfipLgbmHarFomc}{39.2}
\providecommand{\rhfipLgbmHarOther}{}\renewcommand{\rhfipLgbmHarOther}{52.2}
\providecommand{\rhfipLgbmNsFirstBar}{}\renewcommand{\rhfipLgbmNsFirstBar}{23.9}
\providecommand{\rhfipLgbmNsLastBar}{}\renewcommand{\rhfipLgbmNsLastBar}{20.5}
\providecommand{\rhfipLgbmNsAftMin}{}\renewcommand{\rhfipLgbmNsAftMin}{24.3}
\providecommand{\rhfipLgbmNsAftMax}{}\renewcommand{\rhfipLgbmNsAftMax}{25.2}
\providecommand{\rhfipLgbmNsMidMin}{}\renewcommand{\rhfipLgbmNsMidMin}{20.1}
\providecommand{\rhfipLgbmNsMidMax}{}\renewcommand{\rhfipLgbmNsMidMax}{21.9}
\providecommand{\rhfipLgbmFomcBucketAft}{}\renewcommand{\rhfipLgbmFomcBucketAft}{5.5}
\providecommand{\rhfipLgbmFomcBucketOtherBars}{}\renewcommand{\rhfipLgbmFomcBucketOtherBars}{1.4}  % median FOMC-announcement share over the other 11 bars
\providecommand{\rhfipLgbmFomcDiff}{}\renewcommand{\rhfipLgbmFomcDiff}{\ensuremath{+}21.7}
\providecommand{\rhfipLgbmFomcDiffLo}{}\renewcommand{\rhfipLgbmFomcDiffLo}{\ensuremath{+}19.3}
\providecommand{\rhfipLgbmFomcDiffHi}{}\renewcommand{\rhfipLgbmFomcDiffHi}{\ensuremath{+}24.1}
\providecommand{\rhfipLgbmLateDiff}{}\renewcommand{\rhfipLgbmLateDiff}{\ensuremath{-}3.4}
\providecommand{\rhfipLgbmLateDiffLo}{}\renewcommand{\rhfipLgbmLateDiffLo}{\ensuremath{-}4.3}
\providecommand{\rhfipLgbmLateDiffHi}{}\renewcommand{\rhfipLgbmLateDiffHi}{\ensuremath{-}2.5}
\providecommand{\rhfipLgbmVixDiff}{}\renewcommand{\rhfipLgbmVixDiff}{\ensuremath{-}3.2}
\providecommand{\rhfipLgbmVixDiffLo}{}\renewcommand{\rhfipLgbmVixDiffLo}{\ensuremath{-}5.3}
\providecommand{\rhfipLgbmVixDiffHi}{}\renewcommand{\rhfipLgbmVixDiffHi}{\ensuremath{-}1.1}
\providecommand{\rhfipLgbmVixHar}{}\renewcommand{\rhfipLgbmVixHar}{\ensuremath{+}5.1}
\providecommand{\rhfipLgbmVixHarLo}{}\renewcommand{\rhfipLgbmVixHarLo}{\ensuremath{+}1.8}
\providecommand{\rhfipLgbmVixHarHi}{}\renewcommand{\rhfipLgbmVixHarHi}{\ensuremath{+}8.4}
\providecommand{\rhfipLgbmVixDiffEx}{}\renewcommand{\rhfipLgbmVixDiffEx}{\ensuremath{-}2.3}
\providecommand{\rhfipLgbmVixDiffExLo}{}\renewcommand{\rhfipLgbmVixDiffExLo}{\ensuremath{-}4.5}
\providecommand{\rhfipLgbmVixDiffExHi}{}\renewcommand{\rhfipLgbmVixDiffExHi}{\ensuremath{-}0.2}
\providecommand{\rhfipXgbNsFomc}{}\renewcommand{\rhfipXgbNsFomc}{41.1}
\providecommand{\rhfipXgbNsOther}{}\renewcommand{\rhfipXgbNsOther}{19.7}
\providecommand{\rhfipXgbUntilFomc}{}\renewcommand{\rhfipXgbUntilFomc}{23.9}
\providecommand{\rhfipXgbUntilOther}{}\renewcommand{\rhfipXgbUntilOther}{1.10}
\providecommand{\rhfipXgbFomcBucketFomc}{}\renewcommand{\rhfipXgbFomcBucketFomc}{25.6}
\providecommand{\rhfipXgbFomcBucketOther}{}\renewcommand{\rhfipXgbFomcBucketOther}{1.3}
\providecommand{\rhfipXgbHarFomc}{}\renewcommand{\rhfipXgbHarFomc}{39.1}
\providecommand{\rhfipXgbHarOther}{}\renewcommand{\rhfipXgbHarOther}{52.5}
\providecommand{\rhfipXgbNsFirstBar}{}\renewcommand{\rhfipXgbNsFirstBar}{22.5}
\providecommand{\rhfipXgbNsLastBar}{}\renewcommand{\rhfipXgbNsLastBar}{19.5}
\providecommand{\rhfipXgbNsAftMin}{}\renewcommand{\rhfipXgbNsAftMin}{23.5}
\providecommand{\rhfipXgbNsAftMax}{}\renewcommand{\rhfipXgbNsAftMax}{23.7}
\providecommand{\rhfipXgbNsMidMin}{}\renewcommand{\rhfipXgbNsMidMin}{19.1}
\providecommand{\rhfipXgbNsMidMax}{}\renewcommand{\rhfipXgbNsMidMax}{20.4}
\providecommand{\rhfipXgbFomcBucketAft}{}\renewcommand{\rhfipXgbFomcBucketAft}{5.8}
\providecommand{\rhfipXgbFomcBucketOtherBars}{}\renewcommand{\rhfipXgbFomcBucketOtherBars}{1.3}  % median FOMC-announcement share over the other 11 bars
\providecommand{\rhfipXgbFomcDiff}{}\renewcommand{\rhfipXgbFomcDiff}{\ensuremath{+}21.4}
\providecommand{\rhfipXgbFomcDiffLo}{}\renewcommand{\rhfipXgbFomcDiffLo}{\ensuremath{+}19.3}
\providecommand{\rhfipXgbFomcDiffHi}{}\renewcommand{\rhfipXgbFomcDiffHi}{\ensuremath{+}23.5}
\providecommand{\rhfipXgbLateDiff}{}\renewcommand{\rhfipXgbLateDiff}{\ensuremath{-}3.1}
\providecommand{\rhfipXgbLateDiffLo}{}\renewcommand{\rhfipXgbLateDiffLo}{\ensuremath{-}4.0}
\providecommand{\rhfipXgbLateDiffHi}{}\renewcommand{\rhfipXgbLateDiffHi}{\ensuremath{-}2.3}
\providecommand{\rhfipXgbVixDiff}{}\renewcommand{\rhfipXgbVixDiff}{\ensuremath{-}3.0}
\providecommand{\rhfipXgbVixDiffLo}{}\renewcommand{\rhfipXgbVixDiffLo}{\ensuremath{-}5.1}
\providecommand{\rhfipXgbVixDiffHi}{}\renewcommand{\rhfipXgbVixDiffHi}{\ensuremath{-}0.9}
\providecommand{\rhfipXgbVixHar}{}\renewcommand{\rhfipXgbVixHar}{\ensuremath{+}3.9}
\providecommand{\rhfipXgbVixHarLo}{}\renewcommand{\rhfipXgbVixHarLo}{\ensuremath{+}0.4}
\providecommand{\rhfipXgbVixHarHi}{}\renewcommand{\rhfipXgbVixHarHi}{\ensuremath{+}7.2}
\providecommand{\rhfipXgbVixDiffEx}{}\renewcommand{\rhfipXgbVixDiffEx}{\ensuremath{-}3.2}
\providecommand{\rhfipXgbVixDiffExLo}{}\renewcommand{\rhfipXgbVixDiffExLo}{\ensuremath{-}5.7}
\providecommand{\rhfipXgbVixDiffExHi}{}\renewcommand{\rhfipXgbVixDiffExHi}{\ensuremath{-}0.8}
\providecommand{\rhfipRfNsFomc}{}\renewcommand{\rhfipRfNsFomc}{28.9}
\providecommand{\rhfipRfNsOther}{}\renewcommand{\rhfipRfNsOther}{18.3}
\providecommand{\rhfipRfUntilFomc}{}\renewcommand{\rhfipRfUntilFomc}{12.4}
\providecommand{\rhfipRfFomcBucketFomc}{}\renewcommand{\rhfipRfFomcBucketFomc}{13.0}
\providecommand{\rhfipRfFomcBucketOther}{}\renewcommand{\rhfipRfFomcBucketOther}{1.0}
\providecommand{\rhfipRfHarFomc}{}\renewcommand{\rhfipRfHarFomc}{50.8}
\providecommand{\rhfipRfHarOther}{}\renewcommand{\rhfipRfHarOther}{57.9}
\providecommand{\rhfipRfNsFirstBar}{}\renewcommand{\rhfipRfNsFirstBar}{20.1}
\providecommand{\rhfipRfNsLastBar}{}\renewcommand{\rhfipRfNsLastBar}{17.8}
\providecommand{\rhfipRfNsAftMin}{}\renewcommand{\rhfipRfNsAftMin}{19.5}
\providecommand{\rhfipRfNsAftMax}{}\renewcommand{\rhfipRfNsAftMax}{19.9}
\providecommand{\rhfipRfNsMidMin}{}\renewcommand{\rhfipRfNsMidMin}{17.8}
\providecommand{\rhfipRfNsMidMax}{}\renewcommand{\rhfipRfNsMidMax}{18.6}
\providecommand{\rhfipRfFomcBucketAft}{}\renewcommand{\rhfipRfFomcBucketAft}{3.1}
\providecommand{\rhfipRfFomcBucketOtherBars}{}\renewcommand{\rhfipRfFomcBucketOtherBars}{1.0}  % median FOMC-announcement share over the other 11 bars
\providecommand{\rhfipRfFomcDiff}{}\renewcommand{\rhfipRfFomcDiff}{\ensuremath{+}10.6}
\providecommand{\rhfipRfFomcDiffLo}{}\renewcommand{\rhfipRfFomcDiffLo}{\ensuremath{+}8.3}
\providecommand{\rhfipRfFomcDiffHi}{}\renewcommand{\rhfipRfFomcDiffHi}{\ensuremath{+}13.1}
\providecommand{\rhfipRfLateDiff}{}\renewcommand{\rhfipRfLateDiff}{\ensuremath{-}2.3}
\providecommand{\rhfipRfLateDiffLo}{}\renewcommand{\rhfipRfLateDiffLo}{\ensuremath{-}3.1}
\providecommand{\rhfipRfLateDiffHi}{}\renewcommand{\rhfipRfLateDiffHi}{\ensuremath{-}1.6}
\providecommand{\rhfipRfVixDiff}{}\renewcommand{\rhfipRfVixDiff}{\ensuremath{-}3.3}
\providecommand{\rhfipRfVixDiffLo}{}\renewcommand{\rhfipRfVixDiffLo}{\ensuremath{-}5.1}
\providecommand{\rhfipRfVixDiffHi}{}\renewcommand{\rhfipRfVixDiffHi}{\ensuremath{-}1.4}
\providecommand{\rhfipRfVixHar}{}\renewcommand{\rhfipRfVixHar}{\ensuremath{+}4.2}
\providecommand{\rhfipRfVixHarLo}{}\renewcommand{\rhfipRfVixHarLo}{\ensuremath{+}1.0}
\providecommand{\rhfipRfVixHarHi}{}\renewcommand{\rhfipRfVixHarHi}{\ensuremath{+}7.3}
\providecommand{\rhfipRfVixDiffEx}{}\renewcommand{\rhfipRfVixDiffEx}{\ensuremath{-}3.7}
\providecommand{\rhfipRfVixDiffExLo}{}\renewcommand{\rhfipRfVixDiffExLo}{\ensuremath{-}6.3}
\providecommand{\rhfipRfVixDiffExHi}{}\renewcommand{\rhfipRfVixDiffExHi}{\ensuremath{-}1.2}
\providecommand{\rhfipLgbmInterVar}{}\renewcommand{\rhfipLgbmInterVar}{33.1}
\providecommand{\rhfipLgbmInterCol}{}\renewcommand{\rhfipLgbmInterCol}{45.3}
\providecommand{\rhfipLgbmInterNs}{}\renewcommand{\rhfipLgbmInterNs}{49.6}
\providecommand{\rhfipLgbmInterStd}{}\renewcommand{\rhfipLgbmInterStd}{26.8}
\providecommand{\rhfipLgbmInterStdEx}{}\renewcommand{\rhfipLgbmInterStdEx}{40.2}  % descriptive: standard exogenous series only (no HAR, no calendar)
\providecommand{\rhfipLgbmInterHar}{}\renewcommand{\rhfipLgbmInterHar}{17.1}
\providecommand{\rhfipLgbmInterWeakMin}{}\renewcommand{\rhfipLgbmInterWeakMin}{54}
\providecommand{\rhfipLgbmInterWeakMax}{}\renewcommand{\rhfipLgbmInterWeakMax}{66}
\providecommand{\rhfipLgbmPairOne}{}\renewcommand{\rhfipLgbmPairOne}{1.25}
\providecommand{\rhfipLgbmPairTwo}{}\renewcommand{\rhfipLgbmPairTwo}{0.86}
\providecommand{\rhfipLgbmCellHarHar}{}\renewcommand{\rhfipLgbmCellHarHar}{48.0}
\providecommand{\rhfipXgbInterVar}{}\renewcommand{\rhfipXgbInterVar}{30.9}
\providecommand{\rhfipXgbInterCol}{}\renewcommand{\rhfipXgbInterCol}{42.5}
\providecommand{\rhfipXgbInterNs}{}\renewcommand{\rhfipXgbInterNs}{45.7}
\providecommand{\rhfipXgbInterStd}{}\renewcommand{\rhfipXgbInterStd}{26.0}
\providecommand{\rhfipXgbInterStdEx}{}\renewcommand{\rhfipXgbInterStdEx}{37.8}  % descriptive: standard exogenous series only (no HAR, no calendar)
\providecommand{\rhfipXgbInterHar}{}\renewcommand{\rhfipXgbInterHar}{17.6}
\providecommand{\rhfipXgbInterWeakMin}{}\renewcommand{\rhfipXgbInterWeakMin}{54}
\providecommand{\rhfipXgbInterWeakMax}{}\renewcommand{\rhfipXgbInterWeakMax}{63}
\providecommand{\rhfipXgbPairOne}{}\renewcommand{\rhfipXgbPairOne}{1.23}
\providecommand{\rhfipXgbPairTwo}{}\renewcommand{\rhfipXgbPairTwo}{0.83}
\providecommand{\rhfipXgbCellHarHar}{}\renewcommand{\rhfipXgbCellHarHar}{47.4}
\providecommand{\rhfipPairOneA}{}\renewcommand{\rhfipPairOneA}{HAR components}
\providecommand{\rhfipPairOneB}{}\renewcommand{\rhfipPairOneB}{\ensuremath{\ser{sumabsret_vwstock}{\mathrm{RAV}^{\mathrm{VW}}}}}
\providecommand{\rhfipPairTwoA}{}\renewcommand{\rhfipPairTwoA}{HAR components}
\providecommand{\rhfipPairTwoB}{}\renewcommand{\rhfipPairTwoB}{\ensuremath{\ser{sumabsret_ewstock}{\mathrm{RAV}^{\mathrm{EW}}}}}
\providecommand{\rhfipLgbmPairsHar}{}\renewcommand{\rhfipLgbmPairsHar}{7}
\providecommand{\rhfipXgbPairsHar}{}\renewcommand{\rhfipXgbPairsHar}{7}
\providecommand{\rhfipWMin}{}\renewcommand{\rhfipWMin}{1{,}000}
\providecommand{\rhfipWMax}{}\renewcommand{\rhfipWMax}{3{,}500}
\providecommand{\rhfipNWindows}{}\renewcommand{\rhfipNWindows}{6}
\providecommand{\rhfipSOneMeasures}{}\renewcommand{\rhfipSOneMeasures}{105}  % measure x common-penalty fits of the window study (P1 and P2)
\providecommand{\rhfipSOneRises}{}\renewcommand{\rhfipSOneRises}{0}  % ... of which rise with W
\providecommand{\rhfipSOneFalls}{}\renewcommand{\rhfipSOneFalls}{102}
\providecommand{\rhfipSOneFlat}{}\renewcommand{\rhfipSOneFlat}{3}
\providecommand{\rhfipSOneNotSupported}{}\renewcommand{\rhfipSOneNotSupported}{14}  % common-penalty verdicts NOT SUPPORTED
\providecommand{\rhfipSOneVerdicts}{}\renewcommand{\rhfipSOneVerdicts}{14}
\providecommand{\rhfipRidgeNeffWMin}{}\renewcommand{\rhfipRidgeNeffWMin}{29.2}
\providecommand{\rhfipRidgeNeffWMain}{}\renewcommand{\rhfipRidgeNeffWMain}{29.2}
\providecommand{\rhfipRidgeNeffWMax}{}\renewcommand{\rhfipRidgeNeffWMax}{28.2}
\providecommand{\rhfipRidgeNeffSlope}{}\renewcommand{\rhfipRidgeNeffSlope}{\ensuremath{-}0.10}
\providecommand{\rhfipRidgeNeffSlopeLo}{}\renewcommand{\rhfipRidgeNeffSlopeLo}{\ensuremath{-}0.25}
\providecommand{\rhfipRidgeNeffSlopeHi}{}\renewcommand{\rhfipRidgeNeffSlopeHi}{0.12}
\providecommand{\rhfipRidgeSOneLoNeffWMin}{}\renewcommand{\rhfipRidgeSOneLoNeffWMin}{33.9}
\providecommand{\rhfipRidgeSOneLoNeffWMax}{}\renewcommand{\rhfipRidgeSOneLoNeffWMax}{28.2}
\providecommand{\rhfipRidgeSOneHiNeffWMin}{}\renewcommand{\rhfipRidgeSOneHiNeffWMin}{29.2}
\providecommand{\rhfipRidgeSOneHiNeffWMax}{}\renewcommand{\rhfipRidgeSOneHiNeffWMax}{23.5}
\providecommand{\rhfipLassoSOneVWMin}{}\renewcommand{\rhfipLassoSOneVWMin}{38.7}
\providecommand{\rhfipLassoSOneVWMax}{}\renewcommand{\rhfipLassoSOneVWMax}{31.9}
\providecommand{\rhfipLassoSOneCWMin}{}\renewcommand{\rhfipLassoSOneCWMin}{213.3}
\providecommand{\rhfipLassoSOneCWMax}{}\renewcommand{\rhfipLassoSOneCWMax}{177.5}
\providecommand{\rhfipLassoWalkVWMin}{}\renewcommand{\rhfipLassoWalkVWMin}{38.8}
\providecommand{\rhfipLassoWalkVWMax}{}\renewcommand{\rhfipLassoWalkVWMax}{41.9}
\providecommand{\rhfipLassoSOneLo}{}\renewcommand{\rhfipLassoSOneLo}{\ensuremath{10^{-3}}}
\providecommand{\rhfipRidgeSOneLo}{}\renewcommand{\rhfipRidgeSOneLo}{\ensuremath{10^{3}}}
\providecommand{\rhfipRidgeSOneHi}{}\renewcommand{\rhfipRidgeSOneHi}{\ensuremath{10^{4}}}
\providecommand{\rhfipEnetPenWMin}{}\renewcommand{\rhfipEnetPenWMin}{\ensuremath{10^{-2}}}
\providecommand{\rhfipEnetPenWNext}{}\renewcommand{\rhfipEnetPenWNext}{\ensuremath{10^{-3}}}
\providecommand{\rhfipWNext}{}\renewcommand{\rhfipWNext}{1{,}500}
\providecommand{\rhfipRidgePOneWalk}{}\renewcommand{\rhfipRidgePOneWalk}{not supported}
\providecommand{\rhfipRidgePTwoWalk}{}\renewcommand{\rhfipRidgePTwoWalk}{not supported}
\providecommand{\rhfipLassoPOneWalk}{}\renewcommand{\rhfipLassoPOneWalk}{partly}
\providecommand{\rhfipLassoPTwoWalk}{}\renewcommand{\rhfipLassoPTwoWalk}{not supported}
\providecommand{\rhfipEnetPOneWalk}{}\renewcommand{\rhfipEnetPOneWalk}{supported}
\providecommand{\rhfipEnetPTwoWalk}{}\renewcommand{\rhfipEnetPTwoWalk}{supported}
\fi
% 2026-10-26 (integration, JFEc): the ablation macros (\rhab..., generated/rehaul_ablation_numbers.tex) and the
% window x feature-count macros (\rhcx..., generated/rehaul_complexity_numbers.tex) are used below; FRONT loads both
% in the preamble of paper_rehaul.tex (the guards only serve a standalone build).
\ifdefined\rhabLgbmFullAllPct\else% AUTO-GENERATED by writeup/make_rehaul_ablations.py -- do not edit; re-run the script.
% results/ablations_2014start_w3000 (RULE.txt): column-subset arms of the final 13-bar design,
% frozen hyperparameters, every bar scored as Table 1; DM on daily (average) QLIKE.
\newcommand{\rhabEnetAddEwAllQlike}{0.1871}  % QLIKE enet add_ew all
\newcommand{\rhabEnetAddEwAllMse}{1.47}  % MSE x 1e10 enet add_ew all
\newcommand{\rhabEnetAddEwAllPct}{$-$1.1}  % QLIKE change % vs har
\newcommand{\rhabEnetAddEwAllDm}{$-$2.81}  % DM vs the reference arm
\newcommand{\rhabEnetAddEwAllP}{0.005}  % p-value
\newcommand{\rhabEnetAddEwCloseQlike}{0.1249}  % QLIKE enet add_ew c16
\newcommand{\rhabEnetAddEwCloseMse}{2.10}  % MSE x 1e10 enet add_ew c16
\newcommand{\rhabEnetAddEwClosePct}{$-$1.5}  % QLIKE change % vs har
\newcommand{\rhabEnetAddEwCloseDm}{$-$1.37}  % DM vs the reference arm
\newcommand{\rhabEnetAddEwCloseP}{0.170}  % p-value
\newcommand{\rhabEnetAddEwCloseLateQlike}{0.1120}  % QLIKE enet add_ew c16_2020
\newcommand{\rhabEnetAddEwCloseLateMse}{4.56}  % MSE x 1e10 enet add_ew c16_2020
\newcommand{\rhabEnetAddEwCloseLatePct}{$-$1.5}  % QLIKE change % vs har
\newcommand{\rhabEnetAddEwCloseLateDm}{$-$0.90}  % DM vs the reference arm
\newcommand{\rhabEnetAddEwCloseLateP}{0.369}  % p-value
\newcommand{\rhabEnetAddFomcAllQlike}{0.1847}  % QLIKE enet add_fomc all
\newcommand{\rhabEnetAddFomcAllMse}{1.13}  % MSE x 1e10 enet add_fomc all
\newcommand{\rhabEnetAddFomcAllPct}{$-$2.4}  % QLIKE change % vs har
\newcommand{\rhabEnetAddFomcAllDm}{$-$5.85}  % DM vs the reference arm
\newcommand{\rhabEnetAddFomcAllP}{\ensuremath{<}0.001}  % p-value
\newcommand{\rhabEnetAddFomcCloseQlike}{0.1274}  % QLIKE enet add_fomc c16
\newcommand{\rhabEnetAddFomcCloseMse}{1.85}  % MSE x 1e10 enet add_fomc c16
\newcommand{\rhabEnetAddFomcClosePct}{0.4}  % QLIKE change % vs har
\newcommand{\rhabEnetAddFomcCloseDm}{1.66}  % DM vs the reference arm
\newcommand{\rhabEnetAddFomcCloseP}{0.097}  % p-value
\newcommand{\rhabEnetAddFomcCloseLateQlike}{0.1136}  % QLIKE enet add_fomc c16_2020
\newcommand{\rhabEnetAddFomcCloseLateMse}{3.96}  % MSE x 1e10 enet add_fomc c16_2020
\newcommand{\rhabEnetAddFomcCloseLatePct}{0.0}  % QLIKE change % vs har
\newcommand{\rhabEnetAddFomcCloseLateDm}{$-$0.09}  % DM vs the reference arm
\newcommand{\rhabEnetAddFomcCloseLateP}{0.927}  % p-value
\newcommand{\rhabEnetAddIvAllQlike}{0.1894}  % QLIKE enet add_iv all
\newcommand{\rhabEnetAddIvAllMse}{1.32}  % MSE x 1e10 enet add_iv all
\newcommand{\rhabEnetAddIvAllPct}{0.1}  % QLIKE change % vs har
\newcommand{\rhabEnetAddIvAllDm}{0.23}  % DM vs the reference arm
\newcommand{\rhabEnetAddIvAllP}{0.817}  % p-value
\newcommand{\rhabEnetAddIvCloseQlike}{0.1272}  % QLIKE enet add_iv c16
\newcommand{\rhabEnetAddIvCloseMse}{1.76}  % MSE x 1e10 enet add_iv c16
\newcommand{\rhabEnetAddIvClosePct}{0.3}  % QLIKE change % vs har
\newcommand{\rhabEnetAddIvCloseDm}{0.28}  % DM vs the reference arm
\newcommand{\rhabEnetAddIvCloseP}{0.780}  % p-value
\newcommand{\rhabEnetAddIvCloseLateQlike}{0.1152}  % QLIKE enet add_iv c16_2020
\newcommand{\rhabEnetAddIvCloseLateMse}{3.75}  % MSE x 1e10 enet add_iv c16_2020
\newcommand{\rhabEnetAddIvCloseLatePct}{1.4}  % QLIKE change % vs har
\newcommand{\rhabEnetAddIvCloseLateDm}{0.70}  % DM vs the reference arm
\newcommand{\rhabEnetAddIvCloseLateP}{0.485}  % p-value
\newcommand{\rhabEnetAddLiqAllQlike}{0.1881}  % QLIKE enet add_liq all
\newcommand{\rhabEnetAddLiqAllMse}{1.16}  % MSE x 1e10 enet add_liq all
\newcommand{\rhabEnetAddLiqAllPct}{$-$0.6}  % QLIKE change % vs har
\newcommand{\rhabEnetAddLiqAllDm}{$-$1.54}  % DM vs the reference arm
\newcommand{\rhabEnetAddLiqAllP}{0.124}  % p-value
\newcommand{\rhabEnetAddLiqCloseQlike}{0.1299}  % QLIKE enet add_liq c16
\newcommand{\rhabEnetAddLiqCloseMse}{1.99}  % MSE x 1e10 enet add_liq c16
\newcommand{\rhabEnetAddLiqClosePct}{2.4}  % QLIKE change % vs har
\newcommand{\rhabEnetAddLiqCloseDm}{1.77}  % DM vs the reference arm
\newcommand{\rhabEnetAddLiqCloseP}{0.077}  % p-value
\newcommand{\rhabEnetAddLiqCloseLateQlike}{0.1192}  % QLIKE enet add_liq c16_2020
\newcommand{\rhabEnetAddLiqCloseLateMse}{4.29}  % MSE x 1e10 enet add_liq c16_2020
\newcommand{\rhabEnetAddLiqCloseLatePct}{4.9}  % QLIKE change % vs har
\newcommand{\rhabEnetAddLiqCloseLateDm}{1.90}  % DM vs the reference arm
\newcommand{\rhabEnetAddLiqCloseLateP}{0.058}  % p-value
\newcommand{\rhabEnetAddRmAllQlike}{0.1842}  % QLIKE enet add_rm all
\newcommand{\rhabEnetAddRmAllMse}{1.72}  % MSE x 1e10 enet add_rm all
\newcommand{\rhabEnetAddRmAllPct}{$-$2.7}  % QLIKE change % vs har
\newcommand{\rhabEnetAddRmAllDm}{$-$4.25}  % DM vs the reference arm
\newcommand{\rhabEnetAddRmAllP}{\ensuremath{<}0.001}  % p-value
\newcommand{\rhabEnetAddRmCloseQlike}{0.1223}  % QLIKE enet add_rm c16
\newcommand{\rhabEnetAddRmCloseMse}{2.92}  % MSE x 1e10 enet add_rm c16
\newcommand{\rhabEnetAddRmClosePct}{$-$3.6}  % QLIKE change % vs har
\newcommand{\rhabEnetAddRmCloseDm}{$-$2.25}  % DM vs the reference arm
\newcommand{\rhabEnetAddRmCloseP}{0.024}  % p-value
\newcommand{\rhabEnetAddRmCloseLateQlike}{0.1147}  % QLIKE enet add_rm c16_2020
\newcommand{\rhabEnetAddRmCloseLateMse}{6.56}  % MSE x 1e10 enet add_rm c16_2020
\newcommand{\rhabEnetAddRmCloseLatePct}{1.0}  % QLIKE change % vs har
\newcommand{\rhabEnetAddRmCloseLateDm}{0.33}  % DM vs the reference arm
\newcommand{\rhabEnetAddRmCloseLateP}{0.741}  % p-value
\newcommand{\rhabEnetAddSentAllQlike}{0.1893}  % QLIKE enet add_sent all
\newcommand{\rhabEnetAddSentAllMse}{1.11}  % MSE x 1e10 enet add_sent all
\newcommand{\rhabEnetAddSentAllPct}{0.0}  % QLIKE change % vs har
\newcommand{\rhabEnetAddSentAllDm}{0.04}  % DM vs the reference arm
\newcommand{\rhabEnetAddSentAllP}{0.971}  % p-value
\newcommand{\rhabEnetAddSentCloseQlike}{0.1300}  % QLIKE enet add_sent c16
\newcommand{\rhabEnetAddSentCloseMse}{1.98}  % MSE x 1e10 enet add_sent c16
\newcommand{\rhabEnetAddSentClosePct}{2.5}  % QLIKE change % vs har
\newcommand{\rhabEnetAddSentCloseDm}{3.62}  % DM vs the reference arm
\newcommand{\rhabEnetAddSentCloseP}{\ensuremath{<}0.001}  % p-value
\newcommand{\rhabEnetAddSentCloseLateQlike}{0.1194}  % QLIKE enet add_sent c16_2020
\newcommand{\rhabEnetAddSentCloseLateMse}{4.27}  % MSE x 1e10 enet add_sent c16_2020
\newcommand{\rhabEnetAddSentCloseLatePct}{5.0}  % QLIKE change % vs har
\newcommand{\rhabEnetAddSentCloseLateDm}{2.88}  % DM vs the reference arm
\newcommand{\rhabEnetAddSentCloseLateP}{0.004}  % p-value
\newcommand{\rhabEnetAddVdAllQlike}{0.1900}  % QLIKE enet add_vd all
\newcommand{\rhabEnetAddVdAllMse}{1.11}  % MSE x 1e10 enet add_vd all
\newcommand{\rhabEnetAddVdAllPct}{0.4}  % QLIKE change % vs har
\newcommand{\rhabEnetAddVdAllDm}{2.92}  % DM vs the reference arm
\newcommand{\rhabEnetAddVdAllP}{0.003}  % p-value
\newcommand{\rhabEnetAddVdCloseQlike}{0.1277}  % QLIKE enet add_vd c16
\newcommand{\rhabEnetAddVdCloseMse}{1.96}  % MSE x 1e10 enet add_vd c16
\newcommand{\rhabEnetAddVdClosePct}{0.6}  % QLIKE change % vs har
\newcommand{\rhabEnetAddVdCloseDm}{1.85}  % DM vs the reference arm
\newcommand{\rhabEnetAddVdCloseP}{0.065}  % p-value
\newcommand{\rhabEnetAddVdCloseLateQlike}{0.1160}  % QLIKE enet add_vd c16_2020
\newcommand{\rhabEnetAddVdCloseLateMse}{4.21}  % MSE x 1e10 enet add_vd c16_2020
\newcommand{\rhabEnetAddVdCloseLatePct}{2.1}  % QLIKE change % vs har
\newcommand{\rhabEnetAddVdCloseLateDm}{2.57}  % DM vs the reference arm
\newcommand{\rhabEnetAddVdCloseLateP}{0.010}  % p-value
\newcommand{\rhabEnetAddVwAllQlike}{0.1872}  % QLIKE enet add_vw all
\newcommand{\rhabEnetAddVwAllMse}{1.99}  % MSE x 1e10 enet add_vw all
\newcommand{\rhabEnetAddVwAllPct}{$-$1.1}  % QLIKE change % vs har
\newcommand{\rhabEnetAddVwAllDm}{$-$2.03}  % DM vs the reference arm
\newcommand{\rhabEnetAddVwAllP}{0.042}  % p-value
\newcommand{\rhabEnetAddVwCloseQlike}{0.1256}  % QLIKE enet add_vw c16
\newcommand{\rhabEnetAddVwCloseMse}{2.37}  % MSE x 1e10 enet add_vw c16
\newcommand{\rhabEnetAddVwClosePct}{$-$1.0}  % QLIKE change % vs har
\newcommand{\rhabEnetAddVwCloseDm}{$-$0.75}  % DM vs the reference arm
\newcommand{\rhabEnetAddVwCloseP}{0.454}  % p-value
\newcommand{\rhabEnetAddVwCloseLateQlike}{0.1134}  % QLIKE enet add_vw c16_2020
\newcommand{\rhabEnetAddVwCloseLateMse}{5.22}  % MSE x 1e10 enet add_vw c16_2020
\newcommand{\rhabEnetAddVwCloseLatePct}{$-$0.2}  % QLIKE change % vs har
\newcommand{\rhabEnetAddVwCloseLateDm}{$-$0.13}  % DM vs the reference arm
\newcommand{\rhabEnetAddVwCloseLateP}{0.897}  % p-value
\newcommand{\rhabEnetFullAllQlike}{0.1778}  % QLIKE enet full all
\newcommand{\rhabEnetFullAllMse}{2.43}  % MSE x 1e10 enet full all
\newcommand{\rhabEnetFullAllPct}{$-$6.0}  % QLIKE change % vs har
\newcommand{\rhabEnetFullAllDm}{$-$6.98}  % DM vs the reference arm
\newcommand{\rhabEnetFullAllP}{\ensuremath{<}0.001}  % p-value
\newcommand{\rhabEnetFullCloseQlike}{0.1262}  % QLIKE enet full c16
\newcommand{\rhabEnetFullCloseMse}{3.94}  % MSE x 1e10 enet full c16
\newcommand{\rhabEnetFullClosePct}{$-$0.5}  % QLIKE change % vs har
\newcommand{\rhabEnetFullCloseDm}{$-$0.21}  % DM vs the reference arm
\newcommand{\rhabEnetFullCloseP}{0.832}  % p-value
\newcommand{\rhabEnetFullCloseLateQlike}{0.1234}  % QLIKE enet full c16_2020
\newcommand{\rhabEnetFullCloseLateMse}{9.01}  % MSE x 1e10 enet full c16_2020
\newcommand{\rhabEnetFullCloseLatePct}{8.5}  % QLIKE change % vs har
\newcommand{\rhabEnetFullCloseLateDm}{1.91}  % DM vs the reference arm
\newcommand{\rhabEnetFullCloseLateP}{0.056}  % p-value
\newcommand{\rhabEnetHarAllQlike}{0.1893}  % QLIKE enet har all
\newcommand{\rhabEnetHarAllMse}{1.11}  % MSE x 1e10 enet har all
\newcommand{\rhabEnetHarCloseQlike}{0.1269}  % QLIKE enet har c16
\newcommand{\rhabEnetHarCloseMse}{1.97}  % MSE x 1e10 enet har c16
\newcommand{\rhabEnetHarCloseLateQlike}{0.1137}  % QLIKE enet har c16_2020
\newcommand{\rhabEnetHarCloseLateMse}{4.24}  % MSE x 1e10 enet har c16_2020
\newcommand{\rhabEnetHarqAllQlike}{0.1893}  % QLIKE enet harq all
\newcommand{\rhabEnetHarqAllMse}{1.16}  % MSE x 1e10 enet harq all
\newcommand{\rhabEnetHarqAllPct}{0.0}  % QLIKE change % vs har
\newcommand{\rhabEnetHarqAllDm}{0.63}  % DM vs the reference arm
\newcommand{\rhabEnetHarqAllP}{0.531}  % p-value
\newcommand{\rhabEnetHarqCloseQlike}{0.1267}  % QLIKE enet harq c16
\newcommand{\rhabEnetHarqCloseMse}{1.80}  % MSE x 1e10 enet harq c16
\newcommand{\rhabEnetHarqClosePct}{$-$0.1}  % QLIKE change % vs har
\newcommand{\rhabEnetHarqCloseDm}{$-$1.31}  % DM vs the reference arm
\newcommand{\rhabEnetHarqCloseP}{0.191}  % p-value
\newcommand{\rhabEnetHarqCloseLateQlike}{0.1134}  % QLIKE enet harq c16_2020
\newcommand{\rhabEnetHarqCloseLateMse}{3.84}  % MSE x 1e10 enet harq c16_2020
\newcommand{\rhabEnetHarqCloseLatePct}{$-$0.3}  % QLIKE change % vs har
\newcommand{\rhabEnetHarqCloseLateDm}{$-$1.15}  % DM vs the reference arm
\newcommand{\rhabEnetHarqCloseLateP}{0.251}  % p-value
\newcommand{\rhabEnetLharAllQlike}{0.1841}  % QLIKE enet lhar all
\newcommand{\rhabEnetLharAllMse}{1.71}  % MSE x 1e10 enet lhar all
\newcommand{\rhabEnetLharAllPct}{$-$2.7}  % QLIKE change % vs har
\newcommand{\rhabEnetLharAllDm}{$-$6.66}  % DM vs the reference arm
\newcommand{\rhabEnetLharAllP}{\ensuremath{<}0.001}  % p-value
\newcommand{\rhabEnetLharCloseQlike}{0.1217}  % QLIKE enet lhar c16
\newcommand{\rhabEnetLharCloseMse}{2.76}  % MSE x 1e10 enet lhar c16
\newcommand{\rhabEnetLharClosePct}{$-$4.1}  % QLIKE change % vs har
\newcommand{\rhabEnetLharCloseDm}{$-$3.16}  % DM vs the reference arm
\newcommand{\rhabEnetLharCloseP}{0.002}  % p-value
\newcommand{\rhabEnetLharCloseLateQlike}{0.1124}  % QLIKE enet lhar c16_2020
\newcommand{\rhabEnetLharCloseLateMse}{6.19}  % MSE x 1e10 enet lhar c16_2020
\newcommand{\rhabEnetLharCloseLatePct}{$-$1.1}  % QLIKE change % vs har
\newcommand{\rhabEnetLharCloseLateDm}{$-$0.45}  % DM vs the reference arm
\newcommand{\rhabEnetLharCloseLateP}{0.656}  % p-value
\newcommand{\rhabEnetLoEwAllQlike}{0.1776}  % QLIKE enet lo_ew all
\newcommand{\rhabEnetLoEwAllMse}{2.54}  % MSE x 1e10 enet lo_ew all
\newcommand{\rhabEnetLoEwAllPct}{$-$0.1}  % QLIKE change % vs full
\newcommand{\rhabEnetLoEwAllDm}{$-$1.41}  % DM vs the reference arm
\newcommand{\rhabEnetLoEwAllP}{0.157}  % p-value
\newcommand{\rhabEnetLoEwCloseQlike}{0.1267}  % QLIKE enet lo_ew c16
\newcommand{\rhabEnetLoEwCloseMse}{4.02}  % MSE x 1e10 enet lo_ew c16
\newcommand{\rhabEnetLoEwClosePct}{0.3}  % QLIKE change % vs full
\newcommand{\rhabEnetLoEwCloseDm}{1.18}  % DM vs the reference arm
\newcommand{\rhabEnetLoEwCloseP}{0.238}  % p-value
\newcommand{\rhabEnetLoEwCloseLateQlike}{0.1235}  % QLIKE enet lo_ew c16_2020
\newcommand{\rhabEnetLoEwCloseLateMse}{9.21}  % MSE x 1e10 enet lo_ew c16_2020
\newcommand{\rhabEnetLoEwCloseLatePct}{0.1}  % QLIKE change % vs full
\newcommand{\rhabEnetLoEwCloseLateDm}{0.23}  % DM vs the reference arm
\newcommand{\rhabEnetLoEwCloseLateP}{0.822}  % p-value
\newcommand{\rhabEnetLoFomcAllQlike}{0.1820}  % QLIKE enet lo_fomc all
\newcommand{\rhabEnetLoFomcAllMse}{2.42}  % MSE x 1e10 enet lo_fomc all
\newcommand{\rhabEnetLoFomcAllPct}{2.3}  % QLIKE change % vs full
\newcommand{\rhabEnetLoFomcAllDm}{5.74}  % DM vs the reference arm
\newcommand{\rhabEnetLoFomcAllP}{\ensuremath{<}0.001}  % p-value
\newcommand{\rhabEnetLoFomcCloseQlike}{0.1261}  % QLIKE enet lo_fomc c16
\newcommand{\rhabEnetLoFomcCloseMse}{3.99}  % MSE x 1e10 enet lo_fomc c16
\newcommand{\rhabEnetLoFomcClosePct}{$-$0.1}  % QLIKE change % vs full
\newcommand{\rhabEnetLoFomcCloseDm}{$-$0.51}  % DM vs the reference arm
\newcommand{\rhabEnetLoFomcCloseP}{0.607}  % p-value
\newcommand{\rhabEnetLoFomcCloseLateQlike}{0.1234}  % QLIKE enet lo_fomc c16_2020
\newcommand{\rhabEnetLoFomcCloseLateMse}{9.13}  % MSE x 1e10 enet lo_fomc c16_2020
\newcommand{\rhabEnetLoFomcCloseLatePct}{0.1}  % QLIKE change % vs full
\newcommand{\rhabEnetLoFomcCloseLateDm}{0.13}  % DM vs the reference arm
\newcommand{\rhabEnetLoFomcCloseLateP}{0.894}  % p-value
\newcommand{\rhabEnetLoIvAllQlike}{0.1778}  % QLIKE enet lo_iv all
\newcommand{\rhabEnetLoIvAllMse}{2.49}  % MSE x 1e10 enet lo_iv all
\newcommand{\rhabEnetLoIvAllPct}{0.0}  % QLIKE change % vs full
\newcommand{\rhabEnetLoIvAllDm}{$-$0.13}  % DM vs the reference arm
\newcommand{\rhabEnetLoIvAllP}{0.898}  % p-value
\newcommand{\rhabEnetLoIvCloseQlike}{0.1243}  % QLIKE enet lo_iv c16
\newcommand{\rhabEnetLoIvCloseMse}{4.03}  % MSE x 1e10 enet lo_iv c16
\newcommand{\rhabEnetLoIvClosePct}{$-$1.5}  % QLIKE change % vs full
\newcommand{\rhabEnetLoIvCloseDm}{$-$1.51}  % DM vs the reference arm
\newcommand{\rhabEnetLoIvCloseP}{0.132}  % p-value
\newcommand{\rhabEnetLoIvCloseLateQlike}{0.1193}  % QLIKE enet lo_iv c16_2020
\newcommand{\rhabEnetLoIvCloseLateMse}{9.21}  % MSE x 1e10 enet lo_iv c16_2020
\newcommand{\rhabEnetLoIvCloseLatePct}{$-$3.3}  % QLIKE change % vs full
\newcommand{\rhabEnetLoIvCloseLateDm}{$-$1.72}  % DM vs the reference arm
\newcommand{\rhabEnetLoIvCloseLateP}{0.086}  % p-value
\newcommand{\rhabEnetLoLiqAllQlike}{0.1779}  % QLIKE enet lo_liq all
\newcommand{\rhabEnetLoLiqAllMse}{2.43}  % MSE x 1e10 enet lo_liq all
\newcommand{\rhabEnetLoLiqAllPct}{0.1}  % QLIKE change % vs full
\newcommand{\rhabEnetLoLiqAllDm}{0.35}  % DM vs the reference arm
\newcommand{\rhabEnetLoLiqAllP}{0.724}  % p-value
\newcommand{\rhabEnetLoLiqCloseQlike}{0.1251}  % QLIKE enet lo_liq c16
\newcommand{\rhabEnetLoLiqCloseMse}{4.29}  % MSE x 1e10 enet lo_liq c16
\newcommand{\rhabEnetLoLiqClosePct}{$-$0.9}  % QLIKE change % vs full
\newcommand{\rhabEnetLoLiqCloseDm}{$-$1.54}  % DM vs the reference arm
\newcommand{\rhabEnetLoLiqCloseP}{0.123}  % p-value
\newcommand{\rhabEnetLoLiqCloseLateQlike}{0.1223}  % QLIKE enet lo_liq c16_2020
\newcommand{\rhabEnetLoLiqCloseLateMse}{9.84}  % MSE x 1e10 enet lo_liq c16_2020
\newcommand{\rhabEnetLoLiqCloseLatePct}{$-$0.9}  % QLIKE change % vs full
\newcommand{\rhabEnetLoLiqCloseLateDm}{$-$0.92}  % DM vs the reference arm
\newcommand{\rhabEnetLoLiqCloseLateP}{0.358}  % p-value
\newcommand{\rhabEnetLoRmAllQlike}{0.1813}  % QLIKE enet lo_rm all
\newcommand{\rhabEnetLoRmAllMse}{2.00}  % MSE x 1e10 enet lo_rm all
\newcommand{\rhabEnetLoRmAllPct}{1.9}  % QLIKE change % vs full
\newcommand{\rhabEnetLoRmAllDm}{4.38}  % DM vs the reference arm
\newcommand{\rhabEnetLoRmAllP}{\ensuremath{<}0.001}  % p-value
\newcommand{\rhabEnetLoRmCloseQlike}{0.1285}  % QLIKE enet lo_rm c16
\newcommand{\rhabEnetLoRmCloseMse}{2.35}  % MSE x 1e10 enet lo_rm c16
\newcommand{\rhabEnetLoRmClosePct}{1.8}  % QLIKE change % vs full
\newcommand{\rhabEnetLoRmCloseDm}{1.31}  % DM vs the reference arm
\newcommand{\rhabEnetLoRmCloseP}{0.190}  % p-value
\newcommand{\rhabEnetLoRmCloseLateQlike}{0.1217}  % QLIKE enet lo_rm c16_2020
\newcommand{\rhabEnetLoRmCloseLateMse}{5.18}  % MSE x 1e10 enet lo_rm c16_2020
\newcommand{\rhabEnetLoRmCloseLatePct}{$-$1.3}  % QLIKE change % vs full
\newcommand{\rhabEnetLoRmCloseLateDm}{$-$0.62}  % DM vs the reference arm
\newcommand{\rhabEnetLoRmCloseLateP}{0.538}  % p-value
\newcommand{\rhabEnetLoSentAllQlike}{0.1782}  % QLIKE enet lo_sent all
\newcommand{\rhabEnetLoSentAllMse}{2.43}  % MSE x 1e10 enet lo_sent all
\newcommand{\rhabEnetLoSentAllPct}{0.2}  % QLIKE change % vs full
\newcommand{\rhabEnetLoSentAllDm}{1.49}  % DM vs the reference arm
\newcommand{\rhabEnetLoSentAllP}{0.136}  % p-value
\newcommand{\rhabEnetLoSentCloseQlike}{0.1253}  % QLIKE enet lo_sent c16
\newcommand{\rhabEnetLoSentCloseMse}{3.94}  % MSE x 1e10 enet lo_sent c16
\newcommand{\rhabEnetLoSentClosePct}{$-$0.7}  % QLIKE change % vs full
\newcommand{\rhabEnetLoSentCloseDm}{$-$2.12}  % DM vs the reference arm
\newcommand{\rhabEnetLoSentCloseP}{0.034}  % p-value
\newcommand{\rhabEnetLoSentCloseLateQlike}{0.1212}  % QLIKE enet lo_sent c16_2020
\newcommand{\rhabEnetLoSentCloseLateMse}{9.01}  % MSE x 1e10 enet lo_sent c16_2020
\newcommand{\rhabEnetLoSentCloseLatePct}{$-$1.7}  % QLIKE change % vs full
\newcommand{\rhabEnetLoSentCloseLateDm}{$-$2.21}  % DM vs the reference arm
\newcommand{\rhabEnetLoSentCloseLateP}{0.027}  % p-value
\newcommand{\rhabEnetLoVdAllQlike}{0.1777}  % QLIKE enet lo_vd all
\newcommand{\rhabEnetLoVdAllMse}{2.43}  % MSE x 1e10 enet lo_vd all
\newcommand{\rhabEnetLoVdAllPct}{$-$0.1}  % QLIKE change % vs full
\newcommand{\rhabEnetLoVdAllDm}{$-$1.13}  % DM vs the reference arm
\newcommand{\rhabEnetLoVdAllP}{0.260}  % p-value
\newcommand{\rhabEnetLoVdCloseQlike}{0.1252}  % QLIKE enet lo_vd c16
\newcommand{\rhabEnetLoVdCloseMse}{3.94}  % MSE x 1e10 enet lo_vd c16
\newcommand{\rhabEnetLoVdClosePct}{$-$0.8}  % QLIKE change % vs full
\newcommand{\rhabEnetLoVdCloseDm}{$-$1.94}  % DM vs the reference arm
\newcommand{\rhabEnetLoVdCloseP}{0.053}  % p-value
\newcommand{\rhabEnetLoVdCloseLateQlike}{0.1207}  % QLIKE enet lo_vd c16_2020
\newcommand{\rhabEnetLoVdCloseLateMse}{9.01}  % MSE x 1e10 enet lo_vd c16_2020
\newcommand{\rhabEnetLoVdCloseLatePct}{$-$2.1}  % QLIKE change % vs full
\newcommand{\rhabEnetLoVdCloseLateDm}{$-$2.18}  % DM vs the reference arm
\newcommand{\rhabEnetLoVdCloseLateP}{0.029}  % p-value
\newcommand{\rhabEnetLoVwAllQlike}{0.1777}  % QLIKE enet lo_vw all
\newcommand{\rhabEnetLoVwAllMse}{2.08}  % MSE x 1e10 enet lo_vw all
\newcommand{\rhabEnetLoVwAllPct}{$-$0.1}  % QLIKE change % vs full
\newcommand{\rhabEnetLoVwAllDm}{$-$0.65}  % DM vs the reference arm
\newcommand{\rhabEnetLoVwAllP}{0.516}  % p-value
\newcommand{\rhabEnetLoVwCloseQlike}{0.1254}  % QLIKE enet lo_vw c16
\newcommand{\rhabEnetLoVwCloseMse}{3.77}  % MSE x 1e10 enet lo_vw c16
\newcommand{\rhabEnetLoVwClosePct}{$-$0.7}  % QLIKE change % vs full
\newcommand{\rhabEnetLoVwCloseDm}{$-$1.66}  % DM vs the reference arm
\newcommand{\rhabEnetLoVwCloseP}{0.097}  % p-value
\newcommand{\rhabEnetLoVwCloseLateQlike}{0.1218}  % QLIKE enet lo_vw c16_2020
\newcommand{\rhabEnetLoVwCloseLateMse}{8.59}  % MSE x 1e10 enet lo_vw c16_2020
\newcommand{\rhabEnetLoVwCloseLatePct}{$-$1.3}  % QLIKE change % vs full
\newcommand{\rhabEnetLoVwCloseLateDm}{$-$2.65}  % DM vs the reference arm
\newcommand{\rhabEnetLoVwCloseLateP}{0.008}  % p-value
\newcommand{\rhabEnetSharAllQlike}{0.1877}  % QLIKE enet shar all
\newcommand{\rhabEnetSharAllMse}{1.38}  % MSE x 1e10 enet shar all
\newcommand{\rhabEnetSharAllPct}{$-$0.9}  % QLIKE change % vs har
\newcommand{\rhabEnetSharAllDm}{$-$3.68}  % DM vs the reference arm
\newcommand{\rhabEnetSharAllP}{\ensuremath{<}0.001}  % p-value
\newcommand{\rhabEnetSharCloseQlike}{0.1265}  % QLIKE enet shar c16
\newcommand{\rhabEnetSharCloseMse}{1.83}  % MSE x 1e10 enet shar c16
\newcommand{\rhabEnetSharClosePct}{$-$0.3}  % QLIKE change % vs har
\newcommand{\rhabEnetSharCloseDm}{$-$0.54}  % DM vs the reference arm
\newcommand{\rhabEnetSharCloseP}{0.586}  % p-value
\newcommand{\rhabEnetSharCloseLateQlike}{0.1141}  % QLIKE enet shar c16_2020
\newcommand{\rhabEnetSharCloseLateMse}{3.92}  % MSE x 1e10 enet shar c16_2020
\newcommand{\rhabEnetSharCloseLatePct}{0.3}  % QLIKE change % vs har
\newcommand{\rhabEnetSharCloseLateDm}{0.34}  % DM vs the reference arm
\newcommand{\rhabEnetSharCloseLateP}{0.731}  % p-value
\newcommand{\rhabEnetSharqlAllQlike}{0.1840}  % QLIKE enet sharql all
\newcommand{\rhabEnetSharqlAllMse}{1.69}  % MSE x 1e10 enet sharql all
\newcommand{\rhabEnetSharqlAllPct}{$-$2.8}  % QLIKE change % vs har
\newcommand{\rhabEnetSharqlAllDm}{$-$6.95}  % DM vs the reference arm
\newcommand{\rhabEnetSharqlAllP}{\ensuremath{<}0.001}  % p-value
\newcommand{\rhabEnetSharqlCloseQlike}{0.1216}  % QLIKE enet sharql c16
\newcommand{\rhabEnetSharqlCloseMse}{2.64}  % MSE x 1e10 enet sharql c16
\newcommand{\rhabEnetSharqlClosePct}{$-$4.1}  % QLIKE change % vs har
\newcommand{\rhabEnetSharqlCloseDm}{$-$3.46}  % DM vs the reference arm
\newcommand{\rhabEnetSharqlCloseP}{\ensuremath{<}0.001}  % p-value
\newcommand{\rhabEnetSharqlCloseLateQlike}{0.1121}  % QLIKE enet sharql c16_2020
\newcommand{\rhabEnetSharqlCloseLateMse}{5.90}  % MSE x 1e10 enet sharql c16_2020
\newcommand{\rhabEnetSharqlCloseLatePct}{$-$1.4}  % QLIKE change % vs har
\newcommand{\rhabEnetSharqlCloseLateDm}{$-$0.64}  % DM vs the reference arm
\newcommand{\rhabEnetSharqlCloseLateP}{0.524}  % p-value
\newcommand{\rhabLassoAddEwAllQlike}{0.1877}  % QLIKE lasso add_ew all
\newcommand{\rhabLassoAddEwAllMse}{1.47}  % MSE x 1e10 lasso add_ew all
\newcommand{\rhabLassoAddEwAllPct}{$-$0.8}  % QLIKE change % vs har
\newcommand{\rhabLassoAddEwAllDm}{$-$2.01}  % DM vs the reference arm
\newcommand{\rhabLassoAddEwAllP}{0.045}  % p-value
\newcommand{\rhabLassoAddEwCloseQlike}{0.1257}  % QLIKE lasso add_ew c16
\newcommand{\rhabLassoAddEwCloseMse}{2.13}  % MSE x 1e10 lasso add_ew c16
\newcommand{\rhabLassoAddEwClosePct}{$-$1.0}  % QLIKE change % vs har
\newcommand{\rhabLassoAddEwCloseDm}{$-$0.81}  % DM vs the reference arm
\newcommand{\rhabLassoAddEwCloseP}{0.417}  % p-value
\newcommand{\rhabLassoAddEwCloseLateQlike}{0.1123}  % QLIKE lasso add_ew c16_2020
\newcommand{\rhabLassoAddEwCloseLateMse}{4.65}  % MSE x 1e10 lasso add_ew c16_2020
\newcommand{\rhabLassoAddEwCloseLatePct}{$-$1.2}  % QLIKE change % vs har
\newcommand{\rhabLassoAddEwCloseLateDm}{$-$0.70}  % DM vs the reference arm
\newcommand{\rhabLassoAddEwCloseLateP}{0.485}  % p-value
\newcommand{\rhabLassoAddFomcAllQlike}{0.1820}  % QLIKE lasso add_fomc all
\newcommand{\rhabLassoAddFomcAllMse}{1.13}  % MSE x 1e10 lasso add_fomc all
\newcommand{\rhabLassoAddFomcAllPct}{$-$3.8}  % QLIKE change % vs har
\newcommand{\rhabLassoAddFomcAllDm}{$-$5.76}  % DM vs the reference arm
\newcommand{\rhabLassoAddFomcAllP}{\ensuremath{<}0.001}  % p-value
\newcommand{\rhabLassoAddFomcCloseQlike}{0.1296}  % QLIKE lasso add_fomc c16
\newcommand{\rhabLassoAddFomcCloseMse}{1.83}  % MSE x 1e10 lasso add_fomc c16
\newcommand{\rhabLassoAddFomcClosePct}{2.2}  % QLIKE change % vs har
\newcommand{\rhabLassoAddFomcCloseDm}{4.33}  % DM vs the reference arm
\newcommand{\rhabLassoAddFomcCloseP}{\ensuremath{<}0.001}  % p-value
\newcommand{\rhabLassoAddFomcCloseLateQlike}{0.1135}  % QLIKE lasso add_fomc c16_2020
\newcommand{\rhabLassoAddFomcCloseLateMse}{3.92}  % MSE x 1e10 lasso add_fomc c16_2020
\newcommand{\rhabLassoAddFomcCloseLatePct}{$-$0.1}  % QLIKE change % vs har
\newcommand{\rhabLassoAddFomcCloseLateDm}{$-$0.16}  % DM vs the reference arm
\newcommand{\rhabLassoAddFomcCloseLateP}{0.875}  % p-value
\newcommand{\rhabLassoAddIvAllQlike}{0.1887}  % QLIKE lasso add_iv all
\newcommand{\rhabLassoAddIvAllMse}{1.91}  % MSE x 1e10 lasso add_iv all
\newcommand{\rhabLassoAddIvAllPct}{$-$0.3}  % QLIKE change % vs har
\newcommand{\rhabLassoAddIvAllDm}{$-$0.79}  % DM vs the reference arm
\newcommand{\rhabLassoAddIvAllP}{0.430}  % p-value
\newcommand{\rhabLassoAddIvCloseQlike}{0.1263}  % QLIKE lasso add_iv c16
\newcommand{\rhabLassoAddIvCloseMse}{1.82}  % MSE x 1e10 lasso add_iv c16
\newcommand{\rhabLassoAddIvClosePct}{$-$0.4}  % QLIKE change % vs har
\newcommand{\rhabLassoAddIvCloseDm}{$-$0.38}  % DM vs the reference arm
\newcommand{\rhabLassoAddIvCloseP}{0.705}  % p-value
\newcommand{\rhabLassoAddIvCloseLateQlike}{0.1144}  % QLIKE lasso add_iv c16_2020
\newcommand{\rhabLassoAddIvCloseLateMse}{3.92}  % MSE x 1e10 lasso add_iv c16_2020
\newcommand{\rhabLassoAddIvCloseLatePct}{0.6}  % QLIKE change % vs har
\newcommand{\rhabLassoAddIvCloseLateDm}{0.29}  % DM vs the reference arm
\newcommand{\rhabLassoAddIvCloseLateP}{0.768}  % p-value
\newcommand{\rhabLassoAddLiqAllQlike}{0.1888}  % QLIKE lasso add_liq all
\newcommand{\rhabLassoAddLiqAllMse}{1.17}  % MSE x 1e10 lasso add_liq all
\newcommand{\rhabLassoAddLiqAllPct}{$-$0.3}  % QLIKE change % vs har
\newcommand{\rhabLassoAddLiqAllDm}{$-$0.55}  % DM vs the reference arm
\newcommand{\rhabLassoAddLiqAllP}{0.583}  % p-value
\newcommand{\rhabLassoAddLiqCloseQlike}{0.1305}  % QLIKE lasso add_liq c16
\newcommand{\rhabLassoAddLiqCloseMse}{2.02}  % MSE x 1e10 lasso add_liq c16
\newcommand{\rhabLassoAddLiqClosePct}{2.9}  % QLIKE change % vs har
\newcommand{\rhabLassoAddLiqCloseDm}{1.86}  % DM vs the reference arm
\newcommand{\rhabLassoAddLiqCloseP}{0.063}  % p-value
\newcommand{\rhabLassoAddLiqCloseLateQlike}{0.1206}  % QLIKE lasso add_liq c16_2020
\newcommand{\rhabLassoAddLiqCloseLateMse}{4.37}  % MSE x 1e10 lasso add_liq c16_2020
\newcommand{\rhabLassoAddLiqCloseLatePct}{6.1}  % QLIKE change % vs har
\newcommand{\rhabLassoAddLiqCloseLateDm}{2.00}  % DM vs the reference arm
\newcommand{\rhabLassoAddLiqCloseLateP}{0.045}  % p-value
\newcommand{\rhabLassoAddRmAllQlike}{0.1852}  % QLIKE lasso add_rm all
\newcommand{\rhabLassoAddRmAllMse}{1.81}  % MSE x 1e10 lasso add_rm all
\newcommand{\rhabLassoAddRmAllPct}{$-$2.1}  % QLIKE change % vs har
\newcommand{\rhabLassoAddRmAllDm}{$-$2.64}  % DM vs the reference arm
\newcommand{\rhabLassoAddRmAllP}{0.008}  % p-value
\newcommand{\rhabLassoAddRmCloseQlike}{0.1234}  % QLIKE lasso add_rm c16
\newcommand{\rhabLassoAddRmCloseMse}{3.11}  % MSE x 1e10 lasso add_rm c16
\newcommand{\rhabLassoAddRmClosePct}{$-$2.8}  % QLIKE change % vs har
\newcommand{\rhabLassoAddRmCloseDm}{$-$1.41}  % DM vs the reference arm
\newcommand{\rhabLassoAddRmCloseP}{0.160}  % p-value
\newcommand{\rhabLassoAddRmCloseLateQlike}{0.1176}  % QLIKE lasso add_rm c16_2020
\newcommand{\rhabLassoAddRmCloseLateMse}{7.01}  % MSE x 1e10 lasso add_rm c16_2020
\newcommand{\rhabLassoAddRmCloseLatePct}{3.5}  % QLIKE change % vs har
\newcommand{\rhabLassoAddRmCloseLateDm}{0.90}  % DM vs the reference arm
\newcommand{\rhabLassoAddRmCloseLateP}{0.368}  % p-value
\newcommand{\rhabLassoAddSentAllQlike}{0.1895}  % QLIKE lasso add_sent all
\newcommand{\rhabLassoAddSentAllMse}{1.11}  % MSE x 1e10 lasso add_sent all
\newcommand{\rhabLassoAddSentAllPct}{0.1}  % QLIKE change % vs har
\newcommand{\rhabLassoAddSentAllDm}{0.60}  % DM vs the reference arm
\newcommand{\rhabLassoAddSentAllP}{0.550}  % p-value
\newcommand{\rhabLassoAddSentCloseQlike}{0.1310}  % QLIKE lasso add_sent c16
\newcommand{\rhabLassoAddSentCloseMse}{1.98}  % MSE x 1e10 lasso add_sent c16
\newcommand{\rhabLassoAddSentClosePct}{3.2}  % QLIKE change % vs har
\newcommand{\rhabLassoAddSentCloseDm}{4.34}  % DM vs the reference arm
\newcommand{\rhabLassoAddSentCloseP}{\ensuremath{<}0.001}  % p-value
\newcommand{\rhabLassoAddSentCloseLateQlike}{0.1214}  % QLIKE lasso add_sent c16_2020
\newcommand{\rhabLassoAddSentCloseLateMse}{4.27}  % MSE x 1e10 lasso add_sent c16_2020
\newcommand{\rhabLassoAddSentCloseLatePct}{6.8}  % QLIKE change % vs har
\newcommand{\rhabLassoAddSentCloseLateDm}{3.64}  % DM vs the reference arm
\newcommand{\rhabLassoAddSentCloseLateP}{\ensuremath{<}0.001}  % p-value
\newcommand{\rhabLassoAddVdAllQlike}{0.1903}  % QLIKE lasso add_vd all
\newcommand{\rhabLassoAddVdAllMse}{1.12}  % MSE x 1e10 lasso add_vd all
\newcommand{\rhabLassoAddVdAllPct}{0.6}  % QLIKE change % vs har
\newcommand{\rhabLassoAddVdAllDm}{3.05}  % DM vs the reference arm
\newcommand{\rhabLassoAddVdAllP}{0.002}  % p-value
\newcommand{\rhabLassoAddVdCloseQlike}{0.1272}  % QLIKE lasso add_vd c16
\newcommand{\rhabLassoAddVdCloseMse}{1.95}  % MSE x 1e10 lasso add_vd c16
\newcommand{\rhabLassoAddVdClosePct}{0.3}  % QLIKE change % vs har
\newcommand{\rhabLassoAddVdCloseDm}{0.65}  % DM vs the reference arm
\newcommand{\rhabLassoAddVdCloseP}{0.516}  % p-value
\newcommand{\rhabLassoAddVdCloseLateQlike}{0.1169}  % QLIKE lasso add_vd c16_2020
\newcommand{\rhabLassoAddVdCloseLateMse}{4.20}  % MSE x 1e10 lasso add_vd c16_2020
\newcommand{\rhabLassoAddVdCloseLatePct}{2.9}  % QLIKE change % vs har
\newcommand{\rhabLassoAddVdCloseLateDm}{3.03}  % DM vs the reference arm
\newcommand{\rhabLassoAddVdCloseLateP}{0.002}  % p-value
\newcommand{\rhabLassoAddVwAllQlike}{0.1881}  % QLIKE lasso add_vw all
\newcommand{\rhabLassoAddVwAllMse}{2.04}  % MSE x 1e10 lasso add_vw all
\newcommand{\rhabLassoAddVwAllPct}{$-$0.6}  % QLIKE change % vs har
\newcommand{\rhabLassoAddVwAllDm}{$-$1.02}  % DM vs the reference arm
\newcommand{\rhabLassoAddVwAllP}{0.309}  % p-value
\newcommand{\rhabLassoAddVwCloseQlike}{0.1262}  % QLIKE lasso add_vw c16
\newcommand{\rhabLassoAddVwCloseMse}{2.38}  % MSE x 1e10 lasso add_vw c16
\newcommand{\rhabLassoAddVwClosePct}{$-$0.5}  % QLIKE change % vs har
\newcommand{\rhabLassoAddVwCloseDm}{$-$0.35}  % DM vs the reference arm
\newcommand{\rhabLassoAddVwCloseP}{0.726}  % p-value
\newcommand{\rhabLassoAddVwCloseLateQlike}{0.1136}  % QLIKE lasso add_vw c16_2020
\newcommand{\rhabLassoAddVwCloseLateMse}{5.24}  % MSE x 1e10 lasso add_vw c16_2020
\newcommand{\rhabLassoAddVwCloseLatePct}{$-$0.1}  % QLIKE change % vs har
\newcommand{\rhabLassoAddVwCloseLateDm}{$-$0.04}  % DM vs the reference arm
\newcommand{\rhabLassoAddVwCloseLateP}{0.968}  % p-value
\newcommand{\rhabLassoFullAllQlike}{0.1775}  % QLIKE lasso full all
\newcommand{\rhabLassoFullAllMse}{2.37}  % MSE x 1e10 lasso full all
\newcommand{\rhabLassoFullAllPct}{$-$6.2}  % QLIKE change % vs har
\newcommand{\rhabLassoFullAllDm}{$-$4.72}  % DM vs the reference arm
\newcommand{\rhabLassoFullAllP}{\ensuremath{<}0.001}  % p-value
\newcommand{\rhabLassoFullCloseQlike}{0.1310}  % QLIKE lasso full c16
\newcommand{\rhabLassoFullCloseMse}{3.57}  % MSE x 1e10 lasso full c16
\newcommand{\rhabLassoFullClosePct}{3.3}  % QLIKE change % vs har
\newcommand{\rhabLassoFullCloseDm}{1.04}  % DM vs the reference arm
\newcommand{\rhabLassoFullCloseP}{0.300}  % p-value
\newcommand{\rhabLassoFullCloseLateQlike}{0.1301}  % QLIKE lasso full c16_2020
\newcommand{\rhabLassoFullCloseLateMse}{8.13}  % MSE x 1e10 lasso full c16_2020
\newcommand{\rhabLassoFullCloseLatePct}{14.4}  % QLIKE change % vs har
\newcommand{\rhabLassoFullCloseLateDm}{2.51}  % DM vs the reference arm
\newcommand{\rhabLassoFullCloseLateP}{0.012}  % p-value
\newcommand{\rhabLassoHarAllQlike}{0.1893}  % QLIKE lasso har all
\newcommand{\rhabLassoHarAllMse}{1.11}  % MSE x 1e10 lasso har all
\newcommand{\rhabLassoHarCloseQlike}{0.1269}  % QLIKE lasso har c16
\newcommand{\rhabLassoHarCloseMse}{1.97}  % MSE x 1e10 lasso har c16
\newcommand{\rhabLassoHarCloseLateQlike}{0.1137}  % QLIKE lasso har c16_2020
\newcommand{\rhabLassoHarCloseLateMse}{4.24}  % MSE x 1e10 lasso har c16_2020
\newcommand{\rhabLassoHarqAllQlike}{0.1893}  % QLIKE lasso harq all
\newcommand{\rhabLassoHarqAllMse}{1.17}  % MSE x 1e10 lasso harq all
\newcommand{\rhabLassoHarqAllPct}{0.0}  % QLIKE change % vs har
\newcommand{\rhabLassoHarqAllDm}{0.61}  % DM vs the reference arm
\newcommand{\rhabLassoHarqAllP}{0.544}  % p-value
\newcommand{\rhabLassoHarqCloseQlike}{0.1267}  % QLIKE lasso harq c16
\newcommand{\rhabLassoHarqCloseMse}{1.78}  % MSE x 1e10 lasso harq c16
\newcommand{\rhabLassoHarqClosePct}{$-$0.2}  % QLIKE change % vs har
\newcommand{\rhabLassoHarqCloseDm}{$-$1.37}  % DM vs the reference arm
\newcommand{\rhabLassoHarqCloseP}{0.171}  % p-value
\newcommand{\rhabLassoHarqCloseLateQlike}{0.1133}  % QLIKE lasso harq c16_2020
\newcommand{\rhabLassoHarqCloseLateMse}{3.80}  % MSE x 1e10 lasso harq c16_2020
\newcommand{\rhabLassoHarqCloseLatePct}{$-$0.3}  % QLIKE change % vs har
\newcommand{\rhabLassoHarqCloseLateDm}{$-$1.18}  % DM vs the reference arm
\newcommand{\rhabLassoHarqCloseLateP}{0.238}  % p-value
\newcommand{\rhabLassoLharAllQlike}{0.1842}  % QLIKE lasso lhar all
\newcommand{\rhabLassoLharAllMse}{1.71}  % MSE x 1e10 lasso lhar all
\newcommand{\rhabLassoLharAllPct}{$-$2.7}  % QLIKE change % vs har
\newcommand{\rhabLassoLharAllDm}{$-$6.30}  % DM vs the reference arm
\newcommand{\rhabLassoLharAllP}{\ensuremath{<}0.001}  % p-value
\newcommand{\rhabLassoLharCloseQlike}{0.1219}  % QLIKE lasso lhar c16
\newcommand{\rhabLassoLharCloseMse}{2.80}  % MSE x 1e10 lasso lhar c16
\newcommand{\rhabLassoLharClosePct}{$-$3.9}  % QLIKE change % vs har
\newcommand{\rhabLassoLharCloseDm}{$-$2.77}  % DM vs the reference arm
\newcommand{\rhabLassoLharCloseP}{0.006}  % p-value
\newcommand{\rhabLassoLharCloseLateQlike}{0.1130}  % QLIKE lasso lhar c16_2020
\newcommand{\rhabLassoLharCloseLateMse}{6.28}  % MSE x 1e10 lasso lhar c16_2020
\newcommand{\rhabLassoLharCloseLatePct}{$-$0.6}  % QLIKE change % vs har
\newcommand{\rhabLassoLharCloseLateDm}{$-$0.21}  % DM vs the reference arm
\newcommand{\rhabLassoLharCloseLateP}{0.836}  % p-value
\newcommand{\rhabLassoLoEwAllQlike}{0.1770}  % QLIKE lasso lo_ew all
\newcommand{\rhabLassoLoEwAllMse}{2.54}  % MSE x 1e10 lasso lo_ew all
\newcommand{\rhabLassoLoEwAllPct}{$-$0.3}  % QLIKE change % vs full
\newcommand{\rhabLassoLoEwAllDm}{$-$1.46}  % DM vs the reference arm
\newcommand{\rhabLassoLoEwAllP}{0.144}  % p-value
\newcommand{\rhabLassoLoEwCloseQlike}{0.1310}  % QLIKE lasso lo_ew c16
\newcommand{\rhabLassoLoEwCloseMse}{3.82}  % MSE x 1e10 lasso lo_ew c16
\newcommand{\rhabLassoLoEwClosePct}{0.0}  % QLIKE change % vs full
\newcommand{\rhabLassoLoEwCloseDm}{0.00}  % DM vs the reference arm
\newcommand{\rhabLassoLoEwCloseP}{0.999}  % p-value
\newcommand{\rhabLassoLoEwCloseLateQlike}{0.1289}  % QLIKE lasso lo_ew c16_2020
\newcommand{\rhabLassoLoEwCloseLateMse}{8.72}  % MSE x 1e10 lasso lo_ew c16_2020
\newcommand{\rhabLassoLoEwCloseLatePct}{$-$0.9}  % QLIKE change % vs full
\newcommand{\rhabLassoLoEwCloseLateDm}{$-$0.76}  % DM vs the reference arm
\newcommand{\rhabLassoLoEwCloseLateP}{0.446}  % p-value
\newcommand{\rhabLassoLoFomcAllQlike}{0.1839}  % QLIKE lasso lo_fomc all
\newcommand{\rhabLassoLoFomcAllMse}{2.39}  % MSE x 1e10 lasso lo_fomc all
\newcommand{\rhabLassoLoFomcAllPct}{3.6}  % QLIKE change % vs full
\newcommand{\rhabLassoLoFomcAllDm}{5.10}  % DM vs the reference arm
\newcommand{\rhabLassoLoFomcAllP}{\ensuremath{<}0.001}  % p-value
\newcommand{\rhabLassoLoFomcCloseQlike}{0.1293}  % QLIKE lasso lo_fomc c16
\newcommand{\rhabLassoLoFomcCloseMse}{3.62}  % MSE x 1e10 lasso lo_fomc c16
\newcommand{\rhabLassoLoFomcClosePct}{$-$1.4}  % QLIKE change % vs full
\newcommand{\rhabLassoLoFomcCloseDm}{$-$3.21}  % DM vs the reference arm
\newcommand{\rhabLassoLoFomcCloseP}{0.001}  % p-value
\newcommand{\rhabLassoLoFomcCloseLateQlike}{0.1293}  % QLIKE lasso lo_fomc c16_2020
\newcommand{\rhabLassoLoFomcCloseLateMse}{8.25}  % MSE x 1e10 lasso lo_fomc c16_2020
\newcommand{\rhabLassoLoFomcCloseLatePct}{$-$0.6}  % QLIKE change % vs full
\newcommand{\rhabLassoLoFomcCloseLateDm}{$-$0.92}  % DM vs the reference arm
\newcommand{\rhabLassoLoFomcCloseLateP}{0.358}  % p-value
\newcommand{\rhabLassoLoIvAllQlike}{0.1778}  % QLIKE lasso lo_iv all
\newcommand{\rhabLassoLoIvAllMse}{2.43}  % MSE x 1e10 lasso lo_iv all
\newcommand{\rhabLassoLoIvAllPct}{0.2}  % QLIKE change % vs full
\newcommand{\rhabLassoLoIvAllDm}{0.34}  % DM vs the reference arm
\newcommand{\rhabLassoLoIvAllP}{0.731}  % p-value
\newcommand{\rhabLassoLoIvCloseQlike}{0.1284}  % QLIKE lasso lo_iv c16
\newcommand{\rhabLassoLoIvCloseMse}{3.90}  % MSE x 1e10 lasso lo_iv c16
\newcommand{\rhabLassoLoIvClosePct}{$-$2.0}  % QLIKE change % vs full
\newcommand{\rhabLassoLoIvCloseDm}{$-$1.66}  % DM vs the reference arm
\newcommand{\rhabLassoLoIvCloseP}{0.096}  % p-value
\newcommand{\rhabLassoLoIvCloseLateQlike}{0.1251}  % QLIKE lasso lo_iv c16_2020
\newcommand{\rhabLassoLoIvCloseLateMse}{8.92}  % MSE x 1e10 lasso lo_iv c16_2020
\newcommand{\rhabLassoLoIvCloseLatePct}{$-$3.8}  % QLIKE change % vs full
\newcommand{\rhabLassoLoIvCloseLateDm}{$-$1.77}  % DM vs the reference arm
\newcommand{\rhabLassoLoIvCloseLateP}{0.076}  % p-value
\newcommand{\rhabLassoLoLiqAllQlike}{0.1774}  % QLIKE lasso lo_liq all
\newcommand{\rhabLassoLoLiqAllMse}{2.39}  % MSE x 1e10 lasso lo_liq all
\newcommand{\rhabLassoLoLiqAllPct}{$-$0.1}  % QLIKE change % vs full
\newcommand{\rhabLassoLoLiqAllDm}{$-$0.21}  % DM vs the reference arm
\newcommand{\rhabLassoLoLiqAllP}{0.831}  % p-value
\newcommand{\rhabLassoLoLiqCloseQlike}{0.1285}  % QLIKE lasso lo_liq c16
\newcommand{\rhabLassoLoLiqCloseMse}{4.03}  % MSE x 1e10 lasso lo_liq c16
\newcommand{\rhabLassoLoLiqClosePct}{$-$1.9}  % QLIKE change % vs full
\newcommand{\rhabLassoLoLiqCloseDm}{$-$2.17}  % DM vs the reference arm
\newcommand{\rhabLassoLoLiqCloseP}{0.030}  % p-value
\newcommand{\rhabLassoLoLiqCloseLateQlike}{0.1272}  % QLIKE lasso lo_liq c16_2020
\newcommand{\rhabLassoLoLiqCloseLateMse}{9.23}  % MSE x 1e10 lasso lo_liq c16_2020
\newcommand{\rhabLassoLoLiqCloseLatePct}{$-$2.2}  % QLIKE change % vs full
\newcommand{\rhabLassoLoLiqCloseLateDm}{$-$1.99}  % DM vs the reference arm
\newcommand{\rhabLassoLoLiqCloseLateP}{0.047}  % p-value
\newcommand{\rhabLassoLoRmAllQlike}{0.1798}  % QLIKE lasso lo_rm all
\newcommand{\rhabLassoLoRmAllMse}{2.07}  % MSE x 1e10 lasso lo_rm all
\newcommand{\rhabLassoLoRmAllPct}{1.3}  % QLIKE change % vs full
\newcommand{\rhabLassoLoRmAllDm}{2.07}  % DM vs the reference arm
\newcommand{\rhabLassoLoRmAllP}{0.039}  % p-value
\newcommand{\rhabLassoLoRmCloseQlike}{0.1311}  % QLIKE lasso lo_rm c16
\newcommand{\rhabLassoLoRmCloseMse}{2.27}  % MSE x 1e10 lasso lo_rm c16
\newcommand{\rhabLassoLoRmClosePct}{0.0}  % QLIKE change % vs full
\newcommand{\rhabLassoLoRmCloseDm}{0.01}  % DM vs the reference arm
\newcommand{\rhabLassoLoRmCloseP}{0.990}  % p-value
\newcommand{\rhabLassoLoRmCloseLateQlike}{0.1255}  % QLIKE lasso lo_rm c16_2020
\newcommand{\rhabLassoLoRmCloseLateMse}{4.99}  % MSE x 1e10 lasso lo_rm c16_2020
\newcommand{\rhabLassoLoRmCloseLatePct}{$-$3.5}  % QLIKE change % vs full
\newcommand{\rhabLassoLoRmCloseLateDm}{$-$1.28}  % DM vs the reference arm
\newcommand{\rhabLassoLoRmCloseLateP}{0.200}  % p-value
\newcommand{\rhabLassoLoSentAllQlike}{0.1776}  % QLIKE lasso lo_sent all
\newcommand{\rhabLassoLoSentAllMse}{2.37}  % MSE x 1e10 lasso lo_sent all
\newcommand{\rhabLassoLoSentAllPct}{0.1}  % QLIKE change % vs full
\newcommand{\rhabLassoLoSentAllDm}{0.41}  % DM vs the reference arm
\newcommand{\rhabLassoLoSentAllP}{0.679}  % p-value
\newcommand{\rhabLassoLoSentCloseQlike}{0.1294}  % QLIKE lasso lo_sent c16
\newcommand{\rhabLassoLoSentCloseMse}{3.60}  % MSE x 1e10 lasso lo_sent c16
\newcommand{\rhabLassoLoSentClosePct}{$-$1.3}  % QLIKE change % vs full
\newcommand{\rhabLassoLoSentCloseDm}{$-$2.54}  % DM vs the reference arm
\newcommand{\rhabLassoLoSentCloseP}{0.011}  % p-value
\newcommand{\rhabLassoLoSentCloseLateQlike}{0.1262}  % QLIKE lasso lo_sent c16_2020
\newcommand{\rhabLassoLoSentCloseLateMse}{8.21}  % MSE x 1e10 lasso lo_sent c16_2020
\newcommand{\rhabLassoLoSentCloseLatePct}{$-$3.0}  % QLIKE change % vs full
\newcommand{\rhabLassoLoSentCloseLateDm}{$-$2.62}  % DM vs the reference arm
\newcommand{\rhabLassoLoSentCloseLateP}{0.009}  % p-value
\newcommand{\rhabLassoLoVdAllQlike}{0.1773}  % QLIKE lasso lo_vd all
\newcommand{\rhabLassoLoVdAllMse}{2.35}  % MSE x 1e10 lasso lo_vd all
\newcommand{\rhabLassoLoVdAllPct}{$-$0.1}  % QLIKE change % vs full
\newcommand{\rhabLassoLoVdAllDm}{$-$0.63}  % DM vs the reference arm
\newcommand{\rhabLassoLoVdAllP}{0.531}  % p-value
\newcommand{\rhabLassoLoVdCloseQlike}{0.1299}  % QLIKE lasso lo_vd c16
\newcommand{\rhabLassoLoVdCloseMse}{3.56}  % MSE x 1e10 lasso lo_vd c16
\newcommand{\rhabLassoLoVdClosePct}{$-$0.9}  % QLIKE change % vs full
\newcommand{\rhabLassoLoVdCloseDm}{$-$2.10}  % DM vs the reference arm
\newcommand{\rhabLassoLoVdCloseP}{0.036}  % p-value
\newcommand{\rhabLassoLoVdCloseLateQlike}{0.1275}  % QLIKE lasso lo_vd c16_2020
\newcommand{\rhabLassoLoVdCloseLateMse}{8.12}  % MSE x 1e10 lasso lo_vd c16_2020
\newcommand{\rhabLassoLoVdCloseLatePct}{$-$2.0}  % QLIKE change % vs full
\newcommand{\rhabLassoLoVdCloseLateDm}{$-$2.12}  % DM vs the reference arm
\newcommand{\rhabLassoLoVdCloseLateP}{0.034}  % p-value
\newcommand{\rhabLassoLoVwAllQlike}{0.1769}  % QLIKE lasso lo_vw all
\newcommand{\rhabLassoLoVwAllMse}{2.07}  % MSE x 1e10 lasso lo_vw all
\newcommand{\rhabLassoLoVwAllPct}{$-$0.3}  % QLIKE change % vs full
\newcommand{\rhabLassoLoVwAllDm}{$-$1.09}  % DM vs the reference arm
\newcommand{\rhabLassoLoVwAllP}{0.274}  % p-value
\newcommand{\rhabLassoLoVwCloseQlike}{0.1302}  % QLIKE lasso lo_vw c16
\newcommand{\rhabLassoLoVwCloseMse}{3.61}  % MSE x 1e10 lasso lo_vw c16
\newcommand{\rhabLassoLoVwClosePct}{$-$0.6}  % QLIKE change % vs full
\newcommand{\rhabLassoLoVwCloseDm}{$-$1.06}  % DM vs the reference arm
\newcommand{\rhabLassoLoVwCloseP}{0.289}  % p-value
\newcommand{\rhabLassoLoVwCloseLateQlike}{0.1277}  % QLIKE lasso lo_vw c16_2020
\newcommand{\rhabLassoLoVwCloseLateMse}{8.22}  % MSE x 1e10 lasso lo_vw c16_2020
\newcommand{\rhabLassoLoVwCloseLatePct}{$-$1.8}  % QLIKE change % vs full
\newcommand{\rhabLassoLoVwCloseLateDm}{$-$1.96}  % DM vs the reference arm
\newcommand{\rhabLassoLoVwCloseLateP}{0.050}  % p-value
\newcommand{\rhabLassoSharAllQlike}{0.1877}  % QLIKE lasso shar all
\newcommand{\rhabLassoSharAllMse}{1.39}  % MSE x 1e10 lasso shar all
\newcommand{\rhabLassoSharAllPct}{$-$0.8}  % QLIKE change % vs har
\newcommand{\rhabLassoSharAllDm}{$-$3.52}  % DM vs the reference arm
\newcommand{\rhabLassoSharAllP}{\ensuremath{<}0.001}  % p-value
\newcommand{\rhabLassoSharCloseQlike}{0.1265}  % QLIKE lasso shar c16
\newcommand{\rhabLassoSharCloseMse}{1.84}  % MSE x 1e10 lasso shar c16
\newcommand{\rhabLassoSharClosePct}{$-$0.3}  % QLIKE change % vs har
\newcommand{\rhabLassoSharCloseDm}{$-$0.44}  % DM vs the reference arm
\newcommand{\rhabLassoSharCloseP}{0.657}  % p-value
\newcommand{\rhabLassoSharCloseLateQlike}{0.1142}  % QLIKE lasso shar c16_2020
\newcommand{\rhabLassoSharCloseLateMse}{3.95}  % MSE x 1e10 lasso shar c16_2020
\newcommand{\rhabLassoSharCloseLatePct}{0.4}  % QLIKE change % vs har
\newcommand{\rhabLassoSharCloseLateDm}{0.42}  % DM vs the reference arm
\newcommand{\rhabLassoSharCloseLateP}{0.673}  % p-value
\newcommand{\rhabLassoSharqlAllQlike}{0.1840}  % QLIKE lasso sharql all
\newcommand{\rhabLassoSharqlAllMse}{1.69}  % MSE x 1e10 lasso sharql all
\newcommand{\rhabLassoSharqlAllPct}{$-$2.8}  % QLIKE change % vs har
\newcommand{\rhabLassoSharqlAllDm}{$-$6.63}  % DM vs the reference arm
\newcommand{\rhabLassoSharqlAllP}{\ensuremath{<}0.001}  % p-value
\newcommand{\rhabLassoSharqlCloseQlike}{0.1218}  % QLIKE lasso sharql c16
\newcommand{\rhabLassoSharqlCloseMse}{2.68}  % MSE x 1e10 lasso sharql c16
\newcommand{\rhabLassoSharqlClosePct}{$-$4.0}  % QLIKE change % vs har
\newcommand{\rhabLassoSharqlCloseDm}{$-$3.14}  % DM vs the reference arm
\newcommand{\rhabLassoSharqlCloseP}{0.002}  % p-value
\newcommand{\rhabLassoSharqlCloseLateQlike}{0.1125}  % QLIKE lasso sharql c16_2020
\newcommand{\rhabLassoSharqlCloseLateMse}{5.98}  % MSE x 1e10 lasso sharql c16_2020
\newcommand{\rhabLassoSharqlCloseLatePct}{$-$1.0}  % QLIKE change % vs har
\newcommand{\rhabLassoSharqlCloseLateDm}{$-$0.44}  % DM vs the reference arm
\newcommand{\rhabLassoSharqlCloseLateP}{0.662}  % p-value
\newcommand{\rhabLgbmFullAllQlike}{0.1613}  % QLIKE lgbm full all
\newcommand{\rhabLgbmFullAllMse}{1.41}  % MSE x 1e10 lgbm full all
\newcommand{\rhabLgbmFullAllPct}{$-$14.0}  % QLIKE change % vs har
\newcommand{\rhabLgbmFullAllDm}{$-$10.32}  % DM vs the reference arm
\newcommand{\rhabLgbmFullAllP}{\ensuremath{<}0.001}  % p-value
\newcommand{\rhabLgbmFullCloseQlike}{0.1201}  % QLIKE lgbm full c16
\newcommand{\rhabLgbmFullCloseMse}{2.34}  % MSE x 1e10 lgbm full c16
\newcommand{\rhabLgbmFullClosePct}{$-$2.7}  % QLIKE change % vs har
\newcommand{\rhabLgbmFullCloseDm}{$-$1.81}  % DM vs the reference arm
\newcommand{\rhabLgbmFullCloseP}{0.071}  % p-value
\newcommand{\rhabLgbmFullCloseLateQlike}{0.1089}  % QLIKE lgbm full c16_2020
\newcommand{\rhabLgbmFullCloseLateMse}{5.09}  % MSE x 1e10 lgbm full c16_2020
\newcommand{\rhabLgbmFullCloseLatePct}{$-$1.1}  % QLIKE change % vs har
\newcommand{\rhabLgbmFullCloseLateDm}{$-$0.39}  % DM vs the reference arm
\newcommand{\rhabLgbmFullCloseLateP}{0.697}  % p-value
\newcommand{\rhabLgbmHarAllQlike}{0.1875}  % QLIKE lgbm har all
\newcommand{\rhabLgbmHarAllMse}{1.18}  % MSE x 1e10 lgbm har all
\newcommand{\rhabLgbmHarCloseQlike}{0.1235}  % QLIKE lgbm har c16
\newcommand{\rhabLgbmHarCloseMse}{1.77}  % MSE x 1e10 lgbm har c16
\newcommand{\rhabLgbmHarCloseLateQlike}{0.1101}  % QLIKE lgbm har c16_2020
\newcommand{\rhabLgbmHarCloseLateMse}{3.68}  % MSE x 1e10 lgbm har c16_2020
\newcommand{\rhabOlsAddRmAllQlike}{0.1858}  % QLIKE ols add_rm all
\newcommand{\rhabOlsAddRmAllMse}{1.79}  % MSE x 1e10 ols add_rm all
\newcommand{\rhabOlsAddRmAllPct}{$-$1.9}  % QLIKE change % vs har
\newcommand{\rhabOlsAddRmAllDm}{$-$2.00}  % DM vs the reference arm
\newcommand{\rhabOlsAddRmAllP}{0.045}  % p-value
\newcommand{\rhabOlsAddRmCloseQlike}{0.1252}  % QLIKE ols add_rm c16
\newcommand{\rhabOlsAddRmCloseMse}{3.02}  % MSE x 1e10 ols add_rm c16
\newcommand{\rhabOlsAddRmClosePct}{$-$1.3}  % QLIKE change % vs har
\newcommand{\rhabOlsAddRmCloseDm}{$-$0.56}  % DM vs the reference arm
\newcommand{\rhabOlsAddRmCloseP}{0.576}  % p-value
\newcommand{\rhabOlsAddRmCloseLateQlike}{0.1221}  % QLIKE ols add_rm c16_2020
\newcommand{\rhabOlsAddRmCloseLateMse}{6.80}  % MSE x 1e10 ols add_rm c16_2020
\newcommand{\rhabOlsAddRmCloseLatePct}{7.5}  % QLIKE change % vs har
\newcommand{\rhabOlsAddRmCloseLateDm}{1.43}  % DM vs the reference arm
\newcommand{\rhabOlsAddRmCloseLateP}{0.152}  % p-value
\newcommand{\rhabOlsHarAllQlike}{0.1893}  % QLIKE ols har all
\newcommand{\rhabOlsHarAllMse}{1.11}  % MSE x 1e10 ols har all
\newcommand{\rhabOlsHarCloseQlike}{0.1269}  % QLIKE ols har c16
\newcommand{\rhabOlsHarCloseMse}{1.97}  % MSE x 1e10 ols har c16
\newcommand{\rhabOlsHarCloseLateQlike}{0.1137}  % QLIKE ols har c16_2020
\newcommand{\rhabOlsHarCloseLateMse}{4.24}  % MSE x 1e10 ols har c16_2020
\newcommand{\rhabOlsHarqAllQlike}{0.1893}  % QLIKE ols harq all
\newcommand{\rhabOlsHarqAllMse}{1.17}  % MSE x 1e10 ols harq all
\newcommand{\rhabOlsHarqAllPct}{0.0}  % QLIKE change % vs har
\newcommand{\rhabOlsHarqAllDm}{0.61}  % DM vs the reference arm
\newcommand{\rhabOlsHarqAllP}{0.543}  % p-value
\newcommand{\rhabOlsHarqCloseQlike}{0.1267}  % QLIKE ols harq c16
\newcommand{\rhabOlsHarqCloseMse}{1.78}  % MSE x 1e10 ols harq c16
\newcommand{\rhabOlsHarqClosePct}{$-$0.2}  % QLIKE change % vs har
\newcommand{\rhabOlsHarqCloseDm}{$-$1.38}  % DM vs the reference arm
\newcommand{\rhabOlsHarqCloseP}{0.166}  % p-value
\newcommand{\rhabOlsHarqCloseLateQlike}{0.1133}  % QLIKE ols harq c16_2020
\newcommand{\rhabOlsHarqCloseLateMse}{3.80}  % MSE x 1e10 ols harq c16_2020
\newcommand{\rhabOlsHarqCloseLatePct}{$-$0.3}  % QLIKE change % vs har
\newcommand{\rhabOlsHarqCloseLateDm}{$-$1.19}  % DM vs the reference arm
\newcommand{\rhabOlsHarqCloseLateP}{0.235}  % p-value
\newcommand{\rhabOlsLharAllQlike}{0.1842}  % QLIKE ols lhar all
\newcommand{\rhabOlsLharAllMse}{1.71}  % MSE x 1e10 ols lhar all
\newcommand{\rhabOlsLharAllPct}{$-$2.7}  % QLIKE change % vs har
\newcommand{\rhabOlsLharAllDm}{$-$6.19}  % DM vs the reference arm
\newcommand{\rhabOlsLharAllP}{\ensuremath{<}0.001}  % p-value
\newcommand{\rhabOlsLharCloseQlike}{0.1220}  % QLIKE ols lhar c16
\newcommand{\rhabOlsLharCloseMse}{2.81}  % MSE x 1e10 ols lhar c16
\newcommand{\rhabOlsLharClosePct}{$-$3.9}  % QLIKE change % vs har
\newcommand{\rhabOlsLharCloseDm}{$-$2.66}  % DM vs the reference arm
\newcommand{\rhabOlsLharCloseP}{0.008}  % p-value
\newcommand{\rhabOlsLharCloseLateQlike}{0.1132}  % QLIKE ols lhar c16_2020
\newcommand{\rhabOlsLharCloseLateMse}{6.30}  % MSE x 1e10 ols lhar c16_2020
\newcommand{\rhabOlsLharCloseLatePct}{$-$0.4}  % QLIKE change % vs har
\newcommand{\rhabOlsLharCloseLateDm}{$-$0.15}  % DM vs the reference arm
\newcommand{\rhabOlsLharCloseLateP}{0.883}  % p-value
\newcommand{\rhabOlsSharAllQlike}{0.1877}  % QLIKE ols shar all
\newcommand{\rhabOlsSharAllMse}{1.40}  % MSE x 1e10 ols shar all
\newcommand{\rhabOlsSharAllPct}{$-$0.8}  % QLIKE change % vs har
\newcommand{\rhabOlsSharAllDm}{$-$3.48}  % DM vs the reference arm
\newcommand{\rhabOlsSharAllP}{\ensuremath{<}0.001}  % p-value
\newcommand{\rhabOlsSharCloseQlike}{0.1265}  % QLIKE ols shar c16
\newcommand{\rhabOlsSharCloseMse}{1.84}  % MSE x 1e10 ols shar c16
\newcommand{\rhabOlsSharClosePct}{$-$0.3}  % QLIKE change % vs har
\newcommand{\rhabOlsSharCloseDm}{$-$0.42}  % DM vs the reference arm
\newcommand{\rhabOlsSharCloseP}{0.674}  % p-value
\newcommand{\rhabOlsSharCloseLateQlike}{0.1142}  % QLIKE ols shar c16_2020
\newcommand{\rhabOlsSharCloseLateMse}{3.95}  % MSE x 1e10 ols shar c16_2020
\newcommand{\rhabOlsSharCloseLatePct}{0.5}  % QLIKE change % vs har
\newcommand{\rhabOlsSharCloseLateDm}{0.44}  % DM vs the reference arm
\newcommand{\rhabOlsSharCloseLateP}{0.659}  % p-value
\newcommand{\rhabOlsSharqlAllQlike}{0.1841}  % QLIKE ols sharql all
\newcommand{\rhabOlsSharqlAllMse}{1.69}  % MSE x 1e10 ols sharql all
\newcommand{\rhabOlsSharqlAllPct}{$-$2.7}  % QLIKE change % vs har
\newcommand{\rhabOlsSharqlAllDm}{$-$6.53}  % DM vs the reference arm
\newcommand{\rhabOlsSharqlAllP}{\ensuremath{<}0.001}  % p-value
\newcommand{\rhabOlsSharqlCloseQlike}{0.1218}  % QLIKE ols sharql c16
\newcommand{\rhabOlsSharqlCloseMse}{2.68}  % MSE x 1e10 ols sharql c16
\newcommand{\rhabOlsSharqlClosePct}{$-$4.0}  % QLIKE change % vs har
\newcommand{\rhabOlsSharqlCloseDm}{$-$3.05}  % DM vs the reference arm
\newcommand{\rhabOlsSharqlCloseP}{0.002}  % p-value
\newcommand{\rhabOlsSharqlCloseLateQlike}{0.1126}  % QLIKE ols sharql c16_2020
\newcommand{\rhabOlsSharqlCloseLateMse}{5.99}  % MSE x 1e10 ols sharql c16_2020
\newcommand{\rhabOlsSharqlCloseLatePct}{$-$0.9}  % QLIKE change % vs har
\newcommand{\rhabOlsSharqlCloseLateDm}{$-$0.38}  % DM vs the reference arm
\newcommand{\rhabOlsSharqlCloseLateP}{0.703}  % p-value
\newcommand{\rhabRfFullAllQlike}{0.1685}  % QLIKE rf full all
\newcommand{\rhabRfFullAllMse}{1.32}  % MSE x 1e10 rf full all
\newcommand{\rhabRfFullAllPct}{$-$11.1}  % QLIKE change % vs har
\newcommand{\rhabRfFullAllDm}{$-$9.85}  % DM vs the reference arm
\newcommand{\rhabRfFullAllP}{\ensuremath{<}0.001}  % p-value
\newcommand{\rhabRfFullCloseQlike}{0.1248}  % QLIKE rf full c16
\newcommand{\rhabRfFullCloseMse}{1.99}  % MSE x 1e10 rf full c16
\newcommand{\rhabRfFullClosePct}{$-$1.4}  % QLIKE change % vs har
\newcommand{\rhabRfFullCloseDm}{$-$0.92}  % DM vs the reference arm
\newcommand{\rhabRfFullCloseP}{0.357}  % p-value
\newcommand{\rhabRfFullCloseLateQlike}{0.1123}  % QLIKE rf full c16_2020
\newcommand{\rhabRfFullCloseLateMse}{4.26}  % MSE x 1e10 rf full c16_2020
\newcommand{\rhabRfFullCloseLatePct}{$-$0.7}  % QLIKE change % vs har
\newcommand{\rhabRfFullCloseLateDm}{$-$0.26}  % DM vs the reference arm
\newcommand{\rhabRfFullCloseLateP}{0.794}  % p-value
\newcommand{\rhabRfHarAllQlike}{0.1895}  % QLIKE rf har all
\newcommand{\rhabRfHarAllMse}{1.19}  % MSE x 1e10 rf har all
\newcommand{\rhabRfHarCloseQlike}{0.1266}  % QLIKE rf har c16
\newcommand{\rhabRfHarCloseMse}{1.63}  % MSE x 1e10 rf har c16
\newcommand{\rhabRfHarCloseLateQlike}{0.1131}  % QLIKE rf har c16_2020
\newcommand{\rhabRfHarCloseLateMse}{3.39}  % MSE x 1e10 rf har c16_2020
\newcommand{\rhabRidgeAddEwAllQlike}{0.1876}  % QLIKE ridge add_ew all
\newcommand{\rhabRidgeAddEwAllMse}{1.42}  % MSE x 1e10 ridge add_ew all
\newcommand{\rhabRidgeAddEwAllPct}{$-$0.9}  % QLIKE change % vs har
\newcommand{\rhabRidgeAddEwAllDm}{$-$2.26}  % DM vs the reference arm
\newcommand{\rhabRidgeAddEwAllP}{0.024}  % p-value
\newcommand{\rhabRidgeAddEwCloseQlike}{0.1256}  % QLIKE ridge add_ew c16
\newcommand{\rhabRidgeAddEwCloseMse}{2.08}  % MSE x 1e10 ridge add_ew c16
\newcommand{\rhabRidgeAddEwClosePct}{$-$1.0}  % QLIKE change % vs har
\newcommand{\rhabRidgeAddEwCloseDm}{$-$0.89}  % DM vs the reference arm
\newcommand{\rhabRidgeAddEwCloseP}{0.374}  % p-value
\newcommand{\rhabRidgeAddEwCloseLateQlike}{0.1125}  % QLIKE ridge add_ew c16_2020
\newcommand{\rhabRidgeAddEwCloseLateMse}{4.51}  % MSE x 1e10 ridge add_ew c16_2020
\newcommand{\rhabRidgeAddEwCloseLatePct}{$-$1.0}  % QLIKE change % vs har
\newcommand{\rhabRidgeAddEwCloseLateDm}{$-$0.62}  % DM vs the reference arm
\newcommand{\rhabRidgeAddEwCloseLateP}{0.537}  % p-value
\newcommand{\rhabRidgeAddFomcAllQlike}{0.1847}  % QLIKE ridge add_fomc all
\newcommand{\rhabRidgeAddFomcAllMse}{1.13}  % MSE x 1e10 ridge add_fomc all
\newcommand{\rhabRidgeAddFomcAllPct}{$-$2.4}  % QLIKE change % vs har
\newcommand{\rhabRidgeAddFomcAllDm}{$-$5.79}  % DM vs the reference arm
\newcommand{\rhabRidgeAddFomcAllP}{\ensuremath{<}0.001}  % p-value
\newcommand{\rhabRidgeAddFomcCloseQlike}{0.1276}  % QLIKE ridge add_fomc c16
\newcommand{\rhabRidgeAddFomcCloseMse}{1.84}  % MSE x 1e10 ridge add_fomc c16
\newcommand{\rhabRidgeAddFomcClosePct}{0.6}  % QLIKE change % vs har
\newcommand{\rhabRidgeAddFomcCloseDm}{2.17}  % DM vs the reference arm
\newcommand{\rhabRidgeAddFomcCloseP}{0.030}  % p-value
\newcommand{\rhabRidgeAddFomcCloseLateQlike}{0.1135}  % QLIKE ridge add_fomc c16_2020
\newcommand{\rhabRidgeAddFomcCloseLateMse}{3.94}  % MSE x 1e10 ridge add_fomc c16_2020
\newcommand{\rhabRidgeAddFomcCloseLatePct}{$-$0.2}  % QLIKE change % vs har
\newcommand{\rhabRidgeAddFomcCloseLateDm}{$-$0.35}  % DM vs the reference arm
\newcommand{\rhabRidgeAddFomcCloseLateP}{0.729}  % p-value
\newcommand{\rhabRidgeAddIvAllQlike}{0.1897}  % QLIKE ridge add_iv all
\newcommand{\rhabRidgeAddIvAllMse}{1.28}  % MSE x 1e10 ridge add_iv all
\newcommand{\rhabRidgeAddIvAllPct}{0.2}  % QLIKE change % vs har
\newcommand{\rhabRidgeAddIvAllDm}{0.67}  % DM vs the reference arm
\newcommand{\rhabRidgeAddIvAllP}{0.506}  % p-value
\newcommand{\rhabRidgeAddIvCloseQlike}{0.1266}  % QLIKE ridge add_iv c16
\newcommand{\rhabRidgeAddIvCloseMse}{1.79}  % MSE x 1e10 ridge add_iv c16
\newcommand{\rhabRidgeAddIvClosePct}{$-$0.2}  % QLIKE change % vs har
\newcommand{\rhabRidgeAddIvCloseDm}{$-$0.20}  % DM vs the reference arm
\newcommand{\rhabRidgeAddIvCloseP}{0.845}  % p-value
\newcommand{\rhabRidgeAddIvCloseLateQlike}{0.1153}  % QLIKE ridge add_iv c16_2020
\newcommand{\rhabRidgeAddIvCloseLateMse}{3.82}  % MSE x 1e10 ridge add_iv c16_2020
\newcommand{\rhabRidgeAddIvCloseLatePct}{1.4}  % QLIKE change % vs har
\newcommand{\rhabRidgeAddIvCloseLateDm}{0.70}  % DM vs the reference arm
\newcommand{\rhabRidgeAddIvCloseLateP}{0.482}  % p-value
\newcommand{\rhabRidgeAddLiqAllQlike}{0.1885}  % QLIKE ridge add_liq all
\newcommand{\rhabRidgeAddLiqAllMse}{1.18}  % MSE x 1e10 ridge add_liq all
\newcommand{\rhabRidgeAddLiqAllPct}{$-$0.4}  % QLIKE change % vs har
\newcommand{\rhabRidgeAddLiqAllDm}{$-$0.98}  % DM vs the reference arm
\newcommand{\rhabRidgeAddLiqAllP}{0.326}  % p-value
\newcommand{\rhabRidgeAddLiqCloseQlike}{0.1304}  % QLIKE ridge add_liq c16
\newcommand{\rhabRidgeAddLiqCloseMse}{2.01}  % MSE x 1e10 ridge add_liq c16
\newcommand{\rhabRidgeAddLiqClosePct}{2.8}  % QLIKE change % vs har
\newcommand{\rhabRidgeAddLiqCloseDm}{1.91}  % DM vs the reference arm
\newcommand{\rhabRidgeAddLiqCloseP}{0.056}  % p-value
\newcommand{\rhabRidgeAddLiqCloseLateQlike}{0.1204}  % QLIKE ridge add_liq c16_2020
\newcommand{\rhabRidgeAddLiqCloseLateMse}{4.35}  % MSE x 1e10 ridge add_liq c16_2020
\newcommand{\rhabRidgeAddLiqCloseLatePct}{6.0}  % QLIKE change % vs har
\newcommand{\rhabRidgeAddLiqCloseLateDm}{2.09}  % DM vs the reference arm
\newcommand{\rhabRidgeAddLiqCloseLateP}{0.036}  % p-value
\newcommand{\rhabRidgeAddRmAllQlike}{0.1843}  % QLIKE ridge add_rm all
\newcommand{\rhabRidgeAddRmAllMse}{1.69}  % MSE x 1e10 ridge add_rm all
\newcommand{\rhabRidgeAddRmAllPct}{$-$2.6}  % QLIKE change % vs har
\newcommand{\rhabRidgeAddRmAllDm}{$-$4.14}  % DM vs the reference arm
\newcommand{\rhabRidgeAddRmAllP}{\ensuremath{<}0.001}  % p-value
\newcommand{\rhabRidgeAddRmCloseQlike}{0.1230}  % QLIKE ridge add_rm c16
\newcommand{\rhabRidgeAddRmCloseMse}{2.81}  % MSE x 1e10 ridge add_rm c16
\newcommand{\rhabRidgeAddRmClosePct}{$-$3.1}  % QLIKE change % vs har
\newcommand{\rhabRidgeAddRmCloseDm}{$-$1.80}  % DM vs the reference arm
\newcommand{\rhabRidgeAddRmCloseP}{0.071}  % p-value
\newcommand{\rhabRidgeAddRmCloseLateQlike}{0.1156}  % QLIKE ridge add_rm c16_2020
\newcommand{\rhabRidgeAddRmCloseLateMse}{6.29}  % MSE x 1e10 ridge add_rm c16_2020
\newcommand{\rhabRidgeAddRmCloseLatePct}{1.7}  % QLIKE change % vs har
\newcommand{\rhabRidgeAddRmCloseLateDm}{0.56}  % DM vs the reference arm
\newcommand{\rhabRidgeAddRmCloseLateP}{0.577}  % p-value
\newcommand{\rhabRidgeAddSentAllQlike}{0.1897}  % QLIKE ridge add_sent all
\newcommand{\rhabRidgeAddSentAllMse}{1.11}  % MSE x 1e10 ridge add_sent all
\newcommand{\rhabRidgeAddSentAllPct}{0.2}  % QLIKE change % vs har
\newcommand{\rhabRidgeAddSentAllDm}{0.82}  % DM vs the reference arm
\newcommand{\rhabRidgeAddSentAllP}{0.411}  % p-value
\newcommand{\rhabRidgeAddSentCloseQlike}{0.1310}  % QLIKE ridge add_sent c16
\newcommand{\rhabRidgeAddSentCloseMse}{1.98}  % MSE x 1e10 ridge add_sent c16
\newcommand{\rhabRidgeAddSentClosePct}{3.2}  % QLIKE change % vs har
\newcommand{\rhabRidgeAddSentCloseDm}{4.41}  % DM vs the reference arm
\newcommand{\rhabRidgeAddSentCloseP}{\ensuremath{<}0.001}  % p-value
\newcommand{\rhabRidgeAddSentCloseLateQlike}{0.1214}  % QLIKE ridge add_sent c16_2020
\newcommand{\rhabRidgeAddSentCloseLateMse}{4.27}  % MSE x 1e10 ridge add_sent c16_2020
\newcommand{\rhabRidgeAddSentCloseLatePct}{6.8}  % QLIKE change % vs har
\newcommand{\rhabRidgeAddSentCloseLateDm}{3.65}  % DM vs the reference arm
\newcommand{\rhabRidgeAddSentCloseLateP}{\ensuremath{<}0.001}  % p-value
\newcommand{\rhabRidgeAddVdAllQlike}{0.1903}  % QLIKE ridge add_vd all
\newcommand{\rhabRidgeAddVdAllMse}{1.11}  % MSE x 1e10 ridge add_vd all
\newcommand{\rhabRidgeAddVdAllPct}{0.5}  % QLIKE change % vs har
\newcommand{\rhabRidgeAddVdAllDm}{3.10}  % DM vs the reference arm
\newcommand{\rhabRidgeAddVdAllP}{0.002}  % p-value
\newcommand{\rhabRidgeAddVdCloseQlike}{0.1276}  % QLIKE ridge add_vd c16
\newcommand{\rhabRidgeAddVdCloseMse}{1.95}  % MSE x 1e10 ridge add_vd c16
\newcommand{\rhabRidgeAddVdClosePct}{0.6}  % QLIKE change % vs har
\newcommand{\rhabRidgeAddVdCloseDm}{1.55}  % DM vs the reference arm
\newcommand{\rhabRidgeAddVdCloseP}{0.120}  % p-value
\newcommand{\rhabRidgeAddVdCloseLateQlike}{0.1169}  % QLIKE ridge add_vd c16_2020
\newcommand{\rhabRidgeAddVdCloseLateMse}{4.21}  % MSE x 1e10 ridge add_vd c16_2020
\newcommand{\rhabRidgeAddVdCloseLatePct}{2.8}  % QLIKE change % vs har
\newcommand{\rhabRidgeAddVdCloseLateDm}{3.10}  % DM vs the reference arm
\newcommand{\rhabRidgeAddVdCloseLateP}{0.002}  % p-value
\newcommand{\rhabRidgeAddVwAllQlike}{0.1876}  % QLIKE ridge add_vw all
\newcommand{\rhabRidgeAddVwAllMse}{1.87}  % MSE x 1e10 ridge add_vw all
\newcommand{\rhabRidgeAddVwAllPct}{$-$0.9}  % QLIKE change % vs har
\newcommand{\rhabRidgeAddVwAllDm}{$-$1.58}  % DM vs the reference arm
\newcommand{\rhabRidgeAddVwAllP}{0.114}  % p-value
\newcommand{\rhabRidgeAddVwCloseQlike}{0.1264}  % QLIKE ridge add_vw c16
\newcommand{\rhabRidgeAddVwCloseMse}{2.27}  % MSE x 1e10 ridge add_vw c16
\newcommand{\rhabRidgeAddVwClosePct}{$-$0.4}  % QLIKE change % vs har
\newcommand{\rhabRidgeAddVwCloseDm}{$-$0.28}  % DM vs the reference arm
\newcommand{\rhabRidgeAddVwCloseP}{0.781}  % p-value
\newcommand{\rhabRidgeAddVwCloseLateQlike}{0.1139}  % QLIKE ridge add_vw c16_2020
\newcommand{\rhabRidgeAddVwCloseLateMse}{5.00}  % MSE x 1e10 ridge add_vw c16_2020
\newcommand{\rhabRidgeAddVwCloseLatePct}{0.2}  % QLIKE change % vs har
\newcommand{\rhabRidgeAddVwCloseLateDm}{0.12}  % DM vs the reference arm
\newcommand{\rhabRidgeAddVwCloseLateP}{0.907}  % p-value
\newcommand{\rhabRidgeFullAllQlike}{0.1791}  % QLIKE ridge full all
\newcommand{\rhabRidgeFullAllMse}{2.24}  % MSE x 1e10 ridge full all
\newcommand{\rhabRidgeFullAllPct}{$-$5.4}  % QLIKE change % vs har
\newcommand{\rhabRidgeFullAllDm}{$-$5.63}  % DM vs the reference arm
\newcommand{\rhabRidgeFullAllP}{\ensuremath{<}0.001}  % p-value
\newcommand{\rhabRidgeFullCloseQlike}{0.1279}  % QLIKE ridge full c16
\newcommand{\rhabRidgeFullCloseMse}{3.52}  % MSE x 1e10 ridge full c16
\newcommand{\rhabRidgeFullClosePct}{0.8}  % QLIKE change % vs har
\newcommand{\rhabRidgeFullCloseDm}{0.31}  % DM vs the reference arm
\newcommand{\rhabRidgeFullCloseP}{0.759}  % p-value
\newcommand{\rhabRidgeFullCloseLateQlike}{0.1260}  % QLIKE ridge full c16_2020
\newcommand{\rhabRidgeFullCloseLateMse}{7.99}  % MSE x 1e10 ridge full c16_2020
\newcommand{\rhabRidgeFullCloseLatePct}{10.9}  % QLIKE change % vs har
\newcommand{\rhabRidgeFullCloseLateDm}{2.31}  % DM vs the reference arm
\newcommand{\rhabRidgeFullCloseLateP}{0.021}  % p-value
\newcommand{\rhabRidgeHarAllQlike}{0.1893}  % QLIKE ridge har all
\newcommand{\rhabRidgeHarAllMse}{1.11}  % MSE x 1e10 ridge har all
\newcommand{\rhabRidgeHarCloseQlike}{0.1269}  % QLIKE ridge har c16
\newcommand{\rhabRidgeHarCloseMse}{1.97}  % MSE x 1e10 ridge har c16
\newcommand{\rhabRidgeHarCloseLateQlike}{0.1137}  % QLIKE ridge har c16_2020
\newcommand{\rhabRidgeHarCloseLateMse}{4.24}  % MSE x 1e10 ridge har c16_2020
\newcommand{\rhabRidgeHarqAllQlike}{0.1893}  % QLIKE ridge harq all
\newcommand{\rhabRidgeHarqAllMse}{1.16}  % MSE x 1e10 ridge harq all
\newcommand{\rhabRidgeHarqAllPct}{0.0}  % QLIKE change % vs har
\newcommand{\rhabRidgeHarqAllDm}{0.76}  % DM vs the reference arm
\newcommand{\rhabRidgeHarqAllP}{0.446}  % p-value
\newcommand{\rhabRidgeHarqCloseQlike}{0.1267}  % QLIKE ridge harq c16
\newcommand{\rhabRidgeHarqCloseMse}{1.79}  % MSE x 1e10 ridge harq c16
\newcommand{\rhabRidgeHarqClosePct}{$-$0.2}  % QLIKE change % vs har
\newcommand{\rhabRidgeHarqCloseDm}{$-$1.42}  % DM vs the reference arm
\newcommand{\rhabRidgeHarqCloseP}{0.157}  % p-value
\newcommand{\rhabRidgeHarqCloseLateQlike}{0.1133}  % QLIKE ridge harq c16_2020
\newcommand{\rhabRidgeHarqCloseLateMse}{3.81}  % MSE x 1e10 ridge harq c16_2020
\newcommand{\rhabRidgeHarqCloseLatePct}{$-$0.3}  % QLIKE change % vs har
\newcommand{\rhabRidgeHarqCloseLateDm}{$-$1.20}  % DM vs the reference arm
\newcommand{\rhabRidgeHarqCloseLateP}{0.228}  % p-value
\newcommand{\rhabRidgeLharAllQlike}{0.1841}  % QLIKE ridge lhar all
\newcommand{\rhabRidgeLharAllMse}{1.72}  % MSE x 1e10 ridge lhar all
\newcommand{\rhabRidgeLharAllPct}{$-$2.7}  % QLIKE change % vs har
\newcommand{\rhabRidgeLharAllDm}{$-$6.56}  % DM vs the reference arm
\newcommand{\rhabRidgeLharAllP}{\ensuremath{<}0.001}  % p-value
\newcommand{\rhabRidgeLharCloseQlike}{0.1216}  % QLIKE ridge lhar c16
\newcommand{\rhabRidgeLharCloseMse}{2.78}  % MSE x 1e10 ridge lhar c16
\newcommand{\rhabRidgeLharClosePct}{$-$4.2}  % QLIKE change % vs har
\newcommand{\rhabRidgeLharCloseDm}{$-$3.28}  % DM vs the reference arm
\newcommand{\rhabRidgeLharCloseP}{0.001}  % p-value
\newcommand{\rhabRidgeLharCloseLateQlike}{0.1121}  % QLIKE ridge lhar c16_2020
\newcommand{\rhabRidgeLharCloseLateMse}{6.23}  % MSE x 1e10 ridge lhar c16_2020
\newcommand{\rhabRidgeLharCloseLatePct}{$-$1.4}  % QLIKE change % vs har
\newcommand{\rhabRidgeLharCloseLateDm}{$-$0.59}  % DM vs the reference arm
\newcommand{\rhabRidgeLharCloseLateP}{0.554}  % p-value
\newcommand{\rhabRidgeLoEwAllQlike}{0.1789}  % QLIKE ridge lo_ew all
\newcommand{\rhabRidgeLoEwAllMse}{2.38}  % MSE x 1e10 ridge lo_ew all
\newcommand{\rhabRidgeLoEwAllPct}{$-$0.1}  % QLIKE change % vs full
\newcommand{\rhabRidgeLoEwAllDm}{$-$1.20}  % DM vs the reference arm
\newcommand{\rhabRidgeLoEwAllP}{0.230}  % p-value
\newcommand{\rhabRidgeLoEwCloseQlike}{0.1285}  % QLIKE ridge lo_ew c16
\newcommand{\rhabRidgeLoEwCloseMse}{3.74}  % MSE x 1e10 ridge lo_ew c16
\newcommand{\rhabRidgeLoEwClosePct}{0.5}  % QLIKE change % vs full
\newcommand{\rhabRidgeLoEwCloseDm}{1.23}  % DM vs the reference arm
\newcommand{\rhabRidgeLoEwCloseP}{0.218}  % p-value
\newcommand{\rhabRidgeLoEwCloseLateQlike}{0.1262}  % QLIKE ridge lo_ew c16_2020
\newcommand{\rhabRidgeLoEwCloseLateMse}{8.53}  % MSE x 1e10 ridge lo_ew c16_2020
\newcommand{\rhabRidgeLoEwCloseLatePct}{0.2}  % QLIKE change % vs full
\newcommand{\rhabRidgeLoEwCloseLateDm}{0.22}  % DM vs the reference arm
\newcommand{\rhabRidgeLoEwCloseLateP}{0.828}  % p-value
\newcommand{\rhabRidgeLoFomcAllQlike}{0.1833}  % QLIKE ridge lo_fomc all
\newcommand{\rhabRidgeLoFomcAllMse}{2.25}  % MSE x 1e10 ridge lo_fomc all
\newcommand{\rhabRidgeLoFomcAllPct}{2.3}  % QLIKE change % vs full
\newcommand{\rhabRidgeLoFomcAllDm}{5.55}  % DM vs the reference arm
\newcommand{\rhabRidgeLoFomcAllP}{\ensuremath{<}0.001}  % p-value
\newcommand{\rhabRidgeLoFomcCloseQlike}{0.1275}  % QLIKE ridge lo_fomc c16
\newcommand{\rhabRidgeLoFomcCloseMse}{3.58}  % MSE x 1e10 ridge lo_fomc c16
\newcommand{\rhabRidgeLoFomcClosePct}{$-$0.3}  % QLIKE change % vs full
\newcommand{\rhabRidgeLoFomcCloseDm}{$-$1.23}  % DM vs the reference arm
\newcommand{\rhabRidgeLoFomcCloseP}{0.219}  % p-value
\newcommand{\rhabRidgeLoFomcCloseLateQlike}{0.1257}  % QLIKE ridge lo_fomc c16_2020
\newcommand{\rhabRidgeLoFomcCloseLateMse}{8.15}  % MSE x 1e10 ridge lo_fomc c16_2020
\newcommand{\rhabRidgeLoFomcCloseLatePct}{$-$0.2}  % QLIKE change % vs full
\newcommand{\rhabRidgeLoFomcCloseLateDm}{$-$0.49}  % DM vs the reference arm
\newcommand{\rhabRidgeLoFomcCloseLateP}{0.622}  % p-value
\newcommand{\rhabRidgeLoIvAllQlike}{0.1791}  % QLIKE ridge lo_iv all
\newcommand{\rhabRidgeLoIvAllMse}{2.30}  % MSE x 1e10 ridge lo_iv all
\newcommand{\rhabRidgeLoIvAllPct}{0.0}  % QLIKE change % vs full
\newcommand{\rhabRidgeLoIvAllDm}{0.03}  % DM vs the reference arm
\newcommand{\rhabRidgeLoIvAllP}{0.980}  % p-value
\newcommand{\rhabRidgeLoIvCloseQlike}{0.1267}  % QLIKE ridge lo_iv c16
\newcommand{\rhabRidgeLoIvCloseMse}{3.61}  % MSE x 1e10 ridge lo_iv c16
\newcommand{\rhabRidgeLoIvClosePct}{$-$1.0}  % QLIKE change % vs full
\newcommand{\rhabRidgeLoIvCloseDm}{$-$1.02}  % DM vs the reference arm
\newcommand{\rhabRidgeLoIvCloseP}{0.309}  % p-value
\newcommand{\rhabRidgeLoIvCloseLateQlike}{0.1228}  % QLIKE ridge lo_iv c16_2020
\newcommand{\rhabRidgeLoIvCloseLateMse}{8.20}  % MSE x 1e10 ridge lo_iv c16_2020
\newcommand{\rhabRidgeLoIvCloseLatePct}{$-$2.5}  % QLIKE change % vs full
\newcommand{\rhabRidgeLoIvCloseLateDm}{$-$1.41}  % DM vs the reference arm
\newcommand{\rhabRidgeLoIvCloseLateP}{0.159}  % p-value
\newcommand{\rhabRidgeLoLiqAllQlike}{0.1794}  % QLIKE ridge lo_liq all
\newcommand{\rhabRidgeLoLiqAllMse}{2.29}  % MSE x 1e10 ridge lo_liq all
\newcommand{\rhabRidgeLoLiqAllPct}{0.2}  % QLIKE change % vs full
\newcommand{\rhabRidgeLoLiqAllDm}{0.75}  % DM vs the reference arm
\newcommand{\rhabRidgeLoLiqAllP}{0.454}  % p-value
\newcommand{\rhabRidgeLoLiqCloseQlike}{0.1269}  % QLIKE ridge lo_liq c16
\newcommand{\rhabRidgeLoLiqCloseMse}{3.93}  % MSE x 1e10 ridge lo_liq c16
\newcommand{\rhabRidgeLoLiqClosePct}{$-$0.8}  % QLIKE change % vs full
\newcommand{\rhabRidgeLoLiqCloseDm}{$-$1.14}  % DM vs the reference arm
\newcommand{\rhabRidgeLoLiqCloseP}{0.256}  % p-value
\newcommand{\rhabRidgeLoLiqCloseLateQlike}{0.1251}  % QLIKE ridge lo_liq c16_2020
\newcommand{\rhabRidgeLoLiqCloseLateMse}{8.99}  % MSE x 1e10 ridge lo_liq c16_2020
\newcommand{\rhabRidgeLoLiqCloseLatePct}{$-$0.7}  % QLIKE change % vs full
\newcommand{\rhabRidgeLoLiqCloseLateDm}{$-$0.77}  % DM vs the reference arm
\newcommand{\rhabRidgeLoLiqCloseLateP}{0.442}  % p-value
\newcommand{\rhabRidgeLoRmAllQlike}{0.1822}  % QLIKE ridge lo_rm all
\newcommand{\rhabRidgeLoRmAllMse}{1.82}  % MSE x 1e10 ridge lo_rm all
\newcommand{\rhabRidgeLoRmAllPct}{1.7}  % QLIKE change % vs full
\newcommand{\rhabRidgeLoRmAllDm}{3.75}  % DM vs the reference arm
\newcommand{\rhabRidgeLoRmAllP}{\ensuremath{<}0.001}  % p-value
\newcommand{\rhabRidgeLoRmCloseQlike}{0.1292}  % QLIKE ridge lo_rm c16
\newcommand{\rhabRidgeLoRmCloseMse}{2.26}  % MSE x 1e10 ridge lo_rm c16
\newcommand{\rhabRidgeLoRmClosePct}{1.0}  % QLIKE change % vs full
\newcommand{\rhabRidgeLoRmCloseDm}{0.73}  % DM vs the reference arm
\newcommand{\rhabRidgeLoRmCloseP}{0.466}  % p-value
\newcommand{\rhabRidgeLoRmCloseLateQlike}{0.1245}  % QLIKE ridge lo_rm c16_2020
\newcommand{\rhabRidgeLoRmCloseLateMse}{4.97}  % MSE x 1e10 ridge lo_rm c16_2020
\newcommand{\rhabRidgeLoRmCloseLatePct}{$-$1.2}  % QLIKE change % vs full
\newcommand{\rhabRidgeLoRmCloseLateDm}{$-$0.55}  % DM vs the reference arm
\newcommand{\rhabRidgeLoRmCloseLateP}{0.579}  % p-value
\newcommand{\rhabRidgeLoSentAllQlike}{0.1793}  % QLIKE ridge lo_sent all
\newcommand{\rhabRidgeLoSentAllMse}{2.24}  % MSE x 1e10 ridge lo_sent all
\newcommand{\rhabRidgeLoSentAllPct}{0.1}  % QLIKE change % vs full
\newcommand{\rhabRidgeLoSentAllDm}{0.55}  % DM vs the reference arm
\newcommand{\rhabRidgeLoSentAllP}{0.584}  % p-value
\newcommand{\rhabRidgeLoSentCloseQlike}{0.1261}  % QLIKE ridge lo_sent c16
\newcommand{\rhabRidgeLoSentCloseMse}{3.51}  % MSE x 1e10 ridge lo_sent c16
\newcommand{\rhabRidgeLoSentClosePct}{$-$1.4}  % QLIKE change % vs full
\newcommand{\rhabRidgeLoSentCloseDm}{$-$3.16}  % DM vs the reference arm
\newcommand{\rhabRidgeLoSentCloseP}{0.002}  % p-value
\newcommand{\rhabRidgeLoSentCloseLateQlike}{0.1218}  % QLIKE ridge lo_sent c16_2020
\newcommand{\rhabRidgeLoSentCloseLateMse}{7.97}  % MSE x 1e10 ridge lo_sent c16_2020
\newcommand{\rhabRidgeLoSentCloseLatePct}{$-$3.3}  % QLIKE change % vs full
\newcommand{\rhabRidgeLoSentCloseLateDm}{$-$3.21}  % DM vs the reference arm
\newcommand{\rhabRidgeLoSentCloseLateP}{0.001}  % p-value
\newcommand{\rhabRidgeLoVdAllQlike}{0.1789}  % QLIKE ridge lo_vd all
\newcommand{\rhabRidgeLoVdAllMse}{2.24}  % MSE x 1e10 ridge lo_vd all
\newcommand{\rhabRidgeLoVdAllPct}{$-$0.1}  % QLIKE change % vs full
\newcommand{\rhabRidgeLoVdAllDm}{$-$0.98}  % DM vs the reference arm
\newcommand{\rhabRidgeLoVdAllP}{0.326}  % p-value
\newcommand{\rhabRidgeLoVdCloseQlike}{0.1266}  % QLIKE ridge lo_vd c16
\newcommand{\rhabRidgeLoVdCloseMse}{3.51}  % MSE x 1e10 ridge lo_vd c16
\newcommand{\rhabRidgeLoVdClosePct}{$-$1.0}  % QLIKE change % vs full
\newcommand{\rhabRidgeLoVdCloseDm}{$-$2.17}  % DM vs the reference arm
\newcommand{\rhabRidgeLoVdCloseP}{0.030}  % p-value
\newcommand{\rhabRidgeLoVdCloseLateQlike}{0.1232}  % QLIKE ridge lo_vd c16_2020
\newcommand{\rhabRidgeLoVdCloseLateMse}{7.98}  % MSE x 1e10 ridge lo_vd c16_2020
\newcommand{\rhabRidgeLoVdCloseLatePct}{$-$2.2}  % QLIKE change % vs full
\newcommand{\rhabRidgeLoVdCloseLateDm}{$-$2.26}  % DM vs the reference arm
\newcommand{\rhabRidgeLoVdCloseLateP}{0.024}  % p-value
\newcommand{\rhabRidgeLoVwAllQlike}{0.1789}  % QLIKE ridge lo_vw all
\newcommand{\rhabRidgeLoVwAllMse}{1.97}  % MSE x 1e10 ridge lo_vw all
\newcommand{\rhabRidgeLoVwAllPct}{$-$0.1}  % QLIKE change % vs full
\newcommand{\rhabRidgeLoVwAllDm}{$-$0.70}  % DM vs the reference arm
\newcommand{\rhabRidgeLoVwAllP}{0.484}  % p-value
\newcommand{\rhabRidgeLoVwCloseQlike}{0.1274}  % QLIKE ridge lo_vw c16
\newcommand{\rhabRidgeLoVwCloseMse}{3.50}  % MSE x 1e10 ridge lo_vw c16
\newcommand{\rhabRidgeLoVwClosePct}{$-$0.4}  % QLIKE change % vs full
\newcommand{\rhabRidgeLoVwCloseDm}{$-$0.90}  % DM vs the reference arm
\newcommand{\rhabRidgeLoVwCloseP}{0.366}  % p-value
\newcommand{\rhabRidgeLoVwCloseLateQlike}{0.1246}  % QLIKE ridge lo_vw c16_2020
\newcommand{\rhabRidgeLoVwCloseLateMse}{7.95}  % MSE x 1e10 ridge lo_vw c16_2020
\newcommand{\rhabRidgeLoVwCloseLatePct}{$-$1.1}  % QLIKE change % vs full
\newcommand{\rhabRidgeLoVwCloseLateDm}{$-$1.96}  % DM vs the reference arm
\newcommand{\rhabRidgeLoVwCloseLateP}{0.050}  % p-value
\newcommand{\rhabRidgeSharAllQlike}{0.1876}  % QLIKE ridge shar all
\newcommand{\rhabRidgeSharAllMse}{1.35}  % MSE x 1e10 ridge shar all
\newcommand{\rhabRidgeSharAllPct}{$-$0.9}  % QLIKE change % vs har
\newcommand{\rhabRidgeSharAllDm}{$-$3.95}  % DM vs the reference arm
\newcommand{\rhabRidgeSharAllP}{\ensuremath{<}0.001}  % p-value
\newcommand{\rhabRidgeSharCloseQlike}{0.1264}  % QLIKE ridge shar c16
\newcommand{\rhabRidgeSharCloseMse}{1.81}  % MSE x 1e10 ridge shar c16
\newcommand{\rhabRidgeSharClosePct}{$-$0.4}  % QLIKE change % vs har
\newcommand{\rhabRidgeSharCloseDm}{$-$0.69}  % DM vs the reference arm
\newcommand{\rhabRidgeSharCloseP}{0.488}  % p-value
\newcommand{\rhabRidgeSharCloseLateQlike}{0.1139}  % QLIKE ridge shar c16_2020
\newcommand{\rhabRidgeSharCloseLateMse}{3.87}  % MSE x 1e10 ridge shar c16_2020
\newcommand{\rhabRidgeSharCloseLatePct}{0.2}  % QLIKE change % vs har
\newcommand{\rhabRidgeSharCloseLateDm}{0.23}  % DM vs the reference arm
\newcommand{\rhabRidgeSharCloseLateP}{0.821}  % p-value
\newcommand{\rhabRidgeSharqlAllQlike}{0.1840}  % QLIKE ridge sharql all
\newcommand{\rhabRidgeSharqlAllMse}{1.68}  % MSE x 1e10 ridge sharql all
\newcommand{\rhabRidgeSharqlAllPct}{$-$2.8}  % QLIKE change % vs har
\newcommand{\rhabRidgeSharqlAllDm}{$-$6.82}  % DM vs the reference arm
\newcommand{\rhabRidgeSharqlAllP}{\ensuremath{<}0.001}  % p-value
\newcommand{\rhabRidgeSharqlCloseQlike}{0.1216}  % QLIKE ridge sharql c16
\newcommand{\rhabRidgeSharqlCloseMse}{2.65}  % MSE x 1e10 ridge sharql c16
\newcommand{\rhabRidgeSharqlClosePct}{$-$4.2}  % QLIKE change % vs har
\newcommand{\rhabRidgeSharqlCloseDm}{$-$3.49}  % DM vs the reference arm
\newcommand{\rhabRidgeSharqlCloseP}{\ensuremath{<}0.001}  % p-value
\newcommand{\rhabRidgeSharqlCloseLateQlike}{0.1119}  % QLIKE ridge sharql c16_2020
\newcommand{\rhabRidgeSharqlCloseLateMse}{5.92}  % MSE x 1e10 ridge sharql c16_2020
\newcommand{\rhabRidgeSharqlCloseLatePct}{$-$1.6}  % QLIKE change % vs har
\newcommand{\rhabRidgeSharqlCloseLateDm}{$-$0.73}  % DM vs the reference arm
\newcommand{\rhabRidgeSharqlCloseLateP}{0.468}  % p-value
\newcommand{\rhabXgbFullAllQlike}{0.1627}  % QLIKE xgb full all
\newcommand{\rhabXgbFullAllMse}{1.34}  % MSE x 1e10 xgb full all
\newcommand{\rhabXgbFullAllPct}{$-$13.8}  % QLIKE change % vs har
\newcommand{\rhabXgbFullAllDm}{$-$9.94}  % DM vs the reference arm
\newcommand{\rhabXgbFullAllP}{\ensuremath{<}0.001}  % p-value
\newcommand{\rhabXgbFullCloseQlike}{0.1192}  % QLIKE xgb full c16
\newcommand{\rhabXgbFullCloseMse}{2.20}  % MSE x 1e10 xgb full c16
\newcommand{\rhabXgbFullClosePct}{$-$3.4}  % QLIKE change % vs har
\newcommand{\rhabXgbFullCloseDm}{$-$2.18}  % DM vs the reference arm
\newcommand{\rhabXgbFullCloseP}{0.030}  % p-value
\newcommand{\rhabXgbFullCloseLateQlike}{0.1079}  % QLIKE xgb full c16_2020
\newcommand{\rhabXgbFullCloseLateMse}{4.79}  % MSE x 1e10 xgb full c16_2020
\newcommand{\rhabXgbFullCloseLatePct}{$-$2.1}  % QLIKE change % vs har
\newcommand{\rhabXgbFullCloseLateDm}{$-$0.76}  % DM vs the reference arm
\newcommand{\rhabXgbFullCloseLateP}{0.444}  % p-value
\newcommand{\rhabXgbHarAllQlike}{0.1888}  % QLIKE xgb har all
\newcommand{\rhabXgbHarAllMse}{1.16}  % MSE x 1e10 xgb har all
\newcommand{\rhabXgbHarCloseQlike}{0.1234}  % QLIKE xgb har c16
\newcommand{\rhabXgbHarCloseMse}{1.73}  % MSE x 1e10 xgb har c16
\newcommand{\rhabXgbHarCloseLateQlike}{0.1102}  % QLIKE xgb har c16_2020
\newcommand{\rhabXgbHarCloseLateMse}{3.62}  % MSE x 1e10 xgb har c16_2020
\newcommand{\rhabRidgeFullLiveSentQlike}{0.1833}  % QLIKE ridge full live_sent (subsample)
\newcommand{\rhabRidgeFullLiveSentPct}{$-$1.9}  % QLIKE change % vs the reference (subsample)
\newcommand{\rhabRidgeFullLiveSentDm}{$-$0.98}  % DM vs the reference arm (subsample)
\newcommand{\rhabRidgeFullLiveSentP}{0.327}  % p-value
\newcommand{\rhabRidgeLoSentLiveSentQlike}{0.1842}  % QLIKE ridge lo_sent live_sent (subsample)
\newcommand{\rhabRidgeLoSentLiveSentPct}{0.5}  % QLIKE change % vs the reference (subsample)
\newcommand{\rhabRidgeLoSentLiveSentDm}{0.71}  % DM vs the reference arm (subsample)
\newcommand{\rhabRidgeLoSentLiveSentP}{0.477}  % p-value
\newcommand{\rhabRidgeAddSentLiveSentQlike}{0.1846}  % QLIKE ridge add_sent live_sent (subsample)
\newcommand{\rhabRidgeAddSentLiveSentPct}{$-$1.2}  % QLIKE change % vs the reference (subsample)
\newcommand{\rhabRidgeAddSentLiveSentDm}{$-$1.18}  % DM vs the reference arm (subsample)
\newcommand{\rhabRidgeAddSentLiveSentP}{0.238}  % p-value
\newcommand{\rhabRidgeHarLiveSentQlike}{0.1868}  % QLIKE ridge har live_sent (subsample)
\newcommand{\rhabLassoFullLiveSentQlike}{0.1780}  % QLIKE lasso full live_sent (subsample)
\newcommand{\rhabLassoFullLiveSentPct}{$-$4.7}  % QLIKE change % vs the reference (subsample)
\newcommand{\rhabLassoFullLiveSentDm}{$-$1.83}  % DM vs the reference arm (subsample)
\newcommand{\rhabLassoFullLiveSentP}{0.068}  % p-value
\newcommand{\rhabLassoLoSentLiveSentQlike}{0.1786}  % QLIKE lasso lo_sent live_sent (subsample)
\newcommand{\rhabLassoLoSentLiveSentPct}{0.3}  % QLIKE change % vs the reference (subsample)
\newcommand{\rhabLassoLoSentLiveSentDm}{0.41}  % DM vs the reference arm (subsample)
\newcommand{\rhabLassoLoSentLiveSentP}{0.685}  % p-value
\newcommand{\rhabLassoAddSentLiveSentQlike}{0.1850}  % QLIKE lasso add_sent live_sent (subsample)
\newcommand{\rhabLassoAddSentLiveSentPct}{$-$1.0}  % QLIKE change % vs the reference (subsample)
\newcommand{\rhabLassoAddSentLiveSentDm}{$-$0.96}  % DM vs the reference arm (subsample)
\newcommand{\rhabLassoAddSentLiveSentP}{0.339}  % p-value
\newcommand{\rhabLassoHarLiveSentQlike}{0.1868}  % QLIKE lasso har live_sent (subsample)
\newcommand{\rhabEnetFullLiveSentQlike}{0.1818}  % QLIKE enet full live_sent (subsample)
\newcommand{\rhabEnetFullLiveSentPct}{$-$2.7}  % QLIKE change % vs the reference (subsample)
\newcommand{\rhabEnetFullLiveSentDm}{$-$1.55}  % DM vs the reference arm (subsample)
\newcommand{\rhabEnetFullLiveSentP}{0.121}  % p-value
\newcommand{\rhabEnetLoSentLiveSentQlike}{0.1833}  % QLIKE enet lo_sent live_sent (subsample)
\newcommand{\rhabEnetLoSentLiveSentPct}{0.8}  % QLIKE change % vs the reference (subsample)
\newcommand{\rhabEnetLoSentLiveSentDm}{1.43}  % DM vs the reference arm (subsample)
\newcommand{\rhabEnetLoSentLiveSentP}{0.154}  % p-value
\newcommand{\rhabEnetAddSentLiveSentQlike}{0.1845}  % QLIKE enet add_sent live_sent (subsample)
\newcommand{\rhabEnetAddSentLiveSentPct}{$-$1.3}  % QLIKE change % vs the reference (subsample)
\newcommand{\rhabEnetAddSentLiveSentDm}{$-$1.48}  % DM vs the reference arm (subsample)
\newcommand{\rhabEnetAddSentLiveSentP}{0.139}  % p-value
\newcommand{\rhabEnetHarLiveSentQlike}{0.1868}  % QLIKE enet har live_sent (subsample)
\newcommand{\rhabLgbmFullLiveSentQlike}{0.1564}  % QLIKE lgbm full live_sent (subsample)
\newcommand{\rhabLgbmFullLiveSentPct}{$-$13.8}  % QLIKE change % vs the reference (subsample)
\newcommand{\rhabLgbmFullLiveSentDm}{$-$3.76}  % DM vs the reference arm (subsample)
\newcommand{\rhabLgbmFullLiveSentP}{\ensuremath{<}0.001}  % p-value
\newcommand{\rhabLgbmHarLiveSentQlike}{0.1814}  % QLIKE lgbm har live_sent (subsample)
\newcommand{\rhabXgbFullLiveSentQlike}{0.1559}  % QLIKE xgb full live_sent (subsample)
\newcommand{\rhabXgbFullLiveSentPct}{$-$15.6}  % QLIKE change % vs the reference (subsample)
\newcommand{\rhabXgbFullLiveSentDm}{$-$3.69}  % DM vs the reference arm (subsample)
\newcommand{\rhabXgbFullLiveSentP}{\ensuremath{<}0.001}  % p-value
\newcommand{\rhabXgbHarLiveSentQlike}{0.1846}  % QLIKE xgb har live_sent (subsample)
\newcommand{\rhabRfFullLiveSentQlike}{0.1676}  % QLIKE rf full live_sent (subsample)
\newcommand{\rhabRfFullLiveSentPct}{$-$10.8}  % QLIKE change % vs the reference (subsample)
\newcommand{\rhabRfFullLiveSentDm}{$-$4.09}  % DM vs the reference arm (subsample)
\newcommand{\rhabRfFullLiveSentP}{\ensuremath{<}0.001}  % p-value
\newcommand{\rhabRfHarLiveSentQlike}{0.1878}  % QLIKE rf har live_sent (subsample)
\newcommand{\rhabRidgeFullLiveSentCloseQlike}{0.1222}  % QLIKE ridge full live16_sent (subsample)
\newcommand{\rhabRidgeFullLiveSentClosePct}{12.8}  % QLIKE change % vs the reference (subsample)
\newcommand{\rhabRidgeFullLiveSentCloseDm}{2.03}  % DM vs the reference arm (subsample)
\newcommand{\rhabRidgeFullLiveSentCloseP}{0.042}  % p-value
\newcommand{\rhabRidgeLoSentLiveSentCloseQlike}{0.1141}  % QLIKE ridge lo_sent live16_sent (subsample)
\newcommand{\rhabRidgeLoSentLiveSentClosePct}{$-$6.6}  % QLIKE change % vs the reference (subsample)
\newcommand{\rhabRidgeLoSentLiveSentCloseDm}{$-$3.48}  % DM vs the reference arm (subsample)
\newcommand{\rhabRidgeLoSentLiveSentCloseP}{\ensuremath{<}0.001}  % p-value
\newcommand{\rhabRidgeAddSentLiveSentCloseQlike}{0.1188}  % QLIKE ridge add_sent live16_sent (subsample)
\newcommand{\rhabRidgeAddSentLiveSentClosePct}{9.7}  % QLIKE change % vs the reference (subsample)
\newcommand{\rhabRidgeAddSentLiveSentCloseDm}{2.63}  % DM vs the reference arm (subsample)
\newcommand{\rhabRidgeAddSentLiveSentCloseP}{0.009}  % p-value
\newcommand{\rhabRidgeHarLiveSentCloseQlike}{0.1083}  % QLIKE ridge har live16_sent (subsample)
\newcommand{\rhabLassoFullLiveSentCloseQlike}{0.1234}  % QLIKE lasso full live16_sent (subsample)
\newcommand{\rhabLassoFullLiveSentClosePct}{13.9}  % QLIKE change % vs the reference (subsample)
\newcommand{\rhabLassoFullLiveSentCloseDm}{2.17}  % DM vs the reference arm (subsample)
\newcommand{\rhabLassoFullLiveSentCloseP}{0.030}  % p-value
\newcommand{\rhabLassoLoSentLiveSentCloseQlike}{0.1155}  % QLIKE lasso lo_sent live16_sent (subsample)
\newcommand{\rhabLassoLoSentLiveSentClosePct}{$-$6.3}  % QLIKE change % vs the reference (subsample)
\newcommand{\rhabLassoLoSentLiveSentCloseDm}{$-$2.92}  % DM vs the reference arm (subsample)
\newcommand{\rhabLassoLoSentLiveSentCloseP}{0.003}  % p-value
\newcommand{\rhabLassoAddSentLiveSentCloseQlike}{0.1195}  % QLIKE lasso add_sent live16_sent (subsample)
\newcommand{\rhabLassoAddSentLiveSentClosePct}{10.3}  % QLIKE change % vs the reference (subsample)
\newcommand{\rhabLassoAddSentLiveSentCloseDm}{2.76}  % DM vs the reference arm (subsample)
\newcommand{\rhabLassoAddSentLiveSentCloseP}{0.006}  % p-value
\newcommand{\rhabLassoHarLiveSentCloseQlike}{0.1083}  % QLIKE lasso har live16_sent (subsample)
\newcommand{\rhabEnetFullLiveSentCloseQlike}{0.1171}  % QLIKE enet full live16_sent (subsample)
\newcommand{\rhabEnetFullLiveSentClosePct}{8.1}  % QLIKE change % vs the reference (subsample)
\newcommand{\rhabEnetFullLiveSentCloseDm}{1.35}  % DM vs the reference arm (subsample)
\newcommand{\rhabEnetFullLiveSentCloseP}{0.176}  % p-value
\newcommand{\rhabEnetLoSentLiveSentCloseQlike}{0.1134}  % QLIKE enet lo_sent live16_sent (subsample)
\newcommand{\rhabEnetLoSentLiveSentClosePct}{$-$3.2}  % QLIKE change % vs the reference (subsample)
\newcommand{\rhabEnetLoSentLiveSentCloseDm}{$-$2.23}  % DM vs the reference arm (subsample)
\newcommand{\rhabEnetLoSentLiveSentCloseP}{0.026}  % p-value
\newcommand{\rhabEnetAddSentLiveSentCloseQlike}{0.1156}  % QLIKE enet add_sent live16_sent (subsample)
\newcommand{\rhabEnetAddSentLiveSentClosePct}{6.7}  % QLIKE change % vs the reference (subsample)
\newcommand{\rhabEnetAddSentLiveSentCloseDm}{1.91}  % DM vs the reference arm (subsample)
\newcommand{\rhabEnetAddSentLiveSentCloseP}{0.056}  % p-value
\newcommand{\rhabEnetHarLiveSentCloseQlike}{0.1083}  % QLIKE enet har live16_sent (subsample)
\newcommand{\rhabLgbmFullLiveSentCloseQlike}{0.0995}  % QLIKE lgbm full live16_sent (subsample)
\newcommand{\rhabLgbmFullLiveSentClosePct}{$-$2.5}  % QLIKE change % vs the reference (subsample)
\newcommand{\rhabLgbmFullLiveSentCloseDm}{$-$0.57}  % DM vs the reference arm (subsample)
\newcommand{\rhabLgbmFullLiveSentCloseP}{0.565}  % p-value
\newcommand{\rhabLgbmHarLiveSentCloseQlike}{0.1021}  % QLIKE lgbm har live16_sent (subsample)
\newcommand{\rhabXgbFullLiveSentCloseQlike}{0.0997}  % QLIKE xgb full live16_sent (subsample)
\newcommand{\rhabXgbFullLiveSentClosePct}{$-$2.8}  % QLIKE change % vs the reference (subsample)
\newcommand{\rhabXgbFullLiveSentCloseDm}{$-$0.61}  % DM vs the reference arm (subsample)
\newcommand{\rhabXgbFullLiveSentCloseP}{0.541}  % p-value
\newcommand{\rhabXgbHarLiveSentCloseQlike}{0.1025}  % QLIKE xgb har live16_sent (subsample)
\newcommand{\rhabRfFullLiveSentCloseQlike}{0.1026}  % QLIKE rf full live16_sent (subsample)
\newcommand{\rhabRfFullLiveSentClosePct}{$-$2.6}  % QLIKE change % vs the reference (subsample)
\newcommand{\rhabRfFullLiveSentCloseDm}{$-$0.65}  % DM vs the reference arm (subsample)
\newcommand{\rhabRfFullLiveSentCloseP}{0.513}  % p-value
\newcommand{\rhabRfHarLiveSentCloseQlike}{0.1053}  % QLIKE rf har live16_sent (subsample)
\newcommand{\rhabRidgeFullLiveIvQlike}{0.1940}  % QLIKE ridge full live_iv (subsample)
\newcommand{\rhabRidgeFullLiveIvPct}{$-$4.7}  % QLIKE change % vs the reference (subsample)
\newcommand{\rhabRidgeFullLiveIvDm}{$-$3.62}  % DM vs the reference arm (subsample)
\newcommand{\rhabRidgeFullLiveIvP}{\ensuremath{<}0.001}  % p-value
\newcommand{\rhabRidgeLoIvLiveIvQlike}{0.1942}  % QLIKE ridge lo_iv live_iv (subsample)
\newcommand{\rhabRidgeLoIvLiveIvPct}{0.1}  % QLIKE change % vs the reference (subsample)
\newcommand{\rhabRidgeLoIvLiveIvDm}{0.25}  % DM vs the reference arm (subsample)
\newcommand{\rhabRidgeLoIvLiveIvP}{0.806}  % p-value
\newcommand{\rhabRidgeAddIvLiveIvQlike}{0.2037}  % QLIKE ridge add_iv live_iv (subsample)
\newcommand{\rhabRidgeAddIvLiveIvPct}{0.0}  % QLIKE change % vs the reference (subsample)
\newcommand{\rhabRidgeAddIvLiveIvDm}{0.02}  % DM vs the reference arm (subsample)
\newcommand{\rhabRidgeAddIvLiveIvP}{0.986}  % p-value
\newcommand{\rhabRidgeHarLiveIvQlike}{0.2037}  % QLIKE ridge har live_iv (subsample)
\newcommand{\rhabLassoFullLiveIvQlike}{0.1924}  % QLIKE lasso full live_iv (subsample)
\newcommand{\rhabLassoFullLiveIvPct}{$-$5.5}  % QLIKE change % vs the reference (subsample)
\newcommand{\rhabLassoFullLiveIvDm}{$-$2.98}  % DM vs the reference arm (subsample)
\newcommand{\rhabLassoFullLiveIvP}{0.003}  % p-value
\newcommand{\rhabLassoLoIvLiveIvQlike}{0.1931}  % QLIKE lasso lo_iv live_iv (subsample)
\newcommand{\rhabLassoLoIvLiveIvPct}{0.4}  % QLIKE change % vs the reference (subsample)
\newcommand{\rhabLassoLoIvLiveIvDm}{0.47}  % DM vs the reference arm (subsample)
\newcommand{\rhabLassoLoIvLiveIvP}{0.641}  % p-value
\newcommand{\rhabLassoAddIvLiveIvQlike}{0.2019}  % QLIKE lasso add_iv live_iv (subsample)
\newcommand{\rhabLassoAddIvLiveIvPct}{$-$0.9}  % QLIKE change % vs the reference (subsample)
\newcommand{\rhabLassoAddIvLiveIvDm}{$-$1.39}  % DM vs the reference arm (subsample)
\newcommand{\rhabLassoAddIvLiveIvP}{0.165}  % p-value
\newcommand{\rhabLassoHarLiveIvQlike}{0.2037}  % QLIKE lasso har live_iv (subsample)
\newcommand{\rhabEnetFullLiveIvQlike}{0.1921}  % QLIKE enet full live_iv (subsample)
\newcommand{\rhabEnetFullLiveIvPct}{$-$5.7}  % QLIKE change % vs the reference (subsample)
\newcommand{\rhabEnetFullLiveIvDm}{$-$4.87}  % DM vs the reference arm (subsample)
\newcommand{\rhabEnetFullLiveIvP}{\ensuremath{<}0.001}  % p-value
\newcommand{\rhabEnetLoIvLiveIvQlike}{0.1924}  % QLIKE enet lo_iv live_iv (subsample)
\newcommand{\rhabEnetLoIvLiveIvPct}{0.1}  % QLIKE change % vs the reference (subsample)
\newcommand{\rhabEnetLoIvLiveIvDm}{0.23}  % DM vs the reference arm (subsample)
\newcommand{\rhabEnetLoIvLiveIvP}{0.816}  % p-value
\newcommand{\rhabEnetAddIvLiveIvQlike}{0.2033}  % QLIKE enet add_iv live_iv (subsample)
\newcommand{\rhabEnetAddIvLiveIvPct}{$-$0.2}  % QLIKE change % vs the reference (subsample)
\newcommand{\rhabEnetAddIvLiveIvDm}{$-$0.36}  % DM vs the reference arm (subsample)
\newcommand{\rhabEnetAddIvLiveIvP}{0.722}  % p-value
\newcommand{\rhabEnetHarLiveIvQlike}{0.2037}  % QLIKE enet har live_iv (subsample)
\newcommand{\rhabLgbmFullLiveIvQlike}{0.1739}  % QLIKE lgbm full live_iv (subsample)
\newcommand{\rhabLgbmFullLiveIvPct}{$-$13.3}  % QLIKE change % vs the reference (subsample)
\newcommand{\rhabLgbmFullLiveIvDm}{$-$7.31}  % DM vs the reference arm (subsample)
\newcommand{\rhabLgbmFullLiveIvP}{\ensuremath{<}0.001}  % p-value
\newcommand{\rhabLgbmHarLiveIvQlike}{0.2006}  % QLIKE lgbm har live_iv (subsample)
\newcommand{\rhabXgbFullLiveIvQlike}{0.1739}  % QLIKE xgb full live_iv (subsample)
\newcommand{\rhabXgbFullLiveIvPct}{$-$14.1}  % QLIKE change % vs the reference (subsample)
\newcommand{\rhabXgbFullLiveIvDm}{$-$7.16}  % DM vs the reference arm (subsample)
\newcommand{\rhabXgbFullLiveIvP}{\ensuremath{<}0.001}  % p-value
\newcommand{\rhabXgbHarLiveIvQlike}{0.2024}  % QLIKE xgb har live_iv (subsample)
\newcommand{\rhabRfFullLiveIvQlike}{0.1818}  % QLIKE rf full live_iv (subsample)
\newcommand{\rhabRfFullLiveIvPct}{$-$10.4}  % QLIKE change % vs the reference (subsample)
\newcommand{\rhabRfFullLiveIvDm}{$-$7.47}  % DM vs the reference arm (subsample)
\newcommand{\rhabRfFullLiveIvP}{\ensuremath{<}0.001}  % p-value
\newcommand{\rhabRfHarLiveIvQlike}{0.2029}  % QLIKE rf har live_iv (subsample)
\newcommand{\rhabRidgeFullLiveIvCloseQlike}{0.1380}  % QLIKE ridge full live16_iv (subsample)
\newcommand{\rhabRidgeFullLiveIvClosePct}{6.0}  % QLIKE change % vs the reference (subsample)
\newcommand{\rhabRidgeFullLiveIvCloseDm}{1.50}  % DM vs the reference arm (subsample)
\newcommand{\rhabRidgeFullLiveIvCloseP}{0.134}  % p-value
\newcommand{\rhabRidgeLoIvLiveIvCloseQlike}{0.1355}  % QLIKE ridge lo_iv live16_iv (subsample)
\newcommand{\rhabRidgeLoIvLiveIvClosePct}{$-$1.8}  % QLIKE change % vs the reference (subsample)
\newcommand{\rhabRidgeLoIvLiveIvCloseDm}{$-$1.22}  % DM vs the reference arm (subsample)
\newcommand{\rhabRidgeLoIvLiveIvCloseP}{0.222}  % p-value
\newcommand{\rhabRidgeAddIvLiveIvCloseQlike}{0.1306}  % QLIKE ridge add_iv live16_iv (subsample)
\newcommand{\rhabRidgeAddIvLiveIvClosePct}{0.3}  % QLIKE change % vs the reference (subsample)
\newcommand{\rhabRidgeAddIvLiveIvCloseDm}{0.19}  % DM vs the reference arm (subsample)
\newcommand{\rhabRidgeAddIvLiveIvCloseP}{0.853}  % p-value
\newcommand{\rhabRidgeHarLiveIvCloseQlike}{0.1302}  % QLIKE ridge har live16_iv (subsample)
\newcommand{\rhabLassoFullLiveIvCloseQlike}{0.1425}  % QLIKE lasso full live16_iv (subsample)
\newcommand{\rhabLassoFullLiveIvClosePct}{9.4}  % QLIKE change % vs the reference (subsample)
\newcommand{\rhabLassoFullLiveIvCloseDm}{1.93}  % DM vs the reference arm (subsample)
\newcommand{\rhabLassoFullLiveIvCloseP}{0.054}  % p-value
\newcommand{\rhabLassoLoIvLiveIvCloseQlike}{0.1378}  % QLIKE lasso lo_iv live16_iv (subsample)
\newcommand{\rhabLassoLoIvLiveIvClosePct}{$-$3.3}  % QLIKE change % vs the reference (subsample)
\newcommand{\rhabLassoLoIvLiveIvCloseDm}{$-$1.78}  % DM vs the reference arm (subsample)
\newcommand{\rhabLassoLoIvLiveIvCloseP}{0.075}  % p-value
\newcommand{\rhabLassoAddIvLiveIvCloseQlike}{0.1300}  % QLIKE lasso add_iv live16_iv (subsample)
\newcommand{\rhabLassoAddIvLiveIvClosePct}{$-$0.1}  % QLIKE change % vs the reference (subsample)
\newcommand{\rhabLassoAddIvLiveIvCloseDm}{$-$0.08}  % DM vs the reference arm (subsample)
\newcommand{\rhabLassoAddIvLiveIvCloseP}{0.940}  % p-value
\newcommand{\rhabLassoHarLiveIvCloseQlike}{0.1302}  % QLIKE lasso har live16_iv (subsample)
\newcommand{\rhabEnetFullLiveIvCloseQlike}{0.1355}  % QLIKE enet full live16_iv (subsample)
\newcommand{\rhabEnetFullLiveIvClosePct}{4.1}  % QLIKE change % vs the reference (subsample)
\newcommand{\rhabEnetFullLiveIvCloseDm}{1.12}  % DM vs the reference arm (subsample)
\newcommand{\rhabEnetFullLiveIvCloseP}{0.262}  % p-value
\newcommand{\rhabEnetLoIvLiveIvCloseQlike}{0.1321}  % QLIKE enet lo_iv live16_iv (subsample)
\newcommand{\rhabEnetLoIvLiveIvClosePct}{$-$2.5}  % QLIKE change % vs the reference (subsample)
\newcommand{\rhabEnetLoIvLiveIvCloseDm}{$-$1.58}  % DM vs the reference arm (subsample)
\newcommand{\rhabEnetLoIvLiveIvCloseP}{0.115}  % p-value
\newcommand{\rhabEnetAddIvLiveIvCloseQlike}{0.1308}  % QLIKE enet add_iv live16_iv (subsample)
\newcommand{\rhabEnetAddIvLiveIvClosePct}{0.5}  % QLIKE change % vs the reference (subsample)
\newcommand{\rhabEnetAddIvLiveIvCloseDm}{0.31}  % DM vs the reference arm (subsample)
\newcommand{\rhabEnetAddIvLiveIvCloseP}{0.756}  % p-value
\newcommand{\rhabEnetHarLiveIvCloseQlike}{0.1302}  % QLIKE enet har live16_iv (subsample)
\newcommand{\rhabLgbmFullLiveIvCloseQlike}{0.1242}  % QLIKE lgbm full live16_iv (subsample)
\newcommand{\rhabLgbmFullLiveIvClosePct}{$-$1.6}  % QLIKE change % vs the reference (subsample)
\newcommand{\rhabLgbmFullLiveIvCloseDm}{$-$0.74}  % DM vs the reference arm (subsample)
\newcommand{\rhabLgbmFullLiveIvCloseP}{0.461}  % p-value
\newcommand{\rhabLgbmHarLiveIvCloseQlike}{0.1263}  % QLIKE lgbm har live16_iv (subsample)
\newcommand{\rhabXgbFullLiveIvCloseQlike}{0.1234}  % QLIKE xgb full live16_iv (subsample)
\newcommand{\rhabXgbFullLiveIvClosePct}{$-$2.3}  % QLIKE change % vs the reference (subsample)
\newcommand{\rhabXgbFullLiveIvCloseDm}{$-$1.01}  % DM vs the reference arm (subsample)
\newcommand{\rhabXgbFullLiveIvCloseP}{0.312}  % p-value
\newcommand{\rhabXgbHarLiveIvCloseQlike}{0.1263}  % QLIKE xgb har live16_iv (subsample)
\newcommand{\rhabRfFullLiveIvCloseQlike}{0.1304}  % QLIKE rf full live16_iv (subsample)
\newcommand{\rhabRfFullLiveIvClosePct}{$-$0.4}  % QLIKE change % vs the reference (subsample)
\newcommand{\rhabRfFullLiveIvCloseDm}{$-$0.16}  % DM vs the reference arm (subsample)
\newcommand{\rhabRfFullLiveIvCloseP}{0.870}  % p-value
\newcommand{\rhabRfHarLiveIvCloseQlike}{0.1309}  % QLIKE rf har live16_iv (subsample)
\newcommand{\rhabRidgeFullLiveIvbQlike}{0.1757}  % QLIKE ridge full live_ivb (subsample)
\newcommand{\rhabRidgeFullLiveIvbPct}{$-$3.6}  % QLIKE change % vs the reference (subsample)
\newcommand{\rhabRidgeFullLiveIvbDm}{$-$2.12}  % DM vs the reference arm (subsample)
\newcommand{\rhabRidgeFullLiveIvbP}{0.034}  % p-value
\newcommand{\rhabRidgeLoIvLiveIvbQlike}{0.1746}  % QLIKE ridge lo_iv live_ivb (subsample)
\newcommand{\rhabRidgeLoIvLiveIvbPct}{$-$0.6}  % QLIKE change % vs the reference (subsample)
\newcommand{\rhabRidgeLoIvLiveIvbDm}{$-$0.79}  % DM vs the reference arm (subsample)
\newcommand{\rhabRidgeLoIvLiveIvbP}{0.431}  % p-value
\newcommand{\rhabRidgeAddIvLiveIvbQlike}{0.1850}  % QLIKE ridge add_iv live_ivb (subsample)
\newcommand{\rhabRidgeAddIvLiveIvbPct}{1.4}  % QLIKE change % vs the reference (subsample)
\newcommand{\rhabRidgeAddIvLiveIvbDm}{1.50}  % DM vs the reference arm (subsample)
\newcommand{\rhabRidgeAddIvLiveIvbP}{0.134}  % p-value
\newcommand{\rhabRidgeHarLiveIvbQlike}{0.1824}  % QLIKE ridge har live_ivb (subsample)
\newcommand{\rhabLassoFullLiveIvbQlike}{0.1718}  % QLIKE lasso full live_ivb (subsample)
\newcommand{\rhabLassoFullLiveIvbPct}{$-$5.8}  % QLIKE change % vs the reference (subsample)
\newcommand{\rhabLassoFullLiveIvbDm}{$-$2.58}  % DM vs the reference arm (subsample)
\newcommand{\rhabLassoFullLiveIvbP}{0.010}  % p-value
\newcommand{\rhabLassoLoIvLiveIvbQlike}{0.1700}  % QLIKE lasso lo_iv live_ivb (subsample)
\newcommand{\rhabLassoLoIvLiveIvbPct}{$-$1.0}  % QLIKE change % vs the reference (subsample)
\newcommand{\rhabLassoLoIvLiveIvbDm}{$-$0.96}  % DM vs the reference arm (subsample)
\newcommand{\rhabLassoLoIvLiveIvbP}{0.339}  % p-value
\newcommand{\rhabLassoAddIvLiveIvbQlike}{0.1844}  % QLIKE lasso add_iv live_ivb (subsample)
\newcommand{\rhabLassoAddIvLiveIvbPct}{1.1}  % QLIKE change % vs the reference (subsample)
\newcommand{\rhabLassoAddIvLiveIvbDm}{1.04}  % DM vs the reference arm (subsample)
\newcommand{\rhabLassoAddIvLiveIvbP}{0.300}  % p-value
\newcommand{\rhabLassoHarLiveIvbQlike}{0.1824}  % QLIKE lasso har live_ivb (subsample)
\newcommand{\rhabEnetFullLiveIvbQlike}{0.1745}  % QLIKE enet full live_ivb (subsample)
\newcommand{\rhabEnetFullLiveIvbPct}{$-$4.3}  % QLIKE change % vs the reference (subsample)
\newcommand{\rhabEnetFullLiveIvbDm}{$-$2.67}  % DM vs the reference arm (subsample)
\newcommand{\rhabEnetFullLiveIvbP}{0.008}  % p-value
\newcommand{\rhabEnetLoIvLiveIvbQlike}{0.1732}  % QLIKE enet lo_iv live_ivb (subsample)
\newcommand{\rhabEnetLoIvLiveIvbPct}{$-$0.7}  % QLIKE change % vs the reference (subsample)
\newcommand{\rhabEnetLoIvLiveIvbDm}{$-$0.88}  % DM vs the reference arm (subsample)
\newcommand{\rhabEnetLoIvLiveIvbP}{0.377}  % p-value
\newcommand{\rhabEnetAddIvLiveIvbQlike}{0.1844}  % QLIKE enet add_iv live_ivb (subsample)
\newcommand{\rhabEnetAddIvLiveIvbPct}{1.1}  % QLIKE change % vs the reference (subsample)
\newcommand{\rhabEnetAddIvLiveIvbDm}{1.30}  % DM vs the reference arm (subsample)
\newcommand{\rhabEnetAddIvLiveIvbP}{0.195}  % p-value
\newcommand{\rhabEnetHarLiveIvbQlike}{0.1824}  % QLIKE enet har live_ivb (subsample)
\newcommand{\rhabLgbmFullLiveIvbQlike}{0.1518}  % QLIKE lgbm full live_ivb (subsample)
\newcommand{\rhabLgbmFullLiveIvbPct}{$-$14.6}  % QLIKE change % vs the reference (subsample)
\newcommand{\rhabLgbmFullLiveIvbDm}{$-$5.14}  % DM vs the reference arm (subsample)
\newcommand{\rhabLgbmFullLiveIvbP}{\ensuremath{<}0.001}  % p-value
\newcommand{\rhabLgbmHarLiveIvbQlike}{0.1778}  % QLIKE lgbm har live_ivb (subsample)
\newcommand{\rhabXgbFullLiveIvbQlike}{0.1511}  % QLIKE xgb full live_ivb (subsample)
\newcommand{\rhabXgbFullLiveIvbPct}{$-$16.3}  % QLIKE change % vs the reference (subsample)
\newcommand{\rhabXgbFullLiveIvbDm}{$-$4.98}  % DM vs the reference arm (subsample)
\newcommand{\rhabXgbFullLiveIvbP}{\ensuremath{<}0.001}  % p-value
\newcommand{\rhabXgbHarLiveIvbQlike}{0.1805}  % QLIKE xgb har live_ivb (subsample)
\newcommand{\rhabRfFullLiveIvbQlike}{0.1621}  % QLIKE rf full live_ivb (subsample)
\newcommand{\rhabRfFullLiveIvbPct}{$-$11.3}  % QLIKE change % vs the reference (subsample)
\newcommand{\rhabRfFullLiveIvbDm}{$-$5.43}  % DM vs the reference arm (subsample)
\newcommand{\rhabRfFullLiveIvbP}{\ensuremath{<}0.001}  % p-value
\newcommand{\rhabRfHarLiveIvbQlike}{0.1829}  % QLIKE rf har live_ivb (subsample)
\newcommand{\rhabRidgeFullLiveIvbCloseQlike}{0.1180}  % QLIKE ridge full live16_ivb (subsample)
\newcommand{\rhabRidgeFullLiveIvbClosePct}{15.1}  % QLIKE change % vs the reference (subsample)
\newcommand{\rhabRidgeFullLiveIvbCloseDm}{2.54}  % DM vs the reference arm (subsample)
\newcommand{\rhabRidgeFullLiveIvbCloseP}{0.011}  % p-value
\newcommand{\rhabRidgeLoIvLiveIvbCloseQlike}{0.1123}  % QLIKE ridge lo_iv live16_ivb (subsample)
\newcommand{\rhabRidgeLoIvLiveIvbClosePct}{$-$4.9}  % QLIKE change % vs the reference (subsample)
\newcommand{\rhabRidgeLoIvLiveIvbCloseDm}{$-$1.73}  % DM vs the reference arm (subsample)
\newcommand{\rhabRidgeLoIvLiveIvbCloseP}{0.083}  % p-value
\newcommand{\rhabRidgeAddIvLiveIvbCloseQlike}{0.1073}  % QLIKE ridge add_iv live16_ivb (subsample)
\newcommand{\rhabRidgeAddIvLiveIvbClosePct}{4.6}  % QLIKE change % vs the reference (subsample)
\newcommand{\rhabRidgeAddIvLiveIvbCloseDm}{1.47}  % DM vs the reference arm (subsample)
\newcommand{\rhabRidgeAddIvLiveIvbCloseP}{0.142}  % p-value
\newcommand{\rhabRidgeHarLiveIvbCloseQlike}{0.1026}  % QLIKE ridge har live16_ivb (subsample)
\newcommand{\rhabLassoFullLiveIvbCloseQlike}{0.1202}  % QLIKE lasso full live16_ivb (subsample)
\newcommand{\rhabLassoFullLiveIvbClosePct}{17.2}  % QLIKE change % vs the reference (subsample)
\newcommand{\rhabLassoFullLiveIvbCloseDm}{2.74}  % DM vs the reference arm (subsample)
\newcommand{\rhabLassoFullLiveIvbCloseP}{0.006}  % p-value
\newcommand{\rhabLassoLoIvLiveIvbCloseQlike}{0.1120}  % QLIKE lasso lo_iv live16_ivb (subsample)
\newcommand{\rhabLassoLoIvLiveIvbClosePct}{$-$6.8}  % QLIKE change % vs the reference (subsample)
\newcommand{\rhabLassoLoIvLiveIvbCloseDm}{$-$2.01}  % DM vs the reference arm (subsample)
\newcommand{\rhabLassoLoIvLiveIvbCloseP}{0.044}  % p-value
\newcommand{\rhabLassoAddIvLiveIvbCloseQlike}{0.1066}  % QLIKE lasso add_iv live16_ivb (subsample)
\newcommand{\rhabLassoAddIvLiveIvbClosePct}{3.9}  % QLIKE change % vs the reference (subsample)
\newcommand{\rhabLassoAddIvLiveIvbCloseDm}{1.22}  % DM vs the reference arm (subsample)
\newcommand{\rhabLassoAddIvLiveIvbCloseP}{0.221}  % p-value
\newcommand{\rhabLassoHarLiveIvbCloseQlike}{0.1026}  % QLIKE lasso har live16_ivb (subsample)
\newcommand{\rhabEnetFullLiveIvbCloseQlike}{0.1147}  % QLIKE enet full live16_ivb (subsample)
\newcommand{\rhabEnetFullLiveIvbClosePct}{11.8}  % QLIKE change % vs the reference (subsample)
\newcommand{\rhabEnetFullLiveIvbCloseDm}{2.03}  % DM vs the reference arm (subsample)
\newcommand{\rhabEnetFullLiveIvbCloseP}{0.043}  % p-value
\newcommand{\rhabEnetLoIvLiveIvbCloseQlike}{0.1073}  % QLIKE enet lo_iv live16_ivb (subsample)
\newcommand{\rhabEnetLoIvLiveIvbClosePct}{$-$6.5}  % QLIKE change % vs the reference (subsample)
\newcommand{\rhabEnetLoIvLiveIvbCloseDm}{$-$2.14}  % DM vs the reference arm (subsample)
\newcommand{\rhabEnetLoIvLiveIvbCloseP}{0.033}  % p-value
\newcommand{\rhabEnetAddIvLiveIvbCloseQlike}{0.1076}  % QLIKE enet add_iv live16_ivb (subsample)
\newcommand{\rhabEnetAddIvLiveIvbClosePct}{4.9}  % QLIKE change % vs the reference (subsample)
\newcommand{\rhabEnetAddIvLiveIvbCloseDm}{1.60}  % DM vs the reference arm (subsample)
\newcommand{\rhabEnetAddIvLiveIvbCloseP}{0.110}  % p-value
\newcommand{\rhabEnetHarLiveIvbCloseQlike}{0.1026}  % QLIKE enet har live16_ivb (subsample)
\newcommand{\rhabLgbmFullLiveIvbCloseQlike}{0.0964}  % QLIKE lgbm full live16_ivb (subsample)
\newcommand{\rhabLgbmFullLiveIvbClosePct}{$-$1.0}  % QLIKE change % vs the reference (subsample)
\newcommand{\rhabLgbmFullLiveIvbCloseDm}{$-$0.29}  % DM vs the reference arm (subsample)
\newcommand{\rhabLgbmFullLiveIvbCloseP}{0.775}  % p-value
\newcommand{\rhabLgbmHarLiveIvbCloseQlike}{0.0974}  % QLIKE lgbm har live16_ivb (subsample)
\newcommand{\rhabXgbFullLiveIvbCloseQlike}{0.0962}  % QLIKE xgb full live16_ivb (subsample)
\newcommand{\rhabXgbFullLiveIvbClosePct}{$-$1.5}  % QLIKE change % vs the reference (subsample)
\newcommand{\rhabXgbFullLiveIvbCloseDm}{$-$0.40}  % DM vs the reference arm (subsample)
\newcommand{\rhabXgbFullLiveIvbCloseP}{0.689}  % p-value
\newcommand{\rhabXgbHarLiveIvbCloseQlike}{0.0976}  % QLIKE xgb har live16_ivb (subsample)
\newcommand{\rhabRfFullLiveIvbCloseQlike}{0.1005}  % QLIKE rf full live16_ivb (subsample)
\newcommand{\rhabRfFullLiveIvbClosePct}{0.2}  % QLIKE change % vs the reference (subsample)
\newcommand{\rhabRfFullLiveIvbCloseDm}{0.05}  % DM vs the reference arm (subsample)
\newcommand{\rhabRfFullLiveIvbCloseP}{0.958}  % p-value
\newcommand{\rhabRfHarLiveIvbCloseQlike}{0.1003}  % QLIKE rf har live16_ivb (subsample)
\newcommand{\rhabRidgeFullLiveVdQlike}{0.1942}  % QLIKE ridge full live_vd (subsample)
\newcommand{\rhabRidgeFullLiveVdPct}{4.2}  % QLIKE change % vs the reference (subsample)
\newcommand{\rhabRidgeFullLiveVdDm}{1.33}  % DM vs the reference arm (subsample)
\newcommand{\rhabRidgeFullLiveVdP}{0.183}  % p-value
\newcommand{\rhabRidgeLoVdLiveVdQlike}{0.1887}  % QLIKE ridge lo_vd live_vd (subsample)
\newcommand{\rhabRidgeLoVdLiveVdPct}{$-$2.8}  % QLIKE change % vs the reference (subsample)
\newcommand{\rhabRidgeLoVdLiveVdDm}{$-$3.17}  % DM vs the reference arm (subsample)
\newcommand{\rhabRidgeLoVdLiveVdP}{0.002}  % p-value
\newcommand{\rhabRidgeAddVdLiveVdQlike}{0.1944}  % QLIKE ridge add_vd live_vd (subsample)
\newcommand{\rhabRidgeAddVdLiveVdPct}{4.3}  % QLIKE change % vs the reference (subsample)
\newcommand{\rhabRidgeAddVdLiveVdDm}{2.52}  % DM vs the reference arm (subsample)
\newcommand{\rhabRidgeAddVdLiveVdP}{0.012}  % p-value
\newcommand{\rhabRidgeHarLiveVdQlike}{0.1864}  % QLIKE ridge har live_vd (subsample)
\newcommand{\rhabLassoFullLiveVdQlike}{0.1898}  % QLIKE lasso full live_vd (subsample)
\newcommand{\rhabLassoFullLiveVdPct}{1.9}  % QLIKE change % vs the reference (subsample)
\newcommand{\rhabLassoFullLiveVdDm}{0.45}  % DM vs the reference arm (subsample)
\newcommand{\rhabLassoFullLiveVdP}{0.651}  % p-value
\newcommand{\rhabLassoLoVdLiveVdQlike}{0.1846}  % QLIKE lasso lo_vd live_vd (subsample)
\newcommand{\rhabLassoLoVdLiveVdPct}{$-$2.7}  % QLIKE change % vs the reference (subsample)
\newcommand{\rhabLassoLoVdLiveVdDm}{$-$2.51}  % DM vs the reference arm (subsample)
\newcommand{\rhabLassoLoVdLiveVdP}{0.012}  % p-value
\newcommand{\rhabLassoAddVdLiveVdQlike}{0.1941}  % QLIKE lasso add_vd live_vd (subsample)
\newcommand{\rhabLassoAddVdLiveVdPct}{4.2}  % QLIKE change % vs the reference (subsample)
\newcommand{\rhabLassoAddVdLiveVdDm}{2.49}  % DM vs the reference arm (subsample)
\newcommand{\rhabLassoAddVdLiveVdP}{0.013}  % p-value
\newcommand{\rhabLassoHarLiveVdQlike}{0.1864}  % QLIKE lasso har live_vd (subsample)
\newcommand{\rhabEnetFullLiveVdQlike}{0.1908}  % QLIKE enet full live_vd (subsample)
\newcommand{\rhabEnetFullLiveVdPct}{2.4}  % QLIKE change % vs the reference (subsample)
\newcommand{\rhabEnetFullLiveVdDm}{0.85}  % DM vs the reference arm (subsample)
\newcommand{\rhabEnetFullLiveVdP}{0.397}  % p-value
\newcommand{\rhabEnetLoVdLiveVdQlike}{0.1869}  % QLIKE enet lo_vd live_vd (subsample)
\newcommand{\rhabEnetLoVdLiveVdPct}{$-$2.1}  % QLIKE change % vs the reference (subsample)
\newcommand{\rhabEnetLoVdLiveVdDm}{$-$3.47}  % DM vs the reference arm (subsample)
\newcommand{\rhabEnetLoVdLiveVdP}{\ensuremath{<}0.001}  % p-value
\newcommand{\rhabEnetAddVdLiveVdQlike}{0.1916}  % QLIKE enet add_vd live_vd (subsample)
\newcommand{\rhabEnetAddVdLiveVdPct}{2.8}  % QLIKE change % vs the reference (subsample)
\newcommand{\rhabEnetAddVdLiveVdDm}{2.20}  % DM vs the reference arm (subsample)
\newcommand{\rhabEnetAddVdLiveVdP}{0.028}  % p-value
\newcommand{\rhabEnetHarLiveVdQlike}{0.1864}  % QLIKE enet har live_vd (subsample)
\newcommand{\rhabLgbmFullLiveVdQlike}{0.1612}  % QLIKE lgbm full live_vd (subsample)
\newcommand{\rhabLgbmFullLiveVdPct}{$-$10.9}  % QLIKE change % vs the reference (subsample)
\newcommand{\rhabLgbmFullLiveVdDm}{$-$2.01}  % DM vs the reference arm (subsample)
\newcommand{\rhabLgbmFullLiveVdP}{0.044}  % p-value
\newcommand{\rhabLgbmHarLiveVdQlike}{0.1809}  % QLIKE lgbm har live_vd (subsample)
\newcommand{\rhabXgbFullLiveVdQlike}{0.1616}  % QLIKE xgb full live_vd (subsample)
\newcommand{\rhabXgbFullLiveVdPct}{$-$12.1}  % QLIKE change % vs the reference (subsample)
\newcommand{\rhabXgbFullLiveVdDm}{$-$2.10}  % DM vs the reference arm (subsample)
\newcommand{\rhabXgbFullLiveVdP}{0.036}  % p-value
\newcommand{\rhabXgbHarLiveVdQlike}{0.1839}  % QLIKE xgb har live_vd (subsample)
\newcommand{\rhabRfFullLiveVdQlike}{0.1704}  % QLIKE rf full live_vd (subsample)
\newcommand{\rhabRfFullLiveVdPct}{$-$9.1}  % QLIKE change % vs the reference (subsample)
\newcommand{\rhabRfFullLiveVdDm}{$-$2.40}  % DM vs the reference arm (subsample)
\newcommand{\rhabRfFullLiveVdP}{0.016}  % p-value
\newcommand{\rhabRfHarLiveVdQlike}{0.1874}  % QLIKE rf har live_vd (subsample)
\newcommand{\rhabRidgeFullLiveVdCloseQlike}{0.1269}  % QLIKE ridge full live16_vd (subsample)
\newcommand{\rhabRidgeFullLiveVdClosePct}{45.8}  % QLIKE change % vs the reference (subsample)
\newcommand{\rhabRidgeFullLiveVdCloseDm}{3.73}  % DM vs the reference arm (subsample)
\newcommand{\rhabRidgeFullLiveVdCloseP}{\ensuremath{<}0.001}  % p-value
\newcommand{\rhabRidgeLoVdLiveVdCloseQlike}{0.1048}  % QLIKE ridge lo_vd live16_vd (subsample)
\newcommand{\rhabRidgeLoVdLiveVdClosePct}{$-$17.4}  % QLIKE change % vs the reference (subsample)
\newcommand{\rhabRidgeLoVdLiveVdCloseDm}{$-$4.97}  % DM vs the reference arm (subsample)
\newcommand{\rhabRidgeLoVdLiveVdCloseP}{\ensuremath{<}0.001}  % p-value
\newcommand{\rhabRidgeAddVdLiveVdCloseQlike}{0.1018}  % QLIKE ridge add_vd live16_vd (subsample)
\newcommand{\rhabRidgeAddVdLiveVdClosePct}{17.0}  % QLIKE change % vs the reference (subsample)
\newcommand{\rhabRidgeAddVdLiveVdCloseDm}{3.66}  % DM vs the reference arm (subsample)
\newcommand{\rhabRidgeAddVdLiveVdCloseP}{\ensuremath{<}0.001}  % p-value
\newcommand{\rhabRidgeHarLiveVdCloseQlike}{0.0870}  % QLIKE ridge har live16_vd (subsample)
\newcommand{\rhabLassoFullLiveVdCloseQlike}{0.1271}  % QLIKE lasso full live16_vd (subsample)
\newcommand{\rhabLassoFullLiveVdClosePct}{46.0}  % QLIKE change % vs the reference (subsample)
\newcommand{\rhabLassoFullLiveVdCloseDm}{3.87}  % DM vs the reference arm (subsample)
\newcommand{\rhabLassoFullLiveVdCloseP}{\ensuremath{<}0.001}  % p-value
\newcommand{\rhabLassoLoVdLiveVdCloseQlike}{0.1097}  % QLIKE lasso lo_vd live16_vd (subsample)
\newcommand{\rhabLassoLoVdLiveVdClosePct}{$-$13.7}  % QLIKE change % vs the reference (subsample)
\newcommand{\rhabLassoLoVdLiveVdCloseDm}{$-$4.18}  % DM vs the reference arm (subsample)
\newcommand{\rhabLassoLoVdLiveVdCloseP}{\ensuremath{<}0.001}  % p-value
\newcommand{\rhabLassoAddVdLiveVdCloseQlike}{0.1018}  % QLIKE lasso add_vd live16_vd (subsample)
\newcommand{\rhabLassoAddVdLiveVdClosePct}{17.0}  % QLIKE change % vs the reference (subsample)
\newcommand{\rhabLassoAddVdLiveVdCloseDm}{3.75}  % DM vs the reference arm (subsample)
\newcommand{\rhabLassoAddVdLiveVdCloseP}{\ensuremath{<}0.001}  % p-value
\newcommand{\rhabLassoHarLiveVdCloseQlike}{0.0870}  % QLIKE lasso har live16_vd (subsample)
\newcommand{\rhabEnetFullLiveVdCloseQlike}{0.1227}  % QLIKE enet full live16_vd (subsample)
\newcommand{\rhabEnetFullLiveVdClosePct}{41.0}  % QLIKE change % vs the reference (subsample)
\newcommand{\rhabEnetFullLiveVdCloseDm}{3.59}  % DM vs the reference arm (subsample)
\newcommand{\rhabEnetFullLiveVdCloseP}{\ensuremath{<}0.001}  % p-value
\newcommand{\rhabEnetLoVdLiveVdCloseQlike}{0.1029}  % QLIKE enet lo_vd live16_vd (subsample)
\newcommand{\rhabEnetLoVdLiveVdClosePct}{$-$16.2}  % QLIKE change % vs the reference (subsample)
\newcommand{\rhabEnetLoVdLiveVdCloseDm}{$-$4.00}  % DM vs the reference arm (subsample)
\newcommand{\rhabEnetLoVdLiveVdCloseP}{\ensuremath{<}0.001}  % p-value
\newcommand{\rhabEnetAddVdLiveVdCloseQlike}{0.1005}  % QLIKE enet add_vd live16_vd (subsample)
\newcommand{\rhabEnetAddVdLiveVdClosePct}{15.5}  % QLIKE change % vs the reference (subsample)
\newcommand{\rhabEnetAddVdLiveVdCloseDm}{3.55}  % DM vs the reference arm (subsample)
\newcommand{\rhabEnetAddVdLiveVdCloseP}{\ensuremath{<}0.001}  % p-value
\newcommand{\rhabEnetHarLiveVdCloseQlike}{0.0870}  % QLIKE enet har live16_vd (subsample)
\newcommand{\rhabLgbmFullLiveVdCloseQlike}{0.0802}  % QLIKE lgbm full live16_vd (subsample)
\newcommand{\rhabLgbmFullLiveVdClosePct}{$-$1.5}  % QLIKE change % vs the reference (subsample)
\newcommand{\rhabLgbmFullLiveVdCloseDm}{$-$0.29}  % DM vs the reference arm (subsample)
\newcommand{\rhabLgbmFullLiveVdCloseP}{0.769}  % p-value
\newcommand{\rhabLgbmHarLiveVdCloseQlike}{0.0814}  % QLIKE lgbm har live16_vd (subsample)
\newcommand{\rhabXgbFullLiveVdCloseQlike}{0.0785}  % QLIKE xgb full live16_vd (subsample)
\newcommand{\rhabXgbFullLiveVdClosePct}{$-$4.3}  % QLIKE change % vs the reference (subsample)
\newcommand{\rhabXgbFullLiveVdCloseDm}{$-$0.89}  % DM vs the reference arm (subsample)
\newcommand{\rhabXgbFullLiveVdCloseP}{0.375}  % p-value
\newcommand{\rhabXgbHarLiveVdCloseQlike}{0.0820}  % QLIKE xgb har live16_vd (subsample)
\newcommand{\rhabRfFullLiveVdCloseQlike}{0.0822}  % QLIKE rf full live16_vd (subsample)
\newcommand{\rhabRfFullLiveVdClosePct}{$-$1.0}  % QLIKE change % vs the reference (subsample)
\newcommand{\rhabRfFullLiveVdCloseDm}{$-$0.21}  % DM vs the reference arm (subsample)
\newcommand{\rhabRfFullLiveVdCloseP}{0.833}  % p-value
\newcommand{\rhabRfHarLiveVdCloseQlike}{0.0831}  % QLIKE rf har live16_vd (subsample)
\newcommand{\rhabLiveSentStart}{19~April 2022}  % first bar of the live window sent
\newcommand{\rhabLiveSentEnd}{30~April 2024}  % last bar of the live window sent
\newcommand{\rhabLiveSentDays}{508}  % trading days of the live window sent
\newcommand{\rhabLiveIvStart}{19~April 2018}  % first bar of the live window iv
\newcommand{\rhabLiveIvEnd}{30~April 2024}  % last bar of the live window iv
\newcommand{\rhabLiveIvDays}{1{,}506}  % trading days of the live window iv
\newcommand{\rhabLiveIvbStart}{6~July 2021}  % first bar of the live window ivb
\newcommand{\rhabLiveIvbEnd}{30~April 2024}  % last bar of the live window ivb
\newcommand{\rhabLiveIvbDays}{706}  % trading days of the live window ivb
\newcommand{\rhabLiveVdStart}{16~November 2022}  % first bar of the live window vd
\newcommand{\rhabLiveVdEnd}{31~August 2023}  % last bar of the live window vd
\newcommand{\rhabLiveVdDays}{196}  % trading days of the live window vd
\newcommand{\rhabRidgeFullVsOlsAddRmAllDm}{$-$4.70}  % DM ridge full vs OLS add_rm
\newcommand{\rhabRidgeFullVsOlsAddRmAllP}{\ensuremath{<}0.001}  % p-value
\newcommand{\rhabRidgeFullVsOlsAddRmCloseDm}{1.16}  % DM ridge full vs OLS add_rm
\newcommand{\rhabRidgeFullVsOlsAddRmCloseP}{0.244}  % p-value
\newcommand{\rhabRidgeFullVsOlsAddRmCloseLateDm}{0.95}  % DM ridge full vs OLS add_rm
\newcommand{\rhabRidgeFullVsOlsAddRmCloseLateP}{0.341}  % p-value
\newcommand{\rhabRidgeFullVsOlsSharAllDm}{$-$4.77}  % DM ridge full vs OLS shar
\newcommand{\rhabRidgeFullVsOlsSharAllP}{\ensuremath{<}0.001}  % p-value
\newcommand{\rhabRidgeFullVsOlsSharCloseDm}{0.41}  % DM ridge full vs OLS shar
\newcommand{\rhabRidgeFullVsOlsSharCloseP}{0.682}  % p-value
\newcommand{\rhabRidgeFullVsOlsSharCloseLateDm}{2.34}  % DM ridge full vs OLS shar
\newcommand{\rhabRidgeFullVsOlsSharCloseLateP}{0.019}  % p-value
\newcommand{\rhabRidgeFullVsOlsHarqAllDm}{$-$5.66}  % DM ridge full vs OLS harq
\newcommand{\rhabRidgeFullVsOlsHarqAllP}{\ensuremath{<}0.001}  % p-value
\newcommand{\rhabRidgeFullVsOlsHarqCloseDm}{0.37}  % DM ridge full vs OLS harq
\newcommand{\rhabRidgeFullVsOlsHarqCloseP}{0.715}  % p-value
\newcommand{\rhabRidgeFullVsOlsHarqCloseLateDm}{2.35}  % DM ridge full vs OLS harq
\newcommand{\rhabRidgeFullVsOlsHarqCloseLateP}{0.019}  % p-value
\newcommand{\rhabRidgeFullVsOlsLharAllDm}{$-$3.23}  % DM ridge full vs OLS lhar
\newcommand{\rhabRidgeFullVsOlsLharAllP}{0.001}  % p-value
\newcommand{\rhabRidgeFullVsOlsLharCloseDm}{2.40}  % DM ridge full vs OLS lhar
\newcommand{\rhabRidgeFullVsOlsLharCloseP}{0.016}  % p-value
\newcommand{\rhabRidgeFullVsOlsLharCloseLateDm}{2.96}  % DM ridge full vs OLS lhar
\newcommand{\rhabRidgeFullVsOlsLharCloseLateP}{0.003}  % p-value
\newcommand{\rhabRidgeFullVsOlsSharqlAllDm}{$-$3.04}  % DM ridge full vs OLS sharql
\newcommand{\rhabRidgeFullVsOlsSharqlAllP}{0.002}  % p-value
\newcommand{\rhabRidgeFullVsOlsSharqlCloseDm}{2.44}  % DM ridge full vs OLS sharql
\newcommand{\rhabRidgeFullVsOlsSharqlCloseP}{0.015}  % p-value
\newcommand{\rhabRidgeFullVsOlsSharqlCloseLateDm}{3.13}  % DM ridge full vs OLS sharql
\newcommand{\rhabRidgeFullVsOlsSharqlCloseLateP}{0.002}  % p-value
\newcommand{\rhabRidgeFullVsOlsHarAllDm}{$-$5.63}  % DM ridge full vs OLS har
\newcommand{\rhabRidgeFullVsOlsHarAllP}{\ensuremath{<}0.001}  % p-value
\newcommand{\rhabRidgeFullVsOlsHarCloseDm}{0.31}  % DM ridge full vs OLS har
\newcommand{\rhabRidgeFullVsOlsHarCloseP}{0.759}  % p-value
\newcommand{\rhabRidgeFullVsOlsHarCloseLateDm}{2.31}  % DM ridge full vs OLS har
\newcommand{\rhabRidgeFullVsOlsHarCloseLateP}{0.021}  % p-value
\newcommand{\rhabLassoFullVsOlsAddRmAllDm}{$-$4.53}  % DM lasso full vs OLS add_rm
\newcommand{\rhabLassoFullVsOlsAddRmAllP}{\ensuremath{<}0.001}  % p-value
\newcommand{\rhabLassoFullVsOlsAddRmCloseDm}{2.40}  % DM lasso full vs OLS add_rm
\newcommand{\rhabLassoFullVsOlsAddRmCloseP}{0.016}  % p-value
\newcommand{\rhabLassoFullVsOlsAddRmCloseLateDm}{2.26}  % DM lasso full vs OLS add_rm
\newcommand{\rhabLassoFullVsOlsAddRmCloseLateP}{0.024}  % p-value
\newcommand{\rhabLassoFullVsOlsSharAllDm}{$-$4.07}  % DM lasso full vs OLS shar
\newcommand{\rhabLassoFullVsOlsSharAllP}{\ensuremath{<}0.001}  % p-value
\newcommand{\rhabLassoFullVsOlsSharCloseDm}{1.12}  % DM lasso full vs OLS shar
\newcommand{\rhabLassoFullVsOlsSharCloseP}{0.261}  % p-value
\newcommand{\rhabLassoFullVsOlsSharCloseLateDm}{2.54}  % DM lasso full vs OLS shar
\newcommand{\rhabLassoFullVsOlsSharCloseLateP}{0.011}  % p-value
\newcommand{\rhabLassoFullVsOlsHarqAllDm}{$-$4.75}  % DM lasso full vs OLS harq
\newcommand{\rhabLassoFullVsOlsHarqAllP}{\ensuremath{<}0.001}  % p-value
\newcommand{\rhabLassoFullVsOlsHarqCloseDm}{1.08}  % DM lasso full vs OLS harq
\newcommand{\rhabLassoFullVsOlsHarqCloseP}{0.279}  % p-value
\newcommand{\rhabLassoFullVsOlsHarqCloseLateDm}{2.55}  % DM lasso full vs OLS harq
\newcommand{\rhabLassoFullVsOlsHarqCloseLateP}{0.011}  % p-value
\newcommand{\rhabLassoFullVsOlsLharAllDm}{$-$2.95}  % DM lasso full vs OLS lhar
\newcommand{\rhabLassoFullVsOlsLharAllP}{0.003}  % p-value
\newcommand{\rhabLassoFullVsOlsLharCloseDm}{3.00}  % DM lasso full vs OLS lhar
\newcommand{\rhabLassoFullVsOlsLharCloseP}{0.003}  % p-value
\newcommand{\rhabLassoFullVsOlsLharCloseLateDm}{3.28}  % DM lasso full vs OLS lhar
\newcommand{\rhabLassoFullVsOlsLharCloseLateP}{0.001}  % p-value
\newcommand{\rhabLassoFullVsOlsSharqlAllDm}{$-$2.81}  % DM lasso full vs OLS sharql
\newcommand{\rhabLassoFullVsOlsSharqlAllP}{0.005}  % p-value
\newcommand{\rhabLassoFullVsOlsSharqlCloseDm}{2.96}  % DM lasso full vs OLS sharql
\newcommand{\rhabLassoFullVsOlsSharqlCloseP}{0.003}  % p-value
\newcommand{\rhabLassoFullVsOlsSharqlCloseLateDm}{3.33}  % DM lasso full vs OLS sharql
\newcommand{\rhabLassoFullVsOlsSharqlCloseLateP}{\ensuremath{<}0.001}  % p-value
\newcommand{\rhabLassoFullVsOlsHarAllDm}{$-$4.72}  % DM lasso full vs OLS har
\newcommand{\rhabLassoFullVsOlsHarAllP}{\ensuremath{<}0.001}  % p-value
\newcommand{\rhabLassoFullVsOlsHarCloseDm}{1.04}  % DM lasso full vs OLS har
\newcommand{\rhabLassoFullVsOlsHarCloseP}{0.300}  % p-value
\newcommand{\rhabLassoFullVsOlsHarCloseLateDm}{2.51}  % DM lasso full vs OLS har
\newcommand{\rhabLassoFullVsOlsHarCloseLateP}{0.012}  % p-value
\newcommand{\rhabEnetFullVsOlsAddRmAllDm}{$-$5.61}  % DM enet full vs OLS add_rm
\newcommand{\rhabEnetFullVsOlsAddRmAllP}{\ensuremath{<}0.001}  % p-value
\newcommand{\rhabEnetFullVsOlsAddRmCloseDm}{0.46}  % DM enet full vs OLS add_rm
\newcommand{\rhabEnetFullVsOlsAddRmCloseP}{0.649}  % p-value
\newcommand{\rhabEnetFullVsOlsAddRmCloseLateDm}{0.29}  % DM enet full vs OLS add_rm
\newcommand{\rhabEnetFullVsOlsAddRmCloseLateP}{0.770}  % p-value
\newcommand{\rhabEnetFullVsOlsSharAllDm}{$-$5.97}  % DM enet full vs OLS shar
\newcommand{\rhabEnetFullVsOlsSharAllP}{\ensuremath{<}0.001}  % p-value
\newcommand{\rhabEnetFullVsOlsSharCloseDm}{$-$0.11}  % DM enet full vs OLS shar
\newcommand{\rhabEnetFullVsOlsSharCloseP}{0.914}  % p-value
\newcommand{\rhabEnetFullVsOlsSharCloseLateDm}{1.92}  % DM enet full vs OLS shar
\newcommand{\rhabEnetFullVsOlsSharCloseLateP}{0.055}  % p-value
\newcommand{\rhabEnetFullVsOlsHarqAllDm}{$-$7.03}  % DM enet full vs OLS harq
\newcommand{\rhabEnetFullVsOlsHarqAllP}{\ensuremath{<}0.001}  % p-value
\newcommand{\rhabEnetFullVsOlsHarqCloseDm}{$-$0.15}  % DM enet full vs OLS harq
\newcommand{\rhabEnetFullVsOlsHarqCloseP}{0.883}  % p-value
\newcommand{\rhabEnetFullVsOlsHarqCloseLateDm}{1.96}  % DM enet full vs OLS harq
\newcommand{\rhabEnetFullVsOlsHarqCloseLateP}{0.050}  % p-value
\newcommand{\rhabEnetFullVsOlsLharAllDm}{$-$4.58}  % DM enet full vs OLS lhar
\newcommand{\rhabEnetFullVsOlsLharAllP}{\ensuremath{<}0.001}  % p-value
\newcommand{\rhabEnetFullVsOlsLharCloseDm}{1.86}  % DM enet full vs OLS lhar
\newcommand{\rhabEnetFullVsOlsLharCloseP}{0.063}  % p-value
\newcommand{\rhabEnetFullVsOlsLharCloseLateDm}{2.46}  % DM enet full vs OLS lhar
\newcommand{\rhabEnetFullVsOlsLharCloseLateP}{0.014}  % p-value
\newcommand{\rhabEnetFullVsOlsSharqlAllDm}{$-$4.34}  % DM enet full vs OLS sharql
\newcommand{\rhabEnetFullVsOlsSharqlAllP}{\ensuremath{<}0.001}  % p-value
\newcommand{\rhabEnetFullVsOlsSharqlCloseDm}{1.92}  % DM enet full vs OLS sharql
\newcommand{\rhabEnetFullVsOlsSharqlCloseP}{0.055}  % p-value
\newcommand{\rhabEnetFullVsOlsSharqlCloseLateDm}{2.65}  % DM enet full vs OLS sharql
\newcommand{\rhabEnetFullVsOlsSharqlCloseLateP}{0.008}  % p-value
\newcommand{\rhabEnetFullVsOlsHarAllDm}{$-$6.98}  % DM enet full vs OLS har
\newcommand{\rhabEnetFullVsOlsHarAllP}{\ensuremath{<}0.001}  % p-value
\newcommand{\rhabEnetFullVsOlsHarCloseDm}{$-$0.21}  % DM enet full vs OLS har
\newcommand{\rhabEnetFullVsOlsHarCloseP}{0.832}  % p-value
\newcommand{\rhabEnetFullVsOlsHarCloseLateDm}{1.91}  % DM enet full vs OLS har
\newcommand{\rhabEnetFullVsOlsHarCloseLateP}{0.056}  % p-value
\newcommand{\rhabLgbmFullVsOlsAddRmAllDm}{$-$7.94}  % DM lgbm full vs OLS add_rm
\newcommand{\rhabLgbmFullVsOlsAddRmAllP}{\ensuremath{<}0.001}  % p-value
\newcommand{\rhabLgbmFullVsOlsAddRmCloseDm}{$-$1.78}  % DM lgbm full vs OLS add_rm
\newcommand{\rhabLgbmFullVsOlsAddRmCloseP}{0.076}  % p-value
\newcommand{\rhabLgbmFullVsOlsAddRmCloseLateDm}{$-$2.44}  % DM lgbm full vs OLS add_rm
\newcommand{\rhabLgbmFullVsOlsAddRmCloseLateP}{0.015}  % p-value
\newcommand{\rhabLgbmFullVsOlsSharAllDm}{$-$9.26}  % DM lgbm full vs OLS shar
\newcommand{\rhabLgbmFullVsOlsSharAllP}{\ensuremath{<}0.001}  % p-value
\newcommand{\rhabLgbmFullVsOlsSharCloseDm}{$-$3.86}  % DM lgbm full vs OLS shar
\newcommand{\rhabLgbmFullVsOlsSharCloseP}{\ensuremath{<}0.001}  % p-value
\newcommand{\rhabLgbmFullVsOlsSharCloseLateDm}{$-$2.34}  % DM lgbm full vs OLS shar
\newcommand{\rhabLgbmFullVsOlsSharCloseLateP}{0.019}  % p-value
\newcommand{\rhabLgbmFullVsOlsHarqAllDm}{$-$10.44}  % DM lgbm full vs OLS harq
\newcommand{\rhabLgbmFullVsOlsHarqAllP}{\ensuremath{<}0.001}  % p-value
\newcommand{\rhabLgbmFullVsOlsHarqCloseDm}{$-$3.77}  % DM lgbm full vs OLS harq
\newcommand{\rhabLgbmFullVsOlsHarqCloseP}{\ensuremath{<}0.001}  % p-value
\newcommand{\rhabLgbmFullVsOlsHarqCloseLateDm}{$-$1.70}  % DM lgbm full vs OLS harq
\newcommand{\rhabLgbmFullVsOlsHarqCloseLateP}{0.088}  % p-value
\newcommand{\rhabLgbmFullVsOlsLharAllDm}{$-$8.32}  % DM lgbm full vs OLS lhar
\newcommand{\rhabLgbmFullVsOlsLharAllP}{\ensuremath{<}0.001}  % p-value
\newcommand{\rhabLgbmFullVsOlsLharCloseDm}{$-$1.02}  % DM lgbm full vs OLS lhar
\newcommand{\rhabLgbmFullVsOlsLharCloseP}{0.310}  % p-value
\newcommand{\rhabLgbmFullVsOlsLharCloseLateDm}{$-$1.42}  % DM lgbm full vs OLS lhar
\newcommand{\rhabLgbmFullVsOlsLharCloseLateP}{0.155}  % p-value
\newcommand{\rhabLgbmFullVsOlsSharqlAllDm}{$-$8.22}  % DM lgbm full vs OLS sharql
\newcommand{\rhabLgbmFullVsOlsSharqlAllP}{\ensuremath{<}0.001}  % p-value
\newcommand{\rhabLgbmFullVsOlsSharqlCloseDm}{$-$1.03}  % DM lgbm full vs OLS sharql
\newcommand{\rhabLgbmFullVsOlsSharqlCloseP}{0.302}  % p-value
\newcommand{\rhabLgbmFullVsOlsSharqlCloseLateDm}{$-$1.45}  % DM lgbm full vs OLS sharql
\newcommand{\rhabLgbmFullVsOlsSharqlCloseLateP}{0.148}  % p-value
\newcommand{\rhabLgbmFullVsOlsHarAllDm}{$-$10.41}  % DM lgbm full vs OLS har
\newcommand{\rhabLgbmFullVsOlsHarAllP}{\ensuremath{<}0.001}  % p-value
\newcommand{\rhabLgbmFullVsOlsHarCloseDm}{$-$3.89}  % DM lgbm full vs OLS har
\newcommand{\rhabLgbmFullVsOlsHarCloseP}{\ensuremath{<}0.001}  % p-value
\newcommand{\rhabLgbmFullVsOlsHarCloseLateDm}{$-$1.86}  % DM lgbm full vs OLS har
\newcommand{\rhabLgbmFullVsOlsHarCloseLateP}{0.063}  % p-value
\newcommand{\rhabXgbFullVsOlsAddRmAllDm}{$-$7.47}  % DM xgb full vs OLS add_rm
\newcommand{\rhabXgbFullVsOlsAddRmAllP}{\ensuremath{<}0.001}  % p-value
\newcommand{\rhabXgbFullVsOlsAddRmCloseDm}{$-$2.16}  % DM xgb full vs OLS add_rm
\newcommand{\rhabXgbFullVsOlsAddRmCloseP}{0.031}  % p-value
\newcommand{\rhabXgbFullVsOlsAddRmCloseLateDm}{$-$2.61}  % DM xgb full vs OLS add_rm
\newcommand{\rhabXgbFullVsOlsAddRmCloseLateP}{0.009}  % p-value
\newcommand{\rhabXgbFullVsOlsSharAllDm}{$-$8.81}  % DM xgb full vs OLS shar
\newcommand{\rhabXgbFullVsOlsSharAllP}{\ensuremath{<}0.001}  % p-value
\newcommand{\rhabXgbFullVsOlsSharCloseDm}{$-$4.87}  % DM xgb full vs OLS shar
\newcommand{\rhabXgbFullVsOlsSharCloseP}{\ensuremath{<}0.001}  % p-value
\newcommand{\rhabXgbFullVsOlsSharCloseLateDm}{$-$2.99}  % DM xgb full vs OLS shar
\newcommand{\rhabXgbFullVsOlsSharCloseLateP}{0.003}  % p-value
\newcommand{\rhabXgbFullVsOlsHarqAllDm}{$-$9.96}  % DM xgb full vs OLS harq
\newcommand{\rhabXgbFullVsOlsHarqAllP}{\ensuremath{<}0.001}  % p-value
\newcommand{\rhabXgbFullVsOlsHarqCloseDm}{$-$4.77}  % DM xgb full vs OLS harq
\newcommand{\rhabXgbFullVsOlsHarqCloseP}{\ensuremath{<}0.001}  % p-value
\newcommand{\rhabXgbFullVsOlsHarqCloseLateDm}{$-$2.27}  % DM xgb full vs OLS harq
\newcommand{\rhabXgbFullVsOlsHarqCloseLateP}{0.023}  % p-value
\newcommand{\rhabXgbFullVsOlsLharAllDm}{$-$7.85}  % DM xgb full vs OLS lhar
\newcommand{\rhabXgbFullVsOlsLharAllP}{\ensuremath{<}0.001}  % p-value
\newcommand{\rhabXgbFullVsOlsLharCloseDm}{$-$1.70}  % DM xgb full vs OLS lhar
\newcommand{\rhabXgbFullVsOlsLharCloseP}{0.090}  % p-value
\newcommand{\rhabXgbFullVsOlsLharCloseLateDm}{$-$1.85}  % DM xgb full vs OLS lhar
\newcommand{\rhabXgbFullVsOlsLharCloseLateP}{0.065}  % p-value
\newcommand{\rhabXgbFullVsOlsSharqlAllDm}{$-$7.76}  % DM xgb full vs OLS sharql
\newcommand{\rhabXgbFullVsOlsSharqlAllP}{\ensuremath{<}0.001}  % p-value
\newcommand{\rhabXgbFullVsOlsSharqlCloseDm}{$-$1.82}  % DM xgb full vs OLS sharql
\newcommand{\rhabXgbFullVsOlsSharqlCloseP}{0.069}  % p-value
\newcommand{\rhabXgbFullVsOlsSharqlCloseLateDm}{$-$1.95}  % DM xgb full vs OLS sharql
\newcommand{\rhabXgbFullVsOlsSharqlCloseLateP}{0.051}  % p-value
\newcommand{\rhabXgbFullVsOlsHarAllDm}{$-$9.93}  % DM xgb full vs OLS har
\newcommand{\rhabXgbFullVsOlsHarAllP}{\ensuremath{<}0.001}  % p-value
\newcommand{\rhabXgbFullVsOlsHarCloseDm}{$-$4.88}  % DM xgb full vs OLS har
\newcommand{\rhabXgbFullVsOlsHarCloseP}{\ensuremath{<}0.001}  % p-value
\newcommand{\rhabXgbFullVsOlsHarCloseLateDm}{$-$2.44}  % DM xgb full vs OLS har
\newcommand{\rhabXgbFullVsOlsHarCloseLateP}{0.015}  % p-value
\newcommand{\rhabRfFullVsOlsAddRmAllDm}{$-$6.68}  % DM rf full vs OLS add_rm
\newcommand{\rhabRfFullVsOlsAddRmAllP}{\ensuremath{<}0.001}  % p-value
\newcommand{\rhabRfFullVsOlsAddRmCloseDm}{$-$0.14}  % DM rf full vs OLS add_rm
\newcommand{\rhabRfFullVsOlsAddRmCloseP}{0.892}  % p-value
\newcommand{\rhabRfFullVsOlsAddRmCloseLateDm}{$-$1.75}  % DM rf full vs OLS add_rm
\newcommand{\rhabRfFullVsOlsAddRmCloseLateP}{0.080}  % p-value
\newcommand{\rhabRfFullVsOlsSharAllDm}{$-$8.45}  % DM rf full vs OLS shar
\newcommand{\rhabRfFullVsOlsSharAllP}{\ensuremath{<}0.001}  % p-value
\newcommand{\rhabRfFullVsOlsSharCloseDm}{$-$0.98}  % DM rf full vs OLS shar
\newcommand{\rhabRfFullVsOlsSharCloseP}{0.328}  % p-value
\newcommand{\rhabRfFullVsOlsSharCloseLateDm}{$-$0.76}  % DM rf full vs OLS shar
\newcommand{\rhabRfFullVsOlsSharCloseLateP}{0.444}  % p-value
\newcommand{\rhabRfFullVsOlsHarqAllDm}{$-$9.73}  % DM rf full vs OLS harq
\newcommand{\rhabRfFullVsOlsHarqAllP}{\ensuremath{<}0.001}  % p-value
\newcommand{\rhabRfFullVsOlsHarqCloseDm}{$-$1.05}  % DM rf full vs OLS harq
\newcommand{\rhabRfFullVsOlsHarqCloseP}{0.293}  % p-value
\newcommand{\rhabRfFullVsOlsHarqCloseLateDm}{$-$0.36}  % DM rf full vs OLS harq
\newcommand{\rhabRfFullVsOlsHarqCloseLateP}{0.715}  % p-value
\newcommand{\rhabRfFullVsOlsLharAllDm}{$-$7.08}  % DM rf full vs OLS lhar
\newcommand{\rhabRfFullVsOlsLharAllP}{\ensuremath{<}0.001}  % p-value
\newcommand{\rhabRfFullVsOlsLharCloseDm}{1.54}  % DM rf full vs OLS lhar
\newcommand{\rhabRfFullVsOlsLharCloseP}{0.124}  % p-value
\newcommand{\rhabRfFullVsOlsLharCloseLateDm}{$-$0.26}  % DM rf full vs OLS lhar
\newcommand{\rhabRfFullVsOlsLharCloseLateP}{0.794}  % p-value
\newcommand{\rhabRfFullVsOlsSharqlAllDm}{$-$6.99}  % DM rf full vs OLS sharql
\newcommand{\rhabRfFullVsOlsSharqlAllP}{\ensuremath{<}0.001}  % p-value
\newcommand{\rhabRfFullVsOlsSharqlCloseDm}{1.76}  % DM rf full vs OLS sharql
\newcommand{\rhabRfFullVsOlsSharqlCloseP}{0.078}  % p-value
\newcommand{\rhabRfFullVsOlsSharqlCloseLateDm}{$-$0.09}  % DM rf full vs OLS sharql
\newcommand{\rhabRfFullVsOlsSharqlCloseLateP}{0.926}  % p-value
\newcommand{\rhabRfFullVsOlsHarAllDm}{$-$9.69}  % DM rf full vs OLS har
\newcommand{\rhabRfFullVsOlsHarAllP}{\ensuremath{<}0.001}  % p-value
\newcommand{\rhabRfFullVsOlsHarCloseDm}{$-$1.16}  % DM rf full vs OLS har
\newcommand{\rhabRfFullVsOlsHarCloseP}{0.246}  % p-value
\newcommand{\rhabRfFullVsOlsHarCloseLateDm}{$-$0.49}  % DM rf full vs OLS har
\newcommand{\rhabRfFullVsOlsHarCloseLateP}{0.621}  % p-value
\newcommand{\rhabNClassesAddRm}{three}  % classes with an add_rm arm
\newcommand{\rhabNClassesBeatAddRm}{none}  % full set beats its own add_rm, DM p < 0.05
\newcommand{\rhabBestOlsBench}{SHAR + HARQ + LHAR}  % lowest QLIKE OLS augmented HAR
\newcommand{\rhabLiveAttnFirst}{23~March 2009}  % first print stocktwits_attention
\newcommand{\rhabLiveAttnLast}{30~April 2024}  % last print stocktwits_attention
\newcommand{\rhabLiveAttnFromKone}{11~February 2021}  % first non-zeroed value, k = 1
\newcommand{\rhabLiveAttnFromKlast}{13~May 2021}  % the same, k = 3125
\newcommand{\rhabLiveAttnSharePct}{31}  % % of evaluation bars live, k = 1
\newcommand{\rhabLiveNSentFirst}{1~June 2010}  % first print stocktwits_sentcount
\newcommand{\rhabLiveNSentLast}{30~April 2024}  % last print stocktwits_sentcount
\newcommand{\rhabLiveNSentFromKone}{19~April 2022}  % first non-zeroed value, k = 1
\newcommand{\rhabLiveNSentFromKlast}{19~July 2022}  % the same, k = 3125
\newcommand{\rhabLiveNSentSharePct}{20}  % % of evaluation bars live, k = 1
\newcommand{\rhabLiveSentFirst}{1~June 2010}  % first print stocktwits_sentiment
\newcommand{\rhabLiveSentLast}{30~April 2024}  % last print stocktwits_sentiment
\newcommand{\rhabLiveSentFromKone}{19~April 2022}  % first non-zeroed value, k = 1
\newcommand{\rhabLiveSentFromKlast}{19~July 2022}  % the same, k = 3125
\newcommand{\rhabLiveSentSharePct}{20}  % % of evaluation bars live, k = 1
\newcommand{\rhabLiveVixFirst}{22~May 2006}  % first print vix
\newcommand{\rhabLiveVixLast}{30~April 2024}  % last print vix
\newcommand{\rhabLiveVixFromKone}{19~April 2018}  % first non-zeroed value, k = 1
\newcommand{\rhabLiveVixFromKlast}{18~July 2018}  % the same, k = 3125
\newcommand{\rhabLiveVixSharePct}{58}  % % of evaluation bars live, k = 1
\newcommand{\rhabLiveVixthreemFirst}{12~August 2009}  % first print vix3m
\newcommand{\rhabLiveVixthreemLast}{30~April 2024}  % last print vix3m
\newcommand{\rhabLiveVixthreemFromKone}{6~July 2021}  % first non-zeroed value, k = 1
\newcommand{\rhabLiveVixthreemFromKlast}{1~October 2021}  % the same, k = 3125
\newcommand{\rhabLiveVixthreemSharePct}{27}  % % of evaluation bars live, k = 1
\newcommand{\rhabLiveVvixFirst}{26~March 2012}  % first print vvix
\newcommand{\rhabLiveVvixLast}{30~April 2024}  % last print vvix
\newcommand{\rhabLiveVvixFromKone}{12~February 2024}  % first non-zeroed value, k = 1
\newcommand{\rhabLiveVvixFromKlast}{never}  % the same, k = 3125
\newcommand{\rhabLiveVvixSharePct}{2}  % % of evaluation bars live, k = 1
\newcommand{\rhabLiveVdAllOcFirst}{3~January 2011}  % first print voldemand_all_open_and_close
\newcommand{\rhabLiveVdAllOcLast}{31~August 2023}  % last print voldemand_all_open_and_close
\newcommand{\rhabLiveVdAllOcFromKone}{16~November 2022}  % first non-zeroed value, k = 1
\newcommand{\rhabLiveVdAllOcFromKlast}{20~February 2023}  % the same, k = 3125
\newcommand{\rhabLiveVdAllOcSharePct}{8}  % % of evaluation bars live, k = 1
\newcommand{\rhabLiveVdAllOFirst}{3~January 2011}  % first print voldemand_all_open_only
\newcommand{\rhabLiveVdAllOLast}{31~August 2023}  % last print voldemand_all_open_only
\newcommand{\rhabLiveVdAllOFromKone}{16~November 2022}  % first non-zeroed value, k = 1
\newcommand{\rhabLiveVdAllOFromKlast}{20~February 2023}  % the same, k = 3125
\newcommand{\rhabLiveVdAllOSharePct}{8}  % % of evaluation bars live, k = 1
\newcommand{\rhabLiveVdSpxOcFirst}{3~January 2011}  % first print voldemand_spx_open_and_close
\newcommand{\rhabLiveVdSpxOcLast}{31~August 2023}  % last print voldemand_spx_open_and_close
\newcommand{\rhabLiveVdSpxOcFromKone}{16~November 2022}  % first non-zeroed value, k = 1
\newcommand{\rhabLiveVdSpxOcFromKlast}{20~February 2023}  % the same, k = 3125
\newcommand{\rhabLiveVdSpxOcSharePct}{8}  % % of evaluation bars live, k = 1
\newcommand{\rhabLiveVdSpxOFirst}{3~January 2011}  % first print voldemand_spx_open_only
\newcommand{\rhabLiveVdSpxOLast}{31~August 2023}  % last print voldemand_spx_open_only
\newcommand{\rhabLiveVdSpxOFromKone}{16~November 2022}  % first non-zeroed value, k = 1
\newcommand{\rhabLiveVdSpxOFromKlast}{20~February 2023}  % the same, k = 3125
\newcommand{\rhabLiveVdSpxOSharePct}{8}  % % of evaluation bars live, k = 1
\fi
\ifdefined\rhcxWMain\else% AUTO-GENERATED by writeup/make_rehaul_complexity.py -- do not edit; re-run the script.
% \rhcx... macros of the draft complexity / scaling subsection (writeup/drafts/2026-10-26_complexity.tex): the
% window x feature-count grid of results/ladder_2026-10-26/grid (RULE.txt ADDENDUM 1).  Names: Gain<Model><Span>
% <W><k> (QLIKE gain over the same window's k = 0 rung, x 10^3), Kstar / Kninety (best k, k at 90 % of the full
% set's gain), ...Read (the reading across windows), Model Lgbm Xgb Ridge Lasso Enet, Span Bar (every bar) / Close
% (16:00), W One Two Main Long, k One Three Five Ten Twenty Full.
% PRE-PATCH design: the FOMC calendar features before the pending FOMC fix (RULE.txt ADDENDUM 2).
\providecommand{\rhcxWOne}{}\renewcommand{\rhcxWOne}{1{,}000}  % rolling window, trading days (13,000 bars)
\providecommand{\rhcxTOne}{}\renewcommand{\rhcxTOne}{13{,}000}  % window length in bars, T = 13 W
\providecommand{\rhcxWTwo}{}\renewcommand{\rhcxWTwo}{2{,}000}  % rolling window, trading days (26,000 bars)
\providecommand{\rhcxTTwo}{}\renewcommand{\rhcxTTwo}{26{,}000}  % window length in bars, T = 13 W
\providecommand{\rhcxWMain}{}\renewcommand{\rhcxWMain}{3{,}000}  % rolling window, trading days (39,000 bars)
\providecommand{\rhcxTMain}{}\renewcommand{\rhcxTMain}{39{,}000}  % window length in bars, T = 13 W
\providecommand{\rhcxWLong}{}\renewcommand{\rhcxWLong}{3{,}500}  % rolling window, trading days (45,500 bars)
\providecommand{\rhcxTLong}{}\renewcommand{\rhcxTLong}{45{,}500}  % window length in bars, T = 13 W
\providecommand{\rhcxNBars}{}\renewcommand{\rhcxNBars}{33{,}527}  % scored bars (Table 1's), every window
\providecommand{\rhcxNDays}{}\renewcommand{\rhcxNDays}{2{,}579}  % evaluation days
\providecommand{\rhcxNOrigins}{}\renewcommand{\rhcxNOrigins}{126}  % monthly origins (refit, held 21 sessions)
\providecommand{\rhcxNSeries}{}\renewcommand{\rhcxNSeries}{46}  % exogenous series
\providecommand{\rhcxPFull}{}\renewcommand{\rhcxPFull}{564}  % predictors of the full set (12 + 12 x 46)
\providecommand{\rhcxPHar}{}\renewcommand{\rhcxPHar}{12}  % predictors of the k = 0 rung (HAR benchmark set)
\providecommand{\rhcxCMax}{}\renewcommand{\rhcxCMax}{0.043}  % largest complexity c = P / T (full set, W = 1,000)
\providecommand{\rhcxCMin}{}\renewcommand{\rhcxCMin}{\ensuremath{2.6\times 10^{-4}}}  % smallest c (k = 0, W = 3,500)
\providecommand{\rhcxCMaxMain}{}\renewcommand{\rhcxCMaxMain}{0.014}  % c of the full set at W = 3,000
\providecommand{\rhcxNGates}{}\renewcommand{\rhcxNGates}{20}  % gates (window x model)
\providecommand{\rhcxNGatesPass}{}\renewcommand{\rhcxNGatesPass}{20}  % gates passed
\providecommand{\rhcxHacLag}{}\renewcommand{\rhcxHacLag}{8}  % Newey-West lag of the DM tests
\providecommand{\rhcxPenRidgeOne}{}\renewcommand{\rhcxPenRidgeOne}{\ensuremath{10^{4}}}  % ridge frozen penalty at W 1000 (all_features)
\providecommand{\rhcxPenLassoOne}{}\renewcommand{\rhcxPenLassoOne}{\ensuremath{10^{-3}}}  % lasso frozen penalty at W 1000 (all_features)
\providecommand{\rhcxPenEnetOne}{}\renewcommand{\rhcxPenEnetOne}{\ensuremath{10^{-2}}}  % enet frozen penalty at W 1000 (all_features)
\providecommand{\rhcxPenRidgeTwo}{}\renewcommand{\rhcxPenRidgeTwo}{\ensuremath{10^{3}}}  % ridge frozen penalty at W 2000 (all_features)
\providecommand{\rhcxPenLassoTwo}{}\renewcommand{\rhcxPenLassoTwo}{\ensuremath{10^{-4}}}  % lasso frozen penalty at W 2000 (all_features)
\providecommand{\rhcxPenEnetTwo}{}\renewcommand{\rhcxPenEnetTwo}{\ensuremath{10^{-3}}}  % enet frozen penalty at W 2000 (all_features)
\providecommand{\rhcxPenRidgeMain}{}\renewcommand{\rhcxPenRidgeMain}{\ensuremath{10^{3}}}  % ridge frozen penalty at W 3000 (all_features)
\providecommand{\rhcxPenLassoMain}{}\renewcommand{\rhcxPenLassoMain}{\ensuremath{10^{-4}}}  % lasso frozen penalty at W 3000 (all_features)
\providecommand{\rhcxPenEnetMain}{}\renewcommand{\rhcxPenEnetMain}{\ensuremath{10^{-3}}}  % enet frozen penalty at W 3000 (all_features)
\providecommand{\rhcxPenRidgeLong}{}\renewcommand{\rhcxPenRidgeLong}{\ensuremath{10^{3}}}  % ridge frozen penalty at W 3500 (all_features)
\providecommand{\rhcxPenLassoLong}{}\renewcommand{\rhcxPenLassoLong}{\ensuremath{10^{-4}}}  % lasso frozen penalty at W 3500 (all_features)
\providecommand{\rhcxPenEnetLong}{}\renewcommand{\rhcxPenEnetLong}{\ensuremath{10^{-3}}}  % enet frozen penalty at W 3500 (all_features)
\providecommand{\rhcxKstarLgbmBarOne}{}\renewcommand{\rhcxKstarLgbmBarOne}{46}  % lgbm all W 1000: rung with the lowest QLIKE
\providecommand{\rhcxKninetyLgbmBarOne}{}\renewcommand{\rhcxKninetyLgbmBarOne}{20}  % lgbm all W 1000: smallest k with 90% of the full set's gain
\providecommand{\rhcxGainLgbmBarOneOne}{}\renewcommand{\rhcxGainLgbmBarOneOne}{\ensuremath{-}0.92}  % lgbm all W 1000 k 1: QLIKE gain over k = 0, x 10^3
\providecommand{\rhcxGainLgbmBarOneThree}{}\renewcommand{\rhcxGainLgbmBarOneThree}{3.45}  % lgbm all W 1000 k 3: QLIKE gain over k = 0, x 10^3
\providecommand{\rhcxGainLgbmBarOneFive}{}\renewcommand{\rhcxGainLgbmBarOneFive}{6.06}  % lgbm all W 1000 k 5: QLIKE gain over k = 0, x 10^3
\providecommand{\rhcxGainLgbmBarOneTen}{}\renewcommand{\rhcxGainLgbmBarOneTen}{6.00}  % lgbm all W 1000 k 10: QLIKE gain over k = 0, x 10^3
\providecommand{\rhcxGainLgbmBarOneTwenty}{}\renewcommand{\rhcxGainLgbmBarOneTwenty}{15.27}  % lgbm all W 1000 k 20: QLIKE gain over k = 0, x 10^3
\providecommand{\rhcxGainLgbmBarOneFull}{}\renewcommand{\rhcxGainLgbmBarOneFull}{15.85}  % lgbm all W 1000 k 46: QLIKE gain over k = 0, x 10^3
\providecommand{\rhcxGainFullPLgbmBarOne}{}\renewcommand{\rhcxGainFullPLgbmBarOne}{p < 0.001}  % lgbm all W 1000: DM p, full set vs k = 0
\providecommand{\rhcxGainBestPLgbmBarOne}{}\renewcommand{\rhcxGainBestPLgbmBarOne}{p < 0.001}  % lgbm all W 1000: DM p, best rung vs k = 0
\providecommand{\rhcxKstarLgbmBarTwo}{}\renewcommand{\rhcxKstarLgbmBarTwo}{46}  % lgbm all W 2000: rung with the lowest QLIKE
\providecommand{\rhcxKninetyLgbmBarTwo}{}\renewcommand{\rhcxKninetyLgbmBarTwo}{10}  % lgbm all W 2000: smallest k with 90% of the full set's gain
\providecommand{\rhcxGainLgbmBarTwoOne}{}\renewcommand{\rhcxGainLgbmBarTwoOne}{0.25}  % lgbm all W 2000 k 1: QLIKE gain over k = 0, x 10^3
\providecommand{\rhcxGainLgbmBarTwoThree}{}\renewcommand{\rhcxGainLgbmBarTwoThree}{5.22}  % lgbm all W 2000 k 3: QLIKE gain over k = 0, x 10^3
\providecommand{\rhcxGainLgbmBarTwoFive}{}\renewcommand{\rhcxGainLgbmBarTwoFive}{7.21}  % lgbm all W 2000 k 5: QLIKE gain over k = 0, x 10^3
\providecommand{\rhcxGainLgbmBarTwoTen}{}\renewcommand{\rhcxGainLgbmBarTwoTen}{22.21}  % lgbm all W 2000 k 10: QLIKE gain over k = 0, x 10^3
\providecommand{\rhcxGainLgbmBarTwoTwenty}{}\renewcommand{\rhcxGainLgbmBarTwoTwenty}{23.05}  % lgbm all W 2000 k 20: QLIKE gain over k = 0, x 10^3
\providecommand{\rhcxGainLgbmBarTwoFull}{}\renewcommand{\rhcxGainLgbmBarTwoFull}{23.08}  % lgbm all W 2000 k 46: QLIKE gain over k = 0, x 10^3
\providecommand{\rhcxGainFullPLgbmBarTwo}{}\renewcommand{\rhcxGainFullPLgbmBarTwo}{p < 0.001}  % lgbm all W 2000: DM p, full set vs k = 0
\providecommand{\rhcxGainBestPLgbmBarTwo}{}\renewcommand{\rhcxGainBestPLgbmBarTwo}{p < 0.001}  % lgbm all W 2000: DM p, best rung vs k = 0
\providecommand{\rhcxKstarLgbmBarMain}{}\renewcommand{\rhcxKstarLgbmBarMain}{46}  % lgbm all W 3000: rung with the lowest QLIKE
\providecommand{\rhcxKninetyLgbmBarMain}{}\renewcommand{\rhcxKninetyLgbmBarMain}{10}  % lgbm all W 3000: smallest k with 90% of the full set's gain
\providecommand{\rhcxGainLgbmBarMainOne}{}\renewcommand{\rhcxGainLgbmBarMainOne}{0.89}  % lgbm all W 3000 k 1: QLIKE gain over k = 0, x 10^3
\providecommand{\rhcxGainLgbmBarMainThree}{}\renewcommand{\rhcxGainLgbmBarMainThree}{8.58}  % lgbm all W 3000 k 3: QLIKE gain over k = 0, x 10^3
\providecommand{\rhcxGainLgbmBarMainFive}{}\renewcommand{\rhcxGainLgbmBarMainFive}{8.73}  % lgbm all W 3000 k 5: QLIKE gain over k = 0, x 10^3
\providecommand{\rhcxGainLgbmBarMainTen}{}\renewcommand{\rhcxGainLgbmBarMainTen}{24.08}  % lgbm all W 3000 k 10: QLIKE gain over k = 0, x 10^3
\providecommand{\rhcxGainLgbmBarMainTwenty}{}\renewcommand{\rhcxGainLgbmBarMainTwenty}{24.46}  % lgbm all W 3000 k 20: QLIKE gain over k = 0, x 10^3
\providecommand{\rhcxGainLgbmBarMainFull}{}\renewcommand{\rhcxGainLgbmBarMainFull}{25.47}  % lgbm all W 3000 k 46: QLIKE gain over k = 0, x 10^3
\providecommand{\rhcxGainFullPLgbmBarMain}{}\renewcommand{\rhcxGainFullPLgbmBarMain}{p < 0.001}  % lgbm all W 3000: DM p, full set vs k = 0
\providecommand{\rhcxGainBestPLgbmBarMain}{}\renewcommand{\rhcxGainBestPLgbmBarMain}{p < 0.001}  % lgbm all W 3000: DM p, best rung vs k = 0
\providecommand{\rhcxKstarLgbmBarLong}{}\renewcommand{\rhcxKstarLgbmBarLong}{46}  % lgbm all W 3500: rung with the lowest QLIKE
\providecommand{\rhcxKninetyLgbmBarLong}{}\renewcommand{\rhcxKninetyLgbmBarLong}{10}  % lgbm all W 3500: smallest k with 90% of the full set's gain
\providecommand{\rhcxGainLgbmBarLongOne}{}\renewcommand{\rhcxGainLgbmBarLongOne}{0.42}  % lgbm all W 3500 k 1: QLIKE gain over k = 0, x 10^3
\providecommand{\rhcxGainLgbmBarLongThree}{}\renewcommand{\rhcxGainLgbmBarLongThree}{8.22}  % lgbm all W 3500 k 3: QLIKE gain over k = 0, x 10^3
\providecommand{\rhcxGainLgbmBarLongFive}{}\renewcommand{\rhcxGainLgbmBarLongFive}{7.78}  % lgbm all W 3500 k 5: QLIKE gain over k = 0, x 10^3
\providecommand{\rhcxGainLgbmBarLongTen}{}\renewcommand{\rhcxGainLgbmBarLongTen}{24.67}  % lgbm all W 3500 k 10: QLIKE gain over k = 0, x 10^3
\providecommand{\rhcxGainLgbmBarLongTwenty}{}\renewcommand{\rhcxGainLgbmBarLongTwenty}{24.81}  % lgbm all W 3500 k 20: QLIKE gain over k = 0, x 10^3
\providecommand{\rhcxGainLgbmBarLongFull}{}\renewcommand{\rhcxGainLgbmBarLongFull}{25.21}  % lgbm all W 3500 k 46: QLIKE gain over k = 0, x 10^3
\providecommand{\rhcxGainFullPLgbmBarLong}{}\renewcommand{\rhcxGainFullPLgbmBarLong}{p < 0.001}  % lgbm all W 3500: DM p, full set vs k = 0
\providecommand{\rhcxGainBestPLgbmBarLong}{}\renewcommand{\rhcxGainBestPLgbmBarLong}{p < 0.001}  % lgbm all W 3500: DM p, best rung vs k = 0
\providecommand{\rhcxKstarReadLgbmBar}{}\renewcommand{\rhcxKstarReadLgbmBar}{is flat}  % lgbm all: best k across W (RULE add. 1 (7))
\providecommand{\rhcxKninetyReadLgbmBar}{}\renewcommand{\rhcxKninetyReadLgbmBar}{falls}  % lgbm all: k90 across W
\providecommand{\rhcxFullBestLgbmBar}{}\renewcommand{\rhcxFullBestLgbmBar}{4}  % lgbm all: windows (of 4) where k = 46 is best
\providecommand{\rhcxSlopePosLgbmBar}{}\renewcommand{\rhcxSlopePosLgbmBar}{4}  % lgbm all: windows with a positive gain-log(c) slope
\providecommand{\rhcxKstarLgbmCloseOne}{}\renewcommand{\rhcxKstarLgbmCloseOne}{5}  % lgbm c16 W 1000: rung with the lowest QLIKE
\providecommand{\rhcxKninetyLgbmCloseOne}{}\renewcommand{\rhcxKninetyLgbmCloseOne}{5}  % lgbm c16 W 1000: smallest k with 90% of the full set's gain
\providecommand{\rhcxGainLgbmCloseOneOne}{}\renewcommand{\rhcxGainLgbmCloseOneOne}{\ensuremath{-}1.95}  % lgbm c16 W 1000 k 1: QLIKE gain over k = 0, x 10^3
\providecommand{\rhcxGainLgbmCloseOneThree}{}\renewcommand{\rhcxGainLgbmCloseOneThree}{1.45}  % lgbm c16 W 1000 k 3: QLIKE gain over k = 0, x 10^3
\providecommand{\rhcxGainLgbmCloseOneFive}{}\renewcommand{\rhcxGainLgbmCloseOneFive}{3.77}  % lgbm c16 W 1000 k 5: QLIKE gain over k = 0, x 10^3
\providecommand{\rhcxGainLgbmCloseOneTen}{}\renewcommand{\rhcxGainLgbmCloseOneTen}{2.77}  % lgbm c16 W 1000 k 10: QLIKE gain over k = 0, x 10^3
\providecommand{\rhcxGainLgbmCloseOneTwenty}{}\renewcommand{\rhcxGainLgbmCloseOneTwenty}{1.22}  % lgbm c16 W 1000 k 20: QLIKE gain over k = 0, x 10^3
\providecommand{\rhcxGainLgbmCloseOneFull}{}\renewcommand{\rhcxGainLgbmCloseOneFull}{1.62}  % lgbm c16 W 1000 k 46: QLIKE gain over k = 0, x 10^3
\providecommand{\rhcxGainFullPLgbmCloseOne}{}\renewcommand{\rhcxGainFullPLgbmCloseOne}{p = 0.490}  % lgbm c16 W 1000: DM p, full set vs k = 0
\providecommand{\rhcxGainBestPLgbmCloseOne}{}\renewcommand{\rhcxGainBestPLgbmCloseOne}{p = 0.084}  % lgbm c16 W 1000: DM p, best rung vs k = 0
\providecommand{\rhcxKstarLgbmCloseTwo}{}\renewcommand{\rhcxKstarLgbmCloseTwo}{5}  % lgbm c16 W 2000: rung with the lowest QLIKE
\providecommand{\rhcxKninetyLgbmCloseTwo}{}\renewcommand{\rhcxKninetyLgbmCloseTwo}{3}  % lgbm c16 W 2000: smallest k with 90% of the full set's gain
\providecommand{\rhcxGainLgbmCloseTwoOne}{}\renewcommand{\rhcxGainLgbmCloseTwoOne}{\ensuremath{-}0.59}  % lgbm c16 W 2000 k 1: QLIKE gain over k = 0, x 10^3
\providecommand{\rhcxGainLgbmCloseTwoThree}{}\renewcommand{\rhcxGainLgbmCloseTwoThree}{1.95}  % lgbm c16 W 2000 k 3: QLIKE gain over k = 0, x 10^3
\providecommand{\rhcxGainLgbmCloseTwoFive}{}\renewcommand{\rhcxGainLgbmCloseTwoFive}{2.45}  % lgbm c16 W 2000 k 5: QLIKE gain over k = 0, x 10^3
\providecommand{\rhcxGainLgbmCloseTwoTen}{}\renewcommand{\rhcxGainLgbmCloseTwoTen}{0.62}  % lgbm c16 W 2000 k 10: QLIKE gain over k = 0, x 10^3
\providecommand{\rhcxGainLgbmCloseTwoTwenty}{}\renewcommand{\rhcxGainLgbmCloseTwoTwenty}{1.20}  % lgbm c16 W 2000 k 20: QLIKE gain over k = 0, x 10^3
\providecommand{\rhcxGainLgbmCloseTwoFull}{}\renewcommand{\rhcxGainLgbmCloseTwoFull}{0.45}  % lgbm c16 W 2000 k 46: QLIKE gain over k = 0, x 10^3
\providecommand{\rhcxGainFullPLgbmCloseTwo}{}\renewcommand{\rhcxGainFullPLgbmCloseTwo}{p = 0.828}  % lgbm c16 W 2000: DM p, full set vs k = 0
\providecommand{\rhcxGainBestPLgbmCloseTwo}{}\renewcommand{\rhcxGainBestPLgbmCloseTwo}{p = 0.224}  % lgbm c16 W 2000: DM p, best rung vs k = 0
\providecommand{\rhcxKstarLgbmCloseMain}{}\renewcommand{\rhcxKstarLgbmCloseMain}{3}  % lgbm c16 W 3000: rung with the lowest QLIKE
\providecommand{\rhcxKninetyLgbmCloseMain}{}\renewcommand{\rhcxKninetyLgbmCloseMain}{3}  % lgbm c16 W 3000: smallest k with 90% of the full set's gain
\providecommand{\rhcxGainLgbmCloseMainOne}{}\renewcommand{\rhcxGainLgbmCloseMainOne}{0.49}  % lgbm c16 W 3000 k 1: QLIKE gain over k = 0, x 10^3
\providecommand{\rhcxGainLgbmCloseMainThree}{}\renewcommand{\rhcxGainLgbmCloseMainThree}{4.35}  % lgbm c16 W 3000 k 3: QLIKE gain over k = 0, x 10^3
\providecommand{\rhcxGainLgbmCloseMainFive}{}\renewcommand{\rhcxGainLgbmCloseMainFive}{3.90}  % lgbm c16 W 3000 k 5: QLIKE gain over k = 0, x 10^3
\providecommand{\rhcxGainLgbmCloseMainTen}{}\renewcommand{\rhcxGainLgbmCloseMainTen}{2.57}  % lgbm c16 W 3000 k 10: QLIKE gain over k = 0, x 10^3
\providecommand{\rhcxGainLgbmCloseMainTwenty}{}\renewcommand{\rhcxGainLgbmCloseMainTwenty}{2.31}  % lgbm c16 W 3000 k 20: QLIKE gain over k = 0, x 10^3
\providecommand{\rhcxGainLgbmCloseMainFull}{}\renewcommand{\rhcxGainLgbmCloseMainFull}{3.01}  % lgbm c16 W 3000 k 46: QLIKE gain over k = 0, x 10^3
\providecommand{\rhcxGainFullPLgbmCloseMain}{}\renewcommand{\rhcxGainFullPLgbmCloseMain}{p = 0.105}  % lgbm c16 W 3000: DM p, full set vs k = 0
\providecommand{\rhcxGainBestPLgbmCloseMain}{}\renewcommand{\rhcxGainBestPLgbmCloseMain}{p = 0.005}  % lgbm c16 W 3000: DM p, best rung vs k = 0
\providecommand{\rhcxKstarLgbmCloseLong}{}\renewcommand{\rhcxKstarLgbmCloseLong}{3}  % lgbm c16 W 3500: rung with the lowest QLIKE
\providecommand{\rhcxKninetyLgbmCloseLong}{}\renewcommand{\rhcxKninetyLgbmCloseLong}{3}  % lgbm c16 W 3500: smallest k with 90% of the full set's gain
\providecommand{\rhcxGainLgbmCloseLongOne}{}\renewcommand{\rhcxGainLgbmCloseLongOne}{0.17}  % lgbm c16 W 3500 k 1: QLIKE gain over k = 0, x 10^3
\providecommand{\rhcxGainLgbmCloseLongThree}{}\renewcommand{\rhcxGainLgbmCloseLongThree}{3.63}  % lgbm c16 W 3500 k 3: QLIKE gain over k = 0, x 10^3
\providecommand{\rhcxGainLgbmCloseLongFive}{}\renewcommand{\rhcxGainLgbmCloseLongFive}{2.38}  % lgbm c16 W 3500 k 5: QLIKE gain over k = 0, x 10^3
\providecommand{\rhcxGainLgbmCloseLongTen}{}\renewcommand{\rhcxGainLgbmCloseLongTen}{2.32}  % lgbm c16 W 3500 k 10: QLIKE gain over k = 0, x 10^3
\providecommand{\rhcxGainLgbmCloseLongTwenty}{}\renewcommand{\rhcxGainLgbmCloseLongTwenty}{1.76}  % lgbm c16 W 3500 k 20: QLIKE gain over k = 0, x 10^3
\providecommand{\rhcxGainLgbmCloseLongFull}{}\renewcommand{\rhcxGainLgbmCloseLongFull}{2.14}  % lgbm c16 W 3500 k 46: QLIKE gain over k = 0, x 10^3
\providecommand{\rhcxGainFullPLgbmCloseLong}{}\renewcommand{\rhcxGainFullPLgbmCloseLong}{p = 0.241}  % lgbm c16 W 3500: DM p, full set vs k = 0
\providecommand{\rhcxGainBestPLgbmCloseLong}{}\renewcommand{\rhcxGainBestPLgbmCloseLong}{p = 0.019}  % lgbm c16 W 3500: DM p, best rung vs k = 0
\providecommand{\rhcxKstarReadLgbmClose}{}\renewcommand{\rhcxKstarReadLgbmClose}{falls}  % lgbm c16: best k across W (RULE add. 1 (7))
\providecommand{\rhcxKninetyReadLgbmClose}{}\renewcommand{\rhcxKninetyReadLgbmClose}{falls}  % lgbm c16: k90 across W
\providecommand{\rhcxFullBestLgbmClose}{}\renewcommand{\rhcxFullBestLgbmClose}{0}  % lgbm c16: windows (of 4) where k = 46 is best
\providecommand{\rhcxSlopePosLgbmClose}{}\renewcommand{\rhcxSlopePosLgbmClose}{4}  % lgbm c16: windows with a positive gain-log(c) slope
\providecommand{\rhcxKstarXgbBarOne}{}\renewcommand{\rhcxKstarXgbBarOne}{20}  % xgb all W 1000: rung with the lowest QLIKE
\providecommand{\rhcxKninetyXgbBarOne}{}\renewcommand{\rhcxKninetyXgbBarOne}{10}  % xgb all W 1000: smallest k with 90% of the full set's gain
\providecommand{\rhcxGainXgbBarOneOne}{}\renewcommand{\rhcxGainXgbBarOneOne}{\ensuremath{-}0.95}  % xgb all W 1000 k 1: QLIKE gain over k = 0, x 10^3
\providecommand{\rhcxGainXgbBarOneThree}{}\renewcommand{\rhcxGainXgbBarOneThree}{1.85}  % xgb all W 1000 k 3: QLIKE gain over k = 0, x 10^3
\providecommand{\rhcxGainXgbBarOneFive}{}\renewcommand{\rhcxGainXgbBarOneFive}{6.46}  % xgb all W 1000 k 5: QLIKE gain over k = 0, x 10^3
\providecommand{\rhcxGainXgbBarOneTen}{}\renewcommand{\rhcxGainXgbBarOneTen}{20.33}  % xgb all W 1000 k 10: QLIKE gain over k = 0, x 10^3
\providecommand{\rhcxGainXgbBarOneTwenty}{}\renewcommand{\rhcxGainXgbBarOneTwenty}{20.59}  % xgb all W 1000 k 20: QLIKE gain over k = 0, x 10^3
\providecommand{\rhcxGainXgbBarOneFull}{}\renewcommand{\rhcxGainXgbBarOneFull}{20.39}  % xgb all W 1000 k 46: QLIKE gain over k = 0, x 10^3
\providecommand{\rhcxGainFullPXgbBarOne}{}\renewcommand{\rhcxGainFullPXgbBarOne}{p < 0.001}  % xgb all W 1000: DM p, full set vs k = 0
\providecommand{\rhcxGainBestPXgbBarOne}{}\renewcommand{\rhcxGainBestPXgbBarOne}{p < 0.001}  % xgb all W 1000: DM p, best rung vs k = 0
\providecommand{\rhcxKstarXgbBarTwo}{}\renewcommand{\rhcxKstarXgbBarTwo}{20}  % xgb all W 2000: rung with the lowest QLIKE
\providecommand{\rhcxKninetyXgbBarTwo}{}\renewcommand{\rhcxKninetyXgbBarTwo}{10}  % xgb all W 2000: smallest k with 90% of the full set's gain
\providecommand{\rhcxGainXgbBarTwoOne}{}\renewcommand{\rhcxGainXgbBarTwoOne}{0.36}  % xgb all W 2000 k 1: QLIKE gain over k = 0, x 10^3
\providecommand{\rhcxGainXgbBarTwoThree}{}\renewcommand{\rhcxGainXgbBarTwoThree}{5.21}  % xgb all W 2000 k 3: QLIKE gain over k = 0, x 10^3
\providecommand{\rhcxGainXgbBarTwoFive}{}\renewcommand{\rhcxGainXgbBarTwoFive}{8.44}  % xgb all W 2000 k 5: QLIKE gain over k = 0, x 10^3
\providecommand{\rhcxGainXgbBarTwoTen}{}\renewcommand{\rhcxGainXgbBarTwoTen}{22.28}  % xgb all W 2000 k 10: QLIKE gain over k = 0, x 10^3
\providecommand{\rhcxGainXgbBarTwoTwenty}{}\renewcommand{\rhcxGainXgbBarTwoTwenty}{22.89}  % xgb all W 2000 k 20: QLIKE gain over k = 0, x 10^3
\providecommand{\rhcxGainXgbBarTwoFull}{}\renewcommand{\rhcxGainXgbBarTwoFull}{22.87}  % xgb all W 2000 k 46: QLIKE gain over k = 0, x 10^3
\providecommand{\rhcxGainFullPXgbBarTwo}{}\renewcommand{\rhcxGainFullPXgbBarTwo}{p < 0.001}  % xgb all W 2000: DM p, full set vs k = 0
\providecommand{\rhcxGainBestPXgbBarTwo}{}\renewcommand{\rhcxGainBestPXgbBarTwo}{p < 0.001}  % xgb all W 2000: DM p, best rung vs k = 0
\providecommand{\rhcxKstarXgbBarMain}{}\renewcommand{\rhcxKstarXgbBarMain}{46}  % xgb all W 3000: rung with the lowest QLIKE
\providecommand{\rhcxKninetyXgbBarMain}{}\renewcommand{\rhcxKninetyXgbBarMain}{10}  % xgb all W 3000: smallest k with 90% of the full set's gain
\providecommand{\rhcxGainXgbBarMainOne}{}\renewcommand{\rhcxGainXgbBarMainOne}{1.05}  % xgb all W 3000 k 1: QLIKE gain over k = 0, x 10^3
\providecommand{\rhcxGainXgbBarMainThree}{}\renewcommand{\rhcxGainXgbBarMainThree}{6.97}  % xgb all W 3000 k 3: QLIKE gain over k = 0, x 10^3
\providecommand{\rhcxGainXgbBarMainFive}{}\renewcommand{\rhcxGainXgbBarMainFive}{9.36}  % xgb all W 3000 k 5: QLIKE gain over k = 0, x 10^3
\providecommand{\rhcxGainXgbBarMainTen}{}\renewcommand{\rhcxGainXgbBarMainTen}{23.85}  % xgb all W 3000 k 10: QLIKE gain over k = 0, x 10^3
\providecommand{\rhcxGainXgbBarMainTwenty}{}\renewcommand{\rhcxGainXgbBarMainTwenty}{24.37}  % xgb all W 3000 k 20: QLIKE gain over k = 0, x 10^3
\providecommand{\rhcxGainXgbBarMainFull}{}\renewcommand{\rhcxGainXgbBarMainFull}{24.60}  % xgb all W 3000 k 46: QLIKE gain over k = 0, x 10^3
\providecommand{\rhcxGainFullPXgbBarMain}{}\renewcommand{\rhcxGainFullPXgbBarMain}{p < 0.001}  % xgb all W 3000: DM p, full set vs k = 0
\providecommand{\rhcxGainBestPXgbBarMain}{}\renewcommand{\rhcxGainBestPXgbBarMain}{p < 0.001}  % xgb all W 3000: DM p, best rung vs k = 0
\providecommand{\rhcxKstarXgbBarLong}{}\renewcommand{\rhcxKstarXgbBarLong}{46}  % xgb all W 3500: rung with the lowest QLIKE
\providecommand{\rhcxKninetyXgbBarLong}{}\renewcommand{\rhcxKninetyXgbBarLong}{10}  % xgb all W 3500: smallest k with 90% of the full set's gain
\providecommand{\rhcxGainXgbBarLongOne}{}\renewcommand{\rhcxGainXgbBarLongOne}{1.15}  % xgb all W 3500 k 1: QLIKE gain over k = 0, x 10^3
\providecommand{\rhcxGainXgbBarLongThree}{}\renewcommand{\rhcxGainXgbBarLongThree}{8.52}  % xgb all W 3500 k 3: QLIKE gain over k = 0, x 10^3
\providecommand{\rhcxGainXgbBarLongFive}{}\renewcommand{\rhcxGainXgbBarLongFive}{8.23}  % xgb all W 3500 k 5: QLIKE gain over k = 0, x 10^3
\providecommand{\rhcxGainXgbBarLongTen}{}\renewcommand{\rhcxGainXgbBarLongTen}{24.04}  % xgb all W 3500 k 10: QLIKE gain over k = 0, x 10^3
\providecommand{\rhcxGainXgbBarLongTwenty}{}\renewcommand{\rhcxGainXgbBarLongTwenty}{24.01}  % xgb all W 3500 k 20: QLIKE gain over k = 0, x 10^3
\providecommand{\rhcxGainXgbBarLongFull}{}\renewcommand{\rhcxGainXgbBarLongFull}{24.37}  % xgb all W 3500 k 46: QLIKE gain over k = 0, x 10^3
\providecommand{\rhcxGainFullPXgbBarLong}{}\renewcommand{\rhcxGainFullPXgbBarLong}{p < 0.001}  % xgb all W 3500: DM p, full set vs k = 0
\providecommand{\rhcxGainBestPXgbBarLong}{}\renewcommand{\rhcxGainBestPXgbBarLong}{p < 0.001}  % xgb all W 3500: DM p, best rung vs k = 0
\providecommand{\rhcxKstarReadXgbBar}{}\renewcommand{\rhcxKstarReadXgbBar}{rises}  % xgb all: best k across W (RULE add. 1 (7))
\providecommand{\rhcxKninetyReadXgbBar}{}\renewcommand{\rhcxKninetyReadXgbBar}{is flat}  % xgb all: k90 across W
\providecommand{\rhcxFullBestXgbBar}{}\renewcommand{\rhcxFullBestXgbBar}{2}  % xgb all: windows (of 4) where k = 46 is best
\providecommand{\rhcxSlopePosXgbBar}{}\renewcommand{\rhcxSlopePosXgbBar}{4}  % xgb all: windows with a positive gain-log(c) slope
\providecommand{\rhcxKstarXgbCloseOne}{}\renewcommand{\rhcxKstarXgbCloseOne}{5}  % xgb c16 W 1000: rung with the lowest QLIKE
\providecommand{\rhcxKninetyXgbCloseOne}{}\renewcommand{\rhcxKninetyXgbCloseOne}{5}  % xgb c16 W 1000: smallest k with 90% of the full set's gain
\providecommand{\rhcxGainXgbCloseOneOne}{}\renewcommand{\rhcxGainXgbCloseOneOne}{\ensuremath{-}1.91}  % xgb c16 W 1000 k 1: QLIKE gain over k = 0, x 10^3
\providecommand{\rhcxGainXgbCloseOneThree}{}\renewcommand{\rhcxGainXgbCloseOneThree}{\ensuremath{-}1.79}  % xgb c16 W 1000 k 3: QLIKE gain over k = 0, x 10^3
\providecommand{\rhcxGainXgbCloseOneFive}{}\renewcommand{\rhcxGainXgbCloseOneFive}{2.01}  % xgb c16 W 1000 k 5: QLIKE gain over k = 0, x 10^3
\providecommand{\rhcxGainXgbCloseOneTen}{}\renewcommand{\rhcxGainXgbCloseOneTen}{\ensuremath{-}0.19}  % xgb c16 W 1000 k 10: QLIKE gain over k = 0, x 10^3
\providecommand{\rhcxGainXgbCloseOneTwenty}{}\renewcommand{\rhcxGainXgbCloseOneTwenty}{0.19}  % xgb c16 W 1000 k 20: QLIKE gain over k = 0, x 10^3
\providecommand{\rhcxGainXgbCloseOneFull}{}\renewcommand{\rhcxGainXgbCloseOneFull}{0.25}  % xgb c16 W 1000 k 46: QLIKE gain over k = 0, x 10^3
\providecommand{\rhcxGainFullPXgbCloseOne}{}\renewcommand{\rhcxGainFullPXgbCloseOne}{p = 0.909}  % xgb c16 W 1000: DM p, full set vs k = 0
\providecommand{\rhcxGainBestPXgbCloseOne}{}\renewcommand{\rhcxGainBestPXgbCloseOne}{p = 0.325}  % xgb c16 W 1000: DM p, best rung vs k = 0
\providecommand{\rhcxKstarXgbCloseTwo}{}\renewcommand{\rhcxKstarXgbCloseTwo}{5}  % xgb c16 W 2000: rung with the lowest QLIKE
\providecommand{\rhcxKninetyXgbCloseTwo}{}\renewcommand{\rhcxKninetyXgbCloseTwo}{5}  % xgb c16 W 2000: smallest k with 90% of the full set's gain
\providecommand{\rhcxGainXgbCloseTwoOne}{}\renewcommand{\rhcxGainXgbCloseTwoOne}{\ensuremath{-}0.66}  % xgb c16 W 2000 k 1: QLIKE gain over k = 0, x 10^3
\providecommand{\rhcxGainXgbCloseTwoThree}{}\renewcommand{\rhcxGainXgbCloseTwoThree}{1.42}  % xgb c16 W 2000 k 3: QLIKE gain over k = 0, x 10^3
\providecommand{\rhcxGainXgbCloseTwoFive}{}\renewcommand{\rhcxGainXgbCloseTwoFive}{2.08}  % xgb c16 W 2000 k 5: QLIKE gain over k = 0, x 10^3
\providecommand{\rhcxGainXgbCloseTwoTen}{}\renewcommand{\rhcxGainXgbCloseTwoTen}{1.30}  % xgb c16 W 2000 k 10: QLIKE gain over k = 0, x 10^3
\providecommand{\rhcxGainXgbCloseTwoTwenty}{}\renewcommand{\rhcxGainXgbCloseTwoTwenty}{1.90}  % xgb c16 W 2000 k 20: QLIKE gain over k = 0, x 10^3
\providecommand{\rhcxGainXgbCloseTwoFull}{}\renewcommand{\rhcxGainXgbCloseTwoFull}{1.59}  % xgb c16 W 2000 k 46: QLIKE gain over k = 0, x 10^3
\providecommand{\rhcxGainFullPXgbCloseTwo}{}\renewcommand{\rhcxGainFullPXgbCloseTwo}{p = 0.423}  % xgb c16 W 2000: DM p, full set vs k = 0
\providecommand{\rhcxGainBestPXgbCloseTwo}{}\renewcommand{\rhcxGainBestPXgbCloseTwo}{p = 0.291}  % xgb c16 W 2000: DM p, best rung vs k = 0
\providecommand{\rhcxKstarXgbCloseMain}{}\renewcommand{\rhcxKstarXgbCloseMain}{5}  % xgb c16 W 3000: rung with the lowest QLIKE
\providecommand{\rhcxKninetyXgbCloseMain}{}\renewcommand{\rhcxKninetyXgbCloseMain}{5}  % xgb c16 W 3000: smallest k with 90% of the full set's gain
\providecommand{\rhcxGainXgbCloseMainOne}{}\renewcommand{\rhcxGainXgbCloseMainOne}{0.47}  % xgb c16 W 3000 k 1: QLIKE gain over k = 0, x 10^3
\providecommand{\rhcxGainXgbCloseMainThree}{}\renewcommand{\rhcxGainXgbCloseMainThree}{1.29}  % xgb c16 W 3000 k 3: QLIKE gain over k = 0, x 10^3
\providecommand{\rhcxGainXgbCloseMainFive}{}\renewcommand{\rhcxGainXgbCloseMainFive}{3.81}  % xgb c16 W 3000 k 5: QLIKE gain over k = 0, x 10^3
\providecommand{\rhcxGainXgbCloseMainTen}{}\renewcommand{\rhcxGainXgbCloseMainTen}{2.24}  % xgb c16 W 3000 k 10: QLIKE gain over k = 0, x 10^3
\providecommand{\rhcxGainXgbCloseMainTwenty}{}\renewcommand{\rhcxGainXgbCloseMainTwenty}{3.46}  % xgb c16 W 3000 k 20: QLIKE gain over k = 0, x 10^3
\providecommand{\rhcxGainXgbCloseMainFull}{}\renewcommand{\rhcxGainXgbCloseMainFull}{3.06}  % xgb c16 W 3000 k 46: QLIKE gain over k = 0, x 10^3
\providecommand{\rhcxGainFullPXgbCloseMain}{}\renewcommand{\rhcxGainFullPXgbCloseMain}{p = 0.120}  % xgb c16 W 3000: DM p, full set vs k = 0
\providecommand{\rhcxGainBestPXgbCloseMain}{}\renewcommand{\rhcxGainBestPXgbCloseMain}{p = 0.050}  % xgb c16 W 3000: DM p, best rung vs k = 0
\providecommand{\rhcxKstarXgbCloseLong}{}\renewcommand{\rhcxKstarXgbCloseLong}{3}  % xgb c16 W 3500: rung with the lowest QLIKE
\providecommand{\rhcxKninetyXgbCloseLong}{}\renewcommand{\rhcxKninetyXgbCloseLong}{3}  % xgb c16 W 3500: smallest k with 90% of the full set's gain
\providecommand{\rhcxGainXgbCloseLongOne}{}\renewcommand{\rhcxGainXgbCloseLongOne}{1.13}  % xgb c16 W 3500 k 1: QLIKE gain over k = 0, x 10^3
\providecommand{\rhcxGainXgbCloseLongThree}{}\renewcommand{\rhcxGainXgbCloseLongThree}{4.34}  % xgb c16 W 3500 k 3: QLIKE gain over k = 0, x 10^3
\providecommand{\rhcxGainXgbCloseLongFive}{}\renewcommand{\rhcxGainXgbCloseLongFive}{4.22}  % xgb c16 W 3500 k 5: QLIKE gain over k = 0, x 10^3
\providecommand{\rhcxGainXgbCloseLongTen}{}\renewcommand{\rhcxGainXgbCloseLongTen}{3.09}  % xgb c16 W 3500 k 10: QLIKE gain over k = 0, x 10^3
\providecommand{\rhcxGainXgbCloseLongTwenty}{}\renewcommand{\rhcxGainXgbCloseLongTwenty}{3.06}  % xgb c16 W 3500 k 20: QLIKE gain over k = 0, x 10^3
\providecommand{\rhcxGainXgbCloseLongFull}{}\renewcommand{\rhcxGainXgbCloseLongFull}{3.04}  % xgb c16 W 3500 k 46: QLIKE gain over k = 0, x 10^3
\providecommand{\rhcxGainFullPXgbCloseLong}{}\renewcommand{\rhcxGainFullPXgbCloseLong}{p = 0.100}  % xgb c16 W 3500: DM p, full set vs k = 0
\providecommand{\rhcxGainBestPXgbCloseLong}{}\renewcommand{\rhcxGainBestPXgbCloseLong}{p = 0.008}  % xgb c16 W 3500: DM p, best rung vs k = 0
\providecommand{\rhcxKstarReadXgbClose}{}\renewcommand{\rhcxKstarReadXgbClose}{falls}  % xgb c16: best k across W (RULE add. 1 (7))
\providecommand{\rhcxKninetyReadXgbClose}{}\renewcommand{\rhcxKninetyReadXgbClose}{falls}  % xgb c16: k90 across W
\providecommand{\rhcxFullBestXgbClose}{}\renewcommand{\rhcxFullBestXgbClose}{0}  % xgb c16: windows (of 4) where k = 46 is best
\providecommand{\rhcxSlopePosXgbClose}{}\renewcommand{\rhcxSlopePosXgbClose}{4}  % xgb c16: windows with a positive gain-log(c) slope
\providecommand{\rhcxKstarRidgeBarOne}{}\renewcommand{\rhcxKstarRidgeBarOne}{46}  % ridge all W 1000: rung with the lowest QLIKE
\providecommand{\rhcxKninetyRidgeBarOne}{}\renewcommand{\rhcxKninetyRidgeBarOne}{20}  % ridge all W 1000: smallest k with 90% of the full set's gain
\providecommand{\rhcxGainRidgeBarOneOne}{}\renewcommand{\rhcxGainRidgeBarOneOne}{4.23}  % ridge all W 1000 k 1: QLIKE gain over k = 0, x 10^3
\providecommand{\rhcxGainRidgeBarOneThree}{}\renewcommand{\rhcxGainRidgeBarOneThree}{6.12}  % ridge all W 1000 k 3: QLIKE gain over k = 0, x 10^3
\providecommand{\rhcxGainRidgeBarOneFive}{}\renewcommand{\rhcxGainRidgeBarOneFive}{5.78}  % ridge all W 1000 k 5: QLIKE gain over k = 0, x 10^3
\providecommand{\rhcxGainRidgeBarOneTen}{}\renewcommand{\rhcxGainRidgeBarOneTen}{6.26}  % ridge all W 1000 k 10: QLIKE gain over k = 0, x 10^3
\providecommand{\rhcxGainRidgeBarOneTwenty}{}\renewcommand{\rhcxGainRidgeBarOneTwenty}{9.24}  % ridge all W 1000 k 20: QLIKE gain over k = 0, x 10^3
\providecommand{\rhcxGainRidgeBarOneFull}{}\renewcommand{\rhcxGainRidgeBarOneFull}{9.82}  % ridge all W 1000 k 46: QLIKE gain over k = 0, x 10^3
\providecommand{\rhcxGainFullPRidgeBarOne}{}\renewcommand{\rhcxGainFullPRidgeBarOne}{p < 0.001}  % ridge all W 1000: DM p, full set vs k = 0
\providecommand{\rhcxGainBestPRidgeBarOne}{}\renewcommand{\rhcxGainBestPRidgeBarOne}{p < 0.001}  % ridge all W 1000: DM p, best rung vs k = 0
\providecommand{\rhcxKstarRidgeBarTwo}{}\renewcommand{\rhcxKstarRidgeBarTwo}{5}  % ridge all W 2000: rung with the lowest QLIKE
\providecommand{\rhcxKninetyRidgeBarTwo}{}\renewcommand{\rhcxKninetyRidgeBarTwo}{3}  % ridge all W 2000: smallest k with 90% of the full set's gain
\providecommand{\rhcxGainRidgeBarTwoOne}{}\renewcommand{\rhcxGainRidgeBarTwoOne}{2.40}  % ridge all W 2000 k 1: QLIKE gain over k = 0, x 10^3
\providecommand{\rhcxGainRidgeBarTwoThree}{}\renewcommand{\rhcxGainRidgeBarTwoThree}{9.09}  % ridge all W 2000 k 3: QLIKE gain over k = 0, x 10^3
\providecommand{\rhcxGainRidgeBarTwoFive}{}\renewcommand{\rhcxGainRidgeBarTwoFive}{10.30}  % ridge all W 2000 k 5: QLIKE gain over k = 0, x 10^3
\providecommand{\rhcxGainRidgeBarTwoTen}{}\renewcommand{\rhcxGainRidgeBarTwoTen}{8.58}  % ridge all W 2000 k 10: QLIKE gain over k = 0, x 10^3
\providecommand{\rhcxGainRidgeBarTwoTwenty}{}\renewcommand{\rhcxGainRidgeBarTwoTwenty}{9.79}  % ridge all W 2000 k 20: QLIKE gain over k = 0, x 10^3
\providecommand{\rhcxGainRidgeBarTwoFull}{}\renewcommand{\rhcxGainRidgeBarTwoFull}{6.84}  % ridge all W 2000 k 46: QLIKE gain over k = 0, x 10^3
\providecommand{\rhcxGainFullPRidgeBarTwo}{}\renewcommand{\rhcxGainFullPRidgeBarTwo}{p = 0.004}  % ridge all W 2000: DM p, full set vs k = 0
\providecommand{\rhcxGainBestPRidgeBarTwo}{}\renewcommand{\rhcxGainBestPRidgeBarTwo}{p < 0.001}  % ridge all W 2000: DM p, best rung vs k = 0
\providecommand{\rhcxKstarRidgeBarMain}{}\renewcommand{\rhcxKstarRidgeBarMain}{10}  % ridge all W 3000: rung with the lowest QLIKE
\providecommand{\rhcxKninetyRidgeBarMain}{}\renewcommand{\rhcxKninetyRidgeBarMain}{3}  % ridge all W 3000: smallest k with 90% of the full set's gain
\providecommand{\rhcxGainRidgeBarMainOne}{}\renewcommand{\rhcxGainRidgeBarMainOne}{3.80}  % ridge all W 3000 k 1: QLIKE gain over k = 0, x 10^3
\providecommand{\rhcxGainRidgeBarMainThree}{}\renewcommand{\rhcxGainRidgeBarMainThree}{5.07}  % ridge all W 3000 k 3: QLIKE gain over k = 0, x 10^3
\providecommand{\rhcxGainRidgeBarMainFive}{}\renewcommand{\rhcxGainRidgeBarMainFive}{11.42}  % ridge all W 3000 k 5: QLIKE gain over k = 0, x 10^3
\providecommand{\rhcxGainRidgeBarMainTen}{}\renewcommand{\rhcxGainRidgeBarMainTen}{11.71}  % ridge all W 3000 k 10: QLIKE gain over k = 0, x 10^3
\providecommand{\rhcxGainRidgeBarMainTwenty}{}\renewcommand{\rhcxGainRidgeBarMainTwenty}{10.81}  % ridge all W 3000 k 20: QLIKE gain over k = 0, x 10^3
\providecommand{\rhcxGainRidgeBarMainFull}{}\renewcommand{\rhcxGainRidgeBarMainFull}{5.41}  % ridge all W 3000 k 46: QLIKE gain over k = 0, x 10^3
\providecommand{\rhcxGainFullPRidgeBarMain}{}\renewcommand{\rhcxGainFullPRidgeBarMain}{p = 0.015}  % ridge all W 3000: DM p, full set vs k = 0
\providecommand{\rhcxGainBestPRidgeBarMain}{}\renewcommand{\rhcxGainBestPRidgeBarMain}{p < 0.001}  % ridge all W 3000: DM p, best rung vs k = 0
\providecommand{\rhcxKstarRidgeBarLong}{}\renewcommand{\rhcxKstarRidgeBarLong}{10}  % ridge all W 3500: rung with the lowest QLIKE
\providecommand{\rhcxKninetyRidgeBarLong}{}\renewcommand{\rhcxKninetyRidgeBarLong}{10}  % ridge all W 3500: smallest k with 90% of the full set's gain
\providecommand{\rhcxGainRidgeBarLongOne}{}\renewcommand{\rhcxGainRidgeBarLongOne}{2.51}  % ridge all W 3500 k 1: QLIKE gain over k = 0, x 10^3
\providecommand{\rhcxGainRidgeBarLongThree}{}\renewcommand{\rhcxGainRidgeBarLongThree}{5.39}  % ridge all W 3500 k 3: QLIKE gain over k = 0, x 10^3
\providecommand{\rhcxGainRidgeBarLongFive}{}\renewcommand{\rhcxGainRidgeBarLongFive}{8.87}  % ridge all W 3500 k 5: QLIKE gain over k = 0, x 10^3
\providecommand{\rhcxGainRidgeBarLongTen}{}\renewcommand{\rhcxGainRidgeBarLongTen}{12.53}  % ridge all W 3500 k 10: QLIKE gain over k = 0, x 10^3
\providecommand{\rhcxGainRidgeBarLongTwenty}{}\renewcommand{\rhcxGainRidgeBarLongTwenty}{11.90}  % ridge all W 3500 k 20: QLIKE gain over k = 0, x 10^3
\providecommand{\rhcxGainRidgeBarLongFull}{}\renewcommand{\rhcxGainRidgeBarLongFull}{11.15}  % ridge all W 3500 k 46: QLIKE gain over k = 0, x 10^3
\providecommand{\rhcxGainFullPRidgeBarLong}{}\renewcommand{\rhcxGainFullPRidgeBarLong}{p < 0.001}  % ridge all W 3500: DM p, full set vs k = 0
\providecommand{\rhcxGainBestPRidgeBarLong}{}\renewcommand{\rhcxGainBestPRidgeBarLong}{p < 0.001}  % ridge all W 3500: DM p, best rung vs k = 0
\providecommand{\rhcxKstarReadRidgeBar}{}\renewcommand{\rhcxKstarReadRidgeBar}{is not monotone}  % ridge all: best k across W (RULE add. 1 (7))
\providecommand{\rhcxKninetyReadRidgeBar}{}\renewcommand{\rhcxKninetyReadRidgeBar}{is not monotone}  % ridge all: k90 across W
\providecommand{\rhcxFullBestRidgeBar}{}\renewcommand{\rhcxFullBestRidgeBar}{1}  % ridge all: windows (of 4) where k = 46 is best
\providecommand{\rhcxSlopePosRidgeBar}{}\renewcommand{\rhcxSlopePosRidgeBar}{4}  % ridge all: windows with a positive gain-log(c) slope
\providecommand{\rhcxKstarRidgeCloseOne}{}\renewcommand{\rhcxKstarRidgeCloseOne}{10}  % ridge c16 W 1000: rung with the lowest QLIKE
\providecommand{\rhcxKninetyRidgeCloseOne}{}\renewcommand{\rhcxKninetyRidgeCloseOne}{3}  % ridge c16 W 1000: smallest k with 90% of the full set's gain
\providecommand{\rhcxGainRidgeCloseOneOne}{}\renewcommand{\rhcxGainRidgeCloseOneOne}{3.08}  % ridge c16 W 1000 k 1: QLIKE gain over k = 0, x 10^3
\providecommand{\rhcxGainRidgeCloseOneThree}{}\renewcommand{\rhcxGainRidgeCloseOneThree}{5.03}  % ridge c16 W 1000 k 3: QLIKE gain over k = 0, x 10^3
\providecommand{\rhcxGainRidgeCloseOneFive}{}\renewcommand{\rhcxGainRidgeCloseOneFive}{4.92}  % ridge c16 W 1000 k 5: QLIKE gain over k = 0, x 10^3
\providecommand{\rhcxGainRidgeCloseOneTen}{}\renewcommand{\rhcxGainRidgeCloseOneTen}{5.36}  % ridge c16 W 1000 k 10: QLIKE gain over k = 0, x 10^3
\providecommand{\rhcxGainRidgeCloseOneTwenty}{}\renewcommand{\rhcxGainRidgeCloseOneTwenty}{3.56}  % ridge c16 W 1000 k 20: QLIKE gain over k = 0, x 10^3
\providecommand{\rhcxGainRidgeCloseOneFull}{}\renewcommand{\rhcxGainRidgeCloseOneFull}{4.14}  % ridge c16 W 1000 k 46: QLIKE gain over k = 0, x 10^3
\providecommand{\rhcxGainFullPRidgeCloseOne}{}\renewcommand{\rhcxGainFullPRidgeCloseOne}{p = 0.074}  % ridge c16 W 1000: DM p, full set vs k = 0
\providecommand{\rhcxGainBestPRidgeCloseOne}{}\renewcommand{\rhcxGainBestPRidgeCloseOne}{p = 0.003}  % ridge c16 W 1000: DM p, best rung vs k = 0
\providecommand{\rhcxKstarRidgeCloseTwo}{}\renewcommand{\rhcxKstarRidgeCloseTwo}{5}  % ridge c16 W 2000: rung with the lowest QLIKE
\providecommand{\rhcxKninetyRidgeCloseTwo}{}\renewcommand{\rhcxKninetyRidgeCloseTwo}{46}  % ridge c16 W 2000: smallest k with 90% of the full set's gain
\providecommand{\rhcxGainRidgeCloseTwoOne}{}\renewcommand{\rhcxGainRidgeCloseTwoOne}{\ensuremath{-}1.34}  % ridge c16 W 2000 k 1: QLIKE gain over k = 0, x 10^3
\providecommand{\rhcxGainRidgeCloseTwoThree}{}\renewcommand{\rhcxGainRidgeCloseTwoThree}{1.50}  % ridge c16 W 2000 k 3: QLIKE gain over k = 0, x 10^3
\providecommand{\rhcxGainRidgeCloseTwoFive}{}\renewcommand{\rhcxGainRidgeCloseTwoFive}{4.62}  % ridge c16 W 2000 k 5: QLIKE gain over k = 0, x 10^3
\providecommand{\rhcxGainRidgeCloseTwoTen}{}\renewcommand{\rhcxGainRidgeCloseTwoTen}{3.83}  % ridge c16 W 2000 k 10: QLIKE gain over k = 0, x 10^3
\providecommand{\rhcxGainRidgeCloseTwoTwenty}{}\renewcommand{\rhcxGainRidgeCloseTwoTwenty}{2.91}  % ridge c16 W 2000 k 20: QLIKE gain over k = 0, x 10^3
\providecommand{\rhcxGainRidgeCloseTwoFull}{}\renewcommand{\rhcxGainRidgeCloseTwoFull}{\ensuremath{-}8.48}  % ridge c16 W 2000 k 46: QLIKE gain over k = 0, x 10^3
\providecommand{\rhcxGainFullPRidgeCloseTwo}{}\renewcommand{\rhcxGainFullPRidgeCloseTwo}{p = 0.166}  % ridge c16 W 2000: DM p, full set vs k = 0
\providecommand{\rhcxGainBestPRidgeCloseTwo}{}\renewcommand{\rhcxGainBestPRidgeCloseTwo}{p = 0.037}  % ridge c16 W 2000: DM p, best rung vs k = 0
\providecommand{\rhcxKstarRidgeCloseMain}{}\renewcommand{\rhcxKstarRidgeCloseMain}{5}  % ridge c16 W 3000: rung with the lowest QLIKE
\providecommand{\rhcxKninetyRidgeCloseMain}{}\renewcommand{\rhcxKninetyRidgeCloseMain}{46}  % ridge c16 W 3000: smallest k with 90% of the full set's gain
\providecommand{\rhcxGainRidgeCloseMainOne}{}\renewcommand{\rhcxGainRidgeCloseMainOne}{\ensuremath{-}0.81}  % ridge c16 W 3000 k 1: QLIKE gain over k = 0, x 10^3
\providecommand{\rhcxGainRidgeCloseMainThree}{}\renewcommand{\rhcxGainRidgeCloseMainThree}{\ensuremath{-}0.06}  % ridge c16 W 3000 k 3: QLIKE gain over k = 0, x 10^3
\providecommand{\rhcxGainRidgeCloseMainFive}{}\renewcommand{\rhcxGainRidgeCloseMainFive}{5.09}  % ridge c16 W 3000 k 5: QLIKE gain over k = 0, x 10^3
\providecommand{\rhcxGainRidgeCloseMainTen}{}\renewcommand{\rhcxGainRidgeCloseMainTen}{4.70}  % ridge c16 W 3000 k 10: QLIKE gain over k = 0, x 10^3
\providecommand{\rhcxGainRidgeCloseMainTwenty}{}\renewcommand{\rhcxGainRidgeCloseMainTwenty}{3.66}  % ridge c16 W 3000 k 20: QLIKE gain over k = 0, x 10^3
\providecommand{\rhcxGainRidgeCloseMainFull}{}\renewcommand{\rhcxGainRidgeCloseMainFull}{\ensuremath{-}4.75}  % ridge c16 W 3000 k 46: QLIKE gain over k = 0, x 10^3
\providecommand{\rhcxGainFullPRidgeCloseMain}{}\renewcommand{\rhcxGainFullPRidgeCloseMain}{p = 0.236}  % ridge c16 W 3000: DM p, full set vs k = 0
\providecommand{\rhcxGainBestPRidgeCloseMain}{}\renewcommand{\rhcxGainBestPRidgeCloseMain}{p = 0.010}  % ridge c16 W 3000: DM p, best rung vs k = 0
\providecommand{\rhcxKstarRidgeCloseLong}{}\renewcommand{\rhcxKstarRidgeCloseLong}{10}  % ridge c16 W 3500: rung with the lowest QLIKE
\providecommand{\rhcxKninetyRidgeCloseLong}{}\renewcommand{\rhcxKninetyRidgeCloseLong}{10}  % ridge c16 W 3500: smallest k with 90% of the full set's gain
\providecommand{\rhcxGainRidgeCloseLongOne}{}\renewcommand{\rhcxGainRidgeCloseLongOne}{\ensuremath{-}1.26}  % ridge c16 W 3500 k 1: QLIKE gain over k = 0, x 10^3
\providecommand{\rhcxGainRidgeCloseLongThree}{}\renewcommand{\rhcxGainRidgeCloseLongThree}{0.08}  % ridge c16 W 3500 k 3: QLIKE gain over k = 0, x 10^3
\providecommand{\rhcxGainRidgeCloseLongFive}{}\renewcommand{\rhcxGainRidgeCloseLongFive}{0.88}  % ridge c16 W 3500 k 5: QLIKE gain over k = 0, x 10^3
\providecommand{\rhcxGainRidgeCloseLongTen}{}\renewcommand{\rhcxGainRidgeCloseLongTen}{5.10}  % ridge c16 W 3500 k 10: QLIKE gain over k = 0, x 10^3
\providecommand{\rhcxGainRidgeCloseLongTwenty}{}\renewcommand{\rhcxGainRidgeCloseLongTwenty}{4.32}  % ridge c16 W 3500 k 20: QLIKE gain over k = 0, x 10^3
\providecommand{\rhcxGainRidgeCloseLongFull}{}\renewcommand{\rhcxGainRidgeCloseLongFull}{4.25}  % ridge c16 W 3500 k 46: QLIKE gain over k = 0, x 10^3
\providecommand{\rhcxGainFullPRidgeCloseLong}{}\renewcommand{\rhcxGainFullPRidgeCloseLong}{p = 0.079}  % ridge c16 W 3500: DM p, full set vs k = 0
\providecommand{\rhcxGainBestPRidgeCloseLong}{}\renewcommand{\rhcxGainBestPRidgeCloseLong}{p = 0.014}  % ridge c16 W 3500: DM p, best rung vs k = 0
\providecommand{\rhcxKstarReadRidgeClose}{}\renewcommand{\rhcxKstarReadRidgeClose}{is not monotone}  % ridge c16: best k across W (RULE add. 1 (7))
\providecommand{\rhcxKninetyReadRidgeClose}{}\renewcommand{\rhcxKninetyReadRidgeClose}{is not monotone}  % ridge c16: k90 across W
\providecommand{\rhcxFullBestRidgeClose}{}\renewcommand{\rhcxFullBestRidgeClose}{0}  % ridge c16: windows (of 4) where k = 46 is best
\providecommand{\rhcxSlopePosRidgeClose}{}\renewcommand{\rhcxSlopePosRidgeClose}{2}  % ridge c16: windows with a positive gain-log(c) slope
\providecommand{\rhcxKstarLassoBarOne}{}\renewcommand{\rhcxKstarLassoBarOne}{20}  % lasso all W 1000: rung with the lowest QLIKE
\providecommand{\rhcxKninetyLassoBarOne}{}\renewcommand{\rhcxKninetyLassoBarOne}{20}  % lasso all W 1000: smallest k with 90% of the full set's gain
\providecommand{\rhcxGainLassoBarOneOne}{}\renewcommand{\rhcxGainLassoBarOneOne}{3.36}  % lasso all W 1000 k 1: QLIKE gain over k = 0, x 10^3
\providecommand{\rhcxGainLassoBarOneThree}{}\renewcommand{\rhcxGainLassoBarOneThree}{6.43}  % lasso all W 1000 k 3: QLIKE gain over k = 0, x 10^3
\providecommand{\rhcxGainLassoBarOneFive}{}\renewcommand{\rhcxGainLassoBarOneFive}{5.55}  % lasso all W 1000 k 5: QLIKE gain over k = 0, x 10^3
\providecommand{\rhcxGainLassoBarOneTen}{}\renewcommand{\rhcxGainLassoBarOneTen}{4.21}  % lasso all W 1000 k 10: QLIKE gain over k = 0, x 10^3
\providecommand{\rhcxGainLassoBarOneTwenty}{}\renewcommand{\rhcxGainLassoBarOneTwenty}{9.97}  % lasso all W 1000 k 20: QLIKE gain over k = 0, x 10^3
\providecommand{\rhcxGainLassoBarOneFull}{}\renewcommand{\rhcxGainLassoBarOneFull}{8.26}  % lasso all W 1000 k 46: QLIKE gain over k = 0, x 10^3
\providecommand{\rhcxGainFullPLassoBarOne}{}\renewcommand{\rhcxGainFullPLassoBarOne}{p < 0.001}  % lasso all W 1000: DM p, full set vs k = 0
\providecommand{\rhcxGainBestPLassoBarOne}{}\renewcommand{\rhcxGainBestPLassoBarOne}{p < 0.001}  % lasso all W 1000: DM p, best rung vs k = 0
\providecommand{\rhcxKstarLassoBarTwo}{}\renewcommand{\rhcxKstarLassoBarTwo}{10}  % lasso all W 2000: rung with the lowest QLIKE
\providecommand{\rhcxKninetyLassoBarTwo}{}\renewcommand{\rhcxKninetyLassoBarTwo}{1}  % lasso all W 2000: smallest k with 90% of the full set's gain
\providecommand{\rhcxGainLassoBarTwoOne}{}\renewcommand{\rhcxGainLassoBarTwoOne}{6.47}  % lasso all W 2000 k 1: QLIKE gain over k = 0, x 10^3
\providecommand{\rhcxGainLassoBarTwoThree}{}\renewcommand{\rhcxGainLassoBarTwoThree}{7.49}  % lasso all W 2000 k 3: QLIKE gain over k = 0, x 10^3
\providecommand{\rhcxGainLassoBarTwoFive}{}\renewcommand{\rhcxGainLassoBarTwoFive}{8.78}  % lasso all W 2000 k 5: QLIKE gain over k = 0, x 10^3
\providecommand{\rhcxGainLassoBarTwoTen}{}\renewcommand{\rhcxGainLassoBarTwoTen}{11.50}  % lasso all W 2000 k 10: QLIKE gain over k = 0, x 10^3
\providecommand{\rhcxGainLassoBarTwoTwenty}{}\renewcommand{\rhcxGainLassoBarTwoTwenty}{9.80}  % lasso all W 2000 k 20: QLIKE gain over k = 0, x 10^3
\providecommand{\rhcxGainLassoBarTwoFull}{}\renewcommand{\rhcxGainLassoBarTwoFull}{4.37}  % lasso all W 2000 k 46: QLIKE gain over k = 0, x 10^3
\providecommand{\rhcxGainFullPLassoBarTwo}{}\renewcommand{\rhcxGainFullPLassoBarTwo}{p = 0.161}  % lasso all W 2000: DM p, full set vs k = 0
\providecommand{\rhcxGainBestPLassoBarTwo}{}\renewcommand{\rhcxGainBestPLassoBarTwo}{p < 0.001}  % lasso all W 2000: DM p, best rung vs k = 0
\providecommand{\rhcxKstarLassoBarMain}{}\renewcommand{\rhcxKstarLassoBarMain}{10}  % lasso all W 3000: rung with the lowest QLIKE
\providecommand{\rhcxKninetyLassoBarMain}{}\renewcommand{\rhcxKninetyLassoBarMain}{1}  % lasso all W 3000: smallest k with 90% of the full set's gain
\providecommand{\rhcxGainLassoBarMainOne}{}\renewcommand{\rhcxGainLassoBarMainOne}{6.08}  % lasso all W 3000 k 1: QLIKE gain over k = 0, x 10^3
\providecommand{\rhcxGainLassoBarMainThree}{}\renewcommand{\rhcxGainLassoBarMainThree}{7.07}  % lasso all W 3000 k 3: QLIKE gain over k = 0, x 10^3
\providecommand{\rhcxGainLassoBarMainFive}{}\renewcommand{\rhcxGainLassoBarMainFive}{11.05}  % lasso all W 3000 k 5: QLIKE gain over k = 0, x 10^3
\providecommand{\rhcxGainLassoBarMainTen}{}\renewcommand{\rhcxGainLassoBarMainTen}{13.42}  % lasso all W 3000 k 10: QLIKE gain over k = 0, x 10^3
\providecommand{\rhcxGainLassoBarMainTwenty}{}\renewcommand{\rhcxGainLassoBarMainTwenty}{12.74}  % lasso all W 3000 k 20: QLIKE gain over k = 0, x 10^3
\providecommand{\rhcxGainLassoBarMainFull}{}\renewcommand{\rhcxGainLassoBarMainFull}{5.81}  % lasso all W 3000 k 46: QLIKE gain over k = 0, x 10^3
\providecommand{\rhcxGainFullPLassoBarMain}{}\renewcommand{\rhcxGainFullPLassoBarMain}{p = 0.059}  % lasso all W 3000: DM p, full set vs k = 0
\providecommand{\rhcxGainBestPLassoBarMain}{}\renewcommand{\rhcxGainBestPLassoBarMain}{p < 0.001}  % lasso all W 3000: DM p, best rung vs k = 0
\providecommand{\rhcxKstarLassoBarLong}{}\renewcommand{\rhcxKstarLassoBarLong}{10}  % lasso all W 3500: rung with the lowest QLIKE
\providecommand{\rhcxKninetyLassoBarLong}{}\renewcommand{\rhcxKninetyLassoBarLong}{10}  % lasso all W 3500: smallest k with 90% of the full set's gain
\providecommand{\rhcxGainLassoBarLongOne}{}\renewcommand{\rhcxGainLassoBarLongOne}{6.18}  % lasso all W 3500 k 1: QLIKE gain over k = 0, x 10^3
\providecommand{\rhcxGainLassoBarLongThree}{}\renewcommand{\rhcxGainLassoBarLongThree}{6.94}  % lasso all W 3500 k 3: QLIKE gain over k = 0, x 10^3
\providecommand{\rhcxGainLassoBarLongFive}{}\renewcommand{\rhcxGainLassoBarLongFive}{6.66}  % lasso all W 3500 k 5: QLIKE gain over k = 0, x 10^3
\providecommand{\rhcxGainLassoBarLongTen}{}\renewcommand{\rhcxGainLassoBarLongTen}{14.76}  % lasso all W 3500 k 10: QLIKE gain over k = 0, x 10^3
\providecommand{\rhcxGainLassoBarLongTwenty}{}\renewcommand{\rhcxGainLassoBarLongTwenty}{14.27}  % lasso all W 3500 k 20: QLIKE gain over k = 0, x 10^3
\providecommand{\rhcxGainLassoBarLongFull}{}\renewcommand{\rhcxGainLassoBarLongFull}{13.34}  % lasso all W 3500 k 46: QLIKE gain over k = 0, x 10^3
\providecommand{\rhcxGainFullPLassoBarLong}{}\renewcommand{\rhcxGainFullPLassoBarLong}{p < 0.001}  % lasso all W 3500: DM p, full set vs k = 0
\providecommand{\rhcxGainBestPLassoBarLong}{}\renewcommand{\rhcxGainBestPLassoBarLong}{p < 0.001}  % lasso all W 3500: DM p, best rung vs k = 0
\providecommand{\rhcxKstarReadLassoBar}{}\renewcommand{\rhcxKstarReadLassoBar}{falls}  % lasso all: best k across W (RULE add. 1 (7))
\providecommand{\rhcxKninetyReadLassoBar}{}\renewcommand{\rhcxKninetyReadLassoBar}{is not monotone}  % lasso all: k90 across W
\providecommand{\rhcxFullBestLassoBar}{}\renewcommand{\rhcxFullBestLassoBar}{0}  % lasso all: windows (of 4) where k = 46 is best
\providecommand{\rhcxSlopePosLassoBar}{}\renewcommand{\rhcxSlopePosLassoBar}{4}  % lasso all: windows with a positive gain-log(c) slope
\providecommand{\rhcxKstarLassoCloseOne}{}\renewcommand{\rhcxKstarLassoCloseOne}{46}  % lasso c16 W 1000: rung with the lowest QLIKE
\providecommand{\rhcxKninetyLassoCloseOne}{}\renewcommand{\rhcxKninetyLassoCloseOne}{46}  % lasso c16 W 1000: smallest k with 90% of the full set's gain
\providecommand{\rhcxGainLassoCloseOneOne}{}\renewcommand{\rhcxGainLassoCloseOneOne}{1.53}  % lasso c16 W 1000 k 1: QLIKE gain over k = 0, x 10^3
\providecommand{\rhcxGainLassoCloseOneThree}{}\renewcommand{\rhcxGainLassoCloseOneThree}{3.17}  % lasso c16 W 1000 k 3: QLIKE gain over k = 0, x 10^3
\providecommand{\rhcxGainLassoCloseOneFive}{}\renewcommand{\rhcxGainLassoCloseOneFive}{\ensuremath{-}1.77}  % lasso c16 W 1000 k 5: QLIKE gain over k = 0, x 10^3
\providecommand{\rhcxGainLassoCloseOneTen}{}\renewcommand{\rhcxGainLassoCloseOneTen}{2.68}  % lasso c16 W 1000 k 10: QLIKE gain over k = 0, x 10^3
\providecommand{\rhcxGainLassoCloseOneTwenty}{}\renewcommand{\rhcxGainLassoCloseOneTwenty}{3.82}  % lasso c16 W 1000 k 20: QLIKE gain over k = 0, x 10^3
\providecommand{\rhcxGainLassoCloseOneFull}{}\renewcommand{\rhcxGainLassoCloseOneFull}{4.33}  % lasso c16 W 1000 k 46: QLIKE gain over k = 0, x 10^3
\providecommand{\rhcxGainFullPLassoCloseOne}{}\renewcommand{\rhcxGainFullPLassoCloseOne}{p = 0.160}  % lasso c16 W 1000: DM p, full set vs k = 0
\providecommand{\rhcxGainBestPLassoCloseOne}{}\renewcommand{\rhcxGainBestPLassoCloseOne}{p = 0.160}  % lasso c16 W 1000: DM p, best rung vs k = 0
\providecommand{\rhcxKstarLassoCloseTwo}{}\renewcommand{\rhcxKstarLassoCloseTwo}{10}  % lasso c16 W 2000: rung with the lowest QLIKE
\providecommand{\rhcxKninetyLassoCloseTwo}{}\renewcommand{\rhcxKninetyLassoCloseTwo}{46}  % lasso c16 W 2000: smallest k with 90% of the full set's gain
\providecommand{\rhcxGainLassoCloseTwoOne}{}\renewcommand{\rhcxGainLassoCloseTwoOne}{\ensuremath{-}1.20}  % lasso c16 W 2000 k 1: QLIKE gain over k = 0, x 10^3
\providecommand{\rhcxGainLassoCloseTwoThree}{}\renewcommand{\rhcxGainLassoCloseTwoThree}{\ensuremath{-}2.00}  % lasso c16 W 2000 k 3: QLIKE gain over k = 0, x 10^3
\providecommand{\rhcxGainLassoCloseTwoFive}{}\renewcommand{\rhcxGainLassoCloseTwoFive}{\ensuremath{-}1.11}  % lasso c16 W 2000 k 5: QLIKE gain over k = 0, x 10^3
\providecommand{\rhcxGainLassoCloseTwoTen}{}\renewcommand{\rhcxGainLassoCloseTwoTen}{1.69}  % lasso c16 W 2000 k 10: QLIKE gain over k = 0, x 10^3
\providecommand{\rhcxGainLassoCloseTwoTwenty}{}\renewcommand{\rhcxGainLassoCloseTwoTwenty}{1.08}  % lasso c16 W 2000 k 20: QLIKE gain over k = 0, x 10^3
\providecommand{\rhcxGainLassoCloseTwoFull}{}\renewcommand{\rhcxGainLassoCloseTwoFull}{\ensuremath{-}17.71}  % lasso c16 W 2000 k 46: QLIKE gain over k = 0, x 10^3
\providecommand{\rhcxGainFullPLassoCloseTwo}{}\renewcommand{\rhcxGainFullPLassoCloseTwo}{p = 0.040}  % lasso c16 W 2000: DM p, full set vs k = 0
\providecommand{\rhcxGainBestPLassoCloseTwo}{}\renewcommand{\rhcxGainBestPLassoCloseTwo}{p = 0.491}  % lasso c16 W 2000: DM p, best rung vs k = 0
\providecommand{\rhcxKstarLassoCloseMain}{}\renewcommand{\rhcxKstarLassoCloseMain}{20}  % lasso c16 W 3000: rung with the lowest QLIKE
\providecommand{\rhcxKninetyLassoCloseMain}{}\renewcommand{\rhcxKninetyLassoCloseMain}{46}  % lasso c16 W 3000: smallest k with 90% of the full set's gain
\providecommand{\rhcxGainLassoCloseMainOne}{}\renewcommand{\rhcxGainLassoCloseMainOne}{\ensuremath{-}1.77}  % lasso c16 W 3000 k 1: QLIKE gain over k = 0, x 10^3
\providecommand{\rhcxGainLassoCloseMainThree}{}\renewcommand{\rhcxGainLassoCloseMainThree}{\ensuremath{-}3.15}  % lasso c16 W 3000 k 3: QLIKE gain over k = 0, x 10^3
\providecommand{\rhcxGainLassoCloseMainFive}{}\renewcommand{\rhcxGainLassoCloseMainFive}{\ensuremath{-}0.81}  % lasso c16 W 3000 k 5: QLIKE gain over k = 0, x 10^3
\providecommand{\rhcxGainLassoCloseMainTen}{}\renewcommand{\rhcxGainLassoCloseMainTen}{2.73}  % lasso c16 W 3000 k 10: QLIKE gain over k = 0, x 10^3
\providecommand{\rhcxGainLassoCloseMainTwenty}{}\renewcommand{\rhcxGainLassoCloseMainTwenty}{2.86}  % lasso c16 W 3000 k 20: QLIKE gain over k = 0, x 10^3
\providecommand{\rhcxGainLassoCloseMainFull}{}\renewcommand{\rhcxGainLassoCloseMainFull}{\ensuremath{-}8.63}  % lasso c16 W 3000 k 46: QLIKE gain over k = 0, x 10^3
\providecommand{\rhcxGainFullPLassoCloseMain}{}\renewcommand{\rhcxGainFullPLassoCloseMain}{p = 0.086}  % lasso c16 W 3000: DM p, full set vs k = 0
\providecommand{\rhcxGainBestPLassoCloseMain}{}\renewcommand{\rhcxGainBestPLassoCloseMain}{p = 0.233}  % lasso c16 W 3000: DM p, best rung vs k = 0
\providecommand{\rhcxKstarLassoCloseLong}{}\renewcommand{\rhcxKstarLassoCloseLong}{10}  % lasso c16 W 3500: rung with the lowest QLIKE
\providecommand{\rhcxKninetyLassoCloseLong}{}\renewcommand{\rhcxKninetyLassoCloseLong}{10}  % lasso c16 W 3500: smallest k with 90% of the full set's gain
\providecommand{\rhcxGainLassoCloseLongOne}{}\renewcommand{\rhcxGainLassoCloseLongOne}{\ensuremath{-}1.33}  % lasso c16 W 3500 k 1: QLIKE gain over k = 0, x 10^3
\providecommand{\rhcxGainLassoCloseLongThree}{}\renewcommand{\rhcxGainLassoCloseLongThree}{\ensuremath{-}3.66}  % lasso c16 W 3500 k 3: QLIKE gain over k = 0, x 10^3
\providecommand{\rhcxGainLassoCloseLongFive}{}\renewcommand{\rhcxGainLassoCloseLongFive}{\ensuremath{-}3.07}  % lasso c16 W 3500 k 5: QLIKE gain over k = 0, x 10^3
\providecommand{\rhcxGainLassoCloseLongTen}{}\renewcommand{\rhcxGainLassoCloseLongTen}{4.13}  % lasso c16 W 3500 k 10: QLIKE gain over k = 0, x 10^3
\providecommand{\rhcxGainLassoCloseLongTwenty}{}\renewcommand{\rhcxGainLassoCloseLongTwenty}{3.73}  % lasso c16 W 3500 k 20: QLIKE gain over k = 0, x 10^3
\providecommand{\rhcxGainLassoCloseLongFull}{}\renewcommand{\rhcxGainLassoCloseLongFull}{2.78}  % lasso c16 W 3500 k 46: QLIKE gain over k = 0, x 10^3
\providecommand{\rhcxGainFullPLassoCloseLong}{}\renewcommand{\rhcxGainFullPLassoCloseLong}{p = 0.294}  % lasso c16 W 3500: DM p, full set vs k = 0
\providecommand{\rhcxGainBestPLassoCloseLong}{}\renewcommand{\rhcxGainBestPLassoCloseLong}{p = 0.054}  % lasso c16 W 3500: DM p, best rung vs k = 0
\providecommand{\rhcxKstarReadLassoClose}{}\renewcommand{\rhcxKstarReadLassoClose}{is not monotone}  % lasso c16: best k across W (RULE add. 1 (7))
\providecommand{\rhcxKninetyReadLassoClose}{}\renewcommand{\rhcxKninetyReadLassoClose}{falls}  % lasso c16: k90 across W
\providecommand{\rhcxFullBestLassoClose}{}\renewcommand{\rhcxFullBestLassoClose}{1}  % lasso c16: windows (of 4) where k = 46 is best
\providecommand{\rhcxSlopePosLassoClose}{}\renewcommand{\rhcxSlopePosLassoClose}{2}  % lasso c16: windows with a positive gain-log(c) slope
\providecommand{\rhcxKstarEnetBarOne}{}\renewcommand{\rhcxKstarEnetBarOne}{46}  % enet all W 1000: rung with the lowest QLIKE
\providecommand{\rhcxKninetyEnetBarOne}{}\renewcommand{\rhcxKninetyEnetBarOne}{10}  % enet all W 1000: smallest k with 90% of the full set's gain
\providecommand{\rhcxGainEnetBarOneOne}{}\renewcommand{\rhcxGainEnetBarOneOne}{3.92}  % enet all W 1000 k 1: QLIKE gain over k = 0, x 10^3
\providecommand{\rhcxGainEnetBarOneThree}{}\renewcommand{\rhcxGainEnetBarOneThree}{6.72}  % enet all W 1000 k 3: QLIKE gain over k = 0, x 10^3
\providecommand{\rhcxGainEnetBarOneFive}{}\renewcommand{\rhcxGainEnetBarOneFive}{6.98}  % enet all W 1000 k 5: QLIKE gain over k = 0, x 10^3
\providecommand{\rhcxGainEnetBarOneTen}{}\renewcommand{\rhcxGainEnetBarOneTen}{10.72}  % enet all W 1000 k 10: QLIKE gain over k = 0, x 10^3
\providecommand{\rhcxGainEnetBarOneTwenty}{}\renewcommand{\rhcxGainEnetBarOneTwenty}{11.01}  % enet all W 1000 k 20: QLIKE gain over k = 0, x 10^3
\providecommand{\rhcxGainEnetBarOneFull}{}\renewcommand{\rhcxGainEnetBarOneFull}{11.22}  % enet all W 1000 k 46: QLIKE gain over k = 0, x 10^3
\providecommand{\rhcxGainFullPEnetBarOne}{}\renewcommand{\rhcxGainFullPEnetBarOne}{p < 0.001}  % enet all W 1000: DM p, full set vs k = 0
\providecommand{\rhcxGainBestPEnetBarOne}{}\renewcommand{\rhcxGainBestPEnetBarOne}{p < 0.001}  % enet all W 1000: DM p, best rung vs k = 0
\providecommand{\rhcxKstarEnetBarTwo}{}\renewcommand{\rhcxKstarEnetBarTwo}{10}  % enet all W 2000: rung with the lowest QLIKE
\providecommand{\rhcxKninetyEnetBarTwo}{}\renewcommand{\rhcxKninetyEnetBarTwo}{5}  % enet all W 2000: smallest k with 90% of the full set's gain
\providecommand{\rhcxGainEnetBarTwoOne}{}\renewcommand{\rhcxGainEnetBarTwoOne}{4.04}  % enet all W 2000 k 1: QLIKE gain over k = 0, x 10^3
\providecommand{\rhcxGainEnetBarTwoThree}{}\renewcommand{\rhcxGainEnetBarTwoThree}{5.55}  % enet all W 2000 k 3: QLIKE gain over k = 0, x 10^3
\providecommand{\rhcxGainEnetBarTwoFive}{}\renewcommand{\rhcxGainEnetBarTwoFive}{10.57}  % enet all W 2000 k 5: QLIKE gain over k = 0, x 10^3
\providecommand{\rhcxGainEnetBarTwoTen}{}\renewcommand{\rhcxGainEnetBarTwoTen}{11.11}  % enet all W 2000 k 10: QLIKE gain over k = 0, x 10^3
\providecommand{\rhcxGainEnetBarTwoTwenty}{}\renewcommand{\rhcxGainEnetBarTwoTwenty}{10.51}  % enet all W 2000 k 20: QLIKE gain over k = 0, x 10^3
\providecommand{\rhcxGainEnetBarTwoFull}{}\renewcommand{\rhcxGainEnetBarTwoFull}{8.14}  % enet all W 2000 k 46: QLIKE gain over k = 0, x 10^3
\providecommand{\rhcxGainFullPEnetBarTwo}{}\renewcommand{\rhcxGainFullPEnetBarTwo}{p < 0.001}  % enet all W 2000: DM p, full set vs k = 0
\providecommand{\rhcxGainBestPEnetBarTwo}{}\renewcommand{\rhcxGainBestPEnetBarTwo}{p < 0.001}  % enet all W 2000: DM p, best rung vs k = 0
\providecommand{\rhcxKstarEnetBarMain}{}\renewcommand{\rhcxKstarEnetBarMain}{10}  % enet all W 3000: rung with the lowest QLIKE
\providecommand{\rhcxKninetyEnetBarMain}{}\renewcommand{\rhcxKninetyEnetBarMain}{3}  % enet all W 3000: smallest k with 90% of the full set's gain
\providecommand{\rhcxGainEnetBarMainOne}{}\renewcommand{\rhcxGainEnetBarMainOne}{\ensuremath{-}0.16}  % enet all W 3000 k 1: QLIKE gain over k = 0, x 10^3
\providecommand{\rhcxGainEnetBarMainThree}{}\renewcommand{\rhcxGainEnetBarMainThree}{7.98}  % enet all W 3000 k 3: QLIKE gain over k = 0, x 10^3
\providecommand{\rhcxGainEnetBarMainFive}{}\renewcommand{\rhcxGainEnetBarMainFive}{11.58}  % enet all W 3000 k 5: QLIKE gain over k = 0, x 10^3
\providecommand{\rhcxGainEnetBarMainTen}{}\renewcommand{\rhcxGainEnetBarMainTen}{11.78}  % enet all W 3000 k 10: QLIKE gain over k = 0, x 10^3
\providecommand{\rhcxGainEnetBarMainTwenty}{}\renewcommand{\rhcxGainEnetBarMainTwenty}{10.97}  % enet all W 3000 k 20: QLIKE gain over k = 0, x 10^3
\providecommand{\rhcxGainEnetBarMainFull}{}\renewcommand{\rhcxGainEnetBarMainFull}{8.80}  % enet all W 3000 k 46: QLIKE gain over k = 0, x 10^3
\providecommand{\rhcxGainFullPEnetBarMain}{}\renewcommand{\rhcxGainFullPEnetBarMain}{p < 0.001}  % enet all W 3000: DM p, full set vs k = 0
\providecommand{\rhcxGainBestPEnetBarMain}{}\renewcommand{\rhcxGainBestPEnetBarMain}{p < 0.001}  % enet all W 3000: DM p, best rung vs k = 0
\providecommand{\rhcxKstarEnetBarLong}{}\renewcommand{\rhcxKstarEnetBarLong}{10}  % enet all W 3500: rung with the lowest QLIKE
\providecommand{\rhcxKninetyEnetBarLong}{}\renewcommand{\rhcxKninetyEnetBarLong}{5}  % enet all W 3500: smallest k with 90% of the full set's gain
\providecommand{\rhcxGainEnetBarLongOne}{}\renewcommand{\rhcxGainEnetBarLongOne}{3.88}  % enet all W 3500 k 1: QLIKE gain over k = 0, x 10^3
\providecommand{\rhcxGainEnetBarLongThree}{}\renewcommand{\rhcxGainEnetBarLongThree}{8.42}  % enet all W 3500 k 3: QLIKE gain over k = 0, x 10^3
\providecommand{\rhcxGainEnetBarLongFive}{}\renewcommand{\rhcxGainEnetBarLongFive}{11.87}  % enet all W 3500 k 5: QLIKE gain over k = 0, x 10^3
\providecommand{\rhcxGainEnetBarLongTen}{}\renewcommand{\rhcxGainEnetBarLongTen}{12.25}  % enet all W 3500 k 10: QLIKE gain over k = 0, x 10^3
\providecommand{\rhcxGainEnetBarLongTwenty}{}\renewcommand{\rhcxGainEnetBarLongTwenty}{12.04}  % enet all W 3500 k 20: QLIKE gain over k = 0, x 10^3
\providecommand{\rhcxGainEnetBarLongFull}{}\renewcommand{\rhcxGainEnetBarLongFull}{12.08}  % enet all W 3500 k 46: QLIKE gain over k = 0, x 10^3
\providecommand{\rhcxGainFullPEnetBarLong}{}\renewcommand{\rhcxGainFullPEnetBarLong}{p < 0.001}  % enet all W 3500: DM p, full set vs k = 0
\providecommand{\rhcxGainBestPEnetBarLong}{}\renewcommand{\rhcxGainBestPEnetBarLong}{p < 0.001}  % enet all W 3500: DM p, best rung vs k = 0
\providecommand{\rhcxKstarReadEnetBar}{}\renewcommand{\rhcxKstarReadEnetBar}{falls}  % enet all: best k across W (RULE add. 1 (7))
\providecommand{\rhcxKninetyReadEnetBar}{}\renewcommand{\rhcxKninetyReadEnetBar}{is not monotone}  % enet all: k90 across W
\providecommand{\rhcxFullBestEnetBar}{}\renewcommand{\rhcxFullBestEnetBar}{1}  % enet all: windows (of 4) where k = 46 is best
\providecommand{\rhcxSlopePosEnetBar}{}\renewcommand{\rhcxSlopePosEnetBar}{4}  % enet all: windows with a positive gain-log(c) slope
\providecommand{\rhcxKstarEnetCloseOne}{}\renewcommand{\rhcxKstarEnetCloseOne}{46}  % enet c16 W 1000: rung with the lowest QLIKE
\providecommand{\rhcxKninetyEnetCloseOne}{}\renewcommand{\rhcxKninetyEnetCloseOne}{46}  % enet c16 W 1000: smallest k with 90% of the full set's gain
\providecommand{\rhcxGainEnetCloseOneOne}{}\renewcommand{\rhcxGainEnetCloseOneOne}{2.47}  % enet c16 W 1000 k 1: QLIKE gain over k = 0, x 10^3
\providecommand{\rhcxGainEnetCloseOneThree}{}\renewcommand{\rhcxGainEnetCloseOneThree}{4.19}  % enet c16 W 1000 k 3: QLIKE gain over k = 0, x 10^3
\providecommand{\rhcxGainEnetCloseOneFive}{}\renewcommand{\rhcxGainEnetCloseOneFive}{4.46}  % enet c16 W 1000 k 5: QLIKE gain over k = 0, x 10^3
\providecommand{\rhcxGainEnetCloseOneTen}{}\renewcommand{\rhcxGainEnetCloseOneTen}{4.97}  % enet c16 W 1000 k 10: QLIKE gain over k = 0, x 10^3
\providecommand{\rhcxGainEnetCloseOneTwenty}{}\renewcommand{\rhcxGainEnetCloseOneTwenty}{5.20}  % enet c16 W 1000 k 20: QLIKE gain over k = 0, x 10^3
\providecommand{\rhcxGainEnetCloseOneFull}{}\renewcommand{\rhcxGainEnetCloseOneFull}{5.87}  % enet c16 W 1000 k 46: QLIKE gain over k = 0, x 10^3
\providecommand{\rhcxGainFullPEnetCloseOne}{}\renewcommand{\rhcxGainFullPEnetCloseOne}{p = 0.004}  % enet c16 W 1000: DM p, full set vs k = 0
\providecommand{\rhcxGainBestPEnetCloseOne}{}\renewcommand{\rhcxGainBestPEnetCloseOne}{p = 0.004}  % enet c16 W 1000: DM p, best rung vs k = 0
\providecommand{\rhcxKstarEnetCloseTwo}{}\renewcommand{\rhcxKstarEnetCloseTwo}{10}  % enet c16 W 2000: rung with the lowest QLIKE
\providecommand{\rhcxKninetyEnetCloseTwo}{}\renewcommand{\rhcxKninetyEnetCloseTwo}{46}  % enet c16 W 2000: smallest k with 90% of the full set's gain
\providecommand{\rhcxGainEnetCloseTwoOne}{}\renewcommand{\rhcxGainEnetCloseTwoOne}{\ensuremath{-}0.77}  % enet c16 W 2000 k 1: QLIKE gain over k = 0, x 10^3
\providecommand{\rhcxGainEnetCloseTwoThree}{}\renewcommand{\rhcxGainEnetCloseTwoThree}{0.35}  % enet c16 W 2000 k 3: QLIKE gain over k = 0, x 10^3
\providecommand{\rhcxGainEnetCloseTwoFive}{}\renewcommand{\rhcxGainEnetCloseTwoFive}{4.82}  % enet c16 W 2000 k 5: QLIKE gain over k = 0, x 10^3
\providecommand{\rhcxGainEnetCloseTwoTen}{}\renewcommand{\rhcxGainEnetCloseTwoTen}{5.86}  % enet c16 W 2000 k 10: QLIKE gain over k = 0, x 10^3
\providecommand{\rhcxGainEnetCloseTwoTwenty}{}\renewcommand{\rhcxGainEnetCloseTwoTwenty}{2.34}  % enet c16 W 2000 k 20: QLIKE gain over k = 0, x 10^3
\providecommand{\rhcxGainEnetCloseTwoFull}{}\renewcommand{\rhcxGainEnetCloseTwoFull}{\ensuremath{-}4.89}  % enet c16 W 2000 k 46: QLIKE gain over k = 0, x 10^3
\providecommand{\rhcxGainFullPEnetCloseTwo}{}\renewcommand{\rhcxGainFullPEnetCloseTwo}{p = 0.325}  % enet c16 W 2000: DM p, full set vs k = 0
\providecommand{\rhcxGainBestPEnetCloseTwo}{}\renewcommand{\rhcxGainBestPEnetCloseTwo}{p = 0.009}  % enet c16 W 2000: DM p, best rung vs k = 0
\providecommand{\rhcxKstarEnetCloseMain}{}\renewcommand{\rhcxKstarEnetCloseMain}{10}  % enet c16 W 3000: rung with the lowest QLIKE
\providecommand{\rhcxKninetyEnetCloseMain}{}\renewcommand{\rhcxKninetyEnetCloseMain}{46}  % enet c16 W 3000: smallest k with 90% of the full set's gain
\providecommand{\rhcxGainEnetCloseMainOne}{}\renewcommand{\rhcxGainEnetCloseMainOne}{0.12}  % enet c16 W 3000 k 1: QLIKE gain over k = 0, x 10^3
\providecommand{\rhcxGainEnetCloseMainThree}{}\renewcommand{\rhcxGainEnetCloseMainThree}{1.80}  % enet c16 W 3000 k 3: QLIKE gain over k = 0, x 10^3
\providecommand{\rhcxGainEnetCloseMainFive}{}\renewcommand{\rhcxGainEnetCloseMainFive}{5.12}  % enet c16 W 3000 k 5: QLIKE gain over k = 0, x 10^3
\providecommand{\rhcxGainEnetCloseMainTen}{}\renewcommand{\rhcxGainEnetCloseMainTen}{5.55}  % enet c16 W 3000 k 10: QLIKE gain over k = 0, x 10^3
\providecommand{\rhcxGainEnetCloseMainTwenty}{}\renewcommand{\rhcxGainEnetCloseMainTwenty}{2.41}  % enet c16 W 3000 k 20: QLIKE gain over k = 0, x 10^3
\providecommand{\rhcxGainEnetCloseMainFull}{}\renewcommand{\rhcxGainEnetCloseMainFull}{\ensuremath{-}1.68}  % enet c16 W 3000 k 46: QLIKE gain over k = 0, x 10^3
\providecommand{\rhcxGainFullPEnetCloseMain}{}\renewcommand{\rhcxGainFullPEnetCloseMain}{p = 0.631}  % enet c16 W 3000: DM p, full set vs k = 0
\providecommand{\rhcxGainBestPEnetCloseMain}{}\renewcommand{\rhcxGainBestPEnetCloseMain}{p = 0.009}  % enet c16 W 3000: DM p, best rung vs k = 0
\providecommand{\rhcxKstarEnetCloseLong}{}\renewcommand{\rhcxKstarEnetCloseLong}{5}  % enet c16 W 3500: rung with the lowest QLIKE
\providecommand{\rhcxKninetyEnetCloseLong}{}\renewcommand{\rhcxKninetyEnetCloseLong}{5}  % enet c16 W 3500: smallest k with 90% of the full set's gain
\providecommand{\rhcxGainEnetCloseLongOne}{}\renewcommand{\rhcxGainEnetCloseLongOne}{\ensuremath{-}0.90}  % enet c16 W 3500 k 1: QLIKE gain over k = 0, x 10^3
\providecommand{\rhcxGainEnetCloseLongThree}{}\renewcommand{\rhcxGainEnetCloseLongThree}{2.24}  % enet c16 W 3500 k 3: QLIKE gain over k = 0, x 10^3
\providecommand{\rhcxGainEnetCloseLongFive}{}\renewcommand{\rhcxGainEnetCloseLongFive}{5.41}  % enet c16 W 3500 k 5: QLIKE gain over k = 0, x 10^3
\providecommand{\rhcxGainEnetCloseLongTen}{}\renewcommand{\rhcxGainEnetCloseLongTen}{5.04}  % enet c16 W 3500 k 10: QLIKE gain over k = 0, x 10^3
\providecommand{\rhcxGainEnetCloseLongTwenty}{}\renewcommand{\rhcxGainEnetCloseLongTwenty}{4.89}  % enet c16 W 3500 k 20: QLIKE gain over k = 0, x 10^3
\providecommand{\rhcxGainEnetCloseLongFull}{}\renewcommand{\rhcxGainEnetCloseLongFull}{4.83}  % enet c16 W 3500 k 46: QLIKE gain over k = 0, x 10^3
\providecommand{\rhcxGainFullPEnetCloseLong}{}\renewcommand{\rhcxGainFullPEnetCloseLong}{p = 0.037}  % enet c16 W 3500: DM p, full set vs k = 0
\providecommand{\rhcxGainBestPEnetCloseLong}{}\renewcommand{\rhcxGainBestPEnetCloseLong}{p = 0.004}  % enet c16 W 3500: DM p, best rung vs k = 0
\providecommand{\rhcxKstarReadEnetClose}{}\renewcommand{\rhcxKstarReadEnetClose}{falls}  % enet c16: best k across W (RULE add. 1 (7))
\providecommand{\rhcxKninetyReadEnetClose}{}\renewcommand{\rhcxKninetyReadEnetClose}{falls}  % enet c16: k90 across W
\providecommand{\rhcxFullBestEnetClose}{}\renewcommand{\rhcxFullBestEnetClose}{1}  % enet c16: windows (of 4) where k = 46 is best
\providecommand{\rhcxSlopePosEnetClose}{}\renewcommand{\rhcxSlopePosEnetClose}{3}  % enet c16: windows with a positive gain-log(c) slope
\providecommand{\rhcxNKstarRisesBar}{}\renewcommand{\rhcxNKstarRisesBar}{1}  % all: models whose best k RISES with W
\providecommand{\rhcxWhoKstarRisesBar}{}\renewcommand{\rhcxWhoKstarRisesBar}{XGBoost}  % all: those models
\providecommand{\rhcxNKstarFallsBar}{}\renewcommand{\rhcxNKstarFallsBar}{2}  % all: models whose best k FALLS with W
\providecommand{\rhcxWhoKstarFallsBar}{}\renewcommand{\rhcxWhoKstarFallsBar}{lasso and elastic net}  % all: those models
\providecommand{\rhcxNKstarFlatBar}{}\renewcommand{\rhcxNKstarFlatBar}{1}  % all: models whose best k FLAT with W
\providecommand{\rhcxWhoKstarFlatBar}{}\renewcommand{\rhcxWhoKstarFlatBar}{LightGBM}  % all: those models
\providecommand{\rhcxNKstarMixedBar}{}\renewcommand{\rhcxNKstarMixedBar}{1}  % all: models whose best k MIXED with W
\providecommand{\rhcxWhoKstarMixedBar}{}\renewcommand{\rhcxWhoKstarMixedBar}{ridge}  % all: those models
\providecommand{\rhcxNFullBestBar}{}\renewcommand{\rhcxNFullBestBar}{8}  % all: (model, window) cells where the full set is best, of 20
\providecommand{\rhcxNFullBestTreesBar}{}\renewcommand{\rhcxNFullBestTreesBar}{6}  % all: tree cells (of 8) where the full set is best
\providecommand{\rhcxNFullBestLinearBar}{}\renewcommand{\rhcxNFullBestLinearBar}{2}  % all: linear cells (of 12) where the full set is best
\providecommand{\rhcxNKstarRisesClose}{}\renewcommand{\rhcxNKstarRisesClose}{0}  % c16: models whose best k RISES with W
\providecommand{\rhcxWhoKstarRisesClose}{}\renewcommand{\rhcxWhoKstarRisesClose}{none}  % c16: those models
\providecommand{\rhcxNKstarFallsClose}{}\renewcommand{\rhcxNKstarFallsClose}{3}  % c16: models whose best k FALLS with W
\providecommand{\rhcxWhoKstarFallsClose}{}\renewcommand{\rhcxWhoKstarFallsClose}{LightGBM, XGBoost and elastic net}  % c16: those models
\providecommand{\rhcxNKstarFlatClose}{}\renewcommand{\rhcxNKstarFlatClose}{0}  % c16: models whose best k FLAT with W
\providecommand{\rhcxWhoKstarFlatClose}{}\renewcommand{\rhcxWhoKstarFlatClose}{none}  % c16: those models
\providecommand{\rhcxNKstarMixedClose}{}\renewcommand{\rhcxNKstarMixedClose}{2}  % c16: models whose best k MIXED with W
\providecommand{\rhcxWhoKstarMixedClose}{}\renewcommand{\rhcxWhoKstarMixedClose}{ridge and lasso}  % c16: those models
\providecommand{\rhcxNFullBestClose}{}\renewcommand{\rhcxNFullBestClose}{2}  % c16: (model, window) cells where the full set is best, of 20
\providecommand{\rhcxNFullBestTreesClose}{}\renewcommand{\rhcxNFullBestTreesClose}{0}  % c16: tree cells (of 8) where the full set is best
\providecommand{\rhcxNFullBestLinearClose}{}\renewcommand{\rhcxNFullBestLinearClose}{2}  % c16: linear cells (of 12) where the full set is best
\fi
% 2026-10-22 (user): rerun note.
% 2026-10-23 (user, RR): parked -- the pooled variable importance is now computed on the final runs (linear models and
% the three trees, W = 3,000); the preliminary run results/rehaul_fi_pooled13_2026-10-22 (151 re-estimations, earlier
% design) is superseded and not printed. Was:
% "\editc{orange}{RR}{Decision: the pooled linear variable importance is computed on the current design (151 monthly
% re-estimations, results/rehaul\_fi\_pooled13\_2026-10-22: dense across variables, sparse within each) and goes in after
% the rerun of the pooled models; the tree analysis is to be decided.\\Problem: this subsection is still a placeholder.
% See the notes to the authors, ``Rerun pending''.}"
% 2026-10-23 (coordinator, tree learning curve): parked -- the note's "trees at the robustness window" item, now scoped
% to the SHAP values (the trees' accuracy across windows is in, tab:rh_fip_lc). Was: "...; the trees at the
% \rhfipWMax-day robustness window (only the linear models are computed there). The forest's walk gate is complete
% (\rhfipRfGateEq{} of \rhfipNOrigins{} origins).}"
% 2026-10-26 (integration, JFEc): parked -- the RR note. Resolved: the ablations are run (Section
% sec:rh_results_ablation, Table tab:rh_ablation); the MinGRU stays parked (author); the trees' accuracy across windows is in Appendix
% app:rh_fi_window; the trees' SHAP values at the robustness window are not planned. Open, reported to FRONT for the
% notes to the authors: the tree SHAP values, the live-window shares, the feature-count ladder and the window x
% feature-count grid are on the design before the FOMC patch (D13) and are to be rerun on the patched design. Was:
% \editc{orange}{RR}{Decision (user, 2026-10-22/23): variable importance on the final pooled runs at the main window, the linear models' standardized coefficients and the tree ensembles' TreeSHAP values at \rhfipNOrigins{} monthly origins; the preliminary pooled run (results/rehaul\_fi\_pooled13\_2026-10-22) is superseded.\\Still to come: the ablations (leave one category out, then one series out within the categories that matter; not run); the MinGRU network (parked, not covered); the trees' SHAP values at the \rhfipWMax-day robustness window (only the linear models' coefficients are computed there). The trees' accuracy across windows is in: \rhlcNArmsDone{} of the \rhlcNArmsAll{} walks at \rhlcWOne--\rhlcWLong{} days (Table~\ref{tab:rh_fip_lc}). The forest's walk gate is complete (\rhfipRfGateEq{} of \rhfipNOrigins{} origins).}
% 2026-10-23 (user, RR): parked -- the placeholder; the results follow. Was: "Variable importance is not yet computed for
% the pooled intraday model; the closing-interval analysis is parked (note E19)."
% 2026-10-26 (integration, JFEc; VI agent): the section in four parts (INTEGRATION_PLAN_2026-10-26.md, budget 3 pp.:
% about 1.5 pp. of text and three floats): (a) what the models use (Table tab:rh_fip_shares), (b) how many series carry
% the gain: the feature-count ladder (Figure fig:rh_ladder; results/ladder_2026-10-26, RULE.txt) and the ridge--lasso
% comparison, which replace the window test (P1/P2) as the main-text answer to the dense-and-weak question, (c) the
% late-starting series while live (Figure fig:rh_fi_live), (d) use versus accuracy (the ablations, \rhab... of
% generated/rehaul_ablation_numbers.tex, Table tab:rh_ablation). Moved to the online appendix (rehaul/A_fi.tex): the
% construction and checks, the rankings and their time profile, the conditional shares, the response shapes and
% interactions, the window test and the trees across windows, and the window x feature-count grid (framing B of
% drafts/2026-10-26_complexity.tex only).
% \rhPending{<name>}: a ladder or ablation number not yet in a generated file; <name> is the macro proposed to NUMBERS
% (generated/rehaul_ladder_numbers.tex, \rhld...; generated/rehaul_ablation_numbers.tex, \rhab...). Replace
% "\rhPending{rhldX}" by "\rhldX" once the file exists. Rung labels (k = 0, 1, 3, 5, 10, 20, 46) are the ladder RULE's
% design, not results; every result is a macro or \rhPending.
% Design note: the SHAP values, the live-window shares, the ladder and the grid are on the design before the FOMC
% patch (D13); they are to be rerun on the patched design, and the FOMC-day subset of the SHAP reducer is to include
% the three statements the loader adds (audit results-9; Appendix app:rh_fi_cond).
% Audit results-8 and results-20 (fixed in the first bullet): the explained bars include the NYSE-holiday sessions
% that Table 1 does not score; the penalties come from the last pre-sample fold, the tree configurations from an
% earlier development sample.
\ifdefined\rhchNEventDays\else% AUTO-GENERATED by writeup/make_rehaul_channels.py -- do not edit; re-run the script.
% \rhch... macros: the two channels of the pooled gain (results/channels_2026-10-26/RULE.txt items 1-3):
% FOMC statement days (EVENT) and ordinary days, Table 1's scored bars, the FOMC-patched design of record;
% DM (HLN) on daily average loss, lag 8.  Names in the generator's docstring.
% \rhchX... = ADDITIONAL, NOT PRE-REGISTERED (coordinator, 2026-10-26, after the first scoring): the full set
% against HAR + the futures' realized moments (FullVsRm) and HAR + realized moments against HAR (RmVsHar).
\providecommand{\rhPending}{\textit{pending}}
\providecommand{\rhchNBars}{}\renewcommand{\rhchNBars}{33{,}527}  % scored bars (Table 1)
\providecommand{\rhchNDays}{}\renewcommand{\rhchNDays}{2{,}579}  % scored days
\providecommand{\rhchNEventDays}{}\renewcommand{\rhchNEventDays}{81}  % EVENT days: scheduled FOMC statement days of the span
\providecommand{\rhchNOrdDays}{}\renewcommand{\rhchNOrdDays}{2{,}498}  % ORDINARY days
\providecommand{\rhchNEventBars}{}\renewcommand{\rhchNEventBars}{1{,}053}  % bars on EVENT days
\providecommand{\rhchEventBarPct}{}\renewcommand{\rhchEventBarPct}{3.1}  % EVENT-day bars, % of the scored bars
\providecommand{\rhchHacLag}{}\renewcommand{\rhchHacLag}{8}  % Newey-West lag of every DM test
\providecommand{\rhchAgreePp}{}\renewcommand{\rhchAgreePp}{2}  % item 1: the corrections agree within this many percentage points
\providecommand{\rhchLadderDesign}{}\renewcommand{\rhchLadderDesign}{patched}  % design of the ladder
\providecommand{\rhchRidgeTabOnePct}{}\renewcommand{\rhchRidgeTabOnePct}{$-$6.0}  % ridge: full vs HAR, every bar (Table 1): % QLIKE change
\providecommand{\rhchRidgeTabOneDm}{}\renewcommand{\rhchRidgeTabOneDm}{$-$6.28}  % ridge: full vs HAR, every bar (Table 1): DM
\providecommand{\rhchRidgeTabOneP}{}\renewcommand{\rhchRidgeTabOneP}{\ensuremath{<}0.001}  % ridge: full vs HAR, every bar (Table 1): DM p
\providecommand{\rhchRidgeTabOneMsePct}{}\renewcommand{\rhchRidgeTabOneMsePct}{100.6}  % ridge: full vs HAR, every bar (Table 1): % MSE change
\providecommand{\rhchRidgeTabOneMseP}{}\renewcommand{\rhchRidgeTabOneMseP}{0.261}  % ridge: full vs HAR, every bar (Table 1): DM p on MSE
\providecommand{\rhchRidgeCorrAPct}{}\renewcommand{\rhchRidgeCorrAPct}{$-$3.0}  % ridge: full vs HAR+FOMC, every bar: % QLIKE change
\providecommand{\rhchRidgeCorrADm}{}\renewcommand{\rhchRidgeCorrADm}{$-$3.64}  % ridge: full vs HAR+FOMC, every bar: DM
\providecommand{\rhchRidgeCorrAP}{}\renewcommand{\rhchRidgeCorrAP}{\ensuremath{<}0.001}  % ridge: full vs HAR+FOMC, every bar: DM p
\providecommand{\rhchRidgeCorrAMsePct}{}\renewcommand{\rhchRidgeCorrAMsePct}{100.8}  % ridge: full vs HAR+FOMC, every bar: % MSE change
\providecommand{\rhchRidgeCorrAMseP}{}\renewcommand{\rhchRidgeCorrAMseP}{0.258}  % ridge: full vs HAR+FOMC, every bar: DM p on MSE
\providecommand{\rhchRidgeCorrBPct}{}\renewcommand{\rhchRidgeCorrBPct}{$-$3.2}  % ridge: full without FOMC vs HAR, every bar: % QLIKE change
\providecommand{\rhchRidgeCorrBDm}{}\renewcommand{\rhchRidgeCorrBDm}{$-$3.74}  % ridge: full without FOMC vs HAR, every bar: DM
\providecommand{\rhchRidgeCorrBP}{}\renewcommand{\rhchRidgeCorrBP}{\ensuremath{<}0.001}  % ridge: full without FOMC vs HAR, every bar: DM p
\providecommand{\rhchRidgeCorrBMsePct}{}\renewcommand{\rhchRidgeCorrBMsePct}{102.3}  % ridge: full without FOMC vs HAR, every bar: % MSE change
\providecommand{\rhchRidgeCorrBMseP}{}\renewcommand{\rhchRidgeCorrBMseP}{0.258}  % ridge: full without FOMC vs HAR, every bar: DM p on MSE
\providecommand{\rhchRidgeCorrCPct}{}\renewcommand{\rhchRidgeCorrCPct}{$-$3.2}  % ridge: full vs HAR, ORDINARY days: % QLIKE change
\providecommand{\rhchRidgeCorrCDm}{}\renewcommand{\rhchRidgeCorrCDm}{$-$3.51}  % ridge: full vs HAR, ORDINARY days: DM
\providecommand{\rhchRidgeCorrCP}{}\renewcommand{\rhchRidgeCorrCP}{\ensuremath{<}0.001}  % ridge: full vs HAR, ORDINARY days: DM p
\providecommand{\rhchRidgeCorrCMsePct}{}\renewcommand{\rhchRidgeCorrCMsePct}{110.2}  % ridge: full vs HAR, ORDINARY days: % MSE change
\providecommand{\rhchRidgeCorrCMseP}{}\renewcommand{\rhchRidgeCorrCMseP}{0.258}  % ridge: full vs HAR, ORDINARY days: DM p on MSE
\providecommand{\rhchRidgeCorrAgree}{}\renewcommand{\rhchRidgeCorrAgree}{agree}  % ridge: the three corrections within 2.0 pp?
\providecommand{\rhchRidgeCorrSpread}{}\renewcommand{\rhchRidgeCorrSpread}{0.2}  % ridge: max - min of the three (pp)
\providecommand{\rhchLassoTabOnePct}{}\renewcommand{\rhchLassoTabOnePct}{$-$9.0}  % lasso: full vs HAR, every bar (Table 1): % QLIKE change
\providecommand{\rhchLassoTabOneDm}{}\renewcommand{\rhchLassoTabOneDm}{$-$6.34}  % lasso: full vs HAR, every bar (Table 1): DM
\providecommand{\rhchLassoTabOneP}{}\renewcommand{\rhchLassoTabOneP}{\ensuremath{<}0.001}  % lasso: full vs HAR, every bar (Table 1): DM p
\providecommand{\rhchLassoTabOneMsePct}{}\renewcommand{\rhchLassoTabOneMsePct}{114.7}  % lasso: full vs HAR, every bar (Table 1): % MSE change
\providecommand{\rhchLassoTabOneMseP}{}\renewcommand{\rhchLassoTabOneMseP}{0.247}  % lasso: full vs HAR, every bar (Table 1): DM p on MSE
\providecommand{\rhchLassoCorrAPct}{}\renewcommand{\rhchLassoCorrAPct}{$-$3.4}  % lasso: full vs HAR+FOMC, every bar: % QLIKE change
\providecommand{\rhchLassoCorrADm}{}\renewcommand{\rhchLassoCorrADm}{$-$3.20}  % lasso: full vs HAR+FOMC, every bar: DM
\providecommand{\rhchLassoCorrAP}{}\renewcommand{\rhchLassoCorrAP}{0.001}  % lasso: full vs HAR+FOMC, every bar: DM p
\providecommand{\rhchLassoCorrAMsePct}{}\renewcommand{\rhchLassoCorrAMsePct}{114.4}  % lasso: full vs HAR+FOMC, every bar: % MSE change
\providecommand{\rhchLassoCorrAMseP}{}\renewcommand{\rhchLassoCorrAMseP}{0.243}  % lasso: full vs HAR+FOMC, every bar: DM p on MSE
\providecommand{\rhchLassoCorrBPct}{}\renewcommand{\rhchLassoCorrBPct}{$-$2.8}  % lasso: full without FOMC vs HAR, every bar: % QLIKE change
\providecommand{\rhchLassoCorrBDm}{}\renewcommand{\rhchLassoCorrBDm}{$-$2.75}  % lasso: full without FOMC vs HAR, every bar: DM
\providecommand{\rhchLassoCorrBP}{}\renewcommand{\rhchLassoCorrBP}{0.006}  % lasso: full without FOMC vs HAR, every bar: DM p
\providecommand{\rhchLassoCorrBMsePct}{}\renewcommand{\rhchLassoCorrBMsePct}{115.3}  % lasso: full without FOMC vs HAR, every bar: % MSE change
\providecommand{\rhchLassoCorrBMseP}{}\renewcommand{\rhchLassoCorrBMseP}{0.245}  % lasso: full without FOMC vs HAR, every bar: DM p on MSE
\providecommand{\rhchLassoCorrCPct}{}\renewcommand{\rhchLassoCorrCPct}{$-$2.4}  % lasso: full vs HAR, ORDINARY days: % QLIKE change
\providecommand{\rhchLassoCorrCDm}{}\renewcommand{\rhchLassoCorrCDm}{$-$2.08}  % lasso: full vs HAR, ORDINARY days: DM
\providecommand{\rhchLassoCorrCP}{}\renewcommand{\rhchLassoCorrCP}{0.037}  % lasso: full vs HAR, ORDINARY days: DM p
\providecommand{\rhchLassoCorrCMsePct}{}\renewcommand{\rhchLassoCorrCMsePct}{126.0}  % lasso: full vs HAR, ORDINARY days: % MSE change
\providecommand{\rhchLassoCorrCMseP}{}\renewcommand{\rhchLassoCorrCMseP}{0.243}  % lasso: full vs HAR, ORDINARY days: DM p on MSE
\providecommand{\rhchLassoCorrAgree}{}\renewcommand{\rhchLassoCorrAgree}{agree}  % lasso: the three corrections within 2.0 pp?
\providecommand{\rhchLassoCorrSpread}{}\renewcommand{\rhchLassoCorrSpread}{1.0}  % lasso: max - min of the three (pp)
\providecommand{\rhchEnetTabOnePct}{}\renewcommand{\rhchEnetTabOnePct}{$-$6.7}  % enet: full vs HAR, every bar (Table 1): % QLIKE change
\providecommand{\rhchEnetTabOneDm}{}\renewcommand{\rhchEnetTabOneDm}{$-$7.36}  % enet: full vs HAR, every bar (Table 1): DM
\providecommand{\rhchEnetTabOneP}{}\renewcommand{\rhchEnetTabOneP}{\ensuremath{<}0.001}  % enet: full vs HAR, every bar (Table 1): DM p
\providecommand{\rhchEnetTabOneMsePct}{}\renewcommand{\rhchEnetTabOneMsePct}{117.7}  % enet: full vs HAR, every bar (Table 1): % MSE change
\providecommand{\rhchEnetTabOneMseP}{}\renewcommand{\rhchEnetTabOneMseP}{0.263}  % enet: full vs HAR, every bar (Table 1): DM p on MSE
\providecommand{\rhchEnetCorrAPct}{}\renewcommand{\rhchEnetCorrAPct}{$-$3.7}  % enet: full vs HAR+FOMC, every bar: % QLIKE change
\providecommand{\rhchEnetCorrADm}{}\renewcommand{\rhchEnetCorrADm}{$-$4.89}  % enet: full vs HAR+FOMC, every bar: DM
\providecommand{\rhchEnetCorrAP}{}\renewcommand{\rhchEnetCorrAP}{\ensuremath{<}0.001}  % enet: full vs HAR+FOMC, every bar: DM p
\providecommand{\rhchEnetCorrAMsePct}{}\renewcommand{\rhchEnetCorrAMsePct}{117.8}  % enet: full vs HAR+FOMC, every bar: % MSE change
\providecommand{\rhchEnetCorrAMseP}{}\renewcommand{\rhchEnetCorrAMseP}{0.261}  % enet: full vs HAR+FOMC, every bar: DM p on MSE
\providecommand{\rhchEnetCorrBPct}{}\renewcommand{\rhchEnetCorrBPct}{$-$3.8}  % enet: full without FOMC vs HAR, every bar: % QLIKE change
\providecommand{\rhchEnetCorrBDm}{}\renewcommand{\rhchEnetCorrBDm}{$-$5.00}  % enet: full without FOMC vs HAR, every bar: DM
\providecommand{\rhchEnetCorrBP}{}\renewcommand{\rhchEnetCorrBP}{\ensuremath{<}0.001}  % enet: full without FOMC vs HAR, every bar: DM p
\providecommand{\rhchEnetCorrBMsePct}{}\renewcommand{\rhchEnetCorrBMsePct}{117.8}  % enet: full without FOMC vs HAR, every bar: % MSE change
\providecommand{\rhchEnetCorrBMseP}{}\renewcommand{\rhchEnetCorrBMseP}{0.259}  % enet: full without FOMC vs HAR, every bar: DM p on MSE
\providecommand{\rhchEnetCorrCPct}{}\renewcommand{\rhchEnetCorrCPct}{$-$4.0}  % enet: full vs HAR, ORDINARY days: % QLIKE change
\providecommand{\rhchEnetCorrCDm}{}\renewcommand{\rhchEnetCorrCDm}{$-$4.73}  % enet: full vs HAR, ORDINARY days: DM
\providecommand{\rhchEnetCorrCP}{}\renewcommand{\rhchEnetCorrCP}{\ensuremath{<}0.001}  % enet: full vs HAR, ORDINARY days: DM p
\providecommand{\rhchEnetCorrCMsePct}{}\renewcommand{\rhchEnetCorrCMsePct}{128.9}  % enet: full vs HAR, ORDINARY days: % MSE change
\providecommand{\rhchEnetCorrCMseP}{}\renewcommand{\rhchEnetCorrCMseP}{0.260}  % enet: full vs HAR, ORDINARY days: DM p on MSE
\providecommand{\rhchEnetCorrAgree}{}\renewcommand{\rhchEnetCorrAgree}{agree}  % enet: the three corrections within 2.0 pp?
\providecommand{\rhchEnetCorrSpread}{}\renewcommand{\rhchEnetCorrSpread}{0.2}  % enet: max - min of the three (pp)
\providecommand{\rhchLgbmTabOnePct}{}\renewcommand{\rhchLgbmTabOnePct}{$-$14.1}  % lgbm: full vs HAR, every bar (Table 1): % QLIKE change
\providecommand{\rhchLgbmTabOneDm}{}\renewcommand{\rhchLgbmTabOneDm}{$-$10.26}  % lgbm: full vs HAR, every bar (Table 1): DM
\providecommand{\rhchLgbmTabOneP}{}\renewcommand{\rhchLgbmTabOneP}{\ensuremath{<}0.001}  % lgbm: full vs HAR, every bar (Table 1): DM p
\providecommand{\rhchLgbmTabOneMsePct}{}\renewcommand{\rhchLgbmTabOneMsePct}{20.0}  % lgbm: full vs HAR, every bar (Table 1): % MSE change
\providecommand{\rhchLgbmTabOneMseP}{}\renewcommand{\rhchLgbmTabOneMseP}{0.301}  % lgbm: full vs HAR, every bar (Table 1): DM p on MSE
\providecommand{\rhchLgbmCorrAPct}{}\renewcommand{\rhchLgbmCorrAPct}{$-$5.4}  % lgbm: full vs HAR+FOMC, every bar: % QLIKE change
\providecommand{\rhchLgbmCorrADm}{}\renewcommand{\rhchLgbmCorrADm}{$-$11.63}  % lgbm: full vs HAR+FOMC, every bar: DM
\providecommand{\rhchLgbmCorrAP}{}\renewcommand{\rhchLgbmCorrAP}{\ensuremath{<}0.001}  % lgbm: full vs HAR+FOMC, every bar: DM p
\providecommand{\rhchLgbmCorrAMsePct}{}\renewcommand{\rhchLgbmCorrAMsePct}{19.4}  % lgbm: full vs HAR+FOMC, every bar: % MSE change
\providecommand{\rhchLgbmCorrAMseP}{}\renewcommand{\rhchLgbmCorrAMseP}{0.276}  % lgbm: full vs HAR+FOMC, every bar: DM p on MSE
\providecommand{\rhchLgbmCorrBPct}{}\renewcommand{\rhchLgbmCorrBPct}{$-$5.3}  % lgbm: full without FOMC vs HAR, every bar: % QLIKE change
\providecommand{\rhchLgbmCorrBDm}{}\renewcommand{\rhchLgbmCorrBDm}{$-$9.62}  % lgbm: full without FOMC vs HAR, every bar: DM
\providecommand{\rhchLgbmCorrBP}{}\renewcommand{\rhchLgbmCorrBP}{\ensuremath{<}0.001}  % lgbm: full without FOMC vs HAR, every bar: DM p
\providecommand{\rhchLgbmCorrBMsePct}{}\renewcommand{\rhchLgbmCorrBMsePct}{21.6}  % lgbm: full without FOMC vs HAR, every bar: % MSE change
\providecommand{\rhchLgbmCorrBMseP}{}\renewcommand{\rhchLgbmCorrBMseP}{0.269}  % lgbm: full without FOMC vs HAR, every bar: DM p on MSE
\providecommand{\rhchLgbmCorrCPct}{}\renewcommand{\rhchLgbmCorrCPct}{$-$6.0}  % lgbm: full vs HAR, ORDINARY days: % QLIKE change
\providecommand{\rhchLgbmCorrCDm}{}\renewcommand{\rhchLgbmCorrCDm}{$-$10.70}  % lgbm: full vs HAR, ORDINARY days: DM
\providecommand{\rhchLgbmCorrCP}{}\renewcommand{\rhchLgbmCorrCP}{\ensuremath{<}0.001}  % lgbm: full vs HAR, ORDINARY days: DM p
\providecommand{\rhchLgbmCorrCMsePct}{}\renewcommand{\rhchLgbmCorrCMsePct}{22.8}  % lgbm: full vs HAR, ORDINARY days: % MSE change
\providecommand{\rhchLgbmCorrCMseP}{}\renewcommand{\rhchLgbmCorrCMseP}{0.275}  % lgbm: full vs HAR, ORDINARY days: DM p on MSE
\providecommand{\rhchLgbmCorrAgree}{}\renewcommand{\rhchLgbmCorrAgree}{agree}  % lgbm: the three corrections within 2.0 pp?
\providecommand{\rhchLgbmCorrSpread}{}\renewcommand{\rhchLgbmCorrSpread}{0.8}  % lgbm: max - min of the three (pp)
\providecommand{\rhchXgbTabOnePct}{}\renewcommand{\rhchXgbTabOnePct}{$-$14.6}  % xgb: full vs HAR, every bar (Table 1): % QLIKE change
\providecommand{\rhchXgbTabOneDm}{}\renewcommand{\rhchXgbTabOneDm}{$-$9.94}  % xgb: full vs HAR, every bar (Table 1): DM
\providecommand{\rhchXgbTabOneP}{}\renewcommand{\rhchXgbTabOneP}{\ensuremath{<}0.001}  % xgb: full vs HAR, every bar (Table 1): DM p
\providecommand{\rhchXgbTabOneMsePct}{}\renewcommand{\rhchXgbTabOneMsePct}{16.3}  % xgb: full vs HAR, every bar (Table 1): % MSE change
\providecommand{\rhchXgbTabOneMseP}{}\renewcommand{\rhchXgbTabOneMseP}{0.301}  % xgb: full vs HAR, every bar (Table 1): DM p on MSE
\providecommand{\rhchXgbCorrAPct}{}\renewcommand{\rhchXgbCorrAPct}{$-$5.7}  % xgb: full vs HAR+FOMC, every bar: % QLIKE change
\providecommand{\rhchXgbCorrADm}{}\renewcommand{\rhchXgbCorrADm}{$-$12.60}  % xgb: full vs HAR+FOMC, every bar: DM
\providecommand{\rhchXgbCorrAP}{}\renewcommand{\rhchXgbCorrAP}{\ensuremath{<}0.001}  % xgb: full vs HAR+FOMC, every bar: DM p
\providecommand{\rhchXgbCorrAMsePct}{}\renewcommand{\rhchXgbCorrAMsePct}{15.4}  % xgb: full vs HAR+FOMC, every bar: % MSE change
\providecommand{\rhchXgbCorrAMseP}{}\renewcommand{\rhchXgbCorrAMseP}{0.282}  % xgb: full vs HAR+FOMC, every bar: DM p on MSE
\providecommand{\rhchXgbCorrBPct}{}\renewcommand{\rhchXgbCorrBPct}{$-$5.8}  % xgb: full without FOMC vs HAR, every bar: % QLIKE change
\providecommand{\rhchXgbCorrBDm}{}\renewcommand{\rhchXgbCorrBDm}{$-$10.81}  % xgb: full without FOMC vs HAR, every bar: DM
\providecommand{\rhchXgbCorrBP}{}\renewcommand{\rhchXgbCorrBP}{\ensuremath{<}0.001}  % xgb: full without FOMC vs HAR, every bar: DM p
\providecommand{\rhchXgbCorrBMsePct}{}\renewcommand{\rhchXgbCorrBMsePct}{17.9}  % xgb: full without FOMC vs HAR, every bar: % MSE change
\providecommand{\rhchXgbCorrBMseP}{}\renewcommand{\rhchXgbCorrBMseP}{0.268}  % xgb: full without FOMC vs HAR, every bar: DM p on MSE
\providecommand{\rhchXgbCorrCPct}{}\renewcommand{\rhchXgbCorrCPct}{$-$6.1}  % xgb: full vs HAR, ORDINARY days: % QLIKE change
\providecommand{\rhchXgbCorrCDm}{}\renewcommand{\rhchXgbCorrCDm}{$-$10.42}  % xgb: full vs HAR, ORDINARY days: DM
\providecommand{\rhchXgbCorrCP}{}\renewcommand{\rhchXgbCorrCP}{\ensuremath{<}0.001}  % xgb: full vs HAR, ORDINARY days: DM p
\providecommand{\rhchXgbCorrCMsePct}{}\renewcommand{\rhchXgbCorrCMsePct}{18.9}  % xgb: full vs HAR, ORDINARY days: % MSE change
\providecommand{\rhchXgbCorrCMseP}{}\renewcommand{\rhchXgbCorrCMseP}{0.268}  % xgb: full vs HAR, ORDINARY days: DM p on MSE
\providecommand{\rhchXgbCorrAgree}{}\renewcommand{\rhchXgbCorrAgree}{agree}  % xgb: the three corrections within 2.0 pp?
\providecommand{\rhchXgbCorrSpread}{}\renewcommand{\rhchXgbCorrSpread}{0.4}  % xgb: max - min of the three (pp)
\providecommand{\rhchRfTabOnePct}{}\renewcommand{\rhchRfTabOnePct}{$-$11.3}  % rf: full vs HAR, every bar (Table 1): % QLIKE change
\providecommand{\rhchRfTabOneDm}{}\renewcommand{\rhchRfTabOneDm}{$-$9.74}  % rf: full vs HAR, every bar (Table 1): DM
\providecommand{\rhchRfTabOneP}{}\renewcommand{\rhchRfTabOneP}{\ensuremath{<}0.001}  % rf: full vs HAR, every bar (Table 1): DM p
\providecommand{\rhchRfTabOneMsePct}{}\renewcommand{\rhchRfTabOneMsePct}{10.9}  % rf: full vs HAR, every bar (Table 1): % MSE change
\providecommand{\rhchRfTabOneMseP}{}\renewcommand{\rhchRfTabOneMseP}{0.271}  % rf: full vs HAR, every bar (Table 1): DM p on MSE
\providecommand{\rhchRfCorrAPct}{}\renewcommand{\rhchRfCorrAPct}{$-$4.0}  % rf: full vs HAR+FOMC, every bar: % QLIKE change
\providecommand{\rhchRfCorrADm}{}\renewcommand{\rhchRfCorrADm}{$-$9.06}  % rf: full vs HAR+FOMC, every bar: DM
\providecommand{\rhchRfCorrAP}{}\renewcommand{\rhchRfCorrAP}{\ensuremath{<}0.001}  % rf: full vs HAR+FOMC, every bar: DM p
\providecommand{\rhchRfCorrAMsePct}{}\renewcommand{\rhchRfCorrAMsePct}{11.3}  % rf: full vs HAR+FOMC, every bar: % MSE change
\providecommand{\rhchRfCorrAMseP}{}\renewcommand{\rhchRfCorrAMseP}{0.255}  % rf: full vs HAR+FOMC, every bar: DM p on MSE
\providecommand{\rhchRfCorrBPct}{}\renewcommand{\rhchRfCorrBPct}{--}  % rf: full without FOMC vs HAR, every bar: not run
\providecommand{\rhchRfCorrBDm}{}\renewcommand{\rhchRfCorrBDm}{--}  % rf: full without FOMC vs HAR, every bar: not run
\providecommand{\rhchRfCorrBP}{}\renewcommand{\rhchRfCorrBP}{--}  % rf: full without FOMC vs HAR, every bar: not run
\providecommand{\rhchRfCorrBMsePct}{}\renewcommand{\rhchRfCorrBMsePct}{--}  % rf: full without FOMC vs HAR, every bar: not run
\providecommand{\rhchRfCorrBMseP}{}\renewcommand{\rhchRfCorrBMseP}{--}  % rf: full without FOMC vs HAR, every bar: not run
\providecommand{\rhchRfCorrCPct}{}\renewcommand{\rhchRfCorrCPct}{$-$5.3}  % rf: full vs HAR, ORDINARY days: % QLIKE change
\providecommand{\rhchRfCorrCDm}{}\renewcommand{\rhchRfCorrCDm}{$-$11.16}  % rf: full vs HAR, ORDINARY days: DM
\providecommand{\rhchRfCorrCP}{}\renewcommand{\rhchRfCorrCP}{\ensuremath{<}0.001}  % rf: full vs HAR, ORDINARY days: DM p
\providecommand{\rhchRfCorrCMsePct}{}\renewcommand{\rhchRfCorrCMsePct}{12.4}  % rf: full vs HAR, ORDINARY days: % MSE change
\providecommand{\rhchRfCorrCMseP}{}\renewcommand{\rhchRfCorrCMseP}{0.247}  % rf: full vs HAR, ORDINARY days: DM p on MSE
\providecommand{\rhchRfCorrAgree}{}\renewcommand{\rhchRfCorrAgree}{--}  % rf: the three corrections within 2.0 pp?
\providecommand{\rhchRfCorrSpread}{}\renewcommand{\rhchRfCorrSpread}{1.3}  % rf: max - min of the three (pp)
\providecommand{\rhchNCorrAgree}{}\renewcommand{\rhchNCorrAgree}{5}  % classes whose corrections agree
\providecommand{\rhchNCorrDisagree}{}\renewcommand{\rhchNCorrDisagree}{0}  % classes whose corrections disagree
\providecommand{\rhchNCorrPending}{}\renewcommand{\rhchNCorrPending}{0}  % classes with a correction pending
\providecommand{\rhchRidgeShareEvent}{}\renewcommand{\rhchRidgeShareEvent}{54}  % ridge: % of the every-bar gain on EVENT days
\providecommand{\rhchRidgeShareEventPm}{}\renewcommand{\rhchRidgeShareEventPm}{56}  % ridge: % of the gain on EVENT days 14:30-16:00
\providecommand{\rhchRidgeEventPeakBar}{}\renewcommand{\rhchRidgeEventPeakBar}{14:30}  % ridge: event: bar end with the most negative mean diff
\providecommand{\rhchRidgeEventPeakDiff}{}\renewcommand{\rhchRidgeEventPeakDiff}{$-$1583.6}  % ridge: event: that mean diff (full - HAR) x 10^3
\providecommand{\rhchRidgeOrdPeakBar}{}\renewcommand{\rhchRidgeOrdPeakBar}{14:30}  % ridge: ord: bar end with the most negative mean diff
\providecommand{\rhchRidgeOrdPeakDiff}{}\renewcommand{\rhchRidgeOrdPeakDiff}{$-$15.9}  % ridge: ord: that mean diff (full - HAR) x 10^3
\providecommand{\rhchLassoShareEvent}{}\renewcommand{\rhchLassoShareEvent}{77}  % lasso: % of the every-bar gain on EVENT days
\providecommand{\rhchLassoShareEventPm}{}\renewcommand{\rhchLassoShareEventPm}{78}  % lasso: % of the gain on EVENT days 14:30-16:00
\providecommand{\rhchLassoEventPeakBar}{}\renewcommand{\rhchLassoEventPeakBar}{14:30}  % lasso: event: bar end with the most negative mean diff
\providecommand{\rhchLassoEventPeakDiff}{}\renewcommand{\rhchLassoEventPeakDiff}{$-$4302.3}  % lasso: event: that mean diff (full - HAR) x 10^3
\providecommand{\rhchLassoOrdPeakBar}{}\renewcommand{\rhchLassoOrdPeakBar}{14:30}  % lasso: ord: bar end with the most negative mean diff
\providecommand{\rhchLassoOrdPeakDiff}{}\renewcommand{\rhchLassoOrdPeakDiff}{$-$16.4}  % lasso: ord: that mean diff (full - HAR) x 10^3
\providecommand{\rhchEnetShareEvent}{}\renewcommand{\rhchEnetShareEvent}{49}  % enet: % of the every-bar gain on EVENT days
\providecommand{\rhchEnetShareEventPm}{}\renewcommand{\rhchEnetShareEventPm}{51}  % enet: % of the gain on EVENT days 14:30-16:00
\providecommand{\rhchEnetEventPeakBar}{}\renewcommand{\rhchEnetEventPeakBar}{14:30}  % enet: event: bar end with the most negative mean diff
\providecommand{\rhchEnetEventPeakDiff}{}\renewcommand{\rhchEnetEventPeakDiff}{$-$1610.3}  % enet: event: that mean diff (full - HAR) x 10^3
\providecommand{\rhchEnetOrdPeakBar}{}\renewcommand{\rhchEnetOrdPeakBar}{15:00}  % enet: ord: bar end with the most negative mean diff
\providecommand{\rhchEnetOrdPeakDiff}{}\renewcommand{\rhchEnetOrdPeakDiff}{$-$15.8}  % enet: ord: that mean diff (full - HAR) x 10^3
\providecommand{\rhchLgbmShareEvent}{}\renewcommand{\rhchLgbmShareEvent}{63}  % lgbm: % of the every-bar gain on EVENT days
\providecommand{\rhchLgbmShareEventPm}{}\renewcommand{\rhchLgbmShareEventPm}{63}  % lgbm: % of the gain on EVENT days 14:30-16:00
\providecommand{\rhchLgbmEventPeakBar}{}\renewcommand{\rhchLgbmEventPeakBar}{14:30}  % lgbm: event: bar end with the most negative mean diff
\providecommand{\rhchLgbmEventPeakDiff}{}\renewcommand{\rhchLgbmEventPeakDiff}{$-$6012.1}  % lgbm: event: that mean diff (full - HAR) x 10^3
\providecommand{\rhchLgbmOrdPeakBar}{}\renewcommand{\rhchLgbmOrdPeakBar}{12:30}  % lgbm: ord: bar end with the most negative mean diff
\providecommand{\rhchLgbmOrdPeakDiff}{}\renewcommand{\rhchLgbmOrdPeakDiff}{$-$15.9}  % lgbm: ord: that mean diff (full - HAR) x 10^3
\providecommand{\rhchXgbShareEvent}{}\renewcommand{\rhchXgbShareEvent}{64}  % xgb: % of the every-bar gain on EVENT days
\providecommand{\rhchXgbShareEventPm}{}\renewcommand{\rhchXgbShareEventPm}{63}  % xgb: % of the gain on EVENT days 14:30-16:00
\providecommand{\rhchXgbEventPeakBar}{}\renewcommand{\rhchXgbEventPeakBar}{14:30}  % xgb: event: bar end with the most negative mean diff
\providecommand{\rhchXgbEventPeakDiff}{}\renewcommand{\rhchXgbEventPeakDiff}{$-$6274.7}  % xgb: event: that mean diff (full - HAR) x 10^3
\providecommand{\rhchXgbOrdPeakBar}{}\renewcommand{\rhchXgbOrdPeakBar}{12:30}  % xgb: ord: bar end with the most negative mean diff
\providecommand{\rhchXgbOrdPeakDiff}{}\renewcommand{\rhchXgbOrdPeakDiff}{$-$16.3}  % xgb: ord: that mean diff (full - HAR) x 10^3
\providecommand{\rhchRfShareEvent}{}\renewcommand{\rhchRfShareEvent}{60}  % rf: % of the every-bar gain on EVENT days
\providecommand{\rhchRfShareEventPm}{}\renewcommand{\rhchRfShareEventPm}{59}  % rf: % of the gain on EVENT days 14:30-16:00
\providecommand{\rhchRfEventPeakBar}{}\renewcommand{\rhchRfEventPeakBar}{14:30}  % rf: event: bar end with the most negative mean diff
\providecommand{\rhchRfEventPeakDiff}{}\renewcommand{\rhchRfEventPeakDiff}{$-$4773.7}  % rf: event: that mean diff (full - HAR) x 10^3
\providecommand{\rhchRfOrdPeakBar}{}\renewcommand{\rhchRfOrdPeakBar}{12:30}  % rf: ord: bar end with the most negative mean diff
\providecommand{\rhchRfOrdPeakDiff}{}\renewcommand{\rhchRfOrdPeakDiff}{$-$15.2}  % rf: ord: that mean diff (full - HAR) x 10^3
\providecommand{\rhchEventPmBarPct}{}\renewcommand{\rhchEventPmBarPct}{1.0}  % EVENT 14:30-16:00 bars, % of all bars
\providecommand{\rhchLinShareEventMin}{}\renewcommand{\rhchLinShareEventMin}{49}  % Lin: lowest share_event_pct
\providecommand{\rhchLinShareEventMax}{}\renewcommand{\rhchLinShareEventMax}{77}  % Lin: highest share_event_pct
\providecommand{\rhchLinShareEventPmMin}{}\renewcommand{\rhchLinShareEventPmMin}{51}  % Lin: lowest share_event_pm_pct
\providecommand{\rhchLinShareEventPmMax}{}\renewcommand{\rhchLinShareEventPmMax}{78}  % Lin: highest share_event_pm_pct
\providecommand{\rhchTreeShareEventMin}{}\renewcommand{\rhchTreeShareEventMin}{60}  % Tree: lowest share_event_pct
\providecommand{\rhchTreeShareEventMax}{}\renewcommand{\rhchTreeShareEventMax}{64}  % Tree: highest share_event_pct
\providecommand{\rhchTreeShareEventPmMin}{}\renewcommand{\rhchTreeShareEventPmMin}{59}  % Tree: lowest share_event_pm_pct
\providecommand{\rhchTreeShareEventPmMax}{}\renewcommand{\rhchTreeShareEventPmMax}{63}  % Tree: highest share_event_pm_pct
\providecommand{\rhchLassoRidgeBarPct}{}\renewcommand{\rhchLassoRidgeBarPct}{$-$3.1}  % lasso vs ridge, full sets, all: % QLIKE change
\providecommand{\rhchLassoRidgeBarDm}{}\renewcommand{\rhchLassoRidgeBarDm}{$-$4.82}  % lasso vs ridge, all: DM
\providecommand{\rhchLassoRidgeBarP}{}\renewcommand{\rhchLassoRidgeBarP}{\ensuremath{<}0.001}  % lasso vs ridge, all: DM p
\providecommand{\rhchLassoRidgeEventPct}{}\renewcommand{\rhchLassoRidgeEventPct}{$-$31.9}  % lasso vs ridge, full sets, event: % QLIKE change
\providecommand{\rhchLassoRidgeEventDm}{}\renewcommand{\rhchLassoRidgeEventDm}{$-$8.42}  % lasso vs ridge, event: DM
\providecommand{\rhchLassoRidgeEventP}{}\renewcommand{\rhchLassoRidgeEventP}{\ensuremath{<}0.001}  % lasso vs ridge, event: DM p
\providecommand{\rhchLassoRidgeOrdPct}{}\renewcommand{\rhchLassoRidgeOrdPct}{0.8}  % lasso vs ridge, full sets, ord: % QLIKE change
\providecommand{\rhchLassoRidgeOrdDm}{}\renewcommand{\rhchLassoRidgeOrdDm}{2.28}  % lasso vs ridge, ord: DM
\providecommand{\rhchLassoRidgeOrdP}{}\renewcommand{\rhchLassoRidgeOrdP}{0.023}  % lasso vs ridge, ord: DM p
\providecommand{\rhchLassoEnetBarPct}{}\renewcommand{\rhchLassoEnetBarPct}{$-$2.4}  % lasso vs enet, full sets, all: % QLIKE change
\providecommand{\rhchLassoEnetBarDm}{}\renewcommand{\rhchLassoEnetBarDm}{$-$3.49}  % lasso vs enet, all: DM
\providecommand{\rhchLassoEnetBarP}{}\renewcommand{\rhchLassoEnetBarP}{\ensuremath{<}0.001}  % lasso vs enet, all: DM p
\providecommand{\rhchLassoEnetEventPct}{}\renewcommand{\rhchLassoEnetEventPct}{$-$31.8}  % lasso vs enet, full sets, event: % QLIKE change
\providecommand{\rhchLassoEnetEventDm}{}\renewcommand{\rhchLassoEnetEventDm}{$-$7.69}  % lasso vs enet, event: DM
\providecommand{\rhchLassoEnetEventP}{}\renewcommand{\rhchLassoEnetEventP}{\ensuremath{<}0.001}  % lasso vs enet, event: DM p
\providecommand{\rhchLassoEnetOrdPct}{}\renewcommand{\rhchLassoEnetOrdPct}{1.6}  % lasso vs enet, full sets, ord: % QLIKE change
\providecommand{\rhchLassoEnetOrdDm}{}\renewcommand{\rhchLassoEnetOrdDm}{3.35}  % lasso vs enet, ord: DM
\providecommand{\rhchLassoEnetOrdP}{}\renewcommand{\rhchLassoEnetOrdP}{\ensuremath{<}0.001}  % lasso vs enet, ord: DM p
\providecommand{\rhchRidgeBarJointPct}{}\renewcommand{\rhchRidgeBarJointPct}{$-$6.0}  % ridge all: full vs HAR, % QLIKE change
\providecommand{\rhchRidgeBarJointDm}{}\renewcommand{\rhchRidgeBarJointDm}{$-$6.28}  % ridge all: full vs HAR, DM
\providecommand{\rhchRidgeBarJointP}{}\renewcommand{\rhchRidgeBarJointP}{\ensuremath{<}0.001}  % ridge all: full vs HAR, DM p
\providecommand{\rhchRidgeBarAddRmPct}{}\renewcommand{\rhchRidgeBarAddRmPct}{$-$2.6}  % ridge all add_rm vs ols|har: %
\providecommand{\rhchRidgeBarAddRmP}{}\renewcommand{\rhchRidgeBarAddRmP}{\ensuremath{<}0.001}  % ridge all add_rm: DM p
\providecommand{\rhchRidgeBarLoRmPct}{}\renewcommand{\rhchRidgeBarLoRmPct}{1.8}  % ridge all lo_rm vs ridge|full: %
\providecommand{\rhchRidgeBarLoRmP}{}\renewcommand{\rhchRidgeBarLoRmP}{\ensuremath{<}0.001}  % ridge all lo_rm: DM p
\providecommand{\rhchRidgeBarAddLiqPct}{}\renewcommand{\rhchRidgeBarAddLiqPct}{$-$0.4}  % ridge all add_liq vs ols|har: %
\providecommand{\rhchRidgeBarAddLiqP}{}\renewcommand{\rhchRidgeBarAddLiqP}{0.326}  % ridge all add_liq: DM p
\providecommand{\rhchRidgeBarLoLiqPct}{}\renewcommand{\rhchRidgeBarLoLiqPct}{0.2}  % ridge all lo_liq vs ridge|full: %
\providecommand{\rhchRidgeBarLoLiqP}{}\renewcommand{\rhchRidgeBarLoLiqP}{0.467}  % ridge all lo_liq: DM p
\providecommand{\rhchRidgeBarAddEwPct}{}\renewcommand{\rhchRidgeBarAddEwPct}{$-$0.9}  % ridge all add_ew vs ols|har: %
\providecommand{\rhchRidgeBarAddEwP}{}\renewcommand{\rhchRidgeBarAddEwP}{0.024}  % ridge all add_ew: DM p
\providecommand{\rhchRidgeBarLoEwPct}{}\renewcommand{\rhchRidgeBarLoEwPct}{$-$0.1}  % ridge all lo_ew vs ridge|full: %
\providecommand{\rhchRidgeBarLoEwP}{}\renewcommand{\rhchRidgeBarLoEwP}{0.381}  % ridge all lo_ew: DM p
\providecommand{\rhchRidgeBarAddVwPct}{}\renewcommand{\rhchRidgeBarAddVwPct}{$-$0.9}  % ridge all add_vw vs ols|har: %
\providecommand{\rhchRidgeBarAddVwP}{}\renewcommand{\rhchRidgeBarAddVwP}{0.114}  % ridge all add_vw: DM p
\providecommand{\rhchRidgeBarLoVwPct}{}\renewcommand{\rhchRidgeBarLoVwPct}{$-$0.1}  % ridge all lo_vw vs ridge|full: %
\providecommand{\rhchRidgeBarLoVwP}{}\renewcommand{\rhchRidgeBarLoVwP}{0.487}  % ridge all lo_vw: DM p
\providecommand{\rhchRidgeBarAddSentPct}{}\renewcommand{\rhchRidgeBarAddSentPct}{0.2}  % ridge all add_sent vs ols|har: %
\providecommand{\rhchRidgeBarAddSentP}{}\renewcommand{\rhchRidgeBarAddSentP}{0.411}  % ridge all add_sent: DM p
\providecommand{\rhchRidgeBarLoSentPct}{}\renewcommand{\rhchRidgeBarLoSentPct}{0.1}  % ridge all lo_sent vs ridge|full: %
\providecommand{\rhchRidgeBarLoSentP}{}\renewcommand{\rhchRidgeBarLoSentP}{0.637}  % ridge all lo_sent: DM p
\providecommand{\rhchRidgeBarAddIvPct}{}\renewcommand{\rhchRidgeBarAddIvPct}{0.2}  % ridge all add_iv vs ols|har: %
\providecommand{\rhchRidgeBarAddIvP}{}\renewcommand{\rhchRidgeBarAddIvP}{0.506}  % ridge all add_iv: DM p
\providecommand{\rhchRidgeBarLoIvPct}{}\renewcommand{\rhchRidgeBarLoIvPct}{0.1}  % ridge all lo_iv vs ridge|full: %
\providecommand{\rhchRidgeBarLoIvP}{}\renewcommand{\rhchRidgeBarLoIvP}{0.675}  % ridge all lo_iv: DM p
\providecommand{\rhchRidgeBarAddVdPct}{}\renewcommand{\rhchRidgeBarAddVdPct}{0.5}  % ridge all add_vd vs ols|har: %
\providecommand{\rhchRidgeBarAddVdP}{}\renewcommand{\rhchRidgeBarAddVdP}{0.002}  % ridge all add_vd: DM p
\providecommand{\rhchRidgeBarLoVdPct}{}\renewcommand{\rhchRidgeBarLoVdPct}{$-$0.1}  % ridge all lo_vd vs ridge|full: %
\providecommand{\rhchRidgeBarLoVdP}{}\renewcommand{\rhchRidgeBarLoVdP}{0.472}  % ridge all lo_vd: DM p
\providecommand{\rhchRidgeBarAddFomcPct}{}\renewcommand{\rhchRidgeBarAddFomcPct}{$-$3.1}  % ridge all add_fomc vs ols|har: %
\providecommand{\rhchRidgeBarAddFomcP}{}\renewcommand{\rhchRidgeBarAddFomcP}{\ensuremath{<}0.001}  % ridge all add_fomc: DM p
\providecommand{\rhchRidgeBarLoFomcPct}{}\renewcommand{\rhchRidgeBarLoFomcPct}{3.1}  % ridge all lo_fomc vs ridge|full: %
\providecommand{\rhchRidgeBarLoFomcP}{}\renewcommand{\rhchRidgeBarLoFomcP}{\ensuremath{<}0.001}  % ridge all lo_fomc: DM p
\providecommand{\rhchRidgeEventJointPct}{}\renewcommand{\rhchRidgeEventJointPct}{$-$22.4}  % ridge event: full vs HAR, % QLIKE change
\providecommand{\rhchRidgeEventJointDm}{}\renewcommand{\rhchRidgeEventJointDm}{$-$6.45}  % ridge event: full vs HAR, DM
\providecommand{\rhchRidgeEventJointP}{}\renewcommand{\rhchRidgeEventJointP}{\ensuremath{<}0.001}  % ridge event: full vs HAR, DM p
\providecommand{\rhchRidgeEventAddRmPct}{}\renewcommand{\rhchRidgeEventAddRmPct}{4.4}  % ridge event add_rm vs ols|har: %
\providecommand{\rhchRidgeEventAddRmP}{}\renewcommand{\rhchRidgeEventAddRmP}{0.023}  % ridge event add_rm: DM p
\providecommand{\rhchRidgeEventLoRmPct}{}\renewcommand{\rhchRidgeEventLoRmPct}{$-$1.1}  % ridge event lo_rm vs ridge|full: %
\providecommand{\rhchRidgeEventLoRmP}{}\renewcommand{\rhchRidgeEventLoRmP}{0.133}  % ridge event lo_rm: DM p
\providecommand{\rhchRidgeEventAddLiqPct}{}\renewcommand{\rhchRidgeEventAddLiqPct}{5.4}  % ridge event add_liq vs ols|har: %
\providecommand{\rhchRidgeEventAddLiqP}{}\renewcommand{\rhchRidgeEventAddLiqP}{0.010}  % ridge event add_liq: DM p
\providecommand{\rhchRidgeEventLoLiqPct}{}\renewcommand{\rhchRidgeEventLoLiqPct}{$-$0.6}  % ridge event lo_liq vs ridge|full: %
\providecommand{\rhchRidgeEventLoLiqP}{}\renewcommand{\rhchRidgeEventLoLiqP}{0.230}  % ridge event lo_liq: DM p
\providecommand{\rhchRidgeEventAddEwPct}{}\renewcommand{\rhchRidgeEventAddEwPct}{$-$1.8}  % ridge event add_ew vs ols|har: %
\providecommand{\rhchRidgeEventAddEwP}{}\renewcommand{\rhchRidgeEventAddEwP}{0.176}  % ridge event add_ew: DM p
\providecommand{\rhchRidgeEventLoEwPct}{}\renewcommand{\rhchRidgeEventLoEwPct}{0.0}  % ridge event lo_ew vs ridge|full: %
\providecommand{\rhchRidgeEventLoEwP}{}\renewcommand{\rhchRidgeEventLoEwP}{0.968}  % ridge event lo_ew: DM p
\providecommand{\rhchRidgeEventAddVwPct}{}\renewcommand{\rhchRidgeEventAddVwPct}{$-$2.5}  % ridge event add_vw vs ols|har: %
\providecommand{\rhchRidgeEventAddVwP}{}\renewcommand{\rhchRidgeEventAddVwP}{0.143}  % ridge event add_vw: DM p
\providecommand{\rhchRidgeEventLoVwPct}{}\renewcommand{\rhchRidgeEventLoVwPct}{0.1}  % ridge event lo_vw vs ridge|full: %
\providecommand{\rhchRidgeEventLoVwP}{}\renewcommand{\rhchRidgeEventLoVwP}{0.704}  % ridge event lo_vw: DM p
\providecommand{\rhchRidgeEventAddSentPct}{}\renewcommand{\rhchRidgeEventAddSentPct}{$-$0.1}  % ridge event add_sent vs ols|har: %
\providecommand{\rhchRidgeEventAddSentP}{}\renewcommand{\rhchRidgeEventAddSentP}{0.970}  % ridge event add_sent: DM p
\providecommand{\rhchRidgeEventLoSentPct}{}\renewcommand{\rhchRidgeEventLoSentPct}{1.7}  % ridge event lo_sent vs ridge|full: %
\providecommand{\rhchRidgeEventLoSentP}{}\renewcommand{\rhchRidgeEventLoSentP}{0.322}  % ridge event lo_sent: DM p
\providecommand{\rhchRidgeEventAddIvPct}{}\renewcommand{\rhchRidgeEventAddIvPct}{2.8}  % ridge event add_iv vs ols|har: %
\providecommand{\rhchRidgeEventAddIvP}{}\renewcommand{\rhchRidgeEventAddIvP}{0.011}  % ridge event add_iv: DM p
\providecommand{\rhchRidgeEventLoIvPct}{}\renewcommand{\rhchRidgeEventLoIvPct}{$-$0.2}  % ridge event lo_iv vs ridge|full: %
\providecommand{\rhchRidgeEventLoIvP}{}\renewcommand{\rhchRidgeEventLoIvP}{0.823}  % ridge event lo_iv: DM p
\providecommand{\rhchRidgeEventAddVdPct}{}\renewcommand{\rhchRidgeEventAddVdPct}{1.9}  % ridge event add_vd vs ols|har: %
\providecommand{\rhchRidgeEventAddVdP}{}\renewcommand{\rhchRidgeEventAddVdP}{0.028}  % ridge event add_vd: DM p
\providecommand{\rhchRidgeEventLoVdPct}{}\renewcommand{\rhchRidgeEventLoVdPct}{0.0}  % ridge event lo_vd vs ridge|full: %
\providecommand{\rhchRidgeEventLoVdP}{}\renewcommand{\rhchRidgeEventLoVdP}{0.971}  % ridge event lo_vd: DM p
\providecommand{\rhchRidgeEventAddFomcPct}{}\renewcommand{\rhchRidgeEventAddFomcPct}{$-$23.4}  % ridge event add_fomc vs ols|har: %
\providecommand{\rhchRidgeEventAddFomcP}{}\renewcommand{\rhchRidgeEventAddFomcP}{\ensuremath{<}0.001}  % ridge event add_fomc: DM p
\providecommand{\rhchRidgeEventLoFomcPct}{}\renewcommand{\rhchRidgeEventLoFomcPct}{28.6}  % ridge event lo_fomc vs ridge|full: %
\providecommand{\rhchRidgeEventLoFomcP}{}\renewcommand{\rhchRidgeEventLoFomcP}{\ensuremath{<}0.001}  % ridge event lo_fomc: DM p
\providecommand{\rhchRidgeEventReading}{}\renewcommand{\rhchRidgeEventReading}{concentrated}  % ridge event: item 3 reading
\providecommand{\rhchRidgeEventTopAddShare}{}\renewcommand{\rhchRidgeEventTopAddShare}{104}  % ridge event: largest add-alone gain, % of the joint gain (fomc)
\providecommand{\rhchRidgeEventSumAddRatio}{}\renewcommand{\rhchRidgeEventSumAddRatio}{0.60}  % ridge event: add-alone gains summed over groups / joint gain
\providecommand{\rhchRidgeEventNLoSig}{}\renewcommand{\rhchRidgeEventNLoSig}{1}  % ridge event: significant leave-one-outs
\providecommand{\rhchRidgeEventNLoRun}{}\renewcommand{\rhchRidgeEventNLoRun}{8}  % ridge event: leave-one-outs available (of 8)
\providecommand{\rhchRidgeEventConcTest}{}\renewcommand{\rhchRidgeEventConcTest}{yes}  % ridge event: the CONC definition on its own (pending: undetermined)
\providecommand{\rhchRidgeEventRedTest}{}\renewcommand{\rhchRidgeEventRedTest}{no}  % ridge event: the RED definition on its own (pending: undetermined)
\providecommand{\rhchRidgeEventReadingRedFirst}{}\renewcommand{\rhchRidgeEventReadingRedFirst}{concentrated}  % ridge event: reading with REDUNDANT tested before CONCENTRATED (the RULE tests CONCENTRATED first)
\providecommand{\rhchRidgeOrdJointPct}{}\renewcommand{\rhchRidgeOrdJointPct}{$-$3.2}  % ridge ord: full vs HAR, % QLIKE change
\providecommand{\rhchRidgeOrdJointDm}{}\renewcommand{\rhchRidgeOrdJointDm}{$-$3.51}  % ridge ord: full vs HAR, DM
\providecommand{\rhchRidgeOrdJointP}{}\renewcommand{\rhchRidgeOrdJointP}{\ensuremath{<}0.001}  % ridge ord: full vs HAR, DM p
\providecommand{\rhchRidgeOrdAddRmPct}{}\renewcommand{\rhchRidgeOrdAddRmPct}{$-$3.8}  % ridge ord add_rm vs ols|har: %
\providecommand{\rhchRidgeOrdAddRmP}{}\renewcommand{\rhchRidgeOrdAddRmP}{\ensuremath{<}0.001}  % ridge ord add_rm: DM p
\providecommand{\rhchRidgeOrdLoRmPct}{}\renewcommand{\rhchRidgeOrdLoRmPct}{2.2}  % ridge ord lo_rm vs ridge|full: %
\providecommand{\rhchRidgeOrdLoRmP}{}\renewcommand{\rhchRidgeOrdLoRmP}{\ensuremath{<}0.001}  % ridge ord lo_rm: DM p
\providecommand{\rhchRidgeOrdAddLiqPct}{}\renewcommand{\rhchRidgeOrdAddLiqPct}{$-$1.4}  % ridge ord add_liq vs ols|har: %
\providecommand{\rhchRidgeOrdAddLiqP}{}\renewcommand{\rhchRidgeOrdAddLiqP}{\ensuremath{<}0.001}  % ridge ord add_liq: DM p
\providecommand{\rhchRidgeOrdLoLiqPct}{}\renewcommand{\rhchRidgeOrdLoLiqPct}{0.3}  % ridge ord lo_liq vs ridge|full: %
\providecommand{\rhchRidgeOrdLoLiqP}{}\renewcommand{\rhchRidgeOrdLoLiqP}{0.261}  % ridge ord lo_liq: DM p
\providecommand{\rhchRidgeOrdAddEwPct}{}\renewcommand{\rhchRidgeOrdAddEwPct}{$-$0.7}  % ridge ord add_ew vs ols|har: %
\providecommand{\rhchRidgeOrdAddEwP}{}\renewcommand{\rhchRidgeOrdAddEwP}{0.044}  % ridge ord add_ew: DM p
\providecommand{\rhchRidgeOrdLoEwPct}{}\renewcommand{\rhchRidgeOrdLoEwPct}{$-$0.1}  % ridge ord lo_ew vs ridge|full: %
\providecommand{\rhchRidgeOrdLoEwP}{}\renewcommand{\rhchRidgeOrdLoEwP}{0.316}  % ridge ord lo_ew: DM p
\providecommand{\rhchRidgeOrdAddVwPct}{}\renewcommand{\rhchRidgeOrdAddVwPct}{$-$0.6}  % ridge ord add_vw vs ols|har: %
\providecommand{\rhchRidgeOrdAddVwP}{}\renewcommand{\rhchRidgeOrdAddVwP}{0.309}  % ridge ord add_vw: DM p
\providecommand{\rhchRidgeOrdLoVwPct}{}\renewcommand{\rhchRidgeOrdLoVwPct}{$-$0.1}  % ridge ord lo_vw vs ridge|full: %
\providecommand{\rhchRidgeOrdLoVwP}{}\renewcommand{\rhchRidgeOrdLoVwP}{0.405}  % ridge ord lo_vw: DM p
\providecommand{\rhchRidgeOrdAddSentPct}{}\renewcommand{\rhchRidgeOrdAddSentPct}{0.3}  % ridge ord add_sent vs ols|har: %
\providecommand{\rhchRidgeOrdAddSentP}{}\renewcommand{\rhchRidgeOrdAddSentP}{0.246}  % ridge ord add_sent: DM p
\providecommand{\rhchRidgeOrdLoSentPct}{}\renewcommand{\rhchRidgeOrdLoSentPct}{$-$0.1}  % ridge ord lo_sent vs ridge|full: %
\providecommand{\rhchRidgeOrdLoSentP}{}\renewcommand{\rhchRidgeOrdLoSentP}{0.271}  % ridge ord lo_sent: DM p
\providecommand{\rhchRidgeOrdAddIvPct}{}\renewcommand{\rhchRidgeOrdAddIvPct}{$-$0.2}  % ridge ord add_iv vs ols|har: %
\providecommand{\rhchRidgeOrdAddIvP}{}\renewcommand{\rhchRidgeOrdAddIvP}{0.580}  % ridge ord add_iv: DM p
\providecommand{\rhchRidgeOrdLoIvPct}{}\renewcommand{\rhchRidgeOrdLoIvPct}{0.2}  % ridge ord lo_iv vs ridge|full: %
\providecommand{\rhchRidgeOrdLoIvP}{}\renewcommand{\rhchRidgeOrdLoIvP}{0.611}  % ridge ord lo_iv: DM p
\providecommand{\rhchRidgeOrdAddVdPct}{}\renewcommand{\rhchRidgeOrdAddVdPct}{0.3}  % ridge ord add_vd vs ols|har: %
\providecommand{\rhchRidgeOrdAddVdP}{}\renewcommand{\rhchRidgeOrdAddVdP}{0.069}  % ridge ord add_vd: DM p
\providecommand{\rhchRidgeOrdLoVdPct}{}\renewcommand{\rhchRidgeOrdLoVdPct}{$-$0.1}  % ridge ord lo_vd vs ridge|full: %
\providecommand{\rhchRidgeOrdLoVdP}{}\renewcommand{\rhchRidgeOrdLoVdP}{0.437}  % ridge ord lo_vd: DM p
\providecommand{\rhchRidgeOrdAddFomcPct}{}\renewcommand{\rhchRidgeOrdAddFomcPct}{0.4}  % ridge ord add_fomc vs ols|har: %
\providecommand{\rhchRidgeOrdAddFomcP}{}\renewcommand{\rhchRidgeOrdAddFomcP}{\ensuremath{<}0.001}  % ridge ord add_fomc: DM p
\providecommand{\rhchRidgeOrdLoFomcPct}{}\renewcommand{\rhchRidgeOrdLoFomcPct}{$-$0.4}  % ridge ord lo_fomc vs ridge|full: %
\providecommand{\rhchRidgeOrdLoFomcP}{}\renewcommand{\rhchRidgeOrdLoFomcP}{\ensuremath{<}0.001}  % ridge ord lo_fomc: DM p
\providecommand{\rhchRidgeOrdReading}{}\renewcommand{\rhchRidgeOrdReading}{concentrated}  % ridge ord: item 3 reading
\providecommand{\rhchRidgeOrdTopAddShare}{}\renewcommand{\rhchRidgeOrdTopAddShare}{118}  % ridge ord: largest add-alone gain, % of the joint gain (rm)
\providecommand{\rhchRidgeOrdSumAddRatio}{}\renewcommand{\rhchRidgeOrdSumAddRatio}{1.79}  % ridge ord: add-alone gains summed over groups / joint gain
\providecommand{\rhchRidgeOrdNLoSig}{}\renewcommand{\rhchRidgeOrdNLoSig}{1}  % ridge ord: significant leave-one-outs
\providecommand{\rhchRidgeOrdNLoRun}{}\renewcommand{\rhchRidgeOrdNLoRun}{8}  % ridge ord: leave-one-outs available (of 8)
\providecommand{\rhchRidgeOrdConcTest}{}\renewcommand{\rhchRidgeOrdConcTest}{yes}  % ridge ord: the CONC definition on its own (pending: undetermined)
\providecommand{\rhchRidgeOrdRedTest}{}\renewcommand{\rhchRidgeOrdRedTest}{yes}  % ridge ord: the RED definition on its own (pending: undetermined)
\providecommand{\rhchRidgeOrdReadingRedFirst}{}\renewcommand{\rhchRidgeOrdReadingRedFirst}{redundant}  % ridge ord: reading with REDUNDANT tested before CONCENTRATED (the RULE tests CONCENTRATED first)
\providecommand{\rhchLassoBarJointPct}{}\renewcommand{\rhchLassoBarJointPct}{$-$9.0}  % lasso all: full vs HAR, % QLIKE change
\providecommand{\rhchLassoBarJointDm}{}\renewcommand{\rhchLassoBarJointDm}{$-$6.34}  % lasso all: full vs HAR, DM
\providecommand{\rhchLassoBarJointP}{}\renewcommand{\rhchLassoBarJointP}{\ensuremath{<}0.001}  % lasso all: full vs HAR, DM p
\providecommand{\rhchLassoBarAddRmPct}{}\renewcommand{\rhchLassoBarAddRmPct}{$-$2.1}  % lasso all add_rm vs ols|har: %
\providecommand{\rhchLassoBarAddRmP}{}\renewcommand{\rhchLassoBarAddRmP}{0.008}  % lasso all add_rm: DM p
\providecommand{\rhchLassoBarLoRmPct}{}\renewcommand{\rhchLassoBarLoRmPct}{1.6}  % lasso all lo_rm vs lasso|full: %
\providecommand{\rhchLassoBarLoRmP}{}\renewcommand{\rhchLassoBarLoRmP}{0.010}  % lasso all lo_rm: DM p
\providecommand{\rhchLassoBarAddLiqPct}{}\renewcommand{\rhchLassoBarAddLiqPct}{$-$0.3}  % lasso all add_liq vs ols|har: %
\providecommand{\rhchLassoBarAddLiqP}{}\renewcommand{\rhchLassoBarAddLiqP}{0.583}  % lasso all add_liq: DM p
\providecommand{\rhchLassoBarLoLiqPct}{}\renewcommand{\rhchLassoBarLoLiqPct}{0.1}  % lasso all lo_liq vs lasso|full: %
\providecommand{\rhchLassoBarLoLiqP}{}\renewcommand{\rhchLassoBarLoLiqP}{0.783}  % lasso all lo_liq: DM p
\providecommand{\rhchLassoBarAddEwPct}{}\renewcommand{\rhchLassoBarAddEwPct}{$-$0.8}  % lasso all add_ew vs ols|har: %
\providecommand{\rhchLassoBarAddEwP}{}\renewcommand{\rhchLassoBarAddEwP}{0.045}  % lasso all add_ew: DM p
\providecommand{\rhchLassoBarLoEwPct}{}\renewcommand{\rhchLassoBarLoEwPct}{$-$0.1}  % lasso all lo_ew vs lasso|full: %
\providecommand{\rhchLassoBarLoEwP}{}\renewcommand{\rhchLassoBarLoEwP}{0.583}  % lasso all lo_ew: DM p
\providecommand{\rhchLassoBarAddVwPct}{}\renewcommand{\rhchLassoBarAddVwPct}{$-$0.6}  % lasso all add_vw vs ols|har: %
\providecommand{\rhchLassoBarAddVwP}{}\renewcommand{\rhchLassoBarAddVwP}{0.309}  % lasso all add_vw: DM p
\providecommand{\rhchLassoBarLoVwPct}{}\renewcommand{\rhchLassoBarLoVwPct}{$-$0.2}  % lasso all lo_vw vs lasso|full: %
\providecommand{\rhchLassoBarLoVwP}{}\renewcommand{\rhchLassoBarLoVwP}{0.428}  % lasso all lo_vw: DM p
\providecommand{\rhchLassoBarAddSentPct}{}\renewcommand{\rhchLassoBarAddSentPct}{0.1}  % lasso all add_sent vs ols|har: %
\providecommand{\rhchLassoBarAddSentP}{}\renewcommand{\rhchLassoBarAddSentP}{0.550}  % lasso all add_sent: DM p
\providecommand{\rhchLassoBarLoSentPct}{}\renewcommand{\rhchLassoBarLoSentPct}{0.1}  % lasso all lo_sent vs lasso|full: %
\providecommand{\rhchLassoBarLoSentP}{}\renewcommand{\rhchLassoBarLoSentP}{0.693}  % lasso all lo_sent: DM p
\providecommand{\rhchLassoBarAddIvPct}{}\renewcommand{\rhchLassoBarAddIvPct}{$-$0.3}  % lasso all add_iv vs ols|har: %
\providecommand{\rhchLassoBarAddIvP}{}\renewcommand{\rhchLassoBarAddIvP}{0.430}  % lasso all add_iv: DM p
\providecommand{\rhchLassoBarLoIvPct}{}\renewcommand{\rhchLassoBarLoIvPct}{0.5}  % lasso all lo_iv vs lasso|full: %
\providecommand{\rhchLassoBarLoIvP}{}\renewcommand{\rhchLassoBarLoIvP}{0.313}  % lasso all lo_iv: DM p
\providecommand{\rhchLassoBarAddVdPct}{}\renewcommand{\rhchLassoBarAddVdPct}{0.6}  % lasso all add_vd vs ols|har: %
\providecommand{\rhchLassoBarAddVdP}{}\renewcommand{\rhchLassoBarAddVdP}{0.002}  % lasso all add_vd: DM p
\providecommand{\rhchLassoBarLoVdPct}{}\renewcommand{\rhchLassoBarLoVdPct}{0.0}  % lasso all lo_vd vs lasso|full: %
\providecommand{\rhchLassoBarLoVdP}{}\renewcommand{\rhchLassoBarLoVdP}{0.807}  % lasso all lo_vd: DM p
\providecommand{\rhchLassoBarAddFomcPct}{}\renewcommand{\rhchLassoBarAddFomcPct}{$-$5.8}  % lasso all add_fomc vs ols|har: %
\providecommand{\rhchLassoBarAddFomcP}{}\renewcommand{\rhchLassoBarAddFomcP}{\ensuremath{<}0.001}  % lasso all add_fomc: DM p
\providecommand{\rhchLassoBarLoFomcPct}{}\renewcommand{\rhchLassoBarLoFomcPct}{6.8}  % lasso all lo_fomc vs lasso|full: %
\providecommand{\rhchLassoBarLoFomcP}{}\renewcommand{\rhchLassoBarLoFomcP}{\ensuremath{<}0.001}  % lasso all lo_fomc: DM p
\providecommand{\rhchLassoEventJointPct}{}\renewcommand{\rhchLassoEventJointPct}{$-$47.2}  % lasso event: full vs HAR, % QLIKE change
\providecommand{\rhchLassoEventJointDm}{}\renewcommand{\rhchLassoEventJointDm}{$-$7.56}  % lasso event: full vs HAR, DM
\providecommand{\rhchLassoEventJointP}{}\renewcommand{\rhchLassoEventJointP}{\ensuremath{<}0.001}  % lasso event: full vs HAR, DM p
\providecommand{\rhchLassoEventAddRmPct}{}\renewcommand{\rhchLassoEventAddRmPct}{4.7}  % lasso event add_rm vs ols|har: %
\providecommand{\rhchLassoEventAddRmP}{}\renewcommand{\rhchLassoEventAddRmP}{0.024}  % lasso event add_rm: DM p
\providecommand{\rhchLassoEventLoRmPct}{}\renewcommand{\rhchLassoEventLoRmPct}{$-$0.8}  % lasso event lo_rm vs lasso|full: %
\providecommand{\rhchLassoEventLoRmP}{}\renewcommand{\rhchLassoEventLoRmP}{0.265}  % lasso event lo_rm: DM p
\providecommand{\rhchLassoEventAddLiqPct}{}\renewcommand{\rhchLassoEventAddLiqPct}{6.1}  % lasso event add_liq vs ols|har: %
\providecommand{\rhchLassoEventAddLiqP}{}\renewcommand{\rhchLassoEventAddLiqP}{0.007}  % lasso event add_liq: DM p
\providecommand{\rhchLassoEventLoLiqPct}{}\renewcommand{\rhchLassoEventLoLiqPct}{0.1}  % lasso event lo_liq vs lasso|full: %
\providecommand{\rhchLassoEventLoLiqP}{}\renewcommand{\rhchLassoEventLoLiqP}{0.838}  % lasso event lo_liq: DM p
\providecommand{\rhchLassoEventAddEwPct}{}\renewcommand{\rhchLassoEventAddEwPct}{$-$1.7}  % lasso event add_ew vs ols|har: %
\providecommand{\rhchLassoEventAddEwP}{}\renewcommand{\rhchLassoEventAddEwP}{0.235}  % lasso event add_ew: DM p
\providecommand{\rhchLassoEventLoEwPct}{}\renewcommand{\rhchLassoEventLoEwPct}{0.1}  % lasso event lo_ew vs lasso|full: %
\providecommand{\rhchLassoEventLoEwP}{}\renewcommand{\rhchLassoEventLoEwP}{0.717}  % lasso event lo_ew: DM p
\providecommand{\rhchLassoEventAddVwPct}{}\renewcommand{\rhchLassoEventAddVwPct}{$-$2.4}  % lasso event add_vw vs ols|har: %
\providecommand{\rhchLassoEventAddVwP}{}\renewcommand{\rhchLassoEventAddVwP}{0.171}  % lasso event add_vw: DM p
\providecommand{\rhchLassoEventLoVwPct}{}\renewcommand{\rhchLassoEventLoVwPct}{$-$0.3}  % lasso event lo_vw vs lasso|full: %
\providecommand{\rhchLassoEventLoVwP}{}\renewcommand{\rhchLassoEventLoVwP}{0.438}  % lasso event lo_vw: DM p
\providecommand{\rhchLassoEventAddSentPct}{}\renewcommand{\rhchLassoEventAddSentPct}{$-$1.3}  % lasso event add_sent vs ols|har: %
\providecommand{\rhchLassoEventAddSentP}{}\renewcommand{\rhchLassoEventAddSentP}{0.464}  % lasso event add_sent: DM p
\providecommand{\rhchLassoEventLoSentPct}{}\renewcommand{\rhchLassoEventLoSentPct}{2.0}  % lasso event lo_sent vs lasso|full: %
\providecommand{\rhchLassoEventLoSentP}{}\renewcommand{\rhchLassoEventLoSentP}{0.209}  % lasso event lo_sent: DM p
\providecommand{\rhchLassoEventAddIvPct}{}\renewcommand{\rhchLassoEventAddIvPct}{2.9}  % lasso event add_iv vs ols|har: %
\providecommand{\rhchLassoEventAddIvP}{}\renewcommand{\rhchLassoEventAddIvP}{0.004}  % lasso event add_iv: DM p
\providecommand{\rhchLassoEventLoIvPct}{}\renewcommand{\rhchLassoEventLoIvPct}{$-$0.1}  % lasso event lo_iv vs lasso|full: %
\providecommand{\rhchLassoEventLoIvP}{}\renewcommand{\rhchLassoEventLoIvP}{0.868}  % lasso event lo_iv: DM p
\providecommand{\rhchLassoEventAddVdPct}{}\renewcommand{\rhchLassoEventAddVdPct}{1.9}  % lasso event add_vd vs ols|har: %
\providecommand{\rhchLassoEventAddVdP}{}\renewcommand{\rhchLassoEventAddVdP}{0.032}  % lasso event add_vd: DM p
\providecommand{\rhchLassoEventLoVdPct}{}\renewcommand{\rhchLassoEventLoVdPct}{0.3}  % lasso event lo_vd vs lasso|full: %
\providecommand{\rhchLassoEventLoVdP}{}\renewcommand{\rhchLassoEventLoVdP}{0.012}  % lasso event lo_vd: DM p
\providecommand{\rhchLassoEventAddFomcPct}{}\renewcommand{\rhchLassoEventAddFomcPct}{$-$43.9}  % lasso event add_fomc vs ols|har: %
\providecommand{\rhchLassoEventAddFomcP}{}\renewcommand{\rhchLassoEventAddFomcP}{\ensuremath{<}0.001}  % lasso event add_fomc: DM p
\providecommand{\rhchLassoEventLoFomcPct}{}\renewcommand{\rhchLassoEventLoFomcPct}{87.4}  % lasso event lo_fomc vs lasso|full: %
\providecommand{\rhchLassoEventLoFomcP}{}\renewcommand{\rhchLassoEventLoFomcP}{\ensuremath{<}0.001}  % lasso event lo_fomc: DM p
\providecommand{\rhchLassoEventReading}{}\renewcommand{\rhchLassoEventReading}{concentrated}  % lasso event: item 3 reading
\providecommand{\rhchLassoEventTopAddShare}{}\renewcommand{\rhchLassoEventTopAddShare}{93}  % lasso event: largest add-alone gain, % of the joint gain (fomc)
\providecommand{\rhchLassoEventSumAddRatio}{}\renewcommand{\rhchLassoEventSumAddRatio}{0.72}  % lasso event: add-alone gains summed over groups / joint gain
\providecommand{\rhchLassoEventNLoSig}{}\renewcommand{\rhchLassoEventNLoSig}{2}  % lasso event: significant leave-one-outs
\providecommand{\rhchLassoEventNLoRun}{}\renewcommand{\rhchLassoEventNLoRun}{8}  % lasso event: leave-one-outs available (of 8)
\providecommand{\rhchLassoEventConcTest}{}\renewcommand{\rhchLassoEventConcTest}{yes}  % lasso event: the CONC definition on its own (pending: undetermined)
\providecommand{\rhchLassoEventRedTest}{}\renewcommand{\rhchLassoEventRedTest}{no}  % lasso event: the RED definition on its own (pending: undetermined)
\providecommand{\rhchLassoEventReadingRedFirst}{}\renewcommand{\rhchLassoEventReadingRedFirst}{concentrated}  % lasso event: reading with REDUNDANT tested before CONCENTRATED (the RULE tests CONCENTRATED first)
\providecommand{\rhchLassoOrdJointPct}{}\renewcommand{\rhchLassoOrdJointPct}{$-$2.4}  % lasso ord: full vs HAR, % QLIKE change
\providecommand{\rhchLassoOrdJointDm}{}\renewcommand{\rhchLassoOrdJointDm}{$-$2.08}  % lasso ord: full vs HAR, DM
\providecommand{\rhchLassoOrdJointP}{}\renewcommand{\rhchLassoOrdJointP}{0.037}  % lasso ord: full vs HAR, DM p
\providecommand{\rhchLassoOrdAddRmPct}{}\renewcommand{\rhchLassoOrdAddRmPct}{$-$3.3}  % lasso ord add_rm vs ols|har: %
\providecommand{\rhchLassoOrdAddRmP}{}\renewcommand{\rhchLassoOrdAddRmP}{\ensuremath{<}0.001}  % lasso ord add_rm: DM p
\providecommand{\rhchLassoOrdLoRmPct}{}\renewcommand{\rhchLassoOrdLoRmPct}{1.8}  % lasso ord lo_rm vs lasso|full: %
\providecommand{\rhchLassoOrdLoRmP}{}\renewcommand{\rhchLassoOrdLoRmP}{0.007}  % lasso ord lo_rm: DM p
\providecommand{\rhchLassoOrdAddLiqPct}{}\renewcommand{\rhchLassoOrdAddLiqPct}{$-$1.3}  % lasso ord add_liq vs ols|har: %
\providecommand{\rhchLassoOrdAddLiqP}{}\renewcommand{\rhchLassoOrdAddLiqP}{0.003}  % lasso ord add_liq: DM p
\providecommand{\rhchLassoOrdLoLiqPct}{}\renewcommand{\rhchLassoOrdLoLiqPct}{0.1}  % lasso ord lo_liq vs lasso|full: %
\providecommand{\rhchLassoOrdLoLiqP}{}\renewcommand{\rhchLassoOrdLoLiqP}{0.799}  % lasso ord lo_liq: DM p
\providecommand{\rhchLassoOrdAddEwPct}{}\renewcommand{\rhchLassoOrdAddEwPct}{$-$0.7}  % lasso ord add_ew vs ols|har: %
\providecommand{\rhchLassoOrdAddEwP}{}\renewcommand{\rhchLassoOrdAddEwP}{0.077}  % lasso ord add_ew: DM p
\providecommand{\rhchLassoOrdLoEwPct}{}\renewcommand{\rhchLassoOrdLoEwPct}{$-$0.1}  % lasso ord lo_ew vs lasso|full: %
\providecommand{\rhchLassoOrdLoEwP}{}\renewcommand{\rhchLassoOrdLoEwP}{0.530}  % lasso ord lo_ew: DM p
\providecommand{\rhchLassoOrdAddVwPct}{}\renewcommand{\rhchLassoOrdAddVwPct}{$-$0.3}  % lasso ord add_vw vs ols|har: %
\providecommand{\rhchLassoOrdAddVwP}{}\renewcommand{\rhchLassoOrdAddVwP}{0.630}  % lasso ord add_vw: DM p
\providecommand{\rhchLassoOrdLoVwPct}{}\renewcommand{\rhchLassoOrdLoVwPct}{$-$0.2}  % lasso ord lo_vw vs lasso|full: %
\providecommand{\rhchLassoOrdLoVwP}{}\renewcommand{\rhchLassoOrdLoVwP}{0.501}  % lasso ord lo_vw: DM p
\providecommand{\rhchLassoOrdAddSentPct}{}\renewcommand{\rhchLassoOrdAddSentPct}{0.4}  % lasso ord add_sent vs ols|har: %
\providecommand{\rhchLassoOrdAddSentP}{}\renewcommand{\rhchLassoOrdAddSentP}{0.098}  % lasso ord add_sent: DM p
\providecommand{\rhchLassoOrdLoSentPct}{}\renewcommand{\rhchLassoOrdLoSentPct}{$-$0.1}  % lasso ord lo_sent vs lasso|full: %
\providecommand{\rhchLassoOrdLoSentP}{}\renewcommand{\rhchLassoOrdLoSentP}{0.502}  % lasso ord lo_sent: DM p
\providecommand{\rhchLassoOrdAddIvPct}{}\renewcommand{\rhchLassoOrdAddIvPct}{$-$0.9}  % lasso ord add_iv vs ols|har: %
\providecommand{\rhchLassoOrdAddIvP}{}\renewcommand{\rhchLassoOrdAddIvP}{0.042}  % lasso ord add_iv: DM p
\providecommand{\rhchLassoOrdLoIvPct}{}\renewcommand{\rhchLassoOrdLoIvPct}{0.5}  % lasso ord lo_iv vs lasso|full: %
\providecommand{\rhchLassoOrdLoIvP}{}\renewcommand{\rhchLassoOrdLoIvP}{0.292}  % lasso ord lo_iv: DM p
\providecommand{\rhchLassoOrdAddVdPct}{}\renewcommand{\rhchLassoOrdAddVdPct}{0.3}  % lasso ord add_vd vs ols|har: %
\providecommand{\rhchLassoOrdAddVdP}{}\renewcommand{\rhchLassoOrdAddVdP}{0.056}  % lasso ord add_vd: DM p
\providecommand{\rhchLassoOrdLoVdPct}{}\renewcommand{\rhchLassoOrdLoVdPct}{$-$0.1}  % lasso ord lo_vd vs lasso|full: %
\providecommand{\rhchLassoOrdLoVdP}{}\renewcommand{\rhchLassoOrdLoVdP}{0.644}  % lasso ord lo_vd: DM p
\providecommand{\rhchLassoOrdAddFomcPct}{}\renewcommand{\rhchLassoOrdAddFomcPct}{0.7}  % lasso ord add_fomc vs ols|har: %
\providecommand{\rhchLassoOrdAddFomcP}{}\renewcommand{\rhchLassoOrdAddFomcP}{\ensuremath{<}0.001}  % lasso ord add_fomc: DM p
\providecommand{\rhchLassoOrdLoFomcPct}{}\renewcommand{\rhchLassoOrdLoFomcPct}{$-$0.7}  % lasso ord lo_fomc vs lasso|full: %
\providecommand{\rhchLassoOrdLoFomcP}{}\renewcommand{\rhchLassoOrdLoFomcP}{\ensuremath{<}0.001}  % lasso ord lo_fomc: DM p
\providecommand{\rhchLassoOrdReading}{}\renewcommand{\rhchLassoOrdReading}{concentrated}  % lasso ord: item 3 reading
\providecommand{\rhchLassoOrdTopAddShare}{}\renewcommand{\rhchLassoOrdTopAddShare}{136}  % lasso ord: largest add-alone gain, % of the joint gain (rm)
\providecommand{\rhchLassoOrdSumAddRatio}{}\renewcommand{\rhchLassoOrdSumAddRatio}{2.07}  % lasso ord: add-alone gains summed over groups / joint gain
\providecommand{\rhchLassoOrdNLoSig}{}\renewcommand{\rhchLassoOrdNLoSig}{1}  % lasso ord: significant leave-one-outs
\providecommand{\rhchLassoOrdNLoRun}{}\renewcommand{\rhchLassoOrdNLoRun}{8}  % lasso ord: leave-one-outs available (of 8)
\providecommand{\rhchLassoOrdConcTest}{}\renewcommand{\rhchLassoOrdConcTest}{yes}  % lasso ord: the CONC definition on its own (pending: undetermined)
\providecommand{\rhchLassoOrdRedTest}{}\renewcommand{\rhchLassoOrdRedTest}{yes}  % lasso ord: the RED definition on its own (pending: undetermined)
\providecommand{\rhchLassoOrdReadingRedFirst}{}\renewcommand{\rhchLassoOrdReadingRedFirst}{redundant}  % lasso ord: reading with REDUNDANT tested before CONCENTRATED (the RULE tests CONCENTRATED first)
\providecommand{\rhchEnetBarJointPct}{}\renewcommand{\rhchEnetBarJointPct}{$-$6.7}  % enet all: full vs HAR, % QLIKE change
\providecommand{\rhchEnetBarJointDm}{}\renewcommand{\rhchEnetBarJointDm}{$-$7.36}  % enet all: full vs HAR, DM
\providecommand{\rhchEnetBarJointP}{}\renewcommand{\rhchEnetBarJointP}{\ensuremath{<}0.001}  % enet all: full vs HAR, DM p
\providecommand{\rhchEnetBarAddRmPct}{}\renewcommand{\rhchEnetBarAddRmPct}{$-$2.7}  % enet all add_rm vs ols|har: %
\providecommand{\rhchEnetBarAddRmP}{}\renewcommand{\rhchEnetBarAddRmP}{\ensuremath{<}0.001}  % enet all add_rm: DM p
\providecommand{\rhchEnetBarLoRmPct}{}\renewcommand{\rhchEnetBarLoRmPct}{2.0}  % enet all lo_rm vs enet|full: %
\providecommand{\rhchEnetBarLoRmP}{}\renewcommand{\rhchEnetBarLoRmP}{\ensuremath{<}0.001}  % enet all lo_rm: DM p
\providecommand{\rhchEnetBarAddLiqPct}{}\renewcommand{\rhchEnetBarAddLiqPct}{$-$0.6}  % enet all add_liq vs ols|har: %
\providecommand{\rhchEnetBarAddLiqP}{}\renewcommand{\rhchEnetBarAddLiqP}{0.124}  % enet all add_liq: DM p
\providecommand{\rhchEnetBarLoLiqPct}{}\renewcommand{\rhchEnetBarLoLiqPct}{0.1}  % enet all lo_liq vs enet|full: %
\providecommand{\rhchEnetBarLoLiqP}{}\renewcommand{\rhchEnetBarLoLiqP}{0.638}  % enet all lo_liq: DM p
\providecommand{\rhchEnetBarAddEwPct}{}\renewcommand{\rhchEnetBarAddEwPct}{$-$1.1}  % enet all add_ew vs ols|har: %
\providecommand{\rhchEnetBarAddEwP}{}\renewcommand{\rhchEnetBarAddEwP}{0.005}  % enet all add_ew: DM p
\providecommand{\rhchEnetBarLoEwPct}{}\renewcommand{\rhchEnetBarLoEwPct}{$-$0.1}  % enet all lo_ew vs enet|full: %
\providecommand{\rhchEnetBarLoEwP}{}\renewcommand{\rhchEnetBarLoEwP}{0.327}  % enet all lo_ew: DM p
\providecommand{\rhchEnetBarAddVwPct}{}\renewcommand{\rhchEnetBarAddVwPct}{$-$1.1}  % enet all add_vw vs ols|har: %
\providecommand{\rhchEnetBarAddVwP}{}\renewcommand{\rhchEnetBarAddVwP}{0.042}  % enet all add_vw: DM p
\providecommand{\rhchEnetBarLoVwPct}{}\renewcommand{\rhchEnetBarLoVwPct}{$-$0.1}  % enet all lo_vw vs enet|full: %
\providecommand{\rhchEnetBarLoVwP}{}\renewcommand{\rhchEnetBarLoVwP}{0.492}  % enet all lo_vw: DM p
\providecommand{\rhchEnetBarAddSentPct}{}\renewcommand{\rhchEnetBarAddSentPct}{0.0}  % enet all add_sent vs ols|har: %
\providecommand{\rhchEnetBarAddSentP}{}\renewcommand{\rhchEnetBarAddSentP}{0.971}  % enet all add_sent: DM p
\providecommand{\rhchEnetBarLoSentPct}{}\renewcommand{\rhchEnetBarLoSentPct}{0.2}  % enet all lo_sent vs enet|full: %
\providecommand{\rhchEnetBarLoSentP}{}\renewcommand{\rhchEnetBarLoSentP}{0.137}  % enet all lo_sent: DM p
\providecommand{\rhchEnetBarAddIvPct}{}\renewcommand{\rhchEnetBarAddIvPct}{0.1}  % enet all add_iv vs ols|har: %
\providecommand{\rhchEnetBarAddIvP}{}\renewcommand{\rhchEnetBarAddIvP}{0.817}  % enet all add_iv: DM p
\providecommand{\rhchEnetBarLoIvPct}{}\renewcommand{\rhchEnetBarLoIvPct}{0.0}  % enet all lo_iv vs enet|full: %
\providecommand{\rhchEnetBarLoIvP}{}\renewcommand{\rhchEnetBarLoIvP}{0.950}  % enet all lo_iv: DM p
\providecommand{\rhchEnetBarAddVdPct}{}\renewcommand{\rhchEnetBarAddVdPct}{0.4}  % enet all add_vd vs ols|har: %
\providecommand{\rhchEnetBarAddVdP}{}\renewcommand{\rhchEnetBarAddVdP}{0.003}  % enet all add_vd: DM p
\providecommand{\rhchEnetBarLoVdPct}{}\renewcommand{\rhchEnetBarLoVdPct}{$-$0.1}  % enet all lo_vd vs enet|full: %
\providecommand{\rhchEnetBarLoVdP}{}\renewcommand{\rhchEnetBarLoVdP}{0.264}  % enet all lo_vd: DM p
\providecommand{\rhchEnetBarAddFomcPct}{}\renewcommand{\rhchEnetBarAddFomcPct}{$-$3.1}  % enet all add_fomc vs ols|har: %
\providecommand{\rhchEnetBarAddFomcP}{}\renewcommand{\rhchEnetBarAddFomcP}{\ensuremath{<}0.001}  % enet all add_fomc: DM p
\providecommand{\rhchEnetBarLoFomcPct}{}\renewcommand{\rhchEnetBarLoFomcPct}{3.1}  % enet all lo_fomc vs enet|full: %
\providecommand{\rhchEnetBarLoFomcP}{}\renewcommand{\rhchEnetBarLoFomcP}{\ensuremath{<}0.001}  % enet all lo_fomc: DM p
\providecommand{\rhchEnetEventJointPct}{}\renewcommand{\rhchEnetEventJointPct}{$-$22.6}  % enet event: full vs HAR, % QLIKE change
\providecommand{\rhchEnetEventJointDm}{}\renewcommand{\rhchEnetEventJointDm}{$-$6.79}  % enet event: full vs HAR, DM
\providecommand{\rhchEnetEventJointP}{}\renewcommand{\rhchEnetEventJointP}{\ensuremath{<}0.001}  % enet event: full vs HAR, DM p
\providecommand{\rhchEnetEventAddRmPct}{}\renewcommand{\rhchEnetEventAddRmPct}{4.2}  % enet event add_rm vs ols|har: %
\providecommand{\rhchEnetEventAddRmP}{}\renewcommand{\rhchEnetEventAddRmP}{0.031}  % enet event add_rm: DM p
\providecommand{\rhchEnetEventLoRmPct}{}\renewcommand{\rhchEnetEventLoRmPct}{$-$1.2}  % enet event lo_rm vs enet|full: %
\providecommand{\rhchEnetEventLoRmP}{}\renewcommand{\rhchEnetEventLoRmP}{0.167}  % enet event lo_rm: DM p
\providecommand{\rhchEnetEventAddLiqPct}{}\renewcommand{\rhchEnetEventAddLiqPct}{5.3}  % enet event add_liq vs ols|har: %
\providecommand{\rhchEnetEventAddLiqP}{}\renewcommand{\rhchEnetEventAddLiqP}{0.008}  % enet event add_liq: DM p
\providecommand{\rhchEnetEventLoLiqPct}{}\renewcommand{\rhchEnetEventLoLiqPct}{$-$0.7}  % enet event lo_liq vs enet|full: %
\providecommand{\rhchEnetEventLoLiqP}{}\renewcommand{\rhchEnetEventLoLiqP}{0.104}  % enet event lo_liq: DM p
\providecommand{\rhchEnetEventAddEwPct}{}\renewcommand{\rhchEnetEventAddEwPct}{$-$2.7}  % enet event add_ew vs ols|har: %
\providecommand{\rhchEnetEventAddEwP}{}\renewcommand{\rhchEnetEventAddEwP}{0.060}  % enet event add_ew: DM p
\providecommand{\rhchEnetEventLoEwPct}{}\renewcommand{\rhchEnetEventLoEwPct}{0.0}  % enet event lo_ew vs enet|full: %
\providecommand{\rhchEnetEventLoEwP}{}\renewcommand{\rhchEnetEventLoEwP}{0.915}  % enet event lo_ew: DM p
\providecommand{\rhchEnetEventAddVwPct}{}\renewcommand{\rhchEnetEventAddVwPct}{$-$3.4}  % enet event add_vw vs ols|har: %
\providecommand{\rhchEnetEventAddVwP}{}\renewcommand{\rhchEnetEventAddVwP}{0.055}  % enet event add_vw: DM p
\providecommand{\rhchEnetEventLoVwPct}{}\renewcommand{\rhchEnetEventLoVwPct}{0.0}  % enet event lo_vw vs enet|full: %
\providecommand{\rhchEnetEventLoVwP}{}\renewcommand{\rhchEnetEventLoVwP}{0.947}  % enet event lo_vw: DM p
\providecommand{\rhchEnetEventAddSentPct}{}\renewcommand{\rhchEnetEventAddSentPct}{$-$0.9}  % enet event add_sent vs ols|har: %
\providecommand{\rhchEnetEventAddSentP}{}\renewcommand{\rhchEnetEventAddSentP}{0.519}  % enet event add_sent: DM p
\providecommand{\rhchEnetEventLoSentPct}{}\renewcommand{\rhchEnetEventLoSentPct}{1.6}  % enet event lo_sent vs enet|full: %
\providecommand{\rhchEnetEventLoSentP}{}\renewcommand{\rhchEnetEventLoSentP}{0.201}  % enet event lo_sent: DM p
\providecommand{\rhchEnetEventAddIvPct}{}\renewcommand{\rhchEnetEventAddIvPct}{2.5}  % enet event add_iv vs ols|har: %
\providecommand{\rhchEnetEventAddIvP}{}\renewcommand{\rhchEnetEventAddIvP}{0.010}  % enet event add_iv: DM p
\providecommand{\rhchEnetEventLoIvPct}{}\renewcommand{\rhchEnetEventLoIvPct}{$-$0.1}  % enet event lo_iv vs enet|full: %
\providecommand{\rhchEnetEventLoIvP}{}\renewcommand{\rhchEnetEventLoIvP}{0.865}  % enet event lo_iv: DM p
\providecommand{\rhchEnetEventAddVdPct}{}\renewcommand{\rhchEnetEventAddVdPct}{1.6}  % enet event add_vd vs ols|har: %
\providecommand{\rhchEnetEventAddVdP}{}\renewcommand{\rhchEnetEventAddVdP}{0.007}  % enet event add_vd: DM p
\providecommand{\rhchEnetEventLoVdPct}{}\renewcommand{\rhchEnetEventLoVdPct}{$-$0.1}  % enet event lo_vd vs enet|full: %
\providecommand{\rhchEnetEventLoVdP}{}\renewcommand{\rhchEnetEventLoVdP}{0.330}  % enet event lo_vd: DM p
\providecommand{\rhchEnetEventAddFomcPct}{}\renewcommand{\rhchEnetEventAddFomcPct}{$-$23.2}  % enet event add_fomc vs ols|har: %
\providecommand{\rhchEnetEventAddFomcP}{}\renewcommand{\rhchEnetEventAddFomcP}{\ensuremath{<}0.001}  % enet event add_fomc: DM p
\providecommand{\rhchEnetEventLoFomcPct}{}\renewcommand{\rhchEnetEventLoFomcPct}{28.2}  % enet event lo_fomc vs enet|full: %
\providecommand{\rhchEnetEventLoFomcP}{}\renewcommand{\rhchEnetEventLoFomcP}{\ensuremath{<}0.001}  % enet event lo_fomc: DM p
\providecommand{\rhchEnetEventReading}{}\renewcommand{\rhchEnetEventReading}{concentrated}  % enet event: item 3 reading
\providecommand{\rhchEnetEventTopAddShare}{}\renewcommand{\rhchEnetEventTopAddShare}{103}  % enet event: largest add-alone gain, % of the joint gain (fomc)
\providecommand{\rhchEnetEventSumAddRatio}{}\renewcommand{\rhchEnetEventSumAddRatio}{0.74}  % enet event: add-alone gains summed over groups / joint gain
\providecommand{\rhchEnetEventNLoSig}{}\renewcommand{\rhchEnetEventNLoSig}{1}  % enet event: significant leave-one-outs
\providecommand{\rhchEnetEventNLoRun}{}\renewcommand{\rhchEnetEventNLoRun}{8}  % enet event: leave-one-outs available (of 8)
\providecommand{\rhchEnetEventConcTest}{}\renewcommand{\rhchEnetEventConcTest}{yes}  % enet event: the CONC definition on its own (pending: undetermined)
\providecommand{\rhchEnetEventRedTest}{}\renewcommand{\rhchEnetEventRedTest}{no}  % enet event: the RED definition on its own (pending: undetermined)
\providecommand{\rhchEnetEventReadingRedFirst}{}\renewcommand{\rhchEnetEventReadingRedFirst}{concentrated}  % enet event: reading with REDUNDANT tested before CONCENTRATED (the RULE tests CONCENTRATED first)
\providecommand{\rhchEnetOrdJointPct}{}\renewcommand{\rhchEnetOrdJointPct}{$-$4.0}  % enet ord: full vs HAR, % QLIKE change
\providecommand{\rhchEnetOrdJointDm}{}\renewcommand{\rhchEnetOrdJointDm}{$-$4.73}  % enet ord: full vs HAR, DM
\providecommand{\rhchEnetOrdJointP}{}\renewcommand{\rhchEnetOrdJointP}{\ensuremath{<}0.001}  % enet ord: full vs HAR, DM p
\providecommand{\rhchEnetOrdAddRmPct}{}\renewcommand{\rhchEnetOrdAddRmPct}{$-$3.9}  % enet ord add_rm vs ols|har: %
\providecommand{\rhchEnetOrdAddRmP}{}\renewcommand{\rhchEnetOrdAddRmP}{\ensuremath{<}0.001}  % enet ord add_rm: DM p
\providecommand{\rhchEnetOrdLoRmPct}{}\renewcommand{\rhchEnetOrdLoRmPct}{2.5}  % enet ord lo_rm vs enet|full: %
\providecommand{\rhchEnetOrdLoRmP}{}\renewcommand{\rhchEnetOrdLoRmP}{\ensuremath{<}0.001}  % enet ord lo_rm: DM p
\providecommand{\rhchEnetOrdAddLiqPct}{}\renewcommand{\rhchEnetOrdAddLiqPct}{$-$1.6}  % enet ord add_liq vs ols|har: %
\providecommand{\rhchEnetOrdAddLiqP}{}\renewcommand{\rhchEnetOrdAddLiqP}{\ensuremath{<}0.001}  % enet ord add_liq: DM p
\providecommand{\rhchEnetOrdLoLiqPct}{}\renewcommand{\rhchEnetOrdLoLiqPct}{0.2}  % enet ord lo_liq vs enet|full: %
\providecommand{\rhchEnetOrdLoLiqP}{}\renewcommand{\rhchEnetOrdLoLiqP}{0.282}  % enet ord lo_liq: DM p
\providecommand{\rhchEnetOrdAddEwPct}{}\renewcommand{\rhchEnetOrdAddEwPct}{$-$0.9}  % enet ord add_ew vs ols|har: %
\providecommand{\rhchEnetOrdAddEwP}{}\renewcommand{\rhchEnetOrdAddEwP}{0.018}  % enet ord add_ew: DM p
\providecommand{\rhchEnetOrdLoEwPct}{}\renewcommand{\rhchEnetOrdLoEwPct}{$-$0.1}  % enet ord lo_ew vs enet|full: %
\providecommand{\rhchEnetOrdLoEwP}{}\renewcommand{\rhchEnetOrdLoEwP}{0.220}  % enet ord lo_ew: DM p
\providecommand{\rhchEnetOrdAddVwPct}{}\renewcommand{\rhchEnetOrdAddVwPct}{$-$0.7}  % enet ord add_vw vs ols|har: %
\providecommand{\rhchEnetOrdAddVwP}{}\renewcommand{\rhchEnetOrdAddVwP}{0.195}  % enet ord add_vw: DM p
\providecommand{\rhchEnetOrdLoVwPct}{}\renewcommand{\rhchEnetOrdLoVwPct}{$-$0.1}  % enet ord lo_vw vs enet|full: %
\providecommand{\rhchEnetOrdLoVwP}{}\renewcommand{\rhchEnetOrdLoVwP}{0.491}  % enet ord lo_vw: DM p
\providecommand{\rhchEnetOrdAddSentPct}{}\renewcommand{\rhchEnetOrdAddSentPct}{0.2}  % enet ord add_sent vs ols|har: %
\providecommand{\rhchEnetOrdAddSentP}{}\renewcommand{\rhchEnetOrdAddSentP}{0.407}  % enet ord add_sent: DM p
\providecommand{\rhchEnetOrdLoSentPct}{}\renewcommand{\rhchEnetOrdLoSentPct}{0.0}  % enet ord lo_sent vs enet|full: %
\providecommand{\rhchEnetOrdLoSentP}{}\renewcommand{\rhchEnetOrdLoSentP}{0.998}  % enet ord lo_sent: DM p
\providecommand{\rhchEnetOrdAddIvPct}{}\renewcommand{\rhchEnetOrdAddIvPct}{$-$0.3}  % enet ord add_iv vs ols|har: %
\providecommand{\rhchEnetOrdAddIvP}{}\renewcommand{\rhchEnetOrdAddIvP}{0.319}  % enet ord add_iv: DM p
\providecommand{\rhchEnetOrdLoIvPct}{}\renewcommand{\rhchEnetOrdLoIvPct}{0.0}  % enet ord lo_iv vs enet|full: %
\providecommand{\rhchEnetOrdLoIvP}{}\renewcommand{\rhchEnetOrdLoIvP}{0.912}  % enet ord lo_iv: DM p
\providecommand{\rhchEnetOrdAddVdPct}{}\renewcommand{\rhchEnetOrdAddVdPct}{0.2}  % enet ord add_vd vs ols|har: %
\providecommand{\rhchEnetOrdAddVdP}{}\renewcommand{\rhchEnetOrdAddVdP}{0.127}  % enet ord add_vd: DM p
\providecommand{\rhchEnetOrdLoVdPct}{}\renewcommand{\rhchEnetOrdLoVdPct}{$-$0.1}  % enet ord lo_vd vs enet|full: %
\providecommand{\rhchEnetOrdLoVdP}{}\renewcommand{\rhchEnetOrdLoVdP}{0.345}  % enet ord lo_vd: DM p
\providecommand{\rhchEnetOrdAddFomcPct}{}\renewcommand{\rhchEnetOrdAddFomcPct}{0.4}  % enet ord add_fomc vs ols|har: %
\providecommand{\rhchEnetOrdAddFomcP}{}\renewcommand{\rhchEnetOrdAddFomcP}{\ensuremath{<}0.001}  % enet ord add_fomc: DM p
\providecommand{\rhchEnetOrdLoFomcPct}{}\renewcommand{\rhchEnetOrdLoFomcPct}{$-$0.4}  % enet ord lo_fomc vs enet|full: %
\providecommand{\rhchEnetOrdLoFomcP}{}\renewcommand{\rhchEnetOrdLoFomcP}{\ensuremath{<}0.001}  % enet ord lo_fomc: DM p
\providecommand{\rhchEnetOrdReading}{}\renewcommand{\rhchEnetOrdReading}{concentrated}  % enet ord: item 3 reading
\providecommand{\rhchEnetOrdTopAddShare}{}\renewcommand{\rhchEnetOrdTopAddShare}{97}  % enet ord: largest add-alone gain, % of the joint gain (rm)
\providecommand{\rhchEnetOrdSumAddRatio}{}\renewcommand{\rhchEnetOrdSumAddRatio}{1.68}  % enet ord: add-alone gains summed over groups / joint gain
\providecommand{\rhchEnetOrdNLoSig}{}\renewcommand{\rhchEnetOrdNLoSig}{1}  % enet ord: significant leave-one-outs
\providecommand{\rhchEnetOrdNLoRun}{}\renewcommand{\rhchEnetOrdNLoRun}{8}  % enet ord: leave-one-outs available (of 8)
\providecommand{\rhchEnetOrdConcTest}{}\renewcommand{\rhchEnetOrdConcTest}{yes}  % enet ord: the CONC definition on its own (pending: undetermined)
\providecommand{\rhchEnetOrdRedTest}{}\renewcommand{\rhchEnetOrdRedTest}{yes}  % enet ord: the RED definition on its own (pending: undetermined)
\providecommand{\rhchEnetOrdReadingRedFirst}{}\renewcommand{\rhchEnetOrdReadingRedFirst}{redundant}  % enet ord: reading with REDUNDANT tested before CONCENTRATED (the RULE tests CONCENTRATED first)
\providecommand{\rhchLgbmBarJointPct}{}\renewcommand{\rhchLgbmBarJointPct}{$-$14.1}  % lgbm all: full vs HAR, % QLIKE change
\providecommand{\rhchLgbmBarJointDm}{}\renewcommand{\rhchLgbmBarJointDm}{$-$10.26}  % lgbm all: full vs HAR, DM
\providecommand{\rhchLgbmBarJointP}{}\renewcommand{\rhchLgbmBarJointP}{\ensuremath{<}0.001}  % lgbm all: full vs HAR, DM p
\providecommand{\rhchLgbmBarAddRmPct}{}\renewcommand{\rhchLgbmBarAddRmPct}{$-$3.0}  % lgbm all add_rm vs lgbm|har: %
\providecommand{\rhchLgbmBarAddRmP}{}\renewcommand{\rhchLgbmBarAddRmP}{\ensuremath{<}0.001}  % lgbm all add_rm: DM p
\providecommand{\rhchLgbmBarLoRmPct}{}\renewcommand{\rhchLgbmBarLoRmPct}{1.5}  % lgbm all lo_rm vs lgbm|full_pre: % [pre-patch]
\providecommand{\rhchLgbmBarLoRmP}{}\renewcommand{\rhchLgbmBarLoRmP}{\ensuremath{<}0.001}  % lgbm all lo_rm: DM p [pre-patch]
\providecommand{\rhchLgbmBarAddLiqPct}{}\renewcommand{\rhchLgbmBarAddLiqPct}{$-$2.9}  % lgbm all add_liq vs lgbm|har: %
\providecommand{\rhchLgbmBarAddLiqP}{}\renewcommand{\rhchLgbmBarAddLiqP}{\ensuremath{<}0.001}  % lgbm all add_liq: DM p
\providecommand{\rhchLgbmBarLoLiqPct}{}\renewcommand{\rhchLgbmBarLoLiqPct}{0.2}  % lgbm all lo_liq vs lgbm|full_pre: % [pre-patch]
\providecommand{\rhchLgbmBarLoLiqP}{}\renewcommand{\rhchLgbmBarLoLiqP}{0.365}  % lgbm all lo_liq: DM p [pre-patch]
\providecommand{\rhchLgbmBarAddEwPct}{}\renewcommand{\rhchLgbmBarAddEwPct}{$-$2.8}  % lgbm all add_ew vs lgbm|har: %
\providecommand{\rhchLgbmBarAddEwP}{}\renewcommand{\rhchLgbmBarAddEwP}{\ensuremath{<}0.001}  % lgbm all add_ew: DM p
\providecommand{\rhchLgbmBarLoEwPct}{}\renewcommand{\rhchLgbmBarLoEwPct}{0.0}  % lgbm all lo_ew vs lgbm|full_pre: % [pre-patch]
\providecommand{\rhchLgbmBarLoEwP}{}\renewcommand{\rhchLgbmBarLoEwP}{0.815}  % lgbm all lo_ew: DM p [pre-patch]
\providecommand{\rhchLgbmBarAddVwPct}{}\renewcommand{\rhchLgbmBarAddVwPct}{$-$3.1}  % lgbm all add_vw vs lgbm|har: %
\providecommand{\rhchLgbmBarAddVwP}{}\renewcommand{\rhchLgbmBarAddVwP}{\ensuremath{<}0.001}  % lgbm all add_vw: DM p
\providecommand{\rhchLgbmBarLoVwPct}{}\renewcommand{\rhchLgbmBarLoVwPct}{0.1}  % lgbm all lo_vw vs lgbm|full_pre: % [pre-patch]
\providecommand{\rhchLgbmBarLoVwP}{}\renewcommand{\rhchLgbmBarLoVwP}{0.489}  % lgbm all lo_vw: DM p [pre-patch]
\providecommand{\rhchLgbmBarAddSentPct}{}\renewcommand{\rhchLgbmBarAddSentPct}{$-$1.1}  % lgbm all add_sent vs lgbm|har: %
\providecommand{\rhchLgbmBarAddSentP}{}\renewcommand{\rhchLgbmBarAddSentP}{\ensuremath{<}0.001}  % lgbm all add_sent: DM p
\providecommand{\rhchLgbmBarLoSentPct}{}\renewcommand{\rhchLgbmBarLoSentPct}{0.2}  % lgbm all lo_sent vs lgbm|full_pre: % [pre-patch]
\providecommand{\rhchLgbmBarLoSentP}{}\renewcommand{\rhchLgbmBarLoSentP}{0.028}  % lgbm all lo_sent: DM p [pre-patch]
\providecommand{\rhchLgbmBarAddIvPct}{}\renewcommand{\rhchLgbmBarAddIvPct}{$-$1.3}  % lgbm all add_iv vs lgbm|har: %
\providecommand{\rhchLgbmBarAddIvP}{}\renewcommand{\rhchLgbmBarAddIvP}{\ensuremath{<}0.001}  % lgbm all add_iv: DM p
\providecommand{\rhchLgbmBarLoIvPct}{}\renewcommand{\rhchLgbmBarLoIvPct}{0.2}  % lgbm all lo_iv vs lgbm|full_pre: % [pre-patch]
\providecommand{\rhchLgbmBarLoIvP}{}\renewcommand{\rhchLgbmBarLoIvP}{0.084}  % lgbm all lo_iv: DM p [pre-patch]
\providecommand{\rhchLgbmBarAddVdPct}{}\renewcommand{\rhchLgbmBarAddVdPct}{$-$1.2}  % lgbm all add_vd vs lgbm|har: %
\providecommand{\rhchLgbmBarAddVdP}{}\renewcommand{\rhchLgbmBarAddVdP}{\ensuremath{<}0.001}  % lgbm all add_vd: DM p
\providecommand{\rhchLgbmBarLoVdPct}{}\renewcommand{\rhchLgbmBarLoVdPct}{0.0}  % lgbm all lo_vd vs lgbm|full_pre: % [pre-patch]
\providecommand{\rhchLgbmBarLoVdP}{}\renewcommand{\rhchLgbmBarLoVdP}{0.916}  % lgbm all lo_vd: DM p [pre-patch]
\providecommand{\rhchLgbmBarAddFomcPct}{}\renewcommand{\rhchLgbmBarAddFomcPct}{$-$9.2}  % lgbm all add_fomc vs lgbm|har: %
\providecommand{\rhchLgbmBarAddFomcP}{}\renewcommand{\rhchLgbmBarAddFomcP}{\ensuremath{<}0.001}  % lgbm all add_fomc: DM p
\providecommand{\rhchLgbmBarLoFomcPct}{}\renewcommand{\rhchLgbmBarLoFomcPct}{10.3}  % lgbm all lo_fomc vs lgbm|full: %
\providecommand{\rhchLgbmBarLoFomcP}{}\renewcommand{\rhchLgbmBarLoFomcP}{\ensuremath{<}0.001}  % lgbm all lo_fomc: DM p
\providecommand{\rhchLgbmEventJointPct}{}\renewcommand{\rhchLgbmEventJointPct}{$-$62.3}  % lgbm event: full vs HAR, % QLIKE change
\providecommand{\rhchLgbmEventJointDm}{}\renewcommand{\rhchLgbmEventJointDm}{$-$8.62}  % lgbm event: full vs HAR, DM
\providecommand{\rhchLgbmEventJointP}{}\renewcommand{\rhchLgbmEventJointP}{\ensuremath{<}0.001}  % lgbm event: full vs HAR, DM p
\providecommand{\rhchLgbmEventAddRmPct}{}\renewcommand{\rhchLgbmEventAddRmPct}{6.8}  % lgbm event add_rm vs lgbm|har: %
\providecommand{\rhchLgbmEventAddRmP}{}\renewcommand{\rhchLgbmEventAddRmP}{0.010}  % lgbm event add_rm: DM p
\providecommand{\rhchLgbmEventLoRmPct}{}\renewcommand{\rhchLgbmEventLoRmPct}{$-$0.3}  % lgbm event lo_rm vs lgbm|full_pre: % [pre-patch]
\providecommand{\rhchLgbmEventLoRmP}{}\renewcommand{\rhchLgbmEventLoRmP}{0.700}  % lgbm event lo_rm: DM p [pre-patch]
\providecommand{\rhchLgbmEventAddLiqPct}{}\renewcommand{\rhchLgbmEventAddLiqPct}{1.5}  % lgbm event add_liq vs lgbm|har: %
\providecommand{\rhchLgbmEventAddLiqP}{}\renewcommand{\rhchLgbmEventAddLiqP}{0.294}  % lgbm event add_liq: DM p
\providecommand{\rhchLgbmEventLoLiqPct}{}\renewcommand{\rhchLgbmEventLoLiqPct}{$-$0.5}  % lgbm event lo_liq vs lgbm|full_pre: % [pre-patch]
\providecommand{\rhchLgbmEventLoLiqP}{}\renewcommand{\rhchLgbmEventLoLiqP}{0.155}  % lgbm event lo_liq: DM p [pre-patch]
\providecommand{\rhchLgbmEventAddEwPct}{}\renewcommand{\rhchLgbmEventAddEwPct}{$-$3.5}  % lgbm event add_ew vs lgbm|har: %
\providecommand{\rhchLgbmEventAddEwP}{}\renewcommand{\rhchLgbmEventAddEwP}{0.045}  % lgbm event add_ew: DM p
\providecommand{\rhchLgbmEventLoEwPct}{}\renewcommand{\rhchLgbmEventLoEwPct}{0.1}  % lgbm event lo_ew vs lgbm|full_pre: % [pre-patch]
\providecommand{\rhchLgbmEventLoEwP}{}\renewcommand{\rhchLgbmEventLoEwP}{0.888}  % lgbm event lo_ew: DM p [pre-patch]
\providecommand{\rhchLgbmEventAddVwPct}{}\renewcommand{\rhchLgbmEventAddVwPct}{$-$3.5}  % lgbm event add_vw vs lgbm|har: %
\providecommand{\rhchLgbmEventAddVwP}{}\renewcommand{\rhchLgbmEventAddVwP}{0.044}  % lgbm event add_vw: DM p
\providecommand{\rhchLgbmEventLoVwPct}{}\renewcommand{\rhchLgbmEventLoVwPct}{$-$0.1}  % lgbm event lo_vw vs lgbm|full_pre: % [pre-patch]
\providecommand{\rhchLgbmEventLoVwP}{}\renewcommand{\rhchLgbmEventLoVwP}{0.804}  % lgbm event lo_vw: DM p [pre-patch]
\providecommand{\rhchLgbmEventAddSentPct}{}\renewcommand{\rhchLgbmEventAddSentPct}{0.5}  % lgbm event add_sent vs lgbm|har: %
\providecommand{\rhchLgbmEventAddSentP}{}\renewcommand{\rhchLgbmEventAddSentP}{0.574}  % lgbm event add_sent: DM p
\providecommand{\rhchLgbmEventLoSentPct}{}\renewcommand{\rhchLgbmEventLoSentPct}{0.2}  % lgbm event lo_sent vs lgbm|full_pre: % [pre-patch]
\providecommand{\rhchLgbmEventLoSentP}{}\renewcommand{\rhchLgbmEventLoSentP}{0.481}  % lgbm event lo_sent: DM p [pre-patch]
\providecommand{\rhchLgbmEventAddIvPct}{}\renewcommand{\rhchLgbmEventAddIvPct}{$-$1.4}  % lgbm event add_iv vs lgbm|har: %
\providecommand{\rhchLgbmEventAddIvP}{}\renewcommand{\rhchLgbmEventAddIvP}{0.036}  % lgbm event add_iv: DM p
\providecommand{\rhchLgbmEventLoIvPct}{}\renewcommand{\rhchLgbmEventLoIvPct}{0.7}  % lgbm event lo_iv vs lgbm|full_pre: % [pre-patch]
\providecommand{\rhchLgbmEventLoIvP}{}\renewcommand{\rhchLgbmEventLoIvP}{0.129}  % lgbm event lo_iv: DM p [pre-patch]
\providecommand{\rhchLgbmEventAddVdPct}{}\renewcommand{\rhchLgbmEventAddVdPct}{$-$3.5}  % lgbm event add_vd vs lgbm|har: %
\providecommand{\rhchLgbmEventAddVdP}{}\renewcommand{\rhchLgbmEventAddVdP}{0.072}  % lgbm event add_vd: DM p
\providecommand{\rhchLgbmEventLoVdPct}{}\renewcommand{\rhchLgbmEventLoVdPct}{0.7}  % lgbm event lo_vd vs lgbm|full_pre: % [pre-patch]
\providecommand{\rhchLgbmEventLoVdP}{}\renewcommand{\rhchLgbmEventLoVdP}{0.010}  % lgbm event lo_vd: DM p [pre-patch]
\providecommand{\rhchLgbmEventAddFomcPct}{}\renewcommand{\rhchLgbmEventAddFomcPct}{$-$62.6}  % lgbm event add_fomc vs lgbm|har: %
\providecommand{\rhchLgbmEventAddFomcP}{}\renewcommand{\rhchLgbmEventAddFomcP}{\ensuremath{<}0.001}  % lgbm event add_fomc: DM p
\providecommand{\rhchLgbmEventLoFomcPct}{}\renewcommand{\rhchLgbmEventLoFomcPct}{163.2}  % lgbm event lo_fomc vs lgbm|full: %
\providecommand{\rhchLgbmEventLoFomcP}{}\renewcommand{\rhchLgbmEventLoFomcP}{\ensuremath{<}0.001}  % lgbm event lo_fomc: DM p
\providecommand{\rhchLgbmEventReading}{}\renewcommand{\rhchLgbmEventReading}{concentrated}  % lgbm event: item 3 reading
\providecommand{\rhchLgbmEventTopAddShare}{}\renewcommand{\rhchLgbmEventTopAddShare}{101}  % lgbm event: largest add-alone gain, % of the joint gain (fomc)
\providecommand{\rhchLgbmEventSumAddRatio}{}\renewcommand{\rhchLgbmEventSumAddRatio}{1.05}  % lgbm event: add-alone gains summed over groups / joint gain
\providecommand{\rhchLgbmEventNLoSig}{}\renewcommand{\rhchLgbmEventNLoSig}{2}  % lgbm event: significant leave-one-outs
\providecommand{\rhchLgbmEventNLoRun}{}\renewcommand{\rhchLgbmEventNLoRun}{8}  % lgbm event: leave-one-outs available (of 8)
\providecommand{\rhchLgbmEventConcTest}{}\renewcommand{\rhchLgbmEventConcTest}{yes}  % lgbm event: the CONC definition on its own (pending: undetermined)
\providecommand{\rhchLgbmEventRedTest}{}\renewcommand{\rhchLgbmEventRedTest}{no}  % lgbm event: the RED definition on its own (pending: undetermined)
\providecommand{\rhchLgbmEventReadingRedFirst}{}\renewcommand{\rhchLgbmEventReadingRedFirst}{concentrated}  % lgbm event: reading with REDUNDANT tested before CONCENTRATED (the RULE tests CONCENTRATED first)
\providecommand{\rhchLgbmOrdJointPct}{}\renewcommand{\rhchLgbmOrdJointPct}{$-$6.0}  % lgbm ord: full vs HAR, % QLIKE change
\providecommand{\rhchLgbmOrdJointDm}{}\renewcommand{\rhchLgbmOrdJointDm}{$-$10.70}  % lgbm ord: full vs HAR, DM
\providecommand{\rhchLgbmOrdJointP}{}\renewcommand{\rhchLgbmOrdJointP}{\ensuremath{<}0.001}  % lgbm ord: full vs HAR, DM p
\providecommand{\rhchLgbmOrdAddRmPct}{}\renewcommand{\rhchLgbmOrdAddRmPct}{$-$4.6}  % lgbm ord add_rm vs lgbm|har: %
\providecommand{\rhchLgbmOrdAddRmP}{}\renewcommand{\rhchLgbmOrdAddRmP}{\ensuremath{<}0.001}  % lgbm ord add_rm: DM p
\providecommand{\rhchLgbmOrdLoRmPct}{}\renewcommand{\rhchLgbmOrdLoRmPct}{1.6}  % lgbm ord lo_rm vs lgbm|full_pre: % [pre-patch]
\providecommand{\rhchLgbmOrdLoRmP}{}\renewcommand{\rhchLgbmOrdLoRmP}{\ensuremath{<}0.001}  % lgbm ord lo_rm: DM p [pre-patch]
\providecommand{\rhchLgbmOrdAddLiqPct}{}\renewcommand{\rhchLgbmOrdAddLiqPct}{$-$3.7}  % lgbm ord add_liq vs lgbm|har: %
\providecommand{\rhchLgbmOrdAddLiqP}{}\renewcommand{\rhchLgbmOrdAddLiqP}{\ensuremath{<}0.001}  % lgbm ord add_liq: DM p
\providecommand{\rhchLgbmOrdLoLiqPct}{}\renewcommand{\rhchLgbmOrdLoLiqPct}{0.2}  % lgbm ord lo_liq vs lgbm|full_pre: % [pre-patch]
\providecommand{\rhchLgbmOrdLoLiqP}{}\renewcommand{\rhchLgbmOrdLoLiqP}{0.266}  % lgbm ord lo_liq: DM p [pre-patch]
\providecommand{\rhchLgbmOrdAddEwPct}{}\renewcommand{\rhchLgbmOrdAddEwPct}{$-$2.7}  % lgbm ord add_ew vs lgbm|har: %
\providecommand{\rhchLgbmOrdAddEwP}{}\renewcommand{\rhchLgbmOrdAddEwP}{\ensuremath{<}0.001}  % lgbm ord add_ew: DM p
\providecommand{\rhchLgbmOrdLoEwPct}{}\renewcommand{\rhchLgbmOrdLoEwPct}{0.0}  % lgbm ord lo_ew vs lgbm|full_pre: % [pre-patch]
\providecommand{\rhchLgbmOrdLoEwP}{}\renewcommand{\rhchLgbmOrdLoEwP}{0.779}  % lgbm ord lo_ew: DM p [pre-patch]
\providecommand{\rhchLgbmOrdAddVwPct}{}\renewcommand{\rhchLgbmOrdAddVwPct}{$-$3.1}  % lgbm ord add_vw vs lgbm|har: %
\providecommand{\rhchLgbmOrdAddVwP}{}\renewcommand{\rhchLgbmOrdAddVwP}{\ensuremath{<}0.001}  % lgbm ord add_vw: DM p
\providecommand{\rhchLgbmOrdLoVwPct}{}\renewcommand{\rhchLgbmOrdLoVwPct}{0.1}  % lgbm ord lo_vw vs lgbm|full_pre: % [pre-patch]
\providecommand{\rhchLgbmOrdLoVwP}{}\renewcommand{\rhchLgbmOrdLoVwP}{0.433}  % lgbm ord lo_vw: DM p [pre-patch]
\providecommand{\rhchLgbmOrdAddSentPct}{}\renewcommand{\rhchLgbmOrdAddSentPct}{$-$1.4}  % lgbm ord add_sent vs lgbm|har: %
\providecommand{\rhchLgbmOrdAddSentP}{}\renewcommand{\rhchLgbmOrdAddSentP}{\ensuremath{<}0.001}  % lgbm ord add_sent: DM p
\providecommand{\rhchLgbmOrdLoSentPct}{}\renewcommand{\rhchLgbmOrdLoSentPct}{0.2}  % lgbm ord lo_sent vs lgbm|full_pre: % [pre-patch]
\providecommand{\rhchLgbmOrdLoSentP}{}\renewcommand{\rhchLgbmOrdLoSentP}{0.035}  % lgbm ord lo_sent: DM p [pre-patch]
\providecommand{\rhchLgbmOrdAddIvPct}{}\renewcommand{\rhchLgbmOrdAddIvPct}{$-$1.3}  % lgbm ord add_iv vs lgbm|har: %
\providecommand{\rhchLgbmOrdAddIvP}{}\renewcommand{\rhchLgbmOrdAddIvP}{\ensuremath{<}0.001}  % lgbm ord add_iv: DM p
\providecommand{\rhchLgbmOrdLoIvPct}{}\renewcommand{\rhchLgbmOrdLoIvPct}{0.1}  % lgbm ord lo_iv vs lgbm|full_pre: % [pre-patch]
\providecommand{\rhchLgbmOrdLoIvP}{}\renewcommand{\rhchLgbmOrdLoIvP}{0.196}  % lgbm ord lo_iv: DM p [pre-patch]
\providecommand{\rhchLgbmOrdAddVdPct}{}\renewcommand{\rhchLgbmOrdAddVdPct}{$-$0.9}  % lgbm ord add_vd vs lgbm|har: %
\providecommand{\rhchLgbmOrdAddVdP}{}\renewcommand{\rhchLgbmOrdAddVdP}{\ensuremath{<}0.001}  % lgbm ord add_vd: DM p
\providecommand{\rhchLgbmOrdLoVdPct}{}\renewcommand{\rhchLgbmOrdLoVdPct}{$-$0.1}  % lgbm ord lo_vd vs lgbm|full_pre: % [pre-patch]
\providecommand{\rhchLgbmOrdLoVdP}{}\renewcommand{\rhchLgbmOrdLoVdP}{0.499}  % lgbm ord lo_vd: DM p [pre-patch]
\providecommand{\rhchLgbmOrdAddFomcPct}{}\renewcommand{\rhchLgbmOrdAddFomcPct}{$-$0.2}  % lgbm ord add_fomc vs lgbm|har: %
\providecommand{\rhchLgbmOrdAddFomcP}{}\renewcommand{\rhchLgbmOrdAddFomcP}{0.364}  % lgbm ord add_fomc: DM p
\providecommand{\rhchLgbmOrdLoFomcPct}{}\renewcommand{\rhchLgbmOrdLoFomcPct}{0.0}  % lgbm ord lo_fomc vs lgbm|full: %
\providecommand{\rhchLgbmOrdLoFomcP}{}\renewcommand{\rhchLgbmOrdLoFomcP}{0.912}  % lgbm ord lo_fomc: DM p
\providecommand{\rhchLgbmOrdReading}{}\renewcommand{\rhchLgbmOrdReading}{concentrated}  % lgbm ord: item 3 reading
\providecommand{\rhchLgbmOrdTopAddShare}{}\renewcommand{\rhchLgbmOrdTopAddShare}{77}  % lgbm ord: largest add-alone gain, % of the joint gain (rm)
\providecommand{\rhchLgbmOrdSumAddRatio}{}\renewcommand{\rhchLgbmOrdSumAddRatio}{2.95}  % lgbm ord: add-alone gains summed over groups / joint gain
\providecommand{\rhchLgbmOrdNLoSig}{}\renewcommand{\rhchLgbmOrdNLoSig}{2}  % lgbm ord: significant leave-one-outs
\providecommand{\rhchLgbmOrdNLoRun}{}\renewcommand{\rhchLgbmOrdNLoRun}{8}  % lgbm ord: leave-one-outs available (of 8)
\providecommand{\rhchLgbmOrdConcTest}{}\renewcommand{\rhchLgbmOrdConcTest}{yes}  % lgbm ord: the CONC definition on its own (pending: undetermined)
\providecommand{\rhchLgbmOrdRedTest}{}\renewcommand{\rhchLgbmOrdRedTest}{no}  % lgbm ord: the RED definition on its own (pending: undetermined)
\providecommand{\rhchLgbmOrdReadingRedFirst}{}\renewcommand{\rhchLgbmOrdReadingRedFirst}{concentrated}  % lgbm ord: reading with REDUNDANT tested before CONCENTRATED (the RULE tests CONCENTRATED first)
\providecommand{\rhchXgbBarJointPct}{}\renewcommand{\rhchXgbBarJointPct}{$-$14.6}  % xgb all: full vs HAR, % QLIKE change
\providecommand{\rhchXgbBarJointDm}{}\renewcommand{\rhchXgbBarJointDm}{$-$9.94}  % xgb all: full vs HAR, DM
\providecommand{\rhchXgbBarJointP}{}\renewcommand{\rhchXgbBarJointP}{\ensuremath{<}0.001}  % xgb all: full vs HAR, DM p
\providecommand{\rhchXgbBarAddRmPct}{}\renewcommand{\rhchXgbBarAddRmPct}{$-$3.4}  % xgb all add_rm vs xgb|har: %
\providecommand{\rhchXgbBarAddRmP}{}\renewcommand{\rhchXgbBarAddRmP}{\ensuremath{<}0.001}  % xgb all add_rm: DM p
\providecommand{\rhchXgbBarLoRmPct}{}\renewcommand{\rhchXgbBarLoRmPct}{1.6}  % xgb all lo_rm vs xgb|full_pre: % [pre-patch]
\providecommand{\rhchXgbBarLoRmP}{}\renewcommand{\rhchXgbBarLoRmP}{\ensuremath{<}0.001}  % xgb all lo_rm: DM p [pre-patch]
\providecommand{\rhchXgbBarAddLiqPct}{}\renewcommand{\rhchXgbBarAddLiqPct}{$-$3.5}  % xgb all add_liq vs xgb|har: %
\providecommand{\rhchXgbBarAddLiqP}{}\renewcommand{\rhchXgbBarAddLiqP}{\ensuremath{<}0.001}  % xgb all add_liq: DM p
\providecommand{\rhchXgbBarLoLiqPct}{}\renewcommand{\rhchXgbBarLoLiqPct}{0.2}  % xgb all lo_liq vs xgb|full_pre: % [pre-patch]
\providecommand{\rhchXgbBarLoLiqP}{}\renewcommand{\rhchXgbBarLoLiqP}{0.267}  % xgb all lo_liq: DM p [pre-patch]
\providecommand{\rhchXgbBarAddEwPct}{}\renewcommand{\rhchXgbBarAddEwPct}{$-$3.1}  % xgb all add_ew vs xgb|har: %
\providecommand{\rhchXgbBarAddEwP}{}\renewcommand{\rhchXgbBarAddEwP}{\ensuremath{<}0.001}  % xgb all add_ew: DM p
\providecommand{\rhchXgbBarLoEwPct}{}\renewcommand{\rhchXgbBarLoEwPct}{$-$0.1}  % xgb all lo_ew vs xgb|full_pre: % [pre-patch]
\providecommand{\rhchXgbBarLoEwP}{}\renewcommand{\rhchXgbBarLoEwP}{0.220}  % xgb all lo_ew: DM p [pre-patch]
\providecommand{\rhchXgbBarAddVwPct}{}\renewcommand{\rhchXgbBarAddVwPct}{$-$3.3}  % xgb all add_vw vs xgb|har: %
\providecommand{\rhchXgbBarAddVwP}{}\renewcommand{\rhchXgbBarAddVwP}{\ensuremath{<}0.001}  % xgb all add_vw: DM p
\providecommand{\rhchXgbBarLoVwPct}{}\renewcommand{\rhchXgbBarLoVwPct}{0.1}  % xgb all lo_vw vs xgb|full_pre: % [pre-patch]
\providecommand{\rhchXgbBarLoVwP}{}\renewcommand{\rhchXgbBarLoVwP}{0.349}  % xgb all lo_vw: DM p [pre-patch]
\providecommand{\rhchXgbBarAddSentPct}{}\renewcommand{\rhchXgbBarAddSentPct}{$-$1.7}  % xgb all add_sent vs xgb|har: %
\providecommand{\rhchXgbBarAddSentP}{}\renewcommand{\rhchXgbBarAddSentP}{\ensuremath{<}0.001}  % xgb all add_sent: DM p
\providecommand{\rhchXgbBarLoSentPct}{}\renewcommand{\rhchXgbBarLoSentPct}{0.2}  % xgb all lo_sent vs xgb|full_pre: % [pre-patch]
\providecommand{\rhchXgbBarLoSentP}{}\renewcommand{\rhchXgbBarLoSentP}{0.066}  % xgb all lo_sent: DM p [pre-patch]
\providecommand{\rhchXgbBarAddIvPct}{}\renewcommand{\rhchXgbBarAddIvPct}{$-$1.6}  % xgb all add_iv vs xgb|har: %
\providecommand{\rhchXgbBarAddIvP}{}\renewcommand{\rhchXgbBarAddIvP}{\ensuremath{<}0.001}  % xgb all add_iv: DM p
\providecommand{\rhchXgbBarLoIvPct}{}\renewcommand{\rhchXgbBarLoIvPct}{0.1}  % xgb all lo_iv vs xgb|full_pre: % [pre-patch]
\providecommand{\rhchXgbBarLoIvP}{}\renewcommand{\rhchXgbBarLoIvP}{0.414}  % xgb all lo_iv: DM p [pre-patch]
\providecommand{\rhchXgbBarAddVdPct}{}\renewcommand{\rhchXgbBarAddVdPct}{$-$1.2}  % xgb all add_vd vs xgb|har: %
\providecommand{\rhchXgbBarAddVdP}{}\renewcommand{\rhchXgbBarAddVdP}{\ensuremath{<}0.001}  % xgb all add_vd: DM p
\providecommand{\rhchXgbBarLoVdPct}{}\renewcommand{\rhchXgbBarLoVdPct}{0.1}  % xgb all lo_vd vs xgb|full_pre: % [pre-patch]
\providecommand{\rhchXgbBarLoVdP}{}\renewcommand{\rhchXgbBarLoVdP}{0.674}  % xgb all lo_vd: DM p [pre-patch]
\providecommand{\rhchXgbBarAddFomcPct}{}\renewcommand{\rhchXgbBarAddFomcPct}{$-$9.5}  % xgb all add_fomc vs xgb|har: %
\providecommand{\rhchXgbBarAddFomcP}{}\renewcommand{\rhchXgbBarAddFomcP}{\ensuremath{<}0.001}  % xgb all add_fomc: DM p
\providecommand{\rhchXgbBarLoFomcPct}{}\renewcommand{\rhchXgbBarLoFomcPct}{10.3}  % xgb all lo_fomc vs xgb|full: %
\providecommand{\rhchXgbBarLoFomcP}{}\renewcommand{\rhchXgbBarLoFomcP}{\ensuremath{<}0.001}  % xgb all lo_fomc: DM p
\providecommand{\rhchXgbEventJointPct}{}\renewcommand{\rhchXgbEventJointPct}{$-$63.9}  % xgb event: full vs HAR, % QLIKE change
\providecommand{\rhchXgbEventJointDm}{}\renewcommand{\rhchXgbEventJointDm}{$-$7.75}  % xgb event: full vs HAR, DM
\providecommand{\rhchXgbEventJointP}{}\renewcommand{\rhchXgbEventJointP}{\ensuremath{<}0.001}  % xgb event: full vs HAR, DM p
\providecommand{\rhchXgbEventAddRmPct}{}\renewcommand{\rhchXgbEventAddRmPct}{5.5}  % xgb event add_rm vs xgb|har: %
\providecommand{\rhchXgbEventAddRmP}{}\renewcommand{\rhchXgbEventAddRmP}{0.014}  % xgb event add_rm: DM p
\providecommand{\rhchXgbEventLoRmPct}{}\renewcommand{\rhchXgbEventLoRmPct}{$-$1.3}  % xgb event lo_rm vs xgb|full_pre: % [pre-patch]
\providecommand{\rhchXgbEventLoRmP}{}\renewcommand{\rhchXgbEventLoRmP}{0.170}  % xgb event lo_rm: DM p [pre-patch]
\providecommand{\rhchXgbEventAddLiqPct}{}\renewcommand{\rhchXgbEventAddLiqPct}{0.4}  % xgb event add_liq vs xgb|har: %
\providecommand{\rhchXgbEventAddLiqP}{}\renewcommand{\rhchXgbEventAddLiqP}{0.771}  % xgb event add_liq: DM p
\providecommand{\rhchXgbEventLoLiqPct}{}\renewcommand{\rhchXgbEventLoLiqPct}{$-$0.9}  % xgb event lo_liq vs xgb|full_pre: % [pre-patch]
\providecommand{\rhchXgbEventLoLiqP}{}\renewcommand{\rhchXgbEventLoLiqP}{0.069}  % xgb event lo_liq: DM p [pre-patch]
\providecommand{\rhchXgbEventAddEwPct}{}\renewcommand{\rhchXgbEventAddEwPct}{$-$3.6}  % xgb event add_ew vs xgb|har: %
\providecommand{\rhchXgbEventAddEwP}{}\renewcommand{\rhchXgbEventAddEwP}{0.044}  % xgb event add_ew: DM p
\providecommand{\rhchXgbEventLoEwPct}{}\renewcommand{\rhchXgbEventLoEwPct}{0.0}  % xgb event lo_ew vs xgb|full_pre: % [pre-patch]
\providecommand{\rhchXgbEventLoEwP}{}\renewcommand{\rhchXgbEventLoEwP}{0.990}  % xgb event lo_ew: DM p [pre-patch]
\providecommand{\rhchXgbEventAddVwPct}{}\renewcommand{\rhchXgbEventAddVwPct}{$-$3.3}  % xgb event add_vw vs xgb|har: %
\providecommand{\rhchXgbEventAddVwP}{}\renewcommand{\rhchXgbEventAddVwP}{0.063}  % xgb event add_vw: DM p
\providecommand{\rhchXgbEventLoVwPct}{}\renewcommand{\rhchXgbEventLoVwPct}{$-$0.7}  % xgb event lo_vw vs xgb|full_pre: % [pre-patch]
\providecommand{\rhchXgbEventLoVwP}{}\renewcommand{\rhchXgbEventLoVwP}{0.304}  % xgb event lo_vw: DM p [pre-patch]
\providecommand{\rhchXgbEventAddSentPct}{}\renewcommand{\rhchXgbEventAddSentPct}{0.3}  % xgb event add_sent vs xgb|har: %
\providecommand{\rhchXgbEventAddSentP}{}\renewcommand{\rhchXgbEventAddSentP}{0.765}  % xgb event add_sent: DM p
\providecommand{\rhchXgbEventLoSentPct}{}\renewcommand{\rhchXgbEventLoSentPct}{0.0}  % xgb event lo_sent vs xgb|full_pre: % [pre-patch]
\providecommand{\rhchXgbEventLoSentP}{}\renewcommand{\rhchXgbEventLoSentP}{0.994}  % xgb event lo_sent: DM p [pre-patch]
\providecommand{\rhchXgbEventAddIvPct}{}\renewcommand{\rhchXgbEventAddIvPct}{$-$0.7}  % xgb event add_iv vs xgb|har: %
\providecommand{\rhchXgbEventAddIvP}{}\renewcommand{\rhchXgbEventAddIvP}{0.371}  % xgb event add_iv: DM p
\providecommand{\rhchXgbEventLoIvPct}{}\renewcommand{\rhchXgbEventLoIvPct}{0.4}  % xgb event lo_iv vs xgb|full_pre: % [pre-patch]
\providecommand{\rhchXgbEventLoIvP}{}\renewcommand{\rhchXgbEventLoIvP}{0.482}  % xgb event lo_iv: DM p [pre-patch]
\providecommand{\rhchXgbEventAddVdPct}{}\renewcommand{\rhchXgbEventAddVdPct}{$-$1.9}  % xgb event add_vd vs xgb|har: %
\providecommand{\rhchXgbEventAddVdP}{}\renewcommand{\rhchXgbEventAddVdP}{0.076}  % xgb event add_vd: DM p
\providecommand{\rhchXgbEventLoVdPct}{}\renewcommand{\rhchXgbEventLoVdPct}{$-$0.7}  % xgb event lo_vd vs xgb|full_pre: % [pre-patch]
\providecommand{\rhchXgbEventLoVdP}{}\renewcommand{\rhchXgbEventLoVdP}{0.148}  % xgb event lo_vd: DM p [pre-patch]
\providecommand{\rhchXgbEventAddFomcPct}{}\renewcommand{\rhchXgbEventAddFomcPct}{$-$64.0}  % xgb event add_fomc vs xgb|har: %
\providecommand{\rhchXgbEventAddFomcP}{}\renewcommand{\rhchXgbEventAddFomcP}{\ensuremath{<}0.001}  % xgb event add_fomc: DM p
\providecommand{\rhchXgbEventLoFomcPct}{}\renewcommand{\rhchXgbEventLoFomcPct}{173.0}  % xgb event lo_fomc vs xgb|full: %
\providecommand{\rhchXgbEventLoFomcP}{}\renewcommand{\rhchXgbEventLoFomcP}{\ensuremath{<}0.001}  % xgb event lo_fomc: DM p
\providecommand{\rhchXgbEventReading}{}\renewcommand{\rhchXgbEventReading}{concentrated}  % xgb event: item 3 reading
\providecommand{\rhchXgbEventTopAddShare}{}\renewcommand{\rhchXgbEventTopAddShare}{100}  % xgb event: largest add-alone gain, % of the joint gain (fomc)
\providecommand{\rhchXgbEventSumAddRatio}{}\renewcommand{\rhchXgbEventSumAddRatio}{1.05}  % xgb event: add-alone gains summed over groups / joint gain
\providecommand{\rhchXgbEventNLoSig}{}\renewcommand{\rhchXgbEventNLoSig}{1}  % xgb event: significant leave-one-outs
\providecommand{\rhchXgbEventNLoRun}{}\renewcommand{\rhchXgbEventNLoRun}{8}  % xgb event: leave-one-outs available (of 8)
\providecommand{\rhchXgbEventConcTest}{}\renewcommand{\rhchXgbEventConcTest}{yes}  % xgb event: the CONC definition on its own (pending: undetermined)
\providecommand{\rhchXgbEventRedTest}{}\renewcommand{\rhchXgbEventRedTest}{no}  % xgb event: the RED definition on its own (pending: undetermined)
\providecommand{\rhchXgbEventReadingRedFirst}{}\renewcommand{\rhchXgbEventReadingRedFirst}{concentrated}  % xgb event: reading with REDUNDANT tested before CONCENTRATED (the RULE tests CONCENTRATED first)
\providecommand{\rhchXgbOrdJointPct}{}\renewcommand{\rhchXgbOrdJointPct}{$-$6.1}  % xgb ord: full vs HAR, % QLIKE change
\providecommand{\rhchXgbOrdJointDm}{}\renewcommand{\rhchXgbOrdJointDm}{$-$10.42}  % xgb ord: full vs HAR, DM
\providecommand{\rhchXgbOrdJointP}{}\renewcommand{\rhchXgbOrdJointP}{\ensuremath{<}0.001}  % xgb ord: full vs HAR, DM p
\providecommand{\rhchXgbOrdAddRmPct}{}\renewcommand{\rhchXgbOrdAddRmPct}{$-$4.9}  % xgb ord add_rm vs xgb|har: %
\providecommand{\rhchXgbOrdAddRmP}{}\renewcommand{\rhchXgbOrdAddRmP}{\ensuremath{<}0.001}  % xgb ord add_rm: DM p
\providecommand{\rhchXgbOrdLoRmPct}{}\renewcommand{\rhchXgbOrdLoRmPct}{1.8}  % xgb ord lo_rm vs xgb|full_pre: % [pre-patch]
\providecommand{\rhchXgbOrdLoRmP}{}\renewcommand{\rhchXgbOrdLoRmP}{\ensuremath{<}0.001}  % xgb ord lo_rm: DM p [pre-patch]
\providecommand{\rhchXgbOrdAddLiqPct}{}\renewcommand{\rhchXgbOrdAddLiqPct}{$-$4.1}  % xgb ord add_liq vs xgb|har: %
\providecommand{\rhchXgbOrdAddLiqP}{}\renewcommand{\rhchXgbOrdAddLiqP}{\ensuremath{<}0.001}  % xgb ord add_liq: DM p
\providecommand{\rhchXgbOrdLoLiqPct}{}\renewcommand{\rhchXgbOrdLoLiqPct}{0.2}  % xgb ord lo_liq vs xgb|full_pre: % [pre-patch]
\providecommand{\rhchXgbOrdLoLiqP}{}\renewcommand{\rhchXgbOrdLoLiqP}{0.095}  % xgb ord lo_liq: DM p [pre-patch]
\providecommand{\rhchXgbOrdAddEwPct}{}\renewcommand{\rhchXgbOrdAddEwPct}{$-$3.0}  % xgb ord add_ew vs xgb|har: %
\providecommand{\rhchXgbOrdAddEwP}{}\renewcommand{\rhchXgbOrdAddEwP}{\ensuremath{<}0.001}  % xgb ord add_ew: DM p
\providecommand{\rhchXgbOrdLoEwPct}{}\renewcommand{\rhchXgbOrdLoEwPct}{$-$0.1}  % xgb ord lo_ew vs xgb|full_pre: % [pre-patch]
\providecommand{\rhchXgbOrdLoEwP}{}\renewcommand{\rhchXgbOrdLoEwP}{0.172}  % xgb ord lo_ew: DM p [pre-patch]
\providecommand{\rhchXgbOrdAddVwPct}{}\renewcommand{\rhchXgbOrdAddVwPct}{$-$3.3}  % xgb ord add_vw vs xgb|har: %
\providecommand{\rhchXgbOrdAddVwP}{}\renewcommand{\rhchXgbOrdAddVwP}{\ensuremath{<}0.001}  % xgb ord add_vw: DM p
\providecommand{\rhchXgbOrdLoVwPct}{}\renewcommand{\rhchXgbOrdLoVwPct}{0.2}  % xgb ord lo_vw vs xgb|full_pre: % [pre-patch]
\providecommand{\rhchXgbOrdLoVwP}{}\renewcommand{\rhchXgbOrdLoVwP}{0.133}  % xgb ord lo_vw: DM p [pre-patch]
\providecommand{\rhchXgbOrdAddSentPct}{}\renewcommand{\rhchXgbOrdAddSentPct}{$-$2.0}  % xgb ord add_sent vs xgb|har: %
\providecommand{\rhchXgbOrdAddSentP}{}\renewcommand{\rhchXgbOrdAddSentP}{\ensuremath{<}0.001}  % xgb ord add_sent: DM p
\providecommand{\rhchXgbOrdLoSentPct}{}\renewcommand{\rhchXgbOrdLoSentPct}{0.2}  % xgb ord lo_sent vs xgb|full_pre: % [pre-patch]
\providecommand{\rhchXgbOrdLoSentP}{}\renewcommand{\rhchXgbOrdLoSentP}{0.051}  % xgb ord lo_sent: DM p [pre-patch]
\providecommand{\rhchXgbOrdAddIvPct}{}\renewcommand{\rhchXgbOrdAddIvPct}{$-$1.7}  % xgb ord add_iv vs xgb|har: %
\providecommand{\rhchXgbOrdAddIvP}{}\renewcommand{\rhchXgbOrdAddIvP}{\ensuremath{<}0.001}  % xgb ord add_iv: DM p
\providecommand{\rhchXgbOrdLoIvPct}{}\renewcommand{\rhchXgbOrdLoIvPct}{0.0}  % xgb ord lo_iv vs xgb|full_pre: % [pre-patch]
\providecommand{\rhchXgbOrdLoIvP}{}\renewcommand{\rhchXgbOrdLoIvP}{0.611}  % xgb ord lo_iv: DM p [pre-patch]
\providecommand{\rhchXgbOrdAddVdPct}{}\renewcommand{\rhchXgbOrdAddVdPct}{$-$1.1}  % xgb ord add_vd vs xgb|har: %
\providecommand{\rhchXgbOrdAddVdP}{}\renewcommand{\rhchXgbOrdAddVdP}{\ensuremath{<}0.001}  % xgb ord add_vd: DM p
\providecommand{\rhchXgbOrdLoVdPct}{}\renewcommand{\rhchXgbOrdLoVdPct}{0.1}  % xgb ord lo_vd vs xgb|full_pre: % [pre-patch]
\providecommand{\rhchXgbOrdLoVdP}{}\renewcommand{\rhchXgbOrdLoVdP}{0.393}  % xgb ord lo_vd: DM p [pre-patch]
\providecommand{\rhchXgbOrdAddFomcPct}{}\renewcommand{\rhchXgbOrdAddFomcPct}{$-$0.1}  % xgb ord add_fomc vs xgb|har: %
\providecommand{\rhchXgbOrdAddFomcP}{}\renewcommand{\rhchXgbOrdAddFomcP}{0.726}  % xgb ord add_fomc: DM p
\providecommand{\rhchXgbOrdLoFomcPct}{}\renewcommand{\rhchXgbOrdLoFomcPct}{$-$0.4}  % xgb ord lo_fomc vs xgb|full: %
\providecommand{\rhchXgbOrdLoFomcP}{}\renewcommand{\rhchXgbOrdLoFomcP}{0.013}  % xgb ord lo_fomc: DM p
\providecommand{\rhchXgbOrdReading}{}\renewcommand{\rhchXgbOrdReading}{concentrated}  % xgb ord: item 3 reading
\providecommand{\rhchXgbOrdTopAddShare}{}\renewcommand{\rhchXgbOrdTopAddShare}{81}  % xgb ord: largest add-alone gain, % of the joint gain (rm)
\providecommand{\rhchXgbOrdSumAddRatio}{}\renewcommand{\rhchXgbOrdSumAddRatio}{3.32}  % xgb ord: add-alone gains summed over groups / joint gain
\providecommand{\rhchXgbOrdNLoSig}{}\renewcommand{\rhchXgbOrdNLoSig}{1}  % xgb ord: significant leave-one-outs
\providecommand{\rhchXgbOrdNLoRun}{}\renewcommand{\rhchXgbOrdNLoRun}{8}  % xgb ord: leave-one-outs available (of 8)
\providecommand{\rhchXgbOrdConcTest}{}\renewcommand{\rhchXgbOrdConcTest}{yes}  % xgb ord: the CONC definition on its own (pending: undetermined)
\providecommand{\rhchXgbOrdRedTest}{}\renewcommand{\rhchXgbOrdRedTest}{yes}  % xgb ord: the RED definition on its own (pending: undetermined)
\providecommand{\rhchXgbOrdReadingRedFirst}{}\renewcommand{\rhchXgbOrdReadingRedFirst}{redundant}  % xgb ord: reading with REDUNDANT tested before CONCENTRATED (the RULE tests CONCENTRATED first)
\providecommand{\rhchRfBarJointPct}{}\renewcommand{\rhchRfBarJointPct}{$-$11.3}  % rf all: full vs HAR, % QLIKE change
\providecommand{\rhchRfBarJointDm}{}\renewcommand{\rhchRfBarJointDm}{$-$9.74}  % rf all: full vs HAR, DM
\providecommand{\rhchRfBarJointP}{}\renewcommand{\rhchRfBarJointP}{\ensuremath{<}0.001}  % rf all: full vs HAR, DM p
\providecommand{\rhchRfBarAddRmPct}{}\renewcommand{\rhchRfBarAddRmPct}{$-$3.2}  % rf all add_rm vs rf|har: %
\providecommand{\rhchRfBarAddRmP}{}\renewcommand{\rhchRfBarAddRmP}{\ensuremath{<}0.001}  % rf all add_rm: DM p
\providecommand{\rhchRfBarLoRmPct}{}\renewcommand{\rhchRfBarLoRmPct}{1.5}  % rf all lo_rm vs rf|full_pre: % [pre-patch]
\providecommand{\rhchRfBarLoRmP}{}\renewcommand{\rhchRfBarLoRmP}{\ensuremath{<}0.001}  % rf all lo_rm: DM p [pre-patch]
\providecommand{\rhchRfBarAddLiqPct}{}\renewcommand{\rhchRfBarAddLiqPct}{$-$2.4}  % rf all add_liq vs rf|har: %
\providecommand{\rhchRfBarAddLiqP}{}\renewcommand{\rhchRfBarAddLiqP}{\ensuremath{<}0.001}  % rf all add_liq: DM p
\providecommand{\rhchRfBarAddEwPct}{}\renewcommand{\rhchRfBarAddEwPct}{$-$2.3}  % rf all add_ew vs rf|har: %
\providecommand{\rhchRfBarAddEwP}{}\renewcommand{\rhchRfBarAddEwP}{\ensuremath{<}0.001}  % rf all add_ew: DM p
\providecommand{\rhchRfBarAddVwPct}{}\renewcommand{\rhchRfBarAddVwPct}{$-$2.5}  % rf all add_vw vs rf|har: %
\providecommand{\rhchRfBarAddVwP}{}\renewcommand{\rhchRfBarAddVwP}{\ensuremath{<}0.001}  % rf all add_vw: DM p
\providecommand{\rhchRfBarAddSentPct}{}\renewcommand{\rhchRfBarAddSentPct}{$-$0.9}  % rf all add_sent vs rf|har: %
\providecommand{\rhchRfBarAddSentP}{}\renewcommand{\rhchRfBarAddSentP}{\ensuremath{<}0.001}  % rf all add_sent: DM p
\providecommand{\rhchRfBarAddIvPct}{}\renewcommand{\rhchRfBarAddIvPct}{$-$0.2}  % rf all add_iv vs rf|har: %
\providecommand{\rhchRfBarAddIvP}{}\renewcommand{\rhchRfBarAddIvP}{0.601}  % rf all add_iv: DM p
\providecommand{\rhchRfBarAddVdPct}{}\renewcommand{\rhchRfBarAddVdPct}{$-$0.6}  % rf all add_vd vs rf|har: %
\providecommand{\rhchRfBarAddVdP}{}\renewcommand{\rhchRfBarAddVdP}{\ensuremath{<}0.001}  % rf all add_vd: DM p
\providecommand{\rhchRfBarAddFomcPct}{}\renewcommand{\rhchRfBarAddFomcPct}{$-$7.6}  % rf all add_fomc vs rf|har: %
\providecommand{\rhchRfBarAddFomcP}{}\renewcommand{\rhchRfBarAddFomcP}{\ensuremath{<}0.001}  % rf all add_fomc: DM p
\providecommand{\rhchRfEventJointPct}{}\renewcommand{\rhchRfEventJointPct}{$-$44.8}  % rf event: full vs HAR, % QLIKE change
\providecommand{\rhchRfEventJointDm}{}\renewcommand{\rhchRfEventJointDm}{$-$8.56}  % rf event: full vs HAR, DM
\providecommand{\rhchRfEventJointP}{}\renewcommand{\rhchRfEventJointP}{\ensuremath{<}0.001}  % rf event: full vs HAR, DM p
\providecommand{\rhchRfEventAddRmPct}{}\renewcommand{\rhchRfEventAddRmPct}{5.3}  % rf event add_rm vs rf|har: %
\providecommand{\rhchRfEventAddRmP}{}\renewcommand{\rhchRfEventAddRmP}{0.061}  % rf event add_rm: DM p
\providecommand{\rhchRfEventLoRmPct}{}\renewcommand{\rhchRfEventLoRmPct}{$-$2.1}  % rf event lo_rm vs rf|full_pre: % [pre-patch]
\providecommand{\rhchRfEventLoRmP}{}\renewcommand{\rhchRfEventLoRmP}{0.036}  % rf event lo_rm: DM p [pre-patch]
\providecommand{\rhchRfEventAddLiqPct}{}\renewcommand{\rhchRfEventAddLiqPct}{2.0}  % rf event add_liq vs rf|har: %
\providecommand{\rhchRfEventAddLiqP}{}\renewcommand{\rhchRfEventAddLiqP}{0.184}  % rf event add_liq: DM p
\providecommand{\rhchRfEventAddEwPct}{}\renewcommand{\rhchRfEventAddEwPct}{$-$2.4}  % rf event add_ew vs rf|har: %
\providecommand{\rhchRfEventAddEwP}{}\renewcommand{\rhchRfEventAddEwP}{0.073}  % rf event add_ew: DM p
\providecommand{\rhchRfEventAddVwPct}{}\renewcommand{\rhchRfEventAddVwPct}{$-$0.9}  % rf event add_vw vs rf|har: %
\providecommand{\rhchRfEventAddVwP}{}\renewcommand{\rhchRfEventAddVwP}{0.651}  % rf event add_vw: DM p
\providecommand{\rhchRfEventAddSentPct}{}\renewcommand{\rhchRfEventAddSentPct}{0.4}  % rf event add_sent vs rf|har: %
\providecommand{\rhchRfEventAddSentP}{}\renewcommand{\rhchRfEventAddSentP}{0.621}  % rf event add_sent: DM p
\providecommand{\rhchRfEventAddIvPct}{}\renewcommand{\rhchRfEventAddIvPct}{4.1}  % rf event add_iv vs rf|har: %
\providecommand{\rhchRfEventAddIvP}{}\renewcommand{\rhchRfEventAddIvP}{0.001}  % rf event add_iv: DM p
\providecommand{\rhchRfEventAddVdPct}{}\renewcommand{\rhchRfEventAddVdPct}{1.3}  % rf event add_vd vs rf|har: %
\providecommand{\rhchRfEventAddVdP}{}\renewcommand{\rhchRfEventAddVdP}{0.046}  % rf event add_vd: DM p
\providecommand{\rhchRfEventAddFomcPct}{}\renewcommand{\rhchRfEventAddFomcPct}{$-$48.5}  % rf event add_fomc vs rf|har: %
\providecommand{\rhchRfEventAddFomcP}{}\renewcommand{\rhchRfEventAddFomcP}{\ensuremath{<}0.001}  % rf event add_fomc: DM p
\providecommand{\rhchRfEventReading}{}\renewcommand{\rhchRfEventReading}{n/a}  % rf event: item 3 reading
\providecommand{\rhchRfEventTopAddShare}{}\renewcommand{\rhchRfEventTopAddShare}{108}  % rf event: largest add-alone gain, % of the joint gain (fomc)
\providecommand{\rhchRfEventSumAddRatio}{}\renewcommand{\rhchRfEventSumAddRatio}{0.86}  % rf event: add-alone gains summed over groups / joint gain
\providecommand{\rhchRfEventNLoSig}{}\renewcommand{\rhchRfEventNLoSig}{0}  % rf event: significant leave-one-outs
\providecommand{\rhchRfEventNLoRun}{}\renewcommand{\rhchRfEventNLoRun}{1}  % rf event: leave-one-outs available (of 8)
\providecommand{\rhchRfEventConcTest}{}\renewcommand{\rhchRfEventConcTest}{\rhPending{}}  % rf event: the CONC definition on its own (pending: undetermined)
\providecommand{\rhchRfEventRedTest}{}\renewcommand{\rhchRfEventRedTest}{no}  % rf event: the RED definition on its own (pending: undetermined)
\providecommand{\rhchRfEventReadingRedFirst}{}\renewcommand{\rhchRfEventReadingRedFirst}{n/a}  % rf event: reading with REDUNDANT tested before CONCENTRATED (the RULE tests CONCENTRATED first)
\providecommand{\rhchRfOrdJointPct}{}\renewcommand{\rhchRfOrdJointPct}{$-$5.3}  % rf ord: full vs HAR, % QLIKE change
\providecommand{\rhchRfOrdJointDm}{}\renewcommand{\rhchRfOrdJointDm}{$-$11.16}  % rf ord: full vs HAR, DM
\providecommand{\rhchRfOrdJointP}{}\renewcommand{\rhchRfOrdJointP}{\ensuremath{<}0.001}  % rf ord: full vs HAR, DM p
\providecommand{\rhchRfOrdAddRmPct}{}\renewcommand{\rhchRfOrdAddRmPct}{$-$4.7}  % rf ord add_rm vs rf|har: %
\providecommand{\rhchRfOrdAddRmP}{}\renewcommand{\rhchRfOrdAddRmP}{\ensuremath{<}0.001}  % rf ord add_rm: DM p
\providecommand{\rhchRfOrdLoRmPct}{}\renewcommand{\rhchRfOrdLoRmPct}{1.8}  % rf ord lo_rm vs rf|full_pre: % [pre-patch]
\providecommand{\rhchRfOrdLoRmP}{}\renewcommand{\rhchRfOrdLoRmP}{\ensuremath{<}0.001}  % rf ord lo_rm: DM p [pre-patch]
\providecommand{\rhchRfOrdAddLiqPct}{}\renewcommand{\rhchRfOrdAddLiqPct}{$-$3.2}  % rf ord add_liq vs rf|har: %
\providecommand{\rhchRfOrdAddLiqP}{}\renewcommand{\rhchRfOrdAddLiqP}{\ensuremath{<}0.001}  % rf ord add_liq: DM p
\providecommand{\rhchRfOrdAddEwPct}{}\renewcommand{\rhchRfOrdAddEwPct}{$-$2.3}  % rf ord add_ew vs rf|har: %
\providecommand{\rhchRfOrdAddEwP}{}\renewcommand{\rhchRfOrdAddEwP}{\ensuremath{<}0.001}  % rf ord add_ew: DM p
\providecommand{\rhchRfOrdAddVwPct}{}\renewcommand{\rhchRfOrdAddVwPct}{$-$2.8}  % rf ord add_vw vs rf|har: %
\providecommand{\rhchRfOrdAddVwP}{}\renewcommand{\rhchRfOrdAddVwP}{\ensuremath{<}0.001}  % rf ord add_vw: DM p
\providecommand{\rhchRfOrdAddSentPct}{}\renewcommand{\rhchRfOrdAddSentPct}{$-$1.1}  % rf ord add_sent vs rf|har: %
\providecommand{\rhchRfOrdAddSentP}{}\renewcommand{\rhchRfOrdAddSentP}{\ensuremath{<}0.001}  % rf ord add_sent: DM p
\providecommand{\rhchRfOrdAddIvPct}{}\renewcommand{\rhchRfOrdAddIvPct}{$-$1.0}  % rf ord add_iv vs rf|har: %
\providecommand{\rhchRfOrdAddIvP}{}\renewcommand{\rhchRfOrdAddIvP}{0.004}  % rf ord add_iv: DM p
\providecommand{\rhchRfOrdAddVdPct}{}\renewcommand{\rhchRfOrdAddVdPct}{$-$1.0}  % rf ord add_vd vs rf|har: %
\providecommand{\rhchRfOrdAddVdP}{}\renewcommand{\rhchRfOrdAddVdP}{\ensuremath{<}0.001}  % rf ord add_vd: DM p
\providecommand{\rhchRfOrdAddFomcPct}{}\renewcommand{\rhchRfOrdAddFomcPct}{$-$0.3}  % rf ord add_fomc vs rf|har: %
\providecommand{\rhchRfOrdAddFomcP}{}\renewcommand{\rhchRfOrdAddFomcP}{0.141}  % rf ord add_fomc: DM p
\providecommand{\rhchRfOrdReading}{}\renewcommand{\rhchRfOrdReading}{concentrated}  % rf ord: item 3 reading
\providecommand{\rhchRfOrdTopAddShare}{}\renewcommand{\rhchRfOrdTopAddShare}{88}  % rf ord: largest add-alone gain, % of the joint gain (rm)
\providecommand{\rhchRfOrdSumAddRatio}{}\renewcommand{\rhchRfOrdSumAddRatio}{3.07}  % rf ord: add-alone gains summed over groups / joint gain
\providecommand{\rhchRfOrdNLoSig}{}\renewcommand{\rhchRfOrdNLoSig}{1}  % rf ord: significant leave-one-outs
\providecommand{\rhchRfOrdNLoRun}{}\renewcommand{\rhchRfOrdNLoRun}{1}  % rf ord: leave-one-outs available (of 8)
\providecommand{\rhchRfOrdConcTest}{}\renewcommand{\rhchRfOrdConcTest}{yes}  % rf ord: the CONC definition on its own (pending: undetermined)
\providecommand{\rhchRfOrdRedTest}{}\renewcommand{\rhchRfOrdRedTest}{\rhPending{}}  % rf ord: the RED definition on its own (pending: undetermined)
\providecommand{\rhchRfOrdReadingRedFirst}{}\renewcommand{\rhchRfOrdReadingRedFirst}{n/a}  % rf ord: reading with REDUNDANT tested before CONCENTRATED (the RULE tests CONCENTRATED first)
\providecommand{\rhchLinEventReading}{}\renewcommand{\rhchLinEventReading}{concentrated}  % lin event: item 3 reading by model type
\providecommand{\rhchLinEventReadingRedFirst}{}\renewcommand{\rhchLinEventReadingRedFirst}{concentrated}  % lin event: by model type, REDUNDANT tested first
\providecommand{\rhchLinEventNConcTest}{}\renewcommand{\rhchLinEventNConcTest}{3}  % lin event: classes meeting the Conc definition (of 3)
\providecommand{\rhchLinEventNRedTest}{}\renewcommand{\rhchLinEventNRedTest}{0}  % lin event: classes meeting the Red definition (of 3)
\providecommand{\rhchTreeEventReading}{}\renewcommand{\rhchTreeEventReading}{concentrated}  % tree event: item 3 reading by model type
\providecommand{\rhchTreeEventReadingRedFirst}{}\renewcommand{\rhchTreeEventReadingRedFirst}{concentrated}  % tree event: by model type, REDUNDANT tested first
\providecommand{\rhchTreeEventNConcTest}{}\renewcommand{\rhchTreeEventNConcTest}{2}  % tree event: classes meeting the Conc definition (of 3)
\providecommand{\rhchTreeEventNRedTest}{}\renewcommand{\rhchTreeEventNRedTest}{0}  % tree event: classes meeting the Red definition (of 3)
\providecommand{\rhchLinOrdReading}{}\renewcommand{\rhchLinOrdReading}{concentrated}  % lin ord: item 3 reading by model type
\providecommand{\rhchLinOrdReadingRedFirst}{}\renewcommand{\rhchLinOrdReadingRedFirst}{redundant}  % lin ord: by model type, REDUNDANT tested first
\providecommand{\rhchLinOrdNConcTest}{}\renewcommand{\rhchLinOrdNConcTest}{3}  % lin ord: classes meeting the Conc definition (of 3)
\providecommand{\rhchLinOrdNRedTest}{}\renewcommand{\rhchLinOrdNRedTest}{3}  % lin ord: classes meeting the Red definition (of 3)
\providecommand{\rhchTreeOrdReading}{}\renewcommand{\rhchTreeOrdReading}{concentrated}  % tree ord: item 3 reading by model type
\providecommand{\rhchTreeOrdReadingRedFirst}{}\renewcommand{\rhchTreeOrdReadingRedFirst}{split}  % tree ord: by model type, REDUNDANT tested first
\providecommand{\rhchTreeOrdNConcTest}{}\renewcommand{\rhchTreeOrdNConcTest}{3}  % tree ord: classes meeting the Conc definition (of 3)
\providecommand{\rhchTreeOrdNRedTest}{}\renewcommand{\rhchTreeOrdNRedTest}{1}  % tree ord: classes meeting the Red definition (of 3)
\providecommand{\rhchXRidgeBarFullVsRmPct}{}\renewcommand{\rhchXRidgeBarFullVsRmPct}{$-$3.5}  % ADDITIONAL, NOT PRE-REGISTERED: ridge all: full vs HAR+RM: % QLIKE change
\providecommand{\rhchXRidgeBarFullVsRmDm}{}\renewcommand{\rhchXRidgeBarFullVsRmDm}{$-$4.52}  % ADDITIONAL, NOT PRE-REGISTERED: ridge all: full vs HAR+RM: DM
\providecommand{\rhchXRidgeBarFullVsRmP}{}\renewcommand{\rhchXRidgeBarFullVsRmP}{\ensuremath{<}0.001}  % ADDITIONAL, NOT PRE-REGISTERED: ridge all: full vs HAR+RM: DM p
\providecommand{\rhchXRidgeBarFullVsRmShare}{}\renewcommand{\rhchXRidgeBarFullVsRmShare}{56}  % ADDITIONAL, NOT PRE-REGISTERED: ridge all: full vs HAR+RM: % of the joint gain
\providecommand{\rhchXRidgeBarRmVsHarPct}{}\renewcommand{\rhchXRidgeBarRmVsHarPct}{$-$2.6}  % ADDITIONAL, NOT PRE-REGISTERED: ridge all: HAR+RM vs HAR: % QLIKE change
\providecommand{\rhchXRidgeBarRmVsHarDm}{}\renewcommand{\rhchXRidgeBarRmVsHarDm}{$-$4.14}  % ADDITIONAL, NOT PRE-REGISTERED: ridge all: HAR+RM vs HAR: DM
\providecommand{\rhchXRidgeBarRmVsHarP}{}\renewcommand{\rhchXRidgeBarRmVsHarP}{\ensuremath{<}0.001}  % ADDITIONAL, NOT PRE-REGISTERED: ridge all: HAR+RM vs HAR: DM p
\providecommand{\rhchXRidgeBarRmVsHarShare}{}\renewcommand{\rhchXRidgeBarRmVsHarShare}{44}  % ADDITIONAL, NOT PRE-REGISTERED: ridge all: HAR+RM vs HAR: % of the joint gain
\providecommand{\rhchXRidgeEventFullVsRmPct}{}\renewcommand{\rhchXRidgeEventFullVsRmPct}{$-$25.7}  % ADDITIONAL, NOT PRE-REGISTERED: ridge event: full vs HAR+RM: % QLIKE change
\providecommand{\rhchXRidgeEventFullVsRmDm}{}\renewcommand{\rhchXRidgeEventFullVsRmDm}{$-$5.37}  % ADDITIONAL, NOT PRE-REGISTERED: ridge event: full vs HAR+RM: DM
\providecommand{\rhchXRidgeEventFullVsRmP}{}\renewcommand{\rhchXRidgeEventFullVsRmP}{\ensuremath{<}0.001}  % ADDITIONAL, NOT PRE-REGISTERED: ridge event: full vs HAR+RM: DM p
\providecommand{\rhchXRidgeEventFullVsRmShare}{}\renewcommand{\rhchXRidgeEventFullVsRmShare}{120}  % ADDITIONAL, NOT PRE-REGISTERED: ridge event: full vs HAR+RM: % of the joint gain
\providecommand{\rhchXRidgeEventRmVsHarPct}{}\renewcommand{\rhchXRidgeEventRmVsHarPct}{4.4}  % ADDITIONAL, NOT PRE-REGISTERED: ridge event: HAR+RM vs HAR: % QLIKE change
\providecommand{\rhchXRidgeEventRmVsHarDm}{}\renewcommand{\rhchXRidgeEventRmVsHarDm}{2.28}  % ADDITIONAL, NOT PRE-REGISTERED: ridge event: HAR+RM vs HAR: DM
\providecommand{\rhchXRidgeEventRmVsHarP}{}\renewcommand{\rhchXRidgeEventRmVsHarP}{0.023}  % ADDITIONAL, NOT PRE-REGISTERED: ridge event: HAR+RM vs HAR: DM p
\providecommand{\rhchXRidgeEventRmVsHarShare}{}\renewcommand{\rhchXRidgeEventRmVsHarShare}{$-$20}  % ADDITIONAL, NOT PRE-REGISTERED: ridge event: HAR+RM vs HAR: % of the joint gain
\providecommand{\rhchXRidgeOrdFullVsRmPct}{}\renewcommand{\rhchXRidgeOrdFullVsRmPct}{0.6}  % ADDITIONAL, NOT PRE-REGISTERED: ridge ord: full vs HAR+RM: % QLIKE change
\providecommand{\rhchXRidgeOrdFullVsRmDm}{}\renewcommand{\rhchXRidgeOrdFullVsRmDm}{1.24}  % ADDITIONAL, NOT PRE-REGISTERED: ridge ord: full vs HAR+RM: DM
\providecommand{\rhchXRidgeOrdFullVsRmP}{}\renewcommand{\rhchXRidgeOrdFullVsRmP}{0.215}  % ADDITIONAL, NOT PRE-REGISTERED: ridge ord: full vs HAR+RM: DM p
\providecommand{\rhchXRidgeOrdFullVsRmShare}{}\renewcommand{\rhchXRidgeOrdFullVsRmShare}{$-$18}  % ADDITIONAL, NOT PRE-REGISTERED: ridge ord: full vs HAR+RM: % of the joint gain
\providecommand{\rhchXRidgeOrdRmVsHarPct}{}\renewcommand{\rhchXRidgeOrdRmVsHarPct}{$-$3.8}  % ADDITIONAL, NOT PRE-REGISTERED: ridge ord: HAR+RM vs HAR: % QLIKE change
\providecommand{\rhchXRidgeOrdRmVsHarDm}{}\renewcommand{\rhchXRidgeOrdRmVsHarDm}{$-$5.48}  % ADDITIONAL, NOT PRE-REGISTERED: ridge ord: HAR+RM vs HAR: DM
\providecommand{\rhchXRidgeOrdRmVsHarP}{}\renewcommand{\rhchXRidgeOrdRmVsHarP}{\ensuremath{<}0.001}  % ADDITIONAL, NOT PRE-REGISTERED: ridge ord: HAR+RM vs HAR: DM p
\providecommand{\rhchXRidgeOrdRmVsHarShare}{}\renewcommand{\rhchXRidgeOrdRmVsHarShare}{118}  % ADDITIONAL, NOT PRE-REGISTERED: ridge ord: HAR+RM vs HAR: % of the joint gain
\providecommand{\rhchXLassoBarFullVsRmPct}{}\renewcommand{\rhchXLassoBarFullVsRmPct}{$-$7.0}  % ADDITIONAL, NOT PRE-REGISTERED: lasso all: full vs HAR+RM: % QLIKE change
\providecommand{\rhchXLassoBarFullVsRmDm}{}\renewcommand{\rhchXLassoBarFullVsRmDm}{$-$5.66}  % ADDITIONAL, NOT PRE-REGISTERED: lasso all: full vs HAR+RM: DM
\providecommand{\rhchXLassoBarFullVsRmP}{}\renewcommand{\rhchXLassoBarFullVsRmP}{\ensuremath{<}0.001}  % ADDITIONAL, NOT PRE-REGISTERED: lasso all: full vs HAR+RM: DM p
\providecommand{\rhchXLassoBarFullVsRmShare}{}\renewcommand{\rhchXLassoBarFullVsRmShare}{76}  % ADDITIONAL, NOT PRE-REGISTERED: lasso all: full vs HAR+RM: % of the joint gain
\providecommand{\rhchXLassoBarRmVsHarPct}{}\renewcommand{\rhchXLassoBarRmVsHarPct}{$-$2.1}  % ADDITIONAL, NOT PRE-REGISTERED: lasso all: HAR+RM vs HAR: % QLIKE change
\providecommand{\rhchXLassoBarRmVsHarDm}{}\renewcommand{\rhchXLassoBarRmVsHarDm}{$-$2.64}  % ADDITIONAL, NOT PRE-REGISTERED: lasso all: HAR+RM vs HAR: DM
\providecommand{\rhchXLassoBarRmVsHarP}{}\renewcommand{\rhchXLassoBarRmVsHarP}{0.008}  % ADDITIONAL, NOT PRE-REGISTERED: lasso all: HAR+RM vs HAR: DM p
\providecommand{\rhchXLassoBarRmVsHarShare}{}\renewcommand{\rhchXLassoBarRmVsHarShare}{24}  % ADDITIONAL, NOT PRE-REGISTERED: lasso all: HAR+RM vs HAR: % of the joint gain
\providecommand{\rhchXLassoEventFullVsRmPct}{}\renewcommand{\rhchXLassoEventFullVsRmPct}{$-$49.5}  % ADDITIONAL, NOT PRE-REGISTERED: lasso event: full vs HAR+RM: % QLIKE change
\providecommand{\rhchXLassoEventFullVsRmDm}{}\renewcommand{\rhchXLassoEventFullVsRmDm}{$-$6.59}  % ADDITIONAL, NOT PRE-REGISTERED: lasso event: full vs HAR+RM: DM
\providecommand{\rhchXLassoEventFullVsRmP}{}\renewcommand{\rhchXLassoEventFullVsRmP}{\ensuremath{<}0.001}  % ADDITIONAL, NOT PRE-REGISTERED: lasso event: full vs HAR+RM: DM p
\providecommand{\rhchXLassoEventFullVsRmShare}{}\renewcommand{\rhchXLassoEventFullVsRmShare}{110}  % ADDITIONAL, NOT PRE-REGISTERED: lasso event: full vs HAR+RM: % of the joint gain
\providecommand{\rhchXLassoEventRmVsHarPct}{}\renewcommand{\rhchXLassoEventRmVsHarPct}{4.7}  % ADDITIONAL, NOT PRE-REGISTERED: lasso event: HAR+RM vs HAR: % QLIKE change
\providecommand{\rhchXLassoEventRmVsHarDm}{}\renewcommand{\rhchXLassoEventRmVsHarDm}{2.25}  % ADDITIONAL, NOT PRE-REGISTERED: lasso event: HAR+RM vs HAR: DM
\providecommand{\rhchXLassoEventRmVsHarP}{}\renewcommand{\rhchXLassoEventRmVsHarP}{0.024}  % ADDITIONAL, NOT PRE-REGISTERED: lasso event: HAR+RM vs HAR: DM p
\providecommand{\rhchXLassoEventRmVsHarShare}{}\renewcommand{\rhchXLassoEventRmVsHarShare}{$-$10}  % ADDITIONAL, NOT PRE-REGISTERED: lasso event: HAR+RM vs HAR: % of the joint gain
\providecommand{\rhchXLassoOrdFullVsRmPct}{}\renewcommand{\rhchXLassoOrdFullVsRmPct}{0.9}  % ADDITIONAL, NOT PRE-REGISTERED: lasso ord: full vs HAR+RM: % QLIKE change
\providecommand{\rhchXLassoOrdFullVsRmDm}{}\renewcommand{\rhchXLassoOrdFullVsRmDm}{1.44}  % ADDITIONAL, NOT PRE-REGISTERED: lasso ord: full vs HAR+RM: DM
\providecommand{\rhchXLassoOrdFullVsRmP}{}\renewcommand{\rhchXLassoOrdFullVsRmP}{0.151}  % ADDITIONAL, NOT PRE-REGISTERED: lasso ord: full vs HAR+RM: DM p
\providecommand{\rhchXLassoOrdFullVsRmShare}{}\renewcommand{\rhchXLassoOrdFullVsRmShare}{$-$36}  % ADDITIONAL, NOT PRE-REGISTERED: lasso ord: full vs HAR+RM: % of the joint gain
\providecommand{\rhchXLassoOrdRmVsHarPct}{}\renewcommand{\rhchXLassoOrdRmVsHarPct}{$-$3.3}  % ADDITIONAL, NOT PRE-REGISTERED: lasso ord: HAR+RM vs HAR: % QLIKE change
\providecommand{\rhchXLassoOrdRmVsHarDm}{}\renewcommand{\rhchXLassoOrdRmVsHarDm}{$-$3.61}  % ADDITIONAL, NOT PRE-REGISTERED: lasso ord: HAR+RM vs HAR: DM
\providecommand{\rhchXLassoOrdRmVsHarP}{}\renewcommand{\rhchXLassoOrdRmVsHarP}{\ensuremath{<}0.001}  % ADDITIONAL, NOT PRE-REGISTERED: lasso ord: HAR+RM vs HAR: DM p
\providecommand{\rhchXLassoOrdRmVsHarShare}{}\renewcommand{\rhchXLassoOrdRmVsHarShare}{136}  % ADDITIONAL, NOT PRE-REGISTERED: lasso ord: HAR+RM vs HAR: % of the joint gain
\providecommand{\rhchXEnetBarFullVsRmPct}{}\renewcommand{\rhchXEnetBarFullVsRmPct}{$-$4.1}  % ADDITIONAL, NOT PRE-REGISTERED: enet all: full vs HAR+RM: % QLIKE change
\providecommand{\rhchXEnetBarFullVsRmDm}{}\renewcommand{\rhchXEnetBarFullVsRmDm}{$-$5.48}  % ADDITIONAL, NOT PRE-REGISTERED: enet all: full vs HAR+RM: DM
\providecommand{\rhchXEnetBarFullVsRmP}{}\renewcommand{\rhchXEnetBarFullVsRmP}{\ensuremath{<}0.001}  % ADDITIONAL, NOT PRE-REGISTERED: enet all: full vs HAR+RM: DM p
\providecommand{\rhchXEnetBarFullVsRmShare}{}\renewcommand{\rhchXEnetBarFullVsRmShare}{60}  % ADDITIONAL, NOT PRE-REGISTERED: enet all: full vs HAR+RM: % of the joint gain
\providecommand{\rhchXEnetBarRmVsHarPct}{}\renewcommand{\rhchXEnetBarRmVsHarPct}{$-$2.7}  % ADDITIONAL, NOT PRE-REGISTERED: enet all: HAR+RM vs HAR: % QLIKE change
\providecommand{\rhchXEnetBarRmVsHarDm}{}\renewcommand{\rhchXEnetBarRmVsHarDm}{$-$4.25}  % ADDITIONAL, NOT PRE-REGISTERED: enet all: HAR+RM vs HAR: DM
\providecommand{\rhchXEnetBarRmVsHarP}{}\renewcommand{\rhchXEnetBarRmVsHarP}{\ensuremath{<}0.001}  % ADDITIONAL, NOT PRE-REGISTERED: enet all: HAR+RM vs HAR: DM p
\providecommand{\rhchXEnetBarRmVsHarShare}{}\renewcommand{\rhchXEnetBarRmVsHarShare}{40}  % ADDITIONAL, NOT PRE-REGISTERED: enet all: HAR+RM vs HAR: % of the joint gain
\providecommand{\rhchXEnetEventFullVsRmPct}{}\renewcommand{\rhchXEnetEventFullVsRmPct}{$-$25.7}  % ADDITIONAL, NOT PRE-REGISTERED: enet event: full vs HAR+RM: % QLIKE change
\providecommand{\rhchXEnetEventFullVsRmDm}{}\renewcommand{\rhchXEnetEventFullVsRmDm}{$-$5.65}  % ADDITIONAL, NOT PRE-REGISTERED: enet event: full vs HAR+RM: DM
\providecommand{\rhchXEnetEventFullVsRmP}{}\renewcommand{\rhchXEnetEventFullVsRmP}{\ensuremath{<}0.001}  % ADDITIONAL, NOT PRE-REGISTERED: enet event: full vs HAR+RM: DM p
\providecommand{\rhchXEnetEventFullVsRmShare}{}\renewcommand{\rhchXEnetEventFullVsRmShare}{118}  % ADDITIONAL, NOT PRE-REGISTERED: enet event: full vs HAR+RM: % of the joint gain
\providecommand{\rhchXEnetEventRmVsHarPct}{}\renewcommand{\rhchXEnetEventRmVsHarPct}{4.2}  % ADDITIONAL, NOT PRE-REGISTERED: enet event: HAR+RM vs HAR: % QLIKE change
\providecommand{\rhchXEnetEventRmVsHarDm}{}\renewcommand{\rhchXEnetEventRmVsHarDm}{2.16}  % ADDITIONAL, NOT PRE-REGISTERED: enet event: HAR+RM vs HAR: DM
\providecommand{\rhchXEnetEventRmVsHarP}{}\renewcommand{\rhchXEnetEventRmVsHarP}{0.031}  % ADDITIONAL, NOT PRE-REGISTERED: enet event: HAR+RM vs HAR: DM p
\providecommand{\rhchXEnetEventRmVsHarShare}{}\renewcommand{\rhchXEnetEventRmVsHarShare}{$-$18}  % ADDITIONAL, NOT PRE-REGISTERED: enet event: HAR+RM vs HAR: % of the joint gain
\providecommand{\rhchXEnetOrdFullVsRmPct}{}\renewcommand{\rhchXEnetOrdFullVsRmPct}{$-$0.1}  % ADDITIONAL, NOT PRE-REGISTERED: enet ord: full vs HAR+RM: % QLIKE change
\providecommand{\rhchXEnetOrdFullVsRmDm}{}\renewcommand{\rhchXEnetOrdFullVsRmDm}{$-$0.26}  % ADDITIONAL, NOT PRE-REGISTERED: enet ord: full vs HAR+RM: DM
\providecommand{\rhchXEnetOrdFullVsRmP}{}\renewcommand{\rhchXEnetOrdFullVsRmP}{0.793}  % ADDITIONAL, NOT PRE-REGISTERED: enet ord: full vs HAR+RM: DM p
\providecommand{\rhchXEnetOrdFullVsRmShare}{}\renewcommand{\rhchXEnetOrdFullVsRmShare}{3}  % ADDITIONAL, NOT PRE-REGISTERED: enet ord: full vs HAR+RM: % of the joint gain
\providecommand{\rhchXEnetOrdRmVsHarPct}{}\renewcommand{\rhchXEnetOrdRmVsHarPct}{$-$3.9}  % ADDITIONAL, NOT PRE-REGISTERED: enet ord: HAR+RM vs HAR: % QLIKE change
\providecommand{\rhchXEnetOrdRmVsHarDm}{}\renewcommand{\rhchXEnetOrdRmVsHarDm}{$-$5.53}  % ADDITIONAL, NOT PRE-REGISTERED: enet ord: HAR+RM vs HAR: DM
\providecommand{\rhchXEnetOrdRmVsHarP}{}\renewcommand{\rhchXEnetOrdRmVsHarP}{\ensuremath{<}0.001}  % ADDITIONAL, NOT PRE-REGISTERED: enet ord: HAR+RM vs HAR: DM p
\providecommand{\rhchXEnetOrdRmVsHarShare}{}\renewcommand{\rhchXEnetOrdRmVsHarShare}{97}  % ADDITIONAL, NOT PRE-REGISTERED: enet ord: HAR+RM vs HAR: % of the joint gain
\providecommand{\rhchXLgbmBarFullVsRmPct}{}\renewcommand{\rhchXLgbmBarFullVsRmPct}{$-$11.5}  % ADDITIONAL, NOT PRE-REGISTERED: lgbm all: full vs HAR+RM: % QLIKE change
\providecommand{\rhchXLgbmBarFullVsRmDm}{}\renewcommand{\rhchXLgbmBarFullVsRmDm}{$-$7.62}  % ADDITIONAL, NOT PRE-REGISTERED: lgbm all: full vs HAR+RM: DM
\providecommand{\rhchXLgbmBarFullVsRmP}{}\renewcommand{\rhchXLgbmBarFullVsRmP}{\ensuremath{<}0.001}  % ADDITIONAL, NOT PRE-REGISTERED: lgbm all: full vs HAR+RM: DM p
\providecommand{\rhchXLgbmBarFullVsRmShare}{}\renewcommand{\rhchXLgbmBarFullVsRmShare}{79}  % ADDITIONAL, NOT PRE-REGISTERED: lgbm all: full vs HAR+RM: % of the joint gain
\providecommand{\rhchXLgbmBarRmVsHarPct}{}\renewcommand{\rhchXLgbmBarRmVsHarPct}{$-$3.0}  % ADDITIONAL, NOT PRE-REGISTERED: lgbm all: HAR+RM vs HAR: % QLIKE change
\providecommand{\rhchXLgbmBarRmVsHarDm}{}\renewcommand{\rhchXLgbmBarRmVsHarDm}{$-$6.32}  % ADDITIONAL, NOT PRE-REGISTERED: lgbm all: HAR+RM vs HAR: DM
\providecommand{\rhchXLgbmBarRmVsHarP}{}\renewcommand{\rhchXLgbmBarRmVsHarP}{\ensuremath{<}0.001}  % ADDITIONAL, NOT PRE-REGISTERED: lgbm all: HAR+RM vs HAR: DM p
\providecommand{\rhchXLgbmBarRmVsHarShare}{}\renewcommand{\rhchXLgbmBarRmVsHarShare}{21}  % ADDITIONAL, NOT PRE-REGISTERED: lgbm all: HAR+RM vs HAR: % of the joint gain
\providecommand{\rhchXLgbmEventFullVsRmPct}{}\renewcommand{\rhchXLgbmEventFullVsRmPct}{$-$64.7}  % ADDITIONAL, NOT PRE-REGISTERED: lgbm event: full vs HAR+RM: % QLIKE change
\providecommand{\rhchXLgbmEventFullVsRmDm}{}\renewcommand{\rhchXLgbmEventFullVsRmDm}{$-$7.56}  % ADDITIONAL, NOT PRE-REGISTERED: lgbm event: full vs HAR+RM: DM
\providecommand{\rhchXLgbmEventFullVsRmP}{}\renewcommand{\rhchXLgbmEventFullVsRmP}{\ensuremath{<}0.001}  % ADDITIONAL, NOT PRE-REGISTERED: lgbm event: full vs HAR+RM: DM p
\providecommand{\rhchXLgbmEventFullVsRmShare}{}\renewcommand{\rhchXLgbmEventFullVsRmShare}{111}  % ADDITIONAL, NOT PRE-REGISTERED: lgbm event: full vs HAR+RM: % of the joint gain
\providecommand{\rhchXLgbmEventRmVsHarPct}{}\renewcommand{\rhchXLgbmEventRmVsHarPct}{6.8}  % ADDITIONAL, NOT PRE-REGISTERED: lgbm event: HAR+RM vs HAR: % QLIKE change
\providecommand{\rhchXLgbmEventRmVsHarDm}{}\renewcommand{\rhchXLgbmEventRmVsHarDm}{2.58}  % ADDITIONAL, NOT PRE-REGISTERED: lgbm event: HAR+RM vs HAR: DM
\providecommand{\rhchXLgbmEventRmVsHarP}{}\renewcommand{\rhchXLgbmEventRmVsHarP}{0.010}  % ADDITIONAL, NOT PRE-REGISTERED: lgbm event: HAR+RM vs HAR: DM p
\providecommand{\rhchXLgbmEventRmVsHarShare}{}\renewcommand{\rhchXLgbmEventRmVsHarShare}{$-$11}  % ADDITIONAL, NOT PRE-REGISTERED: lgbm event: HAR+RM vs HAR: % of the joint gain
\providecommand{\rhchXLgbmOrdFullVsRmPct}{}\renewcommand{\rhchXLgbmOrdFullVsRmPct}{$-$1.5}  % ADDITIONAL, NOT PRE-REGISTERED: lgbm ord: full vs HAR+RM: % QLIKE change
\providecommand{\rhchXLgbmOrdFullVsRmDm}{}\renewcommand{\rhchXLgbmOrdFullVsRmDm}{$-$4.30}  % ADDITIONAL, NOT PRE-REGISTERED: lgbm ord: full vs HAR+RM: DM
\providecommand{\rhchXLgbmOrdFullVsRmP}{}\renewcommand{\rhchXLgbmOrdFullVsRmP}{\ensuremath{<}0.001}  % ADDITIONAL, NOT PRE-REGISTERED: lgbm ord: full vs HAR+RM: DM p
\providecommand{\rhchXLgbmOrdFullVsRmShare}{}\renewcommand{\rhchXLgbmOrdFullVsRmShare}{23}  % ADDITIONAL, NOT PRE-REGISTERED: lgbm ord: full vs HAR+RM: % of the joint gain
\providecommand{\rhchXLgbmOrdRmVsHarPct}{}\renewcommand{\rhchXLgbmOrdRmVsHarPct}{$-$4.6}  % ADDITIONAL, NOT PRE-REGISTERED: lgbm ord: HAR+RM vs HAR: % QLIKE change
\providecommand{\rhchXLgbmOrdRmVsHarDm}{}\renewcommand{\rhchXLgbmOrdRmVsHarDm}{$-$9.55}  % ADDITIONAL, NOT PRE-REGISTERED: lgbm ord: HAR+RM vs HAR: DM
\providecommand{\rhchXLgbmOrdRmVsHarP}{}\renewcommand{\rhchXLgbmOrdRmVsHarP}{\ensuremath{<}0.001}  % ADDITIONAL, NOT PRE-REGISTERED: lgbm ord: HAR+RM vs HAR: DM p
\providecommand{\rhchXLgbmOrdRmVsHarShare}{}\renewcommand{\rhchXLgbmOrdRmVsHarShare}{77}  % ADDITIONAL, NOT PRE-REGISTERED: lgbm ord: HAR+RM vs HAR: % of the joint gain
\providecommand{\rhchXXgbBarFullVsRmPct}{}\renewcommand{\rhchXXgbBarFullVsRmPct}{$-$11.6}  % ADDITIONAL, NOT PRE-REGISTERED: xgb all: full vs HAR+RM: % QLIKE change
\providecommand{\rhchXXgbBarFullVsRmDm}{}\renewcommand{\rhchXXgbBarFullVsRmDm}{$-$7.19}  % ADDITIONAL, NOT PRE-REGISTERED: xgb all: full vs HAR+RM: DM
\providecommand{\rhchXXgbBarFullVsRmP}{}\renewcommand{\rhchXXgbBarFullVsRmP}{\ensuremath{<}0.001}  % ADDITIONAL, NOT PRE-REGISTERED: xgb all: full vs HAR+RM: DM p
\providecommand{\rhchXXgbBarFullVsRmShare}{}\renewcommand{\rhchXXgbBarFullVsRmShare}{77}  % ADDITIONAL, NOT PRE-REGISTERED: xgb all: full vs HAR+RM: % of the joint gain
\providecommand{\rhchXXgbBarRmVsHarPct}{}\renewcommand{\rhchXXgbBarRmVsHarPct}{$-$3.4}  % ADDITIONAL, NOT PRE-REGISTERED: xgb all: HAR+RM vs HAR: % QLIKE change
\providecommand{\rhchXXgbBarRmVsHarDm}{}\renewcommand{\rhchXXgbBarRmVsHarDm}{$-$7.51}  % ADDITIONAL, NOT PRE-REGISTERED: xgb all: HAR+RM vs HAR: DM
\providecommand{\rhchXXgbBarRmVsHarP}{}\renewcommand{\rhchXXgbBarRmVsHarP}{\ensuremath{<}0.001}  % ADDITIONAL, NOT PRE-REGISTERED: xgb all: HAR+RM vs HAR: DM p
\providecommand{\rhchXXgbBarRmVsHarShare}{}\renewcommand{\rhchXXgbBarRmVsHarShare}{23}  % ADDITIONAL, NOT PRE-REGISTERED: xgb all: HAR+RM vs HAR: % of the joint gain
\providecommand{\rhchXXgbEventFullVsRmPct}{}\renewcommand{\rhchXXgbEventFullVsRmPct}{$-$65.8}  % ADDITIONAL, NOT PRE-REGISTERED: xgb event: full vs HAR+RM: % QLIKE change
\providecommand{\rhchXXgbEventFullVsRmDm}{}\renewcommand{\rhchXXgbEventFullVsRmDm}{$-$6.93}  % ADDITIONAL, NOT PRE-REGISTERED: xgb event: full vs HAR+RM: DM
\providecommand{\rhchXXgbEventFullVsRmP}{}\renewcommand{\rhchXXgbEventFullVsRmP}{\ensuremath{<}0.001}  % ADDITIONAL, NOT PRE-REGISTERED: xgb event: full vs HAR+RM: DM p
\providecommand{\rhchXXgbEventFullVsRmShare}{}\renewcommand{\rhchXXgbEventFullVsRmShare}{109}  % ADDITIONAL, NOT PRE-REGISTERED: xgb event: full vs HAR+RM: % of the joint gain
\providecommand{\rhchXXgbEventRmVsHarPct}{}\renewcommand{\rhchXXgbEventRmVsHarPct}{5.5}  % ADDITIONAL, NOT PRE-REGISTERED: xgb event: HAR+RM vs HAR: % QLIKE change
\providecommand{\rhchXXgbEventRmVsHarDm}{}\renewcommand{\rhchXXgbEventRmVsHarDm}{2.46}  % ADDITIONAL, NOT PRE-REGISTERED: xgb event: HAR+RM vs HAR: DM
\providecommand{\rhchXXgbEventRmVsHarP}{}\renewcommand{\rhchXXgbEventRmVsHarP}{0.014}  % ADDITIONAL, NOT PRE-REGISTERED: xgb event: HAR+RM vs HAR: DM p
\providecommand{\rhchXXgbEventRmVsHarShare}{}\renewcommand{\rhchXXgbEventRmVsHarShare}{$-$9}  % ADDITIONAL, NOT PRE-REGISTERED: xgb event: HAR+RM vs HAR: % of the joint gain
\providecommand{\rhchXXgbOrdFullVsRmPct}{}\renewcommand{\rhchXXgbOrdFullVsRmPct}{$-$1.2}  % ADDITIONAL, NOT PRE-REGISTERED: xgb ord: full vs HAR+RM: % QLIKE change
\providecommand{\rhchXXgbOrdFullVsRmDm}{}\renewcommand{\rhchXXgbOrdFullVsRmDm}{$-$3.23}  % ADDITIONAL, NOT PRE-REGISTERED: xgb ord: full vs HAR+RM: DM
\providecommand{\rhchXXgbOrdFullVsRmP}{}\renewcommand{\rhchXXgbOrdFullVsRmP}{0.001}  % ADDITIONAL, NOT PRE-REGISTERED: xgb ord: full vs HAR+RM: DM p
\providecommand{\rhchXXgbOrdFullVsRmShare}{}\renewcommand{\rhchXXgbOrdFullVsRmShare}{19}  % ADDITIONAL, NOT PRE-REGISTERED: xgb ord: full vs HAR+RM: % of the joint gain
\providecommand{\rhchXXgbOrdRmVsHarPct}{}\renewcommand{\rhchXXgbOrdRmVsHarPct}{$-$4.9}  % ADDITIONAL, NOT PRE-REGISTERED: xgb ord: HAR+RM vs HAR: % QLIKE change
\providecommand{\rhchXXgbOrdRmVsHarDm}{}\renewcommand{\rhchXXgbOrdRmVsHarDm}{$-$10.37}  % ADDITIONAL, NOT PRE-REGISTERED: xgb ord: HAR+RM vs HAR: DM
\providecommand{\rhchXXgbOrdRmVsHarP}{}\renewcommand{\rhchXXgbOrdRmVsHarP}{\ensuremath{<}0.001}  % ADDITIONAL, NOT PRE-REGISTERED: xgb ord: HAR+RM vs HAR: DM p
\providecommand{\rhchXXgbOrdRmVsHarShare}{}\renewcommand{\rhchXXgbOrdRmVsHarShare}{81}  % ADDITIONAL, NOT PRE-REGISTERED: xgb ord: HAR+RM vs HAR: % of the joint gain
\providecommand{\rhchXRfBarFullVsRmPct}{}\renewcommand{\rhchXRfBarFullVsRmPct}{$-$8.4}  % ADDITIONAL, NOT PRE-REGISTERED: rf all: full vs HAR+RM: % QLIKE change
\providecommand{\rhchXRfBarFullVsRmDm}{}\renewcommand{\rhchXRfBarFullVsRmDm}{$-$6.48}  % ADDITIONAL, NOT PRE-REGISTERED: rf all: full vs HAR+RM: DM
\providecommand{\rhchXRfBarFullVsRmP}{}\renewcommand{\rhchXRfBarFullVsRmP}{\ensuremath{<}0.001}  % ADDITIONAL, NOT PRE-REGISTERED: rf all: full vs HAR+RM: DM p
\providecommand{\rhchXRfBarFullVsRmShare}{}\renewcommand{\rhchXRfBarFullVsRmShare}{72}  % ADDITIONAL, NOT PRE-REGISTERED: rf all: full vs HAR+RM: % of the joint gain
\providecommand{\rhchXRfBarRmVsHarPct}{}\renewcommand{\rhchXRfBarRmVsHarPct}{$-$3.2}  % ADDITIONAL, NOT PRE-REGISTERED: rf all: HAR+RM vs HAR: % QLIKE change
\providecommand{\rhchXRfBarRmVsHarDm}{}\renewcommand{\rhchXRfBarRmVsHarDm}{$-$6.79}  % ADDITIONAL, NOT PRE-REGISTERED: rf all: HAR+RM vs HAR: DM
\providecommand{\rhchXRfBarRmVsHarP}{}\renewcommand{\rhchXRfBarRmVsHarP}{\ensuremath{<}0.001}  % ADDITIONAL, NOT PRE-REGISTERED: rf all: HAR+RM vs HAR: DM p
\providecommand{\rhchXRfBarRmVsHarShare}{}\renewcommand{\rhchXRfBarRmVsHarShare}{28}  % ADDITIONAL, NOT PRE-REGISTERED: rf all: HAR+RM vs HAR: % of the joint gain
\providecommand{\rhchXRfEventFullVsRmPct}{}\renewcommand{\rhchXRfEventFullVsRmPct}{$-$47.6}  % ADDITIONAL, NOT PRE-REGISTERED: rf event: full vs HAR+RM: % QLIKE change
\providecommand{\rhchXRfEventFullVsRmDm}{}\renewcommand{\rhchXRfEventFullVsRmDm}{$-$7.12}  % ADDITIONAL, NOT PRE-REGISTERED: rf event: full vs HAR+RM: DM
\providecommand{\rhchXRfEventFullVsRmP}{}\renewcommand{\rhchXRfEventFullVsRmP}{\ensuremath{<}0.001}  % ADDITIONAL, NOT PRE-REGISTERED: rf event: full vs HAR+RM: DM p
\providecommand{\rhchXRfEventFullVsRmShare}{}\renewcommand{\rhchXRfEventFullVsRmShare}{112}  % ADDITIONAL, NOT PRE-REGISTERED: rf event: full vs HAR+RM: % of the joint gain
\providecommand{\rhchXRfEventRmVsHarPct}{}\renewcommand{\rhchXRfEventRmVsHarPct}{5.3}  % ADDITIONAL, NOT PRE-REGISTERED: rf event: HAR+RM vs HAR: % QLIKE change
\providecommand{\rhchXRfEventRmVsHarDm}{}\renewcommand{\rhchXRfEventRmVsHarDm}{1.88}  % ADDITIONAL, NOT PRE-REGISTERED: rf event: HAR+RM vs HAR: DM
\providecommand{\rhchXRfEventRmVsHarP}{}\renewcommand{\rhchXRfEventRmVsHarP}{0.061}  % ADDITIONAL, NOT PRE-REGISTERED: rf event: HAR+RM vs HAR: DM p
\providecommand{\rhchXRfEventRmVsHarShare}{}\renewcommand{\rhchXRfEventRmVsHarShare}{$-$12}  % ADDITIONAL, NOT PRE-REGISTERED: rf event: HAR+RM vs HAR: % of the joint gain
\providecommand{\rhchXRfOrdFullVsRmPct}{}\renewcommand{\rhchXRfOrdFullVsRmPct}{$-$0.7}  % ADDITIONAL, NOT PRE-REGISTERED: rf ord: full vs HAR+RM: % QLIKE change
\providecommand{\rhchXRfOrdFullVsRmDm}{}\renewcommand{\rhchXRfOrdFullVsRmDm}{$-$2.22}  % ADDITIONAL, NOT PRE-REGISTERED: rf ord: full vs HAR+RM: DM
\providecommand{\rhchXRfOrdFullVsRmP}{}\renewcommand{\rhchXRfOrdFullVsRmP}{0.026}  % ADDITIONAL, NOT PRE-REGISTERED: rf ord: full vs HAR+RM: DM p
\providecommand{\rhchXRfOrdFullVsRmShare}{}\renewcommand{\rhchXRfOrdFullVsRmShare}{12}  % ADDITIONAL, NOT PRE-REGISTERED: rf ord: full vs HAR+RM: % of the joint gain
\providecommand{\rhchXRfOrdRmVsHarPct}{}\renewcommand{\rhchXRfOrdRmVsHarPct}{$-$4.7}  % ADDITIONAL, NOT PRE-REGISTERED: rf ord: HAR+RM vs HAR: % QLIKE change
\providecommand{\rhchXRfOrdRmVsHarDm}{}\renewcommand{\rhchXRfOrdRmVsHarDm}{$-$10.68}  % ADDITIONAL, NOT PRE-REGISTERED: rf ord: HAR+RM vs HAR: DM
\providecommand{\rhchXRfOrdRmVsHarP}{}\renewcommand{\rhchXRfOrdRmVsHarP}{\ensuremath{<}0.001}  % ADDITIONAL, NOT PRE-REGISTERED: rf ord: HAR+RM vs HAR: DM p
\providecommand{\rhchXRfOrdRmVsHarShare}{}\renewcommand{\rhchXRfOrdRmVsHarShare}{88}  % ADDITIONAL, NOT PRE-REGISTERED: rf ord: HAR+RM vs HAR: % of the joint gain
\providecommand{\rhchXLinBarFullVsRmPctMin}{}\renewcommand{\rhchXLinBarFullVsRmPctMin}{$-$7.0}  % ADDITIONAL: Lin all: lowest %
\providecommand{\rhchXLinBarFullVsRmPctMax}{}\renewcommand{\rhchXLinBarFullVsRmPctMax}{$-$3.5}  % ADDITIONAL: Lin all: highest %
\providecommand{\rhchXLinBarFullVsRmNSig}{}\renewcommand{\rhchXLinBarFullVsRmNSig}{3}  % ADDITIONAL: Lin all: classes whose full set beats HAR+RM at 5%
\providecommand{\rhchXLinEventFullVsRmPctMin}{}\renewcommand{\rhchXLinEventFullVsRmPctMin}{$-$49.5}  % ADDITIONAL: Lin event: lowest %
\providecommand{\rhchXLinEventFullVsRmPctMax}{}\renewcommand{\rhchXLinEventFullVsRmPctMax}{$-$25.7}  % ADDITIONAL: Lin event: highest %
\providecommand{\rhchXLinEventFullVsRmNSig}{}\renewcommand{\rhchXLinEventFullVsRmNSig}{3}  % ADDITIONAL: Lin event: classes whose full set beats HAR+RM at 5%
\providecommand{\rhchXLinOrdFullVsRmPctMin}{}\renewcommand{\rhchXLinOrdFullVsRmPctMin}{$-$0.1}  % ADDITIONAL: Lin ord: lowest %
\providecommand{\rhchXLinOrdFullVsRmPctMax}{}\renewcommand{\rhchXLinOrdFullVsRmPctMax}{0.9}  % ADDITIONAL: Lin ord: highest %
\providecommand{\rhchXLinOrdFullVsRmNSig}{}\renewcommand{\rhchXLinOrdFullVsRmNSig}{0}  % ADDITIONAL: Lin ord: classes whose full set beats HAR+RM at 5%
\providecommand{\rhchXTreeBarFullVsRmPctMin}{}\renewcommand{\rhchXTreeBarFullVsRmPctMin}{$-$11.6}  % ADDITIONAL: Tree all: lowest %
\providecommand{\rhchXTreeBarFullVsRmPctMax}{}\renewcommand{\rhchXTreeBarFullVsRmPctMax}{$-$8.4}  % ADDITIONAL: Tree all: highest %
\providecommand{\rhchXTreeBarFullVsRmNSig}{}\renewcommand{\rhchXTreeBarFullVsRmNSig}{3}  % ADDITIONAL: Tree all: classes whose full set beats HAR+RM at 5%
\providecommand{\rhchXTreeEventFullVsRmPctMin}{}\renewcommand{\rhchXTreeEventFullVsRmPctMin}{$-$65.8}  % ADDITIONAL: Tree event: lowest %
\providecommand{\rhchXTreeEventFullVsRmPctMax}{}\renewcommand{\rhchXTreeEventFullVsRmPctMax}{$-$47.6}  % ADDITIONAL: Tree event: highest %
\providecommand{\rhchXTreeEventFullVsRmNSig}{}\renewcommand{\rhchXTreeEventFullVsRmNSig}{3}  % ADDITIONAL: Tree event: classes whose full set beats HAR+RM at 5%
\providecommand{\rhchXTreeOrdFullVsRmPctMin}{}\renewcommand{\rhchXTreeOrdFullVsRmPctMin}{$-$1.5}  % ADDITIONAL: Tree ord: lowest %
\providecommand{\rhchXTreeOrdFullVsRmPctMax}{}\renewcommand{\rhchXTreeOrdFullVsRmPctMax}{$-$0.7}  % ADDITIONAL: Tree ord: highest %
\providecommand{\rhchXTreeOrdFullVsRmNSig}{}\renewcommand{\rhchXTreeOrdFullVsRmNSig}{3}  % ADDITIONAL: Tree ord: classes whose full set beats HAR+RM at 5%
\providecommand{\rhchRidgeBarLadFullPct}{}\renewcommand{\rhchRidgeBarLadFullPct}{$-$3.6}  % ridge all: ladder k=46 vs 0 %
\providecommand{\rhchRidgeBarROne}{}\renewcommand{\rhchRidgeBarROne}{56}  % ridge all: ladder R_1 (%)
\providecommand{\rhchRidgeBarRThree}{}\renewcommand{\rhchRidgeBarRThree}{84}  % ridge all: ladder R_3 (%)
\providecommand{\rhchRidgeBarRFive}{}\renewcommand{\rhchRidgeBarRFive}{145}  % ridge all: ladder R_5 (%)
\providecommand{\rhchRidgeBarRTen}{}\renewcommand{\rhchRidgeBarRTen}{188}  % ridge all: ladder R_10 (%)
\providecommand{\rhchRidgeBarRTwenty}{}\renewcommand{\rhchRidgeBarRTwenty}{171}  % ridge all: ladder R_20 (%)
\providecommand{\rhchRidgeEventLadFullPct}{}\renewcommand{\rhchRidgeEventLadFullPct}{$-$23.2}  % ridge event: ladder k=46 vs 0 %
\providecommand{\rhchRidgeEventROne}{}\renewcommand{\rhchRidgeEventROne}{66}  % ridge event: ladder R_1 (%)
\providecommand{\rhchRidgeEventRThree}{}\renewcommand{\rhchRidgeEventRThree}{95}  % ridge event: ladder R_3 (%)
\providecommand{\rhchRidgeEventRFive}{}\renewcommand{\rhchRidgeEventRFive}{96}  % ridge event: ladder R_5 (%)
\providecommand{\rhchRidgeEventRTen}{}\renewcommand{\rhchRidgeEventRTen}{101}  % ridge event: ladder R_10 (%)
\providecommand{\rhchRidgeEventRTwenty}{}\renewcommand{\rhchRidgeEventRTwenty}{99}  % ridge event: ladder R_20 (%)
\providecommand{\rhchRidgeOrdLadFullPct}{}\renewcommand{\rhchRidgeOrdLadFullPct}{$-$0.3}  % ridge ord: ladder k=46 vs 0 %
\providecommand{\rhchRidgeOrdROne}{}\renewcommand{\rhchRidgeOrdROne}{$-$70}  % ridge ord: ladder R_1 (%)
\providecommand{\rhchRidgeOrdRThree}{}\renewcommand{\rhchRidgeOrdRThree}{$-$75}  % ridge ord: ladder R_3 (%)
\providecommand{\rhchRidgeOrdRFive}{}\renewcommand{\rhchRidgeOrdRFive}{824}  % ridge ord: ladder R_5 (%)
\providecommand{\rhchRidgeOrdRTen}{}\renewcommand{\rhchRidgeOrdRTen}{1391}  % ridge ord: ladder R_10 (%)
\providecommand{\rhchRidgeOrdRTwenty}{}\renewcommand{\rhchRidgeOrdRTwenty}{1168}  % ridge ord: ladder R_20 (%)
\providecommand{\rhchLassoBarLadFullPct}{}\renewcommand{\rhchLassoBarLadFullPct}{$-$6.0}  % lasso all: ladder k=46 vs 0 %
\providecommand{\rhchLassoBarROne}{}\renewcommand{\rhchLassoBarROne}{69}  % lasso all: ladder R_1 (%)
\providecommand{\rhchLassoBarRThree}{}\renewcommand{\rhchLassoBarRThree}{63}  % lasso all: ladder R_3 (%)
\providecommand{\rhchLassoBarRFive}{}\renewcommand{\rhchLassoBarRFive}{108}  % lasso all: ladder R_5 (%)
\providecommand{\rhchLassoBarRTen}{}\renewcommand{\rhchLassoBarRTen}{155}  % lasso all: ladder R_10 (%)
\providecommand{\rhchLassoBarRTwenty}{}\renewcommand{\rhchLassoBarRTwenty}{139}  % lasso all: ladder R_20 (%)
\providecommand{\rhchLassoEventLadFullPct}{}\renewcommand{\rhchLassoEventLadFullPct}{$-$46.9}  % lasso event: ladder k=46 vs 0 %
\providecommand{\rhchLassoEventROne}{}\renewcommand{\rhchLassoEventROne}{64}  % lasso event: ladder R_1 (%)
\providecommand{\rhchLassoEventRThree}{}\renewcommand{\rhchLassoEventRThree}{67}  % lasso event: ladder R_3 (%)
\providecommand{\rhchLassoEventRFive}{}\renewcommand{\rhchLassoEventRFive}{92}  % lasso event: ladder R_5 (%)
\providecommand{\rhchLassoEventRTen}{}\renewcommand{\rhchLassoEventRTen}{96}  % lasso event: ladder R_10 (%)
\providecommand{\rhchLassoEventRTwenty}{}\renewcommand{\rhchLassoEventRTwenty}{97}  % lasso event: ladder R_20 (%)
\providecommand{\rhchLassoOrdLadFullPct}{}\renewcommand{\rhchLassoOrdLadFullPct}{1.0}  % lasso ord: ladder k=46 vs 0 %
\providecommand{\rhchLassoOrdROne}{}\renewcommand{\rhchLassoOrdROne}{--}  % lasso ord: ladder R_1 (%)
\providecommand{\rhchLassoOrdRThree}{}\renewcommand{\rhchLassoOrdRThree}{--}  % lasso ord: ladder R_3 (%)
\providecommand{\rhchLassoOrdRFive}{}\renewcommand{\rhchLassoOrdRFive}{--}  % lasso ord: ladder R_5 (%)
\providecommand{\rhchLassoOrdRTen}{}\renewcommand{\rhchLassoOrdRTen}{--}  % lasso ord: ladder R_10 (%)
\providecommand{\rhchLassoOrdRTwenty}{}\renewcommand{\rhchLassoOrdRTwenty}{--}  % lasso ord: ladder R_20 (%)
\providecommand{\rhchEnetBarLadFullPct}{}\renewcommand{\rhchEnetBarLadFullPct}{$-$5.3}  % enet all: ladder k=46 vs 0 %
\providecommand{\rhchEnetBarROne}{}\renewcommand{\rhchEnetBarROne}{55}  % enet all: ladder R_1 (%)
\providecommand{\rhchEnetBarRThree}{}\renewcommand{\rhchEnetBarRThree}{58}  % enet all: ladder R_3 (%)
\providecommand{\rhchEnetBarRFive}{}\renewcommand{\rhchEnetBarRFive}{99}  % enet all: ladder R_5 (%)
\providecommand{\rhchEnetBarRTen}{}\renewcommand{\rhchEnetBarRTen}{130}  % enet all: ladder R_10 (%)
\providecommand{\rhchEnetBarRTwenty}{}\renewcommand{\rhchEnetBarRTwenty}{120}  % enet all: ladder R_20 (%)
\providecommand{\rhchEnetEventLadFullPct}{}\renewcommand{\rhchEnetEventLadFullPct}{$-$23.0}  % enet event: ladder k=46 vs 0 %
\providecommand{\rhchEnetEventROne}{}\renewcommand{\rhchEnetEventROne}{90}  % enet event: ladder R_1 (%)
\providecommand{\rhchEnetEventRThree}{}\renewcommand{\rhchEnetEventRThree}{99}  % enet event: ladder R_3 (%)
\providecommand{\rhchEnetEventRFive}{}\renewcommand{\rhchEnetEventRFive}{97}  % enet event: ladder R_5 (%)
\providecommand{\rhchEnetEventRTen}{}\renewcommand{\rhchEnetEventRTen}{102}  % enet event: ladder R_10 (%)
\providecommand{\rhchEnetEventRTwenty}{}\renewcommand{\rhchEnetEventRTwenty}{98}  % enet event: ladder R_20 (%)
\providecommand{\rhchEnetOrdLadFullPct}{}\renewcommand{\rhchEnetOrdLadFullPct}{$-$2.2}  % enet ord: ladder k=46 vs 0 %
\providecommand{\rhchEnetOrdROne}{}\renewcommand{\rhchEnetOrdROne}{$-$7}  % enet ord: ladder R_1 (%)
\providecommand{\rhchEnetOrdRThree}{}\renewcommand{\rhchEnetOrdRThree}{$-$16}  % enet ord: ladder R_3 (%)
\providecommand{\rhchEnetOrdRFive}{}\renewcommand{\rhchEnetOrdRFive}{102}  % enet ord: ladder R_5 (%)
\providecommand{\rhchEnetOrdRTen}{}\renewcommand{\rhchEnetOrdRTen}{180}  % enet ord: ladder R_10 (%)
\providecommand{\rhchEnetOrdRTwenty}{}\renewcommand{\rhchEnetOrdRTwenty}{159}  % enet ord: ladder R_20 (%)
\providecommand{\rhchLgbmBarLadFullPct}{}\renewcommand{\rhchLgbmBarLadFullPct}{$-$13.6}  % lgbm all: ladder k=46 vs 0 %
\providecommand{\rhchLgbmBarROne}{}\renewcommand{\rhchLgbmBarROne}{3}  % lgbm all: ladder R_1 (%)
\providecommand{\rhchLgbmBarRThree}{}\renewcommand{\rhchLgbmBarRThree}{34}  % lgbm all: ladder R_3 (%)
\providecommand{\rhchLgbmBarRFive}{}\renewcommand{\rhchLgbmBarRFive}{34}  % lgbm all: ladder R_5 (%)
\providecommand{\rhchLgbmBarRTen}{}\renewcommand{\rhchLgbmBarRTen}{95}  % lgbm all: ladder R_10 (%)
\providecommand{\rhchLgbmBarRTwenty}{}\renewcommand{\rhchLgbmBarRTwenty}{96}  % lgbm all: ladder R_20 (%)
\providecommand{\rhchLgbmEventLadFullPct}{}\renewcommand{\rhchLgbmEventLadFullPct}{$-$62.0}  % lgbm event: ladder k=46 vs 0 %
\providecommand{\rhchLgbmEventROne}{}\renewcommand{\rhchLgbmEventROne}{$-$1}  % lgbm event: ladder R_1 (%)
\providecommand{\rhchLgbmEventRThree}{}\renewcommand{\rhchLgbmEventRThree}{5}  % lgbm event: ladder R_3 (%)
\providecommand{\rhchLgbmEventRFive}{}\renewcommand{\rhchLgbmEventRFive}{3}  % lgbm event: ladder R_5 (%)
\providecommand{\rhchLgbmEventRTen}{}\renewcommand{\rhchLgbmEventRTen}{100}  % lgbm event: ladder R_10 (%)
\providecommand{\rhchLgbmEventRTwenty}{}\renewcommand{\rhchLgbmEventRTwenty}{100}  % lgbm event: ladder R_20 (%)
\providecommand{\rhchLgbmOrdLadFullPct}{}\renewcommand{\rhchLgbmOrdLadFullPct}{$-$5.5}  % lgbm ord: ladder k=46 vs 0 %
\providecommand{\rhchLgbmOrdROne}{}\renewcommand{\rhchLgbmOrdROne}{12}  % lgbm ord: ladder R_1 (%)
\providecommand{\rhchLgbmOrdRThree}{}\renewcommand{\rhchLgbmOrdRThree}{88}  % lgbm ord: ladder R_3 (%)
\providecommand{\rhchLgbmOrdRFive}{}\renewcommand{\rhchLgbmOrdRFive}{93}  % lgbm ord: ladder R_5 (%)
\providecommand{\rhchLgbmOrdRTen}{}\renewcommand{\rhchLgbmOrdRTen}{86}  % lgbm ord: ladder R_10 (%)
\providecommand{\rhchLgbmOrdRTwenty}{}\renewcommand{\rhchLgbmOrdRTwenty}{90}  % lgbm ord: ladder R_20 (%)
\providecommand{\rhchXgbBarLadFullPct}{}\renewcommand{\rhchXgbBarLadFullPct}{$-$13.6}  % xgb all: ladder k=46 vs 0 %
\providecommand{\rhchXgbBarROne}{}\renewcommand{\rhchXgbBarROne}{4}  % xgb all: ladder R_1 (%)
\providecommand{\rhchXgbBarRThree}{}\renewcommand{\rhchXgbBarRThree}{27}  % xgb all: ladder R_3 (%)
\providecommand{\rhchXgbBarRFive}{}\renewcommand{\rhchXgbBarRFive}{36}  % xgb all: ladder R_5 (%)
\providecommand{\rhchXgbBarRTen}{}\renewcommand{\rhchXgbBarRTen}{98}  % xgb all: ladder R_10 (%)
\providecommand{\rhchXgbBarRTwenty}{}\renewcommand{\rhchXgbBarRTwenty}{99}  % xgb all: ladder R_20 (%)
\providecommand{\rhchXgbEventLadFullPct}{}\renewcommand{\rhchXgbEventLadFullPct}{$-$63.5}  % xgb event: ladder k=46 vs 0 %
\providecommand{\rhchXgbEventROne}{}\renewcommand{\rhchXgbEventROne}{0}  % xgb event: ladder R_1 (%)
\providecommand{\rhchXgbEventRThree}{}\renewcommand{\rhchXgbEventRThree}{10}  % xgb event: ladder R_3 (%)
\providecommand{\rhchXgbEventRFive}{}\renewcommand{\rhchXgbEventRFive}{3}  % xgb event: ladder R_5 (%)
\providecommand{\rhchXgbEventRTen}{}\renewcommand{\rhchXgbEventRTen}{99}  % xgb event: ladder R_10 (%)
\providecommand{\rhchXgbEventRTwenty}{}\renewcommand{\rhchXgbEventRTwenty}{100}  % xgb event: ladder R_20 (%)
\providecommand{\rhchXgbOrdLadFullPct}{}\renewcommand{\rhchXgbOrdLadFullPct}{$-$5.1}  % xgb ord: ladder k=46 vs 0 %
\providecommand{\rhchXgbOrdROne}{}\renewcommand{\rhchXgbOrdROne}{13}  % xgb ord: ladder R_1 (%)
\providecommand{\rhchXgbOrdRThree}{}\renewcommand{\rhchXgbOrdRThree}{64}  % xgb ord: ladder R_3 (%)
\providecommand{\rhchXgbOrdRFive}{}\renewcommand{\rhchXgbOrdRFive}{108}  % xgb ord: ladder R_5 (%)
\providecommand{\rhchXgbOrdRTen}{}\renewcommand{\rhchXgbOrdRTen}{95}  % xgb ord: ladder R_10 (%)
\providecommand{\rhchXgbOrdRTwenty}{}\renewcommand{\rhchXgbOrdRTwenty}{99}  % xgb ord: ladder R_20 (%)
\providecommand{\rhchRfBarLadFullPct}{}\renewcommand{\rhchRfBarLadFullPct}{$-$10.6}  % rf all: ladder k=46 vs 0 %
\providecommand{\rhchRfBarROne}{}\renewcommand{\rhchRfBarROne}{6}  % rf all: ladder R_1 (%)
\providecommand{\rhchRfBarRThree}{}\renewcommand{\rhchRfBarRThree}{42}  % rf all: ladder R_3 (%)
\providecommand{\rhchRfBarRFive}{}\renewcommand{\rhchRfBarRFive}{40}  % rf all: ladder R_5 (%)
\providecommand{\rhchRfBarRTen}{}\renewcommand{\rhchRfBarRTen}{98}  % rf all: ladder R_10 (%)
\providecommand{\rhchRfBarRTwenty}{}\renewcommand{\rhchRfBarRTwenty}{101}  % rf all: ladder R_20 (%)
\providecommand{\rhchRfEventLadFullPct}{}\renewcommand{\rhchRfEventLadFullPct}{$-$44.4}  % rf event: ladder k=46 vs 0 %
\providecommand{\rhchRfEventROne}{}\renewcommand{\rhchRfEventROne}{$-$3}  % rf event: ladder R_1 (%)
\providecommand{\rhchRfEventRThree}{}\renewcommand{\rhchRfEventRThree}{3}  % rf event: ladder R_3 (%)
\providecommand{\rhchRfEventRFive}{}\renewcommand{\rhchRfEventRFive}{6}  % rf event: ladder R_5 (%)
\providecommand{\rhchRfEventRTen}{}\renewcommand{\rhchRfEventRTen}{100}  % rf event: ladder R_10 (%)
\providecommand{\rhchRfEventRTwenty}{}\renewcommand{\rhchRfEventRTwenty}{100}  % rf event: ladder R_20 (%)
\providecommand{\rhchRfOrdLadFullPct}{}\renewcommand{\rhchRfOrdLadFullPct}{$-$4.6}  % rf ord: ladder k=46 vs 0 %
\providecommand{\rhchRfOrdROne}{}\renewcommand{\rhchRfOrdROne}{20}  % rf ord: ladder R_1 (%)
\providecommand{\rhchRfOrdRThree}{}\renewcommand{\rhchRfOrdRThree}{109}  % rf ord: ladder R_3 (%)
\providecommand{\rhchRfOrdRFive}{}\renewcommand{\rhchRfOrdRFive}{97}  % rf ord: ladder R_5 (%)
\providecommand{\rhchRfOrdRTen}{}\renewcommand{\rhchRfOrdRTen}{95}  % rf ord: ladder R_10 (%)
\providecommand{\rhchRfOrdRTwenty}{}\renewcommand{\rhchRfOrdRTwenty}{102}  % rf ord: ladder R_20 (%)
\fi
\ifdefined\rhfaReading\else% AUTO-GENERATED by writeup/make_rehaul_factor.py -- do not edit; re-run the script.
% \rhfa... macros: factor-augmented HAR (diffusion index) on Table 1's design (results/channels_2026-10-26/
% RULE.txt section 4; results/channels_2026-10-26/factor).  Spans: Ord = ORDINARY days (primary), All = every
% bar, Ev = the 81 FOMC statement days.  Models: Pca<k>, Pls<k>, Full (the span's L_full model).
\providecommand{\rhfaPcaOneOrdPct}{}\renewcommand{\rhfaPcaOneOrdPct}{0.08}  % pca1 ordinary: QLIKE change vs OLS-HAR, %
\providecommand{\rhfaPcaOneOrdMsePct}{}\renewcommand{\rhfaPcaOneOrdMsePct}{3.28}  % pca1 ordinary: MSE change vs OLS-HAR, %
\providecommand{\rhfaPcaOneOrdDm}{}\renewcommand{\rhfaPcaOneOrdDm}{0.90}  % pca1 ordinary: DM vs OLS-HAR (QLIKE)
\providecommand{\rhfaPcaOneOrdP}{}\renewcommand{\rhfaPcaOneOrdP}{0.369}  % pca1 ordinary: p
\providecommand{\rhfaPcaOneOrdDmFull}{}\renewcommand{\rhfaPcaOneOrdDmFull}{4.89}  % pca1 ordinary: DM vs the full set (enet)
\providecommand{\rhfaPcaOneOrdPFull}{}\renewcommand{\rhfaPcaOneOrdPFull}{$<$0.001}  % pca1 ordinary: p
\providecommand{\rhfaPcaOneOrdR}{}\renewcommand{\rhfaPcaOneOrdR}{$-$0.02}  % pca1 ordinary: R_k
\providecommand{\rhfaPcaTwoOrdPct}{}\renewcommand{\rhfaPcaTwoOrdPct}{0.06}  % pca2 ordinary: QLIKE change vs OLS-HAR, %
\providecommand{\rhfaPcaTwoOrdMsePct}{}\renewcommand{\rhfaPcaTwoOrdMsePct}{1.85}  % pca2 ordinary: MSE change vs OLS-HAR, %
\providecommand{\rhfaPcaTwoOrdDm}{}\renewcommand{\rhfaPcaTwoOrdDm}{0.32}  % pca2 ordinary: DM vs OLS-HAR (QLIKE)
\providecommand{\rhfaPcaTwoOrdP}{}\renewcommand{\rhfaPcaTwoOrdP}{0.748}  % pca2 ordinary: p
\providecommand{\rhfaPcaTwoOrdDmFull}{}\renewcommand{\rhfaPcaTwoOrdDmFull}{4.86}  % pca2 ordinary: DM vs the full set (enet)
\providecommand{\rhfaPcaTwoOrdPFull}{}\renewcommand{\rhfaPcaTwoOrdPFull}{$<$0.001}  % pca2 ordinary: p
\providecommand{\rhfaPcaTwoOrdR}{}\renewcommand{\rhfaPcaTwoOrdR}{$-$0.02}  % pca2 ordinary: R_k
\providecommand{\rhfaPcaThreeOrdPct}{}\renewcommand{\rhfaPcaThreeOrdPct}{0.45}  % pca3 ordinary: QLIKE change vs OLS-HAR, %
\providecommand{\rhfaPcaThreeOrdMsePct}{}\renewcommand{\rhfaPcaThreeOrdMsePct}{5.76}  % pca3 ordinary: MSE change vs OLS-HAR, %
\providecommand{\rhfaPcaThreeOrdDm}{}\renewcommand{\rhfaPcaThreeOrdDm}{1.44}  % pca3 ordinary: DM vs OLS-HAR (QLIKE)
\providecommand{\rhfaPcaThreeOrdP}{}\renewcommand{\rhfaPcaThreeOrdP}{0.151}  % pca3 ordinary: p
\providecommand{\rhfaPcaThreeOrdDmFull}{}\renewcommand{\rhfaPcaThreeOrdDmFull}{5.33}  % pca3 ordinary: DM vs the full set (enet)
\providecommand{\rhfaPcaThreeOrdPFull}{}\renewcommand{\rhfaPcaThreeOrdPFull}{$<$0.001}  % pca3 ordinary: p
\providecommand{\rhfaPcaThreeOrdR}{}\renewcommand{\rhfaPcaThreeOrdR}{$-$0.11}  % pca3 ordinary: R_k
\providecommand{\rhfaPcaFiveOrdPct}{}\renewcommand{\rhfaPcaFiveOrdPct}{0.93}  % pca5 ordinary: QLIKE change vs OLS-HAR, %
\providecommand{\rhfaPcaFiveOrdMsePct}{}\renewcommand{\rhfaPcaFiveOrdMsePct}{3.77}  % pca5 ordinary: MSE change vs OLS-HAR, %
\providecommand{\rhfaPcaFiveOrdDm}{}\renewcommand{\rhfaPcaFiveOrdDm}{1.85}  % pca5 ordinary: DM vs OLS-HAR (QLIKE)
\providecommand{\rhfaPcaFiveOrdP}{}\renewcommand{\rhfaPcaFiveOrdP}{0.064}  % pca5 ordinary: p
\providecommand{\rhfaPcaFiveOrdDmFull}{}\renewcommand{\rhfaPcaFiveOrdDmFull}{5.94}  % pca5 ordinary: DM vs the full set (enet)
\providecommand{\rhfaPcaFiveOrdPFull}{}\renewcommand{\rhfaPcaFiveOrdPFull}{$<$0.001}  % pca5 ordinary: p
\providecommand{\rhfaPcaFiveOrdR}{}\renewcommand{\rhfaPcaFiveOrdR}{$-$0.23}  % pca5 ordinary: R_k
\providecommand{\rhfaPcaTenOrdPct}{}\renewcommand{\rhfaPcaTenOrdPct}{$-$1.06}  % pca10 ordinary: QLIKE change vs OLS-HAR, %
\providecommand{\rhfaPcaTenOrdMsePct}{}\renewcommand{\rhfaPcaTenOrdMsePct}{21.39}  % pca10 ordinary: MSE change vs OLS-HAR, %
\providecommand{\rhfaPcaTenOrdDm}{}\renewcommand{\rhfaPcaTenOrdDm}{$-$2.19}  % pca10 ordinary: DM vs OLS-HAR (QLIKE)
\providecommand{\rhfaPcaTenOrdP}{}\renewcommand{\rhfaPcaTenOrdP}{0.028}  % pca10 ordinary: p
\providecommand{\rhfaPcaTenOrdDmFull}{}\renewcommand{\rhfaPcaTenOrdDmFull}{4.45}  % pca10 ordinary: DM vs the full set (enet)
\providecommand{\rhfaPcaTenOrdPFull}{}\renewcommand{\rhfaPcaTenOrdPFull}{$<$0.001}  % pca10 ordinary: p
\providecommand{\rhfaPcaTenOrdR}{}\renewcommand{\rhfaPcaTenOrdR}{0.27}  % pca10 ordinary: R_k
\providecommand{\rhfaPcaTwentyOrdPct}{}\renewcommand{\rhfaPcaTwentyOrdPct}{$-$1.09}  % pca20 ordinary: QLIKE change vs OLS-HAR, %
\providecommand{\rhfaPcaTwentyOrdMsePct}{}\renewcommand{\rhfaPcaTwentyOrdMsePct}{45.47}  % pca20 ordinary: MSE change vs OLS-HAR, %
\providecommand{\rhfaPcaTwentyOrdDm}{}\renewcommand{\rhfaPcaTwentyOrdDm}{$-$1.53}  % pca20 ordinary: DM vs OLS-HAR (QLIKE)
\providecommand{\rhfaPcaTwentyOrdP}{}\renewcommand{\rhfaPcaTwentyOrdP}{0.126}  % pca20 ordinary: p
\providecommand{\rhfaPcaTwentyOrdDmFull}{}\renewcommand{\rhfaPcaTwentyOrdDmFull}{3.66}  % pca20 ordinary: DM vs the full set (enet)
\providecommand{\rhfaPcaTwentyOrdPFull}{}\renewcommand{\rhfaPcaTwentyOrdPFull}{$<$0.001}  % pca20 ordinary: p
\providecommand{\rhfaPcaTwentyOrdR}{}\renewcommand{\rhfaPcaTwentyOrdR}{0.27}  % pca20 ordinary: R_k
\providecommand{\rhfaPlsOneOrdPct}{}\renewcommand{\rhfaPlsOneOrdPct}{0.01}  % pls1 ordinary: QLIKE change vs OLS-HAR, %
\providecommand{\rhfaPlsOneOrdMsePct}{}\renewcommand{\rhfaPlsOneOrdMsePct}{14.31}  % pls1 ordinary: MSE change vs OLS-HAR, %
\providecommand{\rhfaPlsOneOrdDm}{}\renewcommand{\rhfaPlsOneOrdDm}{0.05}  % pls1 ordinary: DM vs OLS-HAR (QLIKE)
\providecommand{\rhfaPlsOneOrdP}{}\renewcommand{\rhfaPlsOneOrdP}{0.962}  % pls1 ordinary: p
\providecommand{\rhfaPlsOneOrdDmFull}{}\renewcommand{\rhfaPlsOneOrdDmFull}{4.95}  % pls1 ordinary: DM vs the full set (enet)
\providecommand{\rhfaPlsOneOrdPFull}{}\renewcommand{\rhfaPlsOneOrdPFull}{$<$0.001}  % pls1 ordinary: p
\providecommand{\rhfaPlsOneOrdR}{}\renewcommand{\rhfaPlsOneOrdR}{0.00}  % pls1 ordinary: R_k
\providecommand{\rhfaPlsTwoOrdPct}{}\renewcommand{\rhfaPlsTwoOrdPct}{1.57}  % pls2 ordinary: QLIKE change vs OLS-HAR, %
\providecommand{\rhfaPlsTwoOrdMsePct}{}\renewcommand{\rhfaPlsTwoOrdMsePct}{24.30}  % pls2 ordinary: MSE change vs OLS-HAR, %
\providecommand{\rhfaPlsTwoOrdDm}{}\renewcommand{\rhfaPlsTwoOrdDm}{1.29}  % pls2 ordinary: DM vs OLS-HAR (QLIKE)
\providecommand{\rhfaPlsTwoOrdP}{}\renewcommand{\rhfaPlsTwoOrdP}{0.196}  % pls2 ordinary: p
\providecommand{\rhfaPlsTwoOrdDmFull}{}\renewcommand{\rhfaPlsTwoOrdDmFull}{4.51}  % pls2 ordinary: DM vs the full set (enet)
\providecommand{\rhfaPlsTwoOrdPFull}{}\renewcommand{\rhfaPlsTwoOrdPFull}{$<$0.001}  % pls2 ordinary: p
\providecommand{\rhfaPlsTwoOrdR}{}\renewcommand{\rhfaPlsTwoOrdR}{$-$0.40}  % pls2 ordinary: R_k
\providecommand{\rhfaPlsThreeOrdPct}{}\renewcommand{\rhfaPlsThreeOrdPct}{0.04}  % pls3 ordinary: QLIKE change vs OLS-HAR, %
\providecommand{\rhfaPlsThreeOrdMsePct}{}\renewcommand{\rhfaPlsThreeOrdMsePct}{99.04}  % pls3 ordinary: MSE change vs OLS-HAR, %
\providecommand{\rhfaPlsThreeOrdDm}{}\renewcommand{\rhfaPlsThreeOrdDm}{0.03}  % pls3 ordinary: DM vs OLS-HAR (QLIKE)
\providecommand{\rhfaPlsThreeOrdP}{}\renewcommand{\rhfaPlsThreeOrdP}{0.974}  % pls3 ordinary: p
\providecommand{\rhfaPlsThreeOrdDmFull}{}\renewcommand{\rhfaPlsThreeOrdDmFull}{3.84}  % pls3 ordinary: DM vs the full set (enet)
\providecommand{\rhfaPlsThreeOrdPFull}{}\renewcommand{\rhfaPlsThreeOrdPFull}{$<$0.001}  % pls3 ordinary: p
\providecommand{\rhfaPlsThreeOrdR}{}\renewcommand{\rhfaPlsThreeOrdR}{$-$0.01}  % pls3 ordinary: R_k
\providecommand{\rhfaFullOrdModel}{}\renewcommand{\rhfaFullOrdModel}{elastic net}  % ordinary: L_full model (lower QLIKE of ridge, elastic net)
\providecommand{\rhfaFullOrdPct}{}\renewcommand{\rhfaFullOrdPct}{$-$3.97}  % ordinary: full set QLIKE change vs OLS-HAR, %
\providecommand{\rhfaFullOrdMsePct}{}\renewcommand{\rhfaFullOrdMsePct}{128.86}  % ordinary: full set MSE change vs OLS-HAR, %
\providecommand{\rhfaFullOrdDm}{}\renewcommand{\rhfaFullOrdDm}{$-$4.73}  % ordinary: full set DM vs OLS-HAR
\providecommand{\rhfaFullOrdP}{}\renewcommand{\rhfaFullOrdP}{$<$0.001}  % ordinary: p
\providecommand{\rhfaNOrdDays}{}\renewcommand{\rhfaNOrdDays}{2{,}498}  % ordinary: days
\providecommand{\rhfaNOrdBars}{}\renewcommand{\rhfaNOrdBars}{32{,}474}  % ordinary: bars
\providecommand{\rhfaPcaOneAllPct}{}\renewcommand{\rhfaPcaOneAllPct}{$-$0.05}  % pca1 all: QLIKE change vs OLS-HAR, %
\providecommand{\rhfaPcaOneAllMsePct}{}\renewcommand{\rhfaPcaOneAllMsePct}{3.00}  % pca1 all: MSE change vs OLS-HAR, %
\providecommand{\rhfaPcaOneAllDm}{}\renewcommand{\rhfaPcaOneAllDm}{$-$0.53}  % pca1 all: DM vs OLS-HAR (QLIKE)
\providecommand{\rhfaPcaOneAllP}{}\renewcommand{\rhfaPcaOneAllP}{0.595}  % pca1 all: p
\providecommand{\rhfaPcaOneAllDmFull}{}\renewcommand{\rhfaPcaOneAllDmFull}{7.44}  % pca1 all: DM vs the full set (enet)
\providecommand{\rhfaPcaOneAllPFull}{}\renewcommand{\rhfaPcaOneAllPFull}{$<$0.001}  % pca1 all: p
\providecommand{\rhfaPcaOneAllR}{}\renewcommand{\rhfaPcaOneAllR}{0.01}  % pca1 all: R_k
\providecommand{\rhfaPcaTwoAllPct}{}\renewcommand{\rhfaPcaTwoAllPct}{0.11}  % pca2 all: QLIKE change vs OLS-HAR, %
\providecommand{\rhfaPcaTwoAllMsePct}{}\renewcommand{\rhfaPcaTwoAllMsePct}{1.71}  % pca2 all: MSE change vs OLS-HAR, %
\providecommand{\rhfaPcaTwoAllDm}{}\renewcommand{\rhfaPcaTwoAllDm}{0.55}  % pca2 all: DM vs OLS-HAR (QLIKE)
\providecommand{\rhfaPcaTwoAllP}{}\renewcommand{\rhfaPcaTwoAllP}{0.579}  % pca2 all: p
\providecommand{\rhfaPcaTwoAllDmFull}{}\renewcommand{\rhfaPcaTwoAllDmFull}{7.48}  % pca2 all: DM vs the full set (enet)
\providecommand{\rhfaPcaTwoAllPFull}{}\renewcommand{\rhfaPcaTwoAllPFull}{$<$0.001}  % pca2 all: p
\providecommand{\rhfaPcaTwoAllR}{}\renewcommand{\rhfaPcaTwoAllR}{$-$0.02}  % pca2 all: R_k
\providecommand{\rhfaPcaThreeAllPct}{}\renewcommand{\rhfaPcaThreeAllPct}{0.69}  % pca3 all: QLIKE change vs OLS-HAR, %
\providecommand{\rhfaPcaThreeAllMsePct}{}\renewcommand{\rhfaPcaThreeAllMsePct}{5.30}  % pca3 all: MSE change vs OLS-HAR, %
\providecommand{\rhfaPcaThreeAllDm}{}\renewcommand{\rhfaPcaThreeAllDm}{2.24}  % pca3 all: DM vs OLS-HAR (QLIKE)
\providecommand{\rhfaPcaThreeAllP}{}\renewcommand{\rhfaPcaThreeAllP}{0.025}  % pca3 all: p
\providecommand{\rhfaPcaThreeAllDmFull}{}\renewcommand{\rhfaPcaThreeAllDmFull}{7.96}  % pca3 all: DM vs the full set (enet)
\providecommand{\rhfaPcaThreeAllPFull}{}\renewcommand{\rhfaPcaThreeAllPFull}{$<$0.001}  % pca3 all: p
\providecommand{\rhfaPcaThreeAllR}{}\renewcommand{\rhfaPcaThreeAllR}{$-$0.10}  % pca3 all: R_k
\providecommand{\rhfaPcaFiveAllPct}{}\renewcommand{\rhfaPcaFiveAllPct}{0.94}  % pca5 all: QLIKE change vs OLS-HAR, %
\providecommand{\rhfaPcaFiveAllMsePct}{}\renewcommand{\rhfaPcaFiveAllMsePct}{3.39}  % pca5 all: MSE change vs OLS-HAR, %
\providecommand{\rhfaPcaFiveAllDm}{}\renewcommand{\rhfaPcaFiveAllDm}{2.05}  % pca5 all: DM vs OLS-HAR (QLIKE)
\providecommand{\rhfaPcaFiveAllP}{}\renewcommand{\rhfaPcaFiveAllP}{0.041}  % pca5 all: p
\providecommand{\rhfaPcaFiveAllDmFull}{}\renewcommand{\rhfaPcaFiveAllDmFull}{8.48}  % pca5 all: DM vs the full set (enet)
\providecommand{\rhfaPcaFiveAllPFull}{}\renewcommand{\rhfaPcaFiveAllPFull}{$<$0.001}  % pca5 all: p
\providecommand{\rhfaPcaFiveAllR}{}\renewcommand{\rhfaPcaFiveAllR}{$-$0.14}  % pca5 all: R_k
\providecommand{\rhfaPcaTenAllPct}{}\renewcommand{\rhfaPcaTenAllPct}{$-$0.66}  % pca10 all: QLIKE change vs OLS-HAR, %
\providecommand{\rhfaPcaTenAllMsePct}{}\renewcommand{\rhfaPcaTenAllMsePct}{19.64}  % pca10 all: MSE change vs OLS-HAR, %
\providecommand{\rhfaPcaTenAllDm}{}\renewcommand{\rhfaPcaTenAllDm}{$-$1.43}  % pca10 all: DM vs OLS-HAR (QLIKE)
\providecommand{\rhfaPcaTenAllP}{}\renewcommand{\rhfaPcaTenAllP}{0.152}  % pca10 all: p
\providecommand{\rhfaPcaTenAllDmFull}{}\renewcommand{\rhfaPcaTenAllDmFull}{7.31}  % pca10 all: DM vs the full set (enet)
\providecommand{\rhfaPcaTenAllPFull}{}\renewcommand{\rhfaPcaTenAllPFull}{$<$0.001}  % pca10 all: p
\providecommand{\rhfaPcaTenAllR}{}\renewcommand{\rhfaPcaTenAllR}{0.10}  % pca10 all: R_k
\providecommand{\rhfaPcaTwentyAllPct}{}\renewcommand{\rhfaPcaTwentyAllPct}{$-$0.87}  % pca20 all: QLIKE change vs OLS-HAR, %
\providecommand{\rhfaPcaTwentyAllMsePct}{}\renewcommand{\rhfaPcaTwentyAllMsePct}{41.73}  % pca20 all: MSE change vs OLS-HAR, %
\providecommand{\rhfaPcaTwentyAllDm}{}\renewcommand{\rhfaPcaTwentyAllDm}{$-$1.28}  % pca20 all: DM vs OLS-HAR (QLIKE)
\providecommand{\rhfaPcaTwentyAllP}{}\renewcommand{\rhfaPcaTwentyAllP}{0.201}  % pca20 all: p
\providecommand{\rhfaPcaTwentyAllDmFull}{}\renewcommand{\rhfaPcaTwentyAllDmFull}{6.40}  % pca20 all: DM vs the full set (enet)
\providecommand{\rhfaPcaTwentyAllPFull}{}\renewcommand{\rhfaPcaTwentyAllPFull}{$<$0.001}  % pca20 all: p
\providecommand{\rhfaPcaTwentyAllR}{}\renewcommand{\rhfaPcaTwentyAllR}{0.13}  % pca20 all: R_k
\providecommand{\rhfaPlsOneAllPct}{}\renewcommand{\rhfaPlsOneAllPct}{$-$0.23}  % pls1 all: QLIKE change vs OLS-HAR, %
\providecommand{\rhfaPlsOneAllMsePct}{}\renewcommand{\rhfaPlsOneAllMsePct}{13.12}  % pls1 all: MSE change vs OLS-HAR, %
\providecommand{\rhfaPlsOneAllDm}{}\renewcommand{\rhfaPlsOneAllDm}{$-$1.07}  % pls1 all: DM vs OLS-HAR (QLIKE)
\providecommand{\rhfaPlsOneAllP}{}\renewcommand{\rhfaPlsOneAllP}{0.284}  % pls1 all: p
\providecommand{\rhfaPlsOneAllDmFull}{}\renewcommand{\rhfaPlsOneAllDmFull}{7.49}  % pls1 all: DM vs the full set (enet)
\providecommand{\rhfaPlsOneAllPFull}{}\renewcommand{\rhfaPlsOneAllPFull}{$<$0.001}  % pls1 all: p
\providecommand{\rhfaPlsOneAllR}{}\renewcommand{\rhfaPlsOneAllR}{0.03}  % pls1 all: R_k
\providecommand{\rhfaPlsTwoAllPct}{}\renewcommand{\rhfaPlsTwoAllPct}{1.66}  % pls2 all: QLIKE change vs OLS-HAR, %
\providecommand{\rhfaPlsTwoAllMsePct}{}\renewcommand{\rhfaPlsTwoAllMsePct}{22.30}  % pls2 all: MSE change vs OLS-HAR, %
\providecommand{\rhfaPlsTwoAllDm}{}\renewcommand{\rhfaPlsTwoAllDm}{1.55}  % pls2 all: DM vs OLS-HAR (QLIKE)
\providecommand{\rhfaPlsTwoAllP}{}\renewcommand{\rhfaPlsTwoAllP}{0.120}  % pls2 all: p
\providecommand{\rhfaPlsTwoAllDmFull}{}\renewcommand{\rhfaPlsTwoAllDmFull}{6.99}  % pls2 all: DM vs the full set (enet)
\providecommand{\rhfaPlsTwoAllPFull}{}\renewcommand{\rhfaPlsTwoAllPFull}{$<$0.001}  % pls2 all: p
\providecommand{\rhfaPlsTwoAllR}{}\renewcommand{\rhfaPlsTwoAllR}{$-$0.25}  % pls2 all: R_k
\providecommand{\rhfaPlsThreeAllPct}{}\renewcommand{\rhfaPlsThreeAllPct}{0.01}  % pls3 all: QLIKE change vs OLS-HAR, %
\providecommand{\rhfaPlsThreeAllMsePct}{}\renewcommand{\rhfaPlsThreeAllMsePct}{90.89}  % pls3 all: MSE change vs OLS-HAR, %
\providecommand{\rhfaPlsThreeAllDm}{}\renewcommand{\rhfaPlsThreeAllDm}{0.01}  % pls3 all: DM vs OLS-HAR (QLIKE)
\providecommand{\rhfaPlsThreeAllP}{}\renewcommand{\rhfaPlsThreeAllP}{0.989}  % pls3 all: p
\providecommand{\rhfaPlsThreeAllDmFull}{}\renewcommand{\rhfaPlsThreeAllDmFull}{6.39}  % pls3 all: DM vs the full set (enet)
\providecommand{\rhfaPlsThreeAllPFull}{}\renewcommand{\rhfaPlsThreeAllPFull}{$<$0.001}  % pls3 all: p
\providecommand{\rhfaPlsThreeAllR}{}\renewcommand{\rhfaPlsThreeAllR}{0.00}  % pls3 all: R_k
\providecommand{\rhfaFullAllModel}{}\renewcommand{\rhfaFullAllModel}{elastic net}  % all: L_full model (lower QLIKE of ridge, elastic net)
\providecommand{\rhfaFullAllPct}{}\renewcommand{\rhfaFullAllPct}{$-$6.69}  % all: full set QLIKE change vs OLS-HAR, %
\providecommand{\rhfaFullAllMsePct}{}\renewcommand{\rhfaFullAllMsePct}{117.67}  % all: full set MSE change vs OLS-HAR, %
\providecommand{\rhfaFullAllDm}{}\renewcommand{\rhfaFullAllDm}{$-$7.36}  % all: full set DM vs OLS-HAR
\providecommand{\rhfaFullAllP}{}\renewcommand{\rhfaFullAllP}{$<$0.001}  % all: p
\providecommand{\rhfaNAllDays}{}\renewcommand{\rhfaNAllDays}{2{,}579}  % all: days
\providecommand{\rhfaNAllBars}{}\renewcommand{\rhfaNAllBars}{33{,}527}  % all: bars
\providecommand{\rhfaPcaOneEvPct}{}\renewcommand{\rhfaPcaOneEvPct}{$-$0.82}  % pca1 event: QLIKE change vs OLS-HAR, %
\providecommand{\rhfaPcaOneEvMsePct}{}\renewcommand{\rhfaPcaOneEvMsePct}{$-$0.06}  % pca1 event: MSE change vs OLS-HAR, %
\providecommand{\rhfaPcaOneEvDm}{}\renewcommand{\rhfaPcaOneEvDm}{$-$1.95}  % pca1 event: DM vs OLS-HAR (QLIKE)
\providecommand{\rhfaPcaOneEvP}{}\renewcommand{\rhfaPcaOneEvP}{0.051}  % pca1 event: p
\providecommand{\rhfaPcaOneEvDmFull}{}\renewcommand{\rhfaPcaOneEvDmFull}{6.60}  % pca1 event: DM vs the full set (enet)
\providecommand{\rhfaPcaOneEvPFull}{}\renewcommand{\rhfaPcaOneEvPFull}{$<$0.001}  % pca1 event: p
\providecommand{\rhfaPcaOneEvR}{}\renewcommand{\rhfaPcaOneEvR}{0.04}  % pca1 event: R_k
\providecommand{\rhfaPcaTwoEvPct}{}\renewcommand{\rhfaPcaTwoEvPct}{0.39}  % pca2 event: QLIKE change vs OLS-HAR, %
\providecommand{\rhfaPcaTwoEvMsePct}{}\renewcommand{\rhfaPcaTwoEvMsePct}{0.12}  % pca2 event: MSE change vs OLS-HAR, %
\providecommand{\rhfaPcaTwoEvDm}{}\renewcommand{\rhfaPcaTwoEvDm}{0.36}  % pca2 event: DM vs OLS-HAR (QLIKE)
\providecommand{\rhfaPcaTwoEvP}{}\renewcommand{\rhfaPcaTwoEvP}{0.715}  % pca2 event: p
\providecommand{\rhfaPcaTwoEvDmFull}{}\renewcommand{\rhfaPcaTwoEvDmFull}{6.20}  % pca2 event: DM vs the full set (enet)
\providecommand{\rhfaPcaTwoEvPFull}{}\renewcommand{\rhfaPcaTwoEvPFull}{$<$0.001}  % pca2 event: p
\providecommand{\rhfaPcaTwoEvR}{}\renewcommand{\rhfaPcaTwoEvR}{$-$0.02}  % pca2 event: R_k
\providecommand{\rhfaPcaThreeEvPct}{}\renewcommand{\rhfaPcaThreeEvPct}{2.06}  % pca3 event: QLIKE change vs OLS-HAR, %
\providecommand{\rhfaPcaThreeEvMsePct}{}\renewcommand{\rhfaPcaThreeEvMsePct}{0.11}  % pca3 event: MSE change vs OLS-HAR, %
\providecommand{\rhfaPcaThreeEvDm}{}\renewcommand{\rhfaPcaThreeEvDm}{2.48}  % pca3 event: DM vs OLS-HAR (QLIKE)
\providecommand{\rhfaPcaThreeEvP}{}\renewcommand{\rhfaPcaThreeEvP}{0.013}  % pca3 event: p
\providecommand{\rhfaPcaThreeEvDmFull}{}\renewcommand{\rhfaPcaThreeEvDmFull}{6.74}  % pca3 event: DM vs the full set (enet)
\providecommand{\rhfaPcaThreeEvPFull}{}\renewcommand{\rhfaPcaThreeEvPFull}{$<$0.001}  % pca3 event: p
\providecommand{\rhfaPcaThreeEvR}{}\renewcommand{\rhfaPcaThreeEvR}{$-$0.09}  % pca3 event: R_k
\providecommand{\rhfaPcaFiveEvPct}{}\renewcommand{\rhfaPcaFiveEvPct}{1.05}  % pca5 event: QLIKE change vs OLS-HAR, %
\providecommand{\rhfaPcaFiveEvMsePct}{}\renewcommand{\rhfaPcaFiveEvMsePct}{$-$0.86}  % pca5 event: MSE change vs OLS-HAR, %
\providecommand{\rhfaPcaFiveEvDm}{}\renewcommand{\rhfaPcaFiveEvDm}{0.75}  % pca5 event: DM vs OLS-HAR (QLIKE)
\providecommand{\rhfaPcaFiveEvP}{}\renewcommand{\rhfaPcaFiveEvP}{0.453}  % pca5 event: p
\providecommand{\rhfaPcaFiveEvDmFull}{}\renewcommand{\rhfaPcaFiveEvDmFull}{6.70}  % pca5 event: DM vs the full set (enet)
\providecommand{\rhfaPcaFiveEvPFull}{}\renewcommand{\rhfaPcaFiveEvPFull}{$<$0.001}  % pca5 event: p
\providecommand{\rhfaPcaFiveEvR}{}\renewcommand{\rhfaPcaFiveEvR}{$-$0.05}  % pca5 event: R_k
\providecommand{\rhfaPcaTenEvPct}{}\renewcommand{\rhfaPcaTenEvPct}{1.72}  % pca10 event: QLIKE change vs OLS-HAR, %
\providecommand{\rhfaPcaTenEvMsePct}{}\renewcommand{\rhfaPcaTenEvMsePct}{$-$0.07}  % pca10 event: MSE change vs OLS-HAR, %
\providecommand{\rhfaPcaTenEvDm}{}\renewcommand{\rhfaPcaTenEvDm}{1.03}  % pca10 event: DM vs OLS-HAR (QLIKE)
\providecommand{\rhfaPcaTenEvP}{}\renewcommand{\rhfaPcaTenEvP}{0.302}  % pca10 event: p
\providecommand{\rhfaPcaTenEvDmFull}{}\renewcommand{\rhfaPcaTenEvDmFull}{5.52}  % pca10 event: DM vs the full set (enet)
\providecommand{\rhfaPcaTenEvPFull}{}\renewcommand{\rhfaPcaTenEvPFull}{$<$0.001}  % pca10 event: p
\providecommand{\rhfaPcaTenEvR}{}\renewcommand{\rhfaPcaTenEvR}{$-$0.08}  % pca10 event: R_k
\providecommand{\rhfaPcaTwentyEvPct}{}\renewcommand{\rhfaPcaTwentyEvPct}{0.39}  % pca20 event: QLIKE change vs OLS-HAR, %
\providecommand{\rhfaPcaTwentyEvMsePct}{}\renewcommand{\rhfaPcaTwentyEvMsePct}{$-$0.48}  % pca20 event: MSE change vs OLS-HAR, %
\providecommand{\rhfaPcaTwentyEvDm}{}\renewcommand{\rhfaPcaTwentyEvDm}{0.30}  % pca20 event: DM vs OLS-HAR (QLIKE)
\providecommand{\rhfaPcaTwentyEvP}{}\renewcommand{\rhfaPcaTwentyEvP}{0.762}  % pca20 event: p
\providecommand{\rhfaPcaTwentyEvDmFull}{}\renewcommand{\rhfaPcaTwentyEvDmFull}{5.91}  % pca20 event: DM vs the full set (enet)
\providecommand{\rhfaPcaTwentyEvPFull}{}\renewcommand{\rhfaPcaTwentyEvPFull}{$<$0.001}  % pca20 event: p
\providecommand{\rhfaPcaTwentyEvR}{}\renewcommand{\rhfaPcaTwentyEvR}{$-$0.02}  % pca20 event: R_k
\providecommand{\rhfaPlsOneEvPct}{}\renewcommand{\rhfaPlsOneEvPct}{$-$1.66}  % pls1 event: QLIKE change vs OLS-HAR, %
\providecommand{\rhfaPlsOneEvMsePct}{}\renewcommand{\rhfaPlsOneEvMsePct}{$-$0.32}  % pls1 event: MSE change vs OLS-HAR, %
\providecommand{\rhfaPlsOneEvDm}{}\renewcommand{\rhfaPlsOneEvDm}{$-$2.59}  % pls1 event: DM vs OLS-HAR (QLIKE)
\providecommand{\rhfaPlsOneEvP}{}\renewcommand{\rhfaPlsOneEvP}{0.010}  % pls1 event: p
\providecommand{\rhfaPlsOneEvDmFull}{}\renewcommand{\rhfaPlsOneEvDmFull}{6.61}  % pls1 event: DM vs the full set (enet)
\providecommand{\rhfaPlsOneEvPFull}{}\renewcommand{\rhfaPlsOneEvPFull}{$<$0.001}  % pls1 event: p
\providecommand{\rhfaPlsOneEvR}{}\renewcommand{\rhfaPlsOneEvR}{0.07}  % pls1 event: R_k
\providecommand{\rhfaPlsTwoEvPct}{}\renewcommand{\rhfaPlsTwoEvPct}{2.22}  % pls2 event: QLIKE change vs OLS-HAR, %
\providecommand{\rhfaPlsTwoEvMsePct}{}\renewcommand{\rhfaPlsTwoEvMsePct}{$-$0.33}  % pls2 event: MSE change vs OLS-HAR, %
\providecommand{\rhfaPlsTwoEvDm}{}\renewcommand{\rhfaPlsTwoEvDm}{1.79}  % pls2 event: DM vs OLS-HAR (QLIKE)
\providecommand{\rhfaPlsTwoEvP}{}\renewcommand{\rhfaPlsTwoEvP}{0.074}  % pls2 event: p
\providecommand{\rhfaPlsTwoEvDmFull}{}\renewcommand{\rhfaPlsTwoEvDmFull}{6.63}  % pls2 event: DM vs the full set (enet)
\providecommand{\rhfaPlsTwoEvPFull}{}\renewcommand{\rhfaPlsTwoEvPFull}{$<$0.001}  % pls2 event: p
\providecommand{\rhfaPlsTwoEvR}{}\renewcommand{\rhfaPlsTwoEvR}{$-$0.10}  % pls2 event: R_k
\providecommand{\rhfaPlsThreeEvPct}{}\renewcommand{\rhfaPlsThreeEvPct}{$-$0.12}  % pls3 event: QLIKE change vs OLS-HAR, %
\providecommand{\rhfaPlsThreeEvMsePct}{}\renewcommand{\rhfaPlsThreeEvMsePct}{$-$1.07}  % pls3 event: MSE change vs OLS-HAR, %
\providecommand{\rhfaPlsThreeEvDm}{}\renewcommand{\rhfaPlsThreeEvDm}{$-$0.06}  % pls3 event: DM vs OLS-HAR (QLIKE)
\providecommand{\rhfaPlsThreeEvP}{}\renewcommand{\rhfaPlsThreeEvP}{0.949}  % pls3 event: p
\providecommand{\rhfaPlsThreeEvDmFull}{}\renewcommand{\rhfaPlsThreeEvDmFull}{5.29}  % pls3 event: DM vs the full set (enet)
\providecommand{\rhfaPlsThreeEvPFull}{}\renewcommand{\rhfaPlsThreeEvPFull}{$<$0.001}  % pls3 event: p
\providecommand{\rhfaPlsThreeEvR}{}\renewcommand{\rhfaPlsThreeEvR}{0.01}  % pls3 event: R_k
\providecommand{\rhfaFullEvModel}{}\renewcommand{\rhfaFullEvModel}{elastic net}  % event: L_full model (lower QLIKE of ridge, elastic net)
\providecommand{\rhfaFullEvPct}{}\renewcommand{\rhfaFullEvPct}{$-$22.61}  % event: full set QLIKE change vs OLS-HAR, %
\providecommand{\rhfaFullEvMsePct}{}\renewcommand{\rhfaFullEvMsePct}{$-$8.54}  % event: full set MSE change vs OLS-HAR, %
\providecommand{\rhfaFullEvDm}{}\renewcommand{\rhfaFullEvDm}{$-$6.79}  % event: full set DM vs OLS-HAR
\providecommand{\rhfaFullEvP}{}\renewcommand{\rhfaFullEvP}{$<$0.001}  % event: p
\providecommand{\rhfaNEvDays}{}\renewcommand{\rhfaNEvDays}{81}  % event: days
\providecommand{\rhfaNEvBars}{}\renewcommand{\rhfaNEvBars}{1{,}053}  % event: bars
\providecommand{\rhfaReading}{}\renewcommand{\rhfaReading}{dense beyond a few factors}  % the RULE's reading, ORDINARY days
\providecommand{\rhfaBestSmall}{}\renewcommand{\rhfaBestSmall}{PLS, $k = 1$}  % the k <= 3 model with the highest R_k on ORDINARY days
\providecommand{\rhfaBestSmallR}{}\renewcommand{\rhfaBestSmallR}{0.00}  % its R_k
\providecommand{\rhfaBestSmallPFull}{}\renewcommand{\rhfaBestSmallPFull}{$<$0.001}  % its DM p against the full set
\providecommand{\rhfaNSharedModels}{}\renewcommand{\rhfaNSharedModels}{0}  % k <= 3 models meeting the SHARED FACTOR condition
\providecommand{\rhfaShareThreshold}{}\renewcommand{\rhfaShareThreshold}{0.75}  % SHARED FACTOR: R_k threshold
\providecommand{\rhfaDenseThreshold}{}\renewcommand{\rhfaDenseThreshold}{0.5}  % DENSE: R_10 threshold
\providecommand{\rhfaNBlock}{}\renewcommand{\rhfaNBlock}{258}  % block columns (value averages, 43 series x 6 horizons)
\providecommand{\rhfaNSeries}{}\renewcommand{\rhfaNSeries}{43}  % non-FOMC exogenous series
\providecommand{\rhfaNActiveMin}{}\renewcommand{\rhfaNActiveMin}{198}  % fewest active block columns on a scored day
\providecommand{\rhfaNActiveMax}{}\renewcommand{\rhfaNActiveMax}{257}  % most active block columns on a scored day
\providecommand{\rhfaSharePcOne}{}\renewcommand{\rhfaSharePcOne}{26.3}  % mean share of block variance, component 1, %
\providecommand{\rhfaSharePcTwo}{}\renewcommand{\rhfaSharePcTwo}{9.7}  % mean share of block variance, component 2, %
\providecommand{\rhfaSharePcThree}{}\renewcommand{\rhfaSharePcThree}{7.6}  % mean share of block variance, component 3, %
\providecommand{\rhfaShareCumOne}{}\renewcommand{\rhfaShareCumOne}{26.3}  % mean share in components 1-1, %
\providecommand{\rhfaShareCumTwo}{}\renewcommand{\rhfaShareCumTwo}{36.0}  % mean share in components 1-2, %
\providecommand{\rhfaShareCumThree}{}\renewcommand{\rhfaShareCumThree}{43.6}  % mean share in components 1-3, %
\providecommand{\rhfaShareCumFive}{}\renewcommand{\rhfaShareCumFive}{52.6}  % mean share in components 1-5, %
\providecommand{\rhfaShareCumTen}{}\renewcommand{\rhfaShareCumTen}{66.2}  % mean share in components 1-10, %
\providecommand{\rhfaShareCumTwenty}{}\renewcommand{\rhfaShareCumTwenty}{78.8}  % mean share in components 1-20, %
\providecommand{\rhfaShareCumThreeMin}{}\renewcommand{\rhfaShareCumThreeMin}{36.6}  % components 1-3: lowest day, %
\providecommand{\rhfaShareCumThreeMax}{}\renewcommand{\rhfaShareCumThreeMax}{47.1}  % components 1-3: highest day, %
\providecommand{\rhfaLoadRm}{}\renewcommand{\rhfaLoadRm}{2.61}  % PC1: sum of |loadings|, Realized moments
\providecommand{\rhfaLoadSqRm}{}\renewcommand{\rhfaLoadSqRm}{22.3}  % PC1: share of squared loadings, Realized moments, %
\providecommand{\rhfaPlsLoadRm}{}\renewcommand{\rhfaPlsLoadRm}{3.04}  % PLS1: sum of |weights|, Realized moments
\providecommand{\rhfaPlsLoadSqRm}{}\renewcommand{\rhfaPlsLoadSqRm}{34.0}  % PLS1: share of squared weights, Realized moments, %
\providecommand{\rhfaNColsRm}{}\renewcommand{\rhfaNColsRm}{42.0}  % active block columns of Realized moments (mean over days)
\providecommand{\rhfaLoadLiq}{}\renewcommand{\rhfaLoadLiq}{3.84}  % PC1: sum of |loadings|, Liquidity
\providecommand{\rhfaLoadSqLiq}{}\renewcommand{\rhfaLoadSqLiq}{21.2}  % PC1: share of squared loadings, Liquidity, %
\providecommand{\rhfaPlsLoadLiq}{}\renewcommand{\rhfaPlsLoadLiq}{3.95}  % PLS1: sum of |weights|, Liquidity
\providecommand{\rhfaPlsLoadSqLiq}{}\renewcommand{\rhfaPlsLoadSqLiq}{26.8}  % PLS1: share of squared weights, Liquidity, %
\providecommand{\rhfaNColsLiq}{}\renewcommand{\rhfaNColsLiq}{84.0}  % active block columns of Liquidity (mean over days)
\providecommand{\rhfaLoadEw}{}\renewcommand{\rhfaLoadEw}{2.97}  % PC1: sum of |loadings|, Equal-weighted market
\providecommand{\rhfaLoadSqEw}{}\renewcommand{\rhfaLoadSqEw}{28.2}  % PC1: share of squared loadings, Equal-weighted market, %
\providecommand{\rhfaPlsLoadEw}{}\renewcommand{\rhfaPlsLoadEw}{2.08}  % PLS1: sum of |weights|, Equal-weighted market
\providecommand{\rhfaPlsLoadSqEw}{}\renewcommand{\rhfaPlsLoadSqEw}{17.7}  % PLS1: share of squared weights, Equal-weighted market, %
\providecommand{\rhfaNColsEw}{}\renewcommand{\rhfaNColsEw}{36.0}  % active block columns of Equal-weighted market (mean over days)
\providecommand{\rhfaLoadVw}{}\renewcommand{\rhfaLoadVw}{2.91}  % PC1: sum of |loadings|, Value-weighted market
\providecommand{\rhfaLoadSqVw}{}\renewcommand{\rhfaLoadSqVw}{27.2}  % PC1: share of squared loadings, Value-weighted market, %
\providecommand{\rhfaPlsLoadVw}{}\renewcommand{\rhfaPlsLoadVw}{2.31}  % PLS1: sum of |weights|, Value-weighted market
\providecommand{\rhfaPlsLoadSqVw}{}\renewcommand{\rhfaPlsLoadSqVw}{21.2}  % PLS1: share of squared weights, Value-weighted market, %
\providecommand{\rhfaNColsVw}{}\renewcommand{\rhfaNColsVw}{36.0}  % active block columns of Value-weighted market (mean over days)
\providecommand{\rhfaLoadSent}{}\renewcommand{\rhfaLoadSent}{0.02}  % PC1: sum of |loadings|, Sentiment
\providecommand{\rhfaLoadSqSent}{}\renewcommand{\rhfaLoadSqSent}{0.0}  % PC1: share of squared loadings, Sentiment, %
\providecommand{\rhfaPlsLoadSent}{}\renewcommand{\rhfaPlsLoadSent}{0.04}  % PLS1: sum of |weights|, Sentiment
\providecommand{\rhfaPlsLoadSqSent}{}\renewcommand{\rhfaPlsLoadSqSent}{0.1}  % PLS1: share of squared weights, Sentiment, %
\providecommand{\rhfaNColsSent}{}\renewcommand{\rhfaNColsSent}{4.0}  % active block columns of Sentiment (mean over days)
\providecommand{\rhfaLoadIv}{}\renewcommand{\rhfaLoadIv}{0.18}  % PC1: sum of |loadings|, Implied volatility
\providecommand{\rhfaLoadSqIv}{}\renewcommand{\rhfaLoadSqIv}{1.1}  % PC1: share of squared loadings, Implied volatility, %
\providecommand{\rhfaPlsLoadIv}{}\renewcommand{\rhfaPlsLoadIv}{0.09}  % PLS1: sum of |weights|, Implied volatility
\providecommand{\rhfaPlsLoadSqIv}{}\renewcommand{\rhfaPlsLoadSqIv}{0.2}  % PLS1: share of squared weights, Implied volatility, %
\providecommand{\rhfaNColsIv}{}\renewcommand{\rhfaNColsIv}{5.1}  % active block columns of Implied volatility (mean over days)
\providecommand{\rhfaLoadVd}{}\renewcommand{\rhfaLoadVd}{0.01}  % PC1: sum of |loadings|, Volatility demand
\providecommand{\rhfaLoadSqVd}{}\renewcommand{\rhfaLoadSqVd}{0.0}  % PC1: share of squared loadings, Volatility demand, %
\providecommand{\rhfaPlsLoadVd}{}\renewcommand{\rhfaPlsLoadVd}{0.01}  % PLS1: sum of |weights|, Volatility demand
\providecommand{\rhfaPlsLoadSqVd}{}\renewcommand{\rhfaPlsLoadSqVd}{0.0}  % PLS1: share of squared weights, Volatility demand, %
\providecommand{\rhfaNColsVd}{}\renewcommand{\rhfaNColsVd}{3.1}  % active block columns of Volatility demand (mean over days)
\providecommand{\rhfaLoadTopGroup}{}\renewcommand{\rhfaLoadTopGroup}{equal-weighted market}  % group with the largest PC1 share of squared loadings
\providecommand{\rhfaGOneMaxRel}{}\renewcommand{\rhfaGOneMaxRel}{$4 \times 10^{-14}$}  % G1: k = 0 vs Table 1's OLS-HAR, max relative difference
\providecommand{\rhfaGTwoMaxRel}{}\renewcommand{\rhfaGTwoMaxRel}{$4 \times 10^{-14}$}  % G2: Gram route vs direct lstsq, max relative difference
\providecommand{\rhfaGTwoNDays}{}\renewcommand{\rhfaGTwoNDays}{7}  % G2 days
\providecommand{\rhfaSdReading}{}\renewcommand{\rhfaSdReading}{dense beyond a few factors}  % sensitivity (s.d.-only entry rule): reading
\providecommand{\rhfaSdPcaOneOrdR}{}\renewcommand{\rhfaSdPcaOneOrdR}{$-$0.02}  % sensitivity: pca1 R_k, ORDINARY days
\providecommand{\rhfaSdPcaTwoOrdR}{}\renewcommand{\rhfaSdPcaTwoOrdR}{$-$0.03}  % sensitivity: pca2 R_k, ORDINARY days
\providecommand{\rhfaSdPcaThreeOrdR}{}\renewcommand{\rhfaSdPcaThreeOrdR}{$-$0.14}  % sensitivity: pca3 R_k, ORDINARY days
\providecommand{\rhfaSdPlsOneOrdR}{}\renewcommand{\rhfaSdPlsOneOrdR}{0.00}  % sensitivity: pls1 R_k, ORDINARY days
\providecommand{\rhfaSdPlsTwoOrdR}{}\renewcommand{\rhfaSdPlsTwoOrdR}{$-$0.40}  % sensitivity: pls2 R_k, ORDINARY days
\providecommand{\rhfaSdPlsThreeOrdR}{}\renewcommand{\rhfaSdPlsThreeOrdR}{$-$0.04}  % sensitivity: pls3 R_k, ORDINARY days
\providecommand{\rhfaSdPcaTenOrdR}{}\renewcommand{\rhfaSdPcaTenOrdR}{0.21}  % sensitivity: pca10 R_k, ORDINARY days
\fi
\ifdefined\rhbyLXgbReading\else% AUTO-GENERATED by writeup/make_rehaul_beyond_ladder.py -- do not edit; re-run the script.
% results/ablations_2014start_w3000/scores_fomcfix; requested by the author after the beyond-base results; not pre-registered
\newcommand{\rhbyLLgbmFullAllPct}{$-$1.6}  % lgbm FULL vs BASE, all
\newcommand{\rhbyLLgbmKOneAllPct}{$-$1.2}  % lgbm k=1 vs BASE, all; requested by the author after the beyond-base results; not pre-registered
\newcommand{\rhbyLLgbmKOneAllDm}{$-$4.76}  % DM vs BASE
\newcommand{\rhbyLLgbmKOneAllP}{\ensuremath{<}0.001}  % p-value
\newcommand{\rhbyLLgbmKOneAllDmFull}{2.59}  % DM vs FULL (positive: FULL lower)
\newcommand{\rhbyLLgbmKOneAllPFull}{0.010}  % p-value
\newcommand{\rhbyLLgbmKOneAllR}{0.74}  % recovery (L_BASE - L_k) / (L_BASE - L_FULL)
\newcommand{\rhbyLLgbmKThreeAllPct}{$-$1.1}  % lgbm k=3 vs BASE, all; requested by the author after the beyond-base results; not pre-registered
\newcommand{\rhbyLLgbmKThreeAllDm}{$-$4.17}  % DM vs BASE
\newcommand{\rhbyLLgbmKThreeAllP}{\ensuremath{<}0.001}  % p-value
\newcommand{\rhbyLLgbmKThreeAllDmFull}{2.97}  % DM vs FULL (positive: FULL lower)
\newcommand{\rhbyLLgbmKThreeAllPFull}{0.003}  % p-value
\newcommand{\rhbyLLgbmKThreeAllR}{0.72}  % recovery (L_BASE - L_k) / (L_BASE - L_FULL)
\newcommand{\rhbyLLgbmKFiveAllPct}{$-$1.2}  % lgbm k=5 vs BASE, all; requested by the author after the beyond-base results; not pre-registered
\newcommand{\rhbyLLgbmKFiveAllDm}{$-$4.52}  % DM vs BASE
\newcommand{\rhbyLLgbmKFiveAllP}{\ensuremath{<}0.001}  % p-value
\newcommand{\rhbyLLgbmKFiveAllDmFull}{2.57}  % DM vs FULL (positive: FULL lower)
\newcommand{\rhbyLLgbmKFiveAllPFull}{0.010}  % p-value
\newcommand{\rhbyLLgbmKFiveAllR}{0.76}  % recovery (L_BASE - L_k) / (L_BASE - L_FULL)
\newcommand{\rhbyLLgbmKTenAllPct}{$-$1.1}  % lgbm k=10 vs BASE, all; requested by the author after the beyond-base results; not pre-registered
\newcommand{\rhbyLLgbmKTenAllDm}{$-$4.20}  % DM vs BASE
\newcommand{\rhbyLLgbmKTenAllP}{\ensuremath{<}0.001}  % p-value
\newcommand{\rhbyLLgbmKTenAllDmFull}{2.78}  % DM vs FULL (positive: FULL lower)
\newcommand{\rhbyLLgbmKTenAllPFull}{0.005}  % p-value
\newcommand{\rhbyLLgbmKTenAllR}{0.72}  % recovery (L_BASE - L_k) / (L_BASE - L_FULL)
\newcommand{\rhbyLLgbmKTwentyAllPct}{$-$1.4}  % lgbm k=20 vs BASE, all; requested by the author after the beyond-base results; not pre-registered
\newcommand{\rhbyLLgbmKTwentyAllDm}{$-$5.25}  % DM vs BASE
\newcommand{\rhbyLLgbmKTwentyAllP}{\ensuremath{<}0.001}  % p-value
\newcommand{\rhbyLLgbmKTwentyAllDmFull}{1.28}  % DM vs FULL (positive: FULL lower)
\newcommand{\rhbyLLgbmKTwentyAllPFull}{0.202}  % p-value
\newcommand{\rhbyLLgbmKTwentyAllR}{0.93}  % recovery (L_BASE - L_k) / (L_BASE - L_FULL)
\newcommand{\rhbyLLgbmFullOrdPct}{$-$1.6}  % lgbm FULL vs BASE, ord
\newcommand{\rhbyLLgbmKOneOrdPct}{$-$1.2}  % lgbm k=1 vs BASE, ord; requested by the author after the beyond-base results; not pre-registered
\newcommand{\rhbyLLgbmKOneOrdDm}{$-$5.13}  % DM vs BASE
\newcommand{\rhbyLLgbmKOneOrdP}{\ensuremath{<}0.001}  % p-value
\newcommand{\rhbyLLgbmKOneOrdDmFull}{2.61}  % DM vs FULL (positive: FULL lower)
\newcommand{\rhbyLLgbmKOneOrdPFull}{0.009}  % p-value
\newcommand{\rhbyLLgbmKOneOrdR}{0.75}  % recovery (L_BASE - L_k) / (L_BASE - L_FULL)
\newcommand{\rhbyLLgbmKThreeOrdPct}{$-$1.1}  % lgbm k=3 vs BASE, ord; requested by the author after the beyond-base results; not pre-registered
\newcommand{\rhbyLLgbmKThreeOrdDm}{$-$4.54}  % DM vs BASE
\newcommand{\rhbyLLgbmKThreeOrdP}{\ensuremath{<}0.001}  % p-value
\newcommand{\rhbyLLgbmKThreeOrdDmFull}{3.13}  % DM vs FULL (positive: FULL lower)
\newcommand{\rhbyLLgbmKThreeOrdPFull}{0.002}  % p-value
\newcommand{\rhbyLLgbmKThreeOrdR}{0.70}  % recovery (L_BASE - L_k) / (L_BASE - L_FULL)
\newcommand{\rhbyLLgbmKFiveOrdPct}{$-$1.2}  % lgbm k=5 vs BASE, ord; requested by the author after the beyond-base results; not pre-registered
\newcommand{\rhbyLLgbmKFiveOrdDm}{$-$4.99}  % DM vs BASE
\newcommand{\rhbyLLgbmKFiveOrdP}{\ensuremath{<}0.001}  % p-value
\newcommand{\rhbyLLgbmKFiveOrdDmFull}{3.03}  % DM vs FULL (positive: FULL lower)
\newcommand{\rhbyLLgbmKFiveOrdPFull}{0.002}  % p-value
\newcommand{\rhbyLLgbmKFiveOrdR}{0.72}  % recovery (L_BASE - L_k) / (L_BASE - L_FULL)
\newcommand{\rhbyLLgbmKTenOrdPct}{$-$1.1}  % lgbm k=10 vs BASE, ord; requested by the author after the beyond-base results; not pre-registered
\newcommand{\rhbyLLgbmKTenOrdDm}{$-$4.68}  % DM vs BASE
\newcommand{\rhbyLLgbmKTenOrdP}{\ensuremath{<}0.001}  % p-value
\newcommand{\rhbyLLgbmKTenOrdDmFull}{2.98}  % DM vs FULL (positive: FULL lower)
\newcommand{\rhbyLLgbmKTenOrdPFull}{0.003}  % p-value
\newcommand{\rhbyLLgbmKTenOrdR}{0.70}  % recovery (L_BASE - L_k) / (L_BASE - L_FULL)
\newcommand{\rhbyLLgbmKTwentyOrdPct}{$-$1.5}  % lgbm k=20 vs BASE, ord; requested by the author after the beyond-base results; not pre-registered
\newcommand{\rhbyLLgbmKTwentyOrdDm}{$-$5.82}  % DM vs BASE
\newcommand{\rhbyLLgbmKTwentyOrdP}{\ensuremath{<}0.001}  % p-value
\newcommand{\rhbyLLgbmKTwentyOrdDmFull}{1.03}  % DM vs FULL (positive: FULL lower)
\newcommand{\rhbyLLgbmKTwentyOrdPFull}{0.303}  % p-value
\newcommand{\rhbyLLgbmKTwentyOrdR}{0.94}  % recovery (L_BASE - L_k) / (L_BASE - L_FULL)
\newcommand{\rhbyLLgbmFullEvPct}{$-$0.6}  % lgbm FULL vs BASE, ev
\newcommand{\rhbyLLgbmKOneEvPct}{$-$0.3}  % lgbm k=1 vs BASE, ev; requested by the author after the beyond-base results; not pre-registered
\newcommand{\rhbyLLgbmKOneEvDm}{$-$0.26}  % DM vs BASE
\newcommand{\rhbyLLgbmKOneEvP}{0.795}  % p-value
\newcommand{\rhbyLLgbmKOneEvDmFull}{0.54}  % DM vs FULL (positive: FULL lower)
\newcommand{\rhbyLLgbmKOneEvPFull}{0.592}  % p-value
\newcommand{\rhbyLLgbmKOneEvR}{0.44}  % recovery (L_BASE - L_k) / (L_BASE - L_FULL)
\newcommand{\rhbyLLgbmKThreeEvPct}{$-$0.8}  % lgbm k=3 vs BASE, ev; requested by the author after the beyond-base results; not pre-registered
\newcommand{\rhbyLLgbmKThreeEvDm}{$-$0.63}  % DM vs BASE
\newcommand{\rhbyLLgbmKThreeEvP}{0.526}  % p-value
\newcommand{\rhbyLLgbmKThreeEvDmFull}{$-$0.38}  % DM vs FULL (positive: FULL lower)
\newcommand{\rhbyLLgbmKThreeEvPFull}{0.701}  % p-value
\newcommand{\rhbyLLgbmKThreeEvR}{1.36}  % recovery (L_BASE - L_k) / (L_BASE - L_FULL)
\newcommand{\rhbyLLgbmKFiveEvPct}{$-$1.6}  % lgbm k=5 vs BASE, ev; requested by the author after the beyond-base results; not pre-registered
\newcommand{\rhbyLLgbmKFiveEvDm}{$-$1.03}  % DM vs BASE
\newcommand{\rhbyLLgbmKFiveEvP}{0.305}  % p-value
\newcommand{\rhbyLLgbmKFiveEvDmFull}{$-$1.72}  % DM vs FULL (positive: FULL lower)
\newcommand{\rhbyLLgbmKFiveEvPFull}{0.086}  % p-value
\newcommand{\rhbyLLgbmKFiveEvR}{2.62}  % recovery (L_BASE - L_k) / (L_BASE - L_FULL)
\newcommand{\rhbyLLgbmKTenEvPct}{$-$1.0}  % lgbm k=10 vs BASE, ev; requested by the author after the beyond-base results; not pre-registered
\newcommand{\rhbyLLgbmKTenEvDm}{$-$0.71}  % DM vs BASE
\newcommand{\rhbyLLgbmKTenEvP}{0.477}  % p-value
\newcommand{\rhbyLLgbmKTenEvDmFull}{$-$1.05}  % DM vs FULL (positive: FULL lower)
\newcommand{\rhbyLLgbmKTenEvPFull}{0.294}  % p-value
\newcommand{\rhbyLLgbmKTenEvR}{1.66}  % recovery (L_BASE - L_k) / (L_BASE - L_FULL)
\newcommand{\rhbyLLgbmKTwentyEvPct}{$-$0.3}  % lgbm k=20 vs BASE, ev; requested by the author after the beyond-base results; not pre-registered
\newcommand{\rhbyLLgbmKTwentyEvDm}{$-$0.18}  % DM vs BASE
\newcommand{\rhbyLLgbmKTwentyEvP}{0.855}  % p-value
\newcommand{\rhbyLLgbmKTwentyEvDmFull}{1.20}  % DM vs FULL (positive: FULL lower)
\newcommand{\rhbyLLgbmKTwentyEvPFull}{0.230}  % p-value
\newcommand{\rhbyLLgbmKTwentyEvR}{0.40}  % recovery (L_BASE - L_k) / (L_BASE - L_FULL)
\newcommand{\rhbyLLgbmReading}{intermediate}  % section 2 reading, ORDINARY days
\newcommand{\rhbyLLgbmTop}{sumabsret\_vwstock, sumabsret\_ewstock, effspread\_vwstock, vix, sumret3\_vwstock}  % the model's top 5 block series
\newcommand{\rhbyLXgbFullAllPct}{$-$1.8}  % xgb FULL vs BASE, all
\newcommand{\rhbyLXgbKOneAllPct}{$-$0.6}  % xgb k=1 vs BASE, all; requested by the author after the beyond-base results; not pre-registered
\newcommand{\rhbyLXgbKOneAllDm}{$-$5.62}  % DM vs BASE
\newcommand{\rhbyLXgbKOneAllP}{\ensuremath{<}0.001}  % p-value
\newcommand{\rhbyLXgbKOneAllDmFull}{4.04}  % DM vs FULL (positive: FULL lower)
\newcommand{\rhbyLXgbKOneAllPFull}{\ensuremath{<}0.001}  % p-value
\newcommand{\rhbyLXgbKOneAllR}{0.33}  % recovery (L_BASE - L_k) / (L_BASE - L_FULL)
\newcommand{\rhbyLXgbKThreeAllPct}{$-$1.4}  % xgb k=3 vs BASE, all; requested by the author after the beyond-base results; not pre-registered
\newcommand{\rhbyLXgbKThreeAllDm}{$-$4.92}  % DM vs BASE
\newcommand{\rhbyLXgbKThreeAllP}{\ensuremath{<}0.001}  % p-value
\newcommand{\rhbyLXgbKThreeAllDmFull}{2.59}  % DM vs FULL (positive: FULL lower)
\newcommand{\rhbyLXgbKThreeAllPFull}{0.010}  % p-value
\newcommand{\rhbyLXgbKThreeAllR}{0.77}  % recovery (L_BASE - L_k) / (L_BASE - L_FULL)
\newcommand{\rhbyLXgbKFiveAllPct}{$-$1.6}  % xgb k=5 vs BASE, all; requested by the author after the beyond-base results; not pre-registered
\newcommand{\rhbyLXgbKFiveAllDm}{$-$5.41}  % DM vs BASE
\newcommand{\rhbyLXgbKFiveAllP}{\ensuremath{<}0.001}  % p-value
\newcommand{\rhbyLXgbKFiveAllDmFull}{1.53}  % DM vs FULL (positive: FULL lower)
\newcommand{\rhbyLXgbKFiveAllPFull}{0.127}  % p-value
\newcommand{\rhbyLXgbKFiveAllR}{0.90}  % recovery (L_BASE - L_k) / (L_BASE - L_FULL)
\newcommand{\rhbyLXgbKTenAllPct}{$-$1.4}  % xgb k=10 vs BASE, all; requested by the author after the beyond-base results; not pre-registered
\newcommand{\rhbyLXgbKTenAllDm}{$-$4.81}  % DM vs BASE
\newcommand{\rhbyLXgbKTenAllP}{\ensuremath{<}0.001}  % p-value
\newcommand{\rhbyLXgbKTenAllDmFull}{3.28}  % DM vs FULL (positive: FULL lower)
\newcommand{\rhbyLXgbKTenAllPFull}{0.001}  % p-value
\newcommand{\rhbyLXgbKTenAllR}{0.77}  % recovery (L_BASE - L_k) / (L_BASE - L_FULL)
\newcommand{\rhbyLXgbKTwentyAllPct}{$-$1.6}  % xgb k=20 vs BASE, all; requested by the author after the beyond-base results; not pre-registered
\newcommand{\rhbyLXgbKTwentyAllDm}{$-$4.87}  % DM vs BASE
\newcommand{\rhbyLXgbKTwentyAllP}{\ensuremath{<}0.001}  % p-value
\newcommand{\rhbyLXgbKTwentyAllDmFull}{1.97}  % DM vs FULL (positive: FULL lower)
\newcommand{\rhbyLXgbKTwentyAllPFull}{0.049}  % p-value
\newcommand{\rhbyLXgbKTwentyAllR}{0.90}  % recovery (L_BASE - L_k) / (L_BASE - L_FULL)
\newcommand{\rhbyLXgbFullOrdPct}{$-$1.8}  % xgb FULL vs BASE, ord
\newcommand{\rhbyLXgbKOneOrdPct}{$-$0.6}  % xgb k=1 vs BASE, ord; requested by the author after the beyond-base results; not pre-registered
\newcommand{\rhbyLXgbKOneOrdDm}{$-$5.89}  % DM vs BASE
\newcommand{\rhbyLXgbKOneOrdP}{\ensuremath{<}0.001}  % p-value
\newcommand{\rhbyLXgbKOneOrdDmFull}{4.42}  % DM vs FULL (positive: FULL lower)
\newcommand{\rhbyLXgbKOneOrdPFull}{\ensuremath{<}0.001}  % p-value
\newcommand{\rhbyLXgbKOneOrdR}{0.33}  % recovery (L_BASE - L_k) / (L_BASE - L_FULL)
\newcommand{\rhbyLXgbKThreeOrdPct}{$-$1.5}  % xgb k=3 vs BASE, ord; requested by the author after the beyond-base results; not pre-registered
\newcommand{\rhbyLXgbKThreeOrdDm}{$-$5.55}  % DM vs BASE
\newcommand{\rhbyLXgbKThreeOrdP}{\ensuremath{<}0.001}  % p-value
\newcommand{\rhbyLXgbKThreeOrdDmFull}{2.32}  % DM vs FULL (positive: FULL lower)
\newcommand{\rhbyLXgbKThreeOrdPFull}{0.020}  % p-value
\newcommand{\rhbyLXgbKThreeOrdR}{0.81}  % recovery (L_BASE - L_k) / (L_BASE - L_FULL)
\newcommand{\rhbyLXgbKFiveOrdPct}{$-$1.6}  % xgb k=5 vs BASE, ord; requested by the author after the beyond-base results; not pre-registered
\newcommand{\rhbyLXgbKFiveOrdDm}{$-$5.95}  % DM vs BASE
\newcommand{\rhbyLXgbKFiveOrdP}{\ensuremath{<}0.001}  % p-value
\newcommand{\rhbyLXgbKFiveOrdDmFull}{1.43}  % DM vs FULL (positive: FULL lower)
\newcommand{\rhbyLXgbKFiveOrdPFull}{0.152}  % p-value
\newcommand{\rhbyLXgbKFiveOrdR}{0.91}  % recovery (L_BASE - L_k) / (L_BASE - L_FULL)
\newcommand{\rhbyLXgbKTenOrdPct}{$-$1.4}  % xgb k=10 vs BASE, ord; requested by the author after the beyond-base results; not pre-registered
\newcommand{\rhbyLXgbKTenOrdDm}{$-$5.42}  % DM vs BASE
\newcommand{\rhbyLXgbKTenOrdP}{\ensuremath{<}0.001}  % p-value
\newcommand{\rhbyLXgbKTenOrdDmFull}{3.18}  % DM vs FULL (positive: FULL lower)
\newcommand{\rhbyLXgbKTenOrdPFull}{0.001}  % p-value
\newcommand{\rhbyLXgbKTenOrdR}{0.77}  % recovery (L_BASE - L_k) / (L_BASE - L_FULL)
\newcommand{\rhbyLXgbKTwentyOrdPct}{$-$1.7}  % xgb k=20 vs BASE, ord; requested by the author after the beyond-base results; not pre-registered
\newcommand{\rhbyLXgbKTwentyOrdDm}{$-$5.48}  % DM vs BASE
\newcommand{\rhbyLXgbKTwentyOrdP}{\ensuremath{<}0.001}  % p-value
\newcommand{\rhbyLXgbKTwentyOrdDmFull}{1.60}  % DM vs FULL (positive: FULL lower)
\newcommand{\rhbyLXgbKTwentyOrdPFull}{0.110}  % p-value
\newcommand{\rhbyLXgbKTwentyOrdR}{0.92}  % recovery (L_BASE - L_k) / (L_BASE - L_FULL)
\newcommand{\rhbyLXgbFullEvPct}{$-$1.3}  % xgb FULL vs BASE, ev
\newcommand{\rhbyLXgbKOneEvPct}{$-$0.5}  % xgb k=1 vs BASE, ev; requested by the author after the beyond-base results; not pre-registered
\newcommand{\rhbyLXgbKOneEvDm}{$-$0.63}  % DM vs BASE
\newcommand{\rhbyLXgbKOneEvP}{0.530}  % p-value
\newcommand{\rhbyLXgbKOneEvDmFull}{0.57}  % DM vs FULL (positive: FULL lower)
\newcommand{\rhbyLXgbKOneEvPFull}{0.568}  % p-value
\newcommand{\rhbyLXgbKOneEvR}{0.35}  % recovery (L_BASE - L_k) / (L_BASE - L_FULL)
\newcommand{\rhbyLXgbKThreeEvPct}{0.0}  % xgb k=3 vs BASE, ev; requested by the author after the beyond-base results; not pre-registered
\newcommand{\rhbyLXgbKThreeEvDm}{0.01}  % DM vs BASE
\newcommand{\rhbyLXgbKThreeEvP}{0.992}  % p-value
\newcommand{\rhbyLXgbKThreeEvDmFull}{1.34}  % DM vs FULL (positive: FULL lower)
\newcommand{\rhbyLXgbKThreeEvPFull}{0.180}  % p-value
\newcommand{\rhbyLXgbKThreeEvR}{$-$0.01}  % recovery (L_BASE - L_k) / (L_BASE - L_FULL)
\newcommand{\rhbyLXgbKFiveEvPct}{$-$0.9}  % xgb k=5 vs BASE, ev; requested by the author after the beyond-base results; not pre-registered
\newcommand{\rhbyLXgbKFiveEvDm}{$-$0.60}  % DM vs BASE
\newcommand{\rhbyLXgbKFiveEvP}{0.549}  % p-value
\newcommand{\rhbyLXgbKFiveEvDmFull}{0.48}  % DM vs FULL (positive: FULL lower)
\newcommand{\rhbyLXgbKFiveEvPFull}{0.632}  % p-value
\newcommand{\rhbyLXgbKFiveEvR}{0.71}  % recovery (L_BASE - L_k) / (L_BASE - L_FULL)
\newcommand{\rhbyLXgbKTenEvPct}{$-$1.0}  % xgb k=10 vs BASE, ev; requested by the author after the beyond-base results; not pre-registered
\newcommand{\rhbyLXgbKTenEvDm}{$-$0.63}  % DM vs BASE
\newcommand{\rhbyLXgbKTenEvP}{0.529}  % p-value
\newcommand{\rhbyLXgbKTenEvDmFull}{0.46}  % DM vs FULL (positive: FULL lower)
\newcommand{\rhbyLXgbKTenEvPFull}{0.648}  % p-value
\newcommand{\rhbyLXgbKTenEvR}{0.79}  % recovery (L_BASE - L_k) / (L_BASE - L_FULL)
\newcommand{\rhbyLXgbKTwentyEvPct}{$-$0.8}  % xgb k=20 vs BASE, ev; requested by the author after the beyond-base results; not pre-registered
\newcommand{\rhbyLXgbKTwentyEvDm}{$-$0.42}  % DM vs BASE
\newcommand{\rhbyLXgbKTwentyEvP}{0.674}  % p-value
\newcommand{\rhbyLXgbKTwentyEvDmFull}{0.79}  % DM vs FULL (positive: FULL lower)
\newcommand{\rhbyLXgbKTwentyEvPFull}{0.427}  % p-value
\newcommand{\rhbyLXgbKTwentyEvR}{0.59}  % recovery (L_BASE - L_k) / (L_BASE - L_FULL)
\newcommand{\rhbyLXgbReading}{concentrated}  % section 2 reading, ORDINARY days
\newcommand{\rhbyLXgbTop}{stocktwits\_sentiment, sumabsret\_vwstock, effspread\_vwstock, sumabsret\_ewstock, vix}  % the model's top 5 block series
\newcommand{\rhbyLAgree}{split}  % the two models' readings
\fi
\ifdefined\rhbyFXgbReading\else% AUTO-GENERATED by writeup/make_rehaul_beyond_factor.py -- do not edit; re-run the script.
% The trees' factor test beyond HAR + RM + FOMC (requested by the author after the beyond-base results; not pre-registered); results/beyond_trees_2026-10-27/factor.
\newcommand{\rhbyFLgbmBaseOrdQlike}{0.1584}  % QLIKE lgbm base ordinary
\newcommand{\rhbyFLgbmBaseOrdDmFull}{5.26}  % DM vs FULL
\newcommand{\rhbyFLgbmBaseOrdPFull}{\ensuremath{<}0.001}  % p vs FULL
\newcommand{\rhbyFLgbmPcaOneOrdQlike}{0.1583}  % QLIKE lgbm pca1 ordinary
\newcommand{\rhbyFLgbmPcaOneOrdPct}{$-$0.04}  % QLIKE change % vs BASE
\newcommand{\rhbyFLgbmPcaOneOrdDm}{$-$0.47}  % DM vs BASE (HLN, daily means, lag 8)
\newcommand{\rhbyFLgbmPcaOneOrdP}{0.641}  % p vs BASE
\newcommand{\rhbyFLgbmPcaOneOrdR}{0.02}  % recovery R
\newcommand{\rhbyFLgbmPcaOneOrdDmFull}{5.26}  % DM vs FULL
\newcommand{\rhbyFLgbmPcaOneOrdPFull}{\ensuremath{<}0.001}  % p vs FULL
\newcommand{\rhbyFLgbmPcaTwoOrdQlike}{0.1583}  % QLIKE lgbm pca2 ordinary
\newcommand{\rhbyFLgbmPcaTwoOrdPct}{$-$0.02}  % QLIKE change % vs BASE
\newcommand{\rhbyFLgbmPcaTwoOrdDm}{$-$0.24}  % DM vs BASE (HLN, daily means, lag 8)
\newcommand{\rhbyFLgbmPcaTwoOrdP}{0.807}  % p vs BASE
\newcommand{\rhbyFLgbmPcaTwoOrdR}{0.02}  % recovery R
\newcommand{\rhbyFLgbmPcaTwoOrdDmFull}{5.53}  % DM vs FULL
\newcommand{\rhbyFLgbmPcaTwoOrdPFull}{\ensuremath{<}0.001}  % p vs FULL
\newcommand{\rhbyFLgbmPcaThreeOrdQlike}{0.1582}  % QLIKE lgbm pca3 ordinary
\newcommand{\rhbyFLgbmPcaThreeOrdPct}{$-$0.13}  % QLIKE change % vs BASE
\newcommand{\rhbyFLgbmPcaThreeOrdDm}{$-$1.13}  % DM vs BASE (HLN, daily means, lag 8)
\newcommand{\rhbyFLgbmPcaThreeOrdP}{0.256}  % p vs BASE
\newcommand{\rhbyFLgbmPcaThreeOrdR}{0.08}  % recovery R
\newcommand{\rhbyFLgbmPcaThreeOrdDmFull}{5.51}  % DM vs FULL
\newcommand{\rhbyFLgbmPcaThreeOrdPFull}{\ensuremath{<}0.001}  % p vs FULL
\newcommand{\rhbyFLgbmPcaFiveOrdQlike}{0.1581}  % QLIKE lgbm pca5 ordinary
\newcommand{\rhbyFLgbmPcaFiveOrdPct}{$-$0.16}  % QLIKE change % vs BASE
\newcommand{\rhbyFLgbmPcaFiveOrdDm}{$-$1.13}  % DM vs BASE (HLN, daily means, lag 8)
\newcommand{\rhbyFLgbmPcaFiveOrdP}{0.257}  % p vs BASE
\newcommand{\rhbyFLgbmPcaFiveOrdR}{0.10}  % recovery R
\newcommand{\rhbyFLgbmPcaFiveOrdDmFull}{5.52}  % DM vs FULL
\newcommand{\rhbyFLgbmPcaFiveOrdPFull}{\ensuremath{<}0.001}  % p vs FULL
\newcommand{\rhbyFLgbmPcaTenOrdQlike}{0.1575}  % QLIKE lgbm pca10 ordinary
\newcommand{\rhbyFLgbmPcaTenOrdPct}{$-$0.57}  % QLIKE change % vs BASE
\newcommand{\rhbyFLgbmPcaTenOrdDm}{$-$2.94}  % DM vs BASE (HLN, daily means, lag 8)
\newcommand{\rhbyFLgbmPcaTenOrdP}{0.003}  % p vs BASE
\newcommand{\rhbyFLgbmPcaTenOrdR}{0.35}  % recovery R
\newcommand{\rhbyFLgbmPcaTenOrdDmFull}{4.66}  % DM vs FULL
\newcommand{\rhbyFLgbmPcaTenOrdPFull}{\ensuremath{<}0.001}  % p vs FULL
\newcommand{\rhbyFLgbmPlsOneOrdQlike}{0.1581}  % QLIKE lgbm pls1 ordinary
\newcommand{\rhbyFLgbmPlsOneOrdPct}{$-$0.19}  % QLIKE change % vs BASE
\newcommand{\rhbyFLgbmPlsOneOrdDm}{$-$1.67}  % DM vs BASE (HLN, daily means, lag 8)
\newcommand{\rhbyFLgbmPlsOneOrdP}{0.096}  % p vs BASE
\newcommand{\rhbyFLgbmPlsOneOrdR}{0.12}  % recovery R
\newcommand{\rhbyFLgbmPlsOneOrdDmFull}{5.31}  % DM vs FULL
\newcommand{\rhbyFLgbmPlsOneOrdPFull}{\ensuremath{<}0.001}  % p vs FULL
\newcommand{\rhbyFLgbmPlsThreeOrdQlike}{0.1575}  % QLIKE lgbm pls3 ordinary
\newcommand{\rhbyFLgbmPlsThreeOrdPct}{$-$0.52}  % QLIKE change % vs BASE
\newcommand{\rhbyFLgbmPlsThreeOrdDm}{$-$2.95}  % DM vs BASE (HLN, daily means, lag 8)
\newcommand{\rhbyFLgbmPlsThreeOrdP}{0.003}  % p vs BASE
\newcommand{\rhbyFLgbmPlsThreeOrdR}{0.32}  % recovery R
\newcommand{\rhbyFLgbmPlsThreeOrdDmFull}{4.30}  % DM vs FULL
\newcommand{\rhbyFLgbmPlsThreeOrdPFull}{\ensuremath{<}0.001}  % p vs FULL
\newcommand{\rhbyFLgbmFullOrdQlike}{0.1558}  % QLIKE lgbm full ordinary
\newcommand{\rhbyFLgbmFullOrdPct}{$-$1.62}  % QLIKE change % vs BASE
\newcommand{\rhbyFLgbmFullOrdDm}{$-$5.26}  % DM vs BASE (HLN, daily means, lag 8)
\newcommand{\rhbyFLgbmFullOrdP}{\ensuremath{<}0.001}  % p vs BASE
\newcommand{\rhbyFLgbmFullOrdR}{1.00}  % recovery R
\newcommand{\rhbyFLgbmBaseAllQlike}{0.1636}  % QLIKE lgbm base all
\newcommand{\rhbyFLgbmBaseAllDmFull}{4.92}  % DM vs FULL
\newcommand{\rhbyFLgbmBaseAllPFull}{\ensuremath{<}0.001}  % p vs FULL
\newcommand{\rhbyFLgbmPcaOneAllQlike}{0.1636}  % QLIKE lgbm pca1 all
\newcommand{\rhbyFLgbmPcaOneAllPct}{0.00}  % QLIKE change % vs BASE
\newcommand{\rhbyFLgbmPcaOneAllDm}{$-$0.01}  % DM vs BASE (HLN, daily means, lag 8)
\newcommand{\rhbyFLgbmPcaOneAllP}{0.992}  % p vs BASE
\newcommand{\rhbyFLgbmPcaOneAllR}{0.00}  % recovery R
\newcommand{\rhbyFLgbmPcaOneAllDmFull}{4.91}  % DM vs FULL
\newcommand{\rhbyFLgbmPcaOneAllPFull}{\ensuremath{<}0.001}  % p vs FULL
\newcommand{\rhbyFLgbmPcaTwoAllQlike}{0.1636}  % QLIKE lgbm pca2 all
\newcommand{\rhbyFLgbmPcaTwoAllPct}{0.00}  % QLIKE change % vs BASE
\newcommand{\rhbyFLgbmPcaTwoAllDm}{$-$0.03}  % DM vs BASE (HLN, daily means, lag 8)
\newcommand{\rhbyFLgbmPcaTwoAllP}{0.974}  % p vs BASE
\newcommand{\rhbyFLgbmPcaTwoAllR}{0.00}  % recovery R
\newcommand{\rhbyFLgbmPcaTwoAllDmFull}{5.17}  % DM vs FULL
\newcommand{\rhbyFLgbmPcaTwoAllPFull}{\ensuremath{<}0.001}  % p vs FULL
\newcommand{\rhbyFLgbmPcaThreeAllQlike}{0.1634}  % QLIKE lgbm pca3 all
\newcommand{\rhbyFLgbmPcaThreeAllPct}{$-$0.10}  % QLIKE change % vs BASE
\newcommand{\rhbyFLgbmPcaThreeAllDm}{$-$0.86}  % DM vs BASE (HLN, daily means, lag 8)
\newcommand{\rhbyFLgbmPcaThreeAllP}{0.391}  % p vs BASE
\newcommand{\rhbyFLgbmPcaThreeAllR}{0.06}  % recovery R
\newcommand{\rhbyFLgbmPcaThreeAllDmFull}{5.26}  % DM vs FULL
\newcommand{\rhbyFLgbmPcaThreeAllPFull}{\ensuremath{<}0.001}  % p vs FULL
\newcommand{\rhbyFLgbmPcaFiveAllQlike}{0.1635}  % QLIKE lgbm pca5 all
\newcommand{\rhbyFLgbmPcaFiveAllPct}{$-$0.09}  % QLIKE change % vs BASE
\newcommand{\rhbyFLgbmPcaFiveAllDm}{$-$0.65}  % DM vs BASE (HLN, daily means, lag 8)
\newcommand{\rhbyFLgbmPcaFiveAllP}{0.515}  % p vs BASE
\newcommand{\rhbyFLgbmPcaFiveAllR}{0.06}  % recovery R
\newcommand{\rhbyFLgbmPcaFiveAllDmFull}{5.14}  % DM vs FULL
\newcommand{\rhbyFLgbmPcaFiveAllPFull}{\ensuremath{<}0.001}  % p vs FULL
\newcommand{\rhbyFLgbmPcaTenAllQlike}{0.1628}  % QLIKE lgbm pca10 all
\newcommand{\rhbyFLgbmPcaTenAllPct}{$-$0.47}  % QLIKE change % vs BASE
\newcommand{\rhbyFLgbmPcaTenAllDm}{$-$2.58}  % DM vs BASE (HLN, daily means, lag 8)
\newcommand{\rhbyFLgbmPcaTenAllP}{0.010}  % p vs BASE
\newcommand{\rhbyFLgbmPcaTenAllR}{0.30}  % recovery R
\newcommand{\rhbyFLgbmPcaTenAllDmFull}{4.46}  % DM vs FULL
\newcommand{\rhbyFLgbmPcaTenAllPFull}{\ensuremath{<}0.001}  % p vs FULL
\newcommand{\rhbyFLgbmPlsOneAllQlike}{0.1634}  % QLIKE lgbm pls1 all
\newcommand{\rhbyFLgbmPlsOneAllPct}{$-$0.15}  % QLIKE change % vs BASE
\newcommand{\rhbyFLgbmPlsOneAllDm}{$-$1.35}  % DM vs BASE (HLN, daily means, lag 8)
\newcommand{\rhbyFLgbmPlsOneAllP}{0.178}  % p vs BASE
\newcommand{\rhbyFLgbmPlsOneAllR}{0.10}  % recovery R
\newcommand{\rhbyFLgbmPlsOneAllDmFull}{4.91}  % DM vs FULL
\newcommand{\rhbyFLgbmPlsOneAllPFull}{\ensuremath{<}0.001}  % p vs FULL
\newcommand{\rhbyFLgbmPlsThreeAllQlike}{0.1628}  % QLIKE lgbm pls3 all
\newcommand{\rhbyFLgbmPlsThreeAllPct}{$-$0.48}  % QLIKE change % vs BASE
\newcommand{\rhbyFLgbmPlsThreeAllDm}{$-$2.86}  % DM vs BASE (HLN, daily means, lag 8)
\newcommand{\rhbyFLgbmPlsThreeAllP}{0.004}  % p vs BASE
\newcommand{\rhbyFLgbmPlsThreeAllR}{0.31}  % recovery R
\newcommand{\rhbyFLgbmPlsThreeAllDmFull}{4.04}  % DM vs FULL
\newcommand{\rhbyFLgbmPlsThreeAllPFull}{\ensuremath{<}0.001}  % p vs FULL
\newcommand{\rhbyFLgbmFullAllQlike}{0.1611}  % QLIKE lgbm full all
\newcommand{\rhbyFLgbmFullAllPct}{$-$1.56}  % QLIKE change % vs BASE
\newcommand{\rhbyFLgbmFullAllDm}{$-$4.92}  % DM vs BASE (HLN, daily means, lag 8)
\newcommand{\rhbyFLgbmFullAllP}{\ensuremath{<}0.001}  % p vs BASE
\newcommand{\rhbyFLgbmFullAllR}{1.00}  % recovery R
\newcommand{\rhbyFLgbmBaseEvQlike}{0.3252}  % QLIKE lgbm base event
\newcommand{\rhbyFLgbmBaseEvDmFull}{0.47}  % DM vs FULL
\newcommand{\rhbyFLgbmBaseEvPFull}{0.642}  % p vs FULL
\newcommand{\rhbyFLgbmPcaOneEvQlike}{0.3271}  % QLIKE lgbm pca1 event
\newcommand{\rhbyFLgbmPcaOneEvPct}{0.58}  % QLIKE change % vs BASE
\newcommand{\rhbyFLgbmPcaOneEvDm}{0.91}  % DM vs BASE (HLN, daily means, lag 8)
\newcommand{\rhbyFLgbmPcaOneEvP}{0.365}  % p vs BASE
\newcommand{\rhbyFLgbmPcaOneEvR}{$-$0.94}  % recovery R
\newcommand{\rhbyFLgbmPcaOneEvDmFull}{0.69}  % DM vs FULL
\newcommand{\rhbyFLgbmPcaOneEvPFull}{0.489}  % p vs FULL
\newcommand{\rhbyFLgbmPcaTwoEvQlike}{0.3263}  % QLIKE lgbm pca2 event
\newcommand{\rhbyFLgbmPcaTwoEvPct}{0.32}  % QLIKE change % vs BASE
\newcommand{\rhbyFLgbmPcaTwoEvDm}{0.77}  % DM vs BASE (HLN, daily means, lag 8)
\newcommand{\rhbyFLgbmPcaTwoEvP}{0.443}  % p vs BASE
\newcommand{\rhbyFLgbmPcaTwoEvR}{$-$0.52}  % recovery R
\newcommand{\rhbyFLgbmPcaTwoEvDmFull}{0.67}  % DM vs FULL
\newcommand{\rhbyFLgbmPcaTwoEvPFull}{0.506}  % p vs FULL
\newcommand{\rhbyFLgbmPcaThreeEvQlike}{0.3267}  % QLIKE lgbm pca3 event
\newcommand{\rhbyFLgbmPcaThreeEvPct}{0.46}  % QLIKE change % vs BASE
\newcommand{\rhbyFLgbmPcaThreeEvDm}{1.00}  % DM vs BASE (HLN, daily means, lag 8)
\newcommand{\rhbyFLgbmPcaThreeEvP}{0.319}  % p vs BASE
\newcommand{\rhbyFLgbmPcaThreeEvR}{$-$0.74}  % recovery R
\newcommand{\rhbyFLgbmPcaThreeEvDmFull}{0.80}  % DM vs FULL
\newcommand{\rhbyFLgbmPcaThreeEvPFull}{0.425}  % p vs FULL
\newcommand{\rhbyFLgbmPcaFiveEvQlike}{0.3286}  % QLIKE lgbm pca5 event
\newcommand{\rhbyFLgbmPcaFiveEvPct}{1.05}  % QLIKE change % vs BASE
\newcommand{\rhbyFLgbmPcaFiveEvDm}{2.06}  % DM vs BASE (HLN, daily means, lag 8)
\newcommand{\rhbyFLgbmPcaFiveEvP}{0.040}  % p vs BASE
\newcommand{\rhbyFLgbmPcaFiveEvR}{$-$1.68}  % recovery R
\newcommand{\rhbyFLgbmPcaFiveEvDmFull}{1.04}  % DM vs FULL
\newcommand{\rhbyFLgbmPcaFiveEvPFull}{0.298}  % p vs FULL
\newcommand{\rhbyFLgbmPcaTenEvQlike}{0.3284}  % QLIKE lgbm pca10 event
\newcommand{\rhbyFLgbmPcaTenEvPct}{0.97}  % QLIKE change % vs BASE
\newcommand{\rhbyFLgbmPcaTenEvDm}{2.23}  % DM vs BASE (HLN, daily means, lag 8)
\newcommand{\rhbyFLgbmPcaTenEvP}{0.026}  % p vs BASE
\newcommand{\rhbyFLgbmPcaTenEvR}{$-$1.56}  % recovery R
\newcommand{\rhbyFLgbmPcaTenEvDmFull}{1.09}  % DM vs FULL
\newcommand{\rhbyFLgbmPcaTenEvPFull}{0.278}  % p vs FULL
\newcommand{\rhbyFLgbmPlsOneEvQlike}{0.3268}  % QLIKE lgbm pls1 event
\newcommand{\rhbyFLgbmPlsOneEvPct}{0.49}  % QLIKE change % vs BASE
\newcommand{\rhbyFLgbmPlsOneEvDm}{0.76}  % DM vs BASE (HLN, daily means, lag 8)
\newcommand{\rhbyFLgbmPlsOneEvP}{0.446}  % p vs BASE
\newcommand{\rhbyFLgbmPlsOneEvR}{$-$0.78}  % recovery R
\newcommand{\rhbyFLgbmPlsOneEvDmFull}{0.67}  % DM vs FULL
\newcommand{\rhbyFLgbmPlsOneEvPFull}{0.505}  % p vs FULL
\newcommand{\rhbyFLgbmPlsThreeEvQlike}{0.3256}  % QLIKE lgbm pls3 event
\newcommand{\rhbyFLgbmPlsThreeEvPct}{0.11}  % QLIKE change % vs BASE
\newcommand{\rhbyFLgbmPlsThreeEvDm}{0.32}  % DM vs BASE (HLN, daily means, lag 8)
\newcommand{\rhbyFLgbmPlsThreeEvP}{0.749}  % p vs BASE
\newcommand{\rhbyFLgbmPlsThreeEvR}{$-$0.18}  % recovery R
\newcommand{\rhbyFLgbmPlsThreeEvDmFull}{0.55}  % DM vs FULL
\newcommand{\rhbyFLgbmPlsThreeEvPFull}{0.583}  % p vs FULL
\newcommand{\rhbyFLgbmFullEvQlike}{0.3232}  % QLIKE lgbm full event
\newcommand{\rhbyFLgbmFullEvPct}{$-$0.62}  % QLIKE change % vs BASE
\newcommand{\rhbyFLgbmFullEvDm}{$-$0.47}  % DM vs BASE (HLN, daily means, lag 8)
\newcommand{\rhbyFLgbmFullEvP}{0.642}  % p vs BASE
\newcommand{\rhbyFLgbmFullEvR}{1.00}  % recovery R
\newcommand{\rhbyFLgbmReading}{dense beyond a few components}  % section 1 reading, ORDINARY days
\newcommand{\rhbyFLgbmRPcaTenOrd}{0.35}  % R of BASE + 10 PCs, ORDINARY
\newcommand{\rhbyFLgbmBestShortOrdR}{0.32}  % best k <= 3 arm (pls3), R
\newcommand{\rhbyFXgbBaseOrdQlike}{0.1590}  % QLIKE xgb base ordinary
\newcommand{\rhbyFXgbBaseOrdDmFull}{5.86}  % DM vs FULL
\newcommand{\rhbyFXgbBaseOrdPFull}{\ensuremath{<}0.001}  % p vs FULL
\newcommand{\rhbyFXgbPcaOneOrdQlike}{0.1587}  % QLIKE xgb pca1 ordinary
\newcommand{\rhbyFXgbPcaOneOrdPct}{$-$0.15}  % QLIKE change % vs BASE
\newcommand{\rhbyFXgbPcaOneOrdDm}{$-$1.96}  % DM vs BASE (HLN, daily means, lag 8)
\newcommand{\rhbyFXgbPcaOneOrdP}{0.049}  % p vs BASE
\newcommand{\rhbyFXgbPcaOneOrdR}{0.08}  % recovery R
\newcommand{\rhbyFXgbPcaOneOrdDmFull}{5.77}  % DM vs FULL
\newcommand{\rhbyFXgbPcaOneOrdPFull}{\ensuremath{<}0.001}  % p vs FULL
\newcommand{\rhbyFXgbPcaTwoOrdQlike}{0.1586}  % QLIKE xgb pca2 ordinary
\newcommand{\rhbyFXgbPcaTwoOrdPct}{$-$0.22}  % QLIKE change % vs BASE
\newcommand{\rhbyFXgbPcaTwoOrdDm}{$-$2.16}  % DM vs BASE (HLN, daily means, lag 8)
\newcommand{\rhbyFXgbPcaTwoOrdP}{0.030}  % p vs BASE
\newcommand{\rhbyFXgbPcaTwoOrdR}{0.12}  % recovery R
\newcommand{\rhbyFXgbPcaTwoOrdDmFull}{5.95}  % DM vs FULL
\newcommand{\rhbyFXgbPcaTwoOrdPFull}{\ensuremath{<}0.001}  % p vs FULL
\newcommand{\rhbyFXgbPcaThreeOrdQlike}{0.1586}  % QLIKE xgb pca3 ordinary
\newcommand{\rhbyFXgbPcaThreeOrdPct}{$-$0.22}  % QLIKE change % vs BASE
\newcommand{\rhbyFXgbPcaThreeOrdDm}{$-$2.17}  % DM vs BASE (HLN, daily means, lag 8)
\newcommand{\rhbyFXgbPcaThreeOrdP}{0.030}  % p vs BASE
\newcommand{\rhbyFXgbPcaThreeOrdR}{0.12}  % recovery R
\newcommand{\rhbyFXgbPcaThreeOrdDmFull}{5.60}  % DM vs FULL
\newcommand{\rhbyFXgbPcaThreeOrdPFull}{\ensuremath{<}0.001}  % p vs FULL
\newcommand{\rhbyFXgbPcaFiveOrdQlike}{0.1585}  % QLIKE xgb pca5 ordinary
\newcommand{\rhbyFXgbPcaFiveOrdPct}{$-$0.33}  % QLIKE change % vs BASE
\newcommand{\rhbyFXgbPcaFiveOrdDm}{$-$2.67}  % DM vs BASE (HLN, daily means, lag 8)
\newcommand{\rhbyFXgbPcaFiveOrdP}{0.007}  % p vs BASE
\newcommand{\rhbyFXgbPcaFiveOrdR}{0.18}  % recovery R
\newcommand{\rhbyFXgbPcaFiveOrdDmFull}{5.81}  % DM vs FULL
\newcommand{\rhbyFXgbPcaFiveOrdPFull}{\ensuremath{<}0.001}  % p vs FULL
\newcommand{\rhbyFXgbPcaTenOrdQlike}{0.1579}  % QLIKE xgb pca10 ordinary
\newcommand{\rhbyFXgbPcaTenOrdPct}{$-$0.67}  % QLIKE change % vs BASE
\newcommand{\rhbyFXgbPcaTenOrdDm}{$-$3.42}  % DM vs BASE (HLN, daily means, lag 8)
\newcommand{\rhbyFXgbPcaTenOrdP}{\ensuremath{<}0.001}  % p vs BASE
\newcommand{\rhbyFXgbPcaTenOrdR}{0.37}  % recovery R
\newcommand{\rhbyFXgbPcaTenOrdDmFull}{4.96}  % DM vs FULL
\newcommand{\rhbyFXgbPcaTenOrdPFull}{\ensuremath{<}0.001}  % p vs FULL
\newcommand{\rhbyFXgbPlsOneOrdQlike}{0.1585}  % QLIKE xgb pls1 ordinary
\newcommand{\rhbyFXgbPlsOneOrdPct}{$-$0.29}  % QLIKE change % vs BASE
\newcommand{\rhbyFXgbPlsOneOrdDm}{$-$2.78}  % DM vs BASE (HLN, daily means, lag 8)
\newcommand{\rhbyFXgbPlsOneOrdP}{0.005}  % p vs BASE
\newcommand{\rhbyFXgbPlsOneOrdR}{0.16}  % recovery R
\newcommand{\rhbyFXgbPlsOneOrdDmFull}{5.71}  % DM vs FULL
\newcommand{\rhbyFXgbPlsOneOrdPFull}{\ensuremath{<}0.001}  % p vs FULL
\newcommand{\rhbyFXgbPlsThreeOrdQlike}{0.1580}  % QLIKE xgb pls3 ordinary
\newcommand{\rhbyFXgbPlsThreeOrdPct}{$-$0.64}  % QLIKE change % vs BASE
\newcommand{\rhbyFXgbPlsThreeOrdDm}{$-$3.56}  % DM vs BASE (HLN, daily means, lag 8)
\newcommand{\rhbyFXgbPlsThreeOrdP}{\ensuremath{<}0.001}  % p vs BASE
\newcommand{\rhbyFXgbPlsThreeOrdR}{0.36}  % recovery R
\newcommand{\rhbyFXgbPlsThreeOrdDmFull}{4.41}  % DM vs FULL
\newcommand{\rhbyFXgbPlsThreeOrdPFull}{\ensuremath{<}0.001}  % p vs FULL
\newcommand{\rhbyFXgbFullOrdQlike}{0.1561}  % QLIKE xgb full ordinary
\newcommand{\rhbyFXgbFullOrdPct}{$-$1.80}  % QLIKE change % vs BASE
\newcommand{\rhbyFXgbFullOrdDm}{$-$5.86}  % DM vs BASE (HLN, daily means, lag 8)
\newcommand{\rhbyFXgbFullOrdP}{\ensuremath{<}0.001}  % p vs BASE
\newcommand{\rhbyFXgbFullOrdR}{1.00}  % recovery R
\newcommand{\rhbyFXgbBaseAllQlike}{0.1641}  % QLIKE xgb base all
\newcommand{\rhbyFXgbBaseAllDmFull}{5.26}  % DM vs FULL
\newcommand{\rhbyFXgbBaseAllPFull}{\ensuremath{<}0.001}  % p vs FULL
\newcommand{\rhbyFXgbPcaOneAllQlike}{0.1638}  % QLIKE xgb pca1 all
\newcommand{\rhbyFXgbPcaOneAllPct}{$-$0.19}  % QLIKE change % vs BASE
\newcommand{\rhbyFXgbPcaOneAllDm}{$-$2.54}  % DM vs BASE (HLN, daily means, lag 8)
\newcommand{\rhbyFXgbPcaOneAllP}{0.011}  % p vs BASE
\newcommand{\rhbyFXgbPcaOneAllR}{0.11}  % recovery R
\newcommand{\rhbyFXgbPcaOneAllDmFull}{4.97}  % DM vs FULL
\newcommand{\rhbyFXgbPcaOneAllPFull}{\ensuremath{<}0.001}  % p vs FULL
\newcommand{\rhbyFXgbPcaTwoAllQlike}{0.1639}  % QLIKE xgb pca2 all
\newcommand{\rhbyFXgbPcaTwoAllPct}{$-$0.15}  % QLIKE change % vs BASE
\newcommand{\rhbyFXgbPcaTwoAllDm}{$-$1.55}  % DM vs BASE (HLN, daily means, lag 8)
\newcommand{\rhbyFXgbPcaTwoAllP}{0.121}  % p vs BASE
\newcommand{\rhbyFXgbPcaTwoAllR}{0.09}  % recovery R
\newcommand{\rhbyFXgbPcaTwoAllDmFull}{5.53}  % DM vs FULL
\newcommand{\rhbyFXgbPcaTwoAllPFull}{\ensuremath{<}0.001}  % p vs FULL
\newcommand{\rhbyFXgbPcaThreeAllQlike}{0.1638}  % QLIKE xgb pca3 all
\newcommand{\rhbyFXgbPcaThreeAllPct}{$-$0.19}  % QLIKE change % vs BASE
\newcommand{\rhbyFXgbPcaThreeAllDm}{$-$1.95}  % DM vs BASE (HLN, daily means, lag 8)
\newcommand{\rhbyFXgbPcaThreeAllP}{0.051}  % p vs BASE
\newcommand{\rhbyFXgbPcaThreeAllR}{0.11}  % recovery R
\newcommand{\rhbyFXgbPcaThreeAllDmFull}{5.19}  % DM vs FULL
\newcommand{\rhbyFXgbPcaThreeAllPFull}{\ensuremath{<}0.001}  % p vs FULL
\newcommand{\rhbyFXgbPcaFiveAllQlike}{0.1636}  % QLIKE xgb pca5 all
\newcommand{\rhbyFXgbPcaFiveAllPct}{$-$0.29}  % QLIKE change % vs BASE
\newcommand{\rhbyFXgbPcaFiveAllDm}{$-$2.40}  % DM vs BASE (HLN, daily means, lag 8)
\newcommand{\rhbyFXgbPcaFiveAllP}{0.016}  % p vs BASE
\newcommand{\rhbyFXgbPcaFiveAllR}{0.17}  % recovery R
\newcommand{\rhbyFXgbPcaFiveAllDmFull}{5.27}  % DM vs FULL
\newcommand{\rhbyFXgbPcaFiveAllPFull}{\ensuremath{<}0.001}  % p vs FULL
\newcommand{\rhbyFXgbPcaTenAllQlike}{0.1631}  % QLIKE xgb pca10 all
\newcommand{\rhbyFXgbPcaTenAllPct}{$-$0.63}  % QLIKE change % vs BASE
\newcommand{\rhbyFXgbPcaTenAllDm}{$-$3.29}  % DM vs BASE (HLN, daily means, lag 8)
\newcommand{\rhbyFXgbPcaTenAllP}{0.001}  % p vs BASE
\newcommand{\rhbyFXgbPcaTenAllR}{0.36}  % recovery R
\newcommand{\rhbyFXgbPcaTenAllDmFull}{4.51}  % DM vs FULL
\newcommand{\rhbyFXgbPcaTenAllPFull}{\ensuremath{<}0.001}  % p vs FULL
\newcommand{\rhbyFXgbPlsOneAllQlike}{0.1637}  % QLIKE xgb pls1 all
\newcommand{\rhbyFXgbPlsOneAllPct}{$-$0.28}  % QLIKE change % vs BASE
\newcommand{\rhbyFXgbPlsOneAllDm}{$-$2.65}  % DM vs BASE (HLN, daily means, lag 8)
\newcommand{\rhbyFXgbPlsOneAllP}{0.008}  % p vs BASE
\newcommand{\rhbyFXgbPlsOneAllR}{0.16}  % recovery R
\newcommand{\rhbyFXgbPlsOneAllDmFull}{5.13}  % DM vs FULL
\newcommand{\rhbyFXgbPlsOneAllPFull}{\ensuremath{<}0.001}  % p vs FULL
\newcommand{\rhbyFXgbPlsThreeAllQlike}{0.1632}  % QLIKE xgb pls3 all
\newcommand{\rhbyFXgbPlsThreeAllPct}{$-$0.59}  % QLIKE change % vs BASE
\newcommand{\rhbyFXgbPlsThreeAllDm}{$-$3.21}  % DM vs BASE (HLN, daily means, lag 8)
\newcommand{\rhbyFXgbPlsThreeAllP}{0.001}  % p vs BASE
\newcommand{\rhbyFXgbPlsThreeAllR}{0.33}  % recovery R
\newcommand{\rhbyFXgbPlsThreeAllDmFull}{4.27}  % DM vs FULL
\newcommand{\rhbyFXgbPlsThreeAllPFull}{\ensuremath{<}0.001}  % p vs FULL
\newcommand{\rhbyFXgbFullAllQlike}{0.1612}  % QLIKE xgb full all
\newcommand{\rhbyFXgbFullAllPct}{$-$1.77}  % QLIKE change % vs BASE
\newcommand{\rhbyFXgbFullAllDm}{$-$5.26}  % DM vs BASE (HLN, daily means, lag 8)
\newcommand{\rhbyFXgbFullAllP}{\ensuremath{<}0.001}  % p vs BASE
\newcommand{\rhbyFXgbFullAllR}{1.00}  % recovery R
\newcommand{\rhbyFXgbBaseEvQlike}{0.3229}  % QLIKE xgb base event
\newcommand{\rhbyFXgbBaseEvDmFull}{0.72}  % DM vs FULL
\newcommand{\rhbyFXgbBaseEvPFull}{0.470}  % p vs FULL
\newcommand{\rhbyFXgbPcaOneEvQlike}{0.3204}  % QLIKE xgb pca1 event
\newcommand{\rhbyFXgbPcaOneEvPct}{$-$0.77}  % QLIKE change % vs BASE
\newcommand{\rhbyFXgbPcaOneEvDm}{$-$2.02}  % DM vs BASE (HLN, daily means, lag 8)
\newcommand{\rhbyFXgbPcaOneEvP}{0.044}  % p vs BASE
\newcommand{\rhbyFXgbPcaOneEvR}{0.59}  % recovery R
\newcommand{\rhbyFXgbPcaOneEvDmFull}{0.26}  % DM vs FULL
\newcommand{\rhbyFXgbPcaOneEvPFull}{0.793}  % p vs FULL
\newcommand{\rhbyFXgbPcaTwoEvQlike}{0.3255}  % QLIKE xgb pca2 event
\newcommand{\rhbyFXgbPcaTwoEvPct}{0.80}  % QLIKE change % vs BASE
\newcommand{\rhbyFXgbPcaTwoEvDm}{1.73}  % DM vs BASE (HLN, daily means, lag 8)
\newcommand{\rhbyFXgbPcaTwoEvP}{0.083}  % p vs BASE
\newcommand{\rhbyFXgbPcaTwoEvR}{$-$0.61}  % recovery R
\newcommand{\rhbyFXgbPcaTwoEvDmFull}{1.36}  % DM vs FULL
\newcommand{\rhbyFXgbPcaTwoEvPFull}{0.173}  % p vs FULL
\newcommand{\rhbyFXgbPcaThreeEvQlike}{0.3234}  % QLIKE xgb pca3 event
\newcommand{\rhbyFXgbPcaThreeEvPct}{0.17}  % QLIKE change % vs BASE
\newcommand{\rhbyFXgbPcaThreeEvDm}{0.23}  % DM vs BASE (HLN, daily means, lag 8)
\newcommand{\rhbyFXgbPcaThreeEvP}{0.817}  % p vs BASE
\newcommand{\rhbyFXgbPcaThreeEvR}{$-$0.13}  % recovery R
\newcommand{\rhbyFXgbPcaThreeEvDmFull}{0.98}  % DM vs FULL
\newcommand{\rhbyFXgbPcaThreeEvPFull}{0.326}  % p vs FULL
\newcommand{\rhbyFXgbPcaFiveEvQlike}{0.3236}  % QLIKE xgb pca5 event
\newcommand{\rhbyFXgbPcaFiveEvPct}{0.24}  % QLIKE change % vs BASE
\newcommand{\rhbyFXgbPcaFiveEvDm}{0.27}  % DM vs BASE (HLN, daily means, lag 8)
\newcommand{\rhbyFXgbPcaFiveEvP}{0.785}  % p vs BASE
\newcommand{\rhbyFXgbPcaFiveEvR}{$-$0.18}  % recovery R
\newcommand{\rhbyFXgbPcaFiveEvDmFull}{0.95}  % DM vs FULL
\newcommand{\rhbyFXgbPcaFiveEvPFull}{0.344}  % p vs FULL
\newcommand{\rhbyFXgbPcaTenEvQlike}{0.3230}  % QLIKE xgb pca10 event
\newcommand{\rhbyFXgbPcaTenEvPct}{0.04}  % QLIKE change % vs BASE
\newcommand{\rhbyFXgbPcaTenEvDm}{0.04}  % DM vs BASE (HLN, daily means, lag 8)
\newcommand{\rhbyFXgbPcaTenEvP}{0.965}  % p vs BASE
\newcommand{\rhbyFXgbPcaTenEvR}{$-$0.03}  % recovery R
\newcommand{\rhbyFXgbPcaTenEvDmFull}{0.98}  % DM vs FULL
\newcommand{\rhbyFXgbPcaTenEvPFull}{0.325}  % p vs FULL
\newcommand{\rhbyFXgbPlsOneEvQlike}{0.3223}  % QLIKE xgb pls1 event
\newcommand{\rhbyFXgbPlsOneEvPct}{$-$0.16}  % QLIKE change % vs BASE
\newcommand{\rhbyFXgbPlsOneEvDm}{$-$0.37}  % DM vs BASE (HLN, daily means, lag 8)
\newcommand{\rhbyFXgbPlsOneEvP}{0.711}  % p vs BASE
\newcommand{\rhbyFXgbPlsOneEvR}{0.12}  % recovery R
\newcommand{\rhbyFXgbPlsOneEvDmFull}{0.65}  % DM vs FULL
\newcommand{\rhbyFXgbPlsOneEvPFull}{0.517}  % p vs FULL
\newcommand{\rhbyFXgbPlsThreeEvQlike}{0.3238}  % QLIKE xgb pls3 event
\newcommand{\rhbyFXgbPlsThreeEvPct}{0.28}  % QLIKE change % vs BASE
\newcommand{\rhbyFXgbPlsThreeEvDm}{0.24}  % DM vs BASE (HLN, daily means, lag 8)
\newcommand{\rhbyFXgbPlsThreeEvP}{0.808}  % p vs BASE
\newcommand{\rhbyFXgbPlsThreeEvR}{$-$0.21}  % recovery R
\newcommand{\rhbyFXgbPlsThreeEvDmFull}{1.06}  % DM vs FULL
\newcommand{\rhbyFXgbPlsThreeEvPFull}{0.289}  % p vs FULL
\newcommand{\rhbyFXgbFullEvQlike}{0.3186}  % QLIKE xgb full event
\newcommand{\rhbyFXgbFullEvPct}{$-$1.31}  % QLIKE change % vs BASE
\newcommand{\rhbyFXgbFullEvDm}{$-$0.72}  % DM vs BASE (HLN, daily means, lag 8)
\newcommand{\rhbyFXgbFullEvP}{0.470}  % p vs BASE
\newcommand{\rhbyFXgbFullEvR}{1.00}  % recovery R
\newcommand{\rhbyFXgbReading}{dense beyond a few components}  % section 1 reading, ORDINARY days
\newcommand{\rhbyFXgbRPcaTenOrd}{0.37}  % R of BASE + 10 PCs, ORDINARY
\newcommand{\rhbyFXgbBestShortOrdR}{0.36}  % best k <= 3 arm (pls3), R
\newcommand{\rhbyFAgreement}{agree}  % the two models' readings
\newcommand{\rhbyFLabel}{requested by the author after the beyond-base results; not pre-registered}
\fi
\begin{itemize}
\item \textbf{Two groups carry the gain.} The FOMC calendar and the futures' realized moments, added to the HAR
  benchmark set, deliver at least \rhabBaseShareMin\% of every class's every-bar gain from the full set:
  \rhabBaseShareTreeMin--\rhabBaseShareTreeMax\% in the tree ensembles and at least \rhabBaseShareLinMin\% in the
  penalized linear models (Table~\ref{tab:rh_ablation}, Panel~D). The calendar carries the statement days
  (Section~\ref{sec:rh_fi_event}), the realized moments the ordinary days (Section~\ref{sec:rh_fi_ordinary}). What
  remains is small, only the boosted trees capture it, and a few series carry it (Section~\ref{sec:rh_fi_rest}). How
% 2026-10-28 (author, alternative-data edit: the claim is about the 36 series beyond HAR + RM + FOMC). Was:
%   remains is small, only the tree ensembles capture it, and a few series carry it
  much a model uses a series is a poor guide to these gains (Section~\ref{sec:rh_fi_use}).
\item \textbf{Evidence.} Three kinds, each reported for every class and predictor group.
  \begin{itemize}
  \item Accuracy: the predictor-group ablations of Table~\ref{tab:rh_ablation}, with the design of
    Table~\ref{tab:rh_oos}, including the full set's frozen hyperparameters, and only the columns changed
    (Appendix~\ref{app:rh_results_ablation}). Panel~A removes one of the \rhfiNBuckets{} predictor groups from the full
    set and Panel~B adds one to the HAR benchmark set; Panel~C fits by OLS the semivariance HAR of
    \citet{PattonSheppard2015}, the HARQ model of \citet{BollerslevPattonQuaedvlieg2016} and the leverage HAR of
    \citet{CorsiReno2012}, built from the same one-minute E-mini returns as the target; Panel~D sets the full set
    against the HAR benchmark set, alone and with the realized moments, the FOMC calendar or both added
    (HAR\,+\,RM\,+\,FOMC); Panel~E scores the late-starting groups on the bars on which their values are live.
    Table~\ref{tab:rh_channels_abl} scores the additions and removals on each day type.
  \item Concentration: the lasso, which selects, against the ridge and the elastic net, which keep and shrink every
    series \citep[see][]{GiannoneLenzaPrimiceri2021}; the feature-count ladder of
    \citet{MaischZhangUlrichZimmer2026} (Table~\ref{tab:rh_channels_sparse}; Appendix~\ref{app:rh_fi_ladder}); and a
    factor-augmented HAR \citep{StockWatson2002,KellyPruitt2015} on ordinary days (Table~\ref{tab:rh_factor}).
  \item Use: each model's attribution, its TreeSHAP shares \citep{lundberg2020local} or standardized coefficient
    shares (Appendices~\ref{app:rh_fi_compute} and~\ref{app:rh_fi_rank}). A share measures how much a series moves
    the forecast, not how much it improves it.
  \end{itemize}
\item \textbf{Reading.} Fixed before the day types were scored (Appendix~\ref{app:rh_channels}). A part of the gain is
  concentrated when one group added alone recovers at least three-quarters of it and its removal raises the loss
  significantly; it is spread across correlated proxies when the groups' separate gains sum to more than one and a
  half times the joint gain and the removal of at most one group raises the loss significantly. The two can hold at
  once, when one group carries the gain and the others repeat its information; we report both.
  % 2026-10-28 (coordinator): "at least one and a half times" -> "more than", as RULE.txt and the table notes.
\end{itemize}
% 2026-10-28 (coordinator, v2): Section 7's tables, input here (not in 08_feature_importance_floats.tex) so that they
% print within the section. 2026-10-28 (one Results section): moved to just before 6.3.2, so that they follow 6.3.1
% instead of three float pages in a row. Was: the three \input lines here.

% 2026-10-28 (author: Sections 6 and 7 are one section). Was: \subsection{The FOMC calendar: one group, few bars}
% 2026-10-28 (author: conventional titles). Was: \subsubsection{The FOMC calendar: one group, few bars}
\subsubsection{FOMC calendar}
\label{sec:rh_fi_event}
\begin{itemize}
\item \textbf{It carries the statement days} (Tables~\ref{tab:rh_ablation} and~\ref{tab:rh_channels_abl}). On every
  bar the FOMC series, added alone to the HAR benchmark set, deliver
  \rhabTreeFomcGainShareMin--\rhabTreeFomcGainShareMax\% of the trees' gain, and removing them from the full set
  raises the loss by \rhabLoFomcAllPctMin--\rhabLoFomcAllPctMax\%. On statement days they recover
  \rhchEvFomcShareLinRng\% of the full set's gain in the linear models and \rhchEvFomcShareTreeRng\% in the trees, and
  removing them raises the loss by \rhchEvLoFomcRng\%. The reading is concentrated for both model types; no other
  group recovers more than \rhchEvOtherMaxShare\% of it.
\item \textbf{The lasso wins.} On statement days the lasso's full-set loss differs from the ridge's by
  \rhchLassoRidgeEventPct\% ($\DM={}$\rhchLassoRidgeEventDm) and from the elastic net's by \rhchLassoEnetEventPct\%
  ($\DM={}$\rhchLassoEnetEventDm), as a signal in one series predicts; the lasso's top-ranked series, \rhldLassoTopOne,
  is the proximity to the next statement.
\item \textbf{The models use it on few bars.} Over the span the FOMC calendar holds only
  \rhfipFomcBucketMin--\rhfipFomcBucketMax\% of the trees' attribution; on statement days \rhfipFomcUntil{} holds
  \rhfipLgbmUntilFomc\% of LightGBM's, on other days \rhfipLgbmUntilOther\% (Appendix~\ref{app:rh_fi_cond}).
\item \textbf{A cost on ordinary days for the linear models.} A linear term in the FOMC series also moves the forecasts
  of ordinary days: added alone, the series raise the ordinary-day loss of the penalized linear models by
  \rhchOrdAddFomcLinRng\%, while the trees, which can confine the series' effect to the statement bars, pay no such
  cost (\rhchOrdAddFomcTreeRng\%).
\end{itemize}

% AUTO-GENERATED by writeup/make_rehaul_channels.py -- do not edit; re-run the script.
\ifdefined\rhchNBars\else% AUTO-GENERATED by writeup/make_rehaul_channels.py -- do not edit; re-run the script.
% \rhch... macros: the two channels of the pooled gain (results/channels_2026-10-26/RULE.txt items 1-3):
% FOMC statement days (EVENT) and ordinary days, Table 1's scored bars, the FOMC-patched design of record;
% DM (HLN) on daily average loss, lag 8.  Names in the generator's docstring.
% \rhchX... = ADDITIONAL, NOT PRE-REGISTERED (coordinator, 2026-10-26, after the first scoring): the full set
% against HAR + the futures' realized moments (FullVsRm) and HAR + realized moments against HAR (RmVsHar).
\providecommand{\rhPending}{\textit{pending}}
\providecommand{\rhchNBars}{}\renewcommand{\rhchNBars}{33{,}527}  % scored bars (Table 1)
\providecommand{\rhchNDays}{}\renewcommand{\rhchNDays}{2{,}579}  % scored days
\providecommand{\rhchNEventDays}{}\renewcommand{\rhchNEventDays}{81}  % EVENT days: scheduled FOMC statement days of the span
\providecommand{\rhchNOrdDays}{}\renewcommand{\rhchNOrdDays}{2{,}498}  % ORDINARY days
\providecommand{\rhchNEventBars}{}\renewcommand{\rhchNEventBars}{1{,}053}  % bars on EVENT days
\providecommand{\rhchEventBarPct}{}\renewcommand{\rhchEventBarPct}{3.1}  % EVENT-day bars, % of the scored bars
\providecommand{\rhchHacLag}{}\renewcommand{\rhchHacLag}{8}  % Newey-West lag of every DM test
\providecommand{\rhchAgreePp}{}\renewcommand{\rhchAgreePp}{2}  % item 1: the corrections agree within this many percentage points
\providecommand{\rhchLadderDesign}{}\renewcommand{\rhchLadderDesign}{patched}  % design of the ladder
\providecommand{\rhchRidgeTabOnePct}{}\renewcommand{\rhchRidgeTabOnePct}{$-$6.0}  % ridge: full vs HAR, every bar (Table 1): % QLIKE change
\providecommand{\rhchRidgeTabOneDm}{}\renewcommand{\rhchRidgeTabOneDm}{$-$6.28}  % ridge: full vs HAR, every bar (Table 1): DM
\providecommand{\rhchRidgeTabOneP}{}\renewcommand{\rhchRidgeTabOneP}{\ensuremath{<}0.001}  % ridge: full vs HAR, every bar (Table 1): DM p
\providecommand{\rhchRidgeTabOneMsePct}{}\renewcommand{\rhchRidgeTabOneMsePct}{100.6}  % ridge: full vs HAR, every bar (Table 1): % MSE change
\providecommand{\rhchRidgeTabOneMseP}{}\renewcommand{\rhchRidgeTabOneMseP}{0.261}  % ridge: full vs HAR, every bar (Table 1): DM p on MSE
\providecommand{\rhchRidgeCorrAPct}{}\renewcommand{\rhchRidgeCorrAPct}{$-$3.0}  % ridge: full vs HAR+FOMC, every bar: % QLIKE change
\providecommand{\rhchRidgeCorrADm}{}\renewcommand{\rhchRidgeCorrADm}{$-$3.64}  % ridge: full vs HAR+FOMC, every bar: DM
\providecommand{\rhchRidgeCorrAP}{}\renewcommand{\rhchRidgeCorrAP}{\ensuremath{<}0.001}  % ridge: full vs HAR+FOMC, every bar: DM p
\providecommand{\rhchRidgeCorrAMsePct}{}\renewcommand{\rhchRidgeCorrAMsePct}{100.8}  % ridge: full vs HAR+FOMC, every bar: % MSE change
\providecommand{\rhchRidgeCorrAMseP}{}\renewcommand{\rhchRidgeCorrAMseP}{0.258}  % ridge: full vs HAR+FOMC, every bar: DM p on MSE
\providecommand{\rhchRidgeCorrBPct}{}\renewcommand{\rhchRidgeCorrBPct}{$-$3.2}  % ridge: full without FOMC vs HAR, every bar: % QLIKE change
\providecommand{\rhchRidgeCorrBDm}{}\renewcommand{\rhchRidgeCorrBDm}{$-$3.74}  % ridge: full without FOMC vs HAR, every bar: DM
\providecommand{\rhchRidgeCorrBP}{}\renewcommand{\rhchRidgeCorrBP}{\ensuremath{<}0.001}  % ridge: full without FOMC vs HAR, every bar: DM p
\providecommand{\rhchRidgeCorrBMsePct}{}\renewcommand{\rhchRidgeCorrBMsePct}{102.3}  % ridge: full without FOMC vs HAR, every bar: % MSE change
\providecommand{\rhchRidgeCorrBMseP}{}\renewcommand{\rhchRidgeCorrBMseP}{0.258}  % ridge: full without FOMC vs HAR, every bar: DM p on MSE
\providecommand{\rhchRidgeCorrCPct}{}\renewcommand{\rhchRidgeCorrCPct}{$-$3.2}  % ridge: full vs HAR, ORDINARY days: % QLIKE change
\providecommand{\rhchRidgeCorrCDm}{}\renewcommand{\rhchRidgeCorrCDm}{$-$3.51}  % ridge: full vs HAR, ORDINARY days: DM
\providecommand{\rhchRidgeCorrCP}{}\renewcommand{\rhchRidgeCorrCP}{\ensuremath{<}0.001}  % ridge: full vs HAR, ORDINARY days: DM p
\providecommand{\rhchRidgeCorrCMsePct}{}\renewcommand{\rhchRidgeCorrCMsePct}{110.2}  % ridge: full vs HAR, ORDINARY days: % MSE change
\providecommand{\rhchRidgeCorrCMseP}{}\renewcommand{\rhchRidgeCorrCMseP}{0.258}  % ridge: full vs HAR, ORDINARY days: DM p on MSE
\providecommand{\rhchRidgeCorrAgree}{}\renewcommand{\rhchRidgeCorrAgree}{agree}  % ridge: the three corrections within 2.0 pp?
\providecommand{\rhchRidgeCorrSpread}{}\renewcommand{\rhchRidgeCorrSpread}{0.2}  % ridge: max - min of the three (pp)
\providecommand{\rhchLassoTabOnePct}{}\renewcommand{\rhchLassoTabOnePct}{$-$9.0}  % lasso: full vs HAR, every bar (Table 1): % QLIKE change
\providecommand{\rhchLassoTabOneDm}{}\renewcommand{\rhchLassoTabOneDm}{$-$6.34}  % lasso: full vs HAR, every bar (Table 1): DM
\providecommand{\rhchLassoTabOneP}{}\renewcommand{\rhchLassoTabOneP}{\ensuremath{<}0.001}  % lasso: full vs HAR, every bar (Table 1): DM p
\providecommand{\rhchLassoTabOneMsePct}{}\renewcommand{\rhchLassoTabOneMsePct}{114.7}  % lasso: full vs HAR, every bar (Table 1): % MSE change
\providecommand{\rhchLassoTabOneMseP}{}\renewcommand{\rhchLassoTabOneMseP}{0.247}  % lasso: full vs HAR, every bar (Table 1): DM p on MSE
\providecommand{\rhchLassoCorrAPct}{}\renewcommand{\rhchLassoCorrAPct}{$-$3.4}  % lasso: full vs HAR+FOMC, every bar: % QLIKE change
\providecommand{\rhchLassoCorrADm}{}\renewcommand{\rhchLassoCorrADm}{$-$3.20}  % lasso: full vs HAR+FOMC, every bar: DM
\providecommand{\rhchLassoCorrAP}{}\renewcommand{\rhchLassoCorrAP}{0.001}  % lasso: full vs HAR+FOMC, every bar: DM p
\providecommand{\rhchLassoCorrAMsePct}{}\renewcommand{\rhchLassoCorrAMsePct}{114.4}  % lasso: full vs HAR+FOMC, every bar: % MSE change
\providecommand{\rhchLassoCorrAMseP}{}\renewcommand{\rhchLassoCorrAMseP}{0.243}  % lasso: full vs HAR+FOMC, every bar: DM p on MSE
\providecommand{\rhchLassoCorrBPct}{}\renewcommand{\rhchLassoCorrBPct}{$-$2.8}  % lasso: full without FOMC vs HAR, every bar: % QLIKE change
\providecommand{\rhchLassoCorrBDm}{}\renewcommand{\rhchLassoCorrBDm}{$-$2.75}  % lasso: full without FOMC vs HAR, every bar: DM
\providecommand{\rhchLassoCorrBP}{}\renewcommand{\rhchLassoCorrBP}{0.006}  % lasso: full without FOMC vs HAR, every bar: DM p
\providecommand{\rhchLassoCorrBMsePct}{}\renewcommand{\rhchLassoCorrBMsePct}{115.3}  % lasso: full without FOMC vs HAR, every bar: % MSE change
\providecommand{\rhchLassoCorrBMseP}{}\renewcommand{\rhchLassoCorrBMseP}{0.245}  % lasso: full without FOMC vs HAR, every bar: DM p on MSE
\providecommand{\rhchLassoCorrCPct}{}\renewcommand{\rhchLassoCorrCPct}{$-$2.4}  % lasso: full vs HAR, ORDINARY days: % QLIKE change
\providecommand{\rhchLassoCorrCDm}{}\renewcommand{\rhchLassoCorrCDm}{$-$2.08}  % lasso: full vs HAR, ORDINARY days: DM
\providecommand{\rhchLassoCorrCP}{}\renewcommand{\rhchLassoCorrCP}{0.037}  % lasso: full vs HAR, ORDINARY days: DM p
\providecommand{\rhchLassoCorrCMsePct}{}\renewcommand{\rhchLassoCorrCMsePct}{126.0}  % lasso: full vs HAR, ORDINARY days: % MSE change
\providecommand{\rhchLassoCorrCMseP}{}\renewcommand{\rhchLassoCorrCMseP}{0.243}  % lasso: full vs HAR, ORDINARY days: DM p on MSE
\providecommand{\rhchLassoCorrAgree}{}\renewcommand{\rhchLassoCorrAgree}{agree}  % lasso: the three corrections within 2.0 pp?
\providecommand{\rhchLassoCorrSpread}{}\renewcommand{\rhchLassoCorrSpread}{1.0}  % lasso: max - min of the three (pp)
\providecommand{\rhchEnetTabOnePct}{}\renewcommand{\rhchEnetTabOnePct}{$-$6.7}  % enet: full vs HAR, every bar (Table 1): % QLIKE change
\providecommand{\rhchEnetTabOneDm}{}\renewcommand{\rhchEnetTabOneDm}{$-$7.36}  % enet: full vs HAR, every bar (Table 1): DM
\providecommand{\rhchEnetTabOneP}{}\renewcommand{\rhchEnetTabOneP}{\ensuremath{<}0.001}  % enet: full vs HAR, every bar (Table 1): DM p
\providecommand{\rhchEnetTabOneMsePct}{}\renewcommand{\rhchEnetTabOneMsePct}{117.7}  % enet: full vs HAR, every bar (Table 1): % MSE change
\providecommand{\rhchEnetTabOneMseP}{}\renewcommand{\rhchEnetTabOneMseP}{0.263}  % enet: full vs HAR, every bar (Table 1): DM p on MSE
\providecommand{\rhchEnetCorrAPct}{}\renewcommand{\rhchEnetCorrAPct}{$-$3.7}  % enet: full vs HAR+FOMC, every bar: % QLIKE change
\providecommand{\rhchEnetCorrADm}{}\renewcommand{\rhchEnetCorrADm}{$-$4.89}  % enet: full vs HAR+FOMC, every bar: DM
\providecommand{\rhchEnetCorrAP}{}\renewcommand{\rhchEnetCorrAP}{\ensuremath{<}0.001}  % enet: full vs HAR+FOMC, every bar: DM p
\providecommand{\rhchEnetCorrAMsePct}{}\renewcommand{\rhchEnetCorrAMsePct}{117.8}  % enet: full vs HAR+FOMC, every bar: % MSE change
\providecommand{\rhchEnetCorrAMseP}{}\renewcommand{\rhchEnetCorrAMseP}{0.261}  % enet: full vs HAR+FOMC, every bar: DM p on MSE
\providecommand{\rhchEnetCorrBPct}{}\renewcommand{\rhchEnetCorrBPct}{$-$3.8}  % enet: full without FOMC vs HAR, every bar: % QLIKE change
\providecommand{\rhchEnetCorrBDm}{}\renewcommand{\rhchEnetCorrBDm}{$-$5.00}  % enet: full without FOMC vs HAR, every bar: DM
\providecommand{\rhchEnetCorrBP}{}\renewcommand{\rhchEnetCorrBP}{\ensuremath{<}0.001}  % enet: full without FOMC vs HAR, every bar: DM p
\providecommand{\rhchEnetCorrBMsePct}{}\renewcommand{\rhchEnetCorrBMsePct}{117.8}  % enet: full without FOMC vs HAR, every bar: % MSE change
\providecommand{\rhchEnetCorrBMseP}{}\renewcommand{\rhchEnetCorrBMseP}{0.259}  % enet: full without FOMC vs HAR, every bar: DM p on MSE
\providecommand{\rhchEnetCorrCPct}{}\renewcommand{\rhchEnetCorrCPct}{$-$4.0}  % enet: full vs HAR, ORDINARY days: % QLIKE change
\providecommand{\rhchEnetCorrCDm}{}\renewcommand{\rhchEnetCorrCDm}{$-$4.73}  % enet: full vs HAR, ORDINARY days: DM
\providecommand{\rhchEnetCorrCP}{}\renewcommand{\rhchEnetCorrCP}{\ensuremath{<}0.001}  % enet: full vs HAR, ORDINARY days: DM p
\providecommand{\rhchEnetCorrCMsePct}{}\renewcommand{\rhchEnetCorrCMsePct}{128.9}  % enet: full vs HAR, ORDINARY days: % MSE change
\providecommand{\rhchEnetCorrCMseP}{}\renewcommand{\rhchEnetCorrCMseP}{0.260}  % enet: full vs HAR, ORDINARY days: DM p on MSE
\providecommand{\rhchEnetCorrAgree}{}\renewcommand{\rhchEnetCorrAgree}{agree}  % enet: the three corrections within 2.0 pp?
\providecommand{\rhchEnetCorrSpread}{}\renewcommand{\rhchEnetCorrSpread}{0.2}  % enet: max - min of the three (pp)
\providecommand{\rhchLgbmTabOnePct}{}\renewcommand{\rhchLgbmTabOnePct}{$-$14.1}  % lgbm: full vs HAR, every bar (Table 1): % QLIKE change
\providecommand{\rhchLgbmTabOneDm}{}\renewcommand{\rhchLgbmTabOneDm}{$-$10.26}  % lgbm: full vs HAR, every bar (Table 1): DM
\providecommand{\rhchLgbmTabOneP}{}\renewcommand{\rhchLgbmTabOneP}{\ensuremath{<}0.001}  % lgbm: full vs HAR, every bar (Table 1): DM p
\providecommand{\rhchLgbmTabOneMsePct}{}\renewcommand{\rhchLgbmTabOneMsePct}{20.0}  % lgbm: full vs HAR, every bar (Table 1): % MSE change
\providecommand{\rhchLgbmTabOneMseP}{}\renewcommand{\rhchLgbmTabOneMseP}{0.301}  % lgbm: full vs HAR, every bar (Table 1): DM p on MSE
\providecommand{\rhchLgbmCorrAPct}{}\renewcommand{\rhchLgbmCorrAPct}{$-$5.4}  % lgbm: full vs HAR+FOMC, every bar: % QLIKE change
\providecommand{\rhchLgbmCorrADm}{}\renewcommand{\rhchLgbmCorrADm}{$-$11.63}  % lgbm: full vs HAR+FOMC, every bar: DM
\providecommand{\rhchLgbmCorrAP}{}\renewcommand{\rhchLgbmCorrAP}{\ensuremath{<}0.001}  % lgbm: full vs HAR+FOMC, every bar: DM p
\providecommand{\rhchLgbmCorrAMsePct}{}\renewcommand{\rhchLgbmCorrAMsePct}{19.4}  % lgbm: full vs HAR+FOMC, every bar: % MSE change
\providecommand{\rhchLgbmCorrAMseP}{}\renewcommand{\rhchLgbmCorrAMseP}{0.276}  % lgbm: full vs HAR+FOMC, every bar: DM p on MSE
\providecommand{\rhchLgbmCorrBPct}{}\renewcommand{\rhchLgbmCorrBPct}{$-$5.3}  % lgbm: full without FOMC vs HAR, every bar: % QLIKE change
\providecommand{\rhchLgbmCorrBDm}{}\renewcommand{\rhchLgbmCorrBDm}{$-$9.62}  % lgbm: full without FOMC vs HAR, every bar: DM
\providecommand{\rhchLgbmCorrBP}{}\renewcommand{\rhchLgbmCorrBP}{\ensuremath{<}0.001}  % lgbm: full without FOMC vs HAR, every bar: DM p
\providecommand{\rhchLgbmCorrBMsePct}{}\renewcommand{\rhchLgbmCorrBMsePct}{21.6}  % lgbm: full without FOMC vs HAR, every bar: % MSE change
\providecommand{\rhchLgbmCorrBMseP}{}\renewcommand{\rhchLgbmCorrBMseP}{0.269}  % lgbm: full without FOMC vs HAR, every bar: DM p on MSE
\providecommand{\rhchLgbmCorrCPct}{}\renewcommand{\rhchLgbmCorrCPct}{$-$6.0}  % lgbm: full vs HAR, ORDINARY days: % QLIKE change
\providecommand{\rhchLgbmCorrCDm}{}\renewcommand{\rhchLgbmCorrCDm}{$-$10.70}  % lgbm: full vs HAR, ORDINARY days: DM
\providecommand{\rhchLgbmCorrCP}{}\renewcommand{\rhchLgbmCorrCP}{\ensuremath{<}0.001}  % lgbm: full vs HAR, ORDINARY days: DM p
\providecommand{\rhchLgbmCorrCMsePct}{}\renewcommand{\rhchLgbmCorrCMsePct}{22.8}  % lgbm: full vs HAR, ORDINARY days: % MSE change
\providecommand{\rhchLgbmCorrCMseP}{}\renewcommand{\rhchLgbmCorrCMseP}{0.275}  % lgbm: full vs HAR, ORDINARY days: DM p on MSE
\providecommand{\rhchLgbmCorrAgree}{}\renewcommand{\rhchLgbmCorrAgree}{agree}  % lgbm: the three corrections within 2.0 pp?
\providecommand{\rhchLgbmCorrSpread}{}\renewcommand{\rhchLgbmCorrSpread}{0.8}  % lgbm: max - min of the three (pp)
\providecommand{\rhchXgbTabOnePct}{}\renewcommand{\rhchXgbTabOnePct}{$-$14.6}  % xgb: full vs HAR, every bar (Table 1): % QLIKE change
\providecommand{\rhchXgbTabOneDm}{}\renewcommand{\rhchXgbTabOneDm}{$-$9.94}  % xgb: full vs HAR, every bar (Table 1): DM
\providecommand{\rhchXgbTabOneP}{}\renewcommand{\rhchXgbTabOneP}{\ensuremath{<}0.001}  % xgb: full vs HAR, every bar (Table 1): DM p
\providecommand{\rhchXgbTabOneMsePct}{}\renewcommand{\rhchXgbTabOneMsePct}{16.3}  % xgb: full vs HAR, every bar (Table 1): % MSE change
\providecommand{\rhchXgbTabOneMseP}{}\renewcommand{\rhchXgbTabOneMseP}{0.301}  % xgb: full vs HAR, every bar (Table 1): DM p on MSE
\providecommand{\rhchXgbCorrAPct}{}\renewcommand{\rhchXgbCorrAPct}{$-$5.7}  % xgb: full vs HAR+FOMC, every bar: % QLIKE change
\providecommand{\rhchXgbCorrADm}{}\renewcommand{\rhchXgbCorrADm}{$-$12.60}  % xgb: full vs HAR+FOMC, every bar: DM
\providecommand{\rhchXgbCorrAP}{}\renewcommand{\rhchXgbCorrAP}{\ensuremath{<}0.001}  % xgb: full vs HAR+FOMC, every bar: DM p
\providecommand{\rhchXgbCorrAMsePct}{}\renewcommand{\rhchXgbCorrAMsePct}{15.4}  % xgb: full vs HAR+FOMC, every bar: % MSE change
\providecommand{\rhchXgbCorrAMseP}{}\renewcommand{\rhchXgbCorrAMseP}{0.282}  % xgb: full vs HAR+FOMC, every bar: DM p on MSE
\providecommand{\rhchXgbCorrBPct}{}\renewcommand{\rhchXgbCorrBPct}{$-$5.8}  % xgb: full without FOMC vs HAR, every bar: % QLIKE change
\providecommand{\rhchXgbCorrBDm}{}\renewcommand{\rhchXgbCorrBDm}{$-$10.81}  % xgb: full without FOMC vs HAR, every bar: DM
\providecommand{\rhchXgbCorrBP}{}\renewcommand{\rhchXgbCorrBP}{\ensuremath{<}0.001}  % xgb: full without FOMC vs HAR, every bar: DM p
\providecommand{\rhchXgbCorrBMsePct}{}\renewcommand{\rhchXgbCorrBMsePct}{17.9}  % xgb: full without FOMC vs HAR, every bar: % MSE change
\providecommand{\rhchXgbCorrBMseP}{}\renewcommand{\rhchXgbCorrBMseP}{0.268}  % xgb: full without FOMC vs HAR, every bar: DM p on MSE
\providecommand{\rhchXgbCorrCPct}{}\renewcommand{\rhchXgbCorrCPct}{$-$6.1}  % xgb: full vs HAR, ORDINARY days: % QLIKE change
\providecommand{\rhchXgbCorrCDm}{}\renewcommand{\rhchXgbCorrCDm}{$-$10.42}  % xgb: full vs HAR, ORDINARY days: DM
\providecommand{\rhchXgbCorrCP}{}\renewcommand{\rhchXgbCorrCP}{\ensuremath{<}0.001}  % xgb: full vs HAR, ORDINARY days: DM p
\providecommand{\rhchXgbCorrCMsePct}{}\renewcommand{\rhchXgbCorrCMsePct}{18.9}  % xgb: full vs HAR, ORDINARY days: % MSE change
\providecommand{\rhchXgbCorrCMseP}{}\renewcommand{\rhchXgbCorrCMseP}{0.268}  % xgb: full vs HAR, ORDINARY days: DM p on MSE
\providecommand{\rhchXgbCorrAgree}{}\renewcommand{\rhchXgbCorrAgree}{agree}  % xgb: the three corrections within 2.0 pp?
\providecommand{\rhchXgbCorrSpread}{}\renewcommand{\rhchXgbCorrSpread}{0.4}  % xgb: max - min of the three (pp)
\providecommand{\rhchRfTabOnePct}{}\renewcommand{\rhchRfTabOnePct}{$-$11.3}  % rf: full vs HAR, every bar (Table 1): % QLIKE change
\providecommand{\rhchRfTabOneDm}{}\renewcommand{\rhchRfTabOneDm}{$-$9.74}  % rf: full vs HAR, every bar (Table 1): DM
\providecommand{\rhchRfTabOneP}{}\renewcommand{\rhchRfTabOneP}{\ensuremath{<}0.001}  % rf: full vs HAR, every bar (Table 1): DM p
\providecommand{\rhchRfTabOneMsePct}{}\renewcommand{\rhchRfTabOneMsePct}{10.9}  % rf: full vs HAR, every bar (Table 1): % MSE change
\providecommand{\rhchRfTabOneMseP}{}\renewcommand{\rhchRfTabOneMseP}{0.271}  % rf: full vs HAR, every bar (Table 1): DM p on MSE
\providecommand{\rhchRfCorrAPct}{}\renewcommand{\rhchRfCorrAPct}{$-$4.0}  % rf: full vs HAR+FOMC, every bar: % QLIKE change
\providecommand{\rhchRfCorrADm}{}\renewcommand{\rhchRfCorrADm}{$-$9.06}  % rf: full vs HAR+FOMC, every bar: DM
\providecommand{\rhchRfCorrAP}{}\renewcommand{\rhchRfCorrAP}{\ensuremath{<}0.001}  % rf: full vs HAR+FOMC, every bar: DM p
\providecommand{\rhchRfCorrAMsePct}{}\renewcommand{\rhchRfCorrAMsePct}{11.3}  % rf: full vs HAR+FOMC, every bar: % MSE change
\providecommand{\rhchRfCorrAMseP}{}\renewcommand{\rhchRfCorrAMseP}{0.255}  % rf: full vs HAR+FOMC, every bar: DM p on MSE
\providecommand{\rhchRfCorrBPct}{}\renewcommand{\rhchRfCorrBPct}{--}  % rf: full without FOMC vs HAR, every bar: not run
\providecommand{\rhchRfCorrBDm}{}\renewcommand{\rhchRfCorrBDm}{--}  % rf: full without FOMC vs HAR, every bar: not run
\providecommand{\rhchRfCorrBP}{}\renewcommand{\rhchRfCorrBP}{--}  % rf: full without FOMC vs HAR, every bar: not run
\providecommand{\rhchRfCorrBMsePct}{}\renewcommand{\rhchRfCorrBMsePct}{--}  % rf: full without FOMC vs HAR, every bar: not run
\providecommand{\rhchRfCorrBMseP}{}\renewcommand{\rhchRfCorrBMseP}{--}  % rf: full without FOMC vs HAR, every bar: not run
\providecommand{\rhchRfCorrCPct}{}\renewcommand{\rhchRfCorrCPct}{$-$5.3}  % rf: full vs HAR, ORDINARY days: % QLIKE change
\providecommand{\rhchRfCorrCDm}{}\renewcommand{\rhchRfCorrCDm}{$-$11.16}  % rf: full vs HAR, ORDINARY days: DM
\providecommand{\rhchRfCorrCP}{}\renewcommand{\rhchRfCorrCP}{\ensuremath{<}0.001}  % rf: full vs HAR, ORDINARY days: DM p
\providecommand{\rhchRfCorrCMsePct}{}\renewcommand{\rhchRfCorrCMsePct}{12.4}  % rf: full vs HAR, ORDINARY days: % MSE change
\providecommand{\rhchRfCorrCMseP}{}\renewcommand{\rhchRfCorrCMseP}{0.247}  % rf: full vs HAR, ORDINARY days: DM p on MSE
\providecommand{\rhchRfCorrAgree}{}\renewcommand{\rhchRfCorrAgree}{--}  % rf: the three corrections within 2.0 pp?
\providecommand{\rhchRfCorrSpread}{}\renewcommand{\rhchRfCorrSpread}{1.3}  % rf: max - min of the three (pp)
\providecommand{\rhchNCorrAgree}{}\renewcommand{\rhchNCorrAgree}{5}  % classes whose corrections agree
\providecommand{\rhchNCorrDisagree}{}\renewcommand{\rhchNCorrDisagree}{0}  % classes whose corrections disagree
\providecommand{\rhchNCorrPending}{}\renewcommand{\rhchNCorrPending}{0}  % classes with a correction pending
\providecommand{\rhchRidgeShareEvent}{}\renewcommand{\rhchRidgeShareEvent}{54}  % ridge: % of the every-bar gain on EVENT days
\providecommand{\rhchRidgeShareEventPm}{}\renewcommand{\rhchRidgeShareEventPm}{56}  % ridge: % of the gain on EVENT days 14:30-16:00
\providecommand{\rhchRidgeEventPeakBar}{}\renewcommand{\rhchRidgeEventPeakBar}{14:30}  % ridge: event: bar end with the most negative mean diff
\providecommand{\rhchRidgeEventPeakDiff}{}\renewcommand{\rhchRidgeEventPeakDiff}{$-$1583.6}  % ridge: event: that mean diff (full - HAR) x 10^3
\providecommand{\rhchRidgeOrdPeakBar}{}\renewcommand{\rhchRidgeOrdPeakBar}{14:30}  % ridge: ord: bar end with the most negative mean diff
\providecommand{\rhchRidgeOrdPeakDiff}{}\renewcommand{\rhchRidgeOrdPeakDiff}{$-$15.9}  % ridge: ord: that mean diff (full - HAR) x 10^3
\providecommand{\rhchLassoShareEvent}{}\renewcommand{\rhchLassoShareEvent}{77}  % lasso: % of the every-bar gain on EVENT days
\providecommand{\rhchLassoShareEventPm}{}\renewcommand{\rhchLassoShareEventPm}{78}  % lasso: % of the gain on EVENT days 14:30-16:00
\providecommand{\rhchLassoEventPeakBar}{}\renewcommand{\rhchLassoEventPeakBar}{14:30}  % lasso: event: bar end with the most negative mean diff
\providecommand{\rhchLassoEventPeakDiff}{}\renewcommand{\rhchLassoEventPeakDiff}{$-$4302.3}  % lasso: event: that mean diff (full - HAR) x 10^3
\providecommand{\rhchLassoOrdPeakBar}{}\renewcommand{\rhchLassoOrdPeakBar}{14:30}  % lasso: ord: bar end with the most negative mean diff
\providecommand{\rhchLassoOrdPeakDiff}{}\renewcommand{\rhchLassoOrdPeakDiff}{$-$16.4}  % lasso: ord: that mean diff (full - HAR) x 10^3
\providecommand{\rhchEnetShareEvent}{}\renewcommand{\rhchEnetShareEvent}{49}  % enet: % of the every-bar gain on EVENT days
\providecommand{\rhchEnetShareEventPm}{}\renewcommand{\rhchEnetShareEventPm}{51}  % enet: % of the gain on EVENT days 14:30-16:00
\providecommand{\rhchEnetEventPeakBar}{}\renewcommand{\rhchEnetEventPeakBar}{14:30}  % enet: event: bar end with the most negative mean diff
\providecommand{\rhchEnetEventPeakDiff}{}\renewcommand{\rhchEnetEventPeakDiff}{$-$1610.3}  % enet: event: that mean diff (full - HAR) x 10^3
\providecommand{\rhchEnetOrdPeakBar}{}\renewcommand{\rhchEnetOrdPeakBar}{15:00}  % enet: ord: bar end with the most negative mean diff
\providecommand{\rhchEnetOrdPeakDiff}{}\renewcommand{\rhchEnetOrdPeakDiff}{$-$15.8}  % enet: ord: that mean diff (full - HAR) x 10^3
\providecommand{\rhchLgbmShareEvent}{}\renewcommand{\rhchLgbmShareEvent}{63}  % lgbm: % of the every-bar gain on EVENT days
\providecommand{\rhchLgbmShareEventPm}{}\renewcommand{\rhchLgbmShareEventPm}{63}  % lgbm: % of the gain on EVENT days 14:30-16:00
\providecommand{\rhchLgbmEventPeakBar}{}\renewcommand{\rhchLgbmEventPeakBar}{14:30}  % lgbm: event: bar end with the most negative mean diff
\providecommand{\rhchLgbmEventPeakDiff}{}\renewcommand{\rhchLgbmEventPeakDiff}{$-$6012.1}  % lgbm: event: that mean diff (full - HAR) x 10^3
\providecommand{\rhchLgbmOrdPeakBar}{}\renewcommand{\rhchLgbmOrdPeakBar}{12:30}  % lgbm: ord: bar end with the most negative mean diff
\providecommand{\rhchLgbmOrdPeakDiff}{}\renewcommand{\rhchLgbmOrdPeakDiff}{$-$15.9}  % lgbm: ord: that mean diff (full - HAR) x 10^3
\providecommand{\rhchXgbShareEvent}{}\renewcommand{\rhchXgbShareEvent}{64}  % xgb: % of the every-bar gain on EVENT days
\providecommand{\rhchXgbShareEventPm}{}\renewcommand{\rhchXgbShareEventPm}{63}  % xgb: % of the gain on EVENT days 14:30-16:00
\providecommand{\rhchXgbEventPeakBar}{}\renewcommand{\rhchXgbEventPeakBar}{14:30}  % xgb: event: bar end with the most negative mean diff
\providecommand{\rhchXgbEventPeakDiff}{}\renewcommand{\rhchXgbEventPeakDiff}{$-$6274.7}  % xgb: event: that mean diff (full - HAR) x 10^3
\providecommand{\rhchXgbOrdPeakBar}{}\renewcommand{\rhchXgbOrdPeakBar}{12:30}  % xgb: ord: bar end with the most negative mean diff
\providecommand{\rhchXgbOrdPeakDiff}{}\renewcommand{\rhchXgbOrdPeakDiff}{$-$16.3}  % xgb: ord: that mean diff (full - HAR) x 10^3
\providecommand{\rhchRfShareEvent}{}\renewcommand{\rhchRfShareEvent}{60}  % rf: % of the every-bar gain on EVENT days
\providecommand{\rhchRfShareEventPm}{}\renewcommand{\rhchRfShareEventPm}{59}  % rf: % of the gain on EVENT days 14:30-16:00
\providecommand{\rhchRfEventPeakBar}{}\renewcommand{\rhchRfEventPeakBar}{14:30}  % rf: event: bar end with the most negative mean diff
\providecommand{\rhchRfEventPeakDiff}{}\renewcommand{\rhchRfEventPeakDiff}{$-$4773.7}  % rf: event: that mean diff (full - HAR) x 10^3
\providecommand{\rhchRfOrdPeakBar}{}\renewcommand{\rhchRfOrdPeakBar}{12:30}  % rf: ord: bar end with the most negative mean diff
\providecommand{\rhchRfOrdPeakDiff}{}\renewcommand{\rhchRfOrdPeakDiff}{$-$15.2}  % rf: ord: that mean diff (full - HAR) x 10^3
\providecommand{\rhchEventPmBarPct}{}\renewcommand{\rhchEventPmBarPct}{1.0}  % EVENT 14:30-16:00 bars, % of all bars
\providecommand{\rhchLinShareEventMin}{}\renewcommand{\rhchLinShareEventMin}{49}  % Lin: lowest share_event_pct
\providecommand{\rhchLinShareEventMax}{}\renewcommand{\rhchLinShareEventMax}{77}  % Lin: highest share_event_pct
\providecommand{\rhchLinShareEventPmMin}{}\renewcommand{\rhchLinShareEventPmMin}{51}  % Lin: lowest share_event_pm_pct
\providecommand{\rhchLinShareEventPmMax}{}\renewcommand{\rhchLinShareEventPmMax}{78}  % Lin: highest share_event_pm_pct
\providecommand{\rhchTreeShareEventMin}{}\renewcommand{\rhchTreeShareEventMin}{60}  % Tree: lowest share_event_pct
\providecommand{\rhchTreeShareEventMax}{}\renewcommand{\rhchTreeShareEventMax}{64}  % Tree: highest share_event_pct
\providecommand{\rhchTreeShareEventPmMin}{}\renewcommand{\rhchTreeShareEventPmMin}{59}  % Tree: lowest share_event_pm_pct
\providecommand{\rhchTreeShareEventPmMax}{}\renewcommand{\rhchTreeShareEventPmMax}{63}  % Tree: highest share_event_pm_pct
\providecommand{\rhchLassoRidgeBarPct}{}\renewcommand{\rhchLassoRidgeBarPct}{$-$3.1}  % lasso vs ridge, full sets, all: % QLIKE change
\providecommand{\rhchLassoRidgeBarDm}{}\renewcommand{\rhchLassoRidgeBarDm}{$-$4.82}  % lasso vs ridge, all: DM
\providecommand{\rhchLassoRidgeBarP}{}\renewcommand{\rhchLassoRidgeBarP}{\ensuremath{<}0.001}  % lasso vs ridge, all: DM p
\providecommand{\rhchLassoRidgeEventPct}{}\renewcommand{\rhchLassoRidgeEventPct}{$-$31.9}  % lasso vs ridge, full sets, event: % QLIKE change
\providecommand{\rhchLassoRidgeEventDm}{}\renewcommand{\rhchLassoRidgeEventDm}{$-$8.42}  % lasso vs ridge, event: DM
\providecommand{\rhchLassoRidgeEventP}{}\renewcommand{\rhchLassoRidgeEventP}{\ensuremath{<}0.001}  % lasso vs ridge, event: DM p
\providecommand{\rhchLassoRidgeOrdPct}{}\renewcommand{\rhchLassoRidgeOrdPct}{0.8}  % lasso vs ridge, full sets, ord: % QLIKE change
\providecommand{\rhchLassoRidgeOrdDm}{}\renewcommand{\rhchLassoRidgeOrdDm}{2.28}  % lasso vs ridge, ord: DM
\providecommand{\rhchLassoRidgeOrdP}{}\renewcommand{\rhchLassoRidgeOrdP}{0.023}  % lasso vs ridge, ord: DM p
\providecommand{\rhchLassoEnetBarPct}{}\renewcommand{\rhchLassoEnetBarPct}{$-$2.4}  % lasso vs enet, full sets, all: % QLIKE change
\providecommand{\rhchLassoEnetBarDm}{}\renewcommand{\rhchLassoEnetBarDm}{$-$3.49}  % lasso vs enet, all: DM
\providecommand{\rhchLassoEnetBarP}{}\renewcommand{\rhchLassoEnetBarP}{\ensuremath{<}0.001}  % lasso vs enet, all: DM p
\providecommand{\rhchLassoEnetEventPct}{}\renewcommand{\rhchLassoEnetEventPct}{$-$31.8}  % lasso vs enet, full sets, event: % QLIKE change
\providecommand{\rhchLassoEnetEventDm}{}\renewcommand{\rhchLassoEnetEventDm}{$-$7.69}  % lasso vs enet, event: DM
\providecommand{\rhchLassoEnetEventP}{}\renewcommand{\rhchLassoEnetEventP}{\ensuremath{<}0.001}  % lasso vs enet, event: DM p
\providecommand{\rhchLassoEnetOrdPct}{}\renewcommand{\rhchLassoEnetOrdPct}{1.6}  % lasso vs enet, full sets, ord: % QLIKE change
\providecommand{\rhchLassoEnetOrdDm}{}\renewcommand{\rhchLassoEnetOrdDm}{3.35}  % lasso vs enet, ord: DM
\providecommand{\rhchLassoEnetOrdP}{}\renewcommand{\rhchLassoEnetOrdP}{\ensuremath{<}0.001}  % lasso vs enet, ord: DM p
\providecommand{\rhchRidgeBarJointPct}{}\renewcommand{\rhchRidgeBarJointPct}{$-$6.0}  % ridge all: full vs HAR, % QLIKE change
\providecommand{\rhchRidgeBarJointDm}{}\renewcommand{\rhchRidgeBarJointDm}{$-$6.28}  % ridge all: full vs HAR, DM
\providecommand{\rhchRidgeBarJointP}{}\renewcommand{\rhchRidgeBarJointP}{\ensuremath{<}0.001}  % ridge all: full vs HAR, DM p
\providecommand{\rhchRidgeBarAddRmPct}{}\renewcommand{\rhchRidgeBarAddRmPct}{$-$2.6}  % ridge all add_rm vs ols|har: %
\providecommand{\rhchRidgeBarAddRmP}{}\renewcommand{\rhchRidgeBarAddRmP}{\ensuremath{<}0.001}  % ridge all add_rm: DM p
\providecommand{\rhchRidgeBarLoRmPct}{}\renewcommand{\rhchRidgeBarLoRmPct}{1.8}  % ridge all lo_rm vs ridge|full: %
\providecommand{\rhchRidgeBarLoRmP}{}\renewcommand{\rhchRidgeBarLoRmP}{\ensuremath{<}0.001}  % ridge all lo_rm: DM p
\providecommand{\rhchRidgeBarAddLiqPct}{}\renewcommand{\rhchRidgeBarAddLiqPct}{$-$0.4}  % ridge all add_liq vs ols|har: %
\providecommand{\rhchRidgeBarAddLiqP}{}\renewcommand{\rhchRidgeBarAddLiqP}{0.326}  % ridge all add_liq: DM p
\providecommand{\rhchRidgeBarLoLiqPct}{}\renewcommand{\rhchRidgeBarLoLiqPct}{0.2}  % ridge all lo_liq vs ridge|full: %
\providecommand{\rhchRidgeBarLoLiqP}{}\renewcommand{\rhchRidgeBarLoLiqP}{0.467}  % ridge all lo_liq: DM p
\providecommand{\rhchRidgeBarAddEwPct}{}\renewcommand{\rhchRidgeBarAddEwPct}{$-$0.9}  % ridge all add_ew vs ols|har: %
\providecommand{\rhchRidgeBarAddEwP}{}\renewcommand{\rhchRidgeBarAddEwP}{0.024}  % ridge all add_ew: DM p
\providecommand{\rhchRidgeBarLoEwPct}{}\renewcommand{\rhchRidgeBarLoEwPct}{$-$0.1}  % ridge all lo_ew vs ridge|full: %
\providecommand{\rhchRidgeBarLoEwP}{}\renewcommand{\rhchRidgeBarLoEwP}{0.381}  % ridge all lo_ew: DM p
\providecommand{\rhchRidgeBarAddVwPct}{}\renewcommand{\rhchRidgeBarAddVwPct}{$-$0.9}  % ridge all add_vw vs ols|har: %
\providecommand{\rhchRidgeBarAddVwP}{}\renewcommand{\rhchRidgeBarAddVwP}{0.114}  % ridge all add_vw: DM p
\providecommand{\rhchRidgeBarLoVwPct}{}\renewcommand{\rhchRidgeBarLoVwPct}{$-$0.1}  % ridge all lo_vw vs ridge|full: %
\providecommand{\rhchRidgeBarLoVwP}{}\renewcommand{\rhchRidgeBarLoVwP}{0.487}  % ridge all lo_vw: DM p
\providecommand{\rhchRidgeBarAddSentPct}{}\renewcommand{\rhchRidgeBarAddSentPct}{0.2}  % ridge all add_sent vs ols|har: %
\providecommand{\rhchRidgeBarAddSentP}{}\renewcommand{\rhchRidgeBarAddSentP}{0.411}  % ridge all add_sent: DM p
\providecommand{\rhchRidgeBarLoSentPct}{}\renewcommand{\rhchRidgeBarLoSentPct}{0.1}  % ridge all lo_sent vs ridge|full: %
\providecommand{\rhchRidgeBarLoSentP}{}\renewcommand{\rhchRidgeBarLoSentP}{0.637}  % ridge all lo_sent: DM p
\providecommand{\rhchRidgeBarAddIvPct}{}\renewcommand{\rhchRidgeBarAddIvPct}{0.2}  % ridge all add_iv vs ols|har: %
\providecommand{\rhchRidgeBarAddIvP}{}\renewcommand{\rhchRidgeBarAddIvP}{0.506}  % ridge all add_iv: DM p
\providecommand{\rhchRidgeBarLoIvPct}{}\renewcommand{\rhchRidgeBarLoIvPct}{0.1}  % ridge all lo_iv vs ridge|full: %
\providecommand{\rhchRidgeBarLoIvP}{}\renewcommand{\rhchRidgeBarLoIvP}{0.675}  % ridge all lo_iv: DM p
\providecommand{\rhchRidgeBarAddVdPct}{}\renewcommand{\rhchRidgeBarAddVdPct}{0.5}  % ridge all add_vd vs ols|har: %
\providecommand{\rhchRidgeBarAddVdP}{}\renewcommand{\rhchRidgeBarAddVdP}{0.002}  % ridge all add_vd: DM p
\providecommand{\rhchRidgeBarLoVdPct}{}\renewcommand{\rhchRidgeBarLoVdPct}{$-$0.1}  % ridge all lo_vd vs ridge|full: %
\providecommand{\rhchRidgeBarLoVdP}{}\renewcommand{\rhchRidgeBarLoVdP}{0.472}  % ridge all lo_vd: DM p
\providecommand{\rhchRidgeBarAddFomcPct}{}\renewcommand{\rhchRidgeBarAddFomcPct}{$-$3.1}  % ridge all add_fomc vs ols|har: %
\providecommand{\rhchRidgeBarAddFomcP}{}\renewcommand{\rhchRidgeBarAddFomcP}{\ensuremath{<}0.001}  % ridge all add_fomc: DM p
\providecommand{\rhchRidgeBarLoFomcPct}{}\renewcommand{\rhchRidgeBarLoFomcPct}{3.1}  % ridge all lo_fomc vs ridge|full: %
\providecommand{\rhchRidgeBarLoFomcP}{}\renewcommand{\rhchRidgeBarLoFomcP}{\ensuremath{<}0.001}  % ridge all lo_fomc: DM p
\providecommand{\rhchRidgeEventJointPct}{}\renewcommand{\rhchRidgeEventJointPct}{$-$22.4}  % ridge event: full vs HAR, % QLIKE change
\providecommand{\rhchRidgeEventJointDm}{}\renewcommand{\rhchRidgeEventJointDm}{$-$6.45}  % ridge event: full vs HAR, DM
\providecommand{\rhchRidgeEventJointP}{}\renewcommand{\rhchRidgeEventJointP}{\ensuremath{<}0.001}  % ridge event: full vs HAR, DM p
\providecommand{\rhchRidgeEventAddRmPct}{}\renewcommand{\rhchRidgeEventAddRmPct}{4.4}  % ridge event add_rm vs ols|har: %
\providecommand{\rhchRidgeEventAddRmP}{}\renewcommand{\rhchRidgeEventAddRmP}{0.023}  % ridge event add_rm: DM p
\providecommand{\rhchRidgeEventLoRmPct}{}\renewcommand{\rhchRidgeEventLoRmPct}{$-$1.1}  % ridge event lo_rm vs ridge|full: %
\providecommand{\rhchRidgeEventLoRmP}{}\renewcommand{\rhchRidgeEventLoRmP}{0.133}  % ridge event lo_rm: DM p
\providecommand{\rhchRidgeEventAddLiqPct}{}\renewcommand{\rhchRidgeEventAddLiqPct}{5.4}  % ridge event add_liq vs ols|har: %
\providecommand{\rhchRidgeEventAddLiqP}{}\renewcommand{\rhchRidgeEventAddLiqP}{0.010}  % ridge event add_liq: DM p
\providecommand{\rhchRidgeEventLoLiqPct}{}\renewcommand{\rhchRidgeEventLoLiqPct}{$-$0.6}  % ridge event lo_liq vs ridge|full: %
\providecommand{\rhchRidgeEventLoLiqP}{}\renewcommand{\rhchRidgeEventLoLiqP}{0.230}  % ridge event lo_liq: DM p
\providecommand{\rhchRidgeEventAddEwPct}{}\renewcommand{\rhchRidgeEventAddEwPct}{$-$1.8}  % ridge event add_ew vs ols|har: %
\providecommand{\rhchRidgeEventAddEwP}{}\renewcommand{\rhchRidgeEventAddEwP}{0.176}  % ridge event add_ew: DM p
\providecommand{\rhchRidgeEventLoEwPct}{}\renewcommand{\rhchRidgeEventLoEwPct}{0.0}  % ridge event lo_ew vs ridge|full: %
\providecommand{\rhchRidgeEventLoEwP}{}\renewcommand{\rhchRidgeEventLoEwP}{0.968}  % ridge event lo_ew: DM p
\providecommand{\rhchRidgeEventAddVwPct}{}\renewcommand{\rhchRidgeEventAddVwPct}{$-$2.5}  % ridge event add_vw vs ols|har: %
\providecommand{\rhchRidgeEventAddVwP}{}\renewcommand{\rhchRidgeEventAddVwP}{0.143}  % ridge event add_vw: DM p
\providecommand{\rhchRidgeEventLoVwPct}{}\renewcommand{\rhchRidgeEventLoVwPct}{0.1}  % ridge event lo_vw vs ridge|full: %
\providecommand{\rhchRidgeEventLoVwP}{}\renewcommand{\rhchRidgeEventLoVwP}{0.704}  % ridge event lo_vw: DM p
\providecommand{\rhchRidgeEventAddSentPct}{}\renewcommand{\rhchRidgeEventAddSentPct}{$-$0.1}  % ridge event add_sent vs ols|har: %
\providecommand{\rhchRidgeEventAddSentP}{}\renewcommand{\rhchRidgeEventAddSentP}{0.970}  % ridge event add_sent: DM p
\providecommand{\rhchRidgeEventLoSentPct}{}\renewcommand{\rhchRidgeEventLoSentPct}{1.7}  % ridge event lo_sent vs ridge|full: %
\providecommand{\rhchRidgeEventLoSentP}{}\renewcommand{\rhchRidgeEventLoSentP}{0.322}  % ridge event lo_sent: DM p
\providecommand{\rhchRidgeEventAddIvPct}{}\renewcommand{\rhchRidgeEventAddIvPct}{2.8}  % ridge event add_iv vs ols|har: %
\providecommand{\rhchRidgeEventAddIvP}{}\renewcommand{\rhchRidgeEventAddIvP}{0.011}  % ridge event add_iv: DM p
\providecommand{\rhchRidgeEventLoIvPct}{}\renewcommand{\rhchRidgeEventLoIvPct}{$-$0.2}  % ridge event lo_iv vs ridge|full: %
\providecommand{\rhchRidgeEventLoIvP}{}\renewcommand{\rhchRidgeEventLoIvP}{0.823}  % ridge event lo_iv: DM p
\providecommand{\rhchRidgeEventAddVdPct}{}\renewcommand{\rhchRidgeEventAddVdPct}{1.9}  % ridge event add_vd vs ols|har: %
\providecommand{\rhchRidgeEventAddVdP}{}\renewcommand{\rhchRidgeEventAddVdP}{0.028}  % ridge event add_vd: DM p
\providecommand{\rhchRidgeEventLoVdPct}{}\renewcommand{\rhchRidgeEventLoVdPct}{0.0}  % ridge event lo_vd vs ridge|full: %
\providecommand{\rhchRidgeEventLoVdP}{}\renewcommand{\rhchRidgeEventLoVdP}{0.971}  % ridge event lo_vd: DM p
\providecommand{\rhchRidgeEventAddFomcPct}{}\renewcommand{\rhchRidgeEventAddFomcPct}{$-$23.4}  % ridge event add_fomc vs ols|har: %
\providecommand{\rhchRidgeEventAddFomcP}{}\renewcommand{\rhchRidgeEventAddFomcP}{\ensuremath{<}0.001}  % ridge event add_fomc: DM p
\providecommand{\rhchRidgeEventLoFomcPct}{}\renewcommand{\rhchRidgeEventLoFomcPct}{28.6}  % ridge event lo_fomc vs ridge|full: %
\providecommand{\rhchRidgeEventLoFomcP}{}\renewcommand{\rhchRidgeEventLoFomcP}{\ensuremath{<}0.001}  % ridge event lo_fomc: DM p
\providecommand{\rhchRidgeEventReading}{}\renewcommand{\rhchRidgeEventReading}{concentrated}  % ridge event: item 3 reading
\providecommand{\rhchRidgeEventTopAddShare}{}\renewcommand{\rhchRidgeEventTopAddShare}{104}  % ridge event: largest add-alone gain, % of the joint gain (fomc)
\providecommand{\rhchRidgeEventSumAddRatio}{}\renewcommand{\rhchRidgeEventSumAddRatio}{0.60}  % ridge event: add-alone gains summed over groups / joint gain
\providecommand{\rhchRidgeEventNLoSig}{}\renewcommand{\rhchRidgeEventNLoSig}{1}  % ridge event: significant leave-one-outs
\providecommand{\rhchRidgeEventNLoRun}{}\renewcommand{\rhchRidgeEventNLoRun}{8}  % ridge event: leave-one-outs available (of 8)
\providecommand{\rhchRidgeEventConcTest}{}\renewcommand{\rhchRidgeEventConcTest}{yes}  % ridge event: the CONC definition on its own (pending: undetermined)
\providecommand{\rhchRidgeEventRedTest}{}\renewcommand{\rhchRidgeEventRedTest}{no}  % ridge event: the RED definition on its own (pending: undetermined)
\providecommand{\rhchRidgeEventReadingRedFirst}{}\renewcommand{\rhchRidgeEventReadingRedFirst}{concentrated}  % ridge event: reading with REDUNDANT tested before CONCENTRATED (the RULE tests CONCENTRATED first)
\providecommand{\rhchRidgeOrdJointPct}{}\renewcommand{\rhchRidgeOrdJointPct}{$-$3.2}  % ridge ord: full vs HAR, % QLIKE change
\providecommand{\rhchRidgeOrdJointDm}{}\renewcommand{\rhchRidgeOrdJointDm}{$-$3.51}  % ridge ord: full vs HAR, DM
\providecommand{\rhchRidgeOrdJointP}{}\renewcommand{\rhchRidgeOrdJointP}{\ensuremath{<}0.001}  % ridge ord: full vs HAR, DM p
\providecommand{\rhchRidgeOrdAddRmPct}{}\renewcommand{\rhchRidgeOrdAddRmPct}{$-$3.8}  % ridge ord add_rm vs ols|har: %
\providecommand{\rhchRidgeOrdAddRmP}{}\renewcommand{\rhchRidgeOrdAddRmP}{\ensuremath{<}0.001}  % ridge ord add_rm: DM p
\providecommand{\rhchRidgeOrdLoRmPct}{}\renewcommand{\rhchRidgeOrdLoRmPct}{2.2}  % ridge ord lo_rm vs ridge|full: %
\providecommand{\rhchRidgeOrdLoRmP}{}\renewcommand{\rhchRidgeOrdLoRmP}{\ensuremath{<}0.001}  % ridge ord lo_rm: DM p
\providecommand{\rhchRidgeOrdAddLiqPct}{}\renewcommand{\rhchRidgeOrdAddLiqPct}{$-$1.4}  % ridge ord add_liq vs ols|har: %
\providecommand{\rhchRidgeOrdAddLiqP}{}\renewcommand{\rhchRidgeOrdAddLiqP}{\ensuremath{<}0.001}  % ridge ord add_liq: DM p
\providecommand{\rhchRidgeOrdLoLiqPct}{}\renewcommand{\rhchRidgeOrdLoLiqPct}{0.3}  % ridge ord lo_liq vs ridge|full: %
\providecommand{\rhchRidgeOrdLoLiqP}{}\renewcommand{\rhchRidgeOrdLoLiqP}{0.261}  % ridge ord lo_liq: DM p
\providecommand{\rhchRidgeOrdAddEwPct}{}\renewcommand{\rhchRidgeOrdAddEwPct}{$-$0.7}  % ridge ord add_ew vs ols|har: %
\providecommand{\rhchRidgeOrdAddEwP}{}\renewcommand{\rhchRidgeOrdAddEwP}{0.044}  % ridge ord add_ew: DM p
\providecommand{\rhchRidgeOrdLoEwPct}{}\renewcommand{\rhchRidgeOrdLoEwPct}{$-$0.1}  % ridge ord lo_ew vs ridge|full: %
\providecommand{\rhchRidgeOrdLoEwP}{}\renewcommand{\rhchRidgeOrdLoEwP}{0.316}  % ridge ord lo_ew: DM p
\providecommand{\rhchRidgeOrdAddVwPct}{}\renewcommand{\rhchRidgeOrdAddVwPct}{$-$0.6}  % ridge ord add_vw vs ols|har: %
\providecommand{\rhchRidgeOrdAddVwP}{}\renewcommand{\rhchRidgeOrdAddVwP}{0.309}  % ridge ord add_vw: DM p
\providecommand{\rhchRidgeOrdLoVwPct}{}\renewcommand{\rhchRidgeOrdLoVwPct}{$-$0.1}  % ridge ord lo_vw vs ridge|full: %
\providecommand{\rhchRidgeOrdLoVwP}{}\renewcommand{\rhchRidgeOrdLoVwP}{0.405}  % ridge ord lo_vw: DM p
\providecommand{\rhchRidgeOrdAddSentPct}{}\renewcommand{\rhchRidgeOrdAddSentPct}{0.3}  % ridge ord add_sent vs ols|har: %
\providecommand{\rhchRidgeOrdAddSentP}{}\renewcommand{\rhchRidgeOrdAddSentP}{0.246}  % ridge ord add_sent: DM p
\providecommand{\rhchRidgeOrdLoSentPct}{}\renewcommand{\rhchRidgeOrdLoSentPct}{$-$0.1}  % ridge ord lo_sent vs ridge|full: %
\providecommand{\rhchRidgeOrdLoSentP}{}\renewcommand{\rhchRidgeOrdLoSentP}{0.271}  % ridge ord lo_sent: DM p
\providecommand{\rhchRidgeOrdAddIvPct}{}\renewcommand{\rhchRidgeOrdAddIvPct}{$-$0.2}  % ridge ord add_iv vs ols|har: %
\providecommand{\rhchRidgeOrdAddIvP}{}\renewcommand{\rhchRidgeOrdAddIvP}{0.580}  % ridge ord add_iv: DM p
\providecommand{\rhchRidgeOrdLoIvPct}{}\renewcommand{\rhchRidgeOrdLoIvPct}{0.2}  % ridge ord lo_iv vs ridge|full: %
\providecommand{\rhchRidgeOrdLoIvP}{}\renewcommand{\rhchRidgeOrdLoIvP}{0.611}  % ridge ord lo_iv: DM p
\providecommand{\rhchRidgeOrdAddVdPct}{}\renewcommand{\rhchRidgeOrdAddVdPct}{0.3}  % ridge ord add_vd vs ols|har: %
\providecommand{\rhchRidgeOrdAddVdP}{}\renewcommand{\rhchRidgeOrdAddVdP}{0.069}  % ridge ord add_vd: DM p
\providecommand{\rhchRidgeOrdLoVdPct}{}\renewcommand{\rhchRidgeOrdLoVdPct}{$-$0.1}  % ridge ord lo_vd vs ridge|full: %
\providecommand{\rhchRidgeOrdLoVdP}{}\renewcommand{\rhchRidgeOrdLoVdP}{0.437}  % ridge ord lo_vd: DM p
\providecommand{\rhchRidgeOrdAddFomcPct}{}\renewcommand{\rhchRidgeOrdAddFomcPct}{0.4}  % ridge ord add_fomc vs ols|har: %
\providecommand{\rhchRidgeOrdAddFomcP}{}\renewcommand{\rhchRidgeOrdAddFomcP}{\ensuremath{<}0.001}  % ridge ord add_fomc: DM p
\providecommand{\rhchRidgeOrdLoFomcPct}{}\renewcommand{\rhchRidgeOrdLoFomcPct}{$-$0.4}  % ridge ord lo_fomc vs ridge|full: %
\providecommand{\rhchRidgeOrdLoFomcP}{}\renewcommand{\rhchRidgeOrdLoFomcP}{\ensuremath{<}0.001}  % ridge ord lo_fomc: DM p
\providecommand{\rhchRidgeOrdReading}{}\renewcommand{\rhchRidgeOrdReading}{concentrated}  % ridge ord: item 3 reading
\providecommand{\rhchRidgeOrdTopAddShare}{}\renewcommand{\rhchRidgeOrdTopAddShare}{118}  % ridge ord: largest add-alone gain, % of the joint gain (rm)
\providecommand{\rhchRidgeOrdSumAddRatio}{}\renewcommand{\rhchRidgeOrdSumAddRatio}{1.79}  % ridge ord: add-alone gains summed over groups / joint gain
\providecommand{\rhchRidgeOrdNLoSig}{}\renewcommand{\rhchRidgeOrdNLoSig}{1}  % ridge ord: significant leave-one-outs
\providecommand{\rhchRidgeOrdNLoRun}{}\renewcommand{\rhchRidgeOrdNLoRun}{8}  % ridge ord: leave-one-outs available (of 8)
\providecommand{\rhchRidgeOrdConcTest}{}\renewcommand{\rhchRidgeOrdConcTest}{yes}  % ridge ord: the CONC definition on its own (pending: undetermined)
\providecommand{\rhchRidgeOrdRedTest}{}\renewcommand{\rhchRidgeOrdRedTest}{yes}  % ridge ord: the RED definition on its own (pending: undetermined)
\providecommand{\rhchRidgeOrdReadingRedFirst}{}\renewcommand{\rhchRidgeOrdReadingRedFirst}{redundant}  % ridge ord: reading with REDUNDANT tested before CONCENTRATED (the RULE tests CONCENTRATED first)
\providecommand{\rhchLassoBarJointPct}{}\renewcommand{\rhchLassoBarJointPct}{$-$9.0}  % lasso all: full vs HAR, % QLIKE change
\providecommand{\rhchLassoBarJointDm}{}\renewcommand{\rhchLassoBarJointDm}{$-$6.34}  % lasso all: full vs HAR, DM
\providecommand{\rhchLassoBarJointP}{}\renewcommand{\rhchLassoBarJointP}{\ensuremath{<}0.001}  % lasso all: full vs HAR, DM p
\providecommand{\rhchLassoBarAddRmPct}{}\renewcommand{\rhchLassoBarAddRmPct}{$-$2.1}  % lasso all add_rm vs ols|har: %
\providecommand{\rhchLassoBarAddRmP}{}\renewcommand{\rhchLassoBarAddRmP}{0.008}  % lasso all add_rm: DM p
\providecommand{\rhchLassoBarLoRmPct}{}\renewcommand{\rhchLassoBarLoRmPct}{1.6}  % lasso all lo_rm vs lasso|full: %
\providecommand{\rhchLassoBarLoRmP}{}\renewcommand{\rhchLassoBarLoRmP}{0.010}  % lasso all lo_rm: DM p
\providecommand{\rhchLassoBarAddLiqPct}{}\renewcommand{\rhchLassoBarAddLiqPct}{$-$0.3}  % lasso all add_liq vs ols|har: %
\providecommand{\rhchLassoBarAddLiqP}{}\renewcommand{\rhchLassoBarAddLiqP}{0.583}  % lasso all add_liq: DM p
\providecommand{\rhchLassoBarLoLiqPct}{}\renewcommand{\rhchLassoBarLoLiqPct}{0.1}  % lasso all lo_liq vs lasso|full: %
\providecommand{\rhchLassoBarLoLiqP}{}\renewcommand{\rhchLassoBarLoLiqP}{0.783}  % lasso all lo_liq: DM p
\providecommand{\rhchLassoBarAddEwPct}{}\renewcommand{\rhchLassoBarAddEwPct}{$-$0.8}  % lasso all add_ew vs ols|har: %
\providecommand{\rhchLassoBarAddEwP}{}\renewcommand{\rhchLassoBarAddEwP}{0.045}  % lasso all add_ew: DM p
\providecommand{\rhchLassoBarLoEwPct}{}\renewcommand{\rhchLassoBarLoEwPct}{$-$0.1}  % lasso all lo_ew vs lasso|full: %
\providecommand{\rhchLassoBarLoEwP}{}\renewcommand{\rhchLassoBarLoEwP}{0.583}  % lasso all lo_ew: DM p
\providecommand{\rhchLassoBarAddVwPct}{}\renewcommand{\rhchLassoBarAddVwPct}{$-$0.6}  % lasso all add_vw vs ols|har: %
\providecommand{\rhchLassoBarAddVwP}{}\renewcommand{\rhchLassoBarAddVwP}{0.309}  % lasso all add_vw: DM p
\providecommand{\rhchLassoBarLoVwPct}{}\renewcommand{\rhchLassoBarLoVwPct}{$-$0.2}  % lasso all lo_vw vs lasso|full: %
\providecommand{\rhchLassoBarLoVwP}{}\renewcommand{\rhchLassoBarLoVwP}{0.428}  % lasso all lo_vw: DM p
\providecommand{\rhchLassoBarAddSentPct}{}\renewcommand{\rhchLassoBarAddSentPct}{0.1}  % lasso all add_sent vs ols|har: %
\providecommand{\rhchLassoBarAddSentP}{}\renewcommand{\rhchLassoBarAddSentP}{0.550}  % lasso all add_sent: DM p
\providecommand{\rhchLassoBarLoSentPct}{}\renewcommand{\rhchLassoBarLoSentPct}{0.1}  % lasso all lo_sent vs lasso|full: %
\providecommand{\rhchLassoBarLoSentP}{}\renewcommand{\rhchLassoBarLoSentP}{0.693}  % lasso all lo_sent: DM p
\providecommand{\rhchLassoBarAddIvPct}{}\renewcommand{\rhchLassoBarAddIvPct}{$-$0.3}  % lasso all add_iv vs ols|har: %
\providecommand{\rhchLassoBarAddIvP}{}\renewcommand{\rhchLassoBarAddIvP}{0.430}  % lasso all add_iv: DM p
\providecommand{\rhchLassoBarLoIvPct}{}\renewcommand{\rhchLassoBarLoIvPct}{0.5}  % lasso all lo_iv vs lasso|full: %
\providecommand{\rhchLassoBarLoIvP}{}\renewcommand{\rhchLassoBarLoIvP}{0.313}  % lasso all lo_iv: DM p
\providecommand{\rhchLassoBarAddVdPct}{}\renewcommand{\rhchLassoBarAddVdPct}{0.6}  % lasso all add_vd vs ols|har: %
\providecommand{\rhchLassoBarAddVdP}{}\renewcommand{\rhchLassoBarAddVdP}{0.002}  % lasso all add_vd: DM p
\providecommand{\rhchLassoBarLoVdPct}{}\renewcommand{\rhchLassoBarLoVdPct}{0.0}  % lasso all lo_vd vs lasso|full: %
\providecommand{\rhchLassoBarLoVdP}{}\renewcommand{\rhchLassoBarLoVdP}{0.807}  % lasso all lo_vd: DM p
\providecommand{\rhchLassoBarAddFomcPct}{}\renewcommand{\rhchLassoBarAddFomcPct}{$-$5.8}  % lasso all add_fomc vs ols|har: %
\providecommand{\rhchLassoBarAddFomcP}{}\renewcommand{\rhchLassoBarAddFomcP}{\ensuremath{<}0.001}  % lasso all add_fomc: DM p
\providecommand{\rhchLassoBarLoFomcPct}{}\renewcommand{\rhchLassoBarLoFomcPct}{6.8}  % lasso all lo_fomc vs lasso|full: %
\providecommand{\rhchLassoBarLoFomcP}{}\renewcommand{\rhchLassoBarLoFomcP}{\ensuremath{<}0.001}  % lasso all lo_fomc: DM p
\providecommand{\rhchLassoEventJointPct}{}\renewcommand{\rhchLassoEventJointPct}{$-$47.2}  % lasso event: full vs HAR, % QLIKE change
\providecommand{\rhchLassoEventJointDm}{}\renewcommand{\rhchLassoEventJointDm}{$-$7.56}  % lasso event: full vs HAR, DM
\providecommand{\rhchLassoEventJointP}{}\renewcommand{\rhchLassoEventJointP}{\ensuremath{<}0.001}  % lasso event: full vs HAR, DM p
\providecommand{\rhchLassoEventAddRmPct}{}\renewcommand{\rhchLassoEventAddRmPct}{4.7}  % lasso event add_rm vs ols|har: %
\providecommand{\rhchLassoEventAddRmP}{}\renewcommand{\rhchLassoEventAddRmP}{0.024}  % lasso event add_rm: DM p
\providecommand{\rhchLassoEventLoRmPct}{}\renewcommand{\rhchLassoEventLoRmPct}{$-$0.8}  % lasso event lo_rm vs lasso|full: %
\providecommand{\rhchLassoEventLoRmP}{}\renewcommand{\rhchLassoEventLoRmP}{0.265}  % lasso event lo_rm: DM p
\providecommand{\rhchLassoEventAddLiqPct}{}\renewcommand{\rhchLassoEventAddLiqPct}{6.1}  % lasso event add_liq vs ols|har: %
\providecommand{\rhchLassoEventAddLiqP}{}\renewcommand{\rhchLassoEventAddLiqP}{0.007}  % lasso event add_liq: DM p
\providecommand{\rhchLassoEventLoLiqPct}{}\renewcommand{\rhchLassoEventLoLiqPct}{0.1}  % lasso event lo_liq vs lasso|full: %
\providecommand{\rhchLassoEventLoLiqP}{}\renewcommand{\rhchLassoEventLoLiqP}{0.838}  % lasso event lo_liq: DM p
\providecommand{\rhchLassoEventAddEwPct}{}\renewcommand{\rhchLassoEventAddEwPct}{$-$1.7}  % lasso event add_ew vs ols|har: %
\providecommand{\rhchLassoEventAddEwP}{}\renewcommand{\rhchLassoEventAddEwP}{0.235}  % lasso event add_ew: DM p
\providecommand{\rhchLassoEventLoEwPct}{}\renewcommand{\rhchLassoEventLoEwPct}{0.1}  % lasso event lo_ew vs lasso|full: %
\providecommand{\rhchLassoEventLoEwP}{}\renewcommand{\rhchLassoEventLoEwP}{0.717}  % lasso event lo_ew: DM p
\providecommand{\rhchLassoEventAddVwPct}{}\renewcommand{\rhchLassoEventAddVwPct}{$-$2.4}  % lasso event add_vw vs ols|har: %
\providecommand{\rhchLassoEventAddVwP}{}\renewcommand{\rhchLassoEventAddVwP}{0.171}  % lasso event add_vw: DM p
\providecommand{\rhchLassoEventLoVwPct}{}\renewcommand{\rhchLassoEventLoVwPct}{$-$0.3}  % lasso event lo_vw vs lasso|full: %
\providecommand{\rhchLassoEventLoVwP}{}\renewcommand{\rhchLassoEventLoVwP}{0.438}  % lasso event lo_vw: DM p
\providecommand{\rhchLassoEventAddSentPct}{}\renewcommand{\rhchLassoEventAddSentPct}{$-$1.3}  % lasso event add_sent vs ols|har: %
\providecommand{\rhchLassoEventAddSentP}{}\renewcommand{\rhchLassoEventAddSentP}{0.464}  % lasso event add_sent: DM p
\providecommand{\rhchLassoEventLoSentPct}{}\renewcommand{\rhchLassoEventLoSentPct}{2.0}  % lasso event lo_sent vs lasso|full: %
\providecommand{\rhchLassoEventLoSentP}{}\renewcommand{\rhchLassoEventLoSentP}{0.209}  % lasso event lo_sent: DM p
\providecommand{\rhchLassoEventAddIvPct}{}\renewcommand{\rhchLassoEventAddIvPct}{2.9}  % lasso event add_iv vs ols|har: %
\providecommand{\rhchLassoEventAddIvP}{}\renewcommand{\rhchLassoEventAddIvP}{0.004}  % lasso event add_iv: DM p
\providecommand{\rhchLassoEventLoIvPct}{}\renewcommand{\rhchLassoEventLoIvPct}{$-$0.1}  % lasso event lo_iv vs lasso|full: %
\providecommand{\rhchLassoEventLoIvP}{}\renewcommand{\rhchLassoEventLoIvP}{0.868}  % lasso event lo_iv: DM p
\providecommand{\rhchLassoEventAddVdPct}{}\renewcommand{\rhchLassoEventAddVdPct}{1.9}  % lasso event add_vd vs ols|har: %
\providecommand{\rhchLassoEventAddVdP}{}\renewcommand{\rhchLassoEventAddVdP}{0.032}  % lasso event add_vd: DM p
\providecommand{\rhchLassoEventLoVdPct}{}\renewcommand{\rhchLassoEventLoVdPct}{0.3}  % lasso event lo_vd vs lasso|full: %
\providecommand{\rhchLassoEventLoVdP}{}\renewcommand{\rhchLassoEventLoVdP}{0.012}  % lasso event lo_vd: DM p
\providecommand{\rhchLassoEventAddFomcPct}{}\renewcommand{\rhchLassoEventAddFomcPct}{$-$43.9}  % lasso event add_fomc vs ols|har: %
\providecommand{\rhchLassoEventAddFomcP}{}\renewcommand{\rhchLassoEventAddFomcP}{\ensuremath{<}0.001}  % lasso event add_fomc: DM p
\providecommand{\rhchLassoEventLoFomcPct}{}\renewcommand{\rhchLassoEventLoFomcPct}{87.4}  % lasso event lo_fomc vs lasso|full: %
\providecommand{\rhchLassoEventLoFomcP}{}\renewcommand{\rhchLassoEventLoFomcP}{\ensuremath{<}0.001}  % lasso event lo_fomc: DM p
\providecommand{\rhchLassoEventReading}{}\renewcommand{\rhchLassoEventReading}{concentrated}  % lasso event: item 3 reading
\providecommand{\rhchLassoEventTopAddShare}{}\renewcommand{\rhchLassoEventTopAddShare}{93}  % lasso event: largest add-alone gain, % of the joint gain (fomc)
\providecommand{\rhchLassoEventSumAddRatio}{}\renewcommand{\rhchLassoEventSumAddRatio}{0.72}  % lasso event: add-alone gains summed over groups / joint gain
\providecommand{\rhchLassoEventNLoSig}{}\renewcommand{\rhchLassoEventNLoSig}{2}  % lasso event: significant leave-one-outs
\providecommand{\rhchLassoEventNLoRun}{}\renewcommand{\rhchLassoEventNLoRun}{8}  % lasso event: leave-one-outs available (of 8)
\providecommand{\rhchLassoEventConcTest}{}\renewcommand{\rhchLassoEventConcTest}{yes}  % lasso event: the CONC definition on its own (pending: undetermined)
\providecommand{\rhchLassoEventRedTest}{}\renewcommand{\rhchLassoEventRedTest}{no}  % lasso event: the RED definition on its own (pending: undetermined)
\providecommand{\rhchLassoEventReadingRedFirst}{}\renewcommand{\rhchLassoEventReadingRedFirst}{concentrated}  % lasso event: reading with REDUNDANT tested before CONCENTRATED (the RULE tests CONCENTRATED first)
\providecommand{\rhchLassoOrdJointPct}{}\renewcommand{\rhchLassoOrdJointPct}{$-$2.4}  % lasso ord: full vs HAR, % QLIKE change
\providecommand{\rhchLassoOrdJointDm}{}\renewcommand{\rhchLassoOrdJointDm}{$-$2.08}  % lasso ord: full vs HAR, DM
\providecommand{\rhchLassoOrdJointP}{}\renewcommand{\rhchLassoOrdJointP}{0.037}  % lasso ord: full vs HAR, DM p
\providecommand{\rhchLassoOrdAddRmPct}{}\renewcommand{\rhchLassoOrdAddRmPct}{$-$3.3}  % lasso ord add_rm vs ols|har: %
\providecommand{\rhchLassoOrdAddRmP}{}\renewcommand{\rhchLassoOrdAddRmP}{\ensuremath{<}0.001}  % lasso ord add_rm: DM p
\providecommand{\rhchLassoOrdLoRmPct}{}\renewcommand{\rhchLassoOrdLoRmPct}{1.8}  % lasso ord lo_rm vs lasso|full: %
\providecommand{\rhchLassoOrdLoRmP}{}\renewcommand{\rhchLassoOrdLoRmP}{0.007}  % lasso ord lo_rm: DM p
\providecommand{\rhchLassoOrdAddLiqPct}{}\renewcommand{\rhchLassoOrdAddLiqPct}{$-$1.3}  % lasso ord add_liq vs ols|har: %
\providecommand{\rhchLassoOrdAddLiqP}{}\renewcommand{\rhchLassoOrdAddLiqP}{0.003}  % lasso ord add_liq: DM p
\providecommand{\rhchLassoOrdLoLiqPct}{}\renewcommand{\rhchLassoOrdLoLiqPct}{0.1}  % lasso ord lo_liq vs lasso|full: %
\providecommand{\rhchLassoOrdLoLiqP}{}\renewcommand{\rhchLassoOrdLoLiqP}{0.799}  % lasso ord lo_liq: DM p
\providecommand{\rhchLassoOrdAddEwPct}{}\renewcommand{\rhchLassoOrdAddEwPct}{$-$0.7}  % lasso ord add_ew vs ols|har: %
\providecommand{\rhchLassoOrdAddEwP}{}\renewcommand{\rhchLassoOrdAddEwP}{0.077}  % lasso ord add_ew: DM p
\providecommand{\rhchLassoOrdLoEwPct}{}\renewcommand{\rhchLassoOrdLoEwPct}{$-$0.1}  % lasso ord lo_ew vs lasso|full: %
\providecommand{\rhchLassoOrdLoEwP}{}\renewcommand{\rhchLassoOrdLoEwP}{0.530}  % lasso ord lo_ew: DM p
\providecommand{\rhchLassoOrdAddVwPct}{}\renewcommand{\rhchLassoOrdAddVwPct}{$-$0.3}  % lasso ord add_vw vs ols|har: %
\providecommand{\rhchLassoOrdAddVwP}{}\renewcommand{\rhchLassoOrdAddVwP}{0.630}  % lasso ord add_vw: DM p
\providecommand{\rhchLassoOrdLoVwPct}{}\renewcommand{\rhchLassoOrdLoVwPct}{$-$0.2}  % lasso ord lo_vw vs lasso|full: %
\providecommand{\rhchLassoOrdLoVwP}{}\renewcommand{\rhchLassoOrdLoVwP}{0.501}  % lasso ord lo_vw: DM p
\providecommand{\rhchLassoOrdAddSentPct}{}\renewcommand{\rhchLassoOrdAddSentPct}{0.4}  % lasso ord add_sent vs ols|har: %
\providecommand{\rhchLassoOrdAddSentP}{}\renewcommand{\rhchLassoOrdAddSentP}{0.098}  % lasso ord add_sent: DM p
\providecommand{\rhchLassoOrdLoSentPct}{}\renewcommand{\rhchLassoOrdLoSentPct}{$-$0.1}  % lasso ord lo_sent vs lasso|full: %
\providecommand{\rhchLassoOrdLoSentP}{}\renewcommand{\rhchLassoOrdLoSentP}{0.502}  % lasso ord lo_sent: DM p
\providecommand{\rhchLassoOrdAddIvPct}{}\renewcommand{\rhchLassoOrdAddIvPct}{$-$0.9}  % lasso ord add_iv vs ols|har: %
\providecommand{\rhchLassoOrdAddIvP}{}\renewcommand{\rhchLassoOrdAddIvP}{0.042}  % lasso ord add_iv: DM p
\providecommand{\rhchLassoOrdLoIvPct}{}\renewcommand{\rhchLassoOrdLoIvPct}{0.5}  % lasso ord lo_iv vs lasso|full: %
\providecommand{\rhchLassoOrdLoIvP}{}\renewcommand{\rhchLassoOrdLoIvP}{0.292}  % lasso ord lo_iv: DM p
\providecommand{\rhchLassoOrdAddVdPct}{}\renewcommand{\rhchLassoOrdAddVdPct}{0.3}  % lasso ord add_vd vs ols|har: %
\providecommand{\rhchLassoOrdAddVdP}{}\renewcommand{\rhchLassoOrdAddVdP}{0.056}  % lasso ord add_vd: DM p
\providecommand{\rhchLassoOrdLoVdPct}{}\renewcommand{\rhchLassoOrdLoVdPct}{$-$0.1}  % lasso ord lo_vd vs lasso|full: %
\providecommand{\rhchLassoOrdLoVdP}{}\renewcommand{\rhchLassoOrdLoVdP}{0.644}  % lasso ord lo_vd: DM p
\providecommand{\rhchLassoOrdAddFomcPct}{}\renewcommand{\rhchLassoOrdAddFomcPct}{0.7}  % lasso ord add_fomc vs ols|har: %
\providecommand{\rhchLassoOrdAddFomcP}{}\renewcommand{\rhchLassoOrdAddFomcP}{\ensuremath{<}0.001}  % lasso ord add_fomc: DM p
\providecommand{\rhchLassoOrdLoFomcPct}{}\renewcommand{\rhchLassoOrdLoFomcPct}{$-$0.7}  % lasso ord lo_fomc vs lasso|full: %
\providecommand{\rhchLassoOrdLoFomcP}{}\renewcommand{\rhchLassoOrdLoFomcP}{\ensuremath{<}0.001}  % lasso ord lo_fomc: DM p
\providecommand{\rhchLassoOrdReading}{}\renewcommand{\rhchLassoOrdReading}{concentrated}  % lasso ord: item 3 reading
\providecommand{\rhchLassoOrdTopAddShare}{}\renewcommand{\rhchLassoOrdTopAddShare}{136}  % lasso ord: largest add-alone gain, % of the joint gain (rm)
\providecommand{\rhchLassoOrdSumAddRatio}{}\renewcommand{\rhchLassoOrdSumAddRatio}{2.07}  % lasso ord: add-alone gains summed over groups / joint gain
\providecommand{\rhchLassoOrdNLoSig}{}\renewcommand{\rhchLassoOrdNLoSig}{1}  % lasso ord: significant leave-one-outs
\providecommand{\rhchLassoOrdNLoRun}{}\renewcommand{\rhchLassoOrdNLoRun}{8}  % lasso ord: leave-one-outs available (of 8)
\providecommand{\rhchLassoOrdConcTest}{}\renewcommand{\rhchLassoOrdConcTest}{yes}  % lasso ord: the CONC definition on its own (pending: undetermined)
\providecommand{\rhchLassoOrdRedTest}{}\renewcommand{\rhchLassoOrdRedTest}{yes}  % lasso ord: the RED definition on its own (pending: undetermined)
\providecommand{\rhchLassoOrdReadingRedFirst}{}\renewcommand{\rhchLassoOrdReadingRedFirst}{redundant}  % lasso ord: reading with REDUNDANT tested before CONCENTRATED (the RULE tests CONCENTRATED first)
\providecommand{\rhchEnetBarJointPct}{}\renewcommand{\rhchEnetBarJointPct}{$-$6.7}  % enet all: full vs HAR, % QLIKE change
\providecommand{\rhchEnetBarJointDm}{}\renewcommand{\rhchEnetBarJointDm}{$-$7.36}  % enet all: full vs HAR, DM
\providecommand{\rhchEnetBarJointP}{}\renewcommand{\rhchEnetBarJointP}{\ensuremath{<}0.001}  % enet all: full vs HAR, DM p
\providecommand{\rhchEnetBarAddRmPct}{}\renewcommand{\rhchEnetBarAddRmPct}{$-$2.7}  % enet all add_rm vs ols|har: %
\providecommand{\rhchEnetBarAddRmP}{}\renewcommand{\rhchEnetBarAddRmP}{\ensuremath{<}0.001}  % enet all add_rm: DM p
\providecommand{\rhchEnetBarLoRmPct}{}\renewcommand{\rhchEnetBarLoRmPct}{2.0}  % enet all lo_rm vs enet|full: %
\providecommand{\rhchEnetBarLoRmP}{}\renewcommand{\rhchEnetBarLoRmP}{\ensuremath{<}0.001}  % enet all lo_rm: DM p
\providecommand{\rhchEnetBarAddLiqPct}{}\renewcommand{\rhchEnetBarAddLiqPct}{$-$0.6}  % enet all add_liq vs ols|har: %
\providecommand{\rhchEnetBarAddLiqP}{}\renewcommand{\rhchEnetBarAddLiqP}{0.124}  % enet all add_liq: DM p
\providecommand{\rhchEnetBarLoLiqPct}{}\renewcommand{\rhchEnetBarLoLiqPct}{0.1}  % enet all lo_liq vs enet|full: %
\providecommand{\rhchEnetBarLoLiqP}{}\renewcommand{\rhchEnetBarLoLiqP}{0.638}  % enet all lo_liq: DM p
\providecommand{\rhchEnetBarAddEwPct}{}\renewcommand{\rhchEnetBarAddEwPct}{$-$1.1}  % enet all add_ew vs ols|har: %
\providecommand{\rhchEnetBarAddEwP}{}\renewcommand{\rhchEnetBarAddEwP}{0.005}  % enet all add_ew: DM p
\providecommand{\rhchEnetBarLoEwPct}{}\renewcommand{\rhchEnetBarLoEwPct}{$-$0.1}  % enet all lo_ew vs enet|full: %
\providecommand{\rhchEnetBarLoEwP}{}\renewcommand{\rhchEnetBarLoEwP}{0.327}  % enet all lo_ew: DM p
\providecommand{\rhchEnetBarAddVwPct}{}\renewcommand{\rhchEnetBarAddVwPct}{$-$1.1}  % enet all add_vw vs ols|har: %
\providecommand{\rhchEnetBarAddVwP}{}\renewcommand{\rhchEnetBarAddVwP}{0.042}  % enet all add_vw: DM p
\providecommand{\rhchEnetBarLoVwPct}{}\renewcommand{\rhchEnetBarLoVwPct}{$-$0.1}  % enet all lo_vw vs enet|full: %
\providecommand{\rhchEnetBarLoVwP}{}\renewcommand{\rhchEnetBarLoVwP}{0.492}  % enet all lo_vw: DM p
\providecommand{\rhchEnetBarAddSentPct}{}\renewcommand{\rhchEnetBarAddSentPct}{0.0}  % enet all add_sent vs ols|har: %
\providecommand{\rhchEnetBarAddSentP}{}\renewcommand{\rhchEnetBarAddSentP}{0.971}  % enet all add_sent: DM p
\providecommand{\rhchEnetBarLoSentPct}{}\renewcommand{\rhchEnetBarLoSentPct}{0.2}  % enet all lo_sent vs enet|full: %
\providecommand{\rhchEnetBarLoSentP}{}\renewcommand{\rhchEnetBarLoSentP}{0.137}  % enet all lo_sent: DM p
\providecommand{\rhchEnetBarAddIvPct}{}\renewcommand{\rhchEnetBarAddIvPct}{0.1}  % enet all add_iv vs ols|har: %
\providecommand{\rhchEnetBarAddIvP}{}\renewcommand{\rhchEnetBarAddIvP}{0.817}  % enet all add_iv: DM p
\providecommand{\rhchEnetBarLoIvPct}{}\renewcommand{\rhchEnetBarLoIvPct}{0.0}  % enet all lo_iv vs enet|full: %
\providecommand{\rhchEnetBarLoIvP}{}\renewcommand{\rhchEnetBarLoIvP}{0.950}  % enet all lo_iv: DM p
\providecommand{\rhchEnetBarAddVdPct}{}\renewcommand{\rhchEnetBarAddVdPct}{0.4}  % enet all add_vd vs ols|har: %
\providecommand{\rhchEnetBarAddVdP}{}\renewcommand{\rhchEnetBarAddVdP}{0.003}  % enet all add_vd: DM p
\providecommand{\rhchEnetBarLoVdPct}{}\renewcommand{\rhchEnetBarLoVdPct}{$-$0.1}  % enet all lo_vd vs enet|full: %
\providecommand{\rhchEnetBarLoVdP}{}\renewcommand{\rhchEnetBarLoVdP}{0.264}  % enet all lo_vd: DM p
\providecommand{\rhchEnetBarAddFomcPct}{}\renewcommand{\rhchEnetBarAddFomcPct}{$-$3.1}  % enet all add_fomc vs ols|har: %
\providecommand{\rhchEnetBarAddFomcP}{}\renewcommand{\rhchEnetBarAddFomcP}{\ensuremath{<}0.001}  % enet all add_fomc: DM p
\providecommand{\rhchEnetBarLoFomcPct}{}\renewcommand{\rhchEnetBarLoFomcPct}{3.1}  % enet all lo_fomc vs enet|full: %
\providecommand{\rhchEnetBarLoFomcP}{}\renewcommand{\rhchEnetBarLoFomcP}{\ensuremath{<}0.001}  % enet all lo_fomc: DM p
\providecommand{\rhchEnetEventJointPct}{}\renewcommand{\rhchEnetEventJointPct}{$-$22.6}  % enet event: full vs HAR, % QLIKE change
\providecommand{\rhchEnetEventJointDm}{}\renewcommand{\rhchEnetEventJointDm}{$-$6.79}  % enet event: full vs HAR, DM
\providecommand{\rhchEnetEventJointP}{}\renewcommand{\rhchEnetEventJointP}{\ensuremath{<}0.001}  % enet event: full vs HAR, DM p
\providecommand{\rhchEnetEventAddRmPct}{}\renewcommand{\rhchEnetEventAddRmPct}{4.2}  % enet event add_rm vs ols|har: %
\providecommand{\rhchEnetEventAddRmP}{}\renewcommand{\rhchEnetEventAddRmP}{0.031}  % enet event add_rm: DM p
\providecommand{\rhchEnetEventLoRmPct}{}\renewcommand{\rhchEnetEventLoRmPct}{$-$1.2}  % enet event lo_rm vs enet|full: %
\providecommand{\rhchEnetEventLoRmP}{}\renewcommand{\rhchEnetEventLoRmP}{0.167}  % enet event lo_rm: DM p
\providecommand{\rhchEnetEventAddLiqPct}{}\renewcommand{\rhchEnetEventAddLiqPct}{5.3}  % enet event add_liq vs ols|har: %
\providecommand{\rhchEnetEventAddLiqP}{}\renewcommand{\rhchEnetEventAddLiqP}{0.008}  % enet event add_liq: DM p
\providecommand{\rhchEnetEventLoLiqPct}{}\renewcommand{\rhchEnetEventLoLiqPct}{$-$0.7}  % enet event lo_liq vs enet|full: %
\providecommand{\rhchEnetEventLoLiqP}{}\renewcommand{\rhchEnetEventLoLiqP}{0.104}  % enet event lo_liq: DM p
\providecommand{\rhchEnetEventAddEwPct}{}\renewcommand{\rhchEnetEventAddEwPct}{$-$2.7}  % enet event add_ew vs ols|har: %
\providecommand{\rhchEnetEventAddEwP}{}\renewcommand{\rhchEnetEventAddEwP}{0.060}  % enet event add_ew: DM p
\providecommand{\rhchEnetEventLoEwPct}{}\renewcommand{\rhchEnetEventLoEwPct}{0.0}  % enet event lo_ew vs enet|full: %
\providecommand{\rhchEnetEventLoEwP}{}\renewcommand{\rhchEnetEventLoEwP}{0.915}  % enet event lo_ew: DM p
\providecommand{\rhchEnetEventAddVwPct}{}\renewcommand{\rhchEnetEventAddVwPct}{$-$3.4}  % enet event add_vw vs ols|har: %
\providecommand{\rhchEnetEventAddVwP}{}\renewcommand{\rhchEnetEventAddVwP}{0.055}  % enet event add_vw: DM p
\providecommand{\rhchEnetEventLoVwPct}{}\renewcommand{\rhchEnetEventLoVwPct}{0.0}  % enet event lo_vw vs enet|full: %
\providecommand{\rhchEnetEventLoVwP}{}\renewcommand{\rhchEnetEventLoVwP}{0.947}  % enet event lo_vw: DM p
\providecommand{\rhchEnetEventAddSentPct}{}\renewcommand{\rhchEnetEventAddSentPct}{$-$0.9}  % enet event add_sent vs ols|har: %
\providecommand{\rhchEnetEventAddSentP}{}\renewcommand{\rhchEnetEventAddSentP}{0.519}  % enet event add_sent: DM p
\providecommand{\rhchEnetEventLoSentPct}{}\renewcommand{\rhchEnetEventLoSentPct}{1.6}  % enet event lo_sent vs enet|full: %
\providecommand{\rhchEnetEventLoSentP}{}\renewcommand{\rhchEnetEventLoSentP}{0.201}  % enet event lo_sent: DM p
\providecommand{\rhchEnetEventAddIvPct}{}\renewcommand{\rhchEnetEventAddIvPct}{2.5}  % enet event add_iv vs ols|har: %
\providecommand{\rhchEnetEventAddIvP}{}\renewcommand{\rhchEnetEventAddIvP}{0.010}  % enet event add_iv: DM p
\providecommand{\rhchEnetEventLoIvPct}{}\renewcommand{\rhchEnetEventLoIvPct}{$-$0.1}  % enet event lo_iv vs enet|full: %
\providecommand{\rhchEnetEventLoIvP}{}\renewcommand{\rhchEnetEventLoIvP}{0.865}  % enet event lo_iv: DM p
\providecommand{\rhchEnetEventAddVdPct}{}\renewcommand{\rhchEnetEventAddVdPct}{1.6}  % enet event add_vd vs ols|har: %
\providecommand{\rhchEnetEventAddVdP}{}\renewcommand{\rhchEnetEventAddVdP}{0.007}  % enet event add_vd: DM p
\providecommand{\rhchEnetEventLoVdPct}{}\renewcommand{\rhchEnetEventLoVdPct}{$-$0.1}  % enet event lo_vd vs enet|full: %
\providecommand{\rhchEnetEventLoVdP}{}\renewcommand{\rhchEnetEventLoVdP}{0.330}  % enet event lo_vd: DM p
\providecommand{\rhchEnetEventAddFomcPct}{}\renewcommand{\rhchEnetEventAddFomcPct}{$-$23.2}  % enet event add_fomc vs ols|har: %
\providecommand{\rhchEnetEventAddFomcP}{}\renewcommand{\rhchEnetEventAddFomcP}{\ensuremath{<}0.001}  % enet event add_fomc: DM p
\providecommand{\rhchEnetEventLoFomcPct}{}\renewcommand{\rhchEnetEventLoFomcPct}{28.2}  % enet event lo_fomc vs enet|full: %
\providecommand{\rhchEnetEventLoFomcP}{}\renewcommand{\rhchEnetEventLoFomcP}{\ensuremath{<}0.001}  % enet event lo_fomc: DM p
\providecommand{\rhchEnetEventReading}{}\renewcommand{\rhchEnetEventReading}{concentrated}  % enet event: item 3 reading
\providecommand{\rhchEnetEventTopAddShare}{}\renewcommand{\rhchEnetEventTopAddShare}{103}  % enet event: largest add-alone gain, % of the joint gain (fomc)
\providecommand{\rhchEnetEventSumAddRatio}{}\renewcommand{\rhchEnetEventSumAddRatio}{0.74}  % enet event: add-alone gains summed over groups / joint gain
\providecommand{\rhchEnetEventNLoSig}{}\renewcommand{\rhchEnetEventNLoSig}{1}  % enet event: significant leave-one-outs
\providecommand{\rhchEnetEventNLoRun}{}\renewcommand{\rhchEnetEventNLoRun}{8}  % enet event: leave-one-outs available (of 8)
\providecommand{\rhchEnetEventConcTest}{}\renewcommand{\rhchEnetEventConcTest}{yes}  % enet event: the CONC definition on its own (pending: undetermined)
\providecommand{\rhchEnetEventRedTest}{}\renewcommand{\rhchEnetEventRedTest}{no}  % enet event: the RED definition on its own (pending: undetermined)
\providecommand{\rhchEnetEventReadingRedFirst}{}\renewcommand{\rhchEnetEventReadingRedFirst}{concentrated}  % enet event: reading with REDUNDANT tested before CONCENTRATED (the RULE tests CONCENTRATED first)
\providecommand{\rhchEnetOrdJointPct}{}\renewcommand{\rhchEnetOrdJointPct}{$-$4.0}  % enet ord: full vs HAR, % QLIKE change
\providecommand{\rhchEnetOrdJointDm}{}\renewcommand{\rhchEnetOrdJointDm}{$-$4.73}  % enet ord: full vs HAR, DM
\providecommand{\rhchEnetOrdJointP}{}\renewcommand{\rhchEnetOrdJointP}{\ensuremath{<}0.001}  % enet ord: full vs HAR, DM p
\providecommand{\rhchEnetOrdAddRmPct}{}\renewcommand{\rhchEnetOrdAddRmPct}{$-$3.9}  % enet ord add_rm vs ols|har: %
\providecommand{\rhchEnetOrdAddRmP}{}\renewcommand{\rhchEnetOrdAddRmP}{\ensuremath{<}0.001}  % enet ord add_rm: DM p
\providecommand{\rhchEnetOrdLoRmPct}{}\renewcommand{\rhchEnetOrdLoRmPct}{2.5}  % enet ord lo_rm vs enet|full: %
\providecommand{\rhchEnetOrdLoRmP}{}\renewcommand{\rhchEnetOrdLoRmP}{\ensuremath{<}0.001}  % enet ord lo_rm: DM p
\providecommand{\rhchEnetOrdAddLiqPct}{}\renewcommand{\rhchEnetOrdAddLiqPct}{$-$1.6}  % enet ord add_liq vs ols|har: %
\providecommand{\rhchEnetOrdAddLiqP}{}\renewcommand{\rhchEnetOrdAddLiqP}{\ensuremath{<}0.001}  % enet ord add_liq: DM p
\providecommand{\rhchEnetOrdLoLiqPct}{}\renewcommand{\rhchEnetOrdLoLiqPct}{0.2}  % enet ord lo_liq vs enet|full: %
\providecommand{\rhchEnetOrdLoLiqP}{}\renewcommand{\rhchEnetOrdLoLiqP}{0.282}  % enet ord lo_liq: DM p
\providecommand{\rhchEnetOrdAddEwPct}{}\renewcommand{\rhchEnetOrdAddEwPct}{$-$0.9}  % enet ord add_ew vs ols|har: %
\providecommand{\rhchEnetOrdAddEwP}{}\renewcommand{\rhchEnetOrdAddEwP}{0.018}  % enet ord add_ew: DM p
\providecommand{\rhchEnetOrdLoEwPct}{}\renewcommand{\rhchEnetOrdLoEwPct}{$-$0.1}  % enet ord lo_ew vs enet|full: %
\providecommand{\rhchEnetOrdLoEwP}{}\renewcommand{\rhchEnetOrdLoEwP}{0.220}  % enet ord lo_ew: DM p
\providecommand{\rhchEnetOrdAddVwPct}{}\renewcommand{\rhchEnetOrdAddVwPct}{$-$0.7}  % enet ord add_vw vs ols|har: %
\providecommand{\rhchEnetOrdAddVwP}{}\renewcommand{\rhchEnetOrdAddVwP}{0.195}  % enet ord add_vw: DM p
\providecommand{\rhchEnetOrdLoVwPct}{}\renewcommand{\rhchEnetOrdLoVwPct}{$-$0.1}  % enet ord lo_vw vs enet|full: %
\providecommand{\rhchEnetOrdLoVwP}{}\renewcommand{\rhchEnetOrdLoVwP}{0.491}  % enet ord lo_vw: DM p
\providecommand{\rhchEnetOrdAddSentPct}{}\renewcommand{\rhchEnetOrdAddSentPct}{0.2}  % enet ord add_sent vs ols|har: %
\providecommand{\rhchEnetOrdAddSentP}{}\renewcommand{\rhchEnetOrdAddSentP}{0.407}  % enet ord add_sent: DM p
\providecommand{\rhchEnetOrdLoSentPct}{}\renewcommand{\rhchEnetOrdLoSentPct}{0.0}  % enet ord lo_sent vs enet|full: %
\providecommand{\rhchEnetOrdLoSentP}{}\renewcommand{\rhchEnetOrdLoSentP}{0.998}  % enet ord lo_sent: DM p
\providecommand{\rhchEnetOrdAddIvPct}{}\renewcommand{\rhchEnetOrdAddIvPct}{$-$0.3}  % enet ord add_iv vs ols|har: %
\providecommand{\rhchEnetOrdAddIvP}{}\renewcommand{\rhchEnetOrdAddIvP}{0.319}  % enet ord add_iv: DM p
\providecommand{\rhchEnetOrdLoIvPct}{}\renewcommand{\rhchEnetOrdLoIvPct}{0.0}  % enet ord lo_iv vs enet|full: %
\providecommand{\rhchEnetOrdLoIvP}{}\renewcommand{\rhchEnetOrdLoIvP}{0.912}  % enet ord lo_iv: DM p
\providecommand{\rhchEnetOrdAddVdPct}{}\renewcommand{\rhchEnetOrdAddVdPct}{0.2}  % enet ord add_vd vs ols|har: %
\providecommand{\rhchEnetOrdAddVdP}{}\renewcommand{\rhchEnetOrdAddVdP}{0.127}  % enet ord add_vd: DM p
\providecommand{\rhchEnetOrdLoVdPct}{}\renewcommand{\rhchEnetOrdLoVdPct}{$-$0.1}  % enet ord lo_vd vs enet|full: %
\providecommand{\rhchEnetOrdLoVdP}{}\renewcommand{\rhchEnetOrdLoVdP}{0.345}  % enet ord lo_vd: DM p
\providecommand{\rhchEnetOrdAddFomcPct}{}\renewcommand{\rhchEnetOrdAddFomcPct}{0.4}  % enet ord add_fomc vs ols|har: %
\providecommand{\rhchEnetOrdAddFomcP}{}\renewcommand{\rhchEnetOrdAddFomcP}{\ensuremath{<}0.001}  % enet ord add_fomc: DM p
\providecommand{\rhchEnetOrdLoFomcPct}{}\renewcommand{\rhchEnetOrdLoFomcPct}{$-$0.4}  % enet ord lo_fomc vs enet|full: %
\providecommand{\rhchEnetOrdLoFomcP}{}\renewcommand{\rhchEnetOrdLoFomcP}{\ensuremath{<}0.001}  % enet ord lo_fomc: DM p
\providecommand{\rhchEnetOrdReading}{}\renewcommand{\rhchEnetOrdReading}{concentrated}  % enet ord: item 3 reading
\providecommand{\rhchEnetOrdTopAddShare}{}\renewcommand{\rhchEnetOrdTopAddShare}{97}  % enet ord: largest add-alone gain, % of the joint gain (rm)
\providecommand{\rhchEnetOrdSumAddRatio}{}\renewcommand{\rhchEnetOrdSumAddRatio}{1.68}  % enet ord: add-alone gains summed over groups / joint gain
\providecommand{\rhchEnetOrdNLoSig}{}\renewcommand{\rhchEnetOrdNLoSig}{1}  % enet ord: significant leave-one-outs
\providecommand{\rhchEnetOrdNLoRun}{}\renewcommand{\rhchEnetOrdNLoRun}{8}  % enet ord: leave-one-outs available (of 8)
\providecommand{\rhchEnetOrdConcTest}{}\renewcommand{\rhchEnetOrdConcTest}{yes}  % enet ord: the CONC definition on its own (pending: undetermined)
\providecommand{\rhchEnetOrdRedTest}{}\renewcommand{\rhchEnetOrdRedTest}{yes}  % enet ord: the RED definition on its own (pending: undetermined)
\providecommand{\rhchEnetOrdReadingRedFirst}{}\renewcommand{\rhchEnetOrdReadingRedFirst}{redundant}  % enet ord: reading with REDUNDANT tested before CONCENTRATED (the RULE tests CONCENTRATED first)
\providecommand{\rhchLgbmBarJointPct}{}\renewcommand{\rhchLgbmBarJointPct}{$-$14.1}  % lgbm all: full vs HAR, % QLIKE change
\providecommand{\rhchLgbmBarJointDm}{}\renewcommand{\rhchLgbmBarJointDm}{$-$10.26}  % lgbm all: full vs HAR, DM
\providecommand{\rhchLgbmBarJointP}{}\renewcommand{\rhchLgbmBarJointP}{\ensuremath{<}0.001}  % lgbm all: full vs HAR, DM p
\providecommand{\rhchLgbmBarAddRmPct}{}\renewcommand{\rhchLgbmBarAddRmPct}{$-$3.0}  % lgbm all add_rm vs lgbm|har: %
\providecommand{\rhchLgbmBarAddRmP}{}\renewcommand{\rhchLgbmBarAddRmP}{\ensuremath{<}0.001}  % lgbm all add_rm: DM p
\providecommand{\rhchLgbmBarLoRmPct}{}\renewcommand{\rhchLgbmBarLoRmPct}{1.5}  % lgbm all lo_rm vs lgbm|full_pre: % [pre-patch]
\providecommand{\rhchLgbmBarLoRmP}{}\renewcommand{\rhchLgbmBarLoRmP}{\ensuremath{<}0.001}  % lgbm all lo_rm: DM p [pre-patch]
\providecommand{\rhchLgbmBarAddLiqPct}{}\renewcommand{\rhchLgbmBarAddLiqPct}{$-$2.9}  % lgbm all add_liq vs lgbm|har: %
\providecommand{\rhchLgbmBarAddLiqP}{}\renewcommand{\rhchLgbmBarAddLiqP}{\ensuremath{<}0.001}  % lgbm all add_liq: DM p
\providecommand{\rhchLgbmBarLoLiqPct}{}\renewcommand{\rhchLgbmBarLoLiqPct}{0.2}  % lgbm all lo_liq vs lgbm|full_pre: % [pre-patch]
\providecommand{\rhchLgbmBarLoLiqP}{}\renewcommand{\rhchLgbmBarLoLiqP}{0.365}  % lgbm all lo_liq: DM p [pre-patch]
\providecommand{\rhchLgbmBarAddEwPct}{}\renewcommand{\rhchLgbmBarAddEwPct}{$-$2.8}  % lgbm all add_ew vs lgbm|har: %
\providecommand{\rhchLgbmBarAddEwP}{}\renewcommand{\rhchLgbmBarAddEwP}{\ensuremath{<}0.001}  % lgbm all add_ew: DM p
\providecommand{\rhchLgbmBarLoEwPct}{}\renewcommand{\rhchLgbmBarLoEwPct}{0.0}  % lgbm all lo_ew vs lgbm|full_pre: % [pre-patch]
\providecommand{\rhchLgbmBarLoEwP}{}\renewcommand{\rhchLgbmBarLoEwP}{0.815}  % lgbm all lo_ew: DM p [pre-patch]
\providecommand{\rhchLgbmBarAddVwPct}{}\renewcommand{\rhchLgbmBarAddVwPct}{$-$3.1}  % lgbm all add_vw vs lgbm|har: %
\providecommand{\rhchLgbmBarAddVwP}{}\renewcommand{\rhchLgbmBarAddVwP}{\ensuremath{<}0.001}  % lgbm all add_vw: DM p
\providecommand{\rhchLgbmBarLoVwPct}{}\renewcommand{\rhchLgbmBarLoVwPct}{0.1}  % lgbm all lo_vw vs lgbm|full_pre: % [pre-patch]
\providecommand{\rhchLgbmBarLoVwP}{}\renewcommand{\rhchLgbmBarLoVwP}{0.489}  % lgbm all lo_vw: DM p [pre-patch]
\providecommand{\rhchLgbmBarAddSentPct}{}\renewcommand{\rhchLgbmBarAddSentPct}{$-$1.1}  % lgbm all add_sent vs lgbm|har: %
\providecommand{\rhchLgbmBarAddSentP}{}\renewcommand{\rhchLgbmBarAddSentP}{\ensuremath{<}0.001}  % lgbm all add_sent: DM p
\providecommand{\rhchLgbmBarLoSentPct}{}\renewcommand{\rhchLgbmBarLoSentPct}{0.2}  % lgbm all lo_sent vs lgbm|full_pre: % [pre-patch]
\providecommand{\rhchLgbmBarLoSentP}{}\renewcommand{\rhchLgbmBarLoSentP}{0.028}  % lgbm all lo_sent: DM p [pre-patch]
\providecommand{\rhchLgbmBarAddIvPct}{}\renewcommand{\rhchLgbmBarAddIvPct}{$-$1.3}  % lgbm all add_iv vs lgbm|har: %
\providecommand{\rhchLgbmBarAddIvP}{}\renewcommand{\rhchLgbmBarAddIvP}{\ensuremath{<}0.001}  % lgbm all add_iv: DM p
\providecommand{\rhchLgbmBarLoIvPct}{}\renewcommand{\rhchLgbmBarLoIvPct}{0.2}  % lgbm all lo_iv vs lgbm|full_pre: % [pre-patch]
\providecommand{\rhchLgbmBarLoIvP}{}\renewcommand{\rhchLgbmBarLoIvP}{0.084}  % lgbm all lo_iv: DM p [pre-patch]
\providecommand{\rhchLgbmBarAddVdPct}{}\renewcommand{\rhchLgbmBarAddVdPct}{$-$1.2}  % lgbm all add_vd vs lgbm|har: %
\providecommand{\rhchLgbmBarAddVdP}{}\renewcommand{\rhchLgbmBarAddVdP}{\ensuremath{<}0.001}  % lgbm all add_vd: DM p
\providecommand{\rhchLgbmBarLoVdPct}{}\renewcommand{\rhchLgbmBarLoVdPct}{0.0}  % lgbm all lo_vd vs lgbm|full_pre: % [pre-patch]
\providecommand{\rhchLgbmBarLoVdP}{}\renewcommand{\rhchLgbmBarLoVdP}{0.916}  % lgbm all lo_vd: DM p [pre-patch]
\providecommand{\rhchLgbmBarAddFomcPct}{}\renewcommand{\rhchLgbmBarAddFomcPct}{$-$9.2}  % lgbm all add_fomc vs lgbm|har: %
\providecommand{\rhchLgbmBarAddFomcP}{}\renewcommand{\rhchLgbmBarAddFomcP}{\ensuremath{<}0.001}  % lgbm all add_fomc: DM p
\providecommand{\rhchLgbmBarLoFomcPct}{}\renewcommand{\rhchLgbmBarLoFomcPct}{10.3}  % lgbm all lo_fomc vs lgbm|full: %
\providecommand{\rhchLgbmBarLoFomcP}{}\renewcommand{\rhchLgbmBarLoFomcP}{\ensuremath{<}0.001}  % lgbm all lo_fomc: DM p
\providecommand{\rhchLgbmEventJointPct}{}\renewcommand{\rhchLgbmEventJointPct}{$-$62.3}  % lgbm event: full vs HAR, % QLIKE change
\providecommand{\rhchLgbmEventJointDm}{}\renewcommand{\rhchLgbmEventJointDm}{$-$8.62}  % lgbm event: full vs HAR, DM
\providecommand{\rhchLgbmEventJointP}{}\renewcommand{\rhchLgbmEventJointP}{\ensuremath{<}0.001}  % lgbm event: full vs HAR, DM p
\providecommand{\rhchLgbmEventAddRmPct}{}\renewcommand{\rhchLgbmEventAddRmPct}{6.8}  % lgbm event add_rm vs lgbm|har: %
\providecommand{\rhchLgbmEventAddRmP}{}\renewcommand{\rhchLgbmEventAddRmP}{0.010}  % lgbm event add_rm: DM p
\providecommand{\rhchLgbmEventLoRmPct}{}\renewcommand{\rhchLgbmEventLoRmPct}{$-$0.3}  % lgbm event lo_rm vs lgbm|full_pre: % [pre-patch]
\providecommand{\rhchLgbmEventLoRmP}{}\renewcommand{\rhchLgbmEventLoRmP}{0.700}  % lgbm event lo_rm: DM p [pre-patch]
\providecommand{\rhchLgbmEventAddLiqPct}{}\renewcommand{\rhchLgbmEventAddLiqPct}{1.5}  % lgbm event add_liq vs lgbm|har: %
\providecommand{\rhchLgbmEventAddLiqP}{}\renewcommand{\rhchLgbmEventAddLiqP}{0.294}  % lgbm event add_liq: DM p
\providecommand{\rhchLgbmEventLoLiqPct}{}\renewcommand{\rhchLgbmEventLoLiqPct}{$-$0.5}  % lgbm event lo_liq vs lgbm|full_pre: % [pre-patch]
\providecommand{\rhchLgbmEventLoLiqP}{}\renewcommand{\rhchLgbmEventLoLiqP}{0.155}  % lgbm event lo_liq: DM p [pre-patch]
\providecommand{\rhchLgbmEventAddEwPct}{}\renewcommand{\rhchLgbmEventAddEwPct}{$-$3.5}  % lgbm event add_ew vs lgbm|har: %
\providecommand{\rhchLgbmEventAddEwP}{}\renewcommand{\rhchLgbmEventAddEwP}{0.045}  % lgbm event add_ew: DM p
\providecommand{\rhchLgbmEventLoEwPct}{}\renewcommand{\rhchLgbmEventLoEwPct}{0.1}  % lgbm event lo_ew vs lgbm|full_pre: % [pre-patch]
\providecommand{\rhchLgbmEventLoEwP}{}\renewcommand{\rhchLgbmEventLoEwP}{0.888}  % lgbm event lo_ew: DM p [pre-patch]
\providecommand{\rhchLgbmEventAddVwPct}{}\renewcommand{\rhchLgbmEventAddVwPct}{$-$3.5}  % lgbm event add_vw vs lgbm|har: %
\providecommand{\rhchLgbmEventAddVwP}{}\renewcommand{\rhchLgbmEventAddVwP}{0.044}  % lgbm event add_vw: DM p
\providecommand{\rhchLgbmEventLoVwPct}{}\renewcommand{\rhchLgbmEventLoVwPct}{$-$0.1}  % lgbm event lo_vw vs lgbm|full_pre: % [pre-patch]
\providecommand{\rhchLgbmEventLoVwP}{}\renewcommand{\rhchLgbmEventLoVwP}{0.804}  % lgbm event lo_vw: DM p [pre-patch]
\providecommand{\rhchLgbmEventAddSentPct}{}\renewcommand{\rhchLgbmEventAddSentPct}{0.5}  % lgbm event add_sent vs lgbm|har: %
\providecommand{\rhchLgbmEventAddSentP}{}\renewcommand{\rhchLgbmEventAddSentP}{0.574}  % lgbm event add_sent: DM p
\providecommand{\rhchLgbmEventLoSentPct}{}\renewcommand{\rhchLgbmEventLoSentPct}{0.2}  % lgbm event lo_sent vs lgbm|full_pre: % [pre-patch]
\providecommand{\rhchLgbmEventLoSentP}{}\renewcommand{\rhchLgbmEventLoSentP}{0.481}  % lgbm event lo_sent: DM p [pre-patch]
\providecommand{\rhchLgbmEventAddIvPct}{}\renewcommand{\rhchLgbmEventAddIvPct}{$-$1.4}  % lgbm event add_iv vs lgbm|har: %
\providecommand{\rhchLgbmEventAddIvP}{}\renewcommand{\rhchLgbmEventAddIvP}{0.036}  % lgbm event add_iv: DM p
\providecommand{\rhchLgbmEventLoIvPct}{}\renewcommand{\rhchLgbmEventLoIvPct}{0.7}  % lgbm event lo_iv vs lgbm|full_pre: % [pre-patch]
\providecommand{\rhchLgbmEventLoIvP}{}\renewcommand{\rhchLgbmEventLoIvP}{0.129}  % lgbm event lo_iv: DM p [pre-patch]
\providecommand{\rhchLgbmEventAddVdPct}{}\renewcommand{\rhchLgbmEventAddVdPct}{$-$3.5}  % lgbm event add_vd vs lgbm|har: %
\providecommand{\rhchLgbmEventAddVdP}{}\renewcommand{\rhchLgbmEventAddVdP}{0.072}  % lgbm event add_vd: DM p
\providecommand{\rhchLgbmEventLoVdPct}{}\renewcommand{\rhchLgbmEventLoVdPct}{0.7}  % lgbm event lo_vd vs lgbm|full_pre: % [pre-patch]
\providecommand{\rhchLgbmEventLoVdP}{}\renewcommand{\rhchLgbmEventLoVdP}{0.010}  % lgbm event lo_vd: DM p [pre-patch]
\providecommand{\rhchLgbmEventAddFomcPct}{}\renewcommand{\rhchLgbmEventAddFomcPct}{$-$62.6}  % lgbm event add_fomc vs lgbm|har: %
\providecommand{\rhchLgbmEventAddFomcP}{}\renewcommand{\rhchLgbmEventAddFomcP}{\ensuremath{<}0.001}  % lgbm event add_fomc: DM p
\providecommand{\rhchLgbmEventLoFomcPct}{}\renewcommand{\rhchLgbmEventLoFomcPct}{163.2}  % lgbm event lo_fomc vs lgbm|full: %
\providecommand{\rhchLgbmEventLoFomcP}{}\renewcommand{\rhchLgbmEventLoFomcP}{\ensuremath{<}0.001}  % lgbm event lo_fomc: DM p
\providecommand{\rhchLgbmEventReading}{}\renewcommand{\rhchLgbmEventReading}{concentrated}  % lgbm event: item 3 reading
\providecommand{\rhchLgbmEventTopAddShare}{}\renewcommand{\rhchLgbmEventTopAddShare}{101}  % lgbm event: largest add-alone gain, % of the joint gain (fomc)
\providecommand{\rhchLgbmEventSumAddRatio}{}\renewcommand{\rhchLgbmEventSumAddRatio}{1.05}  % lgbm event: add-alone gains summed over groups / joint gain
\providecommand{\rhchLgbmEventNLoSig}{}\renewcommand{\rhchLgbmEventNLoSig}{2}  % lgbm event: significant leave-one-outs
\providecommand{\rhchLgbmEventNLoRun}{}\renewcommand{\rhchLgbmEventNLoRun}{8}  % lgbm event: leave-one-outs available (of 8)
\providecommand{\rhchLgbmEventConcTest}{}\renewcommand{\rhchLgbmEventConcTest}{yes}  % lgbm event: the CONC definition on its own (pending: undetermined)
\providecommand{\rhchLgbmEventRedTest}{}\renewcommand{\rhchLgbmEventRedTest}{no}  % lgbm event: the RED definition on its own (pending: undetermined)
\providecommand{\rhchLgbmEventReadingRedFirst}{}\renewcommand{\rhchLgbmEventReadingRedFirst}{concentrated}  % lgbm event: reading with REDUNDANT tested before CONCENTRATED (the RULE tests CONCENTRATED first)
\providecommand{\rhchLgbmOrdJointPct}{}\renewcommand{\rhchLgbmOrdJointPct}{$-$6.0}  % lgbm ord: full vs HAR, % QLIKE change
\providecommand{\rhchLgbmOrdJointDm}{}\renewcommand{\rhchLgbmOrdJointDm}{$-$10.70}  % lgbm ord: full vs HAR, DM
\providecommand{\rhchLgbmOrdJointP}{}\renewcommand{\rhchLgbmOrdJointP}{\ensuremath{<}0.001}  % lgbm ord: full vs HAR, DM p
\providecommand{\rhchLgbmOrdAddRmPct}{}\renewcommand{\rhchLgbmOrdAddRmPct}{$-$4.6}  % lgbm ord add_rm vs lgbm|har: %
\providecommand{\rhchLgbmOrdAddRmP}{}\renewcommand{\rhchLgbmOrdAddRmP}{\ensuremath{<}0.001}  % lgbm ord add_rm: DM p
\providecommand{\rhchLgbmOrdLoRmPct}{}\renewcommand{\rhchLgbmOrdLoRmPct}{1.6}  % lgbm ord lo_rm vs lgbm|full_pre: % [pre-patch]
\providecommand{\rhchLgbmOrdLoRmP}{}\renewcommand{\rhchLgbmOrdLoRmP}{\ensuremath{<}0.001}  % lgbm ord lo_rm: DM p [pre-patch]
\providecommand{\rhchLgbmOrdAddLiqPct}{}\renewcommand{\rhchLgbmOrdAddLiqPct}{$-$3.7}  % lgbm ord add_liq vs lgbm|har: %
\providecommand{\rhchLgbmOrdAddLiqP}{}\renewcommand{\rhchLgbmOrdAddLiqP}{\ensuremath{<}0.001}  % lgbm ord add_liq: DM p
\providecommand{\rhchLgbmOrdLoLiqPct}{}\renewcommand{\rhchLgbmOrdLoLiqPct}{0.2}  % lgbm ord lo_liq vs lgbm|full_pre: % [pre-patch]
\providecommand{\rhchLgbmOrdLoLiqP}{}\renewcommand{\rhchLgbmOrdLoLiqP}{0.266}  % lgbm ord lo_liq: DM p [pre-patch]
\providecommand{\rhchLgbmOrdAddEwPct}{}\renewcommand{\rhchLgbmOrdAddEwPct}{$-$2.7}  % lgbm ord add_ew vs lgbm|har: %
\providecommand{\rhchLgbmOrdAddEwP}{}\renewcommand{\rhchLgbmOrdAddEwP}{\ensuremath{<}0.001}  % lgbm ord add_ew: DM p
\providecommand{\rhchLgbmOrdLoEwPct}{}\renewcommand{\rhchLgbmOrdLoEwPct}{0.0}  % lgbm ord lo_ew vs lgbm|full_pre: % [pre-patch]
\providecommand{\rhchLgbmOrdLoEwP}{}\renewcommand{\rhchLgbmOrdLoEwP}{0.779}  % lgbm ord lo_ew: DM p [pre-patch]
\providecommand{\rhchLgbmOrdAddVwPct}{}\renewcommand{\rhchLgbmOrdAddVwPct}{$-$3.1}  % lgbm ord add_vw vs lgbm|har: %
\providecommand{\rhchLgbmOrdAddVwP}{}\renewcommand{\rhchLgbmOrdAddVwP}{\ensuremath{<}0.001}  % lgbm ord add_vw: DM p
\providecommand{\rhchLgbmOrdLoVwPct}{}\renewcommand{\rhchLgbmOrdLoVwPct}{0.1}  % lgbm ord lo_vw vs lgbm|full_pre: % [pre-patch]
\providecommand{\rhchLgbmOrdLoVwP}{}\renewcommand{\rhchLgbmOrdLoVwP}{0.433}  % lgbm ord lo_vw: DM p [pre-patch]
\providecommand{\rhchLgbmOrdAddSentPct}{}\renewcommand{\rhchLgbmOrdAddSentPct}{$-$1.4}  % lgbm ord add_sent vs lgbm|har: %
\providecommand{\rhchLgbmOrdAddSentP}{}\renewcommand{\rhchLgbmOrdAddSentP}{\ensuremath{<}0.001}  % lgbm ord add_sent: DM p
\providecommand{\rhchLgbmOrdLoSentPct}{}\renewcommand{\rhchLgbmOrdLoSentPct}{0.2}  % lgbm ord lo_sent vs lgbm|full_pre: % [pre-patch]
\providecommand{\rhchLgbmOrdLoSentP}{}\renewcommand{\rhchLgbmOrdLoSentP}{0.035}  % lgbm ord lo_sent: DM p [pre-patch]
\providecommand{\rhchLgbmOrdAddIvPct}{}\renewcommand{\rhchLgbmOrdAddIvPct}{$-$1.3}  % lgbm ord add_iv vs lgbm|har: %
\providecommand{\rhchLgbmOrdAddIvP}{}\renewcommand{\rhchLgbmOrdAddIvP}{\ensuremath{<}0.001}  % lgbm ord add_iv: DM p
\providecommand{\rhchLgbmOrdLoIvPct}{}\renewcommand{\rhchLgbmOrdLoIvPct}{0.1}  % lgbm ord lo_iv vs lgbm|full_pre: % [pre-patch]
\providecommand{\rhchLgbmOrdLoIvP}{}\renewcommand{\rhchLgbmOrdLoIvP}{0.196}  % lgbm ord lo_iv: DM p [pre-patch]
\providecommand{\rhchLgbmOrdAddVdPct}{}\renewcommand{\rhchLgbmOrdAddVdPct}{$-$0.9}  % lgbm ord add_vd vs lgbm|har: %
\providecommand{\rhchLgbmOrdAddVdP}{}\renewcommand{\rhchLgbmOrdAddVdP}{\ensuremath{<}0.001}  % lgbm ord add_vd: DM p
\providecommand{\rhchLgbmOrdLoVdPct}{}\renewcommand{\rhchLgbmOrdLoVdPct}{$-$0.1}  % lgbm ord lo_vd vs lgbm|full_pre: % [pre-patch]
\providecommand{\rhchLgbmOrdLoVdP}{}\renewcommand{\rhchLgbmOrdLoVdP}{0.499}  % lgbm ord lo_vd: DM p [pre-patch]
\providecommand{\rhchLgbmOrdAddFomcPct}{}\renewcommand{\rhchLgbmOrdAddFomcPct}{$-$0.2}  % lgbm ord add_fomc vs lgbm|har: %
\providecommand{\rhchLgbmOrdAddFomcP}{}\renewcommand{\rhchLgbmOrdAddFomcP}{0.364}  % lgbm ord add_fomc: DM p
\providecommand{\rhchLgbmOrdLoFomcPct}{}\renewcommand{\rhchLgbmOrdLoFomcPct}{0.0}  % lgbm ord lo_fomc vs lgbm|full: %
\providecommand{\rhchLgbmOrdLoFomcP}{}\renewcommand{\rhchLgbmOrdLoFomcP}{0.912}  % lgbm ord lo_fomc: DM p
\providecommand{\rhchLgbmOrdReading}{}\renewcommand{\rhchLgbmOrdReading}{concentrated}  % lgbm ord: item 3 reading
\providecommand{\rhchLgbmOrdTopAddShare}{}\renewcommand{\rhchLgbmOrdTopAddShare}{77}  % lgbm ord: largest add-alone gain, % of the joint gain (rm)
\providecommand{\rhchLgbmOrdSumAddRatio}{}\renewcommand{\rhchLgbmOrdSumAddRatio}{2.95}  % lgbm ord: add-alone gains summed over groups / joint gain
\providecommand{\rhchLgbmOrdNLoSig}{}\renewcommand{\rhchLgbmOrdNLoSig}{2}  % lgbm ord: significant leave-one-outs
\providecommand{\rhchLgbmOrdNLoRun}{}\renewcommand{\rhchLgbmOrdNLoRun}{8}  % lgbm ord: leave-one-outs available (of 8)
\providecommand{\rhchLgbmOrdConcTest}{}\renewcommand{\rhchLgbmOrdConcTest}{yes}  % lgbm ord: the CONC definition on its own (pending: undetermined)
\providecommand{\rhchLgbmOrdRedTest}{}\renewcommand{\rhchLgbmOrdRedTest}{no}  % lgbm ord: the RED definition on its own (pending: undetermined)
\providecommand{\rhchLgbmOrdReadingRedFirst}{}\renewcommand{\rhchLgbmOrdReadingRedFirst}{concentrated}  % lgbm ord: reading with REDUNDANT tested before CONCENTRATED (the RULE tests CONCENTRATED first)
\providecommand{\rhchXgbBarJointPct}{}\renewcommand{\rhchXgbBarJointPct}{$-$14.6}  % xgb all: full vs HAR, % QLIKE change
\providecommand{\rhchXgbBarJointDm}{}\renewcommand{\rhchXgbBarJointDm}{$-$9.94}  % xgb all: full vs HAR, DM
\providecommand{\rhchXgbBarJointP}{}\renewcommand{\rhchXgbBarJointP}{\ensuremath{<}0.001}  % xgb all: full vs HAR, DM p
\providecommand{\rhchXgbBarAddRmPct}{}\renewcommand{\rhchXgbBarAddRmPct}{$-$3.4}  % xgb all add_rm vs xgb|har: %
\providecommand{\rhchXgbBarAddRmP}{}\renewcommand{\rhchXgbBarAddRmP}{\ensuremath{<}0.001}  % xgb all add_rm: DM p
\providecommand{\rhchXgbBarLoRmPct}{}\renewcommand{\rhchXgbBarLoRmPct}{1.6}  % xgb all lo_rm vs xgb|full_pre: % [pre-patch]
\providecommand{\rhchXgbBarLoRmP}{}\renewcommand{\rhchXgbBarLoRmP}{\ensuremath{<}0.001}  % xgb all lo_rm: DM p [pre-patch]
\providecommand{\rhchXgbBarAddLiqPct}{}\renewcommand{\rhchXgbBarAddLiqPct}{$-$3.5}  % xgb all add_liq vs xgb|har: %
\providecommand{\rhchXgbBarAddLiqP}{}\renewcommand{\rhchXgbBarAddLiqP}{\ensuremath{<}0.001}  % xgb all add_liq: DM p
\providecommand{\rhchXgbBarLoLiqPct}{}\renewcommand{\rhchXgbBarLoLiqPct}{0.2}  % xgb all lo_liq vs xgb|full_pre: % [pre-patch]
\providecommand{\rhchXgbBarLoLiqP}{}\renewcommand{\rhchXgbBarLoLiqP}{0.267}  % xgb all lo_liq: DM p [pre-patch]
\providecommand{\rhchXgbBarAddEwPct}{}\renewcommand{\rhchXgbBarAddEwPct}{$-$3.1}  % xgb all add_ew vs xgb|har: %
\providecommand{\rhchXgbBarAddEwP}{}\renewcommand{\rhchXgbBarAddEwP}{\ensuremath{<}0.001}  % xgb all add_ew: DM p
\providecommand{\rhchXgbBarLoEwPct}{}\renewcommand{\rhchXgbBarLoEwPct}{$-$0.1}  % xgb all lo_ew vs xgb|full_pre: % [pre-patch]
\providecommand{\rhchXgbBarLoEwP}{}\renewcommand{\rhchXgbBarLoEwP}{0.220}  % xgb all lo_ew: DM p [pre-patch]
\providecommand{\rhchXgbBarAddVwPct}{}\renewcommand{\rhchXgbBarAddVwPct}{$-$3.3}  % xgb all add_vw vs xgb|har: %
\providecommand{\rhchXgbBarAddVwP}{}\renewcommand{\rhchXgbBarAddVwP}{\ensuremath{<}0.001}  % xgb all add_vw: DM p
\providecommand{\rhchXgbBarLoVwPct}{}\renewcommand{\rhchXgbBarLoVwPct}{0.1}  % xgb all lo_vw vs xgb|full_pre: % [pre-patch]
\providecommand{\rhchXgbBarLoVwP}{}\renewcommand{\rhchXgbBarLoVwP}{0.349}  % xgb all lo_vw: DM p [pre-patch]
\providecommand{\rhchXgbBarAddSentPct}{}\renewcommand{\rhchXgbBarAddSentPct}{$-$1.7}  % xgb all add_sent vs xgb|har: %
\providecommand{\rhchXgbBarAddSentP}{}\renewcommand{\rhchXgbBarAddSentP}{\ensuremath{<}0.001}  % xgb all add_sent: DM p
\providecommand{\rhchXgbBarLoSentPct}{}\renewcommand{\rhchXgbBarLoSentPct}{0.2}  % xgb all lo_sent vs xgb|full_pre: % [pre-patch]
\providecommand{\rhchXgbBarLoSentP}{}\renewcommand{\rhchXgbBarLoSentP}{0.066}  % xgb all lo_sent: DM p [pre-patch]
\providecommand{\rhchXgbBarAddIvPct}{}\renewcommand{\rhchXgbBarAddIvPct}{$-$1.6}  % xgb all add_iv vs xgb|har: %
\providecommand{\rhchXgbBarAddIvP}{}\renewcommand{\rhchXgbBarAddIvP}{\ensuremath{<}0.001}  % xgb all add_iv: DM p
\providecommand{\rhchXgbBarLoIvPct}{}\renewcommand{\rhchXgbBarLoIvPct}{0.1}  % xgb all lo_iv vs xgb|full_pre: % [pre-patch]
\providecommand{\rhchXgbBarLoIvP}{}\renewcommand{\rhchXgbBarLoIvP}{0.414}  % xgb all lo_iv: DM p [pre-patch]
\providecommand{\rhchXgbBarAddVdPct}{}\renewcommand{\rhchXgbBarAddVdPct}{$-$1.2}  % xgb all add_vd vs xgb|har: %
\providecommand{\rhchXgbBarAddVdP}{}\renewcommand{\rhchXgbBarAddVdP}{\ensuremath{<}0.001}  % xgb all add_vd: DM p
\providecommand{\rhchXgbBarLoVdPct}{}\renewcommand{\rhchXgbBarLoVdPct}{0.1}  % xgb all lo_vd vs xgb|full_pre: % [pre-patch]
\providecommand{\rhchXgbBarLoVdP}{}\renewcommand{\rhchXgbBarLoVdP}{0.674}  % xgb all lo_vd: DM p [pre-patch]
\providecommand{\rhchXgbBarAddFomcPct}{}\renewcommand{\rhchXgbBarAddFomcPct}{$-$9.5}  % xgb all add_fomc vs xgb|har: %
\providecommand{\rhchXgbBarAddFomcP}{}\renewcommand{\rhchXgbBarAddFomcP}{\ensuremath{<}0.001}  % xgb all add_fomc: DM p
\providecommand{\rhchXgbBarLoFomcPct}{}\renewcommand{\rhchXgbBarLoFomcPct}{10.3}  % xgb all lo_fomc vs xgb|full: %
\providecommand{\rhchXgbBarLoFomcP}{}\renewcommand{\rhchXgbBarLoFomcP}{\ensuremath{<}0.001}  % xgb all lo_fomc: DM p
\providecommand{\rhchXgbEventJointPct}{}\renewcommand{\rhchXgbEventJointPct}{$-$63.9}  % xgb event: full vs HAR, % QLIKE change
\providecommand{\rhchXgbEventJointDm}{}\renewcommand{\rhchXgbEventJointDm}{$-$7.75}  % xgb event: full vs HAR, DM
\providecommand{\rhchXgbEventJointP}{}\renewcommand{\rhchXgbEventJointP}{\ensuremath{<}0.001}  % xgb event: full vs HAR, DM p
\providecommand{\rhchXgbEventAddRmPct}{}\renewcommand{\rhchXgbEventAddRmPct}{5.5}  % xgb event add_rm vs xgb|har: %
\providecommand{\rhchXgbEventAddRmP}{}\renewcommand{\rhchXgbEventAddRmP}{0.014}  % xgb event add_rm: DM p
\providecommand{\rhchXgbEventLoRmPct}{}\renewcommand{\rhchXgbEventLoRmPct}{$-$1.3}  % xgb event lo_rm vs xgb|full_pre: % [pre-patch]
\providecommand{\rhchXgbEventLoRmP}{}\renewcommand{\rhchXgbEventLoRmP}{0.170}  % xgb event lo_rm: DM p [pre-patch]
\providecommand{\rhchXgbEventAddLiqPct}{}\renewcommand{\rhchXgbEventAddLiqPct}{0.4}  % xgb event add_liq vs xgb|har: %
\providecommand{\rhchXgbEventAddLiqP}{}\renewcommand{\rhchXgbEventAddLiqP}{0.771}  % xgb event add_liq: DM p
\providecommand{\rhchXgbEventLoLiqPct}{}\renewcommand{\rhchXgbEventLoLiqPct}{$-$0.9}  % xgb event lo_liq vs xgb|full_pre: % [pre-patch]
\providecommand{\rhchXgbEventLoLiqP}{}\renewcommand{\rhchXgbEventLoLiqP}{0.069}  % xgb event lo_liq: DM p [pre-patch]
\providecommand{\rhchXgbEventAddEwPct}{}\renewcommand{\rhchXgbEventAddEwPct}{$-$3.6}  % xgb event add_ew vs xgb|har: %
\providecommand{\rhchXgbEventAddEwP}{}\renewcommand{\rhchXgbEventAddEwP}{0.044}  % xgb event add_ew: DM p
\providecommand{\rhchXgbEventLoEwPct}{}\renewcommand{\rhchXgbEventLoEwPct}{0.0}  % xgb event lo_ew vs xgb|full_pre: % [pre-patch]
\providecommand{\rhchXgbEventLoEwP}{}\renewcommand{\rhchXgbEventLoEwP}{0.990}  % xgb event lo_ew: DM p [pre-patch]
\providecommand{\rhchXgbEventAddVwPct}{}\renewcommand{\rhchXgbEventAddVwPct}{$-$3.3}  % xgb event add_vw vs xgb|har: %
\providecommand{\rhchXgbEventAddVwP}{}\renewcommand{\rhchXgbEventAddVwP}{0.063}  % xgb event add_vw: DM p
\providecommand{\rhchXgbEventLoVwPct}{}\renewcommand{\rhchXgbEventLoVwPct}{$-$0.7}  % xgb event lo_vw vs xgb|full_pre: % [pre-patch]
\providecommand{\rhchXgbEventLoVwP}{}\renewcommand{\rhchXgbEventLoVwP}{0.304}  % xgb event lo_vw: DM p [pre-patch]
\providecommand{\rhchXgbEventAddSentPct}{}\renewcommand{\rhchXgbEventAddSentPct}{0.3}  % xgb event add_sent vs xgb|har: %
\providecommand{\rhchXgbEventAddSentP}{}\renewcommand{\rhchXgbEventAddSentP}{0.765}  % xgb event add_sent: DM p
\providecommand{\rhchXgbEventLoSentPct}{}\renewcommand{\rhchXgbEventLoSentPct}{0.0}  % xgb event lo_sent vs xgb|full_pre: % [pre-patch]
\providecommand{\rhchXgbEventLoSentP}{}\renewcommand{\rhchXgbEventLoSentP}{0.994}  % xgb event lo_sent: DM p [pre-patch]
\providecommand{\rhchXgbEventAddIvPct}{}\renewcommand{\rhchXgbEventAddIvPct}{$-$0.7}  % xgb event add_iv vs xgb|har: %
\providecommand{\rhchXgbEventAddIvP}{}\renewcommand{\rhchXgbEventAddIvP}{0.371}  % xgb event add_iv: DM p
\providecommand{\rhchXgbEventLoIvPct}{}\renewcommand{\rhchXgbEventLoIvPct}{0.4}  % xgb event lo_iv vs xgb|full_pre: % [pre-patch]
\providecommand{\rhchXgbEventLoIvP}{}\renewcommand{\rhchXgbEventLoIvP}{0.482}  % xgb event lo_iv: DM p [pre-patch]
\providecommand{\rhchXgbEventAddVdPct}{}\renewcommand{\rhchXgbEventAddVdPct}{$-$1.9}  % xgb event add_vd vs xgb|har: %
\providecommand{\rhchXgbEventAddVdP}{}\renewcommand{\rhchXgbEventAddVdP}{0.076}  % xgb event add_vd: DM p
\providecommand{\rhchXgbEventLoVdPct}{}\renewcommand{\rhchXgbEventLoVdPct}{$-$0.7}  % xgb event lo_vd vs xgb|full_pre: % [pre-patch]
\providecommand{\rhchXgbEventLoVdP}{}\renewcommand{\rhchXgbEventLoVdP}{0.148}  % xgb event lo_vd: DM p [pre-patch]
\providecommand{\rhchXgbEventAddFomcPct}{}\renewcommand{\rhchXgbEventAddFomcPct}{$-$64.0}  % xgb event add_fomc vs xgb|har: %
\providecommand{\rhchXgbEventAddFomcP}{}\renewcommand{\rhchXgbEventAddFomcP}{\ensuremath{<}0.001}  % xgb event add_fomc: DM p
\providecommand{\rhchXgbEventLoFomcPct}{}\renewcommand{\rhchXgbEventLoFomcPct}{173.0}  % xgb event lo_fomc vs xgb|full: %
\providecommand{\rhchXgbEventLoFomcP}{}\renewcommand{\rhchXgbEventLoFomcP}{\ensuremath{<}0.001}  % xgb event lo_fomc: DM p
\providecommand{\rhchXgbEventReading}{}\renewcommand{\rhchXgbEventReading}{concentrated}  % xgb event: item 3 reading
\providecommand{\rhchXgbEventTopAddShare}{}\renewcommand{\rhchXgbEventTopAddShare}{100}  % xgb event: largest add-alone gain, % of the joint gain (fomc)
\providecommand{\rhchXgbEventSumAddRatio}{}\renewcommand{\rhchXgbEventSumAddRatio}{1.05}  % xgb event: add-alone gains summed over groups / joint gain
\providecommand{\rhchXgbEventNLoSig}{}\renewcommand{\rhchXgbEventNLoSig}{1}  % xgb event: significant leave-one-outs
\providecommand{\rhchXgbEventNLoRun}{}\renewcommand{\rhchXgbEventNLoRun}{8}  % xgb event: leave-one-outs available (of 8)
\providecommand{\rhchXgbEventConcTest}{}\renewcommand{\rhchXgbEventConcTest}{yes}  % xgb event: the CONC definition on its own (pending: undetermined)
\providecommand{\rhchXgbEventRedTest}{}\renewcommand{\rhchXgbEventRedTest}{no}  % xgb event: the RED definition on its own (pending: undetermined)
\providecommand{\rhchXgbEventReadingRedFirst}{}\renewcommand{\rhchXgbEventReadingRedFirst}{concentrated}  % xgb event: reading with REDUNDANT tested before CONCENTRATED (the RULE tests CONCENTRATED first)
\providecommand{\rhchXgbOrdJointPct}{}\renewcommand{\rhchXgbOrdJointPct}{$-$6.1}  % xgb ord: full vs HAR, % QLIKE change
\providecommand{\rhchXgbOrdJointDm}{}\renewcommand{\rhchXgbOrdJointDm}{$-$10.42}  % xgb ord: full vs HAR, DM
\providecommand{\rhchXgbOrdJointP}{}\renewcommand{\rhchXgbOrdJointP}{\ensuremath{<}0.001}  % xgb ord: full vs HAR, DM p
\providecommand{\rhchXgbOrdAddRmPct}{}\renewcommand{\rhchXgbOrdAddRmPct}{$-$4.9}  % xgb ord add_rm vs xgb|har: %
\providecommand{\rhchXgbOrdAddRmP}{}\renewcommand{\rhchXgbOrdAddRmP}{\ensuremath{<}0.001}  % xgb ord add_rm: DM p
\providecommand{\rhchXgbOrdLoRmPct}{}\renewcommand{\rhchXgbOrdLoRmPct}{1.8}  % xgb ord lo_rm vs xgb|full_pre: % [pre-patch]
\providecommand{\rhchXgbOrdLoRmP}{}\renewcommand{\rhchXgbOrdLoRmP}{\ensuremath{<}0.001}  % xgb ord lo_rm: DM p [pre-patch]
\providecommand{\rhchXgbOrdAddLiqPct}{}\renewcommand{\rhchXgbOrdAddLiqPct}{$-$4.1}  % xgb ord add_liq vs xgb|har: %
\providecommand{\rhchXgbOrdAddLiqP}{}\renewcommand{\rhchXgbOrdAddLiqP}{\ensuremath{<}0.001}  % xgb ord add_liq: DM p
\providecommand{\rhchXgbOrdLoLiqPct}{}\renewcommand{\rhchXgbOrdLoLiqPct}{0.2}  % xgb ord lo_liq vs xgb|full_pre: % [pre-patch]
\providecommand{\rhchXgbOrdLoLiqP}{}\renewcommand{\rhchXgbOrdLoLiqP}{0.095}  % xgb ord lo_liq: DM p [pre-patch]
\providecommand{\rhchXgbOrdAddEwPct}{}\renewcommand{\rhchXgbOrdAddEwPct}{$-$3.0}  % xgb ord add_ew vs xgb|har: %
\providecommand{\rhchXgbOrdAddEwP}{}\renewcommand{\rhchXgbOrdAddEwP}{\ensuremath{<}0.001}  % xgb ord add_ew: DM p
\providecommand{\rhchXgbOrdLoEwPct}{}\renewcommand{\rhchXgbOrdLoEwPct}{$-$0.1}  % xgb ord lo_ew vs xgb|full_pre: % [pre-patch]
\providecommand{\rhchXgbOrdLoEwP}{}\renewcommand{\rhchXgbOrdLoEwP}{0.172}  % xgb ord lo_ew: DM p [pre-patch]
\providecommand{\rhchXgbOrdAddVwPct}{}\renewcommand{\rhchXgbOrdAddVwPct}{$-$3.3}  % xgb ord add_vw vs xgb|har: %
\providecommand{\rhchXgbOrdAddVwP}{}\renewcommand{\rhchXgbOrdAddVwP}{\ensuremath{<}0.001}  % xgb ord add_vw: DM p
\providecommand{\rhchXgbOrdLoVwPct}{}\renewcommand{\rhchXgbOrdLoVwPct}{0.2}  % xgb ord lo_vw vs xgb|full_pre: % [pre-patch]
\providecommand{\rhchXgbOrdLoVwP}{}\renewcommand{\rhchXgbOrdLoVwP}{0.133}  % xgb ord lo_vw: DM p [pre-patch]
\providecommand{\rhchXgbOrdAddSentPct}{}\renewcommand{\rhchXgbOrdAddSentPct}{$-$2.0}  % xgb ord add_sent vs xgb|har: %
\providecommand{\rhchXgbOrdAddSentP}{}\renewcommand{\rhchXgbOrdAddSentP}{\ensuremath{<}0.001}  % xgb ord add_sent: DM p
\providecommand{\rhchXgbOrdLoSentPct}{}\renewcommand{\rhchXgbOrdLoSentPct}{0.2}  % xgb ord lo_sent vs xgb|full_pre: % [pre-patch]
\providecommand{\rhchXgbOrdLoSentP}{}\renewcommand{\rhchXgbOrdLoSentP}{0.051}  % xgb ord lo_sent: DM p [pre-patch]
\providecommand{\rhchXgbOrdAddIvPct}{}\renewcommand{\rhchXgbOrdAddIvPct}{$-$1.7}  % xgb ord add_iv vs xgb|har: %
\providecommand{\rhchXgbOrdAddIvP}{}\renewcommand{\rhchXgbOrdAddIvP}{\ensuremath{<}0.001}  % xgb ord add_iv: DM p
\providecommand{\rhchXgbOrdLoIvPct}{}\renewcommand{\rhchXgbOrdLoIvPct}{0.0}  % xgb ord lo_iv vs xgb|full_pre: % [pre-patch]
\providecommand{\rhchXgbOrdLoIvP}{}\renewcommand{\rhchXgbOrdLoIvP}{0.611}  % xgb ord lo_iv: DM p [pre-patch]
\providecommand{\rhchXgbOrdAddVdPct}{}\renewcommand{\rhchXgbOrdAddVdPct}{$-$1.1}  % xgb ord add_vd vs xgb|har: %
\providecommand{\rhchXgbOrdAddVdP}{}\renewcommand{\rhchXgbOrdAddVdP}{\ensuremath{<}0.001}  % xgb ord add_vd: DM p
\providecommand{\rhchXgbOrdLoVdPct}{}\renewcommand{\rhchXgbOrdLoVdPct}{0.1}  % xgb ord lo_vd vs xgb|full_pre: % [pre-patch]
\providecommand{\rhchXgbOrdLoVdP}{}\renewcommand{\rhchXgbOrdLoVdP}{0.393}  % xgb ord lo_vd: DM p [pre-patch]
\providecommand{\rhchXgbOrdAddFomcPct}{}\renewcommand{\rhchXgbOrdAddFomcPct}{$-$0.1}  % xgb ord add_fomc vs xgb|har: %
\providecommand{\rhchXgbOrdAddFomcP}{}\renewcommand{\rhchXgbOrdAddFomcP}{0.726}  % xgb ord add_fomc: DM p
\providecommand{\rhchXgbOrdLoFomcPct}{}\renewcommand{\rhchXgbOrdLoFomcPct}{$-$0.4}  % xgb ord lo_fomc vs xgb|full: %
\providecommand{\rhchXgbOrdLoFomcP}{}\renewcommand{\rhchXgbOrdLoFomcP}{0.013}  % xgb ord lo_fomc: DM p
\providecommand{\rhchXgbOrdReading}{}\renewcommand{\rhchXgbOrdReading}{concentrated}  % xgb ord: item 3 reading
\providecommand{\rhchXgbOrdTopAddShare}{}\renewcommand{\rhchXgbOrdTopAddShare}{81}  % xgb ord: largest add-alone gain, % of the joint gain (rm)
\providecommand{\rhchXgbOrdSumAddRatio}{}\renewcommand{\rhchXgbOrdSumAddRatio}{3.32}  % xgb ord: add-alone gains summed over groups / joint gain
\providecommand{\rhchXgbOrdNLoSig}{}\renewcommand{\rhchXgbOrdNLoSig}{1}  % xgb ord: significant leave-one-outs
\providecommand{\rhchXgbOrdNLoRun}{}\renewcommand{\rhchXgbOrdNLoRun}{8}  % xgb ord: leave-one-outs available (of 8)
\providecommand{\rhchXgbOrdConcTest}{}\renewcommand{\rhchXgbOrdConcTest}{yes}  % xgb ord: the CONC definition on its own (pending: undetermined)
\providecommand{\rhchXgbOrdRedTest}{}\renewcommand{\rhchXgbOrdRedTest}{yes}  % xgb ord: the RED definition on its own (pending: undetermined)
\providecommand{\rhchXgbOrdReadingRedFirst}{}\renewcommand{\rhchXgbOrdReadingRedFirst}{redundant}  % xgb ord: reading with REDUNDANT tested before CONCENTRATED (the RULE tests CONCENTRATED first)
\providecommand{\rhchRfBarJointPct}{}\renewcommand{\rhchRfBarJointPct}{$-$11.3}  % rf all: full vs HAR, % QLIKE change
\providecommand{\rhchRfBarJointDm}{}\renewcommand{\rhchRfBarJointDm}{$-$9.74}  % rf all: full vs HAR, DM
\providecommand{\rhchRfBarJointP}{}\renewcommand{\rhchRfBarJointP}{\ensuremath{<}0.001}  % rf all: full vs HAR, DM p
\providecommand{\rhchRfBarAddRmPct}{}\renewcommand{\rhchRfBarAddRmPct}{$-$3.2}  % rf all add_rm vs rf|har: %
\providecommand{\rhchRfBarAddRmP}{}\renewcommand{\rhchRfBarAddRmP}{\ensuremath{<}0.001}  % rf all add_rm: DM p
\providecommand{\rhchRfBarLoRmPct}{}\renewcommand{\rhchRfBarLoRmPct}{1.5}  % rf all lo_rm vs rf|full_pre: % [pre-patch]
\providecommand{\rhchRfBarLoRmP}{}\renewcommand{\rhchRfBarLoRmP}{\ensuremath{<}0.001}  % rf all lo_rm: DM p [pre-patch]
\providecommand{\rhchRfBarAddLiqPct}{}\renewcommand{\rhchRfBarAddLiqPct}{$-$2.4}  % rf all add_liq vs rf|har: %
\providecommand{\rhchRfBarAddLiqP}{}\renewcommand{\rhchRfBarAddLiqP}{\ensuremath{<}0.001}  % rf all add_liq: DM p
\providecommand{\rhchRfBarAddEwPct}{}\renewcommand{\rhchRfBarAddEwPct}{$-$2.3}  % rf all add_ew vs rf|har: %
\providecommand{\rhchRfBarAddEwP}{}\renewcommand{\rhchRfBarAddEwP}{\ensuremath{<}0.001}  % rf all add_ew: DM p
\providecommand{\rhchRfBarAddVwPct}{}\renewcommand{\rhchRfBarAddVwPct}{$-$2.5}  % rf all add_vw vs rf|har: %
\providecommand{\rhchRfBarAddVwP}{}\renewcommand{\rhchRfBarAddVwP}{\ensuremath{<}0.001}  % rf all add_vw: DM p
\providecommand{\rhchRfBarAddSentPct}{}\renewcommand{\rhchRfBarAddSentPct}{$-$0.9}  % rf all add_sent vs rf|har: %
\providecommand{\rhchRfBarAddSentP}{}\renewcommand{\rhchRfBarAddSentP}{\ensuremath{<}0.001}  % rf all add_sent: DM p
\providecommand{\rhchRfBarAddIvPct}{}\renewcommand{\rhchRfBarAddIvPct}{$-$0.2}  % rf all add_iv vs rf|har: %
\providecommand{\rhchRfBarAddIvP}{}\renewcommand{\rhchRfBarAddIvP}{0.601}  % rf all add_iv: DM p
\providecommand{\rhchRfBarAddVdPct}{}\renewcommand{\rhchRfBarAddVdPct}{$-$0.6}  % rf all add_vd vs rf|har: %
\providecommand{\rhchRfBarAddVdP}{}\renewcommand{\rhchRfBarAddVdP}{\ensuremath{<}0.001}  % rf all add_vd: DM p
\providecommand{\rhchRfBarAddFomcPct}{}\renewcommand{\rhchRfBarAddFomcPct}{$-$7.6}  % rf all add_fomc vs rf|har: %
\providecommand{\rhchRfBarAddFomcP}{}\renewcommand{\rhchRfBarAddFomcP}{\ensuremath{<}0.001}  % rf all add_fomc: DM p
\providecommand{\rhchRfEventJointPct}{}\renewcommand{\rhchRfEventJointPct}{$-$44.8}  % rf event: full vs HAR, % QLIKE change
\providecommand{\rhchRfEventJointDm}{}\renewcommand{\rhchRfEventJointDm}{$-$8.56}  % rf event: full vs HAR, DM
\providecommand{\rhchRfEventJointP}{}\renewcommand{\rhchRfEventJointP}{\ensuremath{<}0.001}  % rf event: full vs HAR, DM p
\providecommand{\rhchRfEventAddRmPct}{}\renewcommand{\rhchRfEventAddRmPct}{5.3}  % rf event add_rm vs rf|har: %
\providecommand{\rhchRfEventAddRmP}{}\renewcommand{\rhchRfEventAddRmP}{0.061}  % rf event add_rm: DM p
\providecommand{\rhchRfEventLoRmPct}{}\renewcommand{\rhchRfEventLoRmPct}{$-$2.1}  % rf event lo_rm vs rf|full_pre: % [pre-patch]
\providecommand{\rhchRfEventLoRmP}{}\renewcommand{\rhchRfEventLoRmP}{0.036}  % rf event lo_rm: DM p [pre-patch]
\providecommand{\rhchRfEventAddLiqPct}{}\renewcommand{\rhchRfEventAddLiqPct}{2.0}  % rf event add_liq vs rf|har: %
\providecommand{\rhchRfEventAddLiqP}{}\renewcommand{\rhchRfEventAddLiqP}{0.184}  % rf event add_liq: DM p
\providecommand{\rhchRfEventAddEwPct}{}\renewcommand{\rhchRfEventAddEwPct}{$-$2.4}  % rf event add_ew vs rf|har: %
\providecommand{\rhchRfEventAddEwP}{}\renewcommand{\rhchRfEventAddEwP}{0.073}  % rf event add_ew: DM p
\providecommand{\rhchRfEventAddVwPct}{}\renewcommand{\rhchRfEventAddVwPct}{$-$0.9}  % rf event add_vw vs rf|har: %
\providecommand{\rhchRfEventAddVwP}{}\renewcommand{\rhchRfEventAddVwP}{0.651}  % rf event add_vw: DM p
\providecommand{\rhchRfEventAddSentPct}{}\renewcommand{\rhchRfEventAddSentPct}{0.4}  % rf event add_sent vs rf|har: %
\providecommand{\rhchRfEventAddSentP}{}\renewcommand{\rhchRfEventAddSentP}{0.621}  % rf event add_sent: DM p
\providecommand{\rhchRfEventAddIvPct}{}\renewcommand{\rhchRfEventAddIvPct}{4.1}  % rf event add_iv vs rf|har: %
\providecommand{\rhchRfEventAddIvP}{}\renewcommand{\rhchRfEventAddIvP}{0.001}  % rf event add_iv: DM p
\providecommand{\rhchRfEventAddVdPct}{}\renewcommand{\rhchRfEventAddVdPct}{1.3}  % rf event add_vd vs rf|har: %
\providecommand{\rhchRfEventAddVdP}{}\renewcommand{\rhchRfEventAddVdP}{0.046}  % rf event add_vd: DM p
\providecommand{\rhchRfEventAddFomcPct}{}\renewcommand{\rhchRfEventAddFomcPct}{$-$48.5}  % rf event add_fomc vs rf|har: %
\providecommand{\rhchRfEventAddFomcP}{}\renewcommand{\rhchRfEventAddFomcP}{\ensuremath{<}0.001}  % rf event add_fomc: DM p
\providecommand{\rhchRfEventReading}{}\renewcommand{\rhchRfEventReading}{n/a}  % rf event: item 3 reading
\providecommand{\rhchRfEventTopAddShare}{}\renewcommand{\rhchRfEventTopAddShare}{108}  % rf event: largest add-alone gain, % of the joint gain (fomc)
\providecommand{\rhchRfEventSumAddRatio}{}\renewcommand{\rhchRfEventSumAddRatio}{0.86}  % rf event: add-alone gains summed over groups / joint gain
\providecommand{\rhchRfEventNLoSig}{}\renewcommand{\rhchRfEventNLoSig}{0}  % rf event: significant leave-one-outs
\providecommand{\rhchRfEventNLoRun}{}\renewcommand{\rhchRfEventNLoRun}{1}  % rf event: leave-one-outs available (of 8)
\providecommand{\rhchRfEventConcTest}{}\renewcommand{\rhchRfEventConcTest}{\rhPending{}}  % rf event: the CONC definition on its own (pending: undetermined)
\providecommand{\rhchRfEventRedTest}{}\renewcommand{\rhchRfEventRedTest}{no}  % rf event: the RED definition on its own (pending: undetermined)
\providecommand{\rhchRfEventReadingRedFirst}{}\renewcommand{\rhchRfEventReadingRedFirst}{n/a}  % rf event: reading with REDUNDANT tested before CONCENTRATED (the RULE tests CONCENTRATED first)
\providecommand{\rhchRfOrdJointPct}{}\renewcommand{\rhchRfOrdJointPct}{$-$5.3}  % rf ord: full vs HAR, % QLIKE change
\providecommand{\rhchRfOrdJointDm}{}\renewcommand{\rhchRfOrdJointDm}{$-$11.16}  % rf ord: full vs HAR, DM
\providecommand{\rhchRfOrdJointP}{}\renewcommand{\rhchRfOrdJointP}{\ensuremath{<}0.001}  % rf ord: full vs HAR, DM p
\providecommand{\rhchRfOrdAddRmPct}{}\renewcommand{\rhchRfOrdAddRmPct}{$-$4.7}  % rf ord add_rm vs rf|har: %
\providecommand{\rhchRfOrdAddRmP}{}\renewcommand{\rhchRfOrdAddRmP}{\ensuremath{<}0.001}  % rf ord add_rm: DM p
\providecommand{\rhchRfOrdLoRmPct}{}\renewcommand{\rhchRfOrdLoRmPct}{1.8}  % rf ord lo_rm vs rf|full_pre: % [pre-patch]
\providecommand{\rhchRfOrdLoRmP}{}\renewcommand{\rhchRfOrdLoRmP}{\ensuremath{<}0.001}  % rf ord lo_rm: DM p [pre-patch]
\providecommand{\rhchRfOrdAddLiqPct}{}\renewcommand{\rhchRfOrdAddLiqPct}{$-$3.2}  % rf ord add_liq vs rf|har: %
\providecommand{\rhchRfOrdAddLiqP}{}\renewcommand{\rhchRfOrdAddLiqP}{\ensuremath{<}0.001}  % rf ord add_liq: DM p
\providecommand{\rhchRfOrdAddEwPct}{}\renewcommand{\rhchRfOrdAddEwPct}{$-$2.3}  % rf ord add_ew vs rf|har: %
\providecommand{\rhchRfOrdAddEwP}{}\renewcommand{\rhchRfOrdAddEwP}{\ensuremath{<}0.001}  % rf ord add_ew: DM p
\providecommand{\rhchRfOrdAddVwPct}{}\renewcommand{\rhchRfOrdAddVwPct}{$-$2.8}  % rf ord add_vw vs rf|har: %
\providecommand{\rhchRfOrdAddVwP}{}\renewcommand{\rhchRfOrdAddVwP}{\ensuremath{<}0.001}  % rf ord add_vw: DM p
\providecommand{\rhchRfOrdAddSentPct}{}\renewcommand{\rhchRfOrdAddSentPct}{$-$1.1}  % rf ord add_sent vs rf|har: %
\providecommand{\rhchRfOrdAddSentP}{}\renewcommand{\rhchRfOrdAddSentP}{\ensuremath{<}0.001}  % rf ord add_sent: DM p
\providecommand{\rhchRfOrdAddIvPct}{}\renewcommand{\rhchRfOrdAddIvPct}{$-$1.0}  % rf ord add_iv vs rf|har: %
\providecommand{\rhchRfOrdAddIvP}{}\renewcommand{\rhchRfOrdAddIvP}{0.004}  % rf ord add_iv: DM p
\providecommand{\rhchRfOrdAddVdPct}{}\renewcommand{\rhchRfOrdAddVdPct}{$-$1.0}  % rf ord add_vd vs rf|har: %
\providecommand{\rhchRfOrdAddVdP}{}\renewcommand{\rhchRfOrdAddVdP}{\ensuremath{<}0.001}  % rf ord add_vd: DM p
\providecommand{\rhchRfOrdAddFomcPct}{}\renewcommand{\rhchRfOrdAddFomcPct}{$-$0.3}  % rf ord add_fomc vs rf|har: %
\providecommand{\rhchRfOrdAddFomcP}{}\renewcommand{\rhchRfOrdAddFomcP}{0.141}  % rf ord add_fomc: DM p
\providecommand{\rhchRfOrdReading}{}\renewcommand{\rhchRfOrdReading}{concentrated}  % rf ord: item 3 reading
\providecommand{\rhchRfOrdTopAddShare}{}\renewcommand{\rhchRfOrdTopAddShare}{88}  % rf ord: largest add-alone gain, % of the joint gain (rm)
\providecommand{\rhchRfOrdSumAddRatio}{}\renewcommand{\rhchRfOrdSumAddRatio}{3.07}  % rf ord: add-alone gains summed over groups / joint gain
\providecommand{\rhchRfOrdNLoSig}{}\renewcommand{\rhchRfOrdNLoSig}{1}  % rf ord: significant leave-one-outs
\providecommand{\rhchRfOrdNLoRun}{}\renewcommand{\rhchRfOrdNLoRun}{1}  % rf ord: leave-one-outs available (of 8)
\providecommand{\rhchRfOrdConcTest}{}\renewcommand{\rhchRfOrdConcTest}{yes}  % rf ord: the CONC definition on its own (pending: undetermined)
\providecommand{\rhchRfOrdRedTest}{}\renewcommand{\rhchRfOrdRedTest}{\rhPending{}}  % rf ord: the RED definition on its own (pending: undetermined)
\providecommand{\rhchRfOrdReadingRedFirst}{}\renewcommand{\rhchRfOrdReadingRedFirst}{n/a}  % rf ord: reading with REDUNDANT tested before CONCENTRATED (the RULE tests CONCENTRATED first)
\providecommand{\rhchLinEventReading}{}\renewcommand{\rhchLinEventReading}{concentrated}  % lin event: item 3 reading by model type
\providecommand{\rhchLinEventReadingRedFirst}{}\renewcommand{\rhchLinEventReadingRedFirst}{concentrated}  % lin event: by model type, REDUNDANT tested first
\providecommand{\rhchLinEventNConcTest}{}\renewcommand{\rhchLinEventNConcTest}{3}  % lin event: classes meeting the Conc definition (of 3)
\providecommand{\rhchLinEventNRedTest}{}\renewcommand{\rhchLinEventNRedTest}{0}  % lin event: classes meeting the Red definition (of 3)
\providecommand{\rhchTreeEventReading}{}\renewcommand{\rhchTreeEventReading}{concentrated}  % tree event: item 3 reading by model type
\providecommand{\rhchTreeEventReadingRedFirst}{}\renewcommand{\rhchTreeEventReadingRedFirst}{concentrated}  % tree event: by model type, REDUNDANT tested first
\providecommand{\rhchTreeEventNConcTest}{}\renewcommand{\rhchTreeEventNConcTest}{2}  % tree event: classes meeting the Conc definition (of 3)
\providecommand{\rhchTreeEventNRedTest}{}\renewcommand{\rhchTreeEventNRedTest}{0}  % tree event: classes meeting the Red definition (of 3)
\providecommand{\rhchLinOrdReading}{}\renewcommand{\rhchLinOrdReading}{concentrated}  % lin ord: item 3 reading by model type
\providecommand{\rhchLinOrdReadingRedFirst}{}\renewcommand{\rhchLinOrdReadingRedFirst}{redundant}  % lin ord: by model type, REDUNDANT tested first
\providecommand{\rhchLinOrdNConcTest}{}\renewcommand{\rhchLinOrdNConcTest}{3}  % lin ord: classes meeting the Conc definition (of 3)
\providecommand{\rhchLinOrdNRedTest}{}\renewcommand{\rhchLinOrdNRedTest}{3}  % lin ord: classes meeting the Red definition (of 3)
\providecommand{\rhchTreeOrdReading}{}\renewcommand{\rhchTreeOrdReading}{concentrated}  % tree ord: item 3 reading by model type
\providecommand{\rhchTreeOrdReadingRedFirst}{}\renewcommand{\rhchTreeOrdReadingRedFirst}{split}  % tree ord: by model type, REDUNDANT tested first
\providecommand{\rhchTreeOrdNConcTest}{}\renewcommand{\rhchTreeOrdNConcTest}{3}  % tree ord: classes meeting the Conc definition (of 3)
\providecommand{\rhchTreeOrdNRedTest}{}\renewcommand{\rhchTreeOrdNRedTest}{1}  % tree ord: classes meeting the Red definition (of 3)
\providecommand{\rhchXRidgeBarFullVsRmPct}{}\renewcommand{\rhchXRidgeBarFullVsRmPct}{$-$3.5}  % ADDITIONAL, NOT PRE-REGISTERED: ridge all: full vs HAR+RM: % QLIKE change
\providecommand{\rhchXRidgeBarFullVsRmDm}{}\renewcommand{\rhchXRidgeBarFullVsRmDm}{$-$4.52}  % ADDITIONAL, NOT PRE-REGISTERED: ridge all: full vs HAR+RM: DM
\providecommand{\rhchXRidgeBarFullVsRmP}{}\renewcommand{\rhchXRidgeBarFullVsRmP}{\ensuremath{<}0.001}  % ADDITIONAL, NOT PRE-REGISTERED: ridge all: full vs HAR+RM: DM p
\providecommand{\rhchXRidgeBarFullVsRmShare}{}\renewcommand{\rhchXRidgeBarFullVsRmShare}{56}  % ADDITIONAL, NOT PRE-REGISTERED: ridge all: full vs HAR+RM: % of the joint gain
\providecommand{\rhchXRidgeBarRmVsHarPct}{}\renewcommand{\rhchXRidgeBarRmVsHarPct}{$-$2.6}  % ADDITIONAL, NOT PRE-REGISTERED: ridge all: HAR+RM vs HAR: % QLIKE change
\providecommand{\rhchXRidgeBarRmVsHarDm}{}\renewcommand{\rhchXRidgeBarRmVsHarDm}{$-$4.14}  % ADDITIONAL, NOT PRE-REGISTERED: ridge all: HAR+RM vs HAR: DM
\providecommand{\rhchXRidgeBarRmVsHarP}{}\renewcommand{\rhchXRidgeBarRmVsHarP}{\ensuremath{<}0.001}  % ADDITIONAL, NOT PRE-REGISTERED: ridge all: HAR+RM vs HAR: DM p
\providecommand{\rhchXRidgeBarRmVsHarShare}{}\renewcommand{\rhchXRidgeBarRmVsHarShare}{44}  % ADDITIONAL, NOT PRE-REGISTERED: ridge all: HAR+RM vs HAR: % of the joint gain
\providecommand{\rhchXRidgeEventFullVsRmPct}{}\renewcommand{\rhchXRidgeEventFullVsRmPct}{$-$25.7}  % ADDITIONAL, NOT PRE-REGISTERED: ridge event: full vs HAR+RM: % QLIKE change
\providecommand{\rhchXRidgeEventFullVsRmDm}{}\renewcommand{\rhchXRidgeEventFullVsRmDm}{$-$5.37}  % ADDITIONAL, NOT PRE-REGISTERED: ridge event: full vs HAR+RM: DM
\providecommand{\rhchXRidgeEventFullVsRmP}{}\renewcommand{\rhchXRidgeEventFullVsRmP}{\ensuremath{<}0.001}  % ADDITIONAL, NOT PRE-REGISTERED: ridge event: full vs HAR+RM: DM p
\providecommand{\rhchXRidgeEventFullVsRmShare}{}\renewcommand{\rhchXRidgeEventFullVsRmShare}{120}  % ADDITIONAL, NOT PRE-REGISTERED: ridge event: full vs HAR+RM: % of the joint gain
\providecommand{\rhchXRidgeEventRmVsHarPct}{}\renewcommand{\rhchXRidgeEventRmVsHarPct}{4.4}  % ADDITIONAL, NOT PRE-REGISTERED: ridge event: HAR+RM vs HAR: % QLIKE change
\providecommand{\rhchXRidgeEventRmVsHarDm}{}\renewcommand{\rhchXRidgeEventRmVsHarDm}{2.28}  % ADDITIONAL, NOT PRE-REGISTERED: ridge event: HAR+RM vs HAR: DM
\providecommand{\rhchXRidgeEventRmVsHarP}{}\renewcommand{\rhchXRidgeEventRmVsHarP}{0.023}  % ADDITIONAL, NOT PRE-REGISTERED: ridge event: HAR+RM vs HAR: DM p
\providecommand{\rhchXRidgeEventRmVsHarShare}{}\renewcommand{\rhchXRidgeEventRmVsHarShare}{$-$20}  % ADDITIONAL, NOT PRE-REGISTERED: ridge event: HAR+RM vs HAR: % of the joint gain
\providecommand{\rhchXRidgeOrdFullVsRmPct}{}\renewcommand{\rhchXRidgeOrdFullVsRmPct}{0.6}  % ADDITIONAL, NOT PRE-REGISTERED: ridge ord: full vs HAR+RM: % QLIKE change
\providecommand{\rhchXRidgeOrdFullVsRmDm}{}\renewcommand{\rhchXRidgeOrdFullVsRmDm}{1.24}  % ADDITIONAL, NOT PRE-REGISTERED: ridge ord: full vs HAR+RM: DM
\providecommand{\rhchXRidgeOrdFullVsRmP}{}\renewcommand{\rhchXRidgeOrdFullVsRmP}{0.215}  % ADDITIONAL, NOT PRE-REGISTERED: ridge ord: full vs HAR+RM: DM p
\providecommand{\rhchXRidgeOrdFullVsRmShare}{}\renewcommand{\rhchXRidgeOrdFullVsRmShare}{$-$18}  % ADDITIONAL, NOT PRE-REGISTERED: ridge ord: full vs HAR+RM: % of the joint gain
\providecommand{\rhchXRidgeOrdRmVsHarPct}{}\renewcommand{\rhchXRidgeOrdRmVsHarPct}{$-$3.8}  % ADDITIONAL, NOT PRE-REGISTERED: ridge ord: HAR+RM vs HAR: % QLIKE change
\providecommand{\rhchXRidgeOrdRmVsHarDm}{}\renewcommand{\rhchXRidgeOrdRmVsHarDm}{$-$5.48}  % ADDITIONAL, NOT PRE-REGISTERED: ridge ord: HAR+RM vs HAR: DM
\providecommand{\rhchXRidgeOrdRmVsHarP}{}\renewcommand{\rhchXRidgeOrdRmVsHarP}{\ensuremath{<}0.001}  % ADDITIONAL, NOT PRE-REGISTERED: ridge ord: HAR+RM vs HAR: DM p
\providecommand{\rhchXRidgeOrdRmVsHarShare}{}\renewcommand{\rhchXRidgeOrdRmVsHarShare}{118}  % ADDITIONAL, NOT PRE-REGISTERED: ridge ord: HAR+RM vs HAR: % of the joint gain
\providecommand{\rhchXLassoBarFullVsRmPct}{}\renewcommand{\rhchXLassoBarFullVsRmPct}{$-$7.0}  % ADDITIONAL, NOT PRE-REGISTERED: lasso all: full vs HAR+RM: % QLIKE change
\providecommand{\rhchXLassoBarFullVsRmDm}{}\renewcommand{\rhchXLassoBarFullVsRmDm}{$-$5.66}  % ADDITIONAL, NOT PRE-REGISTERED: lasso all: full vs HAR+RM: DM
\providecommand{\rhchXLassoBarFullVsRmP}{}\renewcommand{\rhchXLassoBarFullVsRmP}{\ensuremath{<}0.001}  % ADDITIONAL, NOT PRE-REGISTERED: lasso all: full vs HAR+RM: DM p
\providecommand{\rhchXLassoBarFullVsRmShare}{}\renewcommand{\rhchXLassoBarFullVsRmShare}{76}  % ADDITIONAL, NOT PRE-REGISTERED: lasso all: full vs HAR+RM: % of the joint gain
\providecommand{\rhchXLassoBarRmVsHarPct}{}\renewcommand{\rhchXLassoBarRmVsHarPct}{$-$2.1}  % ADDITIONAL, NOT PRE-REGISTERED: lasso all: HAR+RM vs HAR: % QLIKE change
\providecommand{\rhchXLassoBarRmVsHarDm}{}\renewcommand{\rhchXLassoBarRmVsHarDm}{$-$2.64}  % ADDITIONAL, NOT PRE-REGISTERED: lasso all: HAR+RM vs HAR: DM
\providecommand{\rhchXLassoBarRmVsHarP}{}\renewcommand{\rhchXLassoBarRmVsHarP}{0.008}  % ADDITIONAL, NOT PRE-REGISTERED: lasso all: HAR+RM vs HAR: DM p
\providecommand{\rhchXLassoBarRmVsHarShare}{}\renewcommand{\rhchXLassoBarRmVsHarShare}{24}  % ADDITIONAL, NOT PRE-REGISTERED: lasso all: HAR+RM vs HAR: % of the joint gain
\providecommand{\rhchXLassoEventFullVsRmPct}{}\renewcommand{\rhchXLassoEventFullVsRmPct}{$-$49.5}  % ADDITIONAL, NOT PRE-REGISTERED: lasso event: full vs HAR+RM: % QLIKE change
\providecommand{\rhchXLassoEventFullVsRmDm}{}\renewcommand{\rhchXLassoEventFullVsRmDm}{$-$6.59}  % ADDITIONAL, NOT PRE-REGISTERED: lasso event: full vs HAR+RM: DM
\providecommand{\rhchXLassoEventFullVsRmP}{}\renewcommand{\rhchXLassoEventFullVsRmP}{\ensuremath{<}0.001}  % ADDITIONAL, NOT PRE-REGISTERED: lasso event: full vs HAR+RM: DM p
\providecommand{\rhchXLassoEventFullVsRmShare}{}\renewcommand{\rhchXLassoEventFullVsRmShare}{110}  % ADDITIONAL, NOT PRE-REGISTERED: lasso event: full vs HAR+RM: % of the joint gain
\providecommand{\rhchXLassoEventRmVsHarPct}{}\renewcommand{\rhchXLassoEventRmVsHarPct}{4.7}  % ADDITIONAL, NOT PRE-REGISTERED: lasso event: HAR+RM vs HAR: % QLIKE change
\providecommand{\rhchXLassoEventRmVsHarDm}{}\renewcommand{\rhchXLassoEventRmVsHarDm}{2.25}  % ADDITIONAL, NOT PRE-REGISTERED: lasso event: HAR+RM vs HAR: DM
\providecommand{\rhchXLassoEventRmVsHarP}{}\renewcommand{\rhchXLassoEventRmVsHarP}{0.024}  % ADDITIONAL, NOT PRE-REGISTERED: lasso event: HAR+RM vs HAR: DM p
\providecommand{\rhchXLassoEventRmVsHarShare}{}\renewcommand{\rhchXLassoEventRmVsHarShare}{$-$10}  % ADDITIONAL, NOT PRE-REGISTERED: lasso event: HAR+RM vs HAR: % of the joint gain
\providecommand{\rhchXLassoOrdFullVsRmPct}{}\renewcommand{\rhchXLassoOrdFullVsRmPct}{0.9}  % ADDITIONAL, NOT PRE-REGISTERED: lasso ord: full vs HAR+RM: % QLIKE change
\providecommand{\rhchXLassoOrdFullVsRmDm}{}\renewcommand{\rhchXLassoOrdFullVsRmDm}{1.44}  % ADDITIONAL, NOT PRE-REGISTERED: lasso ord: full vs HAR+RM: DM
\providecommand{\rhchXLassoOrdFullVsRmP}{}\renewcommand{\rhchXLassoOrdFullVsRmP}{0.151}  % ADDITIONAL, NOT PRE-REGISTERED: lasso ord: full vs HAR+RM: DM p
\providecommand{\rhchXLassoOrdFullVsRmShare}{}\renewcommand{\rhchXLassoOrdFullVsRmShare}{$-$36}  % ADDITIONAL, NOT PRE-REGISTERED: lasso ord: full vs HAR+RM: % of the joint gain
\providecommand{\rhchXLassoOrdRmVsHarPct}{}\renewcommand{\rhchXLassoOrdRmVsHarPct}{$-$3.3}  % ADDITIONAL, NOT PRE-REGISTERED: lasso ord: HAR+RM vs HAR: % QLIKE change
\providecommand{\rhchXLassoOrdRmVsHarDm}{}\renewcommand{\rhchXLassoOrdRmVsHarDm}{$-$3.61}  % ADDITIONAL, NOT PRE-REGISTERED: lasso ord: HAR+RM vs HAR: DM
\providecommand{\rhchXLassoOrdRmVsHarP}{}\renewcommand{\rhchXLassoOrdRmVsHarP}{\ensuremath{<}0.001}  % ADDITIONAL, NOT PRE-REGISTERED: lasso ord: HAR+RM vs HAR: DM p
\providecommand{\rhchXLassoOrdRmVsHarShare}{}\renewcommand{\rhchXLassoOrdRmVsHarShare}{136}  % ADDITIONAL, NOT PRE-REGISTERED: lasso ord: HAR+RM vs HAR: % of the joint gain
\providecommand{\rhchXEnetBarFullVsRmPct}{}\renewcommand{\rhchXEnetBarFullVsRmPct}{$-$4.1}  % ADDITIONAL, NOT PRE-REGISTERED: enet all: full vs HAR+RM: % QLIKE change
\providecommand{\rhchXEnetBarFullVsRmDm}{}\renewcommand{\rhchXEnetBarFullVsRmDm}{$-$5.48}  % ADDITIONAL, NOT PRE-REGISTERED: enet all: full vs HAR+RM: DM
\providecommand{\rhchXEnetBarFullVsRmP}{}\renewcommand{\rhchXEnetBarFullVsRmP}{\ensuremath{<}0.001}  % ADDITIONAL, NOT PRE-REGISTERED: enet all: full vs HAR+RM: DM p
\providecommand{\rhchXEnetBarFullVsRmShare}{}\renewcommand{\rhchXEnetBarFullVsRmShare}{60}  % ADDITIONAL, NOT PRE-REGISTERED: enet all: full vs HAR+RM: % of the joint gain
\providecommand{\rhchXEnetBarRmVsHarPct}{}\renewcommand{\rhchXEnetBarRmVsHarPct}{$-$2.7}  % ADDITIONAL, NOT PRE-REGISTERED: enet all: HAR+RM vs HAR: % QLIKE change
\providecommand{\rhchXEnetBarRmVsHarDm}{}\renewcommand{\rhchXEnetBarRmVsHarDm}{$-$4.25}  % ADDITIONAL, NOT PRE-REGISTERED: enet all: HAR+RM vs HAR: DM
\providecommand{\rhchXEnetBarRmVsHarP}{}\renewcommand{\rhchXEnetBarRmVsHarP}{\ensuremath{<}0.001}  % ADDITIONAL, NOT PRE-REGISTERED: enet all: HAR+RM vs HAR: DM p
\providecommand{\rhchXEnetBarRmVsHarShare}{}\renewcommand{\rhchXEnetBarRmVsHarShare}{40}  % ADDITIONAL, NOT PRE-REGISTERED: enet all: HAR+RM vs HAR: % of the joint gain
\providecommand{\rhchXEnetEventFullVsRmPct}{}\renewcommand{\rhchXEnetEventFullVsRmPct}{$-$25.7}  % ADDITIONAL, NOT PRE-REGISTERED: enet event: full vs HAR+RM: % QLIKE change
\providecommand{\rhchXEnetEventFullVsRmDm}{}\renewcommand{\rhchXEnetEventFullVsRmDm}{$-$5.65}  % ADDITIONAL, NOT PRE-REGISTERED: enet event: full vs HAR+RM: DM
\providecommand{\rhchXEnetEventFullVsRmP}{}\renewcommand{\rhchXEnetEventFullVsRmP}{\ensuremath{<}0.001}  % ADDITIONAL, NOT PRE-REGISTERED: enet event: full vs HAR+RM: DM p
\providecommand{\rhchXEnetEventFullVsRmShare}{}\renewcommand{\rhchXEnetEventFullVsRmShare}{118}  % ADDITIONAL, NOT PRE-REGISTERED: enet event: full vs HAR+RM: % of the joint gain
\providecommand{\rhchXEnetEventRmVsHarPct}{}\renewcommand{\rhchXEnetEventRmVsHarPct}{4.2}  % ADDITIONAL, NOT PRE-REGISTERED: enet event: HAR+RM vs HAR: % QLIKE change
\providecommand{\rhchXEnetEventRmVsHarDm}{}\renewcommand{\rhchXEnetEventRmVsHarDm}{2.16}  % ADDITIONAL, NOT PRE-REGISTERED: enet event: HAR+RM vs HAR: DM
\providecommand{\rhchXEnetEventRmVsHarP}{}\renewcommand{\rhchXEnetEventRmVsHarP}{0.031}  % ADDITIONAL, NOT PRE-REGISTERED: enet event: HAR+RM vs HAR: DM p
\providecommand{\rhchXEnetEventRmVsHarShare}{}\renewcommand{\rhchXEnetEventRmVsHarShare}{$-$18}  % ADDITIONAL, NOT PRE-REGISTERED: enet event: HAR+RM vs HAR: % of the joint gain
\providecommand{\rhchXEnetOrdFullVsRmPct}{}\renewcommand{\rhchXEnetOrdFullVsRmPct}{$-$0.1}  % ADDITIONAL, NOT PRE-REGISTERED: enet ord: full vs HAR+RM: % QLIKE change
\providecommand{\rhchXEnetOrdFullVsRmDm}{}\renewcommand{\rhchXEnetOrdFullVsRmDm}{$-$0.26}  % ADDITIONAL, NOT PRE-REGISTERED: enet ord: full vs HAR+RM: DM
\providecommand{\rhchXEnetOrdFullVsRmP}{}\renewcommand{\rhchXEnetOrdFullVsRmP}{0.793}  % ADDITIONAL, NOT PRE-REGISTERED: enet ord: full vs HAR+RM: DM p
\providecommand{\rhchXEnetOrdFullVsRmShare}{}\renewcommand{\rhchXEnetOrdFullVsRmShare}{3}  % ADDITIONAL, NOT PRE-REGISTERED: enet ord: full vs HAR+RM: % of the joint gain
\providecommand{\rhchXEnetOrdRmVsHarPct}{}\renewcommand{\rhchXEnetOrdRmVsHarPct}{$-$3.9}  % ADDITIONAL, NOT PRE-REGISTERED: enet ord: HAR+RM vs HAR: % QLIKE change
\providecommand{\rhchXEnetOrdRmVsHarDm}{}\renewcommand{\rhchXEnetOrdRmVsHarDm}{$-$5.53}  % ADDITIONAL, NOT PRE-REGISTERED: enet ord: HAR+RM vs HAR: DM
\providecommand{\rhchXEnetOrdRmVsHarP}{}\renewcommand{\rhchXEnetOrdRmVsHarP}{\ensuremath{<}0.001}  % ADDITIONAL, NOT PRE-REGISTERED: enet ord: HAR+RM vs HAR: DM p
\providecommand{\rhchXEnetOrdRmVsHarShare}{}\renewcommand{\rhchXEnetOrdRmVsHarShare}{97}  % ADDITIONAL, NOT PRE-REGISTERED: enet ord: HAR+RM vs HAR: % of the joint gain
\providecommand{\rhchXLgbmBarFullVsRmPct}{}\renewcommand{\rhchXLgbmBarFullVsRmPct}{$-$11.5}  % ADDITIONAL, NOT PRE-REGISTERED: lgbm all: full vs HAR+RM: % QLIKE change
\providecommand{\rhchXLgbmBarFullVsRmDm}{}\renewcommand{\rhchXLgbmBarFullVsRmDm}{$-$7.62}  % ADDITIONAL, NOT PRE-REGISTERED: lgbm all: full vs HAR+RM: DM
\providecommand{\rhchXLgbmBarFullVsRmP}{}\renewcommand{\rhchXLgbmBarFullVsRmP}{\ensuremath{<}0.001}  % ADDITIONAL, NOT PRE-REGISTERED: lgbm all: full vs HAR+RM: DM p
\providecommand{\rhchXLgbmBarFullVsRmShare}{}\renewcommand{\rhchXLgbmBarFullVsRmShare}{79}  % ADDITIONAL, NOT PRE-REGISTERED: lgbm all: full vs HAR+RM: % of the joint gain
\providecommand{\rhchXLgbmBarRmVsHarPct}{}\renewcommand{\rhchXLgbmBarRmVsHarPct}{$-$3.0}  % ADDITIONAL, NOT PRE-REGISTERED: lgbm all: HAR+RM vs HAR: % QLIKE change
\providecommand{\rhchXLgbmBarRmVsHarDm}{}\renewcommand{\rhchXLgbmBarRmVsHarDm}{$-$6.32}  % ADDITIONAL, NOT PRE-REGISTERED: lgbm all: HAR+RM vs HAR: DM
\providecommand{\rhchXLgbmBarRmVsHarP}{}\renewcommand{\rhchXLgbmBarRmVsHarP}{\ensuremath{<}0.001}  % ADDITIONAL, NOT PRE-REGISTERED: lgbm all: HAR+RM vs HAR: DM p
\providecommand{\rhchXLgbmBarRmVsHarShare}{}\renewcommand{\rhchXLgbmBarRmVsHarShare}{21}  % ADDITIONAL, NOT PRE-REGISTERED: lgbm all: HAR+RM vs HAR: % of the joint gain
\providecommand{\rhchXLgbmEventFullVsRmPct}{}\renewcommand{\rhchXLgbmEventFullVsRmPct}{$-$64.7}  % ADDITIONAL, NOT PRE-REGISTERED: lgbm event: full vs HAR+RM: % QLIKE change
\providecommand{\rhchXLgbmEventFullVsRmDm}{}\renewcommand{\rhchXLgbmEventFullVsRmDm}{$-$7.56}  % ADDITIONAL, NOT PRE-REGISTERED: lgbm event: full vs HAR+RM: DM
\providecommand{\rhchXLgbmEventFullVsRmP}{}\renewcommand{\rhchXLgbmEventFullVsRmP}{\ensuremath{<}0.001}  % ADDITIONAL, NOT PRE-REGISTERED: lgbm event: full vs HAR+RM: DM p
\providecommand{\rhchXLgbmEventFullVsRmShare}{}\renewcommand{\rhchXLgbmEventFullVsRmShare}{111}  % ADDITIONAL, NOT PRE-REGISTERED: lgbm event: full vs HAR+RM: % of the joint gain
\providecommand{\rhchXLgbmEventRmVsHarPct}{}\renewcommand{\rhchXLgbmEventRmVsHarPct}{6.8}  % ADDITIONAL, NOT PRE-REGISTERED: lgbm event: HAR+RM vs HAR: % QLIKE change
\providecommand{\rhchXLgbmEventRmVsHarDm}{}\renewcommand{\rhchXLgbmEventRmVsHarDm}{2.58}  % ADDITIONAL, NOT PRE-REGISTERED: lgbm event: HAR+RM vs HAR: DM
\providecommand{\rhchXLgbmEventRmVsHarP}{}\renewcommand{\rhchXLgbmEventRmVsHarP}{0.010}  % ADDITIONAL, NOT PRE-REGISTERED: lgbm event: HAR+RM vs HAR: DM p
\providecommand{\rhchXLgbmEventRmVsHarShare}{}\renewcommand{\rhchXLgbmEventRmVsHarShare}{$-$11}  % ADDITIONAL, NOT PRE-REGISTERED: lgbm event: HAR+RM vs HAR: % of the joint gain
\providecommand{\rhchXLgbmOrdFullVsRmPct}{}\renewcommand{\rhchXLgbmOrdFullVsRmPct}{$-$1.5}  % ADDITIONAL, NOT PRE-REGISTERED: lgbm ord: full vs HAR+RM: % QLIKE change
\providecommand{\rhchXLgbmOrdFullVsRmDm}{}\renewcommand{\rhchXLgbmOrdFullVsRmDm}{$-$4.30}  % ADDITIONAL, NOT PRE-REGISTERED: lgbm ord: full vs HAR+RM: DM
\providecommand{\rhchXLgbmOrdFullVsRmP}{}\renewcommand{\rhchXLgbmOrdFullVsRmP}{\ensuremath{<}0.001}  % ADDITIONAL, NOT PRE-REGISTERED: lgbm ord: full vs HAR+RM: DM p
\providecommand{\rhchXLgbmOrdFullVsRmShare}{}\renewcommand{\rhchXLgbmOrdFullVsRmShare}{23}  % ADDITIONAL, NOT PRE-REGISTERED: lgbm ord: full vs HAR+RM: % of the joint gain
\providecommand{\rhchXLgbmOrdRmVsHarPct}{}\renewcommand{\rhchXLgbmOrdRmVsHarPct}{$-$4.6}  % ADDITIONAL, NOT PRE-REGISTERED: lgbm ord: HAR+RM vs HAR: % QLIKE change
\providecommand{\rhchXLgbmOrdRmVsHarDm}{}\renewcommand{\rhchXLgbmOrdRmVsHarDm}{$-$9.55}  % ADDITIONAL, NOT PRE-REGISTERED: lgbm ord: HAR+RM vs HAR: DM
\providecommand{\rhchXLgbmOrdRmVsHarP}{}\renewcommand{\rhchXLgbmOrdRmVsHarP}{\ensuremath{<}0.001}  % ADDITIONAL, NOT PRE-REGISTERED: lgbm ord: HAR+RM vs HAR: DM p
\providecommand{\rhchXLgbmOrdRmVsHarShare}{}\renewcommand{\rhchXLgbmOrdRmVsHarShare}{77}  % ADDITIONAL, NOT PRE-REGISTERED: lgbm ord: HAR+RM vs HAR: % of the joint gain
\providecommand{\rhchXXgbBarFullVsRmPct}{}\renewcommand{\rhchXXgbBarFullVsRmPct}{$-$11.6}  % ADDITIONAL, NOT PRE-REGISTERED: xgb all: full vs HAR+RM: % QLIKE change
\providecommand{\rhchXXgbBarFullVsRmDm}{}\renewcommand{\rhchXXgbBarFullVsRmDm}{$-$7.19}  % ADDITIONAL, NOT PRE-REGISTERED: xgb all: full vs HAR+RM: DM
\providecommand{\rhchXXgbBarFullVsRmP}{}\renewcommand{\rhchXXgbBarFullVsRmP}{\ensuremath{<}0.001}  % ADDITIONAL, NOT PRE-REGISTERED: xgb all: full vs HAR+RM: DM p
\providecommand{\rhchXXgbBarFullVsRmShare}{}\renewcommand{\rhchXXgbBarFullVsRmShare}{77}  % ADDITIONAL, NOT PRE-REGISTERED: xgb all: full vs HAR+RM: % of the joint gain
\providecommand{\rhchXXgbBarRmVsHarPct}{}\renewcommand{\rhchXXgbBarRmVsHarPct}{$-$3.4}  % ADDITIONAL, NOT PRE-REGISTERED: xgb all: HAR+RM vs HAR: % QLIKE change
\providecommand{\rhchXXgbBarRmVsHarDm}{}\renewcommand{\rhchXXgbBarRmVsHarDm}{$-$7.51}  % ADDITIONAL, NOT PRE-REGISTERED: xgb all: HAR+RM vs HAR: DM
\providecommand{\rhchXXgbBarRmVsHarP}{}\renewcommand{\rhchXXgbBarRmVsHarP}{\ensuremath{<}0.001}  % ADDITIONAL, NOT PRE-REGISTERED: xgb all: HAR+RM vs HAR: DM p
\providecommand{\rhchXXgbBarRmVsHarShare}{}\renewcommand{\rhchXXgbBarRmVsHarShare}{23}  % ADDITIONAL, NOT PRE-REGISTERED: xgb all: HAR+RM vs HAR: % of the joint gain
\providecommand{\rhchXXgbEventFullVsRmPct}{}\renewcommand{\rhchXXgbEventFullVsRmPct}{$-$65.8}  % ADDITIONAL, NOT PRE-REGISTERED: xgb event: full vs HAR+RM: % QLIKE change
\providecommand{\rhchXXgbEventFullVsRmDm}{}\renewcommand{\rhchXXgbEventFullVsRmDm}{$-$6.93}  % ADDITIONAL, NOT PRE-REGISTERED: xgb event: full vs HAR+RM: DM
\providecommand{\rhchXXgbEventFullVsRmP}{}\renewcommand{\rhchXXgbEventFullVsRmP}{\ensuremath{<}0.001}  % ADDITIONAL, NOT PRE-REGISTERED: xgb event: full vs HAR+RM: DM p
\providecommand{\rhchXXgbEventFullVsRmShare}{}\renewcommand{\rhchXXgbEventFullVsRmShare}{109}  % ADDITIONAL, NOT PRE-REGISTERED: xgb event: full vs HAR+RM: % of the joint gain
\providecommand{\rhchXXgbEventRmVsHarPct}{}\renewcommand{\rhchXXgbEventRmVsHarPct}{5.5}  % ADDITIONAL, NOT PRE-REGISTERED: xgb event: HAR+RM vs HAR: % QLIKE change
\providecommand{\rhchXXgbEventRmVsHarDm}{}\renewcommand{\rhchXXgbEventRmVsHarDm}{2.46}  % ADDITIONAL, NOT PRE-REGISTERED: xgb event: HAR+RM vs HAR: DM
\providecommand{\rhchXXgbEventRmVsHarP}{}\renewcommand{\rhchXXgbEventRmVsHarP}{0.014}  % ADDITIONAL, NOT PRE-REGISTERED: xgb event: HAR+RM vs HAR: DM p
\providecommand{\rhchXXgbEventRmVsHarShare}{}\renewcommand{\rhchXXgbEventRmVsHarShare}{$-$9}  % ADDITIONAL, NOT PRE-REGISTERED: xgb event: HAR+RM vs HAR: % of the joint gain
\providecommand{\rhchXXgbOrdFullVsRmPct}{}\renewcommand{\rhchXXgbOrdFullVsRmPct}{$-$1.2}  % ADDITIONAL, NOT PRE-REGISTERED: xgb ord: full vs HAR+RM: % QLIKE change
\providecommand{\rhchXXgbOrdFullVsRmDm}{}\renewcommand{\rhchXXgbOrdFullVsRmDm}{$-$3.23}  % ADDITIONAL, NOT PRE-REGISTERED: xgb ord: full vs HAR+RM: DM
\providecommand{\rhchXXgbOrdFullVsRmP}{}\renewcommand{\rhchXXgbOrdFullVsRmP}{0.001}  % ADDITIONAL, NOT PRE-REGISTERED: xgb ord: full vs HAR+RM: DM p
\providecommand{\rhchXXgbOrdFullVsRmShare}{}\renewcommand{\rhchXXgbOrdFullVsRmShare}{19}  % ADDITIONAL, NOT PRE-REGISTERED: xgb ord: full vs HAR+RM: % of the joint gain
\providecommand{\rhchXXgbOrdRmVsHarPct}{}\renewcommand{\rhchXXgbOrdRmVsHarPct}{$-$4.9}  % ADDITIONAL, NOT PRE-REGISTERED: xgb ord: HAR+RM vs HAR: % QLIKE change
\providecommand{\rhchXXgbOrdRmVsHarDm}{}\renewcommand{\rhchXXgbOrdRmVsHarDm}{$-$10.37}  % ADDITIONAL, NOT PRE-REGISTERED: xgb ord: HAR+RM vs HAR: DM
\providecommand{\rhchXXgbOrdRmVsHarP}{}\renewcommand{\rhchXXgbOrdRmVsHarP}{\ensuremath{<}0.001}  % ADDITIONAL, NOT PRE-REGISTERED: xgb ord: HAR+RM vs HAR: DM p
\providecommand{\rhchXXgbOrdRmVsHarShare}{}\renewcommand{\rhchXXgbOrdRmVsHarShare}{81}  % ADDITIONAL, NOT PRE-REGISTERED: xgb ord: HAR+RM vs HAR: % of the joint gain
\providecommand{\rhchXRfBarFullVsRmPct}{}\renewcommand{\rhchXRfBarFullVsRmPct}{$-$8.4}  % ADDITIONAL, NOT PRE-REGISTERED: rf all: full vs HAR+RM: % QLIKE change
\providecommand{\rhchXRfBarFullVsRmDm}{}\renewcommand{\rhchXRfBarFullVsRmDm}{$-$6.48}  % ADDITIONAL, NOT PRE-REGISTERED: rf all: full vs HAR+RM: DM
\providecommand{\rhchXRfBarFullVsRmP}{}\renewcommand{\rhchXRfBarFullVsRmP}{\ensuremath{<}0.001}  % ADDITIONAL, NOT PRE-REGISTERED: rf all: full vs HAR+RM: DM p
\providecommand{\rhchXRfBarFullVsRmShare}{}\renewcommand{\rhchXRfBarFullVsRmShare}{72}  % ADDITIONAL, NOT PRE-REGISTERED: rf all: full vs HAR+RM: % of the joint gain
\providecommand{\rhchXRfBarRmVsHarPct}{}\renewcommand{\rhchXRfBarRmVsHarPct}{$-$3.2}  % ADDITIONAL, NOT PRE-REGISTERED: rf all: HAR+RM vs HAR: % QLIKE change
\providecommand{\rhchXRfBarRmVsHarDm}{}\renewcommand{\rhchXRfBarRmVsHarDm}{$-$6.79}  % ADDITIONAL, NOT PRE-REGISTERED: rf all: HAR+RM vs HAR: DM
\providecommand{\rhchXRfBarRmVsHarP}{}\renewcommand{\rhchXRfBarRmVsHarP}{\ensuremath{<}0.001}  % ADDITIONAL, NOT PRE-REGISTERED: rf all: HAR+RM vs HAR: DM p
\providecommand{\rhchXRfBarRmVsHarShare}{}\renewcommand{\rhchXRfBarRmVsHarShare}{28}  % ADDITIONAL, NOT PRE-REGISTERED: rf all: HAR+RM vs HAR: % of the joint gain
\providecommand{\rhchXRfEventFullVsRmPct}{}\renewcommand{\rhchXRfEventFullVsRmPct}{$-$47.6}  % ADDITIONAL, NOT PRE-REGISTERED: rf event: full vs HAR+RM: % QLIKE change
\providecommand{\rhchXRfEventFullVsRmDm}{}\renewcommand{\rhchXRfEventFullVsRmDm}{$-$7.12}  % ADDITIONAL, NOT PRE-REGISTERED: rf event: full vs HAR+RM: DM
\providecommand{\rhchXRfEventFullVsRmP}{}\renewcommand{\rhchXRfEventFullVsRmP}{\ensuremath{<}0.001}  % ADDITIONAL, NOT PRE-REGISTERED: rf event: full vs HAR+RM: DM p
\providecommand{\rhchXRfEventFullVsRmShare}{}\renewcommand{\rhchXRfEventFullVsRmShare}{112}  % ADDITIONAL, NOT PRE-REGISTERED: rf event: full vs HAR+RM: % of the joint gain
\providecommand{\rhchXRfEventRmVsHarPct}{}\renewcommand{\rhchXRfEventRmVsHarPct}{5.3}  % ADDITIONAL, NOT PRE-REGISTERED: rf event: HAR+RM vs HAR: % QLIKE change
\providecommand{\rhchXRfEventRmVsHarDm}{}\renewcommand{\rhchXRfEventRmVsHarDm}{1.88}  % ADDITIONAL, NOT PRE-REGISTERED: rf event: HAR+RM vs HAR: DM
\providecommand{\rhchXRfEventRmVsHarP}{}\renewcommand{\rhchXRfEventRmVsHarP}{0.061}  % ADDITIONAL, NOT PRE-REGISTERED: rf event: HAR+RM vs HAR: DM p
\providecommand{\rhchXRfEventRmVsHarShare}{}\renewcommand{\rhchXRfEventRmVsHarShare}{$-$12}  % ADDITIONAL, NOT PRE-REGISTERED: rf event: HAR+RM vs HAR: % of the joint gain
\providecommand{\rhchXRfOrdFullVsRmPct}{}\renewcommand{\rhchXRfOrdFullVsRmPct}{$-$0.7}  % ADDITIONAL, NOT PRE-REGISTERED: rf ord: full vs HAR+RM: % QLIKE change
\providecommand{\rhchXRfOrdFullVsRmDm}{}\renewcommand{\rhchXRfOrdFullVsRmDm}{$-$2.22}  % ADDITIONAL, NOT PRE-REGISTERED: rf ord: full vs HAR+RM: DM
\providecommand{\rhchXRfOrdFullVsRmP}{}\renewcommand{\rhchXRfOrdFullVsRmP}{0.026}  % ADDITIONAL, NOT PRE-REGISTERED: rf ord: full vs HAR+RM: DM p
\providecommand{\rhchXRfOrdFullVsRmShare}{}\renewcommand{\rhchXRfOrdFullVsRmShare}{12}  % ADDITIONAL, NOT PRE-REGISTERED: rf ord: full vs HAR+RM: % of the joint gain
\providecommand{\rhchXRfOrdRmVsHarPct}{}\renewcommand{\rhchXRfOrdRmVsHarPct}{$-$4.7}  % ADDITIONAL, NOT PRE-REGISTERED: rf ord: HAR+RM vs HAR: % QLIKE change
\providecommand{\rhchXRfOrdRmVsHarDm}{}\renewcommand{\rhchXRfOrdRmVsHarDm}{$-$10.68}  % ADDITIONAL, NOT PRE-REGISTERED: rf ord: HAR+RM vs HAR: DM
\providecommand{\rhchXRfOrdRmVsHarP}{}\renewcommand{\rhchXRfOrdRmVsHarP}{\ensuremath{<}0.001}  % ADDITIONAL, NOT PRE-REGISTERED: rf ord: HAR+RM vs HAR: DM p
\providecommand{\rhchXRfOrdRmVsHarShare}{}\renewcommand{\rhchXRfOrdRmVsHarShare}{88}  % ADDITIONAL, NOT PRE-REGISTERED: rf ord: HAR+RM vs HAR: % of the joint gain
\providecommand{\rhchXLinBarFullVsRmPctMin}{}\renewcommand{\rhchXLinBarFullVsRmPctMin}{$-$7.0}  % ADDITIONAL: Lin all: lowest %
\providecommand{\rhchXLinBarFullVsRmPctMax}{}\renewcommand{\rhchXLinBarFullVsRmPctMax}{$-$3.5}  % ADDITIONAL: Lin all: highest %
\providecommand{\rhchXLinBarFullVsRmNSig}{}\renewcommand{\rhchXLinBarFullVsRmNSig}{3}  % ADDITIONAL: Lin all: classes whose full set beats HAR+RM at 5%
\providecommand{\rhchXLinEventFullVsRmPctMin}{}\renewcommand{\rhchXLinEventFullVsRmPctMin}{$-$49.5}  % ADDITIONAL: Lin event: lowest %
\providecommand{\rhchXLinEventFullVsRmPctMax}{}\renewcommand{\rhchXLinEventFullVsRmPctMax}{$-$25.7}  % ADDITIONAL: Lin event: highest %
\providecommand{\rhchXLinEventFullVsRmNSig}{}\renewcommand{\rhchXLinEventFullVsRmNSig}{3}  % ADDITIONAL: Lin event: classes whose full set beats HAR+RM at 5%
\providecommand{\rhchXLinOrdFullVsRmPctMin}{}\renewcommand{\rhchXLinOrdFullVsRmPctMin}{$-$0.1}  % ADDITIONAL: Lin ord: lowest %
\providecommand{\rhchXLinOrdFullVsRmPctMax}{}\renewcommand{\rhchXLinOrdFullVsRmPctMax}{0.9}  % ADDITIONAL: Lin ord: highest %
\providecommand{\rhchXLinOrdFullVsRmNSig}{}\renewcommand{\rhchXLinOrdFullVsRmNSig}{0}  % ADDITIONAL: Lin ord: classes whose full set beats HAR+RM at 5%
\providecommand{\rhchXTreeBarFullVsRmPctMin}{}\renewcommand{\rhchXTreeBarFullVsRmPctMin}{$-$11.6}  % ADDITIONAL: Tree all: lowest %
\providecommand{\rhchXTreeBarFullVsRmPctMax}{}\renewcommand{\rhchXTreeBarFullVsRmPctMax}{$-$8.4}  % ADDITIONAL: Tree all: highest %
\providecommand{\rhchXTreeBarFullVsRmNSig}{}\renewcommand{\rhchXTreeBarFullVsRmNSig}{3}  % ADDITIONAL: Tree all: classes whose full set beats HAR+RM at 5%
\providecommand{\rhchXTreeEventFullVsRmPctMin}{}\renewcommand{\rhchXTreeEventFullVsRmPctMin}{$-$65.8}  % ADDITIONAL: Tree event: lowest %
\providecommand{\rhchXTreeEventFullVsRmPctMax}{}\renewcommand{\rhchXTreeEventFullVsRmPctMax}{$-$47.6}  % ADDITIONAL: Tree event: highest %
\providecommand{\rhchXTreeEventFullVsRmNSig}{}\renewcommand{\rhchXTreeEventFullVsRmNSig}{3}  % ADDITIONAL: Tree event: classes whose full set beats HAR+RM at 5%
\providecommand{\rhchXTreeOrdFullVsRmPctMin}{}\renewcommand{\rhchXTreeOrdFullVsRmPctMin}{$-$1.5}  % ADDITIONAL: Tree ord: lowest %
\providecommand{\rhchXTreeOrdFullVsRmPctMax}{}\renewcommand{\rhchXTreeOrdFullVsRmPctMax}{$-$0.7}  % ADDITIONAL: Tree ord: highest %
\providecommand{\rhchXTreeOrdFullVsRmNSig}{}\renewcommand{\rhchXTreeOrdFullVsRmNSig}{3}  % ADDITIONAL: Tree ord: classes whose full set beats HAR+RM at 5%
\providecommand{\rhchRidgeBarLadFullPct}{}\renewcommand{\rhchRidgeBarLadFullPct}{$-$3.6}  % ridge all: ladder k=46 vs 0 %
\providecommand{\rhchRidgeBarROne}{}\renewcommand{\rhchRidgeBarROne}{56}  % ridge all: ladder R_1 (%)
\providecommand{\rhchRidgeBarRThree}{}\renewcommand{\rhchRidgeBarRThree}{84}  % ridge all: ladder R_3 (%)
\providecommand{\rhchRidgeBarRFive}{}\renewcommand{\rhchRidgeBarRFive}{145}  % ridge all: ladder R_5 (%)
\providecommand{\rhchRidgeBarRTen}{}\renewcommand{\rhchRidgeBarRTen}{188}  % ridge all: ladder R_10 (%)
\providecommand{\rhchRidgeBarRTwenty}{}\renewcommand{\rhchRidgeBarRTwenty}{171}  % ridge all: ladder R_20 (%)
\providecommand{\rhchRidgeEventLadFullPct}{}\renewcommand{\rhchRidgeEventLadFullPct}{$-$23.2}  % ridge event: ladder k=46 vs 0 %
\providecommand{\rhchRidgeEventROne}{}\renewcommand{\rhchRidgeEventROne}{66}  % ridge event: ladder R_1 (%)
\providecommand{\rhchRidgeEventRThree}{}\renewcommand{\rhchRidgeEventRThree}{95}  % ridge event: ladder R_3 (%)
\providecommand{\rhchRidgeEventRFive}{}\renewcommand{\rhchRidgeEventRFive}{96}  % ridge event: ladder R_5 (%)
\providecommand{\rhchRidgeEventRTen}{}\renewcommand{\rhchRidgeEventRTen}{101}  % ridge event: ladder R_10 (%)
\providecommand{\rhchRidgeEventRTwenty}{}\renewcommand{\rhchRidgeEventRTwenty}{99}  % ridge event: ladder R_20 (%)
\providecommand{\rhchRidgeOrdLadFullPct}{}\renewcommand{\rhchRidgeOrdLadFullPct}{$-$0.3}  % ridge ord: ladder k=46 vs 0 %
\providecommand{\rhchRidgeOrdROne}{}\renewcommand{\rhchRidgeOrdROne}{$-$70}  % ridge ord: ladder R_1 (%)
\providecommand{\rhchRidgeOrdRThree}{}\renewcommand{\rhchRidgeOrdRThree}{$-$75}  % ridge ord: ladder R_3 (%)
\providecommand{\rhchRidgeOrdRFive}{}\renewcommand{\rhchRidgeOrdRFive}{824}  % ridge ord: ladder R_5 (%)
\providecommand{\rhchRidgeOrdRTen}{}\renewcommand{\rhchRidgeOrdRTen}{1391}  % ridge ord: ladder R_10 (%)
\providecommand{\rhchRidgeOrdRTwenty}{}\renewcommand{\rhchRidgeOrdRTwenty}{1168}  % ridge ord: ladder R_20 (%)
\providecommand{\rhchLassoBarLadFullPct}{}\renewcommand{\rhchLassoBarLadFullPct}{$-$6.0}  % lasso all: ladder k=46 vs 0 %
\providecommand{\rhchLassoBarROne}{}\renewcommand{\rhchLassoBarROne}{69}  % lasso all: ladder R_1 (%)
\providecommand{\rhchLassoBarRThree}{}\renewcommand{\rhchLassoBarRThree}{63}  % lasso all: ladder R_3 (%)
\providecommand{\rhchLassoBarRFive}{}\renewcommand{\rhchLassoBarRFive}{108}  % lasso all: ladder R_5 (%)
\providecommand{\rhchLassoBarRTen}{}\renewcommand{\rhchLassoBarRTen}{155}  % lasso all: ladder R_10 (%)
\providecommand{\rhchLassoBarRTwenty}{}\renewcommand{\rhchLassoBarRTwenty}{139}  % lasso all: ladder R_20 (%)
\providecommand{\rhchLassoEventLadFullPct}{}\renewcommand{\rhchLassoEventLadFullPct}{$-$46.9}  % lasso event: ladder k=46 vs 0 %
\providecommand{\rhchLassoEventROne}{}\renewcommand{\rhchLassoEventROne}{64}  % lasso event: ladder R_1 (%)
\providecommand{\rhchLassoEventRThree}{}\renewcommand{\rhchLassoEventRThree}{67}  % lasso event: ladder R_3 (%)
\providecommand{\rhchLassoEventRFive}{}\renewcommand{\rhchLassoEventRFive}{92}  % lasso event: ladder R_5 (%)
\providecommand{\rhchLassoEventRTen}{}\renewcommand{\rhchLassoEventRTen}{96}  % lasso event: ladder R_10 (%)
\providecommand{\rhchLassoEventRTwenty}{}\renewcommand{\rhchLassoEventRTwenty}{97}  % lasso event: ladder R_20 (%)
\providecommand{\rhchLassoOrdLadFullPct}{}\renewcommand{\rhchLassoOrdLadFullPct}{1.0}  % lasso ord: ladder k=46 vs 0 %
\providecommand{\rhchLassoOrdROne}{}\renewcommand{\rhchLassoOrdROne}{--}  % lasso ord: ladder R_1 (%)
\providecommand{\rhchLassoOrdRThree}{}\renewcommand{\rhchLassoOrdRThree}{--}  % lasso ord: ladder R_3 (%)
\providecommand{\rhchLassoOrdRFive}{}\renewcommand{\rhchLassoOrdRFive}{--}  % lasso ord: ladder R_5 (%)
\providecommand{\rhchLassoOrdRTen}{}\renewcommand{\rhchLassoOrdRTen}{--}  % lasso ord: ladder R_10 (%)
\providecommand{\rhchLassoOrdRTwenty}{}\renewcommand{\rhchLassoOrdRTwenty}{--}  % lasso ord: ladder R_20 (%)
\providecommand{\rhchEnetBarLadFullPct}{}\renewcommand{\rhchEnetBarLadFullPct}{$-$5.3}  % enet all: ladder k=46 vs 0 %
\providecommand{\rhchEnetBarROne}{}\renewcommand{\rhchEnetBarROne}{55}  % enet all: ladder R_1 (%)
\providecommand{\rhchEnetBarRThree}{}\renewcommand{\rhchEnetBarRThree}{58}  % enet all: ladder R_3 (%)
\providecommand{\rhchEnetBarRFive}{}\renewcommand{\rhchEnetBarRFive}{99}  % enet all: ladder R_5 (%)
\providecommand{\rhchEnetBarRTen}{}\renewcommand{\rhchEnetBarRTen}{130}  % enet all: ladder R_10 (%)
\providecommand{\rhchEnetBarRTwenty}{}\renewcommand{\rhchEnetBarRTwenty}{120}  % enet all: ladder R_20 (%)
\providecommand{\rhchEnetEventLadFullPct}{}\renewcommand{\rhchEnetEventLadFullPct}{$-$23.0}  % enet event: ladder k=46 vs 0 %
\providecommand{\rhchEnetEventROne}{}\renewcommand{\rhchEnetEventROne}{90}  % enet event: ladder R_1 (%)
\providecommand{\rhchEnetEventRThree}{}\renewcommand{\rhchEnetEventRThree}{99}  % enet event: ladder R_3 (%)
\providecommand{\rhchEnetEventRFive}{}\renewcommand{\rhchEnetEventRFive}{97}  % enet event: ladder R_5 (%)
\providecommand{\rhchEnetEventRTen}{}\renewcommand{\rhchEnetEventRTen}{102}  % enet event: ladder R_10 (%)
\providecommand{\rhchEnetEventRTwenty}{}\renewcommand{\rhchEnetEventRTwenty}{98}  % enet event: ladder R_20 (%)
\providecommand{\rhchEnetOrdLadFullPct}{}\renewcommand{\rhchEnetOrdLadFullPct}{$-$2.2}  % enet ord: ladder k=46 vs 0 %
\providecommand{\rhchEnetOrdROne}{}\renewcommand{\rhchEnetOrdROne}{$-$7}  % enet ord: ladder R_1 (%)
\providecommand{\rhchEnetOrdRThree}{}\renewcommand{\rhchEnetOrdRThree}{$-$16}  % enet ord: ladder R_3 (%)
\providecommand{\rhchEnetOrdRFive}{}\renewcommand{\rhchEnetOrdRFive}{102}  % enet ord: ladder R_5 (%)
\providecommand{\rhchEnetOrdRTen}{}\renewcommand{\rhchEnetOrdRTen}{180}  % enet ord: ladder R_10 (%)
\providecommand{\rhchEnetOrdRTwenty}{}\renewcommand{\rhchEnetOrdRTwenty}{159}  % enet ord: ladder R_20 (%)
\providecommand{\rhchLgbmBarLadFullPct}{}\renewcommand{\rhchLgbmBarLadFullPct}{$-$13.6}  % lgbm all: ladder k=46 vs 0 %
\providecommand{\rhchLgbmBarROne}{}\renewcommand{\rhchLgbmBarROne}{3}  % lgbm all: ladder R_1 (%)
\providecommand{\rhchLgbmBarRThree}{}\renewcommand{\rhchLgbmBarRThree}{34}  % lgbm all: ladder R_3 (%)
\providecommand{\rhchLgbmBarRFive}{}\renewcommand{\rhchLgbmBarRFive}{34}  % lgbm all: ladder R_5 (%)
\providecommand{\rhchLgbmBarRTen}{}\renewcommand{\rhchLgbmBarRTen}{95}  % lgbm all: ladder R_10 (%)
\providecommand{\rhchLgbmBarRTwenty}{}\renewcommand{\rhchLgbmBarRTwenty}{96}  % lgbm all: ladder R_20 (%)
\providecommand{\rhchLgbmEventLadFullPct}{}\renewcommand{\rhchLgbmEventLadFullPct}{$-$62.0}  % lgbm event: ladder k=46 vs 0 %
\providecommand{\rhchLgbmEventROne}{}\renewcommand{\rhchLgbmEventROne}{$-$1}  % lgbm event: ladder R_1 (%)
\providecommand{\rhchLgbmEventRThree}{}\renewcommand{\rhchLgbmEventRThree}{5}  % lgbm event: ladder R_3 (%)
\providecommand{\rhchLgbmEventRFive}{}\renewcommand{\rhchLgbmEventRFive}{3}  % lgbm event: ladder R_5 (%)
\providecommand{\rhchLgbmEventRTen}{}\renewcommand{\rhchLgbmEventRTen}{100}  % lgbm event: ladder R_10 (%)
\providecommand{\rhchLgbmEventRTwenty}{}\renewcommand{\rhchLgbmEventRTwenty}{100}  % lgbm event: ladder R_20 (%)
\providecommand{\rhchLgbmOrdLadFullPct}{}\renewcommand{\rhchLgbmOrdLadFullPct}{$-$5.5}  % lgbm ord: ladder k=46 vs 0 %
\providecommand{\rhchLgbmOrdROne}{}\renewcommand{\rhchLgbmOrdROne}{12}  % lgbm ord: ladder R_1 (%)
\providecommand{\rhchLgbmOrdRThree}{}\renewcommand{\rhchLgbmOrdRThree}{88}  % lgbm ord: ladder R_3 (%)
\providecommand{\rhchLgbmOrdRFive}{}\renewcommand{\rhchLgbmOrdRFive}{93}  % lgbm ord: ladder R_5 (%)
\providecommand{\rhchLgbmOrdRTen}{}\renewcommand{\rhchLgbmOrdRTen}{86}  % lgbm ord: ladder R_10 (%)
\providecommand{\rhchLgbmOrdRTwenty}{}\renewcommand{\rhchLgbmOrdRTwenty}{90}  % lgbm ord: ladder R_20 (%)
\providecommand{\rhchXgbBarLadFullPct}{}\renewcommand{\rhchXgbBarLadFullPct}{$-$13.6}  % xgb all: ladder k=46 vs 0 %
\providecommand{\rhchXgbBarROne}{}\renewcommand{\rhchXgbBarROne}{4}  % xgb all: ladder R_1 (%)
\providecommand{\rhchXgbBarRThree}{}\renewcommand{\rhchXgbBarRThree}{27}  % xgb all: ladder R_3 (%)
\providecommand{\rhchXgbBarRFive}{}\renewcommand{\rhchXgbBarRFive}{36}  % xgb all: ladder R_5 (%)
\providecommand{\rhchXgbBarRTen}{}\renewcommand{\rhchXgbBarRTen}{98}  % xgb all: ladder R_10 (%)
\providecommand{\rhchXgbBarRTwenty}{}\renewcommand{\rhchXgbBarRTwenty}{99}  % xgb all: ladder R_20 (%)
\providecommand{\rhchXgbEventLadFullPct}{}\renewcommand{\rhchXgbEventLadFullPct}{$-$63.5}  % xgb event: ladder k=46 vs 0 %
\providecommand{\rhchXgbEventROne}{}\renewcommand{\rhchXgbEventROne}{0}  % xgb event: ladder R_1 (%)
\providecommand{\rhchXgbEventRThree}{}\renewcommand{\rhchXgbEventRThree}{10}  % xgb event: ladder R_3 (%)
\providecommand{\rhchXgbEventRFive}{}\renewcommand{\rhchXgbEventRFive}{3}  % xgb event: ladder R_5 (%)
\providecommand{\rhchXgbEventRTen}{}\renewcommand{\rhchXgbEventRTen}{99}  % xgb event: ladder R_10 (%)
\providecommand{\rhchXgbEventRTwenty}{}\renewcommand{\rhchXgbEventRTwenty}{100}  % xgb event: ladder R_20 (%)
\providecommand{\rhchXgbOrdLadFullPct}{}\renewcommand{\rhchXgbOrdLadFullPct}{$-$5.1}  % xgb ord: ladder k=46 vs 0 %
\providecommand{\rhchXgbOrdROne}{}\renewcommand{\rhchXgbOrdROne}{13}  % xgb ord: ladder R_1 (%)
\providecommand{\rhchXgbOrdRThree}{}\renewcommand{\rhchXgbOrdRThree}{64}  % xgb ord: ladder R_3 (%)
\providecommand{\rhchXgbOrdRFive}{}\renewcommand{\rhchXgbOrdRFive}{108}  % xgb ord: ladder R_5 (%)
\providecommand{\rhchXgbOrdRTen}{}\renewcommand{\rhchXgbOrdRTen}{95}  % xgb ord: ladder R_10 (%)
\providecommand{\rhchXgbOrdRTwenty}{}\renewcommand{\rhchXgbOrdRTwenty}{99}  % xgb ord: ladder R_20 (%)
\providecommand{\rhchRfBarLadFullPct}{}\renewcommand{\rhchRfBarLadFullPct}{$-$10.6}  % rf all: ladder k=46 vs 0 %
\providecommand{\rhchRfBarROne}{}\renewcommand{\rhchRfBarROne}{6}  % rf all: ladder R_1 (%)
\providecommand{\rhchRfBarRThree}{}\renewcommand{\rhchRfBarRThree}{42}  % rf all: ladder R_3 (%)
\providecommand{\rhchRfBarRFive}{}\renewcommand{\rhchRfBarRFive}{40}  % rf all: ladder R_5 (%)
\providecommand{\rhchRfBarRTen}{}\renewcommand{\rhchRfBarRTen}{98}  % rf all: ladder R_10 (%)
\providecommand{\rhchRfBarRTwenty}{}\renewcommand{\rhchRfBarRTwenty}{101}  % rf all: ladder R_20 (%)
\providecommand{\rhchRfEventLadFullPct}{}\renewcommand{\rhchRfEventLadFullPct}{$-$44.4}  % rf event: ladder k=46 vs 0 %
\providecommand{\rhchRfEventROne}{}\renewcommand{\rhchRfEventROne}{$-$3}  % rf event: ladder R_1 (%)
\providecommand{\rhchRfEventRThree}{}\renewcommand{\rhchRfEventRThree}{3}  % rf event: ladder R_3 (%)
\providecommand{\rhchRfEventRFive}{}\renewcommand{\rhchRfEventRFive}{6}  % rf event: ladder R_5 (%)
\providecommand{\rhchRfEventRTen}{}\renewcommand{\rhchRfEventRTen}{100}  % rf event: ladder R_10 (%)
\providecommand{\rhchRfEventRTwenty}{}\renewcommand{\rhchRfEventRTwenty}{100}  % rf event: ladder R_20 (%)
\providecommand{\rhchRfOrdLadFullPct}{}\renewcommand{\rhchRfOrdLadFullPct}{$-$4.6}  % rf ord: ladder k=46 vs 0 %
\providecommand{\rhchRfOrdROne}{}\renewcommand{\rhchRfOrdROne}{20}  % rf ord: ladder R_1 (%)
\providecommand{\rhchRfOrdRThree}{}\renewcommand{\rhchRfOrdRThree}{109}  % rf ord: ladder R_3 (%)
\providecommand{\rhchRfOrdRFive}{}\renewcommand{\rhchRfOrdRFive}{97}  % rf ord: ladder R_5 (%)
\providecommand{\rhchRfOrdRTen}{}\renewcommand{\rhchRfOrdRTen}{95}  % rf ord: ladder R_10 (%)
\providecommand{\rhchRfOrdRTwenty}{}\renewcommand{\rhchRfOrdRTwenty}{102}  % rf ord: ladder R_20 (%)
\fi
\begin{table}[htbp]
\centering
\scriptsize
\setlength{\tabcolsep}{1.5pt}
\begin{threeparttable}
\caption{Predictor groups by day type: each group added alone to the HAR benchmark and left out of the full set, on
  FOMC statement days and on ordinary days (\% change in QLIKE)}
\label{tab:rh_channels_abl}
\begin{tabular*}{\linewidth}{@{\extracolsep{\fill}}lrrrrrrrrrrrr@{}}
\toprule
 & \multicolumn{6}{c}{HAR benchmark plus the group, vs HAR} & \multicolumn{6}{c}{Full set without the group, vs full set} \\
 & \multicolumn{6}{c}{\textit{$-$: the group adds accuracy; $+$: it hurts}}
 & \multicolumn{6}{c}{\textit{$+$: the group adds accuracy; $-$: it hurts}} \\
\cmidrule(lr){2-7}\cmidrule(l){8-13}
 & Ridge & Lasso & E.\ net & LGBM & XGB & Forest & Ridge & Lasso & E.\ net & LGBM$^{\dagger}$ & XGB$^{\dagger}$ & Forest$^{\dagger}$ \\
\midrule
\multicolumn{13}{@{}l}{\textit{A. FOMC statement days (\rhchNEventDays{} days, 1{,}053 bars)}} \\
Realized moments & \textbf{\boldmath 4.4} & \textbf{\boldmath 4.7} & \textbf{\boldmath 4.2} & \textbf{\boldmath 6.8} & \textbf{\boldmath 5.5} & 5.3 & $-$1.1 & $-$0.8 & $-$1.2 & $-$0.3 & $-$1.3 & \textbf{\boldmath $-$2.1} \\
Liquidity & \textbf{\boldmath 5.4} & \textbf{\boldmath 6.1} & \textbf{\boldmath 5.3} & 1.5 & 0.4 & 2.0 & $-$0.6 & 0.1 & $-$0.7 & $-$0.5 & $-$0.9 & -- \\
Equal-weighted market & $-$1.8 & $-$1.7 & $-$2.7 & \textbf{\boldmath $-$3.5} & \textbf{\boldmath $-$3.6} & $-$2.4 & 0.0 & 0.1 & 0.0 & 0.1 & 0.0 & -- \\
Value-weighted market & $-$2.5 & $-$2.4 & $-$3.4 & \textbf{\boldmath $-$3.5} & $-$3.3 & $-$0.9 & 0.1 & $-$0.3 & 0.0 & $-$0.1 & $-$0.7 & -- \\
Sentiment & $-$0.1 & $-$1.3 & $-$0.9 & 0.5 & 0.3 & 0.4 & 1.7 & 2.0 & 1.6 & 0.2 & 0.0 & -- \\
Implied volatility & \textbf{\boldmath 2.8} & \textbf{\boldmath 2.9} & \textbf{\boldmath 2.5} & \textbf{\boldmath $-$1.4} & $-$0.7 & \textbf{\boldmath 4.1} & $-$0.2 & $-$0.1 & $-$0.1 & 0.7 & 0.4 & -- \\
Volatility demand & \textbf{\boldmath 1.9} & \textbf{\boldmath 1.9} & \textbf{\boldmath 1.6} & $-$3.5 & $-$1.9 & \textbf{\boldmath 1.3} & 0.0 & \textbf{\boldmath 0.3} & $-$0.1 & \textbf{\boldmath 0.7} & $-$0.7 & -- \\
FOMC & \cellcolor{black!8}\textbf{\boldmath $-$23.4} & \cellcolor{black!8}\textbf{\boldmath $-$43.9} & \cellcolor{black!8}\textbf{\boldmath $-$23.2} & \cellcolor{black!8}\textbf{\boldmath $-$62.6} & \cellcolor{black!8}\textbf{\boldmath $-$64.0} & \cellcolor{black!8}\textbf{\boldmath $-$48.5} & \cellcolor{black!8}\textbf{\boldmath 28.6} & \cellcolor{black!8}\textbf{\boldmath 87.4} & \cellcolor{black!8}\textbf{\boldmath 28.2} & \cellcolor{black!8}\textbf{\boldmath 163.2} & \cellcolor{black!8}\textbf{\boldmath 173.0} & \cellcolor{black!8}-- \\
\addlinespace[2pt]
Full set vs HAR (joint) & \textbf{\boldmath $-$22.4} & \textbf{\boldmath $-$47.2} & \textbf{\boldmath $-$22.6} & \textbf{\boldmath $-$62.3} & \textbf{\boldmath $-$63.9} & \textbf{\boldmath $-$44.8} &  &  &  &  &  &  \\
Reading, by model type & \multicolumn{3}{c}{concentrated} & \multicolumn{3}{c}{concentrated} & & & & & & \\
\addlinespace[4pt]
\multicolumn{13}{@{}l}{\textit{B. Ordinary days (\rhchNOrdDays{} days)}} \\
Realized moments & \cellcolor{black!8}\textbf{\boldmath $-$3.8} & \cellcolor{black!8}\textbf{\boldmath $-$3.3} & \cellcolor{black!8}\textbf{\boldmath $-$3.9} & \cellcolor{black!8}\textbf{\boldmath $-$4.6} & \cellcolor{black!8}\textbf{\boldmath $-$4.9} & \cellcolor{black!8}\textbf{\boldmath $-$4.7} & \cellcolor{black!8}\textbf{\boldmath 2.2} & \cellcolor{black!8}\textbf{\boldmath 1.8} & \cellcolor{black!8}\textbf{\boldmath 2.5} & \cellcolor{black!8}\textbf{\boldmath 1.6} & \cellcolor{black!8}\textbf{\boldmath 1.8} & \cellcolor{black!8}\textbf{\boldmath 1.8} \\
Liquidity & \textbf{\boldmath $-$1.4} & \textbf{\boldmath $-$1.3} & \textbf{\boldmath $-$1.6} & \textbf{\boldmath $-$3.7} & \textbf{\boldmath $-$4.1} & \textbf{\boldmath $-$3.2} & 0.3 & 0.1 & 0.2 & 0.2 & 0.2 & -- \\
Equal-weighted market & \textbf{\boldmath $-$0.7} & $-$0.7 & \textbf{\boldmath $-$0.9} & \textbf{\boldmath $-$2.7} & \textbf{\boldmath $-$3.0} & \textbf{\boldmath $-$2.3} & $-$0.1 & $-$0.1 & $-$0.1 & 0.0 & $-$0.1 & -- \\
Value-weighted market & $-$0.6 & $-$0.3 & $-$0.7 & \textbf{\boldmath $-$3.1} & \textbf{\boldmath $-$3.3} & \textbf{\boldmath $-$2.8} & $-$0.1 & $-$0.2 & $-$0.1 & 0.1 & 0.2 & -- \\
Sentiment & 0.3 & 0.4 & 0.2 & \textbf{\boldmath $-$1.4} & \textbf{\boldmath $-$2.0} & \textbf{\boldmath $-$1.1} & $-$0.1 & $-$0.1 & 0.0 & \textbf{\boldmath 0.2} & 0.2 & -- \\
Implied volatility & $-$0.2 & \textbf{\boldmath $-$0.9} & $-$0.3 & \textbf{\boldmath $-$1.3} & \textbf{\boldmath $-$1.7} & \textbf{\boldmath $-$1.0} & 0.2 & 0.5 & 0.0 & 0.1 & 0.0 & -- \\
Volatility demand & 0.3 & 0.3 & 0.2 & \textbf{\boldmath $-$0.9} & \textbf{\boldmath $-$1.1} & \textbf{\boldmath $-$1.0} & $-$0.1 & $-$0.1 & $-$0.1 & $-$0.1 & 0.1 & -- \\
FOMC & \textbf{\boldmath 0.4} & \textbf{\boldmath 0.7} & \textbf{\boldmath 0.4} & $-$0.2 & $-$0.1 & $-$0.3 & \textbf{\boldmath $-$0.4} & \textbf{\boldmath $-$0.7} & \textbf{\boldmath $-$0.4} & 0.0 & \textbf{\boldmath $-$0.4} & -- \\
\addlinespace[2pt]
Full set vs HAR (joint) & \textbf{\boldmath $-$3.2} & \textbf{\boldmath $-$2.4} & \textbf{\boldmath $-$4.0} & \textbf{\boldmath $-$6.0} & \textbf{\boldmath $-$6.1} & \textbf{\boldmath $-$5.3} &  &  &  &  &  &  \\
Reading, by model type & \multicolumn{3}{c}{concentrated} & \multicolumn{3}{c}{concentrated} & & & & & & \\
\bottomrule
\end{tabular*}
\begin{tablenotes}[flushleft]
\footnotesize
\item \textit{Notes:} \textit{Sign:} every cell is $100\,(L_{\mathrm{row}}/L_{\mathrm{ref}} - 1)$, the
percentage change in mean QLIKE of the row's forecast against its reference on the bars of the day type; the
Diebold--Mariano statistic has the same sign. Left block, the group added to the HAR benchmark set (reference: the HAR
benchmark): negative = the loss falls when the group is added, i.e.\ the group adds accuracy; positive = the group
hurts. Right block, the group removed from the full set (reference: the full set): positive = the loss rises when the
group is dropped, i.e.\ the group adds accuracy; negative = the group hurts. Joint: the full set against the HAR
benchmark on the same days; negative = the full set is more accurate. Bold: the Diebold--Mariano test rejects equal accuracy
at 5\% (two-sided; daily average loss, Newey--West lag \rhchHacLag). HAR benchmark: the OLS on the HAR set for
the linear models, each tree's own HAR-set fit. Shaded: in each class's column, the group with the largest add-alone gain on the day type, in
both blocks. Readings by model type (fixed before scoring, tested in this order): concentrated if one group's add-alone
gain is at least 75\% of the joint gain and its removal raises the loss significantly; redundant if the add-alone gains
sum to more than 1.5 times the joint gain and at most one removal raises the loss significantly; otherwise mixed; the
majority of ridge, lasso and elastic net, and of LightGBM and XGBoost with the forest where its arms exist.
Table~\ref{tab:rh_channels_abl_full} (Appendix~\ref{app:rh_channels}) adds the per-class readings, the add-alone gains
summed over the groups against the joint gain and the counts of significant leave-outs. $^{\dagger}$ On a class: its removals other than the FOMC group are on the unpatched design (the FOMC proximity series before the fix), compared with the unpatched full set; they were not rerun on the patched design. Arms without a FOMC column are the same forecasts on both designs. The forest is run for the
added groups and for leaving out the realized moments only (--: not run).
\end{tablenotes}
\end{threeparttable}
\end{table}

% AUTO-GENERATED by writeup/make_rehaul_channels.py -- do not edit; re-run the script.
\ifdefined\rhchNBars\else% AUTO-GENERATED by writeup/make_rehaul_channels.py -- do not edit; re-run the script.
% \rhch... macros: the two channels of the pooled gain (results/channels_2026-10-26/RULE.txt items 1-3):
% FOMC statement days (EVENT) and ordinary days, Table 1's scored bars, the FOMC-patched design of record;
% DM (HLN) on daily average loss, lag 8.  Names in the generator's docstring.
% \rhchX... = ADDITIONAL, NOT PRE-REGISTERED (coordinator, 2026-10-26, after the first scoring): the full set
% against HAR + the futures' realized moments (FullVsRm) and HAR + realized moments against HAR (RmVsHar).
\providecommand{\rhPending}{\textit{pending}}
\providecommand{\rhchNBars}{}\renewcommand{\rhchNBars}{33{,}527}  % scored bars (Table 1)
\providecommand{\rhchNDays}{}\renewcommand{\rhchNDays}{2{,}579}  % scored days
\providecommand{\rhchNEventDays}{}\renewcommand{\rhchNEventDays}{81}  % EVENT days: scheduled FOMC statement days of the span
\providecommand{\rhchNOrdDays}{}\renewcommand{\rhchNOrdDays}{2{,}498}  % ORDINARY days
\providecommand{\rhchNEventBars}{}\renewcommand{\rhchNEventBars}{1{,}053}  % bars on EVENT days
\providecommand{\rhchEventBarPct}{}\renewcommand{\rhchEventBarPct}{3.1}  % EVENT-day bars, % of the scored bars
\providecommand{\rhchHacLag}{}\renewcommand{\rhchHacLag}{8}  % Newey-West lag of every DM test
\providecommand{\rhchAgreePp}{}\renewcommand{\rhchAgreePp}{2}  % item 1: the corrections agree within this many percentage points
\providecommand{\rhchLadderDesign}{}\renewcommand{\rhchLadderDesign}{patched}  % design of the ladder
\providecommand{\rhchRidgeTabOnePct}{}\renewcommand{\rhchRidgeTabOnePct}{$-$6.0}  % ridge: full vs HAR, every bar (Table 1): % QLIKE change
\providecommand{\rhchRidgeTabOneDm}{}\renewcommand{\rhchRidgeTabOneDm}{$-$6.28}  % ridge: full vs HAR, every bar (Table 1): DM
\providecommand{\rhchRidgeTabOneP}{}\renewcommand{\rhchRidgeTabOneP}{\ensuremath{<}0.001}  % ridge: full vs HAR, every bar (Table 1): DM p
\providecommand{\rhchRidgeTabOneMsePct}{}\renewcommand{\rhchRidgeTabOneMsePct}{100.6}  % ridge: full vs HAR, every bar (Table 1): % MSE change
\providecommand{\rhchRidgeTabOneMseP}{}\renewcommand{\rhchRidgeTabOneMseP}{0.261}  % ridge: full vs HAR, every bar (Table 1): DM p on MSE
\providecommand{\rhchRidgeCorrAPct}{}\renewcommand{\rhchRidgeCorrAPct}{$-$3.0}  % ridge: full vs HAR+FOMC, every bar: % QLIKE change
\providecommand{\rhchRidgeCorrADm}{}\renewcommand{\rhchRidgeCorrADm}{$-$3.64}  % ridge: full vs HAR+FOMC, every bar: DM
\providecommand{\rhchRidgeCorrAP}{}\renewcommand{\rhchRidgeCorrAP}{\ensuremath{<}0.001}  % ridge: full vs HAR+FOMC, every bar: DM p
\providecommand{\rhchRidgeCorrAMsePct}{}\renewcommand{\rhchRidgeCorrAMsePct}{100.8}  % ridge: full vs HAR+FOMC, every bar: % MSE change
\providecommand{\rhchRidgeCorrAMseP}{}\renewcommand{\rhchRidgeCorrAMseP}{0.258}  % ridge: full vs HAR+FOMC, every bar: DM p on MSE
\providecommand{\rhchRidgeCorrBPct}{}\renewcommand{\rhchRidgeCorrBPct}{$-$3.2}  % ridge: full without FOMC vs HAR, every bar: % QLIKE change
\providecommand{\rhchRidgeCorrBDm}{}\renewcommand{\rhchRidgeCorrBDm}{$-$3.74}  % ridge: full without FOMC vs HAR, every bar: DM
\providecommand{\rhchRidgeCorrBP}{}\renewcommand{\rhchRidgeCorrBP}{\ensuremath{<}0.001}  % ridge: full without FOMC vs HAR, every bar: DM p
\providecommand{\rhchRidgeCorrBMsePct}{}\renewcommand{\rhchRidgeCorrBMsePct}{102.3}  % ridge: full without FOMC vs HAR, every bar: % MSE change
\providecommand{\rhchRidgeCorrBMseP}{}\renewcommand{\rhchRidgeCorrBMseP}{0.258}  % ridge: full without FOMC vs HAR, every bar: DM p on MSE
\providecommand{\rhchRidgeCorrCPct}{}\renewcommand{\rhchRidgeCorrCPct}{$-$3.2}  % ridge: full vs HAR, ORDINARY days: % QLIKE change
\providecommand{\rhchRidgeCorrCDm}{}\renewcommand{\rhchRidgeCorrCDm}{$-$3.51}  % ridge: full vs HAR, ORDINARY days: DM
\providecommand{\rhchRidgeCorrCP}{}\renewcommand{\rhchRidgeCorrCP}{\ensuremath{<}0.001}  % ridge: full vs HAR, ORDINARY days: DM p
\providecommand{\rhchRidgeCorrCMsePct}{}\renewcommand{\rhchRidgeCorrCMsePct}{110.2}  % ridge: full vs HAR, ORDINARY days: % MSE change
\providecommand{\rhchRidgeCorrCMseP}{}\renewcommand{\rhchRidgeCorrCMseP}{0.258}  % ridge: full vs HAR, ORDINARY days: DM p on MSE
\providecommand{\rhchRidgeCorrAgree}{}\renewcommand{\rhchRidgeCorrAgree}{agree}  % ridge: the three corrections within 2.0 pp?
\providecommand{\rhchRidgeCorrSpread}{}\renewcommand{\rhchRidgeCorrSpread}{0.2}  % ridge: max - min of the three (pp)
\providecommand{\rhchLassoTabOnePct}{}\renewcommand{\rhchLassoTabOnePct}{$-$9.0}  % lasso: full vs HAR, every bar (Table 1): % QLIKE change
\providecommand{\rhchLassoTabOneDm}{}\renewcommand{\rhchLassoTabOneDm}{$-$6.34}  % lasso: full vs HAR, every bar (Table 1): DM
\providecommand{\rhchLassoTabOneP}{}\renewcommand{\rhchLassoTabOneP}{\ensuremath{<}0.001}  % lasso: full vs HAR, every bar (Table 1): DM p
\providecommand{\rhchLassoTabOneMsePct}{}\renewcommand{\rhchLassoTabOneMsePct}{114.7}  % lasso: full vs HAR, every bar (Table 1): % MSE change
\providecommand{\rhchLassoTabOneMseP}{}\renewcommand{\rhchLassoTabOneMseP}{0.247}  % lasso: full vs HAR, every bar (Table 1): DM p on MSE
\providecommand{\rhchLassoCorrAPct}{}\renewcommand{\rhchLassoCorrAPct}{$-$3.4}  % lasso: full vs HAR+FOMC, every bar: % QLIKE change
\providecommand{\rhchLassoCorrADm}{}\renewcommand{\rhchLassoCorrADm}{$-$3.20}  % lasso: full vs HAR+FOMC, every bar: DM
\providecommand{\rhchLassoCorrAP}{}\renewcommand{\rhchLassoCorrAP}{0.001}  % lasso: full vs HAR+FOMC, every bar: DM p
\providecommand{\rhchLassoCorrAMsePct}{}\renewcommand{\rhchLassoCorrAMsePct}{114.4}  % lasso: full vs HAR+FOMC, every bar: % MSE change
\providecommand{\rhchLassoCorrAMseP}{}\renewcommand{\rhchLassoCorrAMseP}{0.243}  % lasso: full vs HAR+FOMC, every bar: DM p on MSE
\providecommand{\rhchLassoCorrBPct}{}\renewcommand{\rhchLassoCorrBPct}{$-$2.8}  % lasso: full without FOMC vs HAR, every bar: % QLIKE change
\providecommand{\rhchLassoCorrBDm}{}\renewcommand{\rhchLassoCorrBDm}{$-$2.75}  % lasso: full without FOMC vs HAR, every bar: DM
\providecommand{\rhchLassoCorrBP}{}\renewcommand{\rhchLassoCorrBP}{0.006}  % lasso: full without FOMC vs HAR, every bar: DM p
\providecommand{\rhchLassoCorrBMsePct}{}\renewcommand{\rhchLassoCorrBMsePct}{115.3}  % lasso: full without FOMC vs HAR, every bar: % MSE change
\providecommand{\rhchLassoCorrBMseP}{}\renewcommand{\rhchLassoCorrBMseP}{0.245}  % lasso: full without FOMC vs HAR, every bar: DM p on MSE
\providecommand{\rhchLassoCorrCPct}{}\renewcommand{\rhchLassoCorrCPct}{$-$2.4}  % lasso: full vs HAR, ORDINARY days: % QLIKE change
\providecommand{\rhchLassoCorrCDm}{}\renewcommand{\rhchLassoCorrCDm}{$-$2.08}  % lasso: full vs HAR, ORDINARY days: DM
\providecommand{\rhchLassoCorrCP}{}\renewcommand{\rhchLassoCorrCP}{0.037}  % lasso: full vs HAR, ORDINARY days: DM p
\providecommand{\rhchLassoCorrCMsePct}{}\renewcommand{\rhchLassoCorrCMsePct}{126.0}  % lasso: full vs HAR, ORDINARY days: % MSE change
\providecommand{\rhchLassoCorrCMseP}{}\renewcommand{\rhchLassoCorrCMseP}{0.243}  % lasso: full vs HAR, ORDINARY days: DM p on MSE
\providecommand{\rhchLassoCorrAgree}{}\renewcommand{\rhchLassoCorrAgree}{agree}  % lasso: the three corrections within 2.0 pp?
\providecommand{\rhchLassoCorrSpread}{}\renewcommand{\rhchLassoCorrSpread}{1.0}  % lasso: max - min of the three (pp)
\providecommand{\rhchEnetTabOnePct}{}\renewcommand{\rhchEnetTabOnePct}{$-$6.7}  % enet: full vs HAR, every bar (Table 1): % QLIKE change
\providecommand{\rhchEnetTabOneDm}{}\renewcommand{\rhchEnetTabOneDm}{$-$7.36}  % enet: full vs HAR, every bar (Table 1): DM
\providecommand{\rhchEnetTabOneP}{}\renewcommand{\rhchEnetTabOneP}{\ensuremath{<}0.001}  % enet: full vs HAR, every bar (Table 1): DM p
\providecommand{\rhchEnetTabOneMsePct}{}\renewcommand{\rhchEnetTabOneMsePct}{117.7}  % enet: full vs HAR, every bar (Table 1): % MSE change
\providecommand{\rhchEnetTabOneMseP}{}\renewcommand{\rhchEnetTabOneMseP}{0.263}  % enet: full vs HAR, every bar (Table 1): DM p on MSE
\providecommand{\rhchEnetCorrAPct}{}\renewcommand{\rhchEnetCorrAPct}{$-$3.7}  % enet: full vs HAR+FOMC, every bar: % QLIKE change
\providecommand{\rhchEnetCorrADm}{}\renewcommand{\rhchEnetCorrADm}{$-$4.89}  % enet: full vs HAR+FOMC, every bar: DM
\providecommand{\rhchEnetCorrAP}{}\renewcommand{\rhchEnetCorrAP}{\ensuremath{<}0.001}  % enet: full vs HAR+FOMC, every bar: DM p
\providecommand{\rhchEnetCorrAMsePct}{}\renewcommand{\rhchEnetCorrAMsePct}{117.8}  % enet: full vs HAR+FOMC, every bar: % MSE change
\providecommand{\rhchEnetCorrAMseP}{}\renewcommand{\rhchEnetCorrAMseP}{0.261}  % enet: full vs HAR+FOMC, every bar: DM p on MSE
\providecommand{\rhchEnetCorrBPct}{}\renewcommand{\rhchEnetCorrBPct}{$-$3.8}  % enet: full without FOMC vs HAR, every bar: % QLIKE change
\providecommand{\rhchEnetCorrBDm}{}\renewcommand{\rhchEnetCorrBDm}{$-$5.00}  % enet: full without FOMC vs HAR, every bar: DM
\providecommand{\rhchEnetCorrBP}{}\renewcommand{\rhchEnetCorrBP}{\ensuremath{<}0.001}  % enet: full without FOMC vs HAR, every bar: DM p
\providecommand{\rhchEnetCorrBMsePct}{}\renewcommand{\rhchEnetCorrBMsePct}{117.8}  % enet: full without FOMC vs HAR, every bar: % MSE change
\providecommand{\rhchEnetCorrBMseP}{}\renewcommand{\rhchEnetCorrBMseP}{0.259}  % enet: full without FOMC vs HAR, every bar: DM p on MSE
\providecommand{\rhchEnetCorrCPct}{}\renewcommand{\rhchEnetCorrCPct}{$-$4.0}  % enet: full vs HAR, ORDINARY days: % QLIKE change
\providecommand{\rhchEnetCorrCDm}{}\renewcommand{\rhchEnetCorrCDm}{$-$4.73}  % enet: full vs HAR, ORDINARY days: DM
\providecommand{\rhchEnetCorrCP}{}\renewcommand{\rhchEnetCorrCP}{\ensuremath{<}0.001}  % enet: full vs HAR, ORDINARY days: DM p
\providecommand{\rhchEnetCorrCMsePct}{}\renewcommand{\rhchEnetCorrCMsePct}{128.9}  % enet: full vs HAR, ORDINARY days: % MSE change
\providecommand{\rhchEnetCorrCMseP}{}\renewcommand{\rhchEnetCorrCMseP}{0.260}  % enet: full vs HAR, ORDINARY days: DM p on MSE
\providecommand{\rhchEnetCorrAgree}{}\renewcommand{\rhchEnetCorrAgree}{agree}  % enet: the three corrections within 2.0 pp?
\providecommand{\rhchEnetCorrSpread}{}\renewcommand{\rhchEnetCorrSpread}{0.2}  % enet: max - min of the three (pp)
\providecommand{\rhchLgbmTabOnePct}{}\renewcommand{\rhchLgbmTabOnePct}{$-$14.1}  % lgbm: full vs HAR, every bar (Table 1): % QLIKE change
\providecommand{\rhchLgbmTabOneDm}{}\renewcommand{\rhchLgbmTabOneDm}{$-$10.26}  % lgbm: full vs HAR, every bar (Table 1): DM
\providecommand{\rhchLgbmTabOneP}{}\renewcommand{\rhchLgbmTabOneP}{\ensuremath{<}0.001}  % lgbm: full vs HAR, every bar (Table 1): DM p
\providecommand{\rhchLgbmTabOneMsePct}{}\renewcommand{\rhchLgbmTabOneMsePct}{20.0}  % lgbm: full vs HAR, every bar (Table 1): % MSE change
\providecommand{\rhchLgbmTabOneMseP}{}\renewcommand{\rhchLgbmTabOneMseP}{0.301}  % lgbm: full vs HAR, every bar (Table 1): DM p on MSE
\providecommand{\rhchLgbmCorrAPct}{}\renewcommand{\rhchLgbmCorrAPct}{$-$5.4}  % lgbm: full vs HAR+FOMC, every bar: % QLIKE change
\providecommand{\rhchLgbmCorrADm}{}\renewcommand{\rhchLgbmCorrADm}{$-$11.63}  % lgbm: full vs HAR+FOMC, every bar: DM
\providecommand{\rhchLgbmCorrAP}{}\renewcommand{\rhchLgbmCorrAP}{\ensuremath{<}0.001}  % lgbm: full vs HAR+FOMC, every bar: DM p
\providecommand{\rhchLgbmCorrAMsePct}{}\renewcommand{\rhchLgbmCorrAMsePct}{19.4}  % lgbm: full vs HAR+FOMC, every bar: % MSE change
\providecommand{\rhchLgbmCorrAMseP}{}\renewcommand{\rhchLgbmCorrAMseP}{0.276}  % lgbm: full vs HAR+FOMC, every bar: DM p on MSE
\providecommand{\rhchLgbmCorrBPct}{}\renewcommand{\rhchLgbmCorrBPct}{$-$5.3}  % lgbm: full without FOMC vs HAR, every bar: % QLIKE change
\providecommand{\rhchLgbmCorrBDm}{}\renewcommand{\rhchLgbmCorrBDm}{$-$9.62}  % lgbm: full without FOMC vs HAR, every bar: DM
\providecommand{\rhchLgbmCorrBP}{}\renewcommand{\rhchLgbmCorrBP}{\ensuremath{<}0.001}  % lgbm: full without FOMC vs HAR, every bar: DM p
\providecommand{\rhchLgbmCorrBMsePct}{}\renewcommand{\rhchLgbmCorrBMsePct}{21.6}  % lgbm: full without FOMC vs HAR, every bar: % MSE change
\providecommand{\rhchLgbmCorrBMseP}{}\renewcommand{\rhchLgbmCorrBMseP}{0.269}  % lgbm: full without FOMC vs HAR, every bar: DM p on MSE
\providecommand{\rhchLgbmCorrCPct}{}\renewcommand{\rhchLgbmCorrCPct}{$-$6.0}  % lgbm: full vs HAR, ORDINARY days: % QLIKE change
\providecommand{\rhchLgbmCorrCDm}{}\renewcommand{\rhchLgbmCorrCDm}{$-$10.70}  % lgbm: full vs HAR, ORDINARY days: DM
\providecommand{\rhchLgbmCorrCP}{}\renewcommand{\rhchLgbmCorrCP}{\ensuremath{<}0.001}  % lgbm: full vs HAR, ORDINARY days: DM p
\providecommand{\rhchLgbmCorrCMsePct}{}\renewcommand{\rhchLgbmCorrCMsePct}{22.8}  % lgbm: full vs HAR, ORDINARY days: % MSE change
\providecommand{\rhchLgbmCorrCMseP}{}\renewcommand{\rhchLgbmCorrCMseP}{0.275}  % lgbm: full vs HAR, ORDINARY days: DM p on MSE
\providecommand{\rhchLgbmCorrAgree}{}\renewcommand{\rhchLgbmCorrAgree}{agree}  % lgbm: the three corrections within 2.0 pp?
\providecommand{\rhchLgbmCorrSpread}{}\renewcommand{\rhchLgbmCorrSpread}{0.8}  % lgbm: max - min of the three (pp)
\providecommand{\rhchXgbTabOnePct}{}\renewcommand{\rhchXgbTabOnePct}{$-$14.6}  % xgb: full vs HAR, every bar (Table 1): % QLIKE change
\providecommand{\rhchXgbTabOneDm}{}\renewcommand{\rhchXgbTabOneDm}{$-$9.94}  % xgb: full vs HAR, every bar (Table 1): DM
\providecommand{\rhchXgbTabOneP}{}\renewcommand{\rhchXgbTabOneP}{\ensuremath{<}0.001}  % xgb: full vs HAR, every bar (Table 1): DM p
\providecommand{\rhchXgbTabOneMsePct}{}\renewcommand{\rhchXgbTabOneMsePct}{16.3}  % xgb: full vs HAR, every bar (Table 1): % MSE change
\providecommand{\rhchXgbTabOneMseP}{}\renewcommand{\rhchXgbTabOneMseP}{0.301}  % xgb: full vs HAR, every bar (Table 1): DM p on MSE
\providecommand{\rhchXgbCorrAPct}{}\renewcommand{\rhchXgbCorrAPct}{$-$5.7}  % xgb: full vs HAR+FOMC, every bar: % QLIKE change
\providecommand{\rhchXgbCorrADm}{}\renewcommand{\rhchXgbCorrADm}{$-$12.60}  % xgb: full vs HAR+FOMC, every bar: DM
\providecommand{\rhchXgbCorrAP}{}\renewcommand{\rhchXgbCorrAP}{\ensuremath{<}0.001}  % xgb: full vs HAR+FOMC, every bar: DM p
\providecommand{\rhchXgbCorrAMsePct}{}\renewcommand{\rhchXgbCorrAMsePct}{15.4}  % xgb: full vs HAR+FOMC, every bar: % MSE change
\providecommand{\rhchXgbCorrAMseP}{}\renewcommand{\rhchXgbCorrAMseP}{0.282}  % xgb: full vs HAR+FOMC, every bar: DM p on MSE
\providecommand{\rhchXgbCorrBPct}{}\renewcommand{\rhchXgbCorrBPct}{$-$5.8}  % xgb: full without FOMC vs HAR, every bar: % QLIKE change
\providecommand{\rhchXgbCorrBDm}{}\renewcommand{\rhchXgbCorrBDm}{$-$10.81}  % xgb: full without FOMC vs HAR, every bar: DM
\providecommand{\rhchXgbCorrBP}{}\renewcommand{\rhchXgbCorrBP}{\ensuremath{<}0.001}  % xgb: full without FOMC vs HAR, every bar: DM p
\providecommand{\rhchXgbCorrBMsePct}{}\renewcommand{\rhchXgbCorrBMsePct}{17.9}  % xgb: full without FOMC vs HAR, every bar: % MSE change
\providecommand{\rhchXgbCorrBMseP}{}\renewcommand{\rhchXgbCorrBMseP}{0.268}  % xgb: full without FOMC vs HAR, every bar: DM p on MSE
\providecommand{\rhchXgbCorrCPct}{}\renewcommand{\rhchXgbCorrCPct}{$-$6.1}  % xgb: full vs HAR, ORDINARY days: % QLIKE change
\providecommand{\rhchXgbCorrCDm}{}\renewcommand{\rhchXgbCorrCDm}{$-$10.42}  % xgb: full vs HAR, ORDINARY days: DM
\providecommand{\rhchXgbCorrCP}{}\renewcommand{\rhchXgbCorrCP}{\ensuremath{<}0.001}  % xgb: full vs HAR, ORDINARY days: DM p
\providecommand{\rhchXgbCorrCMsePct}{}\renewcommand{\rhchXgbCorrCMsePct}{18.9}  % xgb: full vs HAR, ORDINARY days: % MSE change
\providecommand{\rhchXgbCorrCMseP}{}\renewcommand{\rhchXgbCorrCMseP}{0.268}  % xgb: full vs HAR, ORDINARY days: DM p on MSE
\providecommand{\rhchXgbCorrAgree}{}\renewcommand{\rhchXgbCorrAgree}{agree}  % xgb: the three corrections within 2.0 pp?
\providecommand{\rhchXgbCorrSpread}{}\renewcommand{\rhchXgbCorrSpread}{0.4}  % xgb: max - min of the three (pp)
\providecommand{\rhchRfTabOnePct}{}\renewcommand{\rhchRfTabOnePct}{$-$11.3}  % rf: full vs HAR, every bar (Table 1): % QLIKE change
\providecommand{\rhchRfTabOneDm}{}\renewcommand{\rhchRfTabOneDm}{$-$9.74}  % rf: full vs HAR, every bar (Table 1): DM
\providecommand{\rhchRfTabOneP}{}\renewcommand{\rhchRfTabOneP}{\ensuremath{<}0.001}  % rf: full vs HAR, every bar (Table 1): DM p
\providecommand{\rhchRfTabOneMsePct}{}\renewcommand{\rhchRfTabOneMsePct}{10.9}  % rf: full vs HAR, every bar (Table 1): % MSE change
\providecommand{\rhchRfTabOneMseP}{}\renewcommand{\rhchRfTabOneMseP}{0.271}  % rf: full vs HAR, every bar (Table 1): DM p on MSE
\providecommand{\rhchRfCorrAPct}{}\renewcommand{\rhchRfCorrAPct}{$-$4.0}  % rf: full vs HAR+FOMC, every bar: % QLIKE change
\providecommand{\rhchRfCorrADm}{}\renewcommand{\rhchRfCorrADm}{$-$9.06}  % rf: full vs HAR+FOMC, every bar: DM
\providecommand{\rhchRfCorrAP}{}\renewcommand{\rhchRfCorrAP}{\ensuremath{<}0.001}  % rf: full vs HAR+FOMC, every bar: DM p
\providecommand{\rhchRfCorrAMsePct}{}\renewcommand{\rhchRfCorrAMsePct}{11.3}  % rf: full vs HAR+FOMC, every bar: % MSE change
\providecommand{\rhchRfCorrAMseP}{}\renewcommand{\rhchRfCorrAMseP}{0.255}  % rf: full vs HAR+FOMC, every bar: DM p on MSE
\providecommand{\rhchRfCorrBPct}{}\renewcommand{\rhchRfCorrBPct}{--}  % rf: full without FOMC vs HAR, every bar: not run
\providecommand{\rhchRfCorrBDm}{}\renewcommand{\rhchRfCorrBDm}{--}  % rf: full without FOMC vs HAR, every bar: not run
\providecommand{\rhchRfCorrBP}{}\renewcommand{\rhchRfCorrBP}{--}  % rf: full without FOMC vs HAR, every bar: not run
\providecommand{\rhchRfCorrBMsePct}{}\renewcommand{\rhchRfCorrBMsePct}{--}  % rf: full without FOMC vs HAR, every bar: not run
\providecommand{\rhchRfCorrBMseP}{}\renewcommand{\rhchRfCorrBMseP}{--}  % rf: full without FOMC vs HAR, every bar: not run
\providecommand{\rhchRfCorrCPct}{}\renewcommand{\rhchRfCorrCPct}{$-$5.3}  % rf: full vs HAR, ORDINARY days: % QLIKE change
\providecommand{\rhchRfCorrCDm}{}\renewcommand{\rhchRfCorrCDm}{$-$11.16}  % rf: full vs HAR, ORDINARY days: DM
\providecommand{\rhchRfCorrCP}{}\renewcommand{\rhchRfCorrCP}{\ensuremath{<}0.001}  % rf: full vs HAR, ORDINARY days: DM p
\providecommand{\rhchRfCorrCMsePct}{}\renewcommand{\rhchRfCorrCMsePct}{12.4}  % rf: full vs HAR, ORDINARY days: % MSE change
\providecommand{\rhchRfCorrCMseP}{}\renewcommand{\rhchRfCorrCMseP}{0.247}  % rf: full vs HAR, ORDINARY days: DM p on MSE
\providecommand{\rhchRfCorrAgree}{}\renewcommand{\rhchRfCorrAgree}{--}  % rf: the three corrections within 2.0 pp?
\providecommand{\rhchRfCorrSpread}{}\renewcommand{\rhchRfCorrSpread}{1.3}  % rf: max - min of the three (pp)
\providecommand{\rhchNCorrAgree}{}\renewcommand{\rhchNCorrAgree}{5}  % classes whose corrections agree
\providecommand{\rhchNCorrDisagree}{}\renewcommand{\rhchNCorrDisagree}{0}  % classes whose corrections disagree
\providecommand{\rhchNCorrPending}{}\renewcommand{\rhchNCorrPending}{0}  % classes with a correction pending
\providecommand{\rhchRidgeShareEvent}{}\renewcommand{\rhchRidgeShareEvent}{54}  % ridge: % of the every-bar gain on EVENT days
\providecommand{\rhchRidgeShareEventPm}{}\renewcommand{\rhchRidgeShareEventPm}{56}  % ridge: % of the gain on EVENT days 14:30-16:00
\providecommand{\rhchRidgeEventPeakBar}{}\renewcommand{\rhchRidgeEventPeakBar}{14:30}  % ridge: event: bar end with the most negative mean diff
\providecommand{\rhchRidgeEventPeakDiff}{}\renewcommand{\rhchRidgeEventPeakDiff}{$-$1583.6}  % ridge: event: that mean diff (full - HAR) x 10^3
\providecommand{\rhchRidgeOrdPeakBar}{}\renewcommand{\rhchRidgeOrdPeakBar}{14:30}  % ridge: ord: bar end with the most negative mean diff
\providecommand{\rhchRidgeOrdPeakDiff}{}\renewcommand{\rhchRidgeOrdPeakDiff}{$-$15.9}  % ridge: ord: that mean diff (full - HAR) x 10^3
\providecommand{\rhchLassoShareEvent}{}\renewcommand{\rhchLassoShareEvent}{77}  % lasso: % of the every-bar gain on EVENT days
\providecommand{\rhchLassoShareEventPm}{}\renewcommand{\rhchLassoShareEventPm}{78}  % lasso: % of the gain on EVENT days 14:30-16:00
\providecommand{\rhchLassoEventPeakBar}{}\renewcommand{\rhchLassoEventPeakBar}{14:30}  % lasso: event: bar end with the most negative mean diff
\providecommand{\rhchLassoEventPeakDiff}{}\renewcommand{\rhchLassoEventPeakDiff}{$-$4302.3}  % lasso: event: that mean diff (full - HAR) x 10^3
\providecommand{\rhchLassoOrdPeakBar}{}\renewcommand{\rhchLassoOrdPeakBar}{14:30}  % lasso: ord: bar end with the most negative mean diff
\providecommand{\rhchLassoOrdPeakDiff}{}\renewcommand{\rhchLassoOrdPeakDiff}{$-$16.4}  % lasso: ord: that mean diff (full - HAR) x 10^3
\providecommand{\rhchEnetShareEvent}{}\renewcommand{\rhchEnetShareEvent}{49}  % enet: % of the every-bar gain on EVENT days
\providecommand{\rhchEnetShareEventPm}{}\renewcommand{\rhchEnetShareEventPm}{51}  % enet: % of the gain on EVENT days 14:30-16:00
\providecommand{\rhchEnetEventPeakBar}{}\renewcommand{\rhchEnetEventPeakBar}{14:30}  % enet: event: bar end with the most negative mean diff
\providecommand{\rhchEnetEventPeakDiff}{}\renewcommand{\rhchEnetEventPeakDiff}{$-$1610.3}  % enet: event: that mean diff (full - HAR) x 10^3
\providecommand{\rhchEnetOrdPeakBar}{}\renewcommand{\rhchEnetOrdPeakBar}{15:00}  % enet: ord: bar end with the most negative mean diff
\providecommand{\rhchEnetOrdPeakDiff}{}\renewcommand{\rhchEnetOrdPeakDiff}{$-$15.8}  % enet: ord: that mean diff (full - HAR) x 10^3
\providecommand{\rhchLgbmShareEvent}{}\renewcommand{\rhchLgbmShareEvent}{63}  % lgbm: % of the every-bar gain on EVENT days
\providecommand{\rhchLgbmShareEventPm}{}\renewcommand{\rhchLgbmShareEventPm}{63}  % lgbm: % of the gain on EVENT days 14:30-16:00
\providecommand{\rhchLgbmEventPeakBar}{}\renewcommand{\rhchLgbmEventPeakBar}{14:30}  % lgbm: event: bar end with the most negative mean diff
\providecommand{\rhchLgbmEventPeakDiff}{}\renewcommand{\rhchLgbmEventPeakDiff}{$-$6012.1}  % lgbm: event: that mean diff (full - HAR) x 10^3
\providecommand{\rhchLgbmOrdPeakBar}{}\renewcommand{\rhchLgbmOrdPeakBar}{12:30}  % lgbm: ord: bar end with the most negative mean diff
\providecommand{\rhchLgbmOrdPeakDiff}{}\renewcommand{\rhchLgbmOrdPeakDiff}{$-$15.9}  % lgbm: ord: that mean diff (full - HAR) x 10^3
\providecommand{\rhchXgbShareEvent}{}\renewcommand{\rhchXgbShareEvent}{64}  % xgb: % of the every-bar gain on EVENT days
\providecommand{\rhchXgbShareEventPm}{}\renewcommand{\rhchXgbShareEventPm}{63}  % xgb: % of the gain on EVENT days 14:30-16:00
\providecommand{\rhchXgbEventPeakBar}{}\renewcommand{\rhchXgbEventPeakBar}{14:30}  % xgb: event: bar end with the most negative mean diff
\providecommand{\rhchXgbEventPeakDiff}{}\renewcommand{\rhchXgbEventPeakDiff}{$-$6274.7}  % xgb: event: that mean diff (full - HAR) x 10^3
\providecommand{\rhchXgbOrdPeakBar}{}\renewcommand{\rhchXgbOrdPeakBar}{12:30}  % xgb: ord: bar end with the most negative mean diff
\providecommand{\rhchXgbOrdPeakDiff}{}\renewcommand{\rhchXgbOrdPeakDiff}{$-$16.3}  % xgb: ord: that mean diff (full - HAR) x 10^3
\providecommand{\rhchRfShareEvent}{}\renewcommand{\rhchRfShareEvent}{60}  % rf: % of the every-bar gain on EVENT days
\providecommand{\rhchRfShareEventPm}{}\renewcommand{\rhchRfShareEventPm}{59}  % rf: % of the gain on EVENT days 14:30-16:00
\providecommand{\rhchRfEventPeakBar}{}\renewcommand{\rhchRfEventPeakBar}{14:30}  % rf: event: bar end with the most negative mean diff
\providecommand{\rhchRfEventPeakDiff}{}\renewcommand{\rhchRfEventPeakDiff}{$-$4773.7}  % rf: event: that mean diff (full - HAR) x 10^3
\providecommand{\rhchRfOrdPeakBar}{}\renewcommand{\rhchRfOrdPeakBar}{12:30}  % rf: ord: bar end with the most negative mean diff
\providecommand{\rhchRfOrdPeakDiff}{}\renewcommand{\rhchRfOrdPeakDiff}{$-$15.2}  % rf: ord: that mean diff (full - HAR) x 10^3
\providecommand{\rhchEventPmBarPct}{}\renewcommand{\rhchEventPmBarPct}{1.0}  % EVENT 14:30-16:00 bars, % of all bars
\providecommand{\rhchLinShareEventMin}{}\renewcommand{\rhchLinShareEventMin}{49}  % Lin: lowest share_event_pct
\providecommand{\rhchLinShareEventMax}{}\renewcommand{\rhchLinShareEventMax}{77}  % Lin: highest share_event_pct
\providecommand{\rhchLinShareEventPmMin}{}\renewcommand{\rhchLinShareEventPmMin}{51}  % Lin: lowest share_event_pm_pct
\providecommand{\rhchLinShareEventPmMax}{}\renewcommand{\rhchLinShareEventPmMax}{78}  % Lin: highest share_event_pm_pct
\providecommand{\rhchTreeShareEventMin}{}\renewcommand{\rhchTreeShareEventMin}{60}  % Tree: lowest share_event_pct
\providecommand{\rhchTreeShareEventMax}{}\renewcommand{\rhchTreeShareEventMax}{64}  % Tree: highest share_event_pct
\providecommand{\rhchTreeShareEventPmMin}{}\renewcommand{\rhchTreeShareEventPmMin}{59}  % Tree: lowest share_event_pm_pct
\providecommand{\rhchTreeShareEventPmMax}{}\renewcommand{\rhchTreeShareEventPmMax}{63}  % Tree: highest share_event_pm_pct
\providecommand{\rhchLassoRidgeBarPct}{}\renewcommand{\rhchLassoRidgeBarPct}{$-$3.1}  % lasso vs ridge, full sets, all: % QLIKE change
\providecommand{\rhchLassoRidgeBarDm}{}\renewcommand{\rhchLassoRidgeBarDm}{$-$4.82}  % lasso vs ridge, all: DM
\providecommand{\rhchLassoRidgeBarP}{}\renewcommand{\rhchLassoRidgeBarP}{\ensuremath{<}0.001}  % lasso vs ridge, all: DM p
\providecommand{\rhchLassoRidgeEventPct}{}\renewcommand{\rhchLassoRidgeEventPct}{$-$31.9}  % lasso vs ridge, full sets, event: % QLIKE change
\providecommand{\rhchLassoRidgeEventDm}{}\renewcommand{\rhchLassoRidgeEventDm}{$-$8.42}  % lasso vs ridge, event: DM
\providecommand{\rhchLassoRidgeEventP}{}\renewcommand{\rhchLassoRidgeEventP}{\ensuremath{<}0.001}  % lasso vs ridge, event: DM p
\providecommand{\rhchLassoRidgeOrdPct}{}\renewcommand{\rhchLassoRidgeOrdPct}{0.8}  % lasso vs ridge, full sets, ord: % QLIKE change
\providecommand{\rhchLassoRidgeOrdDm}{}\renewcommand{\rhchLassoRidgeOrdDm}{2.28}  % lasso vs ridge, ord: DM
\providecommand{\rhchLassoRidgeOrdP}{}\renewcommand{\rhchLassoRidgeOrdP}{0.023}  % lasso vs ridge, ord: DM p
\providecommand{\rhchLassoEnetBarPct}{}\renewcommand{\rhchLassoEnetBarPct}{$-$2.4}  % lasso vs enet, full sets, all: % QLIKE change
\providecommand{\rhchLassoEnetBarDm}{}\renewcommand{\rhchLassoEnetBarDm}{$-$3.49}  % lasso vs enet, all: DM
\providecommand{\rhchLassoEnetBarP}{}\renewcommand{\rhchLassoEnetBarP}{\ensuremath{<}0.001}  % lasso vs enet, all: DM p
\providecommand{\rhchLassoEnetEventPct}{}\renewcommand{\rhchLassoEnetEventPct}{$-$31.8}  % lasso vs enet, full sets, event: % QLIKE change
\providecommand{\rhchLassoEnetEventDm}{}\renewcommand{\rhchLassoEnetEventDm}{$-$7.69}  % lasso vs enet, event: DM
\providecommand{\rhchLassoEnetEventP}{}\renewcommand{\rhchLassoEnetEventP}{\ensuremath{<}0.001}  % lasso vs enet, event: DM p
\providecommand{\rhchLassoEnetOrdPct}{}\renewcommand{\rhchLassoEnetOrdPct}{1.6}  % lasso vs enet, full sets, ord: % QLIKE change
\providecommand{\rhchLassoEnetOrdDm}{}\renewcommand{\rhchLassoEnetOrdDm}{3.35}  % lasso vs enet, ord: DM
\providecommand{\rhchLassoEnetOrdP}{}\renewcommand{\rhchLassoEnetOrdP}{\ensuremath{<}0.001}  % lasso vs enet, ord: DM p
\providecommand{\rhchRidgeBarJointPct}{}\renewcommand{\rhchRidgeBarJointPct}{$-$6.0}  % ridge all: full vs HAR, % QLIKE change
\providecommand{\rhchRidgeBarJointDm}{}\renewcommand{\rhchRidgeBarJointDm}{$-$6.28}  % ridge all: full vs HAR, DM
\providecommand{\rhchRidgeBarJointP}{}\renewcommand{\rhchRidgeBarJointP}{\ensuremath{<}0.001}  % ridge all: full vs HAR, DM p
\providecommand{\rhchRidgeBarAddRmPct}{}\renewcommand{\rhchRidgeBarAddRmPct}{$-$2.6}  % ridge all add_rm vs ols|har: %
\providecommand{\rhchRidgeBarAddRmP}{}\renewcommand{\rhchRidgeBarAddRmP}{\ensuremath{<}0.001}  % ridge all add_rm: DM p
\providecommand{\rhchRidgeBarLoRmPct}{}\renewcommand{\rhchRidgeBarLoRmPct}{1.8}  % ridge all lo_rm vs ridge|full: %
\providecommand{\rhchRidgeBarLoRmP}{}\renewcommand{\rhchRidgeBarLoRmP}{\ensuremath{<}0.001}  % ridge all lo_rm: DM p
\providecommand{\rhchRidgeBarAddLiqPct}{}\renewcommand{\rhchRidgeBarAddLiqPct}{$-$0.4}  % ridge all add_liq vs ols|har: %
\providecommand{\rhchRidgeBarAddLiqP}{}\renewcommand{\rhchRidgeBarAddLiqP}{0.326}  % ridge all add_liq: DM p
\providecommand{\rhchRidgeBarLoLiqPct}{}\renewcommand{\rhchRidgeBarLoLiqPct}{0.2}  % ridge all lo_liq vs ridge|full: %
\providecommand{\rhchRidgeBarLoLiqP}{}\renewcommand{\rhchRidgeBarLoLiqP}{0.467}  % ridge all lo_liq: DM p
\providecommand{\rhchRidgeBarAddEwPct}{}\renewcommand{\rhchRidgeBarAddEwPct}{$-$0.9}  % ridge all add_ew vs ols|har: %
\providecommand{\rhchRidgeBarAddEwP}{}\renewcommand{\rhchRidgeBarAddEwP}{0.024}  % ridge all add_ew: DM p
\providecommand{\rhchRidgeBarLoEwPct}{}\renewcommand{\rhchRidgeBarLoEwPct}{$-$0.1}  % ridge all lo_ew vs ridge|full: %
\providecommand{\rhchRidgeBarLoEwP}{}\renewcommand{\rhchRidgeBarLoEwP}{0.381}  % ridge all lo_ew: DM p
\providecommand{\rhchRidgeBarAddVwPct}{}\renewcommand{\rhchRidgeBarAddVwPct}{$-$0.9}  % ridge all add_vw vs ols|har: %
\providecommand{\rhchRidgeBarAddVwP}{}\renewcommand{\rhchRidgeBarAddVwP}{0.114}  % ridge all add_vw: DM p
\providecommand{\rhchRidgeBarLoVwPct}{}\renewcommand{\rhchRidgeBarLoVwPct}{$-$0.1}  % ridge all lo_vw vs ridge|full: %
\providecommand{\rhchRidgeBarLoVwP}{}\renewcommand{\rhchRidgeBarLoVwP}{0.487}  % ridge all lo_vw: DM p
\providecommand{\rhchRidgeBarAddSentPct}{}\renewcommand{\rhchRidgeBarAddSentPct}{0.2}  % ridge all add_sent vs ols|har: %
\providecommand{\rhchRidgeBarAddSentP}{}\renewcommand{\rhchRidgeBarAddSentP}{0.411}  % ridge all add_sent: DM p
\providecommand{\rhchRidgeBarLoSentPct}{}\renewcommand{\rhchRidgeBarLoSentPct}{0.1}  % ridge all lo_sent vs ridge|full: %
\providecommand{\rhchRidgeBarLoSentP}{}\renewcommand{\rhchRidgeBarLoSentP}{0.637}  % ridge all lo_sent: DM p
\providecommand{\rhchRidgeBarAddIvPct}{}\renewcommand{\rhchRidgeBarAddIvPct}{0.2}  % ridge all add_iv vs ols|har: %
\providecommand{\rhchRidgeBarAddIvP}{}\renewcommand{\rhchRidgeBarAddIvP}{0.506}  % ridge all add_iv: DM p
\providecommand{\rhchRidgeBarLoIvPct}{}\renewcommand{\rhchRidgeBarLoIvPct}{0.1}  % ridge all lo_iv vs ridge|full: %
\providecommand{\rhchRidgeBarLoIvP}{}\renewcommand{\rhchRidgeBarLoIvP}{0.675}  % ridge all lo_iv: DM p
\providecommand{\rhchRidgeBarAddVdPct}{}\renewcommand{\rhchRidgeBarAddVdPct}{0.5}  % ridge all add_vd vs ols|har: %
\providecommand{\rhchRidgeBarAddVdP}{}\renewcommand{\rhchRidgeBarAddVdP}{0.002}  % ridge all add_vd: DM p
\providecommand{\rhchRidgeBarLoVdPct}{}\renewcommand{\rhchRidgeBarLoVdPct}{$-$0.1}  % ridge all lo_vd vs ridge|full: %
\providecommand{\rhchRidgeBarLoVdP}{}\renewcommand{\rhchRidgeBarLoVdP}{0.472}  % ridge all lo_vd: DM p
\providecommand{\rhchRidgeBarAddFomcPct}{}\renewcommand{\rhchRidgeBarAddFomcPct}{$-$3.1}  % ridge all add_fomc vs ols|har: %
\providecommand{\rhchRidgeBarAddFomcP}{}\renewcommand{\rhchRidgeBarAddFomcP}{\ensuremath{<}0.001}  % ridge all add_fomc: DM p
\providecommand{\rhchRidgeBarLoFomcPct}{}\renewcommand{\rhchRidgeBarLoFomcPct}{3.1}  % ridge all lo_fomc vs ridge|full: %
\providecommand{\rhchRidgeBarLoFomcP}{}\renewcommand{\rhchRidgeBarLoFomcP}{\ensuremath{<}0.001}  % ridge all lo_fomc: DM p
\providecommand{\rhchRidgeEventJointPct}{}\renewcommand{\rhchRidgeEventJointPct}{$-$22.4}  % ridge event: full vs HAR, % QLIKE change
\providecommand{\rhchRidgeEventJointDm}{}\renewcommand{\rhchRidgeEventJointDm}{$-$6.45}  % ridge event: full vs HAR, DM
\providecommand{\rhchRidgeEventJointP}{}\renewcommand{\rhchRidgeEventJointP}{\ensuremath{<}0.001}  % ridge event: full vs HAR, DM p
\providecommand{\rhchRidgeEventAddRmPct}{}\renewcommand{\rhchRidgeEventAddRmPct}{4.4}  % ridge event add_rm vs ols|har: %
\providecommand{\rhchRidgeEventAddRmP}{}\renewcommand{\rhchRidgeEventAddRmP}{0.023}  % ridge event add_rm: DM p
\providecommand{\rhchRidgeEventLoRmPct}{}\renewcommand{\rhchRidgeEventLoRmPct}{$-$1.1}  % ridge event lo_rm vs ridge|full: %
\providecommand{\rhchRidgeEventLoRmP}{}\renewcommand{\rhchRidgeEventLoRmP}{0.133}  % ridge event lo_rm: DM p
\providecommand{\rhchRidgeEventAddLiqPct}{}\renewcommand{\rhchRidgeEventAddLiqPct}{5.4}  % ridge event add_liq vs ols|har: %
\providecommand{\rhchRidgeEventAddLiqP}{}\renewcommand{\rhchRidgeEventAddLiqP}{0.010}  % ridge event add_liq: DM p
\providecommand{\rhchRidgeEventLoLiqPct}{}\renewcommand{\rhchRidgeEventLoLiqPct}{$-$0.6}  % ridge event lo_liq vs ridge|full: %
\providecommand{\rhchRidgeEventLoLiqP}{}\renewcommand{\rhchRidgeEventLoLiqP}{0.230}  % ridge event lo_liq: DM p
\providecommand{\rhchRidgeEventAddEwPct}{}\renewcommand{\rhchRidgeEventAddEwPct}{$-$1.8}  % ridge event add_ew vs ols|har: %
\providecommand{\rhchRidgeEventAddEwP}{}\renewcommand{\rhchRidgeEventAddEwP}{0.176}  % ridge event add_ew: DM p
\providecommand{\rhchRidgeEventLoEwPct}{}\renewcommand{\rhchRidgeEventLoEwPct}{0.0}  % ridge event lo_ew vs ridge|full: %
\providecommand{\rhchRidgeEventLoEwP}{}\renewcommand{\rhchRidgeEventLoEwP}{0.968}  % ridge event lo_ew: DM p
\providecommand{\rhchRidgeEventAddVwPct}{}\renewcommand{\rhchRidgeEventAddVwPct}{$-$2.5}  % ridge event add_vw vs ols|har: %
\providecommand{\rhchRidgeEventAddVwP}{}\renewcommand{\rhchRidgeEventAddVwP}{0.143}  % ridge event add_vw: DM p
\providecommand{\rhchRidgeEventLoVwPct}{}\renewcommand{\rhchRidgeEventLoVwPct}{0.1}  % ridge event lo_vw vs ridge|full: %
\providecommand{\rhchRidgeEventLoVwP}{}\renewcommand{\rhchRidgeEventLoVwP}{0.704}  % ridge event lo_vw: DM p
\providecommand{\rhchRidgeEventAddSentPct}{}\renewcommand{\rhchRidgeEventAddSentPct}{$-$0.1}  % ridge event add_sent vs ols|har: %
\providecommand{\rhchRidgeEventAddSentP}{}\renewcommand{\rhchRidgeEventAddSentP}{0.970}  % ridge event add_sent: DM p
\providecommand{\rhchRidgeEventLoSentPct}{}\renewcommand{\rhchRidgeEventLoSentPct}{1.7}  % ridge event lo_sent vs ridge|full: %
\providecommand{\rhchRidgeEventLoSentP}{}\renewcommand{\rhchRidgeEventLoSentP}{0.322}  % ridge event lo_sent: DM p
\providecommand{\rhchRidgeEventAddIvPct}{}\renewcommand{\rhchRidgeEventAddIvPct}{2.8}  % ridge event add_iv vs ols|har: %
\providecommand{\rhchRidgeEventAddIvP}{}\renewcommand{\rhchRidgeEventAddIvP}{0.011}  % ridge event add_iv: DM p
\providecommand{\rhchRidgeEventLoIvPct}{}\renewcommand{\rhchRidgeEventLoIvPct}{$-$0.2}  % ridge event lo_iv vs ridge|full: %
\providecommand{\rhchRidgeEventLoIvP}{}\renewcommand{\rhchRidgeEventLoIvP}{0.823}  % ridge event lo_iv: DM p
\providecommand{\rhchRidgeEventAddVdPct}{}\renewcommand{\rhchRidgeEventAddVdPct}{1.9}  % ridge event add_vd vs ols|har: %
\providecommand{\rhchRidgeEventAddVdP}{}\renewcommand{\rhchRidgeEventAddVdP}{0.028}  % ridge event add_vd: DM p
\providecommand{\rhchRidgeEventLoVdPct}{}\renewcommand{\rhchRidgeEventLoVdPct}{0.0}  % ridge event lo_vd vs ridge|full: %
\providecommand{\rhchRidgeEventLoVdP}{}\renewcommand{\rhchRidgeEventLoVdP}{0.971}  % ridge event lo_vd: DM p
\providecommand{\rhchRidgeEventAddFomcPct}{}\renewcommand{\rhchRidgeEventAddFomcPct}{$-$23.4}  % ridge event add_fomc vs ols|har: %
\providecommand{\rhchRidgeEventAddFomcP}{}\renewcommand{\rhchRidgeEventAddFomcP}{\ensuremath{<}0.001}  % ridge event add_fomc: DM p
\providecommand{\rhchRidgeEventLoFomcPct}{}\renewcommand{\rhchRidgeEventLoFomcPct}{28.6}  % ridge event lo_fomc vs ridge|full: %
\providecommand{\rhchRidgeEventLoFomcP}{}\renewcommand{\rhchRidgeEventLoFomcP}{\ensuremath{<}0.001}  % ridge event lo_fomc: DM p
\providecommand{\rhchRidgeEventReading}{}\renewcommand{\rhchRidgeEventReading}{concentrated}  % ridge event: item 3 reading
\providecommand{\rhchRidgeEventTopAddShare}{}\renewcommand{\rhchRidgeEventTopAddShare}{104}  % ridge event: largest add-alone gain, % of the joint gain (fomc)
\providecommand{\rhchRidgeEventSumAddRatio}{}\renewcommand{\rhchRidgeEventSumAddRatio}{0.60}  % ridge event: add-alone gains summed over groups / joint gain
\providecommand{\rhchRidgeEventNLoSig}{}\renewcommand{\rhchRidgeEventNLoSig}{1}  % ridge event: significant leave-one-outs
\providecommand{\rhchRidgeEventNLoRun}{}\renewcommand{\rhchRidgeEventNLoRun}{8}  % ridge event: leave-one-outs available (of 8)
\providecommand{\rhchRidgeEventConcTest}{}\renewcommand{\rhchRidgeEventConcTest}{yes}  % ridge event: the CONC definition on its own (pending: undetermined)
\providecommand{\rhchRidgeEventRedTest}{}\renewcommand{\rhchRidgeEventRedTest}{no}  % ridge event: the RED definition on its own (pending: undetermined)
\providecommand{\rhchRidgeEventReadingRedFirst}{}\renewcommand{\rhchRidgeEventReadingRedFirst}{concentrated}  % ridge event: reading with REDUNDANT tested before CONCENTRATED (the RULE tests CONCENTRATED first)
\providecommand{\rhchRidgeOrdJointPct}{}\renewcommand{\rhchRidgeOrdJointPct}{$-$3.2}  % ridge ord: full vs HAR, % QLIKE change
\providecommand{\rhchRidgeOrdJointDm}{}\renewcommand{\rhchRidgeOrdJointDm}{$-$3.51}  % ridge ord: full vs HAR, DM
\providecommand{\rhchRidgeOrdJointP}{}\renewcommand{\rhchRidgeOrdJointP}{\ensuremath{<}0.001}  % ridge ord: full vs HAR, DM p
\providecommand{\rhchRidgeOrdAddRmPct}{}\renewcommand{\rhchRidgeOrdAddRmPct}{$-$3.8}  % ridge ord add_rm vs ols|har: %
\providecommand{\rhchRidgeOrdAddRmP}{}\renewcommand{\rhchRidgeOrdAddRmP}{\ensuremath{<}0.001}  % ridge ord add_rm: DM p
\providecommand{\rhchRidgeOrdLoRmPct}{}\renewcommand{\rhchRidgeOrdLoRmPct}{2.2}  % ridge ord lo_rm vs ridge|full: %
\providecommand{\rhchRidgeOrdLoRmP}{}\renewcommand{\rhchRidgeOrdLoRmP}{\ensuremath{<}0.001}  % ridge ord lo_rm: DM p
\providecommand{\rhchRidgeOrdAddLiqPct}{}\renewcommand{\rhchRidgeOrdAddLiqPct}{$-$1.4}  % ridge ord add_liq vs ols|har: %
\providecommand{\rhchRidgeOrdAddLiqP}{}\renewcommand{\rhchRidgeOrdAddLiqP}{\ensuremath{<}0.001}  % ridge ord add_liq: DM p
\providecommand{\rhchRidgeOrdLoLiqPct}{}\renewcommand{\rhchRidgeOrdLoLiqPct}{0.3}  % ridge ord lo_liq vs ridge|full: %
\providecommand{\rhchRidgeOrdLoLiqP}{}\renewcommand{\rhchRidgeOrdLoLiqP}{0.261}  % ridge ord lo_liq: DM p
\providecommand{\rhchRidgeOrdAddEwPct}{}\renewcommand{\rhchRidgeOrdAddEwPct}{$-$0.7}  % ridge ord add_ew vs ols|har: %
\providecommand{\rhchRidgeOrdAddEwP}{}\renewcommand{\rhchRidgeOrdAddEwP}{0.044}  % ridge ord add_ew: DM p
\providecommand{\rhchRidgeOrdLoEwPct}{}\renewcommand{\rhchRidgeOrdLoEwPct}{$-$0.1}  % ridge ord lo_ew vs ridge|full: %
\providecommand{\rhchRidgeOrdLoEwP}{}\renewcommand{\rhchRidgeOrdLoEwP}{0.316}  % ridge ord lo_ew: DM p
\providecommand{\rhchRidgeOrdAddVwPct}{}\renewcommand{\rhchRidgeOrdAddVwPct}{$-$0.6}  % ridge ord add_vw vs ols|har: %
\providecommand{\rhchRidgeOrdAddVwP}{}\renewcommand{\rhchRidgeOrdAddVwP}{0.309}  % ridge ord add_vw: DM p
\providecommand{\rhchRidgeOrdLoVwPct}{}\renewcommand{\rhchRidgeOrdLoVwPct}{$-$0.1}  % ridge ord lo_vw vs ridge|full: %
\providecommand{\rhchRidgeOrdLoVwP}{}\renewcommand{\rhchRidgeOrdLoVwP}{0.405}  % ridge ord lo_vw: DM p
\providecommand{\rhchRidgeOrdAddSentPct}{}\renewcommand{\rhchRidgeOrdAddSentPct}{0.3}  % ridge ord add_sent vs ols|har: %
\providecommand{\rhchRidgeOrdAddSentP}{}\renewcommand{\rhchRidgeOrdAddSentP}{0.246}  % ridge ord add_sent: DM p
\providecommand{\rhchRidgeOrdLoSentPct}{}\renewcommand{\rhchRidgeOrdLoSentPct}{$-$0.1}  % ridge ord lo_sent vs ridge|full: %
\providecommand{\rhchRidgeOrdLoSentP}{}\renewcommand{\rhchRidgeOrdLoSentP}{0.271}  % ridge ord lo_sent: DM p
\providecommand{\rhchRidgeOrdAddIvPct}{}\renewcommand{\rhchRidgeOrdAddIvPct}{$-$0.2}  % ridge ord add_iv vs ols|har: %
\providecommand{\rhchRidgeOrdAddIvP}{}\renewcommand{\rhchRidgeOrdAddIvP}{0.580}  % ridge ord add_iv: DM p
\providecommand{\rhchRidgeOrdLoIvPct}{}\renewcommand{\rhchRidgeOrdLoIvPct}{0.2}  % ridge ord lo_iv vs ridge|full: %
\providecommand{\rhchRidgeOrdLoIvP}{}\renewcommand{\rhchRidgeOrdLoIvP}{0.611}  % ridge ord lo_iv: DM p
\providecommand{\rhchRidgeOrdAddVdPct}{}\renewcommand{\rhchRidgeOrdAddVdPct}{0.3}  % ridge ord add_vd vs ols|har: %
\providecommand{\rhchRidgeOrdAddVdP}{}\renewcommand{\rhchRidgeOrdAddVdP}{0.069}  % ridge ord add_vd: DM p
\providecommand{\rhchRidgeOrdLoVdPct}{}\renewcommand{\rhchRidgeOrdLoVdPct}{$-$0.1}  % ridge ord lo_vd vs ridge|full: %
\providecommand{\rhchRidgeOrdLoVdP}{}\renewcommand{\rhchRidgeOrdLoVdP}{0.437}  % ridge ord lo_vd: DM p
\providecommand{\rhchRidgeOrdAddFomcPct}{}\renewcommand{\rhchRidgeOrdAddFomcPct}{0.4}  % ridge ord add_fomc vs ols|har: %
\providecommand{\rhchRidgeOrdAddFomcP}{}\renewcommand{\rhchRidgeOrdAddFomcP}{\ensuremath{<}0.001}  % ridge ord add_fomc: DM p
\providecommand{\rhchRidgeOrdLoFomcPct}{}\renewcommand{\rhchRidgeOrdLoFomcPct}{$-$0.4}  % ridge ord lo_fomc vs ridge|full: %
\providecommand{\rhchRidgeOrdLoFomcP}{}\renewcommand{\rhchRidgeOrdLoFomcP}{\ensuremath{<}0.001}  % ridge ord lo_fomc: DM p
\providecommand{\rhchRidgeOrdReading}{}\renewcommand{\rhchRidgeOrdReading}{concentrated}  % ridge ord: item 3 reading
\providecommand{\rhchRidgeOrdTopAddShare}{}\renewcommand{\rhchRidgeOrdTopAddShare}{118}  % ridge ord: largest add-alone gain, % of the joint gain (rm)
\providecommand{\rhchRidgeOrdSumAddRatio}{}\renewcommand{\rhchRidgeOrdSumAddRatio}{1.79}  % ridge ord: add-alone gains summed over groups / joint gain
\providecommand{\rhchRidgeOrdNLoSig}{}\renewcommand{\rhchRidgeOrdNLoSig}{1}  % ridge ord: significant leave-one-outs
\providecommand{\rhchRidgeOrdNLoRun}{}\renewcommand{\rhchRidgeOrdNLoRun}{8}  % ridge ord: leave-one-outs available (of 8)
\providecommand{\rhchRidgeOrdConcTest}{}\renewcommand{\rhchRidgeOrdConcTest}{yes}  % ridge ord: the CONC definition on its own (pending: undetermined)
\providecommand{\rhchRidgeOrdRedTest}{}\renewcommand{\rhchRidgeOrdRedTest}{yes}  % ridge ord: the RED definition on its own (pending: undetermined)
\providecommand{\rhchRidgeOrdReadingRedFirst}{}\renewcommand{\rhchRidgeOrdReadingRedFirst}{redundant}  % ridge ord: reading with REDUNDANT tested before CONCENTRATED (the RULE tests CONCENTRATED first)
\providecommand{\rhchLassoBarJointPct}{}\renewcommand{\rhchLassoBarJointPct}{$-$9.0}  % lasso all: full vs HAR, % QLIKE change
\providecommand{\rhchLassoBarJointDm}{}\renewcommand{\rhchLassoBarJointDm}{$-$6.34}  % lasso all: full vs HAR, DM
\providecommand{\rhchLassoBarJointP}{}\renewcommand{\rhchLassoBarJointP}{\ensuremath{<}0.001}  % lasso all: full vs HAR, DM p
\providecommand{\rhchLassoBarAddRmPct}{}\renewcommand{\rhchLassoBarAddRmPct}{$-$2.1}  % lasso all add_rm vs ols|har: %
\providecommand{\rhchLassoBarAddRmP}{}\renewcommand{\rhchLassoBarAddRmP}{0.008}  % lasso all add_rm: DM p
\providecommand{\rhchLassoBarLoRmPct}{}\renewcommand{\rhchLassoBarLoRmPct}{1.6}  % lasso all lo_rm vs lasso|full: %
\providecommand{\rhchLassoBarLoRmP}{}\renewcommand{\rhchLassoBarLoRmP}{0.010}  % lasso all lo_rm: DM p
\providecommand{\rhchLassoBarAddLiqPct}{}\renewcommand{\rhchLassoBarAddLiqPct}{$-$0.3}  % lasso all add_liq vs ols|har: %
\providecommand{\rhchLassoBarAddLiqP}{}\renewcommand{\rhchLassoBarAddLiqP}{0.583}  % lasso all add_liq: DM p
\providecommand{\rhchLassoBarLoLiqPct}{}\renewcommand{\rhchLassoBarLoLiqPct}{0.1}  % lasso all lo_liq vs lasso|full: %
\providecommand{\rhchLassoBarLoLiqP}{}\renewcommand{\rhchLassoBarLoLiqP}{0.783}  % lasso all lo_liq: DM p
\providecommand{\rhchLassoBarAddEwPct}{}\renewcommand{\rhchLassoBarAddEwPct}{$-$0.8}  % lasso all add_ew vs ols|har: %
\providecommand{\rhchLassoBarAddEwP}{}\renewcommand{\rhchLassoBarAddEwP}{0.045}  % lasso all add_ew: DM p
\providecommand{\rhchLassoBarLoEwPct}{}\renewcommand{\rhchLassoBarLoEwPct}{$-$0.1}  % lasso all lo_ew vs lasso|full: %
\providecommand{\rhchLassoBarLoEwP}{}\renewcommand{\rhchLassoBarLoEwP}{0.583}  % lasso all lo_ew: DM p
\providecommand{\rhchLassoBarAddVwPct}{}\renewcommand{\rhchLassoBarAddVwPct}{$-$0.6}  % lasso all add_vw vs ols|har: %
\providecommand{\rhchLassoBarAddVwP}{}\renewcommand{\rhchLassoBarAddVwP}{0.309}  % lasso all add_vw: DM p
\providecommand{\rhchLassoBarLoVwPct}{}\renewcommand{\rhchLassoBarLoVwPct}{$-$0.2}  % lasso all lo_vw vs lasso|full: %
\providecommand{\rhchLassoBarLoVwP}{}\renewcommand{\rhchLassoBarLoVwP}{0.428}  % lasso all lo_vw: DM p
\providecommand{\rhchLassoBarAddSentPct}{}\renewcommand{\rhchLassoBarAddSentPct}{0.1}  % lasso all add_sent vs ols|har: %
\providecommand{\rhchLassoBarAddSentP}{}\renewcommand{\rhchLassoBarAddSentP}{0.550}  % lasso all add_sent: DM p
\providecommand{\rhchLassoBarLoSentPct}{}\renewcommand{\rhchLassoBarLoSentPct}{0.1}  % lasso all lo_sent vs lasso|full: %
\providecommand{\rhchLassoBarLoSentP}{}\renewcommand{\rhchLassoBarLoSentP}{0.693}  % lasso all lo_sent: DM p
\providecommand{\rhchLassoBarAddIvPct}{}\renewcommand{\rhchLassoBarAddIvPct}{$-$0.3}  % lasso all add_iv vs ols|har: %
\providecommand{\rhchLassoBarAddIvP}{}\renewcommand{\rhchLassoBarAddIvP}{0.430}  % lasso all add_iv: DM p
\providecommand{\rhchLassoBarLoIvPct}{}\renewcommand{\rhchLassoBarLoIvPct}{0.5}  % lasso all lo_iv vs lasso|full: %
\providecommand{\rhchLassoBarLoIvP}{}\renewcommand{\rhchLassoBarLoIvP}{0.313}  % lasso all lo_iv: DM p
\providecommand{\rhchLassoBarAddVdPct}{}\renewcommand{\rhchLassoBarAddVdPct}{0.6}  % lasso all add_vd vs ols|har: %
\providecommand{\rhchLassoBarAddVdP}{}\renewcommand{\rhchLassoBarAddVdP}{0.002}  % lasso all add_vd: DM p
\providecommand{\rhchLassoBarLoVdPct}{}\renewcommand{\rhchLassoBarLoVdPct}{0.0}  % lasso all lo_vd vs lasso|full: %
\providecommand{\rhchLassoBarLoVdP}{}\renewcommand{\rhchLassoBarLoVdP}{0.807}  % lasso all lo_vd: DM p
\providecommand{\rhchLassoBarAddFomcPct}{}\renewcommand{\rhchLassoBarAddFomcPct}{$-$5.8}  % lasso all add_fomc vs ols|har: %
\providecommand{\rhchLassoBarAddFomcP}{}\renewcommand{\rhchLassoBarAddFomcP}{\ensuremath{<}0.001}  % lasso all add_fomc: DM p
\providecommand{\rhchLassoBarLoFomcPct}{}\renewcommand{\rhchLassoBarLoFomcPct}{6.8}  % lasso all lo_fomc vs lasso|full: %
\providecommand{\rhchLassoBarLoFomcP}{}\renewcommand{\rhchLassoBarLoFomcP}{\ensuremath{<}0.001}  % lasso all lo_fomc: DM p
\providecommand{\rhchLassoEventJointPct}{}\renewcommand{\rhchLassoEventJointPct}{$-$47.2}  % lasso event: full vs HAR, % QLIKE change
\providecommand{\rhchLassoEventJointDm}{}\renewcommand{\rhchLassoEventJointDm}{$-$7.56}  % lasso event: full vs HAR, DM
\providecommand{\rhchLassoEventJointP}{}\renewcommand{\rhchLassoEventJointP}{\ensuremath{<}0.001}  % lasso event: full vs HAR, DM p
\providecommand{\rhchLassoEventAddRmPct}{}\renewcommand{\rhchLassoEventAddRmPct}{4.7}  % lasso event add_rm vs ols|har: %
\providecommand{\rhchLassoEventAddRmP}{}\renewcommand{\rhchLassoEventAddRmP}{0.024}  % lasso event add_rm: DM p
\providecommand{\rhchLassoEventLoRmPct}{}\renewcommand{\rhchLassoEventLoRmPct}{$-$0.8}  % lasso event lo_rm vs lasso|full: %
\providecommand{\rhchLassoEventLoRmP}{}\renewcommand{\rhchLassoEventLoRmP}{0.265}  % lasso event lo_rm: DM p
\providecommand{\rhchLassoEventAddLiqPct}{}\renewcommand{\rhchLassoEventAddLiqPct}{6.1}  % lasso event add_liq vs ols|har: %
\providecommand{\rhchLassoEventAddLiqP}{}\renewcommand{\rhchLassoEventAddLiqP}{0.007}  % lasso event add_liq: DM p
\providecommand{\rhchLassoEventLoLiqPct}{}\renewcommand{\rhchLassoEventLoLiqPct}{0.1}  % lasso event lo_liq vs lasso|full: %
\providecommand{\rhchLassoEventLoLiqP}{}\renewcommand{\rhchLassoEventLoLiqP}{0.838}  % lasso event lo_liq: DM p
\providecommand{\rhchLassoEventAddEwPct}{}\renewcommand{\rhchLassoEventAddEwPct}{$-$1.7}  % lasso event add_ew vs ols|har: %
\providecommand{\rhchLassoEventAddEwP}{}\renewcommand{\rhchLassoEventAddEwP}{0.235}  % lasso event add_ew: DM p
\providecommand{\rhchLassoEventLoEwPct}{}\renewcommand{\rhchLassoEventLoEwPct}{0.1}  % lasso event lo_ew vs lasso|full: %
\providecommand{\rhchLassoEventLoEwP}{}\renewcommand{\rhchLassoEventLoEwP}{0.717}  % lasso event lo_ew: DM p
\providecommand{\rhchLassoEventAddVwPct}{}\renewcommand{\rhchLassoEventAddVwPct}{$-$2.4}  % lasso event add_vw vs ols|har: %
\providecommand{\rhchLassoEventAddVwP}{}\renewcommand{\rhchLassoEventAddVwP}{0.171}  % lasso event add_vw: DM p
\providecommand{\rhchLassoEventLoVwPct}{}\renewcommand{\rhchLassoEventLoVwPct}{$-$0.3}  % lasso event lo_vw vs lasso|full: %
\providecommand{\rhchLassoEventLoVwP}{}\renewcommand{\rhchLassoEventLoVwP}{0.438}  % lasso event lo_vw: DM p
\providecommand{\rhchLassoEventAddSentPct}{}\renewcommand{\rhchLassoEventAddSentPct}{$-$1.3}  % lasso event add_sent vs ols|har: %
\providecommand{\rhchLassoEventAddSentP}{}\renewcommand{\rhchLassoEventAddSentP}{0.464}  % lasso event add_sent: DM p
\providecommand{\rhchLassoEventLoSentPct}{}\renewcommand{\rhchLassoEventLoSentPct}{2.0}  % lasso event lo_sent vs lasso|full: %
\providecommand{\rhchLassoEventLoSentP}{}\renewcommand{\rhchLassoEventLoSentP}{0.209}  % lasso event lo_sent: DM p
\providecommand{\rhchLassoEventAddIvPct}{}\renewcommand{\rhchLassoEventAddIvPct}{2.9}  % lasso event add_iv vs ols|har: %
\providecommand{\rhchLassoEventAddIvP}{}\renewcommand{\rhchLassoEventAddIvP}{0.004}  % lasso event add_iv: DM p
\providecommand{\rhchLassoEventLoIvPct}{}\renewcommand{\rhchLassoEventLoIvPct}{$-$0.1}  % lasso event lo_iv vs lasso|full: %
\providecommand{\rhchLassoEventLoIvP}{}\renewcommand{\rhchLassoEventLoIvP}{0.868}  % lasso event lo_iv: DM p
\providecommand{\rhchLassoEventAddVdPct}{}\renewcommand{\rhchLassoEventAddVdPct}{1.9}  % lasso event add_vd vs ols|har: %
\providecommand{\rhchLassoEventAddVdP}{}\renewcommand{\rhchLassoEventAddVdP}{0.032}  % lasso event add_vd: DM p
\providecommand{\rhchLassoEventLoVdPct}{}\renewcommand{\rhchLassoEventLoVdPct}{0.3}  % lasso event lo_vd vs lasso|full: %
\providecommand{\rhchLassoEventLoVdP}{}\renewcommand{\rhchLassoEventLoVdP}{0.012}  % lasso event lo_vd: DM p
\providecommand{\rhchLassoEventAddFomcPct}{}\renewcommand{\rhchLassoEventAddFomcPct}{$-$43.9}  % lasso event add_fomc vs ols|har: %
\providecommand{\rhchLassoEventAddFomcP}{}\renewcommand{\rhchLassoEventAddFomcP}{\ensuremath{<}0.001}  % lasso event add_fomc: DM p
\providecommand{\rhchLassoEventLoFomcPct}{}\renewcommand{\rhchLassoEventLoFomcPct}{87.4}  % lasso event lo_fomc vs lasso|full: %
\providecommand{\rhchLassoEventLoFomcP}{}\renewcommand{\rhchLassoEventLoFomcP}{\ensuremath{<}0.001}  % lasso event lo_fomc: DM p
\providecommand{\rhchLassoEventReading}{}\renewcommand{\rhchLassoEventReading}{concentrated}  % lasso event: item 3 reading
\providecommand{\rhchLassoEventTopAddShare}{}\renewcommand{\rhchLassoEventTopAddShare}{93}  % lasso event: largest add-alone gain, % of the joint gain (fomc)
\providecommand{\rhchLassoEventSumAddRatio}{}\renewcommand{\rhchLassoEventSumAddRatio}{0.72}  % lasso event: add-alone gains summed over groups / joint gain
\providecommand{\rhchLassoEventNLoSig}{}\renewcommand{\rhchLassoEventNLoSig}{2}  % lasso event: significant leave-one-outs
\providecommand{\rhchLassoEventNLoRun}{}\renewcommand{\rhchLassoEventNLoRun}{8}  % lasso event: leave-one-outs available (of 8)
\providecommand{\rhchLassoEventConcTest}{}\renewcommand{\rhchLassoEventConcTest}{yes}  % lasso event: the CONC definition on its own (pending: undetermined)
\providecommand{\rhchLassoEventRedTest}{}\renewcommand{\rhchLassoEventRedTest}{no}  % lasso event: the RED definition on its own (pending: undetermined)
\providecommand{\rhchLassoEventReadingRedFirst}{}\renewcommand{\rhchLassoEventReadingRedFirst}{concentrated}  % lasso event: reading with REDUNDANT tested before CONCENTRATED (the RULE tests CONCENTRATED first)
\providecommand{\rhchLassoOrdJointPct}{}\renewcommand{\rhchLassoOrdJointPct}{$-$2.4}  % lasso ord: full vs HAR, % QLIKE change
\providecommand{\rhchLassoOrdJointDm}{}\renewcommand{\rhchLassoOrdJointDm}{$-$2.08}  % lasso ord: full vs HAR, DM
\providecommand{\rhchLassoOrdJointP}{}\renewcommand{\rhchLassoOrdJointP}{0.037}  % lasso ord: full vs HAR, DM p
\providecommand{\rhchLassoOrdAddRmPct}{}\renewcommand{\rhchLassoOrdAddRmPct}{$-$3.3}  % lasso ord add_rm vs ols|har: %
\providecommand{\rhchLassoOrdAddRmP}{}\renewcommand{\rhchLassoOrdAddRmP}{\ensuremath{<}0.001}  % lasso ord add_rm: DM p
\providecommand{\rhchLassoOrdLoRmPct}{}\renewcommand{\rhchLassoOrdLoRmPct}{1.8}  % lasso ord lo_rm vs lasso|full: %
\providecommand{\rhchLassoOrdLoRmP}{}\renewcommand{\rhchLassoOrdLoRmP}{0.007}  % lasso ord lo_rm: DM p
\providecommand{\rhchLassoOrdAddLiqPct}{}\renewcommand{\rhchLassoOrdAddLiqPct}{$-$1.3}  % lasso ord add_liq vs ols|har: %
\providecommand{\rhchLassoOrdAddLiqP}{}\renewcommand{\rhchLassoOrdAddLiqP}{0.003}  % lasso ord add_liq: DM p
\providecommand{\rhchLassoOrdLoLiqPct}{}\renewcommand{\rhchLassoOrdLoLiqPct}{0.1}  % lasso ord lo_liq vs lasso|full: %
\providecommand{\rhchLassoOrdLoLiqP}{}\renewcommand{\rhchLassoOrdLoLiqP}{0.799}  % lasso ord lo_liq: DM p
\providecommand{\rhchLassoOrdAddEwPct}{}\renewcommand{\rhchLassoOrdAddEwPct}{$-$0.7}  % lasso ord add_ew vs ols|har: %
\providecommand{\rhchLassoOrdAddEwP}{}\renewcommand{\rhchLassoOrdAddEwP}{0.077}  % lasso ord add_ew: DM p
\providecommand{\rhchLassoOrdLoEwPct}{}\renewcommand{\rhchLassoOrdLoEwPct}{$-$0.1}  % lasso ord lo_ew vs lasso|full: %
\providecommand{\rhchLassoOrdLoEwP}{}\renewcommand{\rhchLassoOrdLoEwP}{0.530}  % lasso ord lo_ew: DM p
\providecommand{\rhchLassoOrdAddVwPct}{}\renewcommand{\rhchLassoOrdAddVwPct}{$-$0.3}  % lasso ord add_vw vs ols|har: %
\providecommand{\rhchLassoOrdAddVwP}{}\renewcommand{\rhchLassoOrdAddVwP}{0.630}  % lasso ord add_vw: DM p
\providecommand{\rhchLassoOrdLoVwPct}{}\renewcommand{\rhchLassoOrdLoVwPct}{$-$0.2}  % lasso ord lo_vw vs lasso|full: %
\providecommand{\rhchLassoOrdLoVwP}{}\renewcommand{\rhchLassoOrdLoVwP}{0.501}  % lasso ord lo_vw: DM p
\providecommand{\rhchLassoOrdAddSentPct}{}\renewcommand{\rhchLassoOrdAddSentPct}{0.4}  % lasso ord add_sent vs ols|har: %
\providecommand{\rhchLassoOrdAddSentP}{}\renewcommand{\rhchLassoOrdAddSentP}{0.098}  % lasso ord add_sent: DM p
\providecommand{\rhchLassoOrdLoSentPct}{}\renewcommand{\rhchLassoOrdLoSentPct}{$-$0.1}  % lasso ord lo_sent vs lasso|full: %
\providecommand{\rhchLassoOrdLoSentP}{}\renewcommand{\rhchLassoOrdLoSentP}{0.502}  % lasso ord lo_sent: DM p
\providecommand{\rhchLassoOrdAddIvPct}{}\renewcommand{\rhchLassoOrdAddIvPct}{$-$0.9}  % lasso ord add_iv vs ols|har: %
\providecommand{\rhchLassoOrdAddIvP}{}\renewcommand{\rhchLassoOrdAddIvP}{0.042}  % lasso ord add_iv: DM p
\providecommand{\rhchLassoOrdLoIvPct}{}\renewcommand{\rhchLassoOrdLoIvPct}{0.5}  % lasso ord lo_iv vs lasso|full: %
\providecommand{\rhchLassoOrdLoIvP}{}\renewcommand{\rhchLassoOrdLoIvP}{0.292}  % lasso ord lo_iv: DM p
\providecommand{\rhchLassoOrdAddVdPct}{}\renewcommand{\rhchLassoOrdAddVdPct}{0.3}  % lasso ord add_vd vs ols|har: %
\providecommand{\rhchLassoOrdAddVdP}{}\renewcommand{\rhchLassoOrdAddVdP}{0.056}  % lasso ord add_vd: DM p
\providecommand{\rhchLassoOrdLoVdPct}{}\renewcommand{\rhchLassoOrdLoVdPct}{$-$0.1}  % lasso ord lo_vd vs lasso|full: %
\providecommand{\rhchLassoOrdLoVdP}{}\renewcommand{\rhchLassoOrdLoVdP}{0.644}  % lasso ord lo_vd: DM p
\providecommand{\rhchLassoOrdAddFomcPct}{}\renewcommand{\rhchLassoOrdAddFomcPct}{0.7}  % lasso ord add_fomc vs ols|har: %
\providecommand{\rhchLassoOrdAddFomcP}{}\renewcommand{\rhchLassoOrdAddFomcP}{\ensuremath{<}0.001}  % lasso ord add_fomc: DM p
\providecommand{\rhchLassoOrdLoFomcPct}{}\renewcommand{\rhchLassoOrdLoFomcPct}{$-$0.7}  % lasso ord lo_fomc vs lasso|full: %
\providecommand{\rhchLassoOrdLoFomcP}{}\renewcommand{\rhchLassoOrdLoFomcP}{\ensuremath{<}0.001}  % lasso ord lo_fomc: DM p
\providecommand{\rhchLassoOrdReading}{}\renewcommand{\rhchLassoOrdReading}{concentrated}  % lasso ord: item 3 reading
\providecommand{\rhchLassoOrdTopAddShare}{}\renewcommand{\rhchLassoOrdTopAddShare}{136}  % lasso ord: largest add-alone gain, % of the joint gain (rm)
\providecommand{\rhchLassoOrdSumAddRatio}{}\renewcommand{\rhchLassoOrdSumAddRatio}{2.07}  % lasso ord: add-alone gains summed over groups / joint gain
\providecommand{\rhchLassoOrdNLoSig}{}\renewcommand{\rhchLassoOrdNLoSig}{1}  % lasso ord: significant leave-one-outs
\providecommand{\rhchLassoOrdNLoRun}{}\renewcommand{\rhchLassoOrdNLoRun}{8}  % lasso ord: leave-one-outs available (of 8)
\providecommand{\rhchLassoOrdConcTest}{}\renewcommand{\rhchLassoOrdConcTest}{yes}  % lasso ord: the CONC definition on its own (pending: undetermined)
\providecommand{\rhchLassoOrdRedTest}{}\renewcommand{\rhchLassoOrdRedTest}{yes}  % lasso ord: the RED definition on its own (pending: undetermined)
\providecommand{\rhchLassoOrdReadingRedFirst}{}\renewcommand{\rhchLassoOrdReadingRedFirst}{redundant}  % lasso ord: reading with REDUNDANT tested before CONCENTRATED (the RULE tests CONCENTRATED first)
\providecommand{\rhchEnetBarJointPct}{}\renewcommand{\rhchEnetBarJointPct}{$-$6.7}  % enet all: full vs HAR, % QLIKE change
\providecommand{\rhchEnetBarJointDm}{}\renewcommand{\rhchEnetBarJointDm}{$-$7.36}  % enet all: full vs HAR, DM
\providecommand{\rhchEnetBarJointP}{}\renewcommand{\rhchEnetBarJointP}{\ensuremath{<}0.001}  % enet all: full vs HAR, DM p
\providecommand{\rhchEnetBarAddRmPct}{}\renewcommand{\rhchEnetBarAddRmPct}{$-$2.7}  % enet all add_rm vs ols|har: %
\providecommand{\rhchEnetBarAddRmP}{}\renewcommand{\rhchEnetBarAddRmP}{\ensuremath{<}0.001}  % enet all add_rm: DM p
\providecommand{\rhchEnetBarLoRmPct}{}\renewcommand{\rhchEnetBarLoRmPct}{2.0}  % enet all lo_rm vs enet|full: %
\providecommand{\rhchEnetBarLoRmP}{}\renewcommand{\rhchEnetBarLoRmP}{\ensuremath{<}0.001}  % enet all lo_rm: DM p
\providecommand{\rhchEnetBarAddLiqPct}{}\renewcommand{\rhchEnetBarAddLiqPct}{$-$0.6}  % enet all add_liq vs ols|har: %
\providecommand{\rhchEnetBarAddLiqP}{}\renewcommand{\rhchEnetBarAddLiqP}{0.124}  % enet all add_liq: DM p
\providecommand{\rhchEnetBarLoLiqPct}{}\renewcommand{\rhchEnetBarLoLiqPct}{0.1}  % enet all lo_liq vs enet|full: %
\providecommand{\rhchEnetBarLoLiqP}{}\renewcommand{\rhchEnetBarLoLiqP}{0.638}  % enet all lo_liq: DM p
\providecommand{\rhchEnetBarAddEwPct}{}\renewcommand{\rhchEnetBarAddEwPct}{$-$1.1}  % enet all add_ew vs ols|har: %
\providecommand{\rhchEnetBarAddEwP}{}\renewcommand{\rhchEnetBarAddEwP}{0.005}  % enet all add_ew: DM p
\providecommand{\rhchEnetBarLoEwPct}{}\renewcommand{\rhchEnetBarLoEwPct}{$-$0.1}  % enet all lo_ew vs enet|full: %
\providecommand{\rhchEnetBarLoEwP}{}\renewcommand{\rhchEnetBarLoEwP}{0.327}  % enet all lo_ew: DM p
\providecommand{\rhchEnetBarAddVwPct}{}\renewcommand{\rhchEnetBarAddVwPct}{$-$1.1}  % enet all add_vw vs ols|har: %
\providecommand{\rhchEnetBarAddVwP}{}\renewcommand{\rhchEnetBarAddVwP}{0.042}  % enet all add_vw: DM p
\providecommand{\rhchEnetBarLoVwPct}{}\renewcommand{\rhchEnetBarLoVwPct}{$-$0.1}  % enet all lo_vw vs enet|full: %
\providecommand{\rhchEnetBarLoVwP}{}\renewcommand{\rhchEnetBarLoVwP}{0.492}  % enet all lo_vw: DM p
\providecommand{\rhchEnetBarAddSentPct}{}\renewcommand{\rhchEnetBarAddSentPct}{0.0}  % enet all add_sent vs ols|har: %
\providecommand{\rhchEnetBarAddSentP}{}\renewcommand{\rhchEnetBarAddSentP}{0.971}  % enet all add_sent: DM p
\providecommand{\rhchEnetBarLoSentPct}{}\renewcommand{\rhchEnetBarLoSentPct}{0.2}  % enet all lo_sent vs enet|full: %
\providecommand{\rhchEnetBarLoSentP}{}\renewcommand{\rhchEnetBarLoSentP}{0.137}  % enet all lo_sent: DM p
\providecommand{\rhchEnetBarAddIvPct}{}\renewcommand{\rhchEnetBarAddIvPct}{0.1}  % enet all add_iv vs ols|har: %
\providecommand{\rhchEnetBarAddIvP}{}\renewcommand{\rhchEnetBarAddIvP}{0.817}  % enet all add_iv: DM p
\providecommand{\rhchEnetBarLoIvPct}{}\renewcommand{\rhchEnetBarLoIvPct}{0.0}  % enet all lo_iv vs enet|full: %
\providecommand{\rhchEnetBarLoIvP}{}\renewcommand{\rhchEnetBarLoIvP}{0.950}  % enet all lo_iv: DM p
\providecommand{\rhchEnetBarAddVdPct}{}\renewcommand{\rhchEnetBarAddVdPct}{0.4}  % enet all add_vd vs ols|har: %
\providecommand{\rhchEnetBarAddVdP}{}\renewcommand{\rhchEnetBarAddVdP}{0.003}  % enet all add_vd: DM p
\providecommand{\rhchEnetBarLoVdPct}{}\renewcommand{\rhchEnetBarLoVdPct}{$-$0.1}  % enet all lo_vd vs enet|full: %
\providecommand{\rhchEnetBarLoVdP}{}\renewcommand{\rhchEnetBarLoVdP}{0.264}  % enet all lo_vd: DM p
\providecommand{\rhchEnetBarAddFomcPct}{}\renewcommand{\rhchEnetBarAddFomcPct}{$-$3.1}  % enet all add_fomc vs ols|har: %
\providecommand{\rhchEnetBarAddFomcP}{}\renewcommand{\rhchEnetBarAddFomcP}{\ensuremath{<}0.001}  % enet all add_fomc: DM p
\providecommand{\rhchEnetBarLoFomcPct}{}\renewcommand{\rhchEnetBarLoFomcPct}{3.1}  % enet all lo_fomc vs enet|full: %
\providecommand{\rhchEnetBarLoFomcP}{}\renewcommand{\rhchEnetBarLoFomcP}{\ensuremath{<}0.001}  % enet all lo_fomc: DM p
\providecommand{\rhchEnetEventJointPct}{}\renewcommand{\rhchEnetEventJointPct}{$-$22.6}  % enet event: full vs HAR, % QLIKE change
\providecommand{\rhchEnetEventJointDm}{}\renewcommand{\rhchEnetEventJointDm}{$-$6.79}  % enet event: full vs HAR, DM
\providecommand{\rhchEnetEventJointP}{}\renewcommand{\rhchEnetEventJointP}{\ensuremath{<}0.001}  % enet event: full vs HAR, DM p
\providecommand{\rhchEnetEventAddRmPct}{}\renewcommand{\rhchEnetEventAddRmPct}{4.2}  % enet event add_rm vs ols|har: %
\providecommand{\rhchEnetEventAddRmP}{}\renewcommand{\rhchEnetEventAddRmP}{0.031}  % enet event add_rm: DM p
\providecommand{\rhchEnetEventLoRmPct}{}\renewcommand{\rhchEnetEventLoRmPct}{$-$1.2}  % enet event lo_rm vs enet|full: %
\providecommand{\rhchEnetEventLoRmP}{}\renewcommand{\rhchEnetEventLoRmP}{0.167}  % enet event lo_rm: DM p
\providecommand{\rhchEnetEventAddLiqPct}{}\renewcommand{\rhchEnetEventAddLiqPct}{5.3}  % enet event add_liq vs ols|har: %
\providecommand{\rhchEnetEventAddLiqP}{}\renewcommand{\rhchEnetEventAddLiqP}{0.008}  % enet event add_liq: DM p
\providecommand{\rhchEnetEventLoLiqPct}{}\renewcommand{\rhchEnetEventLoLiqPct}{$-$0.7}  % enet event lo_liq vs enet|full: %
\providecommand{\rhchEnetEventLoLiqP}{}\renewcommand{\rhchEnetEventLoLiqP}{0.104}  % enet event lo_liq: DM p
\providecommand{\rhchEnetEventAddEwPct}{}\renewcommand{\rhchEnetEventAddEwPct}{$-$2.7}  % enet event add_ew vs ols|har: %
\providecommand{\rhchEnetEventAddEwP}{}\renewcommand{\rhchEnetEventAddEwP}{0.060}  % enet event add_ew: DM p
\providecommand{\rhchEnetEventLoEwPct}{}\renewcommand{\rhchEnetEventLoEwPct}{0.0}  % enet event lo_ew vs enet|full: %
\providecommand{\rhchEnetEventLoEwP}{}\renewcommand{\rhchEnetEventLoEwP}{0.915}  % enet event lo_ew: DM p
\providecommand{\rhchEnetEventAddVwPct}{}\renewcommand{\rhchEnetEventAddVwPct}{$-$3.4}  % enet event add_vw vs ols|har: %
\providecommand{\rhchEnetEventAddVwP}{}\renewcommand{\rhchEnetEventAddVwP}{0.055}  % enet event add_vw: DM p
\providecommand{\rhchEnetEventLoVwPct}{}\renewcommand{\rhchEnetEventLoVwPct}{0.0}  % enet event lo_vw vs enet|full: %
\providecommand{\rhchEnetEventLoVwP}{}\renewcommand{\rhchEnetEventLoVwP}{0.947}  % enet event lo_vw: DM p
\providecommand{\rhchEnetEventAddSentPct}{}\renewcommand{\rhchEnetEventAddSentPct}{$-$0.9}  % enet event add_sent vs ols|har: %
\providecommand{\rhchEnetEventAddSentP}{}\renewcommand{\rhchEnetEventAddSentP}{0.519}  % enet event add_sent: DM p
\providecommand{\rhchEnetEventLoSentPct}{}\renewcommand{\rhchEnetEventLoSentPct}{1.6}  % enet event lo_sent vs enet|full: %
\providecommand{\rhchEnetEventLoSentP}{}\renewcommand{\rhchEnetEventLoSentP}{0.201}  % enet event lo_sent: DM p
\providecommand{\rhchEnetEventAddIvPct}{}\renewcommand{\rhchEnetEventAddIvPct}{2.5}  % enet event add_iv vs ols|har: %
\providecommand{\rhchEnetEventAddIvP}{}\renewcommand{\rhchEnetEventAddIvP}{0.010}  % enet event add_iv: DM p
\providecommand{\rhchEnetEventLoIvPct}{}\renewcommand{\rhchEnetEventLoIvPct}{$-$0.1}  % enet event lo_iv vs enet|full: %
\providecommand{\rhchEnetEventLoIvP}{}\renewcommand{\rhchEnetEventLoIvP}{0.865}  % enet event lo_iv: DM p
\providecommand{\rhchEnetEventAddVdPct}{}\renewcommand{\rhchEnetEventAddVdPct}{1.6}  % enet event add_vd vs ols|har: %
\providecommand{\rhchEnetEventAddVdP}{}\renewcommand{\rhchEnetEventAddVdP}{0.007}  % enet event add_vd: DM p
\providecommand{\rhchEnetEventLoVdPct}{}\renewcommand{\rhchEnetEventLoVdPct}{$-$0.1}  % enet event lo_vd vs enet|full: %
\providecommand{\rhchEnetEventLoVdP}{}\renewcommand{\rhchEnetEventLoVdP}{0.330}  % enet event lo_vd: DM p
\providecommand{\rhchEnetEventAddFomcPct}{}\renewcommand{\rhchEnetEventAddFomcPct}{$-$23.2}  % enet event add_fomc vs ols|har: %
\providecommand{\rhchEnetEventAddFomcP}{}\renewcommand{\rhchEnetEventAddFomcP}{\ensuremath{<}0.001}  % enet event add_fomc: DM p
\providecommand{\rhchEnetEventLoFomcPct}{}\renewcommand{\rhchEnetEventLoFomcPct}{28.2}  % enet event lo_fomc vs enet|full: %
\providecommand{\rhchEnetEventLoFomcP}{}\renewcommand{\rhchEnetEventLoFomcP}{\ensuremath{<}0.001}  % enet event lo_fomc: DM p
\providecommand{\rhchEnetEventReading}{}\renewcommand{\rhchEnetEventReading}{concentrated}  % enet event: item 3 reading
\providecommand{\rhchEnetEventTopAddShare}{}\renewcommand{\rhchEnetEventTopAddShare}{103}  % enet event: largest add-alone gain, % of the joint gain (fomc)
\providecommand{\rhchEnetEventSumAddRatio}{}\renewcommand{\rhchEnetEventSumAddRatio}{0.74}  % enet event: add-alone gains summed over groups / joint gain
\providecommand{\rhchEnetEventNLoSig}{}\renewcommand{\rhchEnetEventNLoSig}{1}  % enet event: significant leave-one-outs
\providecommand{\rhchEnetEventNLoRun}{}\renewcommand{\rhchEnetEventNLoRun}{8}  % enet event: leave-one-outs available (of 8)
\providecommand{\rhchEnetEventConcTest}{}\renewcommand{\rhchEnetEventConcTest}{yes}  % enet event: the CONC definition on its own (pending: undetermined)
\providecommand{\rhchEnetEventRedTest}{}\renewcommand{\rhchEnetEventRedTest}{no}  % enet event: the RED definition on its own (pending: undetermined)
\providecommand{\rhchEnetEventReadingRedFirst}{}\renewcommand{\rhchEnetEventReadingRedFirst}{concentrated}  % enet event: reading with REDUNDANT tested before CONCENTRATED (the RULE tests CONCENTRATED first)
\providecommand{\rhchEnetOrdJointPct}{}\renewcommand{\rhchEnetOrdJointPct}{$-$4.0}  % enet ord: full vs HAR, % QLIKE change
\providecommand{\rhchEnetOrdJointDm}{}\renewcommand{\rhchEnetOrdJointDm}{$-$4.73}  % enet ord: full vs HAR, DM
\providecommand{\rhchEnetOrdJointP}{}\renewcommand{\rhchEnetOrdJointP}{\ensuremath{<}0.001}  % enet ord: full vs HAR, DM p
\providecommand{\rhchEnetOrdAddRmPct}{}\renewcommand{\rhchEnetOrdAddRmPct}{$-$3.9}  % enet ord add_rm vs ols|har: %
\providecommand{\rhchEnetOrdAddRmP}{}\renewcommand{\rhchEnetOrdAddRmP}{\ensuremath{<}0.001}  % enet ord add_rm: DM p
\providecommand{\rhchEnetOrdLoRmPct}{}\renewcommand{\rhchEnetOrdLoRmPct}{2.5}  % enet ord lo_rm vs enet|full: %
\providecommand{\rhchEnetOrdLoRmP}{}\renewcommand{\rhchEnetOrdLoRmP}{\ensuremath{<}0.001}  % enet ord lo_rm: DM p
\providecommand{\rhchEnetOrdAddLiqPct}{}\renewcommand{\rhchEnetOrdAddLiqPct}{$-$1.6}  % enet ord add_liq vs ols|har: %
\providecommand{\rhchEnetOrdAddLiqP}{}\renewcommand{\rhchEnetOrdAddLiqP}{\ensuremath{<}0.001}  % enet ord add_liq: DM p
\providecommand{\rhchEnetOrdLoLiqPct}{}\renewcommand{\rhchEnetOrdLoLiqPct}{0.2}  % enet ord lo_liq vs enet|full: %
\providecommand{\rhchEnetOrdLoLiqP}{}\renewcommand{\rhchEnetOrdLoLiqP}{0.282}  % enet ord lo_liq: DM p
\providecommand{\rhchEnetOrdAddEwPct}{}\renewcommand{\rhchEnetOrdAddEwPct}{$-$0.9}  % enet ord add_ew vs ols|har: %
\providecommand{\rhchEnetOrdAddEwP}{}\renewcommand{\rhchEnetOrdAddEwP}{0.018}  % enet ord add_ew: DM p
\providecommand{\rhchEnetOrdLoEwPct}{}\renewcommand{\rhchEnetOrdLoEwPct}{$-$0.1}  % enet ord lo_ew vs enet|full: %
\providecommand{\rhchEnetOrdLoEwP}{}\renewcommand{\rhchEnetOrdLoEwP}{0.220}  % enet ord lo_ew: DM p
\providecommand{\rhchEnetOrdAddVwPct}{}\renewcommand{\rhchEnetOrdAddVwPct}{$-$0.7}  % enet ord add_vw vs ols|har: %
\providecommand{\rhchEnetOrdAddVwP}{}\renewcommand{\rhchEnetOrdAddVwP}{0.195}  % enet ord add_vw: DM p
\providecommand{\rhchEnetOrdLoVwPct}{}\renewcommand{\rhchEnetOrdLoVwPct}{$-$0.1}  % enet ord lo_vw vs enet|full: %
\providecommand{\rhchEnetOrdLoVwP}{}\renewcommand{\rhchEnetOrdLoVwP}{0.491}  % enet ord lo_vw: DM p
\providecommand{\rhchEnetOrdAddSentPct}{}\renewcommand{\rhchEnetOrdAddSentPct}{0.2}  % enet ord add_sent vs ols|har: %
\providecommand{\rhchEnetOrdAddSentP}{}\renewcommand{\rhchEnetOrdAddSentP}{0.407}  % enet ord add_sent: DM p
\providecommand{\rhchEnetOrdLoSentPct}{}\renewcommand{\rhchEnetOrdLoSentPct}{0.0}  % enet ord lo_sent vs enet|full: %
\providecommand{\rhchEnetOrdLoSentP}{}\renewcommand{\rhchEnetOrdLoSentP}{0.998}  % enet ord lo_sent: DM p
\providecommand{\rhchEnetOrdAddIvPct}{}\renewcommand{\rhchEnetOrdAddIvPct}{$-$0.3}  % enet ord add_iv vs ols|har: %
\providecommand{\rhchEnetOrdAddIvP}{}\renewcommand{\rhchEnetOrdAddIvP}{0.319}  % enet ord add_iv: DM p
\providecommand{\rhchEnetOrdLoIvPct}{}\renewcommand{\rhchEnetOrdLoIvPct}{0.0}  % enet ord lo_iv vs enet|full: %
\providecommand{\rhchEnetOrdLoIvP}{}\renewcommand{\rhchEnetOrdLoIvP}{0.912}  % enet ord lo_iv: DM p
\providecommand{\rhchEnetOrdAddVdPct}{}\renewcommand{\rhchEnetOrdAddVdPct}{0.2}  % enet ord add_vd vs ols|har: %
\providecommand{\rhchEnetOrdAddVdP}{}\renewcommand{\rhchEnetOrdAddVdP}{0.127}  % enet ord add_vd: DM p
\providecommand{\rhchEnetOrdLoVdPct}{}\renewcommand{\rhchEnetOrdLoVdPct}{$-$0.1}  % enet ord lo_vd vs enet|full: %
\providecommand{\rhchEnetOrdLoVdP}{}\renewcommand{\rhchEnetOrdLoVdP}{0.345}  % enet ord lo_vd: DM p
\providecommand{\rhchEnetOrdAddFomcPct}{}\renewcommand{\rhchEnetOrdAddFomcPct}{0.4}  % enet ord add_fomc vs ols|har: %
\providecommand{\rhchEnetOrdAddFomcP}{}\renewcommand{\rhchEnetOrdAddFomcP}{\ensuremath{<}0.001}  % enet ord add_fomc: DM p
\providecommand{\rhchEnetOrdLoFomcPct}{}\renewcommand{\rhchEnetOrdLoFomcPct}{$-$0.4}  % enet ord lo_fomc vs enet|full: %
\providecommand{\rhchEnetOrdLoFomcP}{}\renewcommand{\rhchEnetOrdLoFomcP}{\ensuremath{<}0.001}  % enet ord lo_fomc: DM p
\providecommand{\rhchEnetOrdReading}{}\renewcommand{\rhchEnetOrdReading}{concentrated}  % enet ord: item 3 reading
\providecommand{\rhchEnetOrdTopAddShare}{}\renewcommand{\rhchEnetOrdTopAddShare}{97}  % enet ord: largest add-alone gain, % of the joint gain (rm)
\providecommand{\rhchEnetOrdSumAddRatio}{}\renewcommand{\rhchEnetOrdSumAddRatio}{1.68}  % enet ord: add-alone gains summed over groups / joint gain
\providecommand{\rhchEnetOrdNLoSig}{}\renewcommand{\rhchEnetOrdNLoSig}{1}  % enet ord: significant leave-one-outs
\providecommand{\rhchEnetOrdNLoRun}{}\renewcommand{\rhchEnetOrdNLoRun}{8}  % enet ord: leave-one-outs available (of 8)
\providecommand{\rhchEnetOrdConcTest}{}\renewcommand{\rhchEnetOrdConcTest}{yes}  % enet ord: the CONC definition on its own (pending: undetermined)
\providecommand{\rhchEnetOrdRedTest}{}\renewcommand{\rhchEnetOrdRedTest}{yes}  % enet ord: the RED definition on its own (pending: undetermined)
\providecommand{\rhchEnetOrdReadingRedFirst}{}\renewcommand{\rhchEnetOrdReadingRedFirst}{redundant}  % enet ord: reading with REDUNDANT tested before CONCENTRATED (the RULE tests CONCENTRATED first)
\providecommand{\rhchLgbmBarJointPct}{}\renewcommand{\rhchLgbmBarJointPct}{$-$14.1}  % lgbm all: full vs HAR, % QLIKE change
\providecommand{\rhchLgbmBarJointDm}{}\renewcommand{\rhchLgbmBarJointDm}{$-$10.26}  % lgbm all: full vs HAR, DM
\providecommand{\rhchLgbmBarJointP}{}\renewcommand{\rhchLgbmBarJointP}{\ensuremath{<}0.001}  % lgbm all: full vs HAR, DM p
\providecommand{\rhchLgbmBarAddRmPct}{}\renewcommand{\rhchLgbmBarAddRmPct}{$-$3.0}  % lgbm all add_rm vs lgbm|har: %
\providecommand{\rhchLgbmBarAddRmP}{}\renewcommand{\rhchLgbmBarAddRmP}{\ensuremath{<}0.001}  % lgbm all add_rm: DM p
\providecommand{\rhchLgbmBarLoRmPct}{}\renewcommand{\rhchLgbmBarLoRmPct}{1.5}  % lgbm all lo_rm vs lgbm|full_pre: % [pre-patch]
\providecommand{\rhchLgbmBarLoRmP}{}\renewcommand{\rhchLgbmBarLoRmP}{\ensuremath{<}0.001}  % lgbm all lo_rm: DM p [pre-patch]
\providecommand{\rhchLgbmBarAddLiqPct}{}\renewcommand{\rhchLgbmBarAddLiqPct}{$-$2.9}  % lgbm all add_liq vs lgbm|har: %
\providecommand{\rhchLgbmBarAddLiqP}{}\renewcommand{\rhchLgbmBarAddLiqP}{\ensuremath{<}0.001}  % lgbm all add_liq: DM p
\providecommand{\rhchLgbmBarLoLiqPct}{}\renewcommand{\rhchLgbmBarLoLiqPct}{0.2}  % lgbm all lo_liq vs lgbm|full_pre: % [pre-patch]
\providecommand{\rhchLgbmBarLoLiqP}{}\renewcommand{\rhchLgbmBarLoLiqP}{0.365}  % lgbm all lo_liq: DM p [pre-patch]
\providecommand{\rhchLgbmBarAddEwPct}{}\renewcommand{\rhchLgbmBarAddEwPct}{$-$2.8}  % lgbm all add_ew vs lgbm|har: %
\providecommand{\rhchLgbmBarAddEwP}{}\renewcommand{\rhchLgbmBarAddEwP}{\ensuremath{<}0.001}  % lgbm all add_ew: DM p
\providecommand{\rhchLgbmBarLoEwPct}{}\renewcommand{\rhchLgbmBarLoEwPct}{0.0}  % lgbm all lo_ew vs lgbm|full_pre: % [pre-patch]
\providecommand{\rhchLgbmBarLoEwP}{}\renewcommand{\rhchLgbmBarLoEwP}{0.815}  % lgbm all lo_ew: DM p [pre-patch]
\providecommand{\rhchLgbmBarAddVwPct}{}\renewcommand{\rhchLgbmBarAddVwPct}{$-$3.1}  % lgbm all add_vw vs lgbm|har: %
\providecommand{\rhchLgbmBarAddVwP}{}\renewcommand{\rhchLgbmBarAddVwP}{\ensuremath{<}0.001}  % lgbm all add_vw: DM p
\providecommand{\rhchLgbmBarLoVwPct}{}\renewcommand{\rhchLgbmBarLoVwPct}{0.1}  % lgbm all lo_vw vs lgbm|full_pre: % [pre-patch]
\providecommand{\rhchLgbmBarLoVwP}{}\renewcommand{\rhchLgbmBarLoVwP}{0.489}  % lgbm all lo_vw: DM p [pre-patch]
\providecommand{\rhchLgbmBarAddSentPct}{}\renewcommand{\rhchLgbmBarAddSentPct}{$-$1.1}  % lgbm all add_sent vs lgbm|har: %
\providecommand{\rhchLgbmBarAddSentP}{}\renewcommand{\rhchLgbmBarAddSentP}{\ensuremath{<}0.001}  % lgbm all add_sent: DM p
\providecommand{\rhchLgbmBarLoSentPct}{}\renewcommand{\rhchLgbmBarLoSentPct}{0.2}  % lgbm all lo_sent vs lgbm|full_pre: % [pre-patch]
\providecommand{\rhchLgbmBarLoSentP}{}\renewcommand{\rhchLgbmBarLoSentP}{0.028}  % lgbm all lo_sent: DM p [pre-patch]
\providecommand{\rhchLgbmBarAddIvPct}{}\renewcommand{\rhchLgbmBarAddIvPct}{$-$1.3}  % lgbm all add_iv vs lgbm|har: %
\providecommand{\rhchLgbmBarAddIvP}{}\renewcommand{\rhchLgbmBarAddIvP}{\ensuremath{<}0.001}  % lgbm all add_iv: DM p
\providecommand{\rhchLgbmBarLoIvPct}{}\renewcommand{\rhchLgbmBarLoIvPct}{0.2}  % lgbm all lo_iv vs lgbm|full_pre: % [pre-patch]
\providecommand{\rhchLgbmBarLoIvP}{}\renewcommand{\rhchLgbmBarLoIvP}{0.084}  % lgbm all lo_iv: DM p [pre-patch]
\providecommand{\rhchLgbmBarAddVdPct}{}\renewcommand{\rhchLgbmBarAddVdPct}{$-$1.2}  % lgbm all add_vd vs lgbm|har: %
\providecommand{\rhchLgbmBarAddVdP}{}\renewcommand{\rhchLgbmBarAddVdP}{\ensuremath{<}0.001}  % lgbm all add_vd: DM p
\providecommand{\rhchLgbmBarLoVdPct}{}\renewcommand{\rhchLgbmBarLoVdPct}{0.0}  % lgbm all lo_vd vs lgbm|full_pre: % [pre-patch]
\providecommand{\rhchLgbmBarLoVdP}{}\renewcommand{\rhchLgbmBarLoVdP}{0.916}  % lgbm all lo_vd: DM p [pre-patch]
\providecommand{\rhchLgbmBarAddFomcPct}{}\renewcommand{\rhchLgbmBarAddFomcPct}{$-$9.2}  % lgbm all add_fomc vs lgbm|har: %
\providecommand{\rhchLgbmBarAddFomcP}{}\renewcommand{\rhchLgbmBarAddFomcP}{\ensuremath{<}0.001}  % lgbm all add_fomc: DM p
\providecommand{\rhchLgbmBarLoFomcPct}{}\renewcommand{\rhchLgbmBarLoFomcPct}{10.3}  % lgbm all lo_fomc vs lgbm|full: %
\providecommand{\rhchLgbmBarLoFomcP}{}\renewcommand{\rhchLgbmBarLoFomcP}{\ensuremath{<}0.001}  % lgbm all lo_fomc: DM p
\providecommand{\rhchLgbmEventJointPct}{}\renewcommand{\rhchLgbmEventJointPct}{$-$62.3}  % lgbm event: full vs HAR, % QLIKE change
\providecommand{\rhchLgbmEventJointDm}{}\renewcommand{\rhchLgbmEventJointDm}{$-$8.62}  % lgbm event: full vs HAR, DM
\providecommand{\rhchLgbmEventJointP}{}\renewcommand{\rhchLgbmEventJointP}{\ensuremath{<}0.001}  % lgbm event: full vs HAR, DM p
\providecommand{\rhchLgbmEventAddRmPct}{}\renewcommand{\rhchLgbmEventAddRmPct}{6.8}  % lgbm event add_rm vs lgbm|har: %
\providecommand{\rhchLgbmEventAddRmP}{}\renewcommand{\rhchLgbmEventAddRmP}{0.010}  % lgbm event add_rm: DM p
\providecommand{\rhchLgbmEventLoRmPct}{}\renewcommand{\rhchLgbmEventLoRmPct}{$-$0.3}  % lgbm event lo_rm vs lgbm|full_pre: % [pre-patch]
\providecommand{\rhchLgbmEventLoRmP}{}\renewcommand{\rhchLgbmEventLoRmP}{0.700}  % lgbm event lo_rm: DM p [pre-patch]
\providecommand{\rhchLgbmEventAddLiqPct}{}\renewcommand{\rhchLgbmEventAddLiqPct}{1.5}  % lgbm event add_liq vs lgbm|har: %
\providecommand{\rhchLgbmEventAddLiqP}{}\renewcommand{\rhchLgbmEventAddLiqP}{0.294}  % lgbm event add_liq: DM p
\providecommand{\rhchLgbmEventLoLiqPct}{}\renewcommand{\rhchLgbmEventLoLiqPct}{$-$0.5}  % lgbm event lo_liq vs lgbm|full_pre: % [pre-patch]
\providecommand{\rhchLgbmEventLoLiqP}{}\renewcommand{\rhchLgbmEventLoLiqP}{0.155}  % lgbm event lo_liq: DM p [pre-patch]
\providecommand{\rhchLgbmEventAddEwPct}{}\renewcommand{\rhchLgbmEventAddEwPct}{$-$3.5}  % lgbm event add_ew vs lgbm|har: %
\providecommand{\rhchLgbmEventAddEwP}{}\renewcommand{\rhchLgbmEventAddEwP}{0.045}  % lgbm event add_ew: DM p
\providecommand{\rhchLgbmEventLoEwPct}{}\renewcommand{\rhchLgbmEventLoEwPct}{0.1}  % lgbm event lo_ew vs lgbm|full_pre: % [pre-patch]
\providecommand{\rhchLgbmEventLoEwP}{}\renewcommand{\rhchLgbmEventLoEwP}{0.888}  % lgbm event lo_ew: DM p [pre-patch]
\providecommand{\rhchLgbmEventAddVwPct}{}\renewcommand{\rhchLgbmEventAddVwPct}{$-$3.5}  % lgbm event add_vw vs lgbm|har: %
\providecommand{\rhchLgbmEventAddVwP}{}\renewcommand{\rhchLgbmEventAddVwP}{0.044}  % lgbm event add_vw: DM p
\providecommand{\rhchLgbmEventLoVwPct}{}\renewcommand{\rhchLgbmEventLoVwPct}{$-$0.1}  % lgbm event lo_vw vs lgbm|full_pre: % [pre-patch]
\providecommand{\rhchLgbmEventLoVwP}{}\renewcommand{\rhchLgbmEventLoVwP}{0.804}  % lgbm event lo_vw: DM p [pre-patch]
\providecommand{\rhchLgbmEventAddSentPct}{}\renewcommand{\rhchLgbmEventAddSentPct}{0.5}  % lgbm event add_sent vs lgbm|har: %
\providecommand{\rhchLgbmEventAddSentP}{}\renewcommand{\rhchLgbmEventAddSentP}{0.574}  % lgbm event add_sent: DM p
\providecommand{\rhchLgbmEventLoSentPct}{}\renewcommand{\rhchLgbmEventLoSentPct}{0.2}  % lgbm event lo_sent vs lgbm|full_pre: % [pre-patch]
\providecommand{\rhchLgbmEventLoSentP}{}\renewcommand{\rhchLgbmEventLoSentP}{0.481}  % lgbm event lo_sent: DM p [pre-patch]
\providecommand{\rhchLgbmEventAddIvPct}{}\renewcommand{\rhchLgbmEventAddIvPct}{$-$1.4}  % lgbm event add_iv vs lgbm|har: %
\providecommand{\rhchLgbmEventAddIvP}{}\renewcommand{\rhchLgbmEventAddIvP}{0.036}  % lgbm event add_iv: DM p
\providecommand{\rhchLgbmEventLoIvPct}{}\renewcommand{\rhchLgbmEventLoIvPct}{0.7}  % lgbm event lo_iv vs lgbm|full_pre: % [pre-patch]
\providecommand{\rhchLgbmEventLoIvP}{}\renewcommand{\rhchLgbmEventLoIvP}{0.129}  % lgbm event lo_iv: DM p [pre-patch]
\providecommand{\rhchLgbmEventAddVdPct}{}\renewcommand{\rhchLgbmEventAddVdPct}{$-$3.5}  % lgbm event add_vd vs lgbm|har: %
\providecommand{\rhchLgbmEventAddVdP}{}\renewcommand{\rhchLgbmEventAddVdP}{0.072}  % lgbm event add_vd: DM p
\providecommand{\rhchLgbmEventLoVdPct}{}\renewcommand{\rhchLgbmEventLoVdPct}{0.7}  % lgbm event lo_vd vs lgbm|full_pre: % [pre-patch]
\providecommand{\rhchLgbmEventLoVdP}{}\renewcommand{\rhchLgbmEventLoVdP}{0.010}  % lgbm event lo_vd: DM p [pre-patch]
\providecommand{\rhchLgbmEventAddFomcPct}{}\renewcommand{\rhchLgbmEventAddFomcPct}{$-$62.6}  % lgbm event add_fomc vs lgbm|har: %
\providecommand{\rhchLgbmEventAddFomcP}{}\renewcommand{\rhchLgbmEventAddFomcP}{\ensuremath{<}0.001}  % lgbm event add_fomc: DM p
\providecommand{\rhchLgbmEventLoFomcPct}{}\renewcommand{\rhchLgbmEventLoFomcPct}{163.2}  % lgbm event lo_fomc vs lgbm|full: %
\providecommand{\rhchLgbmEventLoFomcP}{}\renewcommand{\rhchLgbmEventLoFomcP}{\ensuremath{<}0.001}  % lgbm event lo_fomc: DM p
\providecommand{\rhchLgbmEventReading}{}\renewcommand{\rhchLgbmEventReading}{concentrated}  % lgbm event: item 3 reading
\providecommand{\rhchLgbmEventTopAddShare}{}\renewcommand{\rhchLgbmEventTopAddShare}{101}  % lgbm event: largest add-alone gain, % of the joint gain (fomc)
\providecommand{\rhchLgbmEventSumAddRatio}{}\renewcommand{\rhchLgbmEventSumAddRatio}{1.05}  % lgbm event: add-alone gains summed over groups / joint gain
\providecommand{\rhchLgbmEventNLoSig}{}\renewcommand{\rhchLgbmEventNLoSig}{2}  % lgbm event: significant leave-one-outs
\providecommand{\rhchLgbmEventNLoRun}{}\renewcommand{\rhchLgbmEventNLoRun}{8}  % lgbm event: leave-one-outs available (of 8)
\providecommand{\rhchLgbmEventConcTest}{}\renewcommand{\rhchLgbmEventConcTest}{yes}  % lgbm event: the CONC definition on its own (pending: undetermined)
\providecommand{\rhchLgbmEventRedTest}{}\renewcommand{\rhchLgbmEventRedTest}{no}  % lgbm event: the RED definition on its own (pending: undetermined)
\providecommand{\rhchLgbmEventReadingRedFirst}{}\renewcommand{\rhchLgbmEventReadingRedFirst}{concentrated}  % lgbm event: reading with REDUNDANT tested before CONCENTRATED (the RULE tests CONCENTRATED first)
\providecommand{\rhchLgbmOrdJointPct}{}\renewcommand{\rhchLgbmOrdJointPct}{$-$6.0}  % lgbm ord: full vs HAR, % QLIKE change
\providecommand{\rhchLgbmOrdJointDm}{}\renewcommand{\rhchLgbmOrdJointDm}{$-$10.70}  % lgbm ord: full vs HAR, DM
\providecommand{\rhchLgbmOrdJointP}{}\renewcommand{\rhchLgbmOrdJointP}{\ensuremath{<}0.001}  % lgbm ord: full vs HAR, DM p
\providecommand{\rhchLgbmOrdAddRmPct}{}\renewcommand{\rhchLgbmOrdAddRmPct}{$-$4.6}  % lgbm ord add_rm vs lgbm|har: %
\providecommand{\rhchLgbmOrdAddRmP}{}\renewcommand{\rhchLgbmOrdAddRmP}{\ensuremath{<}0.001}  % lgbm ord add_rm: DM p
\providecommand{\rhchLgbmOrdLoRmPct}{}\renewcommand{\rhchLgbmOrdLoRmPct}{1.6}  % lgbm ord lo_rm vs lgbm|full_pre: % [pre-patch]
\providecommand{\rhchLgbmOrdLoRmP}{}\renewcommand{\rhchLgbmOrdLoRmP}{\ensuremath{<}0.001}  % lgbm ord lo_rm: DM p [pre-patch]
\providecommand{\rhchLgbmOrdAddLiqPct}{}\renewcommand{\rhchLgbmOrdAddLiqPct}{$-$3.7}  % lgbm ord add_liq vs lgbm|har: %
\providecommand{\rhchLgbmOrdAddLiqP}{}\renewcommand{\rhchLgbmOrdAddLiqP}{\ensuremath{<}0.001}  % lgbm ord add_liq: DM p
\providecommand{\rhchLgbmOrdLoLiqPct}{}\renewcommand{\rhchLgbmOrdLoLiqPct}{0.2}  % lgbm ord lo_liq vs lgbm|full_pre: % [pre-patch]
\providecommand{\rhchLgbmOrdLoLiqP}{}\renewcommand{\rhchLgbmOrdLoLiqP}{0.266}  % lgbm ord lo_liq: DM p [pre-patch]
\providecommand{\rhchLgbmOrdAddEwPct}{}\renewcommand{\rhchLgbmOrdAddEwPct}{$-$2.7}  % lgbm ord add_ew vs lgbm|har: %
\providecommand{\rhchLgbmOrdAddEwP}{}\renewcommand{\rhchLgbmOrdAddEwP}{\ensuremath{<}0.001}  % lgbm ord add_ew: DM p
\providecommand{\rhchLgbmOrdLoEwPct}{}\renewcommand{\rhchLgbmOrdLoEwPct}{0.0}  % lgbm ord lo_ew vs lgbm|full_pre: % [pre-patch]
\providecommand{\rhchLgbmOrdLoEwP}{}\renewcommand{\rhchLgbmOrdLoEwP}{0.779}  % lgbm ord lo_ew: DM p [pre-patch]
\providecommand{\rhchLgbmOrdAddVwPct}{}\renewcommand{\rhchLgbmOrdAddVwPct}{$-$3.1}  % lgbm ord add_vw vs lgbm|har: %
\providecommand{\rhchLgbmOrdAddVwP}{}\renewcommand{\rhchLgbmOrdAddVwP}{\ensuremath{<}0.001}  % lgbm ord add_vw: DM p
\providecommand{\rhchLgbmOrdLoVwPct}{}\renewcommand{\rhchLgbmOrdLoVwPct}{0.1}  % lgbm ord lo_vw vs lgbm|full_pre: % [pre-patch]
\providecommand{\rhchLgbmOrdLoVwP}{}\renewcommand{\rhchLgbmOrdLoVwP}{0.433}  % lgbm ord lo_vw: DM p [pre-patch]
\providecommand{\rhchLgbmOrdAddSentPct}{}\renewcommand{\rhchLgbmOrdAddSentPct}{$-$1.4}  % lgbm ord add_sent vs lgbm|har: %
\providecommand{\rhchLgbmOrdAddSentP}{}\renewcommand{\rhchLgbmOrdAddSentP}{\ensuremath{<}0.001}  % lgbm ord add_sent: DM p
\providecommand{\rhchLgbmOrdLoSentPct}{}\renewcommand{\rhchLgbmOrdLoSentPct}{0.2}  % lgbm ord lo_sent vs lgbm|full_pre: % [pre-patch]
\providecommand{\rhchLgbmOrdLoSentP}{}\renewcommand{\rhchLgbmOrdLoSentP}{0.035}  % lgbm ord lo_sent: DM p [pre-patch]
\providecommand{\rhchLgbmOrdAddIvPct}{}\renewcommand{\rhchLgbmOrdAddIvPct}{$-$1.3}  % lgbm ord add_iv vs lgbm|har: %
\providecommand{\rhchLgbmOrdAddIvP}{}\renewcommand{\rhchLgbmOrdAddIvP}{\ensuremath{<}0.001}  % lgbm ord add_iv: DM p
\providecommand{\rhchLgbmOrdLoIvPct}{}\renewcommand{\rhchLgbmOrdLoIvPct}{0.1}  % lgbm ord lo_iv vs lgbm|full_pre: % [pre-patch]
\providecommand{\rhchLgbmOrdLoIvP}{}\renewcommand{\rhchLgbmOrdLoIvP}{0.196}  % lgbm ord lo_iv: DM p [pre-patch]
\providecommand{\rhchLgbmOrdAddVdPct}{}\renewcommand{\rhchLgbmOrdAddVdPct}{$-$0.9}  % lgbm ord add_vd vs lgbm|har: %
\providecommand{\rhchLgbmOrdAddVdP}{}\renewcommand{\rhchLgbmOrdAddVdP}{\ensuremath{<}0.001}  % lgbm ord add_vd: DM p
\providecommand{\rhchLgbmOrdLoVdPct}{}\renewcommand{\rhchLgbmOrdLoVdPct}{$-$0.1}  % lgbm ord lo_vd vs lgbm|full_pre: % [pre-patch]
\providecommand{\rhchLgbmOrdLoVdP}{}\renewcommand{\rhchLgbmOrdLoVdP}{0.499}  % lgbm ord lo_vd: DM p [pre-patch]
\providecommand{\rhchLgbmOrdAddFomcPct}{}\renewcommand{\rhchLgbmOrdAddFomcPct}{$-$0.2}  % lgbm ord add_fomc vs lgbm|har: %
\providecommand{\rhchLgbmOrdAddFomcP}{}\renewcommand{\rhchLgbmOrdAddFomcP}{0.364}  % lgbm ord add_fomc: DM p
\providecommand{\rhchLgbmOrdLoFomcPct}{}\renewcommand{\rhchLgbmOrdLoFomcPct}{0.0}  % lgbm ord lo_fomc vs lgbm|full: %
\providecommand{\rhchLgbmOrdLoFomcP}{}\renewcommand{\rhchLgbmOrdLoFomcP}{0.912}  % lgbm ord lo_fomc: DM p
\providecommand{\rhchLgbmOrdReading}{}\renewcommand{\rhchLgbmOrdReading}{concentrated}  % lgbm ord: item 3 reading
\providecommand{\rhchLgbmOrdTopAddShare}{}\renewcommand{\rhchLgbmOrdTopAddShare}{77}  % lgbm ord: largest add-alone gain, % of the joint gain (rm)
\providecommand{\rhchLgbmOrdSumAddRatio}{}\renewcommand{\rhchLgbmOrdSumAddRatio}{2.95}  % lgbm ord: add-alone gains summed over groups / joint gain
\providecommand{\rhchLgbmOrdNLoSig}{}\renewcommand{\rhchLgbmOrdNLoSig}{2}  % lgbm ord: significant leave-one-outs
\providecommand{\rhchLgbmOrdNLoRun}{}\renewcommand{\rhchLgbmOrdNLoRun}{8}  % lgbm ord: leave-one-outs available (of 8)
\providecommand{\rhchLgbmOrdConcTest}{}\renewcommand{\rhchLgbmOrdConcTest}{yes}  % lgbm ord: the CONC definition on its own (pending: undetermined)
\providecommand{\rhchLgbmOrdRedTest}{}\renewcommand{\rhchLgbmOrdRedTest}{no}  % lgbm ord: the RED definition on its own (pending: undetermined)
\providecommand{\rhchLgbmOrdReadingRedFirst}{}\renewcommand{\rhchLgbmOrdReadingRedFirst}{concentrated}  % lgbm ord: reading with REDUNDANT tested before CONCENTRATED (the RULE tests CONCENTRATED first)
\providecommand{\rhchXgbBarJointPct}{}\renewcommand{\rhchXgbBarJointPct}{$-$14.6}  % xgb all: full vs HAR, % QLIKE change
\providecommand{\rhchXgbBarJointDm}{}\renewcommand{\rhchXgbBarJointDm}{$-$9.94}  % xgb all: full vs HAR, DM
\providecommand{\rhchXgbBarJointP}{}\renewcommand{\rhchXgbBarJointP}{\ensuremath{<}0.001}  % xgb all: full vs HAR, DM p
\providecommand{\rhchXgbBarAddRmPct}{}\renewcommand{\rhchXgbBarAddRmPct}{$-$3.4}  % xgb all add_rm vs xgb|har: %
\providecommand{\rhchXgbBarAddRmP}{}\renewcommand{\rhchXgbBarAddRmP}{\ensuremath{<}0.001}  % xgb all add_rm: DM p
\providecommand{\rhchXgbBarLoRmPct}{}\renewcommand{\rhchXgbBarLoRmPct}{1.6}  % xgb all lo_rm vs xgb|full_pre: % [pre-patch]
\providecommand{\rhchXgbBarLoRmP}{}\renewcommand{\rhchXgbBarLoRmP}{\ensuremath{<}0.001}  % xgb all lo_rm: DM p [pre-patch]
\providecommand{\rhchXgbBarAddLiqPct}{}\renewcommand{\rhchXgbBarAddLiqPct}{$-$3.5}  % xgb all add_liq vs xgb|har: %
\providecommand{\rhchXgbBarAddLiqP}{}\renewcommand{\rhchXgbBarAddLiqP}{\ensuremath{<}0.001}  % xgb all add_liq: DM p
\providecommand{\rhchXgbBarLoLiqPct}{}\renewcommand{\rhchXgbBarLoLiqPct}{0.2}  % xgb all lo_liq vs xgb|full_pre: % [pre-patch]
\providecommand{\rhchXgbBarLoLiqP}{}\renewcommand{\rhchXgbBarLoLiqP}{0.267}  % xgb all lo_liq: DM p [pre-patch]
\providecommand{\rhchXgbBarAddEwPct}{}\renewcommand{\rhchXgbBarAddEwPct}{$-$3.1}  % xgb all add_ew vs xgb|har: %
\providecommand{\rhchXgbBarAddEwP}{}\renewcommand{\rhchXgbBarAddEwP}{\ensuremath{<}0.001}  % xgb all add_ew: DM p
\providecommand{\rhchXgbBarLoEwPct}{}\renewcommand{\rhchXgbBarLoEwPct}{$-$0.1}  % xgb all lo_ew vs xgb|full_pre: % [pre-patch]
\providecommand{\rhchXgbBarLoEwP}{}\renewcommand{\rhchXgbBarLoEwP}{0.220}  % xgb all lo_ew: DM p [pre-patch]
\providecommand{\rhchXgbBarAddVwPct}{}\renewcommand{\rhchXgbBarAddVwPct}{$-$3.3}  % xgb all add_vw vs xgb|har: %
\providecommand{\rhchXgbBarAddVwP}{}\renewcommand{\rhchXgbBarAddVwP}{\ensuremath{<}0.001}  % xgb all add_vw: DM p
\providecommand{\rhchXgbBarLoVwPct}{}\renewcommand{\rhchXgbBarLoVwPct}{0.1}  % xgb all lo_vw vs xgb|full_pre: % [pre-patch]
\providecommand{\rhchXgbBarLoVwP}{}\renewcommand{\rhchXgbBarLoVwP}{0.349}  % xgb all lo_vw: DM p [pre-patch]
\providecommand{\rhchXgbBarAddSentPct}{}\renewcommand{\rhchXgbBarAddSentPct}{$-$1.7}  % xgb all add_sent vs xgb|har: %
\providecommand{\rhchXgbBarAddSentP}{}\renewcommand{\rhchXgbBarAddSentP}{\ensuremath{<}0.001}  % xgb all add_sent: DM p
\providecommand{\rhchXgbBarLoSentPct}{}\renewcommand{\rhchXgbBarLoSentPct}{0.2}  % xgb all lo_sent vs xgb|full_pre: % [pre-patch]
\providecommand{\rhchXgbBarLoSentP}{}\renewcommand{\rhchXgbBarLoSentP}{0.066}  % xgb all lo_sent: DM p [pre-patch]
\providecommand{\rhchXgbBarAddIvPct}{}\renewcommand{\rhchXgbBarAddIvPct}{$-$1.6}  % xgb all add_iv vs xgb|har: %
\providecommand{\rhchXgbBarAddIvP}{}\renewcommand{\rhchXgbBarAddIvP}{\ensuremath{<}0.001}  % xgb all add_iv: DM p
\providecommand{\rhchXgbBarLoIvPct}{}\renewcommand{\rhchXgbBarLoIvPct}{0.1}  % xgb all lo_iv vs xgb|full_pre: % [pre-patch]
\providecommand{\rhchXgbBarLoIvP}{}\renewcommand{\rhchXgbBarLoIvP}{0.414}  % xgb all lo_iv: DM p [pre-patch]
\providecommand{\rhchXgbBarAddVdPct}{}\renewcommand{\rhchXgbBarAddVdPct}{$-$1.2}  % xgb all add_vd vs xgb|har: %
\providecommand{\rhchXgbBarAddVdP}{}\renewcommand{\rhchXgbBarAddVdP}{\ensuremath{<}0.001}  % xgb all add_vd: DM p
\providecommand{\rhchXgbBarLoVdPct}{}\renewcommand{\rhchXgbBarLoVdPct}{0.1}  % xgb all lo_vd vs xgb|full_pre: % [pre-patch]
\providecommand{\rhchXgbBarLoVdP}{}\renewcommand{\rhchXgbBarLoVdP}{0.674}  % xgb all lo_vd: DM p [pre-patch]
\providecommand{\rhchXgbBarAddFomcPct}{}\renewcommand{\rhchXgbBarAddFomcPct}{$-$9.5}  % xgb all add_fomc vs xgb|har: %
\providecommand{\rhchXgbBarAddFomcP}{}\renewcommand{\rhchXgbBarAddFomcP}{\ensuremath{<}0.001}  % xgb all add_fomc: DM p
\providecommand{\rhchXgbBarLoFomcPct}{}\renewcommand{\rhchXgbBarLoFomcPct}{10.3}  % xgb all lo_fomc vs xgb|full: %
\providecommand{\rhchXgbBarLoFomcP}{}\renewcommand{\rhchXgbBarLoFomcP}{\ensuremath{<}0.001}  % xgb all lo_fomc: DM p
\providecommand{\rhchXgbEventJointPct}{}\renewcommand{\rhchXgbEventJointPct}{$-$63.9}  % xgb event: full vs HAR, % QLIKE change
\providecommand{\rhchXgbEventJointDm}{}\renewcommand{\rhchXgbEventJointDm}{$-$7.75}  % xgb event: full vs HAR, DM
\providecommand{\rhchXgbEventJointP}{}\renewcommand{\rhchXgbEventJointP}{\ensuremath{<}0.001}  % xgb event: full vs HAR, DM p
\providecommand{\rhchXgbEventAddRmPct}{}\renewcommand{\rhchXgbEventAddRmPct}{5.5}  % xgb event add_rm vs xgb|har: %
\providecommand{\rhchXgbEventAddRmP}{}\renewcommand{\rhchXgbEventAddRmP}{0.014}  % xgb event add_rm: DM p
\providecommand{\rhchXgbEventLoRmPct}{}\renewcommand{\rhchXgbEventLoRmPct}{$-$1.3}  % xgb event lo_rm vs xgb|full_pre: % [pre-patch]
\providecommand{\rhchXgbEventLoRmP}{}\renewcommand{\rhchXgbEventLoRmP}{0.170}  % xgb event lo_rm: DM p [pre-patch]
\providecommand{\rhchXgbEventAddLiqPct}{}\renewcommand{\rhchXgbEventAddLiqPct}{0.4}  % xgb event add_liq vs xgb|har: %
\providecommand{\rhchXgbEventAddLiqP}{}\renewcommand{\rhchXgbEventAddLiqP}{0.771}  % xgb event add_liq: DM p
\providecommand{\rhchXgbEventLoLiqPct}{}\renewcommand{\rhchXgbEventLoLiqPct}{$-$0.9}  % xgb event lo_liq vs xgb|full_pre: % [pre-patch]
\providecommand{\rhchXgbEventLoLiqP}{}\renewcommand{\rhchXgbEventLoLiqP}{0.069}  % xgb event lo_liq: DM p [pre-patch]
\providecommand{\rhchXgbEventAddEwPct}{}\renewcommand{\rhchXgbEventAddEwPct}{$-$3.6}  % xgb event add_ew vs xgb|har: %
\providecommand{\rhchXgbEventAddEwP}{}\renewcommand{\rhchXgbEventAddEwP}{0.044}  % xgb event add_ew: DM p
\providecommand{\rhchXgbEventLoEwPct}{}\renewcommand{\rhchXgbEventLoEwPct}{0.0}  % xgb event lo_ew vs xgb|full_pre: % [pre-patch]
\providecommand{\rhchXgbEventLoEwP}{}\renewcommand{\rhchXgbEventLoEwP}{0.990}  % xgb event lo_ew: DM p [pre-patch]
\providecommand{\rhchXgbEventAddVwPct}{}\renewcommand{\rhchXgbEventAddVwPct}{$-$3.3}  % xgb event add_vw vs xgb|har: %
\providecommand{\rhchXgbEventAddVwP}{}\renewcommand{\rhchXgbEventAddVwP}{0.063}  % xgb event add_vw: DM p
\providecommand{\rhchXgbEventLoVwPct}{}\renewcommand{\rhchXgbEventLoVwPct}{$-$0.7}  % xgb event lo_vw vs xgb|full_pre: % [pre-patch]
\providecommand{\rhchXgbEventLoVwP}{}\renewcommand{\rhchXgbEventLoVwP}{0.304}  % xgb event lo_vw: DM p [pre-patch]
\providecommand{\rhchXgbEventAddSentPct}{}\renewcommand{\rhchXgbEventAddSentPct}{0.3}  % xgb event add_sent vs xgb|har: %
\providecommand{\rhchXgbEventAddSentP}{}\renewcommand{\rhchXgbEventAddSentP}{0.765}  % xgb event add_sent: DM p
\providecommand{\rhchXgbEventLoSentPct}{}\renewcommand{\rhchXgbEventLoSentPct}{0.0}  % xgb event lo_sent vs xgb|full_pre: % [pre-patch]
\providecommand{\rhchXgbEventLoSentP}{}\renewcommand{\rhchXgbEventLoSentP}{0.994}  % xgb event lo_sent: DM p [pre-patch]
\providecommand{\rhchXgbEventAddIvPct}{}\renewcommand{\rhchXgbEventAddIvPct}{$-$0.7}  % xgb event add_iv vs xgb|har: %
\providecommand{\rhchXgbEventAddIvP}{}\renewcommand{\rhchXgbEventAddIvP}{0.371}  % xgb event add_iv: DM p
\providecommand{\rhchXgbEventLoIvPct}{}\renewcommand{\rhchXgbEventLoIvPct}{0.4}  % xgb event lo_iv vs xgb|full_pre: % [pre-patch]
\providecommand{\rhchXgbEventLoIvP}{}\renewcommand{\rhchXgbEventLoIvP}{0.482}  % xgb event lo_iv: DM p [pre-patch]
\providecommand{\rhchXgbEventAddVdPct}{}\renewcommand{\rhchXgbEventAddVdPct}{$-$1.9}  % xgb event add_vd vs xgb|har: %
\providecommand{\rhchXgbEventAddVdP}{}\renewcommand{\rhchXgbEventAddVdP}{0.076}  % xgb event add_vd: DM p
\providecommand{\rhchXgbEventLoVdPct}{}\renewcommand{\rhchXgbEventLoVdPct}{$-$0.7}  % xgb event lo_vd vs xgb|full_pre: % [pre-patch]
\providecommand{\rhchXgbEventLoVdP}{}\renewcommand{\rhchXgbEventLoVdP}{0.148}  % xgb event lo_vd: DM p [pre-patch]
\providecommand{\rhchXgbEventAddFomcPct}{}\renewcommand{\rhchXgbEventAddFomcPct}{$-$64.0}  % xgb event add_fomc vs xgb|har: %
\providecommand{\rhchXgbEventAddFomcP}{}\renewcommand{\rhchXgbEventAddFomcP}{\ensuremath{<}0.001}  % xgb event add_fomc: DM p
\providecommand{\rhchXgbEventLoFomcPct}{}\renewcommand{\rhchXgbEventLoFomcPct}{173.0}  % xgb event lo_fomc vs xgb|full: %
\providecommand{\rhchXgbEventLoFomcP}{}\renewcommand{\rhchXgbEventLoFomcP}{\ensuremath{<}0.001}  % xgb event lo_fomc: DM p
\providecommand{\rhchXgbEventReading}{}\renewcommand{\rhchXgbEventReading}{concentrated}  % xgb event: item 3 reading
\providecommand{\rhchXgbEventTopAddShare}{}\renewcommand{\rhchXgbEventTopAddShare}{100}  % xgb event: largest add-alone gain, % of the joint gain (fomc)
\providecommand{\rhchXgbEventSumAddRatio}{}\renewcommand{\rhchXgbEventSumAddRatio}{1.05}  % xgb event: add-alone gains summed over groups / joint gain
\providecommand{\rhchXgbEventNLoSig}{}\renewcommand{\rhchXgbEventNLoSig}{1}  % xgb event: significant leave-one-outs
\providecommand{\rhchXgbEventNLoRun}{}\renewcommand{\rhchXgbEventNLoRun}{8}  % xgb event: leave-one-outs available (of 8)
\providecommand{\rhchXgbEventConcTest}{}\renewcommand{\rhchXgbEventConcTest}{yes}  % xgb event: the CONC definition on its own (pending: undetermined)
\providecommand{\rhchXgbEventRedTest}{}\renewcommand{\rhchXgbEventRedTest}{no}  % xgb event: the RED definition on its own (pending: undetermined)
\providecommand{\rhchXgbEventReadingRedFirst}{}\renewcommand{\rhchXgbEventReadingRedFirst}{concentrated}  % xgb event: reading with REDUNDANT tested before CONCENTRATED (the RULE tests CONCENTRATED first)
\providecommand{\rhchXgbOrdJointPct}{}\renewcommand{\rhchXgbOrdJointPct}{$-$6.1}  % xgb ord: full vs HAR, % QLIKE change
\providecommand{\rhchXgbOrdJointDm}{}\renewcommand{\rhchXgbOrdJointDm}{$-$10.42}  % xgb ord: full vs HAR, DM
\providecommand{\rhchXgbOrdJointP}{}\renewcommand{\rhchXgbOrdJointP}{\ensuremath{<}0.001}  % xgb ord: full vs HAR, DM p
\providecommand{\rhchXgbOrdAddRmPct}{}\renewcommand{\rhchXgbOrdAddRmPct}{$-$4.9}  % xgb ord add_rm vs xgb|har: %
\providecommand{\rhchXgbOrdAddRmP}{}\renewcommand{\rhchXgbOrdAddRmP}{\ensuremath{<}0.001}  % xgb ord add_rm: DM p
\providecommand{\rhchXgbOrdLoRmPct}{}\renewcommand{\rhchXgbOrdLoRmPct}{1.8}  % xgb ord lo_rm vs xgb|full_pre: % [pre-patch]
\providecommand{\rhchXgbOrdLoRmP}{}\renewcommand{\rhchXgbOrdLoRmP}{\ensuremath{<}0.001}  % xgb ord lo_rm: DM p [pre-patch]
\providecommand{\rhchXgbOrdAddLiqPct}{}\renewcommand{\rhchXgbOrdAddLiqPct}{$-$4.1}  % xgb ord add_liq vs xgb|har: %
\providecommand{\rhchXgbOrdAddLiqP}{}\renewcommand{\rhchXgbOrdAddLiqP}{\ensuremath{<}0.001}  % xgb ord add_liq: DM p
\providecommand{\rhchXgbOrdLoLiqPct}{}\renewcommand{\rhchXgbOrdLoLiqPct}{0.2}  % xgb ord lo_liq vs xgb|full_pre: % [pre-patch]
\providecommand{\rhchXgbOrdLoLiqP}{}\renewcommand{\rhchXgbOrdLoLiqP}{0.095}  % xgb ord lo_liq: DM p [pre-patch]
\providecommand{\rhchXgbOrdAddEwPct}{}\renewcommand{\rhchXgbOrdAddEwPct}{$-$3.0}  % xgb ord add_ew vs xgb|har: %
\providecommand{\rhchXgbOrdAddEwP}{}\renewcommand{\rhchXgbOrdAddEwP}{\ensuremath{<}0.001}  % xgb ord add_ew: DM p
\providecommand{\rhchXgbOrdLoEwPct}{}\renewcommand{\rhchXgbOrdLoEwPct}{$-$0.1}  % xgb ord lo_ew vs xgb|full_pre: % [pre-patch]
\providecommand{\rhchXgbOrdLoEwP}{}\renewcommand{\rhchXgbOrdLoEwP}{0.172}  % xgb ord lo_ew: DM p [pre-patch]
\providecommand{\rhchXgbOrdAddVwPct}{}\renewcommand{\rhchXgbOrdAddVwPct}{$-$3.3}  % xgb ord add_vw vs xgb|har: %
\providecommand{\rhchXgbOrdAddVwP}{}\renewcommand{\rhchXgbOrdAddVwP}{\ensuremath{<}0.001}  % xgb ord add_vw: DM p
\providecommand{\rhchXgbOrdLoVwPct}{}\renewcommand{\rhchXgbOrdLoVwPct}{0.2}  % xgb ord lo_vw vs xgb|full_pre: % [pre-patch]
\providecommand{\rhchXgbOrdLoVwP}{}\renewcommand{\rhchXgbOrdLoVwP}{0.133}  % xgb ord lo_vw: DM p [pre-patch]
\providecommand{\rhchXgbOrdAddSentPct}{}\renewcommand{\rhchXgbOrdAddSentPct}{$-$2.0}  % xgb ord add_sent vs xgb|har: %
\providecommand{\rhchXgbOrdAddSentP}{}\renewcommand{\rhchXgbOrdAddSentP}{\ensuremath{<}0.001}  % xgb ord add_sent: DM p
\providecommand{\rhchXgbOrdLoSentPct}{}\renewcommand{\rhchXgbOrdLoSentPct}{0.2}  % xgb ord lo_sent vs xgb|full_pre: % [pre-patch]
\providecommand{\rhchXgbOrdLoSentP}{}\renewcommand{\rhchXgbOrdLoSentP}{0.051}  % xgb ord lo_sent: DM p [pre-patch]
\providecommand{\rhchXgbOrdAddIvPct}{}\renewcommand{\rhchXgbOrdAddIvPct}{$-$1.7}  % xgb ord add_iv vs xgb|har: %
\providecommand{\rhchXgbOrdAddIvP}{}\renewcommand{\rhchXgbOrdAddIvP}{\ensuremath{<}0.001}  % xgb ord add_iv: DM p
\providecommand{\rhchXgbOrdLoIvPct}{}\renewcommand{\rhchXgbOrdLoIvPct}{0.0}  % xgb ord lo_iv vs xgb|full_pre: % [pre-patch]
\providecommand{\rhchXgbOrdLoIvP}{}\renewcommand{\rhchXgbOrdLoIvP}{0.611}  % xgb ord lo_iv: DM p [pre-patch]
\providecommand{\rhchXgbOrdAddVdPct}{}\renewcommand{\rhchXgbOrdAddVdPct}{$-$1.1}  % xgb ord add_vd vs xgb|har: %
\providecommand{\rhchXgbOrdAddVdP}{}\renewcommand{\rhchXgbOrdAddVdP}{\ensuremath{<}0.001}  % xgb ord add_vd: DM p
\providecommand{\rhchXgbOrdLoVdPct}{}\renewcommand{\rhchXgbOrdLoVdPct}{0.1}  % xgb ord lo_vd vs xgb|full_pre: % [pre-patch]
\providecommand{\rhchXgbOrdLoVdP}{}\renewcommand{\rhchXgbOrdLoVdP}{0.393}  % xgb ord lo_vd: DM p [pre-patch]
\providecommand{\rhchXgbOrdAddFomcPct}{}\renewcommand{\rhchXgbOrdAddFomcPct}{$-$0.1}  % xgb ord add_fomc vs xgb|har: %
\providecommand{\rhchXgbOrdAddFomcP}{}\renewcommand{\rhchXgbOrdAddFomcP}{0.726}  % xgb ord add_fomc: DM p
\providecommand{\rhchXgbOrdLoFomcPct}{}\renewcommand{\rhchXgbOrdLoFomcPct}{$-$0.4}  % xgb ord lo_fomc vs xgb|full: %
\providecommand{\rhchXgbOrdLoFomcP}{}\renewcommand{\rhchXgbOrdLoFomcP}{0.013}  % xgb ord lo_fomc: DM p
\providecommand{\rhchXgbOrdReading}{}\renewcommand{\rhchXgbOrdReading}{concentrated}  % xgb ord: item 3 reading
\providecommand{\rhchXgbOrdTopAddShare}{}\renewcommand{\rhchXgbOrdTopAddShare}{81}  % xgb ord: largest add-alone gain, % of the joint gain (rm)
\providecommand{\rhchXgbOrdSumAddRatio}{}\renewcommand{\rhchXgbOrdSumAddRatio}{3.32}  % xgb ord: add-alone gains summed over groups / joint gain
\providecommand{\rhchXgbOrdNLoSig}{}\renewcommand{\rhchXgbOrdNLoSig}{1}  % xgb ord: significant leave-one-outs
\providecommand{\rhchXgbOrdNLoRun}{}\renewcommand{\rhchXgbOrdNLoRun}{8}  % xgb ord: leave-one-outs available (of 8)
\providecommand{\rhchXgbOrdConcTest}{}\renewcommand{\rhchXgbOrdConcTest}{yes}  % xgb ord: the CONC definition on its own (pending: undetermined)
\providecommand{\rhchXgbOrdRedTest}{}\renewcommand{\rhchXgbOrdRedTest}{yes}  % xgb ord: the RED definition on its own (pending: undetermined)
\providecommand{\rhchXgbOrdReadingRedFirst}{}\renewcommand{\rhchXgbOrdReadingRedFirst}{redundant}  % xgb ord: reading with REDUNDANT tested before CONCENTRATED (the RULE tests CONCENTRATED first)
\providecommand{\rhchRfBarJointPct}{}\renewcommand{\rhchRfBarJointPct}{$-$11.3}  % rf all: full vs HAR, % QLIKE change
\providecommand{\rhchRfBarJointDm}{}\renewcommand{\rhchRfBarJointDm}{$-$9.74}  % rf all: full vs HAR, DM
\providecommand{\rhchRfBarJointP}{}\renewcommand{\rhchRfBarJointP}{\ensuremath{<}0.001}  % rf all: full vs HAR, DM p
\providecommand{\rhchRfBarAddRmPct}{}\renewcommand{\rhchRfBarAddRmPct}{$-$3.2}  % rf all add_rm vs rf|har: %
\providecommand{\rhchRfBarAddRmP}{}\renewcommand{\rhchRfBarAddRmP}{\ensuremath{<}0.001}  % rf all add_rm: DM p
\providecommand{\rhchRfBarLoRmPct}{}\renewcommand{\rhchRfBarLoRmPct}{1.5}  % rf all lo_rm vs rf|full_pre: % [pre-patch]
\providecommand{\rhchRfBarLoRmP}{}\renewcommand{\rhchRfBarLoRmP}{\ensuremath{<}0.001}  % rf all lo_rm: DM p [pre-patch]
\providecommand{\rhchRfBarAddLiqPct}{}\renewcommand{\rhchRfBarAddLiqPct}{$-$2.4}  % rf all add_liq vs rf|har: %
\providecommand{\rhchRfBarAddLiqP}{}\renewcommand{\rhchRfBarAddLiqP}{\ensuremath{<}0.001}  % rf all add_liq: DM p
\providecommand{\rhchRfBarAddEwPct}{}\renewcommand{\rhchRfBarAddEwPct}{$-$2.3}  % rf all add_ew vs rf|har: %
\providecommand{\rhchRfBarAddEwP}{}\renewcommand{\rhchRfBarAddEwP}{\ensuremath{<}0.001}  % rf all add_ew: DM p
\providecommand{\rhchRfBarAddVwPct}{}\renewcommand{\rhchRfBarAddVwPct}{$-$2.5}  % rf all add_vw vs rf|har: %
\providecommand{\rhchRfBarAddVwP}{}\renewcommand{\rhchRfBarAddVwP}{\ensuremath{<}0.001}  % rf all add_vw: DM p
\providecommand{\rhchRfBarAddSentPct}{}\renewcommand{\rhchRfBarAddSentPct}{$-$0.9}  % rf all add_sent vs rf|har: %
\providecommand{\rhchRfBarAddSentP}{}\renewcommand{\rhchRfBarAddSentP}{\ensuremath{<}0.001}  % rf all add_sent: DM p
\providecommand{\rhchRfBarAddIvPct}{}\renewcommand{\rhchRfBarAddIvPct}{$-$0.2}  % rf all add_iv vs rf|har: %
\providecommand{\rhchRfBarAddIvP}{}\renewcommand{\rhchRfBarAddIvP}{0.601}  % rf all add_iv: DM p
\providecommand{\rhchRfBarAddVdPct}{}\renewcommand{\rhchRfBarAddVdPct}{$-$0.6}  % rf all add_vd vs rf|har: %
\providecommand{\rhchRfBarAddVdP}{}\renewcommand{\rhchRfBarAddVdP}{\ensuremath{<}0.001}  % rf all add_vd: DM p
\providecommand{\rhchRfBarAddFomcPct}{}\renewcommand{\rhchRfBarAddFomcPct}{$-$7.6}  % rf all add_fomc vs rf|har: %
\providecommand{\rhchRfBarAddFomcP}{}\renewcommand{\rhchRfBarAddFomcP}{\ensuremath{<}0.001}  % rf all add_fomc: DM p
\providecommand{\rhchRfEventJointPct}{}\renewcommand{\rhchRfEventJointPct}{$-$44.8}  % rf event: full vs HAR, % QLIKE change
\providecommand{\rhchRfEventJointDm}{}\renewcommand{\rhchRfEventJointDm}{$-$8.56}  % rf event: full vs HAR, DM
\providecommand{\rhchRfEventJointP}{}\renewcommand{\rhchRfEventJointP}{\ensuremath{<}0.001}  % rf event: full vs HAR, DM p
\providecommand{\rhchRfEventAddRmPct}{}\renewcommand{\rhchRfEventAddRmPct}{5.3}  % rf event add_rm vs rf|har: %
\providecommand{\rhchRfEventAddRmP}{}\renewcommand{\rhchRfEventAddRmP}{0.061}  % rf event add_rm: DM p
\providecommand{\rhchRfEventLoRmPct}{}\renewcommand{\rhchRfEventLoRmPct}{$-$2.1}  % rf event lo_rm vs rf|full_pre: % [pre-patch]
\providecommand{\rhchRfEventLoRmP}{}\renewcommand{\rhchRfEventLoRmP}{0.036}  % rf event lo_rm: DM p [pre-patch]
\providecommand{\rhchRfEventAddLiqPct}{}\renewcommand{\rhchRfEventAddLiqPct}{2.0}  % rf event add_liq vs rf|har: %
\providecommand{\rhchRfEventAddLiqP}{}\renewcommand{\rhchRfEventAddLiqP}{0.184}  % rf event add_liq: DM p
\providecommand{\rhchRfEventAddEwPct}{}\renewcommand{\rhchRfEventAddEwPct}{$-$2.4}  % rf event add_ew vs rf|har: %
\providecommand{\rhchRfEventAddEwP}{}\renewcommand{\rhchRfEventAddEwP}{0.073}  % rf event add_ew: DM p
\providecommand{\rhchRfEventAddVwPct}{}\renewcommand{\rhchRfEventAddVwPct}{$-$0.9}  % rf event add_vw vs rf|har: %
\providecommand{\rhchRfEventAddVwP}{}\renewcommand{\rhchRfEventAddVwP}{0.651}  % rf event add_vw: DM p
\providecommand{\rhchRfEventAddSentPct}{}\renewcommand{\rhchRfEventAddSentPct}{0.4}  % rf event add_sent vs rf|har: %
\providecommand{\rhchRfEventAddSentP}{}\renewcommand{\rhchRfEventAddSentP}{0.621}  % rf event add_sent: DM p
\providecommand{\rhchRfEventAddIvPct}{}\renewcommand{\rhchRfEventAddIvPct}{4.1}  % rf event add_iv vs rf|har: %
\providecommand{\rhchRfEventAddIvP}{}\renewcommand{\rhchRfEventAddIvP}{0.001}  % rf event add_iv: DM p
\providecommand{\rhchRfEventAddVdPct}{}\renewcommand{\rhchRfEventAddVdPct}{1.3}  % rf event add_vd vs rf|har: %
\providecommand{\rhchRfEventAddVdP}{}\renewcommand{\rhchRfEventAddVdP}{0.046}  % rf event add_vd: DM p
\providecommand{\rhchRfEventAddFomcPct}{}\renewcommand{\rhchRfEventAddFomcPct}{$-$48.5}  % rf event add_fomc vs rf|har: %
\providecommand{\rhchRfEventAddFomcP}{}\renewcommand{\rhchRfEventAddFomcP}{\ensuremath{<}0.001}  % rf event add_fomc: DM p
\providecommand{\rhchRfEventReading}{}\renewcommand{\rhchRfEventReading}{n/a}  % rf event: item 3 reading
\providecommand{\rhchRfEventTopAddShare}{}\renewcommand{\rhchRfEventTopAddShare}{108}  % rf event: largest add-alone gain, % of the joint gain (fomc)
\providecommand{\rhchRfEventSumAddRatio}{}\renewcommand{\rhchRfEventSumAddRatio}{0.86}  % rf event: add-alone gains summed over groups / joint gain
\providecommand{\rhchRfEventNLoSig}{}\renewcommand{\rhchRfEventNLoSig}{0}  % rf event: significant leave-one-outs
\providecommand{\rhchRfEventNLoRun}{}\renewcommand{\rhchRfEventNLoRun}{1}  % rf event: leave-one-outs available (of 8)
\providecommand{\rhchRfEventConcTest}{}\renewcommand{\rhchRfEventConcTest}{\rhPending{}}  % rf event: the CONC definition on its own (pending: undetermined)
\providecommand{\rhchRfEventRedTest}{}\renewcommand{\rhchRfEventRedTest}{no}  % rf event: the RED definition on its own (pending: undetermined)
\providecommand{\rhchRfEventReadingRedFirst}{}\renewcommand{\rhchRfEventReadingRedFirst}{n/a}  % rf event: reading with REDUNDANT tested before CONCENTRATED (the RULE tests CONCENTRATED first)
\providecommand{\rhchRfOrdJointPct}{}\renewcommand{\rhchRfOrdJointPct}{$-$5.3}  % rf ord: full vs HAR, % QLIKE change
\providecommand{\rhchRfOrdJointDm}{}\renewcommand{\rhchRfOrdJointDm}{$-$11.16}  % rf ord: full vs HAR, DM
\providecommand{\rhchRfOrdJointP}{}\renewcommand{\rhchRfOrdJointP}{\ensuremath{<}0.001}  % rf ord: full vs HAR, DM p
\providecommand{\rhchRfOrdAddRmPct}{}\renewcommand{\rhchRfOrdAddRmPct}{$-$4.7}  % rf ord add_rm vs rf|har: %
\providecommand{\rhchRfOrdAddRmP}{}\renewcommand{\rhchRfOrdAddRmP}{\ensuremath{<}0.001}  % rf ord add_rm: DM p
\providecommand{\rhchRfOrdLoRmPct}{}\renewcommand{\rhchRfOrdLoRmPct}{1.8}  % rf ord lo_rm vs rf|full_pre: % [pre-patch]
\providecommand{\rhchRfOrdLoRmP}{}\renewcommand{\rhchRfOrdLoRmP}{\ensuremath{<}0.001}  % rf ord lo_rm: DM p [pre-patch]
\providecommand{\rhchRfOrdAddLiqPct}{}\renewcommand{\rhchRfOrdAddLiqPct}{$-$3.2}  % rf ord add_liq vs rf|har: %
\providecommand{\rhchRfOrdAddLiqP}{}\renewcommand{\rhchRfOrdAddLiqP}{\ensuremath{<}0.001}  % rf ord add_liq: DM p
\providecommand{\rhchRfOrdAddEwPct}{}\renewcommand{\rhchRfOrdAddEwPct}{$-$2.3}  % rf ord add_ew vs rf|har: %
\providecommand{\rhchRfOrdAddEwP}{}\renewcommand{\rhchRfOrdAddEwP}{\ensuremath{<}0.001}  % rf ord add_ew: DM p
\providecommand{\rhchRfOrdAddVwPct}{}\renewcommand{\rhchRfOrdAddVwPct}{$-$2.8}  % rf ord add_vw vs rf|har: %
\providecommand{\rhchRfOrdAddVwP}{}\renewcommand{\rhchRfOrdAddVwP}{\ensuremath{<}0.001}  % rf ord add_vw: DM p
\providecommand{\rhchRfOrdAddSentPct}{}\renewcommand{\rhchRfOrdAddSentPct}{$-$1.1}  % rf ord add_sent vs rf|har: %
\providecommand{\rhchRfOrdAddSentP}{}\renewcommand{\rhchRfOrdAddSentP}{\ensuremath{<}0.001}  % rf ord add_sent: DM p
\providecommand{\rhchRfOrdAddIvPct}{}\renewcommand{\rhchRfOrdAddIvPct}{$-$1.0}  % rf ord add_iv vs rf|har: %
\providecommand{\rhchRfOrdAddIvP}{}\renewcommand{\rhchRfOrdAddIvP}{0.004}  % rf ord add_iv: DM p
\providecommand{\rhchRfOrdAddVdPct}{}\renewcommand{\rhchRfOrdAddVdPct}{$-$1.0}  % rf ord add_vd vs rf|har: %
\providecommand{\rhchRfOrdAddVdP}{}\renewcommand{\rhchRfOrdAddVdP}{\ensuremath{<}0.001}  % rf ord add_vd: DM p
\providecommand{\rhchRfOrdAddFomcPct}{}\renewcommand{\rhchRfOrdAddFomcPct}{$-$0.3}  % rf ord add_fomc vs rf|har: %
\providecommand{\rhchRfOrdAddFomcP}{}\renewcommand{\rhchRfOrdAddFomcP}{0.141}  % rf ord add_fomc: DM p
\providecommand{\rhchRfOrdReading}{}\renewcommand{\rhchRfOrdReading}{concentrated}  % rf ord: item 3 reading
\providecommand{\rhchRfOrdTopAddShare}{}\renewcommand{\rhchRfOrdTopAddShare}{88}  % rf ord: largest add-alone gain, % of the joint gain (rm)
\providecommand{\rhchRfOrdSumAddRatio}{}\renewcommand{\rhchRfOrdSumAddRatio}{3.07}  % rf ord: add-alone gains summed over groups / joint gain
\providecommand{\rhchRfOrdNLoSig}{}\renewcommand{\rhchRfOrdNLoSig}{1}  % rf ord: significant leave-one-outs
\providecommand{\rhchRfOrdNLoRun}{}\renewcommand{\rhchRfOrdNLoRun}{1}  % rf ord: leave-one-outs available (of 8)
\providecommand{\rhchRfOrdConcTest}{}\renewcommand{\rhchRfOrdConcTest}{yes}  % rf ord: the CONC definition on its own (pending: undetermined)
\providecommand{\rhchRfOrdRedTest}{}\renewcommand{\rhchRfOrdRedTest}{\rhPending{}}  % rf ord: the RED definition on its own (pending: undetermined)
\providecommand{\rhchRfOrdReadingRedFirst}{}\renewcommand{\rhchRfOrdReadingRedFirst}{n/a}  % rf ord: reading with REDUNDANT tested before CONCENTRATED (the RULE tests CONCENTRATED first)
\providecommand{\rhchLinEventReading}{}\renewcommand{\rhchLinEventReading}{concentrated}  % lin event: item 3 reading by model type
\providecommand{\rhchLinEventReadingRedFirst}{}\renewcommand{\rhchLinEventReadingRedFirst}{concentrated}  % lin event: by model type, REDUNDANT tested first
\providecommand{\rhchLinEventNConcTest}{}\renewcommand{\rhchLinEventNConcTest}{3}  % lin event: classes meeting the Conc definition (of 3)
\providecommand{\rhchLinEventNRedTest}{}\renewcommand{\rhchLinEventNRedTest}{0}  % lin event: classes meeting the Red definition (of 3)
\providecommand{\rhchTreeEventReading}{}\renewcommand{\rhchTreeEventReading}{concentrated}  % tree event: item 3 reading by model type
\providecommand{\rhchTreeEventReadingRedFirst}{}\renewcommand{\rhchTreeEventReadingRedFirst}{concentrated}  % tree event: by model type, REDUNDANT tested first
\providecommand{\rhchTreeEventNConcTest}{}\renewcommand{\rhchTreeEventNConcTest}{2}  % tree event: classes meeting the Conc definition (of 3)
\providecommand{\rhchTreeEventNRedTest}{}\renewcommand{\rhchTreeEventNRedTest}{0}  % tree event: classes meeting the Red definition (of 3)
\providecommand{\rhchLinOrdReading}{}\renewcommand{\rhchLinOrdReading}{concentrated}  % lin ord: item 3 reading by model type
\providecommand{\rhchLinOrdReadingRedFirst}{}\renewcommand{\rhchLinOrdReadingRedFirst}{redundant}  % lin ord: by model type, REDUNDANT tested first
\providecommand{\rhchLinOrdNConcTest}{}\renewcommand{\rhchLinOrdNConcTest}{3}  % lin ord: classes meeting the Conc definition (of 3)
\providecommand{\rhchLinOrdNRedTest}{}\renewcommand{\rhchLinOrdNRedTest}{3}  % lin ord: classes meeting the Red definition (of 3)
\providecommand{\rhchTreeOrdReading}{}\renewcommand{\rhchTreeOrdReading}{concentrated}  % tree ord: item 3 reading by model type
\providecommand{\rhchTreeOrdReadingRedFirst}{}\renewcommand{\rhchTreeOrdReadingRedFirst}{split}  % tree ord: by model type, REDUNDANT tested first
\providecommand{\rhchTreeOrdNConcTest}{}\renewcommand{\rhchTreeOrdNConcTest}{3}  % tree ord: classes meeting the Conc definition (of 3)
\providecommand{\rhchTreeOrdNRedTest}{}\renewcommand{\rhchTreeOrdNRedTest}{1}  % tree ord: classes meeting the Red definition (of 3)
\providecommand{\rhchXRidgeBarFullVsRmPct}{}\renewcommand{\rhchXRidgeBarFullVsRmPct}{$-$3.5}  % ADDITIONAL, NOT PRE-REGISTERED: ridge all: full vs HAR+RM: % QLIKE change
\providecommand{\rhchXRidgeBarFullVsRmDm}{}\renewcommand{\rhchXRidgeBarFullVsRmDm}{$-$4.52}  % ADDITIONAL, NOT PRE-REGISTERED: ridge all: full vs HAR+RM: DM
\providecommand{\rhchXRidgeBarFullVsRmP}{}\renewcommand{\rhchXRidgeBarFullVsRmP}{\ensuremath{<}0.001}  % ADDITIONAL, NOT PRE-REGISTERED: ridge all: full vs HAR+RM: DM p
\providecommand{\rhchXRidgeBarFullVsRmShare}{}\renewcommand{\rhchXRidgeBarFullVsRmShare}{56}  % ADDITIONAL, NOT PRE-REGISTERED: ridge all: full vs HAR+RM: % of the joint gain
\providecommand{\rhchXRidgeBarRmVsHarPct}{}\renewcommand{\rhchXRidgeBarRmVsHarPct}{$-$2.6}  % ADDITIONAL, NOT PRE-REGISTERED: ridge all: HAR+RM vs HAR: % QLIKE change
\providecommand{\rhchXRidgeBarRmVsHarDm}{}\renewcommand{\rhchXRidgeBarRmVsHarDm}{$-$4.14}  % ADDITIONAL, NOT PRE-REGISTERED: ridge all: HAR+RM vs HAR: DM
\providecommand{\rhchXRidgeBarRmVsHarP}{}\renewcommand{\rhchXRidgeBarRmVsHarP}{\ensuremath{<}0.001}  % ADDITIONAL, NOT PRE-REGISTERED: ridge all: HAR+RM vs HAR: DM p
\providecommand{\rhchXRidgeBarRmVsHarShare}{}\renewcommand{\rhchXRidgeBarRmVsHarShare}{44}  % ADDITIONAL, NOT PRE-REGISTERED: ridge all: HAR+RM vs HAR: % of the joint gain
\providecommand{\rhchXRidgeEventFullVsRmPct}{}\renewcommand{\rhchXRidgeEventFullVsRmPct}{$-$25.7}  % ADDITIONAL, NOT PRE-REGISTERED: ridge event: full vs HAR+RM: % QLIKE change
\providecommand{\rhchXRidgeEventFullVsRmDm}{}\renewcommand{\rhchXRidgeEventFullVsRmDm}{$-$5.37}  % ADDITIONAL, NOT PRE-REGISTERED: ridge event: full vs HAR+RM: DM
\providecommand{\rhchXRidgeEventFullVsRmP}{}\renewcommand{\rhchXRidgeEventFullVsRmP}{\ensuremath{<}0.001}  % ADDITIONAL, NOT PRE-REGISTERED: ridge event: full vs HAR+RM: DM p
\providecommand{\rhchXRidgeEventFullVsRmShare}{}\renewcommand{\rhchXRidgeEventFullVsRmShare}{120}  % ADDITIONAL, NOT PRE-REGISTERED: ridge event: full vs HAR+RM: % of the joint gain
\providecommand{\rhchXRidgeEventRmVsHarPct}{}\renewcommand{\rhchXRidgeEventRmVsHarPct}{4.4}  % ADDITIONAL, NOT PRE-REGISTERED: ridge event: HAR+RM vs HAR: % QLIKE change
\providecommand{\rhchXRidgeEventRmVsHarDm}{}\renewcommand{\rhchXRidgeEventRmVsHarDm}{2.28}  % ADDITIONAL, NOT PRE-REGISTERED: ridge event: HAR+RM vs HAR: DM
\providecommand{\rhchXRidgeEventRmVsHarP}{}\renewcommand{\rhchXRidgeEventRmVsHarP}{0.023}  % ADDITIONAL, NOT PRE-REGISTERED: ridge event: HAR+RM vs HAR: DM p
\providecommand{\rhchXRidgeEventRmVsHarShare}{}\renewcommand{\rhchXRidgeEventRmVsHarShare}{$-$20}  % ADDITIONAL, NOT PRE-REGISTERED: ridge event: HAR+RM vs HAR: % of the joint gain
\providecommand{\rhchXRidgeOrdFullVsRmPct}{}\renewcommand{\rhchXRidgeOrdFullVsRmPct}{0.6}  % ADDITIONAL, NOT PRE-REGISTERED: ridge ord: full vs HAR+RM: % QLIKE change
\providecommand{\rhchXRidgeOrdFullVsRmDm}{}\renewcommand{\rhchXRidgeOrdFullVsRmDm}{1.24}  % ADDITIONAL, NOT PRE-REGISTERED: ridge ord: full vs HAR+RM: DM
\providecommand{\rhchXRidgeOrdFullVsRmP}{}\renewcommand{\rhchXRidgeOrdFullVsRmP}{0.215}  % ADDITIONAL, NOT PRE-REGISTERED: ridge ord: full vs HAR+RM: DM p
\providecommand{\rhchXRidgeOrdFullVsRmShare}{}\renewcommand{\rhchXRidgeOrdFullVsRmShare}{$-$18}  % ADDITIONAL, NOT PRE-REGISTERED: ridge ord: full vs HAR+RM: % of the joint gain
\providecommand{\rhchXRidgeOrdRmVsHarPct}{}\renewcommand{\rhchXRidgeOrdRmVsHarPct}{$-$3.8}  % ADDITIONAL, NOT PRE-REGISTERED: ridge ord: HAR+RM vs HAR: % QLIKE change
\providecommand{\rhchXRidgeOrdRmVsHarDm}{}\renewcommand{\rhchXRidgeOrdRmVsHarDm}{$-$5.48}  % ADDITIONAL, NOT PRE-REGISTERED: ridge ord: HAR+RM vs HAR: DM
\providecommand{\rhchXRidgeOrdRmVsHarP}{}\renewcommand{\rhchXRidgeOrdRmVsHarP}{\ensuremath{<}0.001}  % ADDITIONAL, NOT PRE-REGISTERED: ridge ord: HAR+RM vs HAR: DM p
\providecommand{\rhchXRidgeOrdRmVsHarShare}{}\renewcommand{\rhchXRidgeOrdRmVsHarShare}{118}  % ADDITIONAL, NOT PRE-REGISTERED: ridge ord: HAR+RM vs HAR: % of the joint gain
\providecommand{\rhchXLassoBarFullVsRmPct}{}\renewcommand{\rhchXLassoBarFullVsRmPct}{$-$7.0}  % ADDITIONAL, NOT PRE-REGISTERED: lasso all: full vs HAR+RM: % QLIKE change
\providecommand{\rhchXLassoBarFullVsRmDm}{}\renewcommand{\rhchXLassoBarFullVsRmDm}{$-$5.66}  % ADDITIONAL, NOT PRE-REGISTERED: lasso all: full vs HAR+RM: DM
\providecommand{\rhchXLassoBarFullVsRmP}{}\renewcommand{\rhchXLassoBarFullVsRmP}{\ensuremath{<}0.001}  % ADDITIONAL, NOT PRE-REGISTERED: lasso all: full vs HAR+RM: DM p
\providecommand{\rhchXLassoBarFullVsRmShare}{}\renewcommand{\rhchXLassoBarFullVsRmShare}{76}  % ADDITIONAL, NOT PRE-REGISTERED: lasso all: full vs HAR+RM: % of the joint gain
\providecommand{\rhchXLassoBarRmVsHarPct}{}\renewcommand{\rhchXLassoBarRmVsHarPct}{$-$2.1}  % ADDITIONAL, NOT PRE-REGISTERED: lasso all: HAR+RM vs HAR: % QLIKE change
\providecommand{\rhchXLassoBarRmVsHarDm}{}\renewcommand{\rhchXLassoBarRmVsHarDm}{$-$2.64}  % ADDITIONAL, NOT PRE-REGISTERED: lasso all: HAR+RM vs HAR: DM
\providecommand{\rhchXLassoBarRmVsHarP}{}\renewcommand{\rhchXLassoBarRmVsHarP}{0.008}  % ADDITIONAL, NOT PRE-REGISTERED: lasso all: HAR+RM vs HAR: DM p
\providecommand{\rhchXLassoBarRmVsHarShare}{}\renewcommand{\rhchXLassoBarRmVsHarShare}{24}  % ADDITIONAL, NOT PRE-REGISTERED: lasso all: HAR+RM vs HAR: % of the joint gain
\providecommand{\rhchXLassoEventFullVsRmPct}{}\renewcommand{\rhchXLassoEventFullVsRmPct}{$-$49.5}  % ADDITIONAL, NOT PRE-REGISTERED: lasso event: full vs HAR+RM: % QLIKE change
\providecommand{\rhchXLassoEventFullVsRmDm}{}\renewcommand{\rhchXLassoEventFullVsRmDm}{$-$6.59}  % ADDITIONAL, NOT PRE-REGISTERED: lasso event: full vs HAR+RM: DM
\providecommand{\rhchXLassoEventFullVsRmP}{}\renewcommand{\rhchXLassoEventFullVsRmP}{\ensuremath{<}0.001}  % ADDITIONAL, NOT PRE-REGISTERED: lasso event: full vs HAR+RM: DM p
\providecommand{\rhchXLassoEventFullVsRmShare}{}\renewcommand{\rhchXLassoEventFullVsRmShare}{110}  % ADDITIONAL, NOT PRE-REGISTERED: lasso event: full vs HAR+RM: % of the joint gain
\providecommand{\rhchXLassoEventRmVsHarPct}{}\renewcommand{\rhchXLassoEventRmVsHarPct}{4.7}  % ADDITIONAL, NOT PRE-REGISTERED: lasso event: HAR+RM vs HAR: % QLIKE change
\providecommand{\rhchXLassoEventRmVsHarDm}{}\renewcommand{\rhchXLassoEventRmVsHarDm}{2.25}  % ADDITIONAL, NOT PRE-REGISTERED: lasso event: HAR+RM vs HAR: DM
\providecommand{\rhchXLassoEventRmVsHarP}{}\renewcommand{\rhchXLassoEventRmVsHarP}{0.024}  % ADDITIONAL, NOT PRE-REGISTERED: lasso event: HAR+RM vs HAR: DM p
\providecommand{\rhchXLassoEventRmVsHarShare}{}\renewcommand{\rhchXLassoEventRmVsHarShare}{$-$10}  % ADDITIONAL, NOT PRE-REGISTERED: lasso event: HAR+RM vs HAR: % of the joint gain
\providecommand{\rhchXLassoOrdFullVsRmPct}{}\renewcommand{\rhchXLassoOrdFullVsRmPct}{0.9}  % ADDITIONAL, NOT PRE-REGISTERED: lasso ord: full vs HAR+RM: % QLIKE change
\providecommand{\rhchXLassoOrdFullVsRmDm}{}\renewcommand{\rhchXLassoOrdFullVsRmDm}{1.44}  % ADDITIONAL, NOT PRE-REGISTERED: lasso ord: full vs HAR+RM: DM
\providecommand{\rhchXLassoOrdFullVsRmP}{}\renewcommand{\rhchXLassoOrdFullVsRmP}{0.151}  % ADDITIONAL, NOT PRE-REGISTERED: lasso ord: full vs HAR+RM: DM p
\providecommand{\rhchXLassoOrdFullVsRmShare}{}\renewcommand{\rhchXLassoOrdFullVsRmShare}{$-$36}  % ADDITIONAL, NOT PRE-REGISTERED: lasso ord: full vs HAR+RM: % of the joint gain
\providecommand{\rhchXLassoOrdRmVsHarPct}{}\renewcommand{\rhchXLassoOrdRmVsHarPct}{$-$3.3}  % ADDITIONAL, NOT PRE-REGISTERED: lasso ord: HAR+RM vs HAR: % QLIKE change
\providecommand{\rhchXLassoOrdRmVsHarDm}{}\renewcommand{\rhchXLassoOrdRmVsHarDm}{$-$3.61}  % ADDITIONAL, NOT PRE-REGISTERED: lasso ord: HAR+RM vs HAR: DM
\providecommand{\rhchXLassoOrdRmVsHarP}{}\renewcommand{\rhchXLassoOrdRmVsHarP}{\ensuremath{<}0.001}  % ADDITIONAL, NOT PRE-REGISTERED: lasso ord: HAR+RM vs HAR: DM p
\providecommand{\rhchXLassoOrdRmVsHarShare}{}\renewcommand{\rhchXLassoOrdRmVsHarShare}{136}  % ADDITIONAL, NOT PRE-REGISTERED: lasso ord: HAR+RM vs HAR: % of the joint gain
\providecommand{\rhchXEnetBarFullVsRmPct}{}\renewcommand{\rhchXEnetBarFullVsRmPct}{$-$4.1}  % ADDITIONAL, NOT PRE-REGISTERED: enet all: full vs HAR+RM: % QLIKE change
\providecommand{\rhchXEnetBarFullVsRmDm}{}\renewcommand{\rhchXEnetBarFullVsRmDm}{$-$5.48}  % ADDITIONAL, NOT PRE-REGISTERED: enet all: full vs HAR+RM: DM
\providecommand{\rhchXEnetBarFullVsRmP}{}\renewcommand{\rhchXEnetBarFullVsRmP}{\ensuremath{<}0.001}  % ADDITIONAL, NOT PRE-REGISTERED: enet all: full vs HAR+RM: DM p
\providecommand{\rhchXEnetBarFullVsRmShare}{}\renewcommand{\rhchXEnetBarFullVsRmShare}{60}  % ADDITIONAL, NOT PRE-REGISTERED: enet all: full vs HAR+RM: % of the joint gain
\providecommand{\rhchXEnetBarRmVsHarPct}{}\renewcommand{\rhchXEnetBarRmVsHarPct}{$-$2.7}  % ADDITIONAL, NOT PRE-REGISTERED: enet all: HAR+RM vs HAR: % QLIKE change
\providecommand{\rhchXEnetBarRmVsHarDm}{}\renewcommand{\rhchXEnetBarRmVsHarDm}{$-$4.25}  % ADDITIONAL, NOT PRE-REGISTERED: enet all: HAR+RM vs HAR: DM
\providecommand{\rhchXEnetBarRmVsHarP}{}\renewcommand{\rhchXEnetBarRmVsHarP}{\ensuremath{<}0.001}  % ADDITIONAL, NOT PRE-REGISTERED: enet all: HAR+RM vs HAR: DM p
\providecommand{\rhchXEnetBarRmVsHarShare}{}\renewcommand{\rhchXEnetBarRmVsHarShare}{40}  % ADDITIONAL, NOT PRE-REGISTERED: enet all: HAR+RM vs HAR: % of the joint gain
\providecommand{\rhchXEnetEventFullVsRmPct}{}\renewcommand{\rhchXEnetEventFullVsRmPct}{$-$25.7}  % ADDITIONAL, NOT PRE-REGISTERED: enet event: full vs HAR+RM: % QLIKE change
\providecommand{\rhchXEnetEventFullVsRmDm}{}\renewcommand{\rhchXEnetEventFullVsRmDm}{$-$5.65}  % ADDITIONAL, NOT PRE-REGISTERED: enet event: full vs HAR+RM: DM
\providecommand{\rhchXEnetEventFullVsRmP}{}\renewcommand{\rhchXEnetEventFullVsRmP}{\ensuremath{<}0.001}  % ADDITIONAL, NOT PRE-REGISTERED: enet event: full vs HAR+RM: DM p
\providecommand{\rhchXEnetEventFullVsRmShare}{}\renewcommand{\rhchXEnetEventFullVsRmShare}{118}  % ADDITIONAL, NOT PRE-REGISTERED: enet event: full vs HAR+RM: % of the joint gain
\providecommand{\rhchXEnetEventRmVsHarPct}{}\renewcommand{\rhchXEnetEventRmVsHarPct}{4.2}  % ADDITIONAL, NOT PRE-REGISTERED: enet event: HAR+RM vs HAR: % QLIKE change
\providecommand{\rhchXEnetEventRmVsHarDm}{}\renewcommand{\rhchXEnetEventRmVsHarDm}{2.16}  % ADDITIONAL, NOT PRE-REGISTERED: enet event: HAR+RM vs HAR: DM
\providecommand{\rhchXEnetEventRmVsHarP}{}\renewcommand{\rhchXEnetEventRmVsHarP}{0.031}  % ADDITIONAL, NOT PRE-REGISTERED: enet event: HAR+RM vs HAR: DM p
\providecommand{\rhchXEnetEventRmVsHarShare}{}\renewcommand{\rhchXEnetEventRmVsHarShare}{$-$18}  % ADDITIONAL, NOT PRE-REGISTERED: enet event: HAR+RM vs HAR: % of the joint gain
\providecommand{\rhchXEnetOrdFullVsRmPct}{}\renewcommand{\rhchXEnetOrdFullVsRmPct}{$-$0.1}  % ADDITIONAL, NOT PRE-REGISTERED: enet ord: full vs HAR+RM: % QLIKE change
\providecommand{\rhchXEnetOrdFullVsRmDm}{}\renewcommand{\rhchXEnetOrdFullVsRmDm}{$-$0.26}  % ADDITIONAL, NOT PRE-REGISTERED: enet ord: full vs HAR+RM: DM
\providecommand{\rhchXEnetOrdFullVsRmP}{}\renewcommand{\rhchXEnetOrdFullVsRmP}{0.793}  % ADDITIONAL, NOT PRE-REGISTERED: enet ord: full vs HAR+RM: DM p
\providecommand{\rhchXEnetOrdFullVsRmShare}{}\renewcommand{\rhchXEnetOrdFullVsRmShare}{3}  % ADDITIONAL, NOT PRE-REGISTERED: enet ord: full vs HAR+RM: % of the joint gain
\providecommand{\rhchXEnetOrdRmVsHarPct}{}\renewcommand{\rhchXEnetOrdRmVsHarPct}{$-$3.9}  % ADDITIONAL, NOT PRE-REGISTERED: enet ord: HAR+RM vs HAR: % QLIKE change
\providecommand{\rhchXEnetOrdRmVsHarDm}{}\renewcommand{\rhchXEnetOrdRmVsHarDm}{$-$5.53}  % ADDITIONAL, NOT PRE-REGISTERED: enet ord: HAR+RM vs HAR: DM
\providecommand{\rhchXEnetOrdRmVsHarP}{}\renewcommand{\rhchXEnetOrdRmVsHarP}{\ensuremath{<}0.001}  % ADDITIONAL, NOT PRE-REGISTERED: enet ord: HAR+RM vs HAR: DM p
\providecommand{\rhchXEnetOrdRmVsHarShare}{}\renewcommand{\rhchXEnetOrdRmVsHarShare}{97}  % ADDITIONAL, NOT PRE-REGISTERED: enet ord: HAR+RM vs HAR: % of the joint gain
\providecommand{\rhchXLgbmBarFullVsRmPct}{}\renewcommand{\rhchXLgbmBarFullVsRmPct}{$-$11.5}  % ADDITIONAL, NOT PRE-REGISTERED: lgbm all: full vs HAR+RM: % QLIKE change
\providecommand{\rhchXLgbmBarFullVsRmDm}{}\renewcommand{\rhchXLgbmBarFullVsRmDm}{$-$7.62}  % ADDITIONAL, NOT PRE-REGISTERED: lgbm all: full vs HAR+RM: DM
\providecommand{\rhchXLgbmBarFullVsRmP}{}\renewcommand{\rhchXLgbmBarFullVsRmP}{\ensuremath{<}0.001}  % ADDITIONAL, NOT PRE-REGISTERED: lgbm all: full vs HAR+RM: DM p
\providecommand{\rhchXLgbmBarFullVsRmShare}{}\renewcommand{\rhchXLgbmBarFullVsRmShare}{79}  % ADDITIONAL, NOT PRE-REGISTERED: lgbm all: full vs HAR+RM: % of the joint gain
\providecommand{\rhchXLgbmBarRmVsHarPct}{}\renewcommand{\rhchXLgbmBarRmVsHarPct}{$-$3.0}  % ADDITIONAL, NOT PRE-REGISTERED: lgbm all: HAR+RM vs HAR: % QLIKE change
\providecommand{\rhchXLgbmBarRmVsHarDm}{}\renewcommand{\rhchXLgbmBarRmVsHarDm}{$-$6.32}  % ADDITIONAL, NOT PRE-REGISTERED: lgbm all: HAR+RM vs HAR: DM
\providecommand{\rhchXLgbmBarRmVsHarP}{}\renewcommand{\rhchXLgbmBarRmVsHarP}{\ensuremath{<}0.001}  % ADDITIONAL, NOT PRE-REGISTERED: lgbm all: HAR+RM vs HAR: DM p
\providecommand{\rhchXLgbmBarRmVsHarShare}{}\renewcommand{\rhchXLgbmBarRmVsHarShare}{21}  % ADDITIONAL, NOT PRE-REGISTERED: lgbm all: HAR+RM vs HAR: % of the joint gain
\providecommand{\rhchXLgbmEventFullVsRmPct}{}\renewcommand{\rhchXLgbmEventFullVsRmPct}{$-$64.7}  % ADDITIONAL, NOT PRE-REGISTERED: lgbm event: full vs HAR+RM: % QLIKE change
\providecommand{\rhchXLgbmEventFullVsRmDm}{}\renewcommand{\rhchXLgbmEventFullVsRmDm}{$-$7.56}  % ADDITIONAL, NOT PRE-REGISTERED: lgbm event: full vs HAR+RM: DM
\providecommand{\rhchXLgbmEventFullVsRmP}{}\renewcommand{\rhchXLgbmEventFullVsRmP}{\ensuremath{<}0.001}  % ADDITIONAL, NOT PRE-REGISTERED: lgbm event: full vs HAR+RM: DM p
\providecommand{\rhchXLgbmEventFullVsRmShare}{}\renewcommand{\rhchXLgbmEventFullVsRmShare}{111}  % ADDITIONAL, NOT PRE-REGISTERED: lgbm event: full vs HAR+RM: % of the joint gain
\providecommand{\rhchXLgbmEventRmVsHarPct}{}\renewcommand{\rhchXLgbmEventRmVsHarPct}{6.8}  % ADDITIONAL, NOT PRE-REGISTERED: lgbm event: HAR+RM vs HAR: % QLIKE change
\providecommand{\rhchXLgbmEventRmVsHarDm}{}\renewcommand{\rhchXLgbmEventRmVsHarDm}{2.58}  % ADDITIONAL, NOT PRE-REGISTERED: lgbm event: HAR+RM vs HAR: DM
\providecommand{\rhchXLgbmEventRmVsHarP}{}\renewcommand{\rhchXLgbmEventRmVsHarP}{0.010}  % ADDITIONAL, NOT PRE-REGISTERED: lgbm event: HAR+RM vs HAR: DM p
\providecommand{\rhchXLgbmEventRmVsHarShare}{}\renewcommand{\rhchXLgbmEventRmVsHarShare}{$-$11}  % ADDITIONAL, NOT PRE-REGISTERED: lgbm event: HAR+RM vs HAR: % of the joint gain
\providecommand{\rhchXLgbmOrdFullVsRmPct}{}\renewcommand{\rhchXLgbmOrdFullVsRmPct}{$-$1.5}  % ADDITIONAL, NOT PRE-REGISTERED: lgbm ord: full vs HAR+RM: % QLIKE change
\providecommand{\rhchXLgbmOrdFullVsRmDm}{}\renewcommand{\rhchXLgbmOrdFullVsRmDm}{$-$4.30}  % ADDITIONAL, NOT PRE-REGISTERED: lgbm ord: full vs HAR+RM: DM
\providecommand{\rhchXLgbmOrdFullVsRmP}{}\renewcommand{\rhchXLgbmOrdFullVsRmP}{\ensuremath{<}0.001}  % ADDITIONAL, NOT PRE-REGISTERED: lgbm ord: full vs HAR+RM: DM p
\providecommand{\rhchXLgbmOrdFullVsRmShare}{}\renewcommand{\rhchXLgbmOrdFullVsRmShare}{23}  % ADDITIONAL, NOT PRE-REGISTERED: lgbm ord: full vs HAR+RM: % of the joint gain
\providecommand{\rhchXLgbmOrdRmVsHarPct}{}\renewcommand{\rhchXLgbmOrdRmVsHarPct}{$-$4.6}  % ADDITIONAL, NOT PRE-REGISTERED: lgbm ord: HAR+RM vs HAR: % QLIKE change
\providecommand{\rhchXLgbmOrdRmVsHarDm}{}\renewcommand{\rhchXLgbmOrdRmVsHarDm}{$-$9.55}  % ADDITIONAL, NOT PRE-REGISTERED: lgbm ord: HAR+RM vs HAR: DM
\providecommand{\rhchXLgbmOrdRmVsHarP}{}\renewcommand{\rhchXLgbmOrdRmVsHarP}{\ensuremath{<}0.001}  % ADDITIONAL, NOT PRE-REGISTERED: lgbm ord: HAR+RM vs HAR: DM p
\providecommand{\rhchXLgbmOrdRmVsHarShare}{}\renewcommand{\rhchXLgbmOrdRmVsHarShare}{77}  % ADDITIONAL, NOT PRE-REGISTERED: lgbm ord: HAR+RM vs HAR: % of the joint gain
\providecommand{\rhchXXgbBarFullVsRmPct}{}\renewcommand{\rhchXXgbBarFullVsRmPct}{$-$11.6}  % ADDITIONAL, NOT PRE-REGISTERED: xgb all: full vs HAR+RM: % QLIKE change
\providecommand{\rhchXXgbBarFullVsRmDm}{}\renewcommand{\rhchXXgbBarFullVsRmDm}{$-$7.19}  % ADDITIONAL, NOT PRE-REGISTERED: xgb all: full vs HAR+RM: DM
\providecommand{\rhchXXgbBarFullVsRmP}{}\renewcommand{\rhchXXgbBarFullVsRmP}{\ensuremath{<}0.001}  % ADDITIONAL, NOT PRE-REGISTERED: xgb all: full vs HAR+RM: DM p
\providecommand{\rhchXXgbBarFullVsRmShare}{}\renewcommand{\rhchXXgbBarFullVsRmShare}{77}  % ADDITIONAL, NOT PRE-REGISTERED: xgb all: full vs HAR+RM: % of the joint gain
\providecommand{\rhchXXgbBarRmVsHarPct}{}\renewcommand{\rhchXXgbBarRmVsHarPct}{$-$3.4}  % ADDITIONAL, NOT PRE-REGISTERED: xgb all: HAR+RM vs HAR: % QLIKE change
\providecommand{\rhchXXgbBarRmVsHarDm}{}\renewcommand{\rhchXXgbBarRmVsHarDm}{$-$7.51}  % ADDITIONAL, NOT PRE-REGISTERED: xgb all: HAR+RM vs HAR: DM
\providecommand{\rhchXXgbBarRmVsHarP}{}\renewcommand{\rhchXXgbBarRmVsHarP}{\ensuremath{<}0.001}  % ADDITIONAL, NOT PRE-REGISTERED: xgb all: HAR+RM vs HAR: DM p
\providecommand{\rhchXXgbBarRmVsHarShare}{}\renewcommand{\rhchXXgbBarRmVsHarShare}{23}  % ADDITIONAL, NOT PRE-REGISTERED: xgb all: HAR+RM vs HAR: % of the joint gain
\providecommand{\rhchXXgbEventFullVsRmPct}{}\renewcommand{\rhchXXgbEventFullVsRmPct}{$-$65.8}  % ADDITIONAL, NOT PRE-REGISTERED: xgb event: full vs HAR+RM: % QLIKE change
\providecommand{\rhchXXgbEventFullVsRmDm}{}\renewcommand{\rhchXXgbEventFullVsRmDm}{$-$6.93}  % ADDITIONAL, NOT PRE-REGISTERED: xgb event: full vs HAR+RM: DM
\providecommand{\rhchXXgbEventFullVsRmP}{}\renewcommand{\rhchXXgbEventFullVsRmP}{\ensuremath{<}0.001}  % ADDITIONAL, NOT PRE-REGISTERED: xgb event: full vs HAR+RM: DM p
\providecommand{\rhchXXgbEventFullVsRmShare}{}\renewcommand{\rhchXXgbEventFullVsRmShare}{109}  % ADDITIONAL, NOT PRE-REGISTERED: xgb event: full vs HAR+RM: % of the joint gain
\providecommand{\rhchXXgbEventRmVsHarPct}{}\renewcommand{\rhchXXgbEventRmVsHarPct}{5.5}  % ADDITIONAL, NOT PRE-REGISTERED: xgb event: HAR+RM vs HAR: % QLIKE change
\providecommand{\rhchXXgbEventRmVsHarDm}{}\renewcommand{\rhchXXgbEventRmVsHarDm}{2.46}  % ADDITIONAL, NOT PRE-REGISTERED: xgb event: HAR+RM vs HAR: DM
\providecommand{\rhchXXgbEventRmVsHarP}{}\renewcommand{\rhchXXgbEventRmVsHarP}{0.014}  % ADDITIONAL, NOT PRE-REGISTERED: xgb event: HAR+RM vs HAR: DM p
\providecommand{\rhchXXgbEventRmVsHarShare}{}\renewcommand{\rhchXXgbEventRmVsHarShare}{$-$9}  % ADDITIONAL, NOT PRE-REGISTERED: xgb event: HAR+RM vs HAR: % of the joint gain
\providecommand{\rhchXXgbOrdFullVsRmPct}{}\renewcommand{\rhchXXgbOrdFullVsRmPct}{$-$1.2}  % ADDITIONAL, NOT PRE-REGISTERED: xgb ord: full vs HAR+RM: % QLIKE change
\providecommand{\rhchXXgbOrdFullVsRmDm}{}\renewcommand{\rhchXXgbOrdFullVsRmDm}{$-$3.23}  % ADDITIONAL, NOT PRE-REGISTERED: xgb ord: full vs HAR+RM: DM
\providecommand{\rhchXXgbOrdFullVsRmP}{}\renewcommand{\rhchXXgbOrdFullVsRmP}{0.001}  % ADDITIONAL, NOT PRE-REGISTERED: xgb ord: full vs HAR+RM: DM p
\providecommand{\rhchXXgbOrdFullVsRmShare}{}\renewcommand{\rhchXXgbOrdFullVsRmShare}{19}  % ADDITIONAL, NOT PRE-REGISTERED: xgb ord: full vs HAR+RM: % of the joint gain
\providecommand{\rhchXXgbOrdRmVsHarPct}{}\renewcommand{\rhchXXgbOrdRmVsHarPct}{$-$4.9}  % ADDITIONAL, NOT PRE-REGISTERED: xgb ord: HAR+RM vs HAR: % QLIKE change
\providecommand{\rhchXXgbOrdRmVsHarDm}{}\renewcommand{\rhchXXgbOrdRmVsHarDm}{$-$10.37}  % ADDITIONAL, NOT PRE-REGISTERED: xgb ord: HAR+RM vs HAR: DM
\providecommand{\rhchXXgbOrdRmVsHarP}{}\renewcommand{\rhchXXgbOrdRmVsHarP}{\ensuremath{<}0.001}  % ADDITIONAL, NOT PRE-REGISTERED: xgb ord: HAR+RM vs HAR: DM p
\providecommand{\rhchXXgbOrdRmVsHarShare}{}\renewcommand{\rhchXXgbOrdRmVsHarShare}{81}  % ADDITIONAL, NOT PRE-REGISTERED: xgb ord: HAR+RM vs HAR: % of the joint gain
\providecommand{\rhchXRfBarFullVsRmPct}{}\renewcommand{\rhchXRfBarFullVsRmPct}{$-$8.4}  % ADDITIONAL, NOT PRE-REGISTERED: rf all: full vs HAR+RM: % QLIKE change
\providecommand{\rhchXRfBarFullVsRmDm}{}\renewcommand{\rhchXRfBarFullVsRmDm}{$-$6.48}  % ADDITIONAL, NOT PRE-REGISTERED: rf all: full vs HAR+RM: DM
\providecommand{\rhchXRfBarFullVsRmP}{}\renewcommand{\rhchXRfBarFullVsRmP}{\ensuremath{<}0.001}  % ADDITIONAL, NOT PRE-REGISTERED: rf all: full vs HAR+RM: DM p
\providecommand{\rhchXRfBarFullVsRmShare}{}\renewcommand{\rhchXRfBarFullVsRmShare}{72}  % ADDITIONAL, NOT PRE-REGISTERED: rf all: full vs HAR+RM: % of the joint gain
\providecommand{\rhchXRfBarRmVsHarPct}{}\renewcommand{\rhchXRfBarRmVsHarPct}{$-$3.2}  % ADDITIONAL, NOT PRE-REGISTERED: rf all: HAR+RM vs HAR: % QLIKE change
\providecommand{\rhchXRfBarRmVsHarDm}{}\renewcommand{\rhchXRfBarRmVsHarDm}{$-$6.79}  % ADDITIONAL, NOT PRE-REGISTERED: rf all: HAR+RM vs HAR: DM
\providecommand{\rhchXRfBarRmVsHarP}{}\renewcommand{\rhchXRfBarRmVsHarP}{\ensuremath{<}0.001}  % ADDITIONAL, NOT PRE-REGISTERED: rf all: HAR+RM vs HAR: DM p
\providecommand{\rhchXRfBarRmVsHarShare}{}\renewcommand{\rhchXRfBarRmVsHarShare}{28}  % ADDITIONAL, NOT PRE-REGISTERED: rf all: HAR+RM vs HAR: % of the joint gain
\providecommand{\rhchXRfEventFullVsRmPct}{}\renewcommand{\rhchXRfEventFullVsRmPct}{$-$47.6}  % ADDITIONAL, NOT PRE-REGISTERED: rf event: full vs HAR+RM: % QLIKE change
\providecommand{\rhchXRfEventFullVsRmDm}{}\renewcommand{\rhchXRfEventFullVsRmDm}{$-$7.12}  % ADDITIONAL, NOT PRE-REGISTERED: rf event: full vs HAR+RM: DM
\providecommand{\rhchXRfEventFullVsRmP}{}\renewcommand{\rhchXRfEventFullVsRmP}{\ensuremath{<}0.001}  % ADDITIONAL, NOT PRE-REGISTERED: rf event: full vs HAR+RM: DM p
\providecommand{\rhchXRfEventFullVsRmShare}{}\renewcommand{\rhchXRfEventFullVsRmShare}{112}  % ADDITIONAL, NOT PRE-REGISTERED: rf event: full vs HAR+RM: % of the joint gain
\providecommand{\rhchXRfEventRmVsHarPct}{}\renewcommand{\rhchXRfEventRmVsHarPct}{5.3}  % ADDITIONAL, NOT PRE-REGISTERED: rf event: HAR+RM vs HAR: % QLIKE change
\providecommand{\rhchXRfEventRmVsHarDm}{}\renewcommand{\rhchXRfEventRmVsHarDm}{1.88}  % ADDITIONAL, NOT PRE-REGISTERED: rf event: HAR+RM vs HAR: DM
\providecommand{\rhchXRfEventRmVsHarP}{}\renewcommand{\rhchXRfEventRmVsHarP}{0.061}  % ADDITIONAL, NOT PRE-REGISTERED: rf event: HAR+RM vs HAR: DM p
\providecommand{\rhchXRfEventRmVsHarShare}{}\renewcommand{\rhchXRfEventRmVsHarShare}{$-$12}  % ADDITIONAL, NOT PRE-REGISTERED: rf event: HAR+RM vs HAR: % of the joint gain
\providecommand{\rhchXRfOrdFullVsRmPct}{}\renewcommand{\rhchXRfOrdFullVsRmPct}{$-$0.7}  % ADDITIONAL, NOT PRE-REGISTERED: rf ord: full vs HAR+RM: % QLIKE change
\providecommand{\rhchXRfOrdFullVsRmDm}{}\renewcommand{\rhchXRfOrdFullVsRmDm}{$-$2.22}  % ADDITIONAL, NOT PRE-REGISTERED: rf ord: full vs HAR+RM: DM
\providecommand{\rhchXRfOrdFullVsRmP}{}\renewcommand{\rhchXRfOrdFullVsRmP}{0.026}  % ADDITIONAL, NOT PRE-REGISTERED: rf ord: full vs HAR+RM: DM p
\providecommand{\rhchXRfOrdFullVsRmShare}{}\renewcommand{\rhchXRfOrdFullVsRmShare}{12}  % ADDITIONAL, NOT PRE-REGISTERED: rf ord: full vs HAR+RM: % of the joint gain
\providecommand{\rhchXRfOrdRmVsHarPct}{}\renewcommand{\rhchXRfOrdRmVsHarPct}{$-$4.7}  % ADDITIONAL, NOT PRE-REGISTERED: rf ord: HAR+RM vs HAR: % QLIKE change
\providecommand{\rhchXRfOrdRmVsHarDm}{}\renewcommand{\rhchXRfOrdRmVsHarDm}{$-$10.68}  % ADDITIONAL, NOT PRE-REGISTERED: rf ord: HAR+RM vs HAR: DM
\providecommand{\rhchXRfOrdRmVsHarP}{}\renewcommand{\rhchXRfOrdRmVsHarP}{\ensuremath{<}0.001}  % ADDITIONAL, NOT PRE-REGISTERED: rf ord: HAR+RM vs HAR: DM p
\providecommand{\rhchXRfOrdRmVsHarShare}{}\renewcommand{\rhchXRfOrdRmVsHarShare}{88}  % ADDITIONAL, NOT PRE-REGISTERED: rf ord: HAR+RM vs HAR: % of the joint gain
\providecommand{\rhchXLinBarFullVsRmPctMin}{}\renewcommand{\rhchXLinBarFullVsRmPctMin}{$-$7.0}  % ADDITIONAL: Lin all: lowest %
\providecommand{\rhchXLinBarFullVsRmPctMax}{}\renewcommand{\rhchXLinBarFullVsRmPctMax}{$-$3.5}  % ADDITIONAL: Lin all: highest %
\providecommand{\rhchXLinBarFullVsRmNSig}{}\renewcommand{\rhchXLinBarFullVsRmNSig}{3}  % ADDITIONAL: Lin all: classes whose full set beats HAR+RM at 5%
\providecommand{\rhchXLinEventFullVsRmPctMin}{}\renewcommand{\rhchXLinEventFullVsRmPctMin}{$-$49.5}  % ADDITIONAL: Lin event: lowest %
\providecommand{\rhchXLinEventFullVsRmPctMax}{}\renewcommand{\rhchXLinEventFullVsRmPctMax}{$-$25.7}  % ADDITIONAL: Lin event: highest %
\providecommand{\rhchXLinEventFullVsRmNSig}{}\renewcommand{\rhchXLinEventFullVsRmNSig}{3}  % ADDITIONAL: Lin event: classes whose full set beats HAR+RM at 5%
\providecommand{\rhchXLinOrdFullVsRmPctMin}{}\renewcommand{\rhchXLinOrdFullVsRmPctMin}{$-$0.1}  % ADDITIONAL: Lin ord: lowest %
\providecommand{\rhchXLinOrdFullVsRmPctMax}{}\renewcommand{\rhchXLinOrdFullVsRmPctMax}{0.9}  % ADDITIONAL: Lin ord: highest %
\providecommand{\rhchXLinOrdFullVsRmNSig}{}\renewcommand{\rhchXLinOrdFullVsRmNSig}{0}  % ADDITIONAL: Lin ord: classes whose full set beats HAR+RM at 5%
\providecommand{\rhchXTreeBarFullVsRmPctMin}{}\renewcommand{\rhchXTreeBarFullVsRmPctMin}{$-$11.6}  % ADDITIONAL: Tree all: lowest %
\providecommand{\rhchXTreeBarFullVsRmPctMax}{}\renewcommand{\rhchXTreeBarFullVsRmPctMax}{$-$8.4}  % ADDITIONAL: Tree all: highest %
\providecommand{\rhchXTreeBarFullVsRmNSig}{}\renewcommand{\rhchXTreeBarFullVsRmNSig}{3}  % ADDITIONAL: Tree all: classes whose full set beats HAR+RM at 5%
\providecommand{\rhchXTreeEventFullVsRmPctMin}{}\renewcommand{\rhchXTreeEventFullVsRmPctMin}{$-$65.8}  % ADDITIONAL: Tree event: lowest %
\providecommand{\rhchXTreeEventFullVsRmPctMax}{}\renewcommand{\rhchXTreeEventFullVsRmPctMax}{$-$47.6}  % ADDITIONAL: Tree event: highest %
\providecommand{\rhchXTreeEventFullVsRmNSig}{}\renewcommand{\rhchXTreeEventFullVsRmNSig}{3}  % ADDITIONAL: Tree event: classes whose full set beats HAR+RM at 5%
\providecommand{\rhchXTreeOrdFullVsRmPctMin}{}\renewcommand{\rhchXTreeOrdFullVsRmPctMin}{$-$1.5}  % ADDITIONAL: Tree ord: lowest %
\providecommand{\rhchXTreeOrdFullVsRmPctMax}{}\renewcommand{\rhchXTreeOrdFullVsRmPctMax}{$-$0.7}  % ADDITIONAL: Tree ord: highest %
\providecommand{\rhchXTreeOrdFullVsRmNSig}{}\renewcommand{\rhchXTreeOrdFullVsRmNSig}{3}  % ADDITIONAL: Tree ord: classes whose full set beats HAR+RM at 5%
\providecommand{\rhchRidgeBarLadFullPct}{}\renewcommand{\rhchRidgeBarLadFullPct}{$-$3.6}  % ridge all: ladder k=46 vs 0 %
\providecommand{\rhchRidgeBarROne}{}\renewcommand{\rhchRidgeBarROne}{56}  % ridge all: ladder R_1 (%)
\providecommand{\rhchRidgeBarRThree}{}\renewcommand{\rhchRidgeBarRThree}{84}  % ridge all: ladder R_3 (%)
\providecommand{\rhchRidgeBarRFive}{}\renewcommand{\rhchRidgeBarRFive}{145}  % ridge all: ladder R_5 (%)
\providecommand{\rhchRidgeBarRTen}{}\renewcommand{\rhchRidgeBarRTen}{188}  % ridge all: ladder R_10 (%)
\providecommand{\rhchRidgeBarRTwenty}{}\renewcommand{\rhchRidgeBarRTwenty}{171}  % ridge all: ladder R_20 (%)
\providecommand{\rhchRidgeEventLadFullPct}{}\renewcommand{\rhchRidgeEventLadFullPct}{$-$23.2}  % ridge event: ladder k=46 vs 0 %
\providecommand{\rhchRidgeEventROne}{}\renewcommand{\rhchRidgeEventROne}{66}  % ridge event: ladder R_1 (%)
\providecommand{\rhchRidgeEventRThree}{}\renewcommand{\rhchRidgeEventRThree}{95}  % ridge event: ladder R_3 (%)
\providecommand{\rhchRidgeEventRFive}{}\renewcommand{\rhchRidgeEventRFive}{96}  % ridge event: ladder R_5 (%)
\providecommand{\rhchRidgeEventRTen}{}\renewcommand{\rhchRidgeEventRTen}{101}  % ridge event: ladder R_10 (%)
\providecommand{\rhchRidgeEventRTwenty}{}\renewcommand{\rhchRidgeEventRTwenty}{99}  % ridge event: ladder R_20 (%)
\providecommand{\rhchRidgeOrdLadFullPct}{}\renewcommand{\rhchRidgeOrdLadFullPct}{$-$0.3}  % ridge ord: ladder k=46 vs 0 %
\providecommand{\rhchRidgeOrdROne}{}\renewcommand{\rhchRidgeOrdROne}{$-$70}  % ridge ord: ladder R_1 (%)
\providecommand{\rhchRidgeOrdRThree}{}\renewcommand{\rhchRidgeOrdRThree}{$-$75}  % ridge ord: ladder R_3 (%)
\providecommand{\rhchRidgeOrdRFive}{}\renewcommand{\rhchRidgeOrdRFive}{824}  % ridge ord: ladder R_5 (%)
\providecommand{\rhchRidgeOrdRTen}{}\renewcommand{\rhchRidgeOrdRTen}{1391}  % ridge ord: ladder R_10 (%)
\providecommand{\rhchRidgeOrdRTwenty}{}\renewcommand{\rhchRidgeOrdRTwenty}{1168}  % ridge ord: ladder R_20 (%)
\providecommand{\rhchLassoBarLadFullPct}{}\renewcommand{\rhchLassoBarLadFullPct}{$-$6.0}  % lasso all: ladder k=46 vs 0 %
\providecommand{\rhchLassoBarROne}{}\renewcommand{\rhchLassoBarROne}{69}  % lasso all: ladder R_1 (%)
\providecommand{\rhchLassoBarRThree}{}\renewcommand{\rhchLassoBarRThree}{63}  % lasso all: ladder R_3 (%)
\providecommand{\rhchLassoBarRFive}{}\renewcommand{\rhchLassoBarRFive}{108}  % lasso all: ladder R_5 (%)
\providecommand{\rhchLassoBarRTen}{}\renewcommand{\rhchLassoBarRTen}{155}  % lasso all: ladder R_10 (%)
\providecommand{\rhchLassoBarRTwenty}{}\renewcommand{\rhchLassoBarRTwenty}{139}  % lasso all: ladder R_20 (%)
\providecommand{\rhchLassoEventLadFullPct}{}\renewcommand{\rhchLassoEventLadFullPct}{$-$46.9}  % lasso event: ladder k=46 vs 0 %
\providecommand{\rhchLassoEventROne}{}\renewcommand{\rhchLassoEventROne}{64}  % lasso event: ladder R_1 (%)
\providecommand{\rhchLassoEventRThree}{}\renewcommand{\rhchLassoEventRThree}{67}  % lasso event: ladder R_3 (%)
\providecommand{\rhchLassoEventRFive}{}\renewcommand{\rhchLassoEventRFive}{92}  % lasso event: ladder R_5 (%)
\providecommand{\rhchLassoEventRTen}{}\renewcommand{\rhchLassoEventRTen}{96}  % lasso event: ladder R_10 (%)
\providecommand{\rhchLassoEventRTwenty}{}\renewcommand{\rhchLassoEventRTwenty}{97}  % lasso event: ladder R_20 (%)
\providecommand{\rhchLassoOrdLadFullPct}{}\renewcommand{\rhchLassoOrdLadFullPct}{1.0}  % lasso ord: ladder k=46 vs 0 %
\providecommand{\rhchLassoOrdROne}{}\renewcommand{\rhchLassoOrdROne}{--}  % lasso ord: ladder R_1 (%)
\providecommand{\rhchLassoOrdRThree}{}\renewcommand{\rhchLassoOrdRThree}{--}  % lasso ord: ladder R_3 (%)
\providecommand{\rhchLassoOrdRFive}{}\renewcommand{\rhchLassoOrdRFive}{--}  % lasso ord: ladder R_5 (%)
\providecommand{\rhchLassoOrdRTen}{}\renewcommand{\rhchLassoOrdRTen}{--}  % lasso ord: ladder R_10 (%)
\providecommand{\rhchLassoOrdRTwenty}{}\renewcommand{\rhchLassoOrdRTwenty}{--}  % lasso ord: ladder R_20 (%)
\providecommand{\rhchEnetBarLadFullPct}{}\renewcommand{\rhchEnetBarLadFullPct}{$-$5.3}  % enet all: ladder k=46 vs 0 %
\providecommand{\rhchEnetBarROne}{}\renewcommand{\rhchEnetBarROne}{55}  % enet all: ladder R_1 (%)
\providecommand{\rhchEnetBarRThree}{}\renewcommand{\rhchEnetBarRThree}{58}  % enet all: ladder R_3 (%)
\providecommand{\rhchEnetBarRFive}{}\renewcommand{\rhchEnetBarRFive}{99}  % enet all: ladder R_5 (%)
\providecommand{\rhchEnetBarRTen}{}\renewcommand{\rhchEnetBarRTen}{130}  % enet all: ladder R_10 (%)
\providecommand{\rhchEnetBarRTwenty}{}\renewcommand{\rhchEnetBarRTwenty}{120}  % enet all: ladder R_20 (%)
\providecommand{\rhchEnetEventLadFullPct}{}\renewcommand{\rhchEnetEventLadFullPct}{$-$23.0}  % enet event: ladder k=46 vs 0 %
\providecommand{\rhchEnetEventROne}{}\renewcommand{\rhchEnetEventROne}{90}  % enet event: ladder R_1 (%)
\providecommand{\rhchEnetEventRThree}{}\renewcommand{\rhchEnetEventRThree}{99}  % enet event: ladder R_3 (%)
\providecommand{\rhchEnetEventRFive}{}\renewcommand{\rhchEnetEventRFive}{97}  % enet event: ladder R_5 (%)
\providecommand{\rhchEnetEventRTen}{}\renewcommand{\rhchEnetEventRTen}{102}  % enet event: ladder R_10 (%)
\providecommand{\rhchEnetEventRTwenty}{}\renewcommand{\rhchEnetEventRTwenty}{98}  % enet event: ladder R_20 (%)
\providecommand{\rhchEnetOrdLadFullPct}{}\renewcommand{\rhchEnetOrdLadFullPct}{$-$2.2}  % enet ord: ladder k=46 vs 0 %
\providecommand{\rhchEnetOrdROne}{}\renewcommand{\rhchEnetOrdROne}{$-$7}  % enet ord: ladder R_1 (%)
\providecommand{\rhchEnetOrdRThree}{}\renewcommand{\rhchEnetOrdRThree}{$-$16}  % enet ord: ladder R_3 (%)
\providecommand{\rhchEnetOrdRFive}{}\renewcommand{\rhchEnetOrdRFive}{102}  % enet ord: ladder R_5 (%)
\providecommand{\rhchEnetOrdRTen}{}\renewcommand{\rhchEnetOrdRTen}{180}  % enet ord: ladder R_10 (%)
\providecommand{\rhchEnetOrdRTwenty}{}\renewcommand{\rhchEnetOrdRTwenty}{159}  % enet ord: ladder R_20 (%)
\providecommand{\rhchLgbmBarLadFullPct}{}\renewcommand{\rhchLgbmBarLadFullPct}{$-$13.6}  % lgbm all: ladder k=46 vs 0 %
\providecommand{\rhchLgbmBarROne}{}\renewcommand{\rhchLgbmBarROne}{3}  % lgbm all: ladder R_1 (%)
\providecommand{\rhchLgbmBarRThree}{}\renewcommand{\rhchLgbmBarRThree}{34}  % lgbm all: ladder R_3 (%)
\providecommand{\rhchLgbmBarRFive}{}\renewcommand{\rhchLgbmBarRFive}{34}  % lgbm all: ladder R_5 (%)
\providecommand{\rhchLgbmBarRTen}{}\renewcommand{\rhchLgbmBarRTen}{95}  % lgbm all: ladder R_10 (%)
\providecommand{\rhchLgbmBarRTwenty}{}\renewcommand{\rhchLgbmBarRTwenty}{96}  % lgbm all: ladder R_20 (%)
\providecommand{\rhchLgbmEventLadFullPct}{}\renewcommand{\rhchLgbmEventLadFullPct}{$-$62.0}  % lgbm event: ladder k=46 vs 0 %
\providecommand{\rhchLgbmEventROne}{}\renewcommand{\rhchLgbmEventROne}{$-$1}  % lgbm event: ladder R_1 (%)
\providecommand{\rhchLgbmEventRThree}{}\renewcommand{\rhchLgbmEventRThree}{5}  % lgbm event: ladder R_3 (%)
\providecommand{\rhchLgbmEventRFive}{}\renewcommand{\rhchLgbmEventRFive}{3}  % lgbm event: ladder R_5 (%)
\providecommand{\rhchLgbmEventRTen}{}\renewcommand{\rhchLgbmEventRTen}{100}  % lgbm event: ladder R_10 (%)
\providecommand{\rhchLgbmEventRTwenty}{}\renewcommand{\rhchLgbmEventRTwenty}{100}  % lgbm event: ladder R_20 (%)
\providecommand{\rhchLgbmOrdLadFullPct}{}\renewcommand{\rhchLgbmOrdLadFullPct}{$-$5.5}  % lgbm ord: ladder k=46 vs 0 %
\providecommand{\rhchLgbmOrdROne}{}\renewcommand{\rhchLgbmOrdROne}{12}  % lgbm ord: ladder R_1 (%)
\providecommand{\rhchLgbmOrdRThree}{}\renewcommand{\rhchLgbmOrdRThree}{88}  % lgbm ord: ladder R_3 (%)
\providecommand{\rhchLgbmOrdRFive}{}\renewcommand{\rhchLgbmOrdRFive}{93}  % lgbm ord: ladder R_5 (%)
\providecommand{\rhchLgbmOrdRTen}{}\renewcommand{\rhchLgbmOrdRTen}{86}  % lgbm ord: ladder R_10 (%)
\providecommand{\rhchLgbmOrdRTwenty}{}\renewcommand{\rhchLgbmOrdRTwenty}{90}  % lgbm ord: ladder R_20 (%)
\providecommand{\rhchXgbBarLadFullPct}{}\renewcommand{\rhchXgbBarLadFullPct}{$-$13.6}  % xgb all: ladder k=46 vs 0 %
\providecommand{\rhchXgbBarROne}{}\renewcommand{\rhchXgbBarROne}{4}  % xgb all: ladder R_1 (%)
\providecommand{\rhchXgbBarRThree}{}\renewcommand{\rhchXgbBarRThree}{27}  % xgb all: ladder R_3 (%)
\providecommand{\rhchXgbBarRFive}{}\renewcommand{\rhchXgbBarRFive}{36}  % xgb all: ladder R_5 (%)
\providecommand{\rhchXgbBarRTen}{}\renewcommand{\rhchXgbBarRTen}{98}  % xgb all: ladder R_10 (%)
\providecommand{\rhchXgbBarRTwenty}{}\renewcommand{\rhchXgbBarRTwenty}{99}  % xgb all: ladder R_20 (%)
\providecommand{\rhchXgbEventLadFullPct}{}\renewcommand{\rhchXgbEventLadFullPct}{$-$63.5}  % xgb event: ladder k=46 vs 0 %
\providecommand{\rhchXgbEventROne}{}\renewcommand{\rhchXgbEventROne}{0}  % xgb event: ladder R_1 (%)
\providecommand{\rhchXgbEventRThree}{}\renewcommand{\rhchXgbEventRThree}{10}  % xgb event: ladder R_3 (%)
\providecommand{\rhchXgbEventRFive}{}\renewcommand{\rhchXgbEventRFive}{3}  % xgb event: ladder R_5 (%)
\providecommand{\rhchXgbEventRTen}{}\renewcommand{\rhchXgbEventRTen}{99}  % xgb event: ladder R_10 (%)
\providecommand{\rhchXgbEventRTwenty}{}\renewcommand{\rhchXgbEventRTwenty}{100}  % xgb event: ladder R_20 (%)
\providecommand{\rhchXgbOrdLadFullPct}{}\renewcommand{\rhchXgbOrdLadFullPct}{$-$5.1}  % xgb ord: ladder k=46 vs 0 %
\providecommand{\rhchXgbOrdROne}{}\renewcommand{\rhchXgbOrdROne}{13}  % xgb ord: ladder R_1 (%)
\providecommand{\rhchXgbOrdRThree}{}\renewcommand{\rhchXgbOrdRThree}{64}  % xgb ord: ladder R_3 (%)
\providecommand{\rhchXgbOrdRFive}{}\renewcommand{\rhchXgbOrdRFive}{108}  % xgb ord: ladder R_5 (%)
\providecommand{\rhchXgbOrdRTen}{}\renewcommand{\rhchXgbOrdRTen}{95}  % xgb ord: ladder R_10 (%)
\providecommand{\rhchXgbOrdRTwenty}{}\renewcommand{\rhchXgbOrdRTwenty}{99}  % xgb ord: ladder R_20 (%)
\providecommand{\rhchRfBarLadFullPct}{}\renewcommand{\rhchRfBarLadFullPct}{$-$10.6}  % rf all: ladder k=46 vs 0 %
\providecommand{\rhchRfBarROne}{}\renewcommand{\rhchRfBarROne}{6}  % rf all: ladder R_1 (%)
\providecommand{\rhchRfBarRThree}{}\renewcommand{\rhchRfBarRThree}{42}  % rf all: ladder R_3 (%)
\providecommand{\rhchRfBarRFive}{}\renewcommand{\rhchRfBarRFive}{40}  % rf all: ladder R_5 (%)
\providecommand{\rhchRfBarRTen}{}\renewcommand{\rhchRfBarRTen}{98}  % rf all: ladder R_10 (%)
\providecommand{\rhchRfBarRTwenty}{}\renewcommand{\rhchRfBarRTwenty}{101}  % rf all: ladder R_20 (%)
\providecommand{\rhchRfEventLadFullPct}{}\renewcommand{\rhchRfEventLadFullPct}{$-$44.4}  % rf event: ladder k=46 vs 0 %
\providecommand{\rhchRfEventROne}{}\renewcommand{\rhchRfEventROne}{$-$3}  % rf event: ladder R_1 (%)
\providecommand{\rhchRfEventRThree}{}\renewcommand{\rhchRfEventRThree}{3}  % rf event: ladder R_3 (%)
\providecommand{\rhchRfEventRFive}{}\renewcommand{\rhchRfEventRFive}{6}  % rf event: ladder R_5 (%)
\providecommand{\rhchRfEventRTen}{}\renewcommand{\rhchRfEventRTen}{100}  % rf event: ladder R_10 (%)
\providecommand{\rhchRfEventRTwenty}{}\renewcommand{\rhchRfEventRTwenty}{100}  % rf event: ladder R_20 (%)
\providecommand{\rhchRfOrdLadFullPct}{}\renewcommand{\rhchRfOrdLadFullPct}{$-$4.6}  % rf ord: ladder k=46 vs 0 %
\providecommand{\rhchRfOrdROne}{}\renewcommand{\rhchRfOrdROne}{20}  % rf ord: ladder R_1 (%)
\providecommand{\rhchRfOrdRThree}{}\renewcommand{\rhchRfOrdRThree}{109}  % rf ord: ladder R_3 (%)
\providecommand{\rhchRfOrdRFive}{}\renewcommand{\rhchRfOrdRFive}{97}  % rf ord: ladder R_5 (%)
\providecommand{\rhchRfOrdRTen}{}\renewcommand{\rhchRfOrdRTen}{95}  % rf ord: ladder R_10 (%)
\providecommand{\rhchRfOrdRTwenty}{}\renewcommand{\rhchRfOrdRTwenty}{102}  % rf ord: ladder R_20 (%)
\fi
\begin{table}[htbp]
\centering
\footnotesize
\setlength{\tabcolsep}{2pt}
\caption{Concentrated or spread, by day type: the lasso against the dense penalties, and the share of the full set's
  gain kept by the top $k$ series}
\label{tab:rh_channels_sparse}
\begin{tabular*}{\linewidth}{@{\extracolsep{\fill}}lrrr@{}}
\toprule
\multicolumn{4}{@{}l}{\textit{A. Full sets: \% change in QLIKE (DM statistic); $-$: the lasso is more accurate, $+$: less accurate}} \\
 & Every bar & FOMC statement days & Ordinary days \\
\midrule
Lasso vs ridge & $-$3.1 ($-$4.82) & $-$31.9 ($-$8.42) & 0.8 (2.28) \\
Lasso vs elastic net & $-$2.4 ($-$3.49) & $-$31.8 ($-$7.69) & 1.6 (3.35) \\
\bottomrule
\end{tabular*}

\medskip
\begin{tabular*}{\linewidth}{@{\extracolsep{\fill}}lrrrrrrrrrrrr@{}}
\toprule
\multicolumn{13}{@{}l}{\textit{B. Feature-count ladder: retained gain $R_k$ (\%) by day type}} \\
\multicolumn{13}{@{}l}{\textit{Full: $-$ = the full set beats HAR. $R_k$: 100 = all of that gain kept, 0 = none, below 0 = worse than HAR}} \\
 & \multicolumn{6}{c}{FOMC statement days} & \multicolumn{6}{c}{Ordinary days} \\
\cmidrule(lr){2-7}\cmidrule(l){8-13}
 & Full (\%) & $k{=}1$ & $k{=}3$ & $k{=}5$ & $k{=}10$ & $k{=}20$ & Full (\%) & $k{=}1$ & $k{=}3$ & $k{=}5$ & $k{=}10$ & $k{=}20$ \\
\midrule
Ridge & $-$23.2 & 66 & 95 & 96 & 101 & 99 & $-$0.3 & $-$70 & $-$75 & 824 & 1391 & 1168 \\
Lasso & $-$46.9 & 64 & 67 & 92 & 96 & 97 & 1.0 & -- & -- & -- & -- & -- \\
Elastic net & $-$23.0 & 90 & 99 & 97 & 102 & 98 & $-$2.2 & $-$7 & $-$16 & 102 & 180 & 159 \\
\addlinespace[2pt]
LightGBM & $-$62.0 & $-$1 & 5 & 3 & 100 & 100 & $-$5.5 & 12 & 88 & 93 & 86 & 90 \\
XGBoost & $-$63.5 & 0 & 10 & 3 & 99 & 100 & $-$5.1 & 13 & 64 & 108 & 95 & 99 \\
Random forest & $-$44.4 & $-$3 & 3 & 6 & 100 & 100 & $-$4.6 & 20 & 109 & 97 & 95 & 102 \\
\bottomrule
\end{tabular*}

\smallskip
\begin{minipage}{\linewidth}
\footnotesize
\textit{Notes:} \textit{Sign:} a \% change is $100\,(L_{\mathrm{row}}/L_{\mathrm{ref}} - 1)$, the percentage
change in mean QLIKE of the row's forecast against its reference, and Panel A's Diebold--Mariano statistic has the same
sign. Panel A (reference: the ridge's or the elastic net's full set): negative = the lasso is more accurate. Panel B,
Full (the full set against the HAR benchmark set): negative = the full set is more accurate; $R_k$: 100 = the top $k$
series keep all of the full set's gain over the HAR benchmark set, 0 = none of it, below 0 = worse than the HAR
benchmark set, above 100 = better than the full set; --: no gain to keep. Panel A: the lasso's full-set forecast against
the ridge's and the elastic net's, percentage change in QLIKE with the Diebold--Mariano statistic in parentheses (HLN
correction, daily average loss, Newey--West lag \rhchHacLag). Panel B: each model refitted monthly on its $k$
highest-ranked exogenous series (ranking from a pre-sample fit; the feature-count ladder of
Section~\ref{sec:rh_fi}), the FOMC-patched design. Full: the $k = 46$ rung (the full set) against the $k = 0$ rung (the HAR benchmark
set), \% change in QLIKE on the day type; $R_k = (L_0 - L_k) / (L_0 - L_{46})$, the share of that gain the top $k$
series keep.
\end{minipage}
\end{table}

% AUTO-GENERATED by writeup/make_rehaul_factor.py -- do not edit; re-run the script.
% tab:rh_factor: factor-augmented HAR, ORDINARY days and every bar (results/channels_2026-10-26/factor).
\ifdefined\rhfaReading\else% AUTO-GENERATED by writeup/make_rehaul_factor.py -- do not edit; re-run the script.
% \rhfa... macros: factor-augmented HAR (diffusion index) on Table 1's design (results/channels_2026-10-26/
% RULE.txt section 4; results/channels_2026-10-26/factor).  Spans: Ord = ORDINARY days (primary), All = every
% bar, Ev = the 81 FOMC statement days.  Models: Pca<k>, Pls<k>, Full (the span's L_full model).
\providecommand{\rhfaPcaOneOrdPct}{}\renewcommand{\rhfaPcaOneOrdPct}{0.08}  % pca1 ordinary: QLIKE change vs OLS-HAR, %
\providecommand{\rhfaPcaOneOrdMsePct}{}\renewcommand{\rhfaPcaOneOrdMsePct}{3.28}  % pca1 ordinary: MSE change vs OLS-HAR, %
\providecommand{\rhfaPcaOneOrdDm}{}\renewcommand{\rhfaPcaOneOrdDm}{0.90}  % pca1 ordinary: DM vs OLS-HAR (QLIKE)
\providecommand{\rhfaPcaOneOrdP}{}\renewcommand{\rhfaPcaOneOrdP}{0.369}  % pca1 ordinary: p
\providecommand{\rhfaPcaOneOrdDmFull}{}\renewcommand{\rhfaPcaOneOrdDmFull}{4.89}  % pca1 ordinary: DM vs the full set (enet)
\providecommand{\rhfaPcaOneOrdPFull}{}\renewcommand{\rhfaPcaOneOrdPFull}{$<$0.001}  % pca1 ordinary: p
\providecommand{\rhfaPcaOneOrdR}{}\renewcommand{\rhfaPcaOneOrdR}{$-$0.02}  % pca1 ordinary: R_k
\providecommand{\rhfaPcaTwoOrdPct}{}\renewcommand{\rhfaPcaTwoOrdPct}{0.06}  % pca2 ordinary: QLIKE change vs OLS-HAR, %
\providecommand{\rhfaPcaTwoOrdMsePct}{}\renewcommand{\rhfaPcaTwoOrdMsePct}{1.85}  % pca2 ordinary: MSE change vs OLS-HAR, %
\providecommand{\rhfaPcaTwoOrdDm}{}\renewcommand{\rhfaPcaTwoOrdDm}{0.32}  % pca2 ordinary: DM vs OLS-HAR (QLIKE)
\providecommand{\rhfaPcaTwoOrdP}{}\renewcommand{\rhfaPcaTwoOrdP}{0.748}  % pca2 ordinary: p
\providecommand{\rhfaPcaTwoOrdDmFull}{}\renewcommand{\rhfaPcaTwoOrdDmFull}{4.86}  % pca2 ordinary: DM vs the full set (enet)
\providecommand{\rhfaPcaTwoOrdPFull}{}\renewcommand{\rhfaPcaTwoOrdPFull}{$<$0.001}  % pca2 ordinary: p
\providecommand{\rhfaPcaTwoOrdR}{}\renewcommand{\rhfaPcaTwoOrdR}{$-$0.02}  % pca2 ordinary: R_k
\providecommand{\rhfaPcaThreeOrdPct}{}\renewcommand{\rhfaPcaThreeOrdPct}{0.45}  % pca3 ordinary: QLIKE change vs OLS-HAR, %
\providecommand{\rhfaPcaThreeOrdMsePct}{}\renewcommand{\rhfaPcaThreeOrdMsePct}{5.76}  % pca3 ordinary: MSE change vs OLS-HAR, %
\providecommand{\rhfaPcaThreeOrdDm}{}\renewcommand{\rhfaPcaThreeOrdDm}{1.44}  % pca3 ordinary: DM vs OLS-HAR (QLIKE)
\providecommand{\rhfaPcaThreeOrdP}{}\renewcommand{\rhfaPcaThreeOrdP}{0.151}  % pca3 ordinary: p
\providecommand{\rhfaPcaThreeOrdDmFull}{}\renewcommand{\rhfaPcaThreeOrdDmFull}{5.33}  % pca3 ordinary: DM vs the full set (enet)
\providecommand{\rhfaPcaThreeOrdPFull}{}\renewcommand{\rhfaPcaThreeOrdPFull}{$<$0.001}  % pca3 ordinary: p
\providecommand{\rhfaPcaThreeOrdR}{}\renewcommand{\rhfaPcaThreeOrdR}{$-$0.11}  % pca3 ordinary: R_k
\providecommand{\rhfaPcaFiveOrdPct}{}\renewcommand{\rhfaPcaFiveOrdPct}{0.93}  % pca5 ordinary: QLIKE change vs OLS-HAR, %
\providecommand{\rhfaPcaFiveOrdMsePct}{}\renewcommand{\rhfaPcaFiveOrdMsePct}{3.77}  % pca5 ordinary: MSE change vs OLS-HAR, %
\providecommand{\rhfaPcaFiveOrdDm}{}\renewcommand{\rhfaPcaFiveOrdDm}{1.85}  % pca5 ordinary: DM vs OLS-HAR (QLIKE)
\providecommand{\rhfaPcaFiveOrdP}{}\renewcommand{\rhfaPcaFiveOrdP}{0.064}  % pca5 ordinary: p
\providecommand{\rhfaPcaFiveOrdDmFull}{}\renewcommand{\rhfaPcaFiveOrdDmFull}{5.94}  % pca5 ordinary: DM vs the full set (enet)
\providecommand{\rhfaPcaFiveOrdPFull}{}\renewcommand{\rhfaPcaFiveOrdPFull}{$<$0.001}  % pca5 ordinary: p
\providecommand{\rhfaPcaFiveOrdR}{}\renewcommand{\rhfaPcaFiveOrdR}{$-$0.23}  % pca5 ordinary: R_k
\providecommand{\rhfaPcaTenOrdPct}{}\renewcommand{\rhfaPcaTenOrdPct}{$-$1.06}  % pca10 ordinary: QLIKE change vs OLS-HAR, %
\providecommand{\rhfaPcaTenOrdMsePct}{}\renewcommand{\rhfaPcaTenOrdMsePct}{21.39}  % pca10 ordinary: MSE change vs OLS-HAR, %
\providecommand{\rhfaPcaTenOrdDm}{}\renewcommand{\rhfaPcaTenOrdDm}{$-$2.19}  % pca10 ordinary: DM vs OLS-HAR (QLIKE)
\providecommand{\rhfaPcaTenOrdP}{}\renewcommand{\rhfaPcaTenOrdP}{0.028}  % pca10 ordinary: p
\providecommand{\rhfaPcaTenOrdDmFull}{}\renewcommand{\rhfaPcaTenOrdDmFull}{4.45}  % pca10 ordinary: DM vs the full set (enet)
\providecommand{\rhfaPcaTenOrdPFull}{}\renewcommand{\rhfaPcaTenOrdPFull}{$<$0.001}  % pca10 ordinary: p
\providecommand{\rhfaPcaTenOrdR}{}\renewcommand{\rhfaPcaTenOrdR}{0.27}  % pca10 ordinary: R_k
\providecommand{\rhfaPcaTwentyOrdPct}{}\renewcommand{\rhfaPcaTwentyOrdPct}{$-$1.09}  % pca20 ordinary: QLIKE change vs OLS-HAR, %
\providecommand{\rhfaPcaTwentyOrdMsePct}{}\renewcommand{\rhfaPcaTwentyOrdMsePct}{45.47}  % pca20 ordinary: MSE change vs OLS-HAR, %
\providecommand{\rhfaPcaTwentyOrdDm}{}\renewcommand{\rhfaPcaTwentyOrdDm}{$-$1.53}  % pca20 ordinary: DM vs OLS-HAR (QLIKE)
\providecommand{\rhfaPcaTwentyOrdP}{}\renewcommand{\rhfaPcaTwentyOrdP}{0.126}  % pca20 ordinary: p
\providecommand{\rhfaPcaTwentyOrdDmFull}{}\renewcommand{\rhfaPcaTwentyOrdDmFull}{3.66}  % pca20 ordinary: DM vs the full set (enet)
\providecommand{\rhfaPcaTwentyOrdPFull}{}\renewcommand{\rhfaPcaTwentyOrdPFull}{$<$0.001}  % pca20 ordinary: p
\providecommand{\rhfaPcaTwentyOrdR}{}\renewcommand{\rhfaPcaTwentyOrdR}{0.27}  % pca20 ordinary: R_k
\providecommand{\rhfaPlsOneOrdPct}{}\renewcommand{\rhfaPlsOneOrdPct}{0.01}  % pls1 ordinary: QLIKE change vs OLS-HAR, %
\providecommand{\rhfaPlsOneOrdMsePct}{}\renewcommand{\rhfaPlsOneOrdMsePct}{14.31}  % pls1 ordinary: MSE change vs OLS-HAR, %
\providecommand{\rhfaPlsOneOrdDm}{}\renewcommand{\rhfaPlsOneOrdDm}{0.05}  % pls1 ordinary: DM vs OLS-HAR (QLIKE)
\providecommand{\rhfaPlsOneOrdP}{}\renewcommand{\rhfaPlsOneOrdP}{0.962}  % pls1 ordinary: p
\providecommand{\rhfaPlsOneOrdDmFull}{}\renewcommand{\rhfaPlsOneOrdDmFull}{4.95}  % pls1 ordinary: DM vs the full set (enet)
\providecommand{\rhfaPlsOneOrdPFull}{}\renewcommand{\rhfaPlsOneOrdPFull}{$<$0.001}  % pls1 ordinary: p
\providecommand{\rhfaPlsOneOrdR}{}\renewcommand{\rhfaPlsOneOrdR}{0.00}  % pls1 ordinary: R_k
\providecommand{\rhfaPlsTwoOrdPct}{}\renewcommand{\rhfaPlsTwoOrdPct}{1.57}  % pls2 ordinary: QLIKE change vs OLS-HAR, %
\providecommand{\rhfaPlsTwoOrdMsePct}{}\renewcommand{\rhfaPlsTwoOrdMsePct}{24.30}  % pls2 ordinary: MSE change vs OLS-HAR, %
\providecommand{\rhfaPlsTwoOrdDm}{}\renewcommand{\rhfaPlsTwoOrdDm}{1.29}  % pls2 ordinary: DM vs OLS-HAR (QLIKE)
\providecommand{\rhfaPlsTwoOrdP}{}\renewcommand{\rhfaPlsTwoOrdP}{0.196}  % pls2 ordinary: p
\providecommand{\rhfaPlsTwoOrdDmFull}{}\renewcommand{\rhfaPlsTwoOrdDmFull}{4.51}  % pls2 ordinary: DM vs the full set (enet)
\providecommand{\rhfaPlsTwoOrdPFull}{}\renewcommand{\rhfaPlsTwoOrdPFull}{$<$0.001}  % pls2 ordinary: p
\providecommand{\rhfaPlsTwoOrdR}{}\renewcommand{\rhfaPlsTwoOrdR}{$-$0.40}  % pls2 ordinary: R_k
\providecommand{\rhfaPlsThreeOrdPct}{}\renewcommand{\rhfaPlsThreeOrdPct}{0.04}  % pls3 ordinary: QLIKE change vs OLS-HAR, %
\providecommand{\rhfaPlsThreeOrdMsePct}{}\renewcommand{\rhfaPlsThreeOrdMsePct}{99.04}  % pls3 ordinary: MSE change vs OLS-HAR, %
\providecommand{\rhfaPlsThreeOrdDm}{}\renewcommand{\rhfaPlsThreeOrdDm}{0.03}  % pls3 ordinary: DM vs OLS-HAR (QLIKE)
\providecommand{\rhfaPlsThreeOrdP}{}\renewcommand{\rhfaPlsThreeOrdP}{0.974}  % pls3 ordinary: p
\providecommand{\rhfaPlsThreeOrdDmFull}{}\renewcommand{\rhfaPlsThreeOrdDmFull}{3.84}  % pls3 ordinary: DM vs the full set (enet)
\providecommand{\rhfaPlsThreeOrdPFull}{}\renewcommand{\rhfaPlsThreeOrdPFull}{$<$0.001}  % pls3 ordinary: p
\providecommand{\rhfaPlsThreeOrdR}{}\renewcommand{\rhfaPlsThreeOrdR}{$-$0.01}  % pls3 ordinary: R_k
\providecommand{\rhfaFullOrdModel}{}\renewcommand{\rhfaFullOrdModel}{elastic net}  % ordinary: L_full model (lower QLIKE of ridge, elastic net)
\providecommand{\rhfaFullOrdPct}{}\renewcommand{\rhfaFullOrdPct}{$-$3.97}  % ordinary: full set QLIKE change vs OLS-HAR, %
\providecommand{\rhfaFullOrdMsePct}{}\renewcommand{\rhfaFullOrdMsePct}{128.86}  % ordinary: full set MSE change vs OLS-HAR, %
\providecommand{\rhfaFullOrdDm}{}\renewcommand{\rhfaFullOrdDm}{$-$4.73}  % ordinary: full set DM vs OLS-HAR
\providecommand{\rhfaFullOrdP}{}\renewcommand{\rhfaFullOrdP}{$<$0.001}  % ordinary: p
\providecommand{\rhfaNOrdDays}{}\renewcommand{\rhfaNOrdDays}{2{,}498}  % ordinary: days
\providecommand{\rhfaNOrdBars}{}\renewcommand{\rhfaNOrdBars}{32{,}474}  % ordinary: bars
\providecommand{\rhfaPcaOneAllPct}{}\renewcommand{\rhfaPcaOneAllPct}{$-$0.05}  % pca1 all: QLIKE change vs OLS-HAR, %
\providecommand{\rhfaPcaOneAllMsePct}{}\renewcommand{\rhfaPcaOneAllMsePct}{3.00}  % pca1 all: MSE change vs OLS-HAR, %
\providecommand{\rhfaPcaOneAllDm}{}\renewcommand{\rhfaPcaOneAllDm}{$-$0.53}  % pca1 all: DM vs OLS-HAR (QLIKE)
\providecommand{\rhfaPcaOneAllP}{}\renewcommand{\rhfaPcaOneAllP}{0.595}  % pca1 all: p
\providecommand{\rhfaPcaOneAllDmFull}{}\renewcommand{\rhfaPcaOneAllDmFull}{7.44}  % pca1 all: DM vs the full set (enet)
\providecommand{\rhfaPcaOneAllPFull}{}\renewcommand{\rhfaPcaOneAllPFull}{$<$0.001}  % pca1 all: p
\providecommand{\rhfaPcaOneAllR}{}\renewcommand{\rhfaPcaOneAllR}{0.01}  % pca1 all: R_k
\providecommand{\rhfaPcaTwoAllPct}{}\renewcommand{\rhfaPcaTwoAllPct}{0.11}  % pca2 all: QLIKE change vs OLS-HAR, %
\providecommand{\rhfaPcaTwoAllMsePct}{}\renewcommand{\rhfaPcaTwoAllMsePct}{1.71}  % pca2 all: MSE change vs OLS-HAR, %
\providecommand{\rhfaPcaTwoAllDm}{}\renewcommand{\rhfaPcaTwoAllDm}{0.55}  % pca2 all: DM vs OLS-HAR (QLIKE)
\providecommand{\rhfaPcaTwoAllP}{}\renewcommand{\rhfaPcaTwoAllP}{0.579}  % pca2 all: p
\providecommand{\rhfaPcaTwoAllDmFull}{}\renewcommand{\rhfaPcaTwoAllDmFull}{7.48}  % pca2 all: DM vs the full set (enet)
\providecommand{\rhfaPcaTwoAllPFull}{}\renewcommand{\rhfaPcaTwoAllPFull}{$<$0.001}  % pca2 all: p
\providecommand{\rhfaPcaTwoAllR}{}\renewcommand{\rhfaPcaTwoAllR}{$-$0.02}  % pca2 all: R_k
\providecommand{\rhfaPcaThreeAllPct}{}\renewcommand{\rhfaPcaThreeAllPct}{0.69}  % pca3 all: QLIKE change vs OLS-HAR, %
\providecommand{\rhfaPcaThreeAllMsePct}{}\renewcommand{\rhfaPcaThreeAllMsePct}{5.30}  % pca3 all: MSE change vs OLS-HAR, %
\providecommand{\rhfaPcaThreeAllDm}{}\renewcommand{\rhfaPcaThreeAllDm}{2.24}  % pca3 all: DM vs OLS-HAR (QLIKE)
\providecommand{\rhfaPcaThreeAllP}{}\renewcommand{\rhfaPcaThreeAllP}{0.025}  % pca3 all: p
\providecommand{\rhfaPcaThreeAllDmFull}{}\renewcommand{\rhfaPcaThreeAllDmFull}{7.96}  % pca3 all: DM vs the full set (enet)
\providecommand{\rhfaPcaThreeAllPFull}{}\renewcommand{\rhfaPcaThreeAllPFull}{$<$0.001}  % pca3 all: p
\providecommand{\rhfaPcaThreeAllR}{}\renewcommand{\rhfaPcaThreeAllR}{$-$0.10}  % pca3 all: R_k
\providecommand{\rhfaPcaFiveAllPct}{}\renewcommand{\rhfaPcaFiveAllPct}{0.94}  % pca5 all: QLIKE change vs OLS-HAR, %
\providecommand{\rhfaPcaFiveAllMsePct}{}\renewcommand{\rhfaPcaFiveAllMsePct}{3.39}  % pca5 all: MSE change vs OLS-HAR, %
\providecommand{\rhfaPcaFiveAllDm}{}\renewcommand{\rhfaPcaFiveAllDm}{2.05}  % pca5 all: DM vs OLS-HAR (QLIKE)
\providecommand{\rhfaPcaFiveAllP}{}\renewcommand{\rhfaPcaFiveAllP}{0.041}  % pca5 all: p
\providecommand{\rhfaPcaFiveAllDmFull}{}\renewcommand{\rhfaPcaFiveAllDmFull}{8.48}  % pca5 all: DM vs the full set (enet)
\providecommand{\rhfaPcaFiveAllPFull}{}\renewcommand{\rhfaPcaFiveAllPFull}{$<$0.001}  % pca5 all: p
\providecommand{\rhfaPcaFiveAllR}{}\renewcommand{\rhfaPcaFiveAllR}{$-$0.14}  % pca5 all: R_k
\providecommand{\rhfaPcaTenAllPct}{}\renewcommand{\rhfaPcaTenAllPct}{$-$0.66}  % pca10 all: QLIKE change vs OLS-HAR, %
\providecommand{\rhfaPcaTenAllMsePct}{}\renewcommand{\rhfaPcaTenAllMsePct}{19.64}  % pca10 all: MSE change vs OLS-HAR, %
\providecommand{\rhfaPcaTenAllDm}{}\renewcommand{\rhfaPcaTenAllDm}{$-$1.43}  % pca10 all: DM vs OLS-HAR (QLIKE)
\providecommand{\rhfaPcaTenAllP}{}\renewcommand{\rhfaPcaTenAllP}{0.152}  % pca10 all: p
\providecommand{\rhfaPcaTenAllDmFull}{}\renewcommand{\rhfaPcaTenAllDmFull}{7.31}  % pca10 all: DM vs the full set (enet)
\providecommand{\rhfaPcaTenAllPFull}{}\renewcommand{\rhfaPcaTenAllPFull}{$<$0.001}  % pca10 all: p
\providecommand{\rhfaPcaTenAllR}{}\renewcommand{\rhfaPcaTenAllR}{0.10}  % pca10 all: R_k
\providecommand{\rhfaPcaTwentyAllPct}{}\renewcommand{\rhfaPcaTwentyAllPct}{$-$0.87}  % pca20 all: QLIKE change vs OLS-HAR, %
\providecommand{\rhfaPcaTwentyAllMsePct}{}\renewcommand{\rhfaPcaTwentyAllMsePct}{41.73}  % pca20 all: MSE change vs OLS-HAR, %
\providecommand{\rhfaPcaTwentyAllDm}{}\renewcommand{\rhfaPcaTwentyAllDm}{$-$1.28}  % pca20 all: DM vs OLS-HAR (QLIKE)
\providecommand{\rhfaPcaTwentyAllP}{}\renewcommand{\rhfaPcaTwentyAllP}{0.201}  % pca20 all: p
\providecommand{\rhfaPcaTwentyAllDmFull}{}\renewcommand{\rhfaPcaTwentyAllDmFull}{6.40}  % pca20 all: DM vs the full set (enet)
\providecommand{\rhfaPcaTwentyAllPFull}{}\renewcommand{\rhfaPcaTwentyAllPFull}{$<$0.001}  % pca20 all: p
\providecommand{\rhfaPcaTwentyAllR}{}\renewcommand{\rhfaPcaTwentyAllR}{0.13}  % pca20 all: R_k
\providecommand{\rhfaPlsOneAllPct}{}\renewcommand{\rhfaPlsOneAllPct}{$-$0.23}  % pls1 all: QLIKE change vs OLS-HAR, %
\providecommand{\rhfaPlsOneAllMsePct}{}\renewcommand{\rhfaPlsOneAllMsePct}{13.12}  % pls1 all: MSE change vs OLS-HAR, %
\providecommand{\rhfaPlsOneAllDm}{}\renewcommand{\rhfaPlsOneAllDm}{$-$1.07}  % pls1 all: DM vs OLS-HAR (QLIKE)
\providecommand{\rhfaPlsOneAllP}{}\renewcommand{\rhfaPlsOneAllP}{0.284}  % pls1 all: p
\providecommand{\rhfaPlsOneAllDmFull}{}\renewcommand{\rhfaPlsOneAllDmFull}{7.49}  % pls1 all: DM vs the full set (enet)
\providecommand{\rhfaPlsOneAllPFull}{}\renewcommand{\rhfaPlsOneAllPFull}{$<$0.001}  % pls1 all: p
\providecommand{\rhfaPlsOneAllR}{}\renewcommand{\rhfaPlsOneAllR}{0.03}  % pls1 all: R_k
\providecommand{\rhfaPlsTwoAllPct}{}\renewcommand{\rhfaPlsTwoAllPct}{1.66}  % pls2 all: QLIKE change vs OLS-HAR, %
\providecommand{\rhfaPlsTwoAllMsePct}{}\renewcommand{\rhfaPlsTwoAllMsePct}{22.30}  % pls2 all: MSE change vs OLS-HAR, %
\providecommand{\rhfaPlsTwoAllDm}{}\renewcommand{\rhfaPlsTwoAllDm}{1.55}  % pls2 all: DM vs OLS-HAR (QLIKE)
\providecommand{\rhfaPlsTwoAllP}{}\renewcommand{\rhfaPlsTwoAllP}{0.120}  % pls2 all: p
\providecommand{\rhfaPlsTwoAllDmFull}{}\renewcommand{\rhfaPlsTwoAllDmFull}{6.99}  % pls2 all: DM vs the full set (enet)
\providecommand{\rhfaPlsTwoAllPFull}{}\renewcommand{\rhfaPlsTwoAllPFull}{$<$0.001}  % pls2 all: p
\providecommand{\rhfaPlsTwoAllR}{}\renewcommand{\rhfaPlsTwoAllR}{$-$0.25}  % pls2 all: R_k
\providecommand{\rhfaPlsThreeAllPct}{}\renewcommand{\rhfaPlsThreeAllPct}{0.01}  % pls3 all: QLIKE change vs OLS-HAR, %
\providecommand{\rhfaPlsThreeAllMsePct}{}\renewcommand{\rhfaPlsThreeAllMsePct}{90.89}  % pls3 all: MSE change vs OLS-HAR, %
\providecommand{\rhfaPlsThreeAllDm}{}\renewcommand{\rhfaPlsThreeAllDm}{0.01}  % pls3 all: DM vs OLS-HAR (QLIKE)
\providecommand{\rhfaPlsThreeAllP}{}\renewcommand{\rhfaPlsThreeAllP}{0.989}  % pls3 all: p
\providecommand{\rhfaPlsThreeAllDmFull}{}\renewcommand{\rhfaPlsThreeAllDmFull}{6.39}  % pls3 all: DM vs the full set (enet)
\providecommand{\rhfaPlsThreeAllPFull}{}\renewcommand{\rhfaPlsThreeAllPFull}{$<$0.001}  % pls3 all: p
\providecommand{\rhfaPlsThreeAllR}{}\renewcommand{\rhfaPlsThreeAllR}{0.00}  % pls3 all: R_k
\providecommand{\rhfaFullAllModel}{}\renewcommand{\rhfaFullAllModel}{elastic net}  % all: L_full model (lower QLIKE of ridge, elastic net)
\providecommand{\rhfaFullAllPct}{}\renewcommand{\rhfaFullAllPct}{$-$6.69}  % all: full set QLIKE change vs OLS-HAR, %
\providecommand{\rhfaFullAllMsePct}{}\renewcommand{\rhfaFullAllMsePct}{117.67}  % all: full set MSE change vs OLS-HAR, %
\providecommand{\rhfaFullAllDm}{}\renewcommand{\rhfaFullAllDm}{$-$7.36}  % all: full set DM vs OLS-HAR
\providecommand{\rhfaFullAllP}{}\renewcommand{\rhfaFullAllP}{$<$0.001}  % all: p
\providecommand{\rhfaNAllDays}{}\renewcommand{\rhfaNAllDays}{2{,}579}  % all: days
\providecommand{\rhfaNAllBars}{}\renewcommand{\rhfaNAllBars}{33{,}527}  % all: bars
\providecommand{\rhfaPcaOneEvPct}{}\renewcommand{\rhfaPcaOneEvPct}{$-$0.82}  % pca1 event: QLIKE change vs OLS-HAR, %
\providecommand{\rhfaPcaOneEvMsePct}{}\renewcommand{\rhfaPcaOneEvMsePct}{$-$0.06}  % pca1 event: MSE change vs OLS-HAR, %
\providecommand{\rhfaPcaOneEvDm}{}\renewcommand{\rhfaPcaOneEvDm}{$-$1.95}  % pca1 event: DM vs OLS-HAR (QLIKE)
\providecommand{\rhfaPcaOneEvP}{}\renewcommand{\rhfaPcaOneEvP}{0.051}  % pca1 event: p
\providecommand{\rhfaPcaOneEvDmFull}{}\renewcommand{\rhfaPcaOneEvDmFull}{6.60}  % pca1 event: DM vs the full set (enet)
\providecommand{\rhfaPcaOneEvPFull}{}\renewcommand{\rhfaPcaOneEvPFull}{$<$0.001}  % pca1 event: p
\providecommand{\rhfaPcaOneEvR}{}\renewcommand{\rhfaPcaOneEvR}{0.04}  % pca1 event: R_k
\providecommand{\rhfaPcaTwoEvPct}{}\renewcommand{\rhfaPcaTwoEvPct}{0.39}  % pca2 event: QLIKE change vs OLS-HAR, %
\providecommand{\rhfaPcaTwoEvMsePct}{}\renewcommand{\rhfaPcaTwoEvMsePct}{0.12}  % pca2 event: MSE change vs OLS-HAR, %
\providecommand{\rhfaPcaTwoEvDm}{}\renewcommand{\rhfaPcaTwoEvDm}{0.36}  % pca2 event: DM vs OLS-HAR (QLIKE)
\providecommand{\rhfaPcaTwoEvP}{}\renewcommand{\rhfaPcaTwoEvP}{0.715}  % pca2 event: p
\providecommand{\rhfaPcaTwoEvDmFull}{}\renewcommand{\rhfaPcaTwoEvDmFull}{6.20}  % pca2 event: DM vs the full set (enet)
\providecommand{\rhfaPcaTwoEvPFull}{}\renewcommand{\rhfaPcaTwoEvPFull}{$<$0.001}  % pca2 event: p
\providecommand{\rhfaPcaTwoEvR}{}\renewcommand{\rhfaPcaTwoEvR}{$-$0.02}  % pca2 event: R_k
\providecommand{\rhfaPcaThreeEvPct}{}\renewcommand{\rhfaPcaThreeEvPct}{2.06}  % pca3 event: QLIKE change vs OLS-HAR, %
\providecommand{\rhfaPcaThreeEvMsePct}{}\renewcommand{\rhfaPcaThreeEvMsePct}{0.11}  % pca3 event: MSE change vs OLS-HAR, %
\providecommand{\rhfaPcaThreeEvDm}{}\renewcommand{\rhfaPcaThreeEvDm}{2.48}  % pca3 event: DM vs OLS-HAR (QLIKE)
\providecommand{\rhfaPcaThreeEvP}{}\renewcommand{\rhfaPcaThreeEvP}{0.013}  % pca3 event: p
\providecommand{\rhfaPcaThreeEvDmFull}{}\renewcommand{\rhfaPcaThreeEvDmFull}{6.74}  % pca3 event: DM vs the full set (enet)
\providecommand{\rhfaPcaThreeEvPFull}{}\renewcommand{\rhfaPcaThreeEvPFull}{$<$0.001}  % pca3 event: p
\providecommand{\rhfaPcaThreeEvR}{}\renewcommand{\rhfaPcaThreeEvR}{$-$0.09}  % pca3 event: R_k
\providecommand{\rhfaPcaFiveEvPct}{}\renewcommand{\rhfaPcaFiveEvPct}{1.05}  % pca5 event: QLIKE change vs OLS-HAR, %
\providecommand{\rhfaPcaFiveEvMsePct}{}\renewcommand{\rhfaPcaFiveEvMsePct}{$-$0.86}  % pca5 event: MSE change vs OLS-HAR, %
\providecommand{\rhfaPcaFiveEvDm}{}\renewcommand{\rhfaPcaFiveEvDm}{0.75}  % pca5 event: DM vs OLS-HAR (QLIKE)
\providecommand{\rhfaPcaFiveEvP}{}\renewcommand{\rhfaPcaFiveEvP}{0.453}  % pca5 event: p
\providecommand{\rhfaPcaFiveEvDmFull}{}\renewcommand{\rhfaPcaFiveEvDmFull}{6.70}  % pca5 event: DM vs the full set (enet)
\providecommand{\rhfaPcaFiveEvPFull}{}\renewcommand{\rhfaPcaFiveEvPFull}{$<$0.001}  % pca5 event: p
\providecommand{\rhfaPcaFiveEvR}{}\renewcommand{\rhfaPcaFiveEvR}{$-$0.05}  % pca5 event: R_k
\providecommand{\rhfaPcaTenEvPct}{}\renewcommand{\rhfaPcaTenEvPct}{1.72}  % pca10 event: QLIKE change vs OLS-HAR, %
\providecommand{\rhfaPcaTenEvMsePct}{}\renewcommand{\rhfaPcaTenEvMsePct}{$-$0.07}  % pca10 event: MSE change vs OLS-HAR, %
\providecommand{\rhfaPcaTenEvDm}{}\renewcommand{\rhfaPcaTenEvDm}{1.03}  % pca10 event: DM vs OLS-HAR (QLIKE)
\providecommand{\rhfaPcaTenEvP}{}\renewcommand{\rhfaPcaTenEvP}{0.302}  % pca10 event: p
\providecommand{\rhfaPcaTenEvDmFull}{}\renewcommand{\rhfaPcaTenEvDmFull}{5.52}  % pca10 event: DM vs the full set (enet)
\providecommand{\rhfaPcaTenEvPFull}{}\renewcommand{\rhfaPcaTenEvPFull}{$<$0.001}  % pca10 event: p
\providecommand{\rhfaPcaTenEvR}{}\renewcommand{\rhfaPcaTenEvR}{$-$0.08}  % pca10 event: R_k
\providecommand{\rhfaPcaTwentyEvPct}{}\renewcommand{\rhfaPcaTwentyEvPct}{0.39}  % pca20 event: QLIKE change vs OLS-HAR, %
\providecommand{\rhfaPcaTwentyEvMsePct}{}\renewcommand{\rhfaPcaTwentyEvMsePct}{$-$0.48}  % pca20 event: MSE change vs OLS-HAR, %
\providecommand{\rhfaPcaTwentyEvDm}{}\renewcommand{\rhfaPcaTwentyEvDm}{0.30}  % pca20 event: DM vs OLS-HAR (QLIKE)
\providecommand{\rhfaPcaTwentyEvP}{}\renewcommand{\rhfaPcaTwentyEvP}{0.762}  % pca20 event: p
\providecommand{\rhfaPcaTwentyEvDmFull}{}\renewcommand{\rhfaPcaTwentyEvDmFull}{5.91}  % pca20 event: DM vs the full set (enet)
\providecommand{\rhfaPcaTwentyEvPFull}{}\renewcommand{\rhfaPcaTwentyEvPFull}{$<$0.001}  % pca20 event: p
\providecommand{\rhfaPcaTwentyEvR}{}\renewcommand{\rhfaPcaTwentyEvR}{$-$0.02}  % pca20 event: R_k
\providecommand{\rhfaPlsOneEvPct}{}\renewcommand{\rhfaPlsOneEvPct}{$-$1.66}  % pls1 event: QLIKE change vs OLS-HAR, %
\providecommand{\rhfaPlsOneEvMsePct}{}\renewcommand{\rhfaPlsOneEvMsePct}{$-$0.32}  % pls1 event: MSE change vs OLS-HAR, %
\providecommand{\rhfaPlsOneEvDm}{}\renewcommand{\rhfaPlsOneEvDm}{$-$2.59}  % pls1 event: DM vs OLS-HAR (QLIKE)
\providecommand{\rhfaPlsOneEvP}{}\renewcommand{\rhfaPlsOneEvP}{0.010}  % pls1 event: p
\providecommand{\rhfaPlsOneEvDmFull}{}\renewcommand{\rhfaPlsOneEvDmFull}{6.61}  % pls1 event: DM vs the full set (enet)
\providecommand{\rhfaPlsOneEvPFull}{}\renewcommand{\rhfaPlsOneEvPFull}{$<$0.001}  % pls1 event: p
\providecommand{\rhfaPlsOneEvR}{}\renewcommand{\rhfaPlsOneEvR}{0.07}  % pls1 event: R_k
\providecommand{\rhfaPlsTwoEvPct}{}\renewcommand{\rhfaPlsTwoEvPct}{2.22}  % pls2 event: QLIKE change vs OLS-HAR, %
\providecommand{\rhfaPlsTwoEvMsePct}{}\renewcommand{\rhfaPlsTwoEvMsePct}{$-$0.33}  % pls2 event: MSE change vs OLS-HAR, %
\providecommand{\rhfaPlsTwoEvDm}{}\renewcommand{\rhfaPlsTwoEvDm}{1.79}  % pls2 event: DM vs OLS-HAR (QLIKE)
\providecommand{\rhfaPlsTwoEvP}{}\renewcommand{\rhfaPlsTwoEvP}{0.074}  % pls2 event: p
\providecommand{\rhfaPlsTwoEvDmFull}{}\renewcommand{\rhfaPlsTwoEvDmFull}{6.63}  % pls2 event: DM vs the full set (enet)
\providecommand{\rhfaPlsTwoEvPFull}{}\renewcommand{\rhfaPlsTwoEvPFull}{$<$0.001}  % pls2 event: p
\providecommand{\rhfaPlsTwoEvR}{}\renewcommand{\rhfaPlsTwoEvR}{$-$0.10}  % pls2 event: R_k
\providecommand{\rhfaPlsThreeEvPct}{}\renewcommand{\rhfaPlsThreeEvPct}{$-$0.12}  % pls3 event: QLIKE change vs OLS-HAR, %
\providecommand{\rhfaPlsThreeEvMsePct}{}\renewcommand{\rhfaPlsThreeEvMsePct}{$-$1.07}  % pls3 event: MSE change vs OLS-HAR, %
\providecommand{\rhfaPlsThreeEvDm}{}\renewcommand{\rhfaPlsThreeEvDm}{$-$0.06}  % pls3 event: DM vs OLS-HAR (QLIKE)
\providecommand{\rhfaPlsThreeEvP}{}\renewcommand{\rhfaPlsThreeEvP}{0.949}  % pls3 event: p
\providecommand{\rhfaPlsThreeEvDmFull}{}\renewcommand{\rhfaPlsThreeEvDmFull}{5.29}  % pls3 event: DM vs the full set (enet)
\providecommand{\rhfaPlsThreeEvPFull}{}\renewcommand{\rhfaPlsThreeEvPFull}{$<$0.001}  % pls3 event: p
\providecommand{\rhfaPlsThreeEvR}{}\renewcommand{\rhfaPlsThreeEvR}{0.01}  % pls3 event: R_k
\providecommand{\rhfaFullEvModel}{}\renewcommand{\rhfaFullEvModel}{elastic net}  % event: L_full model (lower QLIKE of ridge, elastic net)
\providecommand{\rhfaFullEvPct}{}\renewcommand{\rhfaFullEvPct}{$-$22.61}  % event: full set QLIKE change vs OLS-HAR, %
\providecommand{\rhfaFullEvMsePct}{}\renewcommand{\rhfaFullEvMsePct}{$-$8.54}  % event: full set MSE change vs OLS-HAR, %
\providecommand{\rhfaFullEvDm}{}\renewcommand{\rhfaFullEvDm}{$-$6.79}  % event: full set DM vs OLS-HAR
\providecommand{\rhfaFullEvP}{}\renewcommand{\rhfaFullEvP}{$<$0.001}  % event: p
\providecommand{\rhfaNEvDays}{}\renewcommand{\rhfaNEvDays}{81}  % event: days
\providecommand{\rhfaNEvBars}{}\renewcommand{\rhfaNEvBars}{1{,}053}  % event: bars
\providecommand{\rhfaReading}{}\renewcommand{\rhfaReading}{dense beyond a few factors}  % the RULE's reading, ORDINARY days
\providecommand{\rhfaBestSmall}{}\renewcommand{\rhfaBestSmall}{PLS, $k = 1$}  % the k <= 3 model with the highest R_k on ORDINARY days
\providecommand{\rhfaBestSmallR}{}\renewcommand{\rhfaBestSmallR}{0.00}  % its R_k
\providecommand{\rhfaBestSmallPFull}{}\renewcommand{\rhfaBestSmallPFull}{$<$0.001}  % its DM p against the full set
\providecommand{\rhfaNSharedModels}{}\renewcommand{\rhfaNSharedModels}{0}  % k <= 3 models meeting the SHARED FACTOR condition
\providecommand{\rhfaShareThreshold}{}\renewcommand{\rhfaShareThreshold}{0.75}  % SHARED FACTOR: R_k threshold
\providecommand{\rhfaDenseThreshold}{}\renewcommand{\rhfaDenseThreshold}{0.5}  % DENSE: R_10 threshold
\providecommand{\rhfaNBlock}{}\renewcommand{\rhfaNBlock}{258}  % block columns (value averages, 43 series x 6 horizons)
\providecommand{\rhfaNSeries}{}\renewcommand{\rhfaNSeries}{43}  % non-FOMC exogenous series
\providecommand{\rhfaNActiveMin}{}\renewcommand{\rhfaNActiveMin}{198}  % fewest active block columns on a scored day
\providecommand{\rhfaNActiveMax}{}\renewcommand{\rhfaNActiveMax}{257}  % most active block columns on a scored day
\providecommand{\rhfaSharePcOne}{}\renewcommand{\rhfaSharePcOne}{26.3}  % mean share of block variance, component 1, %
\providecommand{\rhfaSharePcTwo}{}\renewcommand{\rhfaSharePcTwo}{9.7}  % mean share of block variance, component 2, %
\providecommand{\rhfaSharePcThree}{}\renewcommand{\rhfaSharePcThree}{7.6}  % mean share of block variance, component 3, %
\providecommand{\rhfaShareCumOne}{}\renewcommand{\rhfaShareCumOne}{26.3}  % mean share in components 1-1, %
\providecommand{\rhfaShareCumTwo}{}\renewcommand{\rhfaShareCumTwo}{36.0}  % mean share in components 1-2, %
\providecommand{\rhfaShareCumThree}{}\renewcommand{\rhfaShareCumThree}{43.6}  % mean share in components 1-3, %
\providecommand{\rhfaShareCumFive}{}\renewcommand{\rhfaShareCumFive}{52.6}  % mean share in components 1-5, %
\providecommand{\rhfaShareCumTen}{}\renewcommand{\rhfaShareCumTen}{66.2}  % mean share in components 1-10, %
\providecommand{\rhfaShareCumTwenty}{}\renewcommand{\rhfaShareCumTwenty}{78.8}  % mean share in components 1-20, %
\providecommand{\rhfaShareCumThreeMin}{}\renewcommand{\rhfaShareCumThreeMin}{36.6}  % components 1-3: lowest day, %
\providecommand{\rhfaShareCumThreeMax}{}\renewcommand{\rhfaShareCumThreeMax}{47.1}  % components 1-3: highest day, %
\providecommand{\rhfaLoadRm}{}\renewcommand{\rhfaLoadRm}{2.61}  % PC1: sum of |loadings|, Realized moments
\providecommand{\rhfaLoadSqRm}{}\renewcommand{\rhfaLoadSqRm}{22.3}  % PC1: share of squared loadings, Realized moments, %
\providecommand{\rhfaPlsLoadRm}{}\renewcommand{\rhfaPlsLoadRm}{3.04}  % PLS1: sum of |weights|, Realized moments
\providecommand{\rhfaPlsLoadSqRm}{}\renewcommand{\rhfaPlsLoadSqRm}{34.0}  % PLS1: share of squared weights, Realized moments, %
\providecommand{\rhfaNColsRm}{}\renewcommand{\rhfaNColsRm}{42.0}  % active block columns of Realized moments (mean over days)
\providecommand{\rhfaLoadLiq}{}\renewcommand{\rhfaLoadLiq}{3.84}  % PC1: sum of |loadings|, Liquidity
\providecommand{\rhfaLoadSqLiq}{}\renewcommand{\rhfaLoadSqLiq}{21.2}  % PC1: share of squared loadings, Liquidity, %
\providecommand{\rhfaPlsLoadLiq}{}\renewcommand{\rhfaPlsLoadLiq}{3.95}  % PLS1: sum of |weights|, Liquidity
\providecommand{\rhfaPlsLoadSqLiq}{}\renewcommand{\rhfaPlsLoadSqLiq}{26.8}  % PLS1: share of squared weights, Liquidity, %
\providecommand{\rhfaNColsLiq}{}\renewcommand{\rhfaNColsLiq}{84.0}  % active block columns of Liquidity (mean over days)
\providecommand{\rhfaLoadEw}{}\renewcommand{\rhfaLoadEw}{2.97}  % PC1: sum of |loadings|, Equal-weighted market
\providecommand{\rhfaLoadSqEw}{}\renewcommand{\rhfaLoadSqEw}{28.2}  % PC1: share of squared loadings, Equal-weighted market, %
\providecommand{\rhfaPlsLoadEw}{}\renewcommand{\rhfaPlsLoadEw}{2.08}  % PLS1: sum of |weights|, Equal-weighted market
\providecommand{\rhfaPlsLoadSqEw}{}\renewcommand{\rhfaPlsLoadSqEw}{17.7}  % PLS1: share of squared weights, Equal-weighted market, %
\providecommand{\rhfaNColsEw}{}\renewcommand{\rhfaNColsEw}{36.0}  % active block columns of Equal-weighted market (mean over days)
\providecommand{\rhfaLoadVw}{}\renewcommand{\rhfaLoadVw}{2.91}  % PC1: sum of |loadings|, Value-weighted market
\providecommand{\rhfaLoadSqVw}{}\renewcommand{\rhfaLoadSqVw}{27.2}  % PC1: share of squared loadings, Value-weighted market, %
\providecommand{\rhfaPlsLoadVw}{}\renewcommand{\rhfaPlsLoadVw}{2.31}  % PLS1: sum of |weights|, Value-weighted market
\providecommand{\rhfaPlsLoadSqVw}{}\renewcommand{\rhfaPlsLoadSqVw}{21.2}  % PLS1: share of squared weights, Value-weighted market, %
\providecommand{\rhfaNColsVw}{}\renewcommand{\rhfaNColsVw}{36.0}  % active block columns of Value-weighted market (mean over days)
\providecommand{\rhfaLoadSent}{}\renewcommand{\rhfaLoadSent}{0.02}  % PC1: sum of |loadings|, Sentiment
\providecommand{\rhfaLoadSqSent}{}\renewcommand{\rhfaLoadSqSent}{0.0}  % PC1: share of squared loadings, Sentiment, %
\providecommand{\rhfaPlsLoadSent}{}\renewcommand{\rhfaPlsLoadSent}{0.04}  % PLS1: sum of |weights|, Sentiment
\providecommand{\rhfaPlsLoadSqSent}{}\renewcommand{\rhfaPlsLoadSqSent}{0.1}  % PLS1: share of squared weights, Sentiment, %
\providecommand{\rhfaNColsSent}{}\renewcommand{\rhfaNColsSent}{4.0}  % active block columns of Sentiment (mean over days)
\providecommand{\rhfaLoadIv}{}\renewcommand{\rhfaLoadIv}{0.18}  % PC1: sum of |loadings|, Implied volatility
\providecommand{\rhfaLoadSqIv}{}\renewcommand{\rhfaLoadSqIv}{1.1}  % PC1: share of squared loadings, Implied volatility, %
\providecommand{\rhfaPlsLoadIv}{}\renewcommand{\rhfaPlsLoadIv}{0.09}  % PLS1: sum of |weights|, Implied volatility
\providecommand{\rhfaPlsLoadSqIv}{}\renewcommand{\rhfaPlsLoadSqIv}{0.2}  % PLS1: share of squared weights, Implied volatility, %
\providecommand{\rhfaNColsIv}{}\renewcommand{\rhfaNColsIv}{5.1}  % active block columns of Implied volatility (mean over days)
\providecommand{\rhfaLoadVd}{}\renewcommand{\rhfaLoadVd}{0.01}  % PC1: sum of |loadings|, Volatility demand
\providecommand{\rhfaLoadSqVd}{}\renewcommand{\rhfaLoadSqVd}{0.0}  % PC1: share of squared loadings, Volatility demand, %
\providecommand{\rhfaPlsLoadVd}{}\renewcommand{\rhfaPlsLoadVd}{0.01}  % PLS1: sum of |weights|, Volatility demand
\providecommand{\rhfaPlsLoadSqVd}{}\renewcommand{\rhfaPlsLoadSqVd}{0.0}  % PLS1: share of squared weights, Volatility demand, %
\providecommand{\rhfaNColsVd}{}\renewcommand{\rhfaNColsVd}{3.1}  % active block columns of Volatility demand (mean over days)
\providecommand{\rhfaLoadTopGroup}{}\renewcommand{\rhfaLoadTopGroup}{equal-weighted market}  % group with the largest PC1 share of squared loadings
\providecommand{\rhfaGOneMaxRel}{}\renewcommand{\rhfaGOneMaxRel}{$4 \times 10^{-14}$}  % G1: k = 0 vs Table 1's OLS-HAR, max relative difference
\providecommand{\rhfaGTwoMaxRel}{}\renewcommand{\rhfaGTwoMaxRel}{$4 \times 10^{-14}$}  % G2: Gram route vs direct lstsq, max relative difference
\providecommand{\rhfaGTwoNDays}{}\renewcommand{\rhfaGTwoNDays}{7}  % G2 days
\providecommand{\rhfaSdReading}{}\renewcommand{\rhfaSdReading}{dense beyond a few factors}  % sensitivity (s.d.-only entry rule): reading
\providecommand{\rhfaSdPcaOneOrdR}{}\renewcommand{\rhfaSdPcaOneOrdR}{$-$0.02}  % sensitivity: pca1 R_k, ORDINARY days
\providecommand{\rhfaSdPcaTwoOrdR}{}\renewcommand{\rhfaSdPcaTwoOrdR}{$-$0.03}  % sensitivity: pca2 R_k, ORDINARY days
\providecommand{\rhfaSdPcaThreeOrdR}{}\renewcommand{\rhfaSdPcaThreeOrdR}{$-$0.14}  % sensitivity: pca3 R_k, ORDINARY days
\providecommand{\rhfaSdPlsOneOrdR}{}\renewcommand{\rhfaSdPlsOneOrdR}{0.00}  % sensitivity: pls1 R_k, ORDINARY days
\providecommand{\rhfaSdPlsTwoOrdR}{}\renewcommand{\rhfaSdPlsTwoOrdR}{$-$0.40}  % sensitivity: pls2 R_k, ORDINARY days
\providecommand{\rhfaSdPlsThreeOrdR}{}\renewcommand{\rhfaSdPlsThreeOrdR}{$-$0.04}  % sensitivity: pls3 R_k, ORDINARY days
\providecommand{\rhfaSdPcaTenOrdR}{}\renewcommand{\rhfaSdPcaTenOrdR}{0.21}  % sensitivity: pca10 R_k, ORDINARY days
\fi
\begin{table}[htbp]
\centering
\footnotesize
\setlength{\tabcolsep}{3pt}
\begin{threeparttable}
\caption{Factor-augmented HAR: how much of the full set's gain a few components recover}
\label{tab:rh_factor}
\begin{tabular*}{\linewidth}{@{\extracolsep{\fill}}lrrrrrr@{}}
\toprule
 & \multicolumn{3}{c}{Ordinary days (\rhfaNOrdDays{} days)} & \multicolumn{3}{c}{Every bar (\rhfaNAllDays{} days)} \\
\cmidrule(lr){2-4}\cmidrule(l){5-7}
 & QLIKE \% (DM) & $R_k$ & DM vs full & QLIKE \% (DM) & $R_k$ & DM vs full \\
\midrule
\multicolumn{7}{@{}l}{\textit{Principal components (\%: $-$, they add accuracy; DM vs full: $+$, the full set is better)}} \\
$k = 1$ & 0.08 (0.90) & $-$0.02 & 4.89 & $-$0.05 ($-$0.53) & 0.01 & 7.44 \\
$k = 2$ & 0.06 (0.32) & $-$0.02 & 4.86 & 0.11 (0.55) & $-$0.02 & 7.48 \\
$k = 3$ & 0.45 (1.44) & $-$0.11 & 5.33 & 0.69 (2.24) & $-$0.10 & 7.96 \\
$k = 5$ & 0.93 (1.85) & $-$0.23 & 5.94 & 0.94 (2.05) & $-$0.14 & 8.48 \\
$k = 10$ & $-$1.06 ($-$2.19) & 0.27 & 4.45 & $-$0.66 ($-$1.43) & 0.10 & 7.31 \\
$k = 20$ & $-$1.09 ($-$1.53) & 0.27 & 3.66 & $-$0.87 ($-$1.28) & 0.13 & 6.40 \\
\addlinespace[3pt]
\multicolumn{7}{@{}l}{\textit{Partial least squares (\%: $-$, they add accuracy; DM vs full: $+$, the full set is better)}} \\
$k = 1$ & 0.01 (0.05) & 0.00 & 4.95 & $-$0.23 ($-$1.07) & 0.03 & 7.49 \\
$k = 2$ & 1.57 (1.29) & $-$0.40 & 4.51 & 1.66 (1.55) & $-$0.25 & 6.99 \\
$k = 3$ & 0.04 (0.03) & $-$0.01 & 3.84 & 0.01 (0.01) & 0.00 & 6.39 \\
\addlinespace[3pt]
Full set (elastic net) & $-$3.97 ($-$4.73) & 1.00 & -- & $-$6.69 ($-$7.36) & 1.00 & -- \\
\bottomrule
\end{tabular*}
\begin{tablenotes}[flushleft]
\footnotesize
\item \textit{Sign:} QLIKE \%: $100\,(L_k/L_{\mathrm{HAR}} - 1)$, the percentage change in mean QLIKE of the row's forecast against Table~\ref{tab:rh_oos}'s OLS on the HAR set, with the Diebold--Mariano statistic (HLN, daily mean losses, lag 8) in parentheses (same sign): negative, the loss falls when the components are added to the HAR set, i.e.\ they add accuracy; positive, they hurt; in the full-set row, negative: the full set is more accurate. $R_k = (L_{\mathrm{HAR}} - L_k)/(L_{\mathrm{HAR}} - L_{\mathrm{full}})$, $L_{\mathrm{full}}$ the lower QLIKE of Table~\ref{tab:rh_oos}'s full-set ridge and elastic net on the same days: 1, all of the full set's gain over the HAR set; 0, none; below 0, worse than the HAR set. DM vs full: the row against that full set; positive, the full set is more accurate.
\item OLS on the HAR benchmark set plus the first $k$ components of the \rhfaNBlock{} value averages of the \rhfaNSeries{} exogenous series outside the FOMC group, standardized in each estimation window; principal components of the window's correlation matrix or partial least squares components against the target, both recomputed once a trading day from the window ending at the previous bar and held for the day; the regression re-estimated at every bar on Table~\ref{tab:rh_oos}'s 39{,}000-row window. Ordinary days: the scored sessions that are not one of the \rhfaNEvDays{} scheduled FOMC statement days. The full sets are penalized and include availability averages; the component regressions are OLS on value averages only, so the comparison is of forecast accuracy, not a nesting.
\end{tablenotes}
\end{threeparttable}
\end{table}


% 2026-10-28 (author: Sections 6 and 7 are one section). Was: \subsection{The futures' realized moments: the ordinary days}
% 2026-10-28 (author: conventional titles). Was: \subsubsection{The futures' realized moments: the ordinary days}
\subsubsection{Realized moments}
\label{sec:rh_fi_ordinary}
\label{sec:rh_results_bench}
\begin{itemize}
\item \textbf{Not an augmented HAR} (Table~\ref{tab:rh_ablation}, Panels~B--D). The futures' realized moments,
  computed from the same one-minute returns that define the target, are the natural augmented HAR. Added to the HAR
  benchmark set they lower the every-bar loss in every class (\rhabOlsAddRmAllPct\% for the OLS), by about as much as
  the three classical extensions estimated together (\rhabOlsSharqlAllPct\%; Panel~C). Yet the full set is
  significantly more accurate than the HAR benchmark set plus the realized moments in \rhabNClassesBeatAddRm{} of the
  \rhabNClassesAddRm{} classes (Panel~D).
\item \textbf{They carry most of the ordinary-day gain} (Table~\ref{tab:rh_channels_abl}). Added alone to the HAR
  benchmark set, the futures' realized moments lower the ordinary-day loss by \rhchRidgeOrdAddRmPct,
  \rhchLassoOrdAddRmPct{} and \rhchEnetOrdAddRmPct\% in the ridge, lasso and elastic net, nearly all or more of the
  full set's ordinary-day gain, and by \rhchOrdAddRmTreeRng\% in the trees, \rhchOrdRmShareTreeRng\% of theirs. Removing
  them from the full set is the only removal that raises the ordinary-day loss significantly in every class; on every
  bar it raises the loss by \rhabLoRmAllPctMin--\rhabLoRmAllPctMax\% ($\DM={}$\rhabLoRmAllDmMin{} or larger). The
  reading is concentrated in all \rhchLinOrdNConcTest{} linear classes and all \rhchTreeOrdNConcTest{} trees.
\item \textbf{The other groups repeat their information.} The constituents' realized moments and the liquidity
  measures each lower the trees' ordinary-day loss on their own, by \rhchOrdAddOtherTreeRng\% (on every bar,
  \rhabXgbAddVwAllPct\% for the value-weighted market and XGBoost), and the groups' separate gains sum to
  \rhchOrdSumRatioTreeRng{} times the trees' joint gain, but removing any one of them from the full set changes the
  gradient-boosted trees' ordinary-day loss by at most \rhchOrdLoOtherTreeMaxAbs\%, none significantly. The
  redundancy reading holds in all \rhchLinOrdNRedTest{} linear classes and in XGBoost; it fails for LightGBM, which
  has a second significant removal, and cannot be applied to the random forest, whose other removals were not run
  (the trees' removals other than the FOMC calendar's are on the design before the FOMC fix;
  Appendix~\ref{app:rh_channels}). The groups are largely correlated proxies of the volatility information that the
  futures' own moments measure most directly.
  % 2026-10-28 (coordinator; claim check on the patched numbers): was "... sum to \rhchOrdSumRatioRng{} times the joint
  % gain, but removing any one of them from the full set costs little. The redundancy reading holds in
  % \rhchLinOrdNRedTest{} of the three linear classes and \rhchTreeOrdNRedTest{} of the three trees: the groups are
  % largely correlated proxies ..." (the tree count 1 of 3 read as a rejection for the forest, which was not tested).
\item \textbf{Averaging beats selection on ordinary days.} On ordinary days the lasso's full-set loss lies
  \rhchLassoRidgeOrdPct\% above the ridge's and \rhchLassoEnetOrdPct\% above the elastic net's
  ($\DM={}$\rhchLassoRidgeOrdDm{} and \rhchLassoEnetOrdDm), the reverse of the statement days: the realized moments
  enter as several correlated columns, which averaging uses better than selection.
\item \textbf{Not a common factor of all the series} (Table~\ref{tab:rh_factor}). We add to the HAR benchmark set the
  first $k$ principal components, or partial-least-squares components, of the value averages of the \rhfaNSeries{}
  non-FOMC exogenous series, standardized and re-estimated each day in the rolling window, and fit by OLS at every bar.
  With $k \le 3$ the factor-augmented HAR recovers none of the full set's ordinary-day gain (at best
  $R_k = \rhfaBestSmallR$, \rhfaBestSmall), although the first three components hold \rhfaShareCumThree\% of the
  series' variance, and ten components recover $R_{10} = \rhfaPcaTenOrdR$. The leading components largely repeat the
  HAR terms; the gain sits in particular series, the realized moments, rather than in the series' common variation.
  % The RULE's mechanical reading is "\rhfaReading" (R_10 < 0.5); the text gives the reason rather than the label.
\end{itemize}

% 2026-10-28 (author: Sections 6 and 7 are one section). Was: \subsection{What remains}
% 2026-10-28 (author: conventional titles). Was: \subsubsection{What remains}
\subsubsection{Other predictor groups}
\label{sec:rh_fi_rest}
\label{sec:rh_results_beyond}
\begin{itemize}
\item \textbf{Small, and only in the boosted trees} (Table~\ref{tab:rh_ablation}, Panel~D, last row;
  Appendix~\ref{app:rh_results_beyond}). Against HAR\,+\,RM\,+\,FOMC the full set changes the every-bar loss of the
  ridge, lasso and elastic net by \rhabRidgeFullVsAddRmFomcAllPct, \rhabLassoFullVsAddRmFomcAllPct{} and
  \rhabEnetFullVsAddRmFomcAllPct\%, none significantly ($p \geq{}$\rhabLinFullVsAddRmFomcAllPMin), and that of
  LightGBM and XGBoost by \rhabLgbmFullVsAddRmFomcAllPct{} and \rhabXgbFullVsAddRmFomcAllPct\%
  ($\DM={}$\rhabLgbmFullVsAddRmFomcAllDm{} and \rhabXgbFullVsAddRmFomcAllDm); the random forest's
  \rhabRfFullVsAddRmFomcAllPct\% is not significant ($p$-value \rhabRfFullVsAddRmFomcAllP). On ordinary days the trees'
  remainder is \rhabLgbmFullVsAddRmFomcOrdPct, \rhabXgbFullVsAddRmFomcOrdPct{} and \rhabRfFullVsAddRmFomcOrdPct\%.
\item \textbf{A few series, from several groups} (Table~\ref{tab:rh_beyond_ladder}). No single remaining group is
  indispensable: removing any group other than the FOMC calendar and the realized moments changes the every-bar loss
  by at most \rhabLoOtherAllMaxAbs\%. Yet a few series carry the boosted trees' remainder, and they come from several
  groups. LightGBM's top five are the value- and equal-weighted constituents' absolute returns, the value-weighted
  effective spread, the VIX and the value-weighted constituents' third moment; XGBoost ranks StockTwits sentiment
  first, then the same constituent series and the VIX. In follow-up tests outside the pre-specified design, scored
  on ordinary days, LightGBM and XGBoost refitted on HAR\,+\,RM\,+\,FOMC plus their own top five series recover
  shares of \rhbyLLgbmKFiveOrdR{} and \rhbyLXgbKFiveOrdR{} of the remaining gain (LightGBM's top series alone
  \rhbyLLgbmKOneOrdR), and ten principal components of the six remaining groups \rhbyFLgbmPcaTenOrdR{} and
  \rhbyFXgbPcaTenOrdR.
\item \textbf{Sparse, not dense.} In the terms of \citet{GiannoneLenzaPrimiceri2021}, the remainder is sparse in
  series and scattered across groups.
  \begin{itemize}
  \item A dense linear signal would show in the ridge, which keeps and shrinks every series; yet the ridge's full set
    changes its ordinary-day loss against HAR\,+\,RM\,+\,FOMC by \rhabRidgeFullVsAddRmFomcOrdPct\%
    ($p = \rhabRidgeFullVsAddRmFomcOrdP$). In the boosted trees five series recover more of the remainder than ten
    principal components of all of them, and under the reading fixed before the ladder was scored XGBoost's
    remainder is \rhbyLXgbReading{} and LightGBM's \rhbyLLgbmReading{} ($R = \rhbyLLgbmKTwentyOrdR$ with twenty
    series): neither is dense.
  \item No group is indispensable, yet the top series come from four or five groups: series in different groups
    appear to measure much the same thing (both models rank the equal- and the value-weighted constituents'
    absolute returns among their top five). Where averaging beats selection, it does so among the correlated
    columns of the futures' realized moments (Section~\ref{sec:rh_fi_ordinary}).
  \end{itemize}
% 2026-10-28 (author, alternative-data edit: the claim is about the 36 series beyond HAR + RM + FOMC). Was:
% \item \textbf{In a few series, not a common factor.} No single remaining group is indispensable: removing any group
%   other than the FOMC calendar and the realized moments changes the every-bar loss by at most
%   \rhabLoOtherAllMaxAbs\%. Yet a few series carry the trees' remainder. In follow-up tests outside the pre-specified
%   design, scored on ordinary days, LightGBM and XGBoost refitted on HAR\,+\,RM\,+\,FOMC plus their own top five series
%   recover shares of \rhbyLLgbmKFiveOrdR{} and \rhbyLXgbKFiveOrdR{} of the remaining gain; ten principal components of
%   the six remaining groups recover \rhbyFLgbmPcaTenOrdR{} and \rhbyFXgbPcaTenOrdR. The gain sits in specific series,
%   not in a factor common to them, and the linear models do not capture it.
\end{itemize}

% 2026-10-28 (author: Sections 6 and 7 are one section). Was: \subsection{StockTwits, net vega demand and the VIX}
% 2026-10-28 (author: conventional titles). Was: \subsubsection{StockTwits, net vega demand and the VIX}
\subsubsection{Sentiment, volatility demand and implied volatility}
\label{sec:rh_fi_alt}
\begin{itemize}
\item \textbf{Little incremental gain} (Table~\ref{tab:rh_channels_abl}). Removing the StockTwits, volatility-demand
  or implied-volatility groups from the full set changes the loss by at most \rhchAltLoMaxAbs\% on either day type,
  significantly in only a few cells, and added alone to the HAR benchmark set the volatility-demand series raise the
  linear models' loss on both day types.
  % 2026-10-28 (coordinator; claim check on the patched numbers): was "\textbf{No incremental gain.} Removing ... changes
  % the loss by at most \rhchAltLoMaxAbs\% on either day type, and added alone to the HAR benchmark set they raise the
  % linear models' loss" -- three removals are significant (volatility demand on statement days for the lasso and
  % LightGBM, StockTwits on ordinary days for LightGBM), and added alone the implied volatility lowers the linear
  % models' ordinary-day loss.
\item \textbf{StockTwits sentiment in XGBoost} (Table~\ref{tab:rh_beyond_ladder}). XGBoost ranks StockTwits
  sentiment first among the remaining series: added alone to HAR\,+\,RM\,+\,FOMC it changes XGBoost's ordinary-day
  loss by \rhbyLXgbKOneOrdPct\% ($\DM={}$\rhbyLXgbKOneOrdDm), a share of \rhbyLXgbKOneOrdR{} of the remainder. Its
  values are live only from \rhabLiveSentStart, so over most of the span the model sees its availability averages,
  when messages print, not what they say.
\item \textbf{On their live windows} (Table~\ref{tab:rh_ablation}, Panel~E). These series start after the data and
% 2026-10-28 (author, alternative-data edit: the claim is about the 36 series beyond HAR + RM + FOMC). Was:
% \item \textbf{On their live windows} (Table~\ref{tab:rh_ablation}, Panel~E).
  enter as zero until the rolling window holds \rhorWindowDays{} trading days of their history
  (Section~\ref{sec:rh_predictors}), so over most of the span Panels~A and~B measure them through their availability
  shares. On the bars on which their values are live the gains remain small:
  \begin{itemize}
  \item Adding the VIX changes the loss of XGBoost by \rhabXgbAddIvLiveIvPct\% from \rhabLiveIvStart{}
    ($\DM={}$\rhabXgbAddIvLiveIvDm).
  \item The StockTwits series change no class's loss significantly from \rhabLiveSentStart{} (the largest change is
    the elastic net's, \rhabEnetAddSentLiveSentPct\%).
  \item The volatility-demand series raise the loss of the penalized linear models between \rhabLiveVdStart{} and
    \rhabLiveVdEnd{} (\rhabRidgeAddVdLiveVdPct\% for ridge regression, $\DM={}$\rhabRidgeAddVdLiveVdDm).
  \end{itemize}
\item \textbf{A weak test of volatility demand.} The net vega bought by non-market-makers is live only from
  \rhabLiveVdStart{} to \rhabLiveVdEnd, when its feed ends: \rhabLiveVdDays{} trading days,
  \rhlvVdEvLo--\rhlvVdEvHi\% of the scored bars. The sample shows that it does not help there, not that it cannot
  help.
\item \textbf{Used more while live, without a gain} (Figure~\ref{fig:rh_fi_live}). Before their values are live the
  trees see mostly the availability averages, which record when a series prints, not what it prints. In the full
  years in which their values are live, the trees' yearly attribution shares rise: StockTwits from
  \rhfilStwPreMin--\rhfilStwPreMax\% (\rhfilStwPreYears) to \rhfilStwLiveMin--\rhfilStwLiveMax\%
  (\rhfilStwLiveYears), VIX3M from \rhfilVixThreePreMin--\rhfilVixThreePreMax\% to
  \rhfilVixThreeLiveMin--\rhfilVixThreeLiveMax\%, volatility demand from \rhfilVdPreMin--\rhfilVdPreMax\% to
  \rhfilVdLiveMin--\rhfilVdLiveMax\% (\rhfilVdLiveYears). The use buys little accuracy (Panel~E).
\end{itemize}

% 2026-10-28 (author: Sections 6 and 7 are one section). Was: \subsection{Use is not accuracy}
% 2026-10-28 (author: conventional titles). Was: \subsubsection{Use is not accuracy}
\subsubsection{Variable importance}
\label{sec:rh_fi_use}
\begin{itemize}
\item \textbf{Attribution ranks the parts the wrong way round.} The trees attribute \rhfipHarMin--\rhfipHarMax\% of
  their forecasts to the HAR terms and \rhfipMomMin--\rhfipMomMax\% to the futures' realized moments, and only
  \rhfipFomcBucketMin--\rhfipFomcBucketMax\% to the FOMC calendar (Appendix~\ref{app:rh_fi_rank}). Yet the calendar
  carries most of their accuracy gain: added alone to the HAR benchmark set it changes QLIKE by
  \rhabLgbmAddFomcAllPct, \rhabXgbAddFomcAllPct{} and \rhabRfAddFomcAllPct\%, against \rhabLgbmFullAllPct,
  \rhabXgbFullAllPct{} and \rhabRfFullAllPct\% for the full set, by moving the forecast strongly on few bars
  (Section~\ref{sec:rh_fi_event}). The realized moments are the converse: second in attribution, yet removing them
  raises the trees' QLIKE by only \rhabLgbmLoRmAllPct, \rhabXgbLoRmAllPct{} and \rhabRfLoRmAllPct\%, as other
  predictors carry much of the same information. What a series adds to accuracy is measured by the ablations, not by
  its share.
\item \textbf{In short.} Two groups carry the pooled gain: the FOMC calendar on the statement days, in one group and on
  few bars, and the futures' realized moments on ordinary days, whose information the constituents' moments and the
  liquidity measures largely repeat and a few common factors of all the series do not recover. Beyond them and the HAR
  terms, the other \rhbyFLNSeries{} series add nothing to the linear models and a small gain to the boosted trees,
  carried by a few series scattered across groups: sparse, not dense. Attribution is a poor guide to all of this: the series a model uses most are not the ones that
% 2026-10-28 (author, alternative-data edit: the claim is about the 36 series beyond HAR + RM + FOMC). Was:
% What remains beyond
%   them and the HAR terms is small, only the trees capture it, and a few series carry it; the alternative data add
%   little.
  make it more accurate.
\end{itemize}
% 2026-10-28 (author: v2 of Sections 6-7 approved preliminarily): the previous body of Section 7 (Variable importance). Questions and measures, What the models use, and How many series carry the gain (with Ridge against lasso and the window scaling) moved to Appendix rehaul/A_fi.tex; the late-starting series, Use is not accuracy and In short are rewritten in Sections 7.4 and 7.5. Was:
% \begin{itemize}
% \item \textbf{Questions and measures.} Which predictors do the pooled models use, is the exogenous signal dense and
%   weak or concentrated, and does a series that moves the forecasts also make them more accurate? We explain the
%   full-set models of Table~\ref{tab:rh_oos} at \rhfipNOrigins{} monthly origins, \rhfipFirstOrigin{} to
%   \rhfipLastOrigin: each origin's refit is held for its block of \rhfipOriginEvery{} trading days, and the blocks
%   cover the \rhfipNRows{} bars of the span, including the shortened futures sessions on NYSE holidays, which
%   Table~\ref{tab:rh_oos} does not score (it scores \rhorNIntB). A group's share is $s_A$, its share of a tree
%   ensemble's mean absolute TreeSHAP value \citep{lundberg2020local}, or $S_A$, its share of a linear model's absolute
%   coefficients on the window-standardized predictors, $|\hat\beta_{k,j}|\,\mathrm{sd}_{k,j}$. Standard groups: the
%   HAR components, the calendar dummies, the futures' \rhfipNMom{} realized moments, the VIX and the futures' volume
%   and trade count; non-standard: the other \rhfipNNonStd{} series. Models are compared by their own shares, across
%   model types on the exogenous total; a share measures how much a series moves the forecast, not how much it
%   improves its accuracy. No explained bar lies in a tuning sample: the penalties were chosen on the last pre-sample
%   fold, ending \rhfipTuneLast, the tree configurations on an earlier development sample
%   (Appendix~\ref{app:rh_hyper}). Construction, intervals and checks: Appendix~\ref{app:rh_fi_compute}.
% \item \textbf{What the models use} (Table~\ref{tab:rh_fip_shares}). In the trees the HAR components carry
%   \rhfipHarMin--\rhfipHarMax\% of the attribution and the futures' realized moments \rhfipMomMin--\rhfipMomMax\%; the
%   leading exogenous series are the futures' own \rhfipTreeTopOne{} and \rhfipTreeTopTwo. The non-standard series hold
%   \rhfipLgbmNs, \rhfipXgbNs{} and \rhfipRfNs\% (LightGBM, XGBoost, random forest), most of it in five series:
%   \rhfipNsOne{} (\rhfipNsOneMin--\rhfipNsOneMax\%), \rhfipNsTwo{} (\rhfipNsTwoMin--\rhfipNsTwoMax\%), \rhfipNsThree{}
%   (\rhfipNsThreeMin--\rhfipNsThreeMax\%), \rhfipNsFour{} (\rhfipNsFourMin--\rhfipNsFourMax\%) and \rhfipNsFive{}
%   (\rhfipNsFiveMin--\rhfipNsFiveMax\%). The trees lean on them less than the linear models do: of the exogenous total
%   they hold \rhfipLgbmNsEx, \rhfipXgbNsEx{} and \rhfipRfNsEx\% in the trees against \rhfipRidgeNsEx, \rhfipLassoNsEx{}
%   and \rhfipEnetNsEx\% in the ridge, lasso and elastic net, the bootstrap interval of the difference lying below zero
%   in \rhfipNDLess{} of \rhfipNDAll{} comparisons, and the linear models rank the FOMC calendar near the top (the ridge
%   puts \rhfipFomcUntil{} at rank \rhfipRidgeRankFomcUntil{} of \rhfipNExog, the trees at
%   \rhfipTreeRankFomcUntilMin--\rhfipTreeRankFomcUntilMax). The trees' exogenous attribution is concentrated: its
%   effective number of series, $N_{\mathrm{eff}} = 1/\sum_A s_A^2$, is \rhfipLgbmNeff, \rhfipXgbNeff{} and
%   \rhfipRfNeff, against \rhfipRidgeNeff{} for the ridge, and \rhfipLgbmNHalf, \rhfipXgbNHalf{} and \rhfipRfNHalf{}
%   series hold half of it. Rankings over time, conditional shares, response shapes and interactions:
%   Appendices~\ref{app:rh_fi_rank}--\ref{app:rh_fi_shapes}.
% % 2026-10-26 (integration, JFEc): added -- the feature-count ladder (results/ladder_2026-10-26: RULE.txt sections 3, 5
% % and 6; tables.txt, readings.csv, posthoc_vs_walk.csv) and the ridge--lasso comparison replace the window test (P1/P2,
% % now Appendix app:rh_fi_window) as the main-text evidence on the dense-and-weak question.
% \item \textbf{How many series carry the gain} (Figure~\ref{fig:rh_ladder}). Following
%   \citet{MaischZhangUlrichZimmer2026}, we refit each model on its top $k$ exogenous series only, $k = 0$ (the HAR
%   benchmark set), 1, 3, 5, 10, 20 and \rhfipNExog{} (the full set), at the same origins and hyperparameters, and score
%   every rung on the \rhorNIntB{} bars of Table~\ref{tab:rh_oos}. Each model ranks the series once, by its own measure
%   on a fit explained on \rhorPenYears, before the span, so the ranking has no look-ahead. Rung $k$ retains
%   $R_k = (L_0 - L_k)/(L_0 - L_{\rhfipNExog})$ of the full set's QLIKE gain; the reading rule was fixed before scoring.
%   \begin{itemize}
%   \item Trees: the top 10 series keep \rhldTreeRTenMin{}--\rhldTreeRTenMax{}\% of the gain, the top
%     5 only \rhldTreeRFiveMin{}--\rhldTreeRFiveMax{}\%; \rhfipFomcUntil{}, ranked
%     \rhldTreeRankFomcUntilMin{}--\rhldTreeRankFomcUntilMax{} by the trees, enters between the two
%     rungs. The other series add little (top 10 against the full set: DM $p = \rhldLgbmPTenFull{}$,
%     \rhldXgbPTenFull{} and \rhldRfPTenFull{}).
%   \item Linear models: with 5 or 10 series they keep \rhldLinRFiveMin{}--\rhldLinRFiveMax{}\% and
%     \rhldLinRTenMin{}--\rhldLinRTenMax{}\% of the gain, more than the full set (top 10 against the
%     full set: DM $p \le \rhldLinPTenFullMax{}$), and the lasso's single top series,
%     \rhldLassoTopOne{}, matches the full set. Part of the shortfall is the cost of holding a fit on
%     \rhfipNCols{} predictors for a month: re-estimated at every bar (Table~\ref{tab:rh_oos}), the full sets are as
%     accurate as the monthly top-10 rungs (post hoc; DM $p$ from \rhldLinPTenVsWalkMin{} to
%     \rhldLinPTenVsWalkMax{}).
%   \item Ridge against lasso: a dense and weak signal favors the ridge, which keeps and shrinks every series, over the
%     lasso, which selects \citep[see][]{GiannoneLenzaPrimiceri2021}. In Table~\ref{tab:rh_oos} the full set changes
%     QLIKE relative to the HAR benchmark by \rhorVsOlsPctLassoFull\% in the lasso and \rhorVsOlsPctRidgeFull\% in the
%     ridge (lasso against ridge: DM $= \rhldDmLassoRidge{}$, $p = \rhldPLassoRidge{}$).
%     % 2026-10-26 (coordinator, D13): percentages from the patched Table 1; were \rhabLassoFullAllPct, \rhabRidgeFullAllPct.
%     % PENDING the author's framing decision (event vs ordinary days): on non-FOMC days ridge and elastic net beat the
%     % lasso (exploratory rescoring, 26 Oct), so this bullet may be split by day type.
%   \end{itemize}
%   The exogenous signal is concentrated, not dense and weak, in both model classes. Across windows of \rhcxWOne{} to
%   \rhcxWLong{} trading days the best number of series rises with the window, as a dense and weak signal predicts,
%   only for \rhcxWhoKstarRisesBar{} (Appendix~\ref{app:rh_fi_scaling}); the window test of the earlier draft is in
%   Appendix~\ref{app:rh_fi_window}.
% % 2026-10-26 (integration, JFEc): the late-starting series, shortened from the bullet of 2026-10-26 (parked below);
% % the live-window ablation (Table tab:rh_ablation, Panel E) added in one sentence.
% \item \textbf{The late-starting series while live} (Figure~\ref{fig:rh_fi_live}). Under the late-start rule
%   (Section~\ref{sec:rh_predictors}) the trees see mostly the availability averages of StockTwits, VIX3M and
%   volatility demand, which record when a series prints, not what it prints. In the full years in which their values
%   are live, the yearly shares rise: StockTwits from \rhfilStwPreMin--\rhfilStwPreMax\% (\rhfilStwPreYears) to
%   \rhfilStwLiveMin--\rhfilStwLiveMax\% (\rhfilStwLiveYears), VIX3M from \rhfilVixThreePreMin--\rhfilVixThreePreMax\%
%   to \rhfilVixThreeLiveMin--\rhfilVixThreeLiveMax\%, volatility demand from \rhfilVdPreMin--\rhfilVdPreMax\% to
%   \rhfilVdLiveMin--\rhfilVdLiveMax\% (\rhfilVdLiveYears). The pooled shares thus understate their use, but the use
%   buys little accuracy: on its live window, adding StockTwits to the HAR benchmark changes the trees' QLIKE by
%   \rhabLgbmAddSentLiveSentPct, \rhabXgbAddSentLiveSentPct{} and \rhabRfAddSentLiveSentPct\%, none significantly
%   (Table~\ref{tab:rh_ablation}, Panel~E).
% % 2026-10-26 (integration, JFEc): added -- use against accuracy: the SHAP share of the FOMC calendar (Table
% % tab:rh_fip_shares, Panel A) against its ablations (Table tab:rh_ablation, Panels A and B; the forest's leave-out arm
% % for the FOMC calendar was not run). The gain-share ratio (add-FOMC change over full-set change, trees) is not in a
% % generated file yet: \rhPending{rhabTreeFomcGainShareMin/Max}.
% \item \textbf{Use is not accuracy.} The FOMC calendar holds \rhfipFomcBucketMin--\rhfipFomcBucketMax\% of the trees'
%   attribution (Table~\ref{tab:rh_fip_shares}) yet carries most of their accuracy gain (Table~\ref{tab:rh_ablation}):
%   added alone to the HAR benchmark it changes QLIKE by \rhabLgbmAddFomcAllPct, \rhabXgbAddFomcAllPct{} and
%   \rhabRfAddFomcAllPct\%, against \rhabLgbmFullAllPct, \rhabXgbFullAllPct{} and \rhabRfFullAllPct\% for the full set
%   (\rhabTreeFomcGainShareMin{}--\rhabTreeFomcGainShareMax{}\% of the gain), and removing it raises
%   LightGBM's and XGBoost's QLIKE by \rhabLgbmLoFomcAllPct{} and \rhabXgbLoFomcAllPct\%. It moves the forecast strongly
%   on few bars: on FOMC statement days \rhfipFomcUntil{} holds \rhfipLgbmUntilFomc\% of LightGBM's attribution, on
%   other days \rhfipLgbmUntilOther\% (Appendix~\ref{app:rh_fi_cond}). The futures' realized moments are the converse:
%   second in attribution, yet removing them raises the trees' QLIKE by only \rhabLgbmLoRmAllPct, \rhabXgbLoRmAllPct{}
%   and \rhabRfLoRmAllPct\%, as other predictors carry much of the same information. What a series adds to accuracy is
%   measured by the ablations (Section~\ref{sec:rh_results_ablation}), not by its share.
% \item \textbf{In short.} The trees rely on the HAR components and the futures' realized moments, and their exogenous
%   attribution is concentrated. The signal is not dense and weak: about 10 series carry the trees' gain and 5 to 10
%   the linear models', and the lasso is at least as accurate as the ridge. The alternative data are used more while
%   live, without a measurable gain, and the FOMC calendar carries most of the trees' gain on a small share of their
%   attribution.
% \end{itemize}
% ---------------------------------------------------------------------------------------------------------------------
% 2026-10-26 (integration, JFEc): parked -- the body of the Section 6.5 version (2026-10-23/26), replaced by the four-part section above.
% Bullets moved to the online appendix are not repeated here; a pointer line names where each went. Was:
% \begin{itemize}
% \item \textbf{Questions.} Which predictors carry the pooled intraday forecasts, and does that change over time? Is the
%   exogenous signal dense and weak, spread thinly over many series, or concentrated in a few? Do the series that matter
%   do so only on some days, at some bars of the day, or through interactions with other predictors? And how do the
%   nonlinear models, the tree ensembles, respond to the non-standard series, the ones outside the HAR benchmark, the
%   futures' own realized moments and trading activity, and the VIX?
% 2026-10-26 (integration, JFEc): moved -- "What is computed" (origins, TreeSHAP, linear coefficients, classes, intervals; condensed in the "Measures" bullet above) -> rehaul/A_fi.tex, app:rh_fi_compute.
% \item \textbf{Scope.}
%   \begin{itemize}
%   \item The trees and the linear models are measured differently (mean absolute SHAP value against $|\gamma|$), so
%     the comparison is of rankings and of shares of each model's own measure, not of magnitudes. With the penalty
%     factor on the HAR block, the linear models' HAR shares are not comparable with the trees'; comparisons across
%     model types use the exogenous total.
%   \item A SHAP share measures how much a series moves the forecast, not how much it improves its accuracy.
%   \item Between correlated predictors (the equal- and value-weighted versions of a series, a series' value and
%     availability averages) the credit follows the trees' split choices; summing over a group removes the split
%     within the group, not the split between groups.
% 2026-10-26 (integration, JFEc): moved -- the Scope sub-bullets on the origin refits against the walks (bit-for-bit gates) and on XGBoost's single-precision additivity -> rehaul/A_fi.tex, app:rh_fi_compute.
%   \item Not covered yet: the MinGRU network (parked), and the ablations (leave one category out, then one series out
%     within the categories that matter), which are not run.
%   \end{itemize}
% \item \textbf{Which series matter} (Table~\ref{tab:rh_fip_shares}, Figure~\ref{fig:rh_fip_rank}).
%   \begin{itemize}
%   \item In the trees the HAR components carry \rhfipHarMin--\rhfipHarMax\% of the total attribution and the
%     futures' realized moments \rhfipMomMin--\rhfipMomMax\%. The two leading exogenous series are the same in all
%     three, \rhfipTreeTopOne{} and \rhfipTreeTopTwo; \rhfipTreeTopOne{} alone holds \rhfipLgbmTopOneEx,
%     \rhfipXgbTopOneEx{} and \rhfipRfTopOneEx\% of the exogenous total.
%   \item The non-standard series hold \rhfipLgbmNs\% [\rhfipLgbmNsLo, \rhfipLgbmNsHi], \rhfipXgbNs\%
%     [\rhfipXgbNsLo, \rhfipXgbNsHi] and \rhfipRfNs\% [\rhfipRfNsLo, \rhfipRfNsHi] of the total (LightGBM, XGBoost,
%     random forest), \rhfipNsExMin--\rhfipNsExMax\% of the exogenous total. Their weight sits in a few series; the five
%     largest by the mean over the three trees are \rhfipNsOne{} (\rhfipNsOneMin--\rhfipNsOneMax\% of the total), \rhfipNsTwo{}
%     (\rhfipNsTwoMin--\rhfipNsTwoMax\%), \rhfipNsThree{} (\rhfipNsThreeMin--\rhfipNsThreeMax\%), \rhfipNsFour{}
%     (\rhfipNsFourMin--\rhfipNsFourMax\%) and \rhfipNsFive{} (\rhfipNsFiveMin--\rhfipNsFiveMax\%). The rest carry
%     little: the StockTwits series together \rhfipSentMin--\rhfipSentMax\%, the volatility-demand series together
%     \rhfipVdMin--\rhfipVdMax\%, VVIX \rhfipVvixMin--\rhfipVvixMax\%, VIX3M \rhfipVixThreeMMin--\rhfipVixThreeMMax\%.
%   % 2026-10-26 (user: "make this explicit"): added -- the late-starting series in the years their value averages
%   % are live (writeup/make_rehaul_fi_live.py: \rhfil..., Figure fig:rh_fi_live; yearly shares of the TreeSHAP study).
%   \item The pooled shares understate the late-starting series (Figure~\ref{fig:rh_fi_live}). Under the late-start
%     rule (Section~\ref{sec:rh_predictors}) their value averages stay at the neutral zero until a full window
%     postdates their first print, so for most of the span the trees can use only their availability averages, which
%     record when a series prints, not what it prints. Over the full calendar years in which the value averages are
%     live, the yearly shares rise: StockTwits from \rhfilStwPreMin--\rhfilStwPreMax\% of the total
%     (\rhfilStwPreYears) to \rhfilStwLiveMin--\rhfilStwLiveMax\% (\rhfilStwLiveYears), VIX3M from
%     \rhfilVixThreePreMin--\rhfilVixThreePreMax\% to \rhfilVixThreeLiveMin--\rhfilVixThreeLiveMax\%, and volatility
%     demand from \rhfilVdPreMin--\rhfilVdPreMax\% to \rhfilVdLiveMin--\rhfilVdLiveMax\% in \rhfilVdLiveYears, its only
%     full live year before its feed stops. The VIX's share does not rise (\rhfilVixPreMin--\rhfilVixPreMax\% in
%     \rhfilVixPreYears, \rhfilVixLiveMin--\rhfilVixLiveMax\% in \rhfilVixLiveYears), and VVIX is live only from
%     \rhfilVvixLiveFrom. Even in the live years the models have seen these values on a small part of their training
%     window: at the last origin, \rhfilLastOrigin, the \rhfilWinDays-day window holds live values on
%     \rhfilStwWinPct\% of its days for StockTwits, \rhfilVixThreeWinPct\% for VIX3M and \rhfilVdWinPct\% for
%     volatility demand (\rhfilVixWinPct\% for the VIX). The live-year shares are the fairer measure of how much the
%     trees use the alternative data; like every SHAP share here, they measure use, not a gain in accuracy.
%   \item The ridge ranks the series differently. Its largest exogenous series is \rhfipRidgeTopOne{}
%     (\rhfipRidgeTopOneEx\% of its exogenous total), and it puts the FOMC calendar near the top: \rhfipFomcUntil{} at
%     rank \rhfipRidgeRankFomcUntil{} and \rhfipFomcDay{} at rank \rhfipRidgeRankFomcDay{} of \rhfipNExog, where the
%     trees rank them \rhfipTreeRankFomcUntilMin--\rhfipTreeRankFomcUntilMax{} and
%     \rhfipTreeRankFomcDayMin--\rhfipTreeRankFomcDayMax. Spearman rank correlation of the \rhfipNExog{} shares with the
%     ridge's: \rhfipLgbmSpRidge, \rhfipXgbSpRidge, \rhfipRfSpRidge{} (the two top-\rhfipTopN{} lists share
%     \rhfipLgbmTopRidge, \rhfipXgbTopRidge{} and \rhfipRfTopRidge{} series); among the trees: \rhfipSpLgbmXgb{} (LightGBM, XGBoost),
%     \rhfipSpLgbmRf{} and \rhfipSpXgbRf{} (each against the forest).
%   \end{itemize}
% \item \textbf{Dense and weak, or concentrated?}
%   \begin{itemize}
% 2026-10-26 (integration, JFEc): moved -- the linear models across windows (the P1/P2 window test, Table tab:rh_fip_window), with the audit fix results-7 -> rehaul/A_fi.tex, app:rh_fi_window.
% 2026-10-26 (integration, JFEc): moved -- the tree ensembles across windows (Table tab:rh_fip_lc) and its source comment -> rehaul/A_fi.tex, app:rh_fi_window.
%   \item Concentration at the main window (Table~\ref{tab:rh_fip_shares}, Panel D; descriptive, computed from the
%     pooled shares). The effective number of exogenous series, $N_{\mathrm{eff}} = 1/\sum_A s_A^2$ over the exogenous
%     shares, is \rhfipLgbmNeff, \rhfipXgbNeff{} and \rhfipRfNeff{} for the trees against \rhfipRidgeNeff{} for the
%     ridge (lasso \rhfipLassoNeff, elastic net \rhfipEnetNeff); half of the trees' exogenous attribution sits in
%     \rhfipLgbmNHalf, \rhfipXgbNHalf{} and \rhfipRfNHalf{} series, half of the ridge's in \rhfipRidgeNHalf. The trees
%     lean on a handful of series; the ridge spreads its weight more widely, as an $L_2$ penalty tends to across
%     correlated predictors.
%   \end{itemize}
% \item \textbf{Trees against linear models on the non-standard series} (Table~\ref{tab:rh_fip_shares}, Panels B
%   and C). Of the exogenous total, the non-standard series hold \rhfipLgbmNsEx, \rhfipXgbNsEx{} and \rhfipRfNsEx\% in
%   the trees, against \rhfipRidgeNsEx\% in the ridge, \rhfipLassoNsEx\% in the lasso and \rhfipEnetNsEx\% in the
%   elastic net. The difference, tree minus ridge, is $D = \rhfipLgbmDRidge$ [\rhfipLgbmDRidgeLo, \rhfipLgbmDRidgeHi],
%   \rhfipXgbDRidge{} [\rhfipXgbDRidgeLo, \rhfipXgbDRidgeHi] and \rhfipRfDRidge{} [\rhfipRfDRidgeLo, \rhfipRfDRidgeHi]
%   percentage points, and the interval lies below zero in all \rhfipNDLess{} of \rhfipNDAll{} comparisons (three trees,
%   three linear models, three denominators: the exogenous total, the total, and the exogenous total with the futures'
%   volume and trade count counted as non-standard). The trees lean on the non-standard series less than the linear
%   models do, not more.
% 2026-10-26 (integration, JFEc): moved -- "Over time" (Figure fig:rh_fip_time) -> rehaul/A_fi.tex, app:rh_fi_rank.
% 2026-10-26 (integration, JFEc): moved -- "Only sometimes, or conditionally?" (Figure fig:rh_fip_cond), with the FOMC-day note of audit results-9 -> rehaul/A_fi.tex, app:rh_fi_cond.
% 2026-10-26 (integration, JFEc): moved -- "Shape of the response" (Table tab:rh_fip_nonstd, Figure fig:rh_fip_shapes) -> rehaul/A_fi.tex, app:rh_fi_shapes.
% 2026-10-26 (integration, JFEc): moved -- "Interactions" (Table tab:rh_fip_inter) -> rehaul/A_fi.tex, app:rh_fi_shapes.
% \item \textbf{In short.} In the trees, the HAR components and the futures' realized moments carry most of the
%   attribution.
%   % 2026-10-23 (coordinator, tree learning curve): parked -- the dense-and-weak sentence now separates the linear
%   % models from the trees (the trees' accuracy across windows, tab:rh_fip_lc). Was: "The exogenous signal is not dense
%   % and weak in the sense tested: a longer window does not bring more series into the linear models, and the trees
%   % concentrate their exogenous attribution on a few series."
%   For the linear models the exogenous signal is not dense and weak in the sense tested: a longer window does not
%   bring more series in. The trees concentrate their exogenous attribution on a few series, yet on every bar their
%   full predictor set's lead over the HAR benchmark widens with the window until \rhlcWMain{} trading days and then
%   stops (at the 16:00 bar the change has no consistent sign), as a dense and weak signal would predict. The
%   non-standard series
%   matter conditionally, on FOMC statement days, at the first bar and at the bars where the FOMC calendar peaks, and
%   in calm markets, and about half of their attribution runs through interactions, the largest of them with the HAR
%   components. The tree ensembles use them less than the linear models do, through monotone, one-sided responses.
% \end{itemize}
% ---------------------------------------------------------------------------------------------------------------------
    % 2026-10-17 (user, E19): parked -- replaced by the pooled intraday feature-importance proposal (weekly vs monthly refit). Was:
    % \editc{red}{E19}{Decision: one cadence, and label the sample on every rank.\\Problem: the model bullet refits every trading day, the notation refits every \rhfiRefitEvery{} trading days, and a rank on \rhfiFirst{}--\rhfiLast{} comes from the closing-interval fits rather than from Table~\ref{tab:rh_oos}.\\
    % The model bullet says lasso and ridge are re-estimated every trading day, and that $\lambda$ stays at the pre-sample values $\rhfiLassoAlpha$ and $\rhfiRidgeAlpha$. The notation bullet then counts $K=\rhfiNRefits$ re-estimations, spaced \rhfiRefitEvery{} trading days apart, each one fit on the \rhfiWindow{} days before its first forecast day $t_k$ and then held for the block of days up to the next re-estimation. A fit every day and a fit every \rhfiRefitEvery{} days produce different coefficients on the days in between. The days those fits are asked to forecast run from \rhfiFirst{} to \rhfiLast, $|\mathcal O|=\rhfiNOos$ days in total. Table~\ref{tab:rh_oos} is another sample: the intraday scheme on $T_{\mathrm{oos}}=\rhoosDays$ days, from \rhoosSpanStart{} to \rhoosSpanEnd. A rank in this section is a summary of the closing-interval fits on \rhfiFirst{}--\rhfiLast. It is not a row of that table.}
    % 2026-10-17 (user, new-notes 10): parked -- E19 resolved: monthly feature-importance cadence by default (the user did not pick between weekly and monthly). Was:
    % \editc{red}{E19}{Decision to consider (not final): compute feature importance for the pooled intraday model, comparing a weekly and a monthly refit cadence, with monthly as the default if the check passes.\\Problem: the pooled intraday model has no feature-importance analysis; every rank in this section comes from the closing-interval models.\\ Feature importance should also be computed for the pooled intraday model, which currently has none. Two refit cadences are candidates to compare: weekly and monthly. At each cadence, re-estimate the pooled model and rank the predictor groups on each refit's held-out block. The comparison checks whether the ranking needs a finer granularity than monthly. If the weekly and monthly rankings are stable against each other, monthly is enough and becomes the default, because it costs less compute; if they differ, keep weekly. Both cadences are proposals to consider, not yet a decision.}
    % 2026-10-23 (user, E19, new-notes 10b): parked -- resolved: the pooled intraday model's variable importance is
    % computed at the monthly cadence (\rhfipNOrigins{} origins, one every \rhfipOriginEvery{} trading days) for the
    % linear models and the three trees, above; new-notes 22 (no explained day in a tuning fold) holds too: the origins
    % start on the first evaluation day, after the last tuning fold. Was:
    % \editc{red}{E19}{Decision (user, 2026-10-17): compute feature importance for the pooled intraday model with a monthly refit cadence, the default; the user has not chosen between weekly and monthly.\\Problem: the pooled intraday model has no feature-importance analysis, and the parked closing-interval analysis refit every \rhfiRefitEvery{} trading days (writeup/make\_rehaul\_fi.py:446).\\ User decision 2026-10-17: switch to a monthly feature-importance refit; requires a run of feature importance for the pooled intraday model. No monthly cadence has been run, so this section reports none.}
\iffalse
\begin{itemize}
% 2026-10-15 (user): parked -- last-30-min scheme moved to the end of the economics section; one clause says which
% scheme these models use. Was: "\item \textbf{Models:} the closing-interval models of record on the full predictor
% set ($p = \rhfiP$)."
\item \textbf{Models:} the closing-interval models of record on the full predictor set ($p = \rhfiP$), trained on
  the last 30 minutes of each trading day (Section~\ref{sec:rh_econ_last30}), not on the intraday scheme.
  \begin{itemize}
  \item Lasso and ridge with the penalty strength $\lambda$ selected on a pre-sample: $\lambda = \rhfiLassoAlpha$ and
    $\rhfiRidgeAlpha$ (scikit-learn parametrization). Penalty weight $\omega_j$, the factor on predictor $j$'s penalty
    term: $\underline{\omega} = \rhfiBackbonePf$ on the HAR benchmark block, $1$ on the exogenous predictors.
    Re-estimated every trading day; the scheduled re-selection of $\lambda$ never replaced the pre-sample choice.
  \item Not the earlier study's single-penalty specification (one $\lambda$ for every predictor, re-selected every
    \rhfiOldBlock{} trading days): the paper does not use it.
  \item Tree ensembles at fixed hyperparameters, the best entries of the development-sample search (LightGBM and
    XGBoost: the top Optuna trial; random forest: trial \rhdmRfTrial); constant and duplicate predictors dropped in
    each estimation window.
    % 2026-10-15 (finish, phase 2): typed 255 -> \rhdmRfTrial (same value).
    % Code names: LightGBM lgbm_opt_00, XGBoost xgb_opt_00, random forest trial 255, window-masked
    % (FEATIMP_WINDOW_MASK=1 in experiments/rehaul_fi_compute.py).
  \end{itemize}
\item \textbf{Notation.}
  \begin{itemize}
  % 2026-10-15 (finish): parked -- "refit" -> "re-estimation" (guide terms: re-estimated daily, re-estimation
  % frequency); the index k is unchanged. Was: "Refit $k = 1, \dots, K$, ... Refit $k$ is estimated on ... to the
  % next refit ($\le \rhfiRefitEvery$)."
  \item Re-estimation $k = 1, \dots, K$, $K = \rhfiNRefits$, one every \rhfiRefitEvery{} trading days. Re-estimation
    $k$ uses the window $\mathcal W_{t_k}$, the \rhfiWindow{} trading days before its first forecast day $t_k$, and
    forecasts its out-of-sample block $\mathcal O_k$, the trading days from $t_k$ to the next re-estimation
    ($\le \rhfiRefitEvery$).
    % 2026-10-15 (user): parked -- "out-of-sample period" (every forecast is out of sample). Was: "The blocks make up
    % this subsection's out-of-sample period $\mathcal O = \bigcup_k \mathcal O_k$:"
    The blocks make up this section's forecast days $\mathcal O = \bigcup_k \mathcal O_k$:
    % 2026-10-18 (user, new-notes 22): parked -- FI out-of-sample start moved to 3 Jan 2020 so it doesn't overlap the lambda folds; numbers in this block are from the 2018-start run until the FI rerun. Was: "$|\mathcal O| = \rhfiNOos$ trading days, \rhfiFirst{} to \rhfiLast." (the line below is unchanged: the dates and the day count are the \rhfi macros, rewritten by make_rehaul_fi.py from summary.json; after the rerun it reads 3 January 2020 to 30 April 2024, 1,087 days)
    $|\mathcal O| = \rhfiNOos$ trading days, \rhfiFirst{} to \rhfiLast.
  \item Predictors $\bm x_t \in \mathbb{R}^p$; forecast $\hat y_t = \hat f_k(\bm x_t)$, $t \in \mathcal O_k$, on the
    scale of the target $y_t$, the diurnally adjusted realized volatility of the closing half hour.
  \item Predictor group $A \subseteq \{1, \dots, p\}$: all predictors built from one series (its averages at every
    horizon, its missing-data and activity indicators). \rhfiNUnits{} groups: \rhfiNExog{} exogenous series in
    \rhfiNBuckets{} categories (Section~\ref{sec:rh_predictors}), the HAR components, the calendar dummies.
  \item Linear models: $\hat{\bm\beta}_k \in \mathbb{R}^p$, the coefficient estimates behind the forecast for $t_k$;
    $\hat{\bm\beta}_{k,A}$, its subvector on group $A$.
  \item Trees: SHAP values $\phi_{t,j}$ \citep{lundberg2017unified} and SHAP interaction values $\phi_{t,jl}$
    \citep{lundberg2020local}; they add up to the forecast,
    \[
      \hat y_t = \phi_0 + \sum_{j=1}^{p} \phi_{t,j}, \qquad
      \phi_{t,j} = \sum_{l=1}^{p} \phi_{t,jl}, \qquad
      \phi_{t,A} = \sum_{j \in A} \phi_{t,j}.
    \]
    For $j \ne l$, $\phi_{t,jl} = \tfrac12\bigl(\phi_{t,l} \mid j \text{ present} - \phi_{t,l} \mid j \text{ absent}\bigr)$,
    and $\phi_{t,jj} = \phi_{t,j} - \sum_{l \ne j} \phi_{t,jl}$.
  \item Path-dependent TreeSHAP \citep{lundberg2020local}: the training observations per node (node covers) stand in
    for the background sample.
  \item Number triples below: LightGBM, XGBoost, random forest.
  \end{itemize}
% 2026-10-19 (user): "a table in the feature importance section that is standard in the literature and showcases the
% quantities from the notes" -- the generator's overview table (writeup/make_rehaul_fi.py, parked from
% 08b_fi_measures_floats.tex on 2026-10-15), placed inside this parked block: the standalone
% (writeup/make_fi_bullets_standalone.py) prints it, the paper does not (note E19).
\item \textbf{Summary table:} Table~\ref{tab:rh_fi_summary} ranks the predictor groups model by model, as in the
  variable-importance rankings of \citet{gu2020empirical}, each model by its own aggregate measure below: the lasso's
  selection count $K_A$, the ridge's coefficient shares $S_A$ and $S^{\mathrm{sd}}_A$, and the trees' mean SHAP rank
  $\bar r_A$. Under each column: Kendall's $W$ of the rankings across re-estimations, and the number of exogenous
  series with a predictor that beats the random probe.
% AUTO-GENERATED by writeup/make_rehaul_fi.py -- do not edit; re-run the script.

\begin{table}[tbp]
\centering
\small
\setlength{\tabcolsep}{4pt}
\begin{threeparttable}
\caption{Feature importance by the coefficient and SHAP measures: the leading units per model, all features.}
\label{tab:rh_fi_summary}
\begin{tabular}{@{}llp{8.4cm}rr@{}}
\toprule
Model & Measure & Leading units & $W$ & Above noise \\
\midrule
Lasso & $N_A$ & Calendar 147; HAR ladder 147; \texttt{sumret} 147 & 0.74 & 23 \\
Ridge & $S_A$ (\%) & Calendar 31.0; HAR ladder 22.8; \texttt{sumret} 3.6 & 0.85 & 38 \\
LightGBM & $\bar r_A$ & HAR ladder 1.0; \texttt{sumret} 2.9; \texttt{sumabsret} 4.8 & 0.57 & 27 \\
XGBoost & $\bar r_A$ & HAR ladder 1.0; \texttt{sumret} 3.0; \texttt{sumabsret} 3.9 & 0.60 & 35 \\
Random forest & $\bar r_A$ & HAR ladder 1.0; \texttt{sumabsret} 2.6; \texttt{sumret} 4.0 & 0.62 & 29 \\
\bottomrule
\end{tabular}
\begin{tablenotes}[flushleft]
\footnotesize
\item \textit{Notes:} Per-bar 16:00 models of record, all features ($p = 628$ columns in 49 units: 47 exogenous series, the HAR ladder and the calendar block), 147 refits, one every 10 sessions. Lasso and ridge: pre-tuned penalties $\alpha = {10^{-2}}$ and ${10^{3}}$, backbone penalised at 1/100. $N_A$: refits in which a column of unit $A$ has $\beta_j \ne 0$; $S_A$: share of $\sum_k \sum_{j} |\beta^{(k)}_j|$ on $A$; $\bar r_A$: mean over refits of $A$'s rank by mean $|\phi_{t,A}|$ on the refit's held-out rows (1 = largest). $W$: Kendall's coefficient of concordance of the per-refit rankings (by $\sum_{j \in A} |\beta_j|$ for the linear models). Noise bar, exogenous series clearing it: lasso, a column selected in a larger share of the refits than the continuous probe in its 49 probe fits (1); ridge, a column with $|\beta_j|$ above the continuous probe's in at least half of the probe fits; trees, a column with mean $|\phi_j|$ above the continuous probe's in at least half of the probe fits.
\end{tablenotes}
\end{threeparttable}
\end{table}

\item \textbf{Lasso, aggregate:} the selection count of a group and of a predictor, the number of the $K$
  % 2026-10-15 (user): parked -- "park all tables that I haven't mentioned in this chat". Was: re-estimations with a nonzero coefficient (Tables~\ref{tab:rh_fi_summary} and~\ref{tab:rh_fi_linear}),
  re-estimations with a nonzero coefficient,
  \[
    K_A = \sum_{k=1}^{K} \mathbf{1}\bigl\{\exists j \in A:\ \hat\beta_{k,j} \ne 0\bigr\}, \qquad
    K_j = \sum_{k=1}^{K} \mathbf{1}\bigl\{\hat\beta_{k,j} \ne 0\bigr\};
  \]
  $K_A / K$ is the selection frequency, the share of estimation windows with a nonzero coefficient.
  % 2026-10-15 (consolidation): parked -- N_A, N_j (N is the one-minute count and the numobs series' symbol).
  % N_A = \sum_{k=1}^{K} \mathbf{1}\{\exists j \in A: \hat\beta_{k,j} \ne 0\}, N_j = \sum_k \mathbf{1}\{\hat\beta_{k,j} \ne 0\}.
  \begin{itemize}
  \item Per re-estimation: \rhfiLassoNzMed{} nonzero coefficients (median; range \rhfiLassoNzMin--\rhfiLassoNzMax), of
    which \rhfiLassoExogNzMed{} on exogenous predictors.
  \item Predictors: \rhfiLassoColsEver{} of \rhfiP{} selected at least once, \rhfiLassoColsMajority{} in at least half
    of the re-estimations.
  \item Exogenous series: \rhfiLassoExogEver{} of \rhfiNExog{} selected at least once; \rhfiLassoExogMajority{} in at
    least half of the re-estimations; \rhfiLassoExogAlways{} in every re-estimation (\rhfiLassoAlwaysList).
  \item HAR benchmark block: selected in every re-estimation ($\omega_j = \rhfiBackbonePf$);
    \rhfiLassoBackboneShare\,\% of $\sum_j |\hat\beta_{k,j}|$ (median re-estimation).
  \item Random-probe benchmark: the continuous probe, an i.i.d.\ uniform noise predictor (defined under Known issues),
    is selected in \rhfiLassoProbeRate\,\% of its fits; \rhfiLassoColsAboveProbe{} predictors, in
    \rhfiLassoSeriesAboveProbe{} exogenous series, are selected more often.
  \end{itemize}
\item \textbf{Lasso, time-varying:} the magnitude and the selection indicator of each group at each re-estimation,
  \[
    \lVert\hat{\bm\beta}_{k,A}\rVert_1 = \sum_{j \in A} \bigl|\hat\beta_{k,j}\bigr|, \qquad
    \mathbf{1}\bigl\{\lVert\hat{\bm\beta}_{k,A}\rVert_1 > 0\bigr\}
  \]
  % 2026-10-15 (notation pass): fig:rh_fi_linear parked by the user (guide section 7); pointer moved to the table
  % with the same numbers. Was: (Figure~\ref{fig:rh_fi_linear}, left).
  % 2026-10-15 (user): parked -- "park all tables that I haven't mentioned in this chat". Was: (Table~\ref{tab:rh_fi_linear}, lasso columns). Rank stability across re-estimations: Kendall's coefficient of
  Rank stability across re-estimations: Kendall's coefficient of
  concordance $W$ of the rankings by $\lVert\hat{\bm\beta}_{k,A}\rVert_1$ at each re-estimation ($W = 1$: the same
  ranking at every re-estimation; formula below), $W = \rhfiLassoW$.
\item \textbf{Ridge, aggregate:} the group's share of the absolute coefficient mass,
  \[
    S_A = \frac{\sum_{k} \sum_{j \in A} |\hat\beta_{k,j}|}{\sum_{k} \sum_{j=1}^{p} |\hat\beta_{k,j}|} .
  \]
  \begin{itemize}
  \item Scale check: only the averages of series values are robust-scaled.
    \begin{itemize}
    \item Averages of series values (\rhfiScaleValueN{} predictors): standard deviation over the estimation window
      \rhfiScaleValueSd{} (median; 5th--95th percentiles: \rhfiScaleValueSdLo--\rhfiScaleValueSdHi), interquartile
      range \rhfiScaleValueIqr.
    \item Averaged missing-data and activity indicators (\rhfiScaleFlagN): standard deviation \rhfiScaleFlagSd;
      calendar dummies \rhfiScaleCalSd; HAR components \rhfiScaleHarSd.
    \item So the raw $|\hat\beta_{k,j}|$ are in different units. Check with the group's share of the absolute
      standardized coefficients,
      \[
        S^{\mathrm{sd}}_A = \frac{\sum_{k} \sum_{j \in A} |\hat\beta_{k,j}|\, \mathrm{sd}_{k,j}}
          {\sum_{k} \sum_{j=1}^{p} |\hat\beta_{k,j}|\, \mathrm{sd}_{k,j}},
      \]
      % 2026-10-15 (user): parked -- "park all tables that I haven't mentioned in this chat". Was: with $\mathrm{sd}_{k,j}$ the standard deviation of predictor $j$ over $\mathcal W_{t_k}$
      with $\mathrm{sd}_{k,j}$ the standard deviation of predictor $j$ over $\mathcal W_{t_k}$;
      % 2026-10-15 (user): parked -- "park all tables that I haven't mentioned in this chat". Was: (Tables~\ref{tab:rh_fi_summary} and~\ref{tab:rh_fi_linear}); Spearman rank correlation of $S^{\mathrm{sd}}_A$
      Spearman rank correlation of $S^{\mathrm{sd}}_A$
      with $S_A$ across groups \rhfiRidgeSpearmanSd.
      % 2026-10-15 (consolidation): parked -- the unnamed sum; S^{sd}_A is the generator's symbol (make_rehaul_fi.py).
      % Check with absolute standardized coefficients: $\sum_k \sum_{j\in A} |\hat\beta_{k,j}|\, \mathrm{sd}_{k,j}$, ...
    \item Leading groups on that scale: \rhfiRidgeSdTopOne{} \rhfiRidgeSdTopOneShare\,\%, \rhfiRidgeSdTopTwo{}
      \rhfiRidgeSdTopTwoShare\,\%, \rhfiRidgeSdTopThree{} \rhfiRidgeSdTopThreeShare\,\%.
    \end{itemize}
  \item Leading groups: \rhfiRidgeTopOne{} \rhfiRidgeTopOneShare\,\%, \rhfiRidgeTopTwo{} \rhfiRidgeTopTwoShare\,\%,
    \rhfiRidgeTopThree{} \rhfiRidgeTopThreeShare\,\%; HAR benchmark block \rhfiRidgeBackboneShare\,\% (median
    re-estimation).
  \item Ridge sets no coefficient to zero: all \rhfiRidgeAllSelected{} groups have $K_A = K$, so the selection frequency
    does not discriminate.
  \item Random-probe benchmark: \rhfiRidgeProbeBeat{} predictors, in \rhfiRidgeProbeBeatSeries{} exogenous series, have
    $|\hat\beta_{k,j}|$ above the continuous probe's (robust-scaled) in most probe fits.
  \end{itemize}
\item \textbf{Ridge, time-varying:} the share at re-estimation $k$,
  $S_{k,A} = \lVert\hat{\bm\beta}_{k,A}\rVert_1 / \sum_{A'} \lVert\hat{\bm\beta}_{k,A'}\rVert_1$, $A'$ over all groups
  % 2026-10-15 (notation pass): fig:rh_fi_linear parked (guide section 7). Was: (Figure~\ref{fig:rh_fi_linear}, right).
  % 2026-10-15 (user): parked -- "park all tables that I haven't mentioned in this chat". Was: (Table~\ref{tab:rh_fi_linear}, ridge columns, mean rank by $\lVert\hat{\bm\beta}_{k,A}\rVert_1$); Kendall's
  (mean rank by $\lVert\hat{\bm\beta}_{k,A}\rVert_1$); Kendall's
  $W = \rhfiRidgeW$.
\item \textbf{Trees, ranking at each re-estimation:} the group's mean absolute SHAP value on the out-of-sample block
  % 2026-10-15 (user): parked -- "park all tables that I haven't mentioned in this chat". Was: of re-estimation $k$, then its rank (Table~\ref{tab:rh_fi_trees}),
  of re-estimation $k$, then its rank,
  \[
    I_{k,A} = \frac{1}{|\mathcal O_k|} \sum_{t \in \mathcal O_k} |\phi_{t,A}|, \qquad
    r_{k,A} = 1 + \#\{A' : I_{k,A'} > I_{k,A}\} \quad \text{(ties averaged)}.
  \]
\item \textbf{Trees, global ranking:} aggregate the $K$ rankings over time,
  \[
    \bar r_A = \frac{1}{K} \sum_{k=1}^{K} r_{k,A}, \qquad
    \mathrm{BC}_A = \sum_{k=1}^{K} \max\bigl(0,\ 11 - r_{k,A}\bigr).
  \]
  \begin{itemize}
  \item $\mathrm{BC}_A$: truncated Borda count (a full Borda count is affine in $\bar r_A$).
  \item Kendall's coefficient of concordance \citep{KendallBabingtonSmith1939} over the $K$ rankings of the $n$
    groups, with rank sums
    $R_A = \sum_k r_{k,A}$ and their mean $\bar R$:
    \[
      W = \frac{12 \sum_{A} \bigl(R_A - \bar R\bigr)^2}
        {K^2 (n^3 - n) - K \sum_{k} \sum_{g \in \mathcal E_k} \bigl(|g|^3 - |g|\bigr)},
    \]
    $\mathcal E_k$ the sets of tied (equally ranked) groups at re-estimation $k$; $W = 1$: the same ranking at every
    re-estimation.
    % 2026-10-15 (consolidation): \mathcal G_k -> \mathcal E_k (\mathcal G is the interval grid of the Data section).
  \item LightGBM: \rhfiLgbmTopOne{} ($\bar r_A = \rhfiLgbmTopOneRank$), \rhfiLgbmTopTwo{} (\rhfiLgbmTopTwoRank),
    \rhfiLgbmTopThree{} (\rhfiLgbmTopThreeRank); Kendall's $W = \rhfiLgbmW$.
  \item XGBoost: \rhfiXgbTopOne{} (\rhfiXgbTopOneRank), \rhfiXgbTopTwo{} (\rhfiXgbTopTwoRank), \rhfiXgbTopThree{}
    (\rhfiXgbTopThreeRank); Kendall's $W = \rhfiXgbW$.
  \item Random forest: \rhfiRfTopOne{} (\rhfiRfTopOneRank), \rhfiRfTopTwo{} (\rhfiRfTopTwoRank), \rhfiRfTopThree{}
    (\rhfiRfTopThreeRank); Kendall's $W = \rhfiRfW$.
  \item Mean rank vs truncated Borda count, Spearman rank correlation over groups: \rhfiLgbmSpBorda, \rhfiXgbSpBorda,
    \rhfiRfSpBorda.
  % 2026-10-15 (user): parked -- "over the out-of-sample period" -> "over $\mathcal O$"
  \item Aggregate of the ranks vs rank of the aggregate, the mean absolute SHAP value over $\mathcal O$,
    $\bar I_A = |\mathcal O|^{-1} \sum_{t \in \mathcal O} |\phi_{t,A}|$: \rhfiLgbmSpGlobal, \rhfiXgbSpGlobal, \rhfiRfSpGlobal.
  \item Predictor level ($n = \rhfiP$): Kendall's $W = \rhfiLgbmWCols$, \rhfiXgbWCols, \rhfiRfWCols.
  \item Agreement of $\bar r_A$ across models (Spearman rank correlation): trees \rhfiSpTreesMin--\rhfiSpTreesMax;
    tree vs linear \rhfiSpTreeLinMin--\rhfiSpTreeLinMax; lasso vs ridge \rhfiSpLassoRidge.
  \end{itemize}
\item \textbf{Interaction effects, global:} the mean absolute SHAP interaction value of two groups over
  % 2026-10-15 (user): parked -- "park all tables that I haven't mentioned in this chat". Was: out-of-sample period (Table~\ref{tab:rh_fi_inter}),
  % 2026-10-15 (user): parked -- "the out-of-sample period" -> $\mathcal O$. Was: "over the out-of-sample period,"
  $\mathcal O$,
  \[
    \phi_{t,AB} = \sum_{j \in A} \sum_{l \in B} \bigl(\phi_{t,jl} + \phi_{t,lj}\bigr), \qquad
    \bar I_{AB} = \frac{1}{|\mathcal O|} \sum_{t \in \mathcal O} |\phi_{t,AB}|, \quad A \ne B;
  \]
  % 2026-10-15 (consolidation): parked -- \Phi_{t,AB} (Phi is the standard normal distribution function, Methods).
  % \Phi_{t,AB} = \sum_{j \in A} \sum_{l \in B} (\phi_{t,jl} + \phi_{t,lj}), \bar I_{AB} = |O|^{-1} \sum_t |\Phi_{t,AB}|.
  for two predictors, $|\phi_{t,jl} + \phi_{t,lj}|$, $j \ne l$.
  \begin{itemize}
  % 2026-10-16 (formula pass): parked -- phi_{t,jl} and phi_{t,jj} are defined with the SHAP decomposition above.
  % Was: phi_{t,jl} = 1/2 (phi_{t,l} | j present - phi_{t,l} | j absent), j != l;
  % phi_{t,jj} = phi_{t,j} - sum_{l != j} phi_{t,jl}.
  \item Top pair: LightGBM \rhfiLgbmInterTopA{} $\times$ \rhfiLgbmInterTopB{} (\rhfiLgbmInterTopVal);
    XGBoost \rhfiXgbInterTopA{} $\times$ \rhfiXgbInterTopB{} (\rhfiXgbInterTopVal); random forest
    \rhfiRfInterTopA{} $\times$ \rhfiRfInterTopB{} (\rhfiRfInterTopVal); $\times 10^{-3}$, in units of $y_t$.
  \item Scale: largest $\bar I_A$ \rhfiLgbmMainTopVal, \rhfiXgbMainTopVal, \rhfiRfMainTopVal{} ($\times 10^{-3}$); all
    series-pair $\bar I_{AB}$ summed, as \% of $\sum_A \bar I_A$: \rhfiLgbmInterRatio, \rhfiXgbInterRatio,
    \rhfiRfInterRatio.
  \item Within a group ($\sum_{j \ne l \in A} \phi_{t,jl}$): largest \rhfiLgbmWithinTop{} (\rhfiLgbmWithinTopVal),
    \rhfiXgbWithinTop{} (\rhfiXgbWithinTopVal), \rhfiRfWithinTop{} (\rhfiRfWithinTopVal); $\times 10^{-3}$.
  \item Noise floor: the continuous probe's largest $\bar I$ with any group, median over probe fits:
    \rhfiLgbmInterFloor, \rhfiXgbInterFloor, \rhfiRfInterFloor{} ($\times 10^{-3}$); series pairs above it:
    \rhfiLgbmInterAboveFloor, \rhfiXgbInterAboveFloor, \rhfiRfInterAboveFloor{} of \rhfiNPairs.
  \end{itemize}
\item \textbf{Interaction effects, cross-check:} the $H$-statistic of \citet{FriedmanPopescu2008} on each model's top
  \rhfiHTopSeries{} series pairs, from the partial dependence functions $\mathrm{PD}_S$, $S \in \{A, B, AB\}$,
  $AB \equiv A \cup B$:
  \[
    H^2_{AB} = \frac{\sum_{t} \bigl[\mathrm{PD}_{AB}(\bm x_t) - \mathrm{PD}_A(\bm x_t) - \mathrm{PD}_B(\bm x_t)\bigr]^2}
                    {\sum_{t} \mathrm{PD}_{AB}(\bm x_t)^2}, \qquad
    \mathrm{PD}_S(\bm x) = \frac{1}{|\mathcal H_k|} \sum_{t' \in \mathcal H_k} \hat f_k\bigl(\bm x_S, \bm x_{t',-S}\bigr),
  \]
  $t \in \mathcal H_k$; each $\mathrm{PD}_S$ centered over $\mathcal H_k$.
  \begin{itemize}
  \item Points: $\mathcal H_k$, the \rhfiHRows{} trading days from $t_k$ on, both evaluation points and background
    sample, at \rhfiHRefits{} re-estimations (every tenth); sums pooled over re-estimations.
  \item The denominator contains the HAR components' own partial dependence whenever they are in the pair. A second
    normalization, by the smaller single-group partial dependence:
    \[
      H^2_{\min,AB} = \frac{\sum_{t} \bigl[\mathrm{PD}_{AB}(\bm x_t) - \mathrm{PD}_A(\bm x_t) - \mathrm{PD}_B(\bm x_t)\bigr]^2}
                           {\min\bigl\{\sum_{t} \mathrm{PD}_A(\bm x_t)^2,\ \sum_{t} \mathrm{PD}_B(\bm x_t)^2\bigr\}} .
    \]
  \item Top pair, $H^2$ (\%): \rhfiLgbmHTop, \rhfiXgbHTop, \rhfiRfHTop; over the top \rhfiHTopSeries:
    \rhfiLgbmHMin--\rhfiLgbmHMax, \rhfiXgbHMin--\rhfiXgbHMax, \rhfiRfHMin--\rhfiRfHMax.
  \item Top pair, $H^2_{\min}$: \rhfiLgbmHMinTop, \rhfiXgbHMinTop, \rhfiRfHMinTop; over the top
    \rhfiHTopSeries: \rhfiLgbmHMinLo--\rhfiLgbmHMinHi, \rhfiXgbHMinLo--\rhfiXgbHMinHi, \rhfiRfHMinLo--\rhfiRfHMinHi.
  \item Top predictor pair, $H^2$ (\%) and $H^2_{\min}$: LightGBM \rhfiLgbmHColTop, \rhfiLgbmHMinColTop; XGBoost
    \rhfiXgbHColTop, \rhfiXgbHMinColTop; random forest \rhfiRfHColTop, \rhfiRfHMinColTop.
  \end{itemize}
\item \textbf{Interaction effects, time-varying:} $\phi_{t,AB}$ on the out-of-sample block of each re-estimation,
  ranked at each re-estimation and aggregated as the SHAP values are,
  \[
    I_{k,AB} = \frac{1}{|\mathcal O_k|} \sum_{t \in \mathcal O_k} |\phi_{t,AB}| \ \to\ r_{k,AB} \ \to\ \bar r_{AB},\ W .
  \]
  \begin{itemize}
  \item Kendall's $W$ over each model's top \rhfiTopWPairs{} pairs: \rhfiLgbmInterW, \rhfiXgbInterW, \rhfiRfInterW.
  \item Mean rank vs $\bar I_{AB}$ (Spearman rank correlation, all pairs): \rhfiLgbmInterSpGlobal,
    \rhfiXgbInterSpGlobal, \rhfiRfInterSpGlobal.
  \end{itemize}
\item \textbf{Known issues, reconciled:}
  \begin{itemize}
  \item Mean decrease in impurity (MDI) favors predictors with many distinct values
    \citep{StroblBoulesteixZeileisHothorn2007}: MDI is not used; SHAP values and coefficients only.
  \item Random probes, at every third re-estimation (\rhfiNProbe{} fits): three i.i.d.\ noise predictors, continuous
    ($\mathrm{U}(0,1)$), discrete uniform on \rhfiProbeLevels{} levels, and binary (Bernoulli, probability one half).
    % experiments/rehaul_fi_compute.py:219-228, 381-389 (rng.random; floor(rng.random * PROBE_LEVELS); rng.random < 0.5).
    \begin{itemize}
    \item Lasso: the probes are selected in \rhfiLassoProbeSel, \rhfiLassoProbeSelLevels, \rhfiLassoProbeSelBinary{}
      of \rhfiNProbe{} fits (robust-scaled, $\omega_j = 1$).
    \item Ridge: \rhfiRidgeProbeBeat{} predictors have $|\hat\beta_{k,j}|$ above the continuous probe's in most probe
      fits.
    \item Trees: the continuous probe's median rank among \rhfiNProbeCols{} predictors by mean $|\phi_{t,j}|$:
      \rhfiLgbmProbeRank,
      \rhfiXgbProbeRank, \rhfiRfProbeRank; exogenous series with a predictor above it in most probe fits:
      \rhfiLgbmProbeSeries, \rhfiXgbProbeSeries, \rhfiRfProbeSeries.
    \end{itemize}
  \item Dilution across correlated predictors: $\phi_{t,A}$ and $\lVert\hat{\bm\beta}_{k,A}\rVert_1$ sum over the
    group's predictors, so credit split among a series' averages at different horizons stays inside the group.
  \item The earlier study's single-penalty lasso ($\lambda$ re-selected): median \rhfiOldLassoNzMed{} nonzero
    coefficients per re-estimation; \rhfiOldLassoExogMajority{} exogenous series selected in at least half of the
    re-estimations;
    Spearman rank correlation of its $K_A$ with that of the lasso of record \rhfiOldLassoSpearman.
  % 2026-10-15 (finish): parked -- "refit" -> "re-estimation" makes "Re-estimation: ... the archived refits ... the
  % re-estimated ones" ambiguous; the bullet is about reproducing the archived models, so it says so. Was:
  % \item Re-estimation: LightGBM reproduces the archived refits exactly; XGBoost and the random forest differ across
  %   library versions (maximum absolute forecast difference \rhfiXgbPredGap{} and \rhfiRfPredGap, in units of $y_t$).
  %   SHAP values: the archived refits'; interaction values: the re-estimated ones' (Spearman rank correlation of
  %   $\bar r_A$, re-estimated vs archived: \rhfiLgbmSpLocal, \rhfiXgbSpLocal, \rhfiRfSpLocal).
  \item Reproduction: re-running the estimation reproduces the archived LightGBM models exactly; XGBoost and the
    random forest differ across library versions (maximum absolute forecast difference \rhfiXgbPredGap{} and
    \rhfiRfPredGap, in units of $y_t$). SHAP values: from the archived models; interaction values: from the reproduced
    ones (Spearman rank correlation of $\bar r_A$, reproduced vs archived: \rhfiLgbmSpLocal, \rhfiXgbSpLocal,
    \rhfiRfSpLocal).
  \end{itemize}
\item \textbf{Compute} (one laptop, \rhfiWorkers{} processes):
  \begin{itemize}
  \item Reproducing the tree models and computing the SHAP interaction values: \rhfiWallTreesMin{} minutes. Per
    re-estimation: random forest
    estimation \rhfiFitSecRf\,s, its SHAP and interaction values \rhfiShapSecRf\,s; LightGBM \rhfiShapSecLgbm\,s,
    XGBoost \rhfiShapSecXgb\,s.
  \item Linear coefficient paths and probes: \rhfiWallLinMin{} minutes; $H^2$: \rhfiWallHMin{} minutes.
  \end{itemize}
\end{itemize}
\fi

% ---------------------------------------------------------------------------------------------------------------------
% 2026-10-15 parked (the measures pass above replaces them; the old study's floats remain in
% 08_feature_importance_floats.tex): the first skeleton's bullets on permutation, MDI, SHAP, drop-column and the probe.
% \section{Feature importance}
% \begin{itemize}
% \item \textbf{Scope:} per-bar 16:00 models, $d = \text{all}$ ($p = 628$); 147 refits.
%   \begin{itemize}
%   \item Refit $k$: fit on window $W_k$ (2,000 sessions); score on its causal tail $T_k$ ($\le 10$ sessions).
%   \item Unit $j$: a column, a series (all its columns) or a cluster ($|\mathrm{corr}| \ge 0.8$, complete linkage).
%   \item Tail loss $L(\hat f; X) = |T_k|^{-1}\sum_{t \in T_k} \ell_t$, with
%     $\ell_t \in \{\QLIKE(RV_t, \hat\sigma^2_t),\ (y_t - \hat y_t)^2\}$.
%   \end{itemize}
% \item \textbf{Permutation} (trees and linear):
%   \[
%     PI_j = L(\hat f; X^{\pi_j}) - L(\hat f; X),
%   \]
%   averaged over 10 draws, the same draws for every model.
%   \begin{itemize}
%   \item P1: $x^{\pi_j}_{t,j} = x_{\pi(t),j}$, $\pi$ a random permutation of $T_k$.
%   \item P2: $x^{\pi_j}_{t,j} = x_{u_t,j}$, $u_t \sim \mathrm{Unif}(W_k)$ i.i.d.\ (past rows only).
%   \item P1 is near-blind to slow inputs: they barely move within $T_k$.
%   \item Group units: all columns take the same permuted rows.
%   \end{itemize}
% \item \textbf{MDI / gain} (trees):
%   \[
%     \mathrm{MDI}_j = \sum_{b}\ \sum_{v \in b:\ j(v) = j} \Delta I_v ,
%   \]
%   the impurity (gain) decrease $\Delta I_v$ summed over every split $v$ on $j$ in every tree $b$.
%   \begin{itemize}
%   \item Reported as a share, $\mathrm{MDI}_j / \sum_i \mathrm{MDI}_i$, per refit.
%   \item Split count: $S_j = \#\{v : j(v) = j\}$.
%   \end{itemize}
% \item \textbf{SHAP:}
%   \[
%     \hat y_t = \phi_0 + \sum_{j=1}^{p} \phi_{t,j}, \qquad
%     \phi_{t,j}^{\text{lin}} = \beta_j (x_{t,j} - \bar x_j),
%   \]
%   importance: mean of $|\phi_{t,j}|$ over $t \in T_k$.
%   \begin{itemize}
%   \item Trees: path-dependent TreeSHAP.
%   \item Linear: $\bar x_j$ the window mean, so $\phi_0 = \beta_0 + \sum_j \beta_j \bar x_j$.
%   \end{itemize}
% \item \textbf{Standardised weight} (linear): $|\beta_j|\,\mathrm{sd}_j$, with $\mathrm{sd}_j$ over $W_k$.
% \item \textbf{Drop-column} (linear):
%   \[
%     D_j = L(\hat f_{-j}; X_{-j}) - L(\hat f; X),
%   \]
%   $\hat f_{-j}$ re-solved on $W_k$ without unit $j$, penalty $\lambda$ held.
% \item \textbf{Noise probe:} every 3rd refit, a second fit with three noise columns:
%   $z^{c}_t \sim U(0,1)$, $z^{20}_t \in \{0, \dots, 19\}$, $z^{b}_t \sim \mathrm{Bernoulli}(1/2)$.
%   \begin{itemize}
%   \item Criterion: unit $j$ counts only if $I_j > I_{z^{c}}$ under the same measure $I$.
%   \end{itemize}
% \item \textbf{Known issues to resolve:}
%   \begin{itemize}
%   \item Cardinality bias: $\mathrm{MDI}_j$ grows with the number of distinct values of $x_j$.
%   \item Noise probes also receive credit, under TreeSHAP too.
%   \item Dilution: for a correlated cluster $C$, $PI_C \ne \sum_{j \in C} PI_j$.
%   \item Linear permutation is heavy-tailed: unbounded $\hat y^{\pi}_t \to 0$ sends $\QLIKE \to \infty$.
%   \end{itemize}
% \item \todo{Results to come: primary measure, inclusion bar, pooled models.}
% \end{itemize}

